% Publication candidate PV4, 23 October 2026: Opus 5.5/max prose bridges,
% with Astra/max opening, integration and precision clarifications.
% Exact-version acceptance and mathematical lineage remain in companion records.
% neural-construction.tex -- AC3R-v2: root Astra/max scope clarification of AC3R-v1; candidate.
% Standalone manuscript written by Claude (Opus 5.5), task A, attempt 1, and
% corrected by Claude (Opus 5.5), task AC, attempt 1 (correction of the frozen A
% bytes, SHA256 dfaa7ca1...138cd, for R-001..R-008, Q-001..Q-005, ROOT-A-001,
% ROOT-A-002, ROOT-A-Q01..Q04 and ROOT-A-FORMAT-001; see CHANGES.md),
% frontier-foundations-20261021, 22 October 2026. Private candidate; exact-version
% review status is recorded in the companion decision. The frozen A bytes and the
% earlier conditional assembly MS-C3-v1 (SHA256 745b1b0e...7eb7cd) are not edited.
% AC3 integration: GPT-6.1 Sol/max; actual SB5/SB6 derivations and SMH repair.
% Exact-version review/acceptance is external; see companion records.
% Local executable/build status: CHECKS.md; coordinator owns rendered-document QA.
\documentclass[11pt]{article}
\usepackage[margin=1in]{geometry}
\usepackage[T1]{fontenc}
\usepackage{lmodern}
\usepackage{amsmath,amssymb,amsthm}
\usepackage{booktabs,longtable,array}
\usepackage{enumitem}
\usepackage{microtype}
\usepackage[numbers,sort&compress]{natbib}
\usepackage[hidelinks]{hyperref}
\hypersetup{
  pdftitle={Neurons Through the Looking Glass: Constructing an Interacting Quantum Field in Two Dimensions},
  pdfauthor={Cody Mitchell},
  pdfsubject={Version 2: process-neuron construction, conditional quantitative bounds and renormalized local squares},
  pdfkeywords={process neurons, reflection positivity, Wick-quartic field, constructive quantum field theory}
}

\setlength{\parskip}{0.35em}
\setlength{\emergencystretch}{2em}
\allowdisplaybreaks

% ---------------------------------------------------------------------------
\theoremstyle{plain}
\newtheorem{theorem}{Theorem}[section]
\newtheorem{proposition}[theorem]{Proposition}
\newtheorem{lemma}[theorem]{Lemma}
\newtheorem{corollary}[theorem]{Corollary}
\theoremstyle{definition}
\newtheorem{definition}[theorem]{Definition}
\newtheorem{classical}{Classical input}
\renewcommand{\theclassical}{C\arabic{classical}}
\theoremstyle{remark}
\newtheorem{remark}[theorem]{Remark}

% ---------------------------------------------------------------------------
\newcommand{\R}{\mathbb{R}}
\newcommand{\C}{\mathbb{C}}
\newcommand{\N}{\mathbb{N}}
\newcommand{\Z}{\mathbb{Z}}
\newcommand{\Q}{\mathbb{Q}}
\newcommand{\E}{\mathbb{E}}
\newcommand{\PP}{\mathbb{P}}
\newcommand{\T}{\mathbb{T}}
\newcommand{\cB}{\mathcal{B}}
\newcommand{\cC}{\mathcal{C}}
\newcommand{\cD}{\mathcal{D}}
\newcommand{\cE}{\mathcal{E}}
\newcommand{\cF}{\mathcal{F}}
\newcommand{\cG}{\mathcal{G}}
\newcommand{\cH}{\mathcal{H}}
\newcommand{\cI}{\mathcal{I}}
\newcommand{\cL}{\mathcal{L}}
\newcommand{\cN}{\mathcal{N}}
\newcommand{\cO}{\mathcal{O}}
\newcommand{\cP}{\mathcal{P}}
\newcommand{\cS}{\mathcal{S}}
\newcommand{\cA}{\mathcal{A}}
\newcommand{\ind}{\mathbf{1}}
\newcommand{\Cend}{C^{\mathrm{end}}_{\mu}}
\newcommand{\hCend}{\hat C^{\mathrm{end}}_{\mu}}
\newcommand{\Cfree}{C_{\mathrm{free}}}
\newcommand{\Res}{\mathrm{Res}}
\newcommand{\Det}{\mathrm{Det}}
\newcommand{\dist}{\mathrm{dist}}
\newcommand{\Schw}{\mathrm{Schw}}
\newcommand{\Hloc}{H^{-2}_{\mathrm{loc}}(\R^2;\R)}
\newcommand{\Lfour}{L^{4}_{\mathrm{loc}}(\R^2;\R)}
\newcommand{\bprod}{\mathrm{prod}}
\newcommand{\bMon}{\mathrm{Mon}}
\newcommand{\Part}{\mathrm{Part}}
\newcommand{\Mcol}{M^{\mathrm{col}}}
\newcommand{\fld}{\mathrm{fld}}
\newcommand{\resB}{\mathrm{res}}
\newcommand{\trd}{\mathrm{tr}}
\newcommand{\Ball}{\mathrm{Ball}}
\newcommand{\clip}{\mathrm{clip}}
\newcommand{\conj}[1]{\overline{#1}}
\DeclareMathOperator{\supp}{supp}
\DeclareMathOperator{\Law}{Law}
\DeclareMathOperator{\trace}{tr}
\DeclareMathOperator{\Var}{Var}

\title{Neurons Through the Looking Glass:\\[0.3em]
Constructing an Interacting Quantum Field in Two Dimensions}
\author{Cody Mitchell\\[0.2em]
\small Fractal Research Group (independent researcher)\\
\small \href{https://orcid.org/0009-0003-1546-0552}{ORCID: 0009-0003-1546-0552}}
\date{\fbox{\parbox{0.92\textwidth}{\centering
\textbf{6 November 2026. Version 2.}\\[0.2em]
Construction, conditional quantitative bounds and renormalized local squares.\\[0.2em]
Dated revision guides retain their historical status.\\[0.2em]
Authorship and provenance: Appendix~\ref{app:provenance}.}}}

% Current integration: globally unique PDF anchors.
\hypersetup{hypertexnames=false}
\begin{document}
\maketitle

\begin{abstract}
We construct an interacting quantum field from random fields called
\emph{process neurons}. We combine these building blocks and change their
probability weights through a local quartic interaction. Two explicit
architectures, built from random sign changes or randomly colored cells, reach
the established massive, weakly coupled two-dimensional Wick-quartic field along
one deterministic parameter sequence. This sequence is fixed before any
measurement is chosen and works for both architectures.

The construction preserves \emph{reflection positivity}, the positivity
condition used to recover a quantum Hilbert space from Euclidean fields.
At every finite stage the normalized interacting law has this property across
the two central coordinate lines. We prove weak convergence of the field laws on
$H^{-2}_{\mathrm{loc}}(\R^2)$, convergence of every power of the local seminorms,
and convergence of all polynomial moments of any finite collection of smooth,
compactly supported measurements chosen afterwards. The joint compact-test
correlation distributions also converge at every order.

We establish the process-neuron, Gaussian-reference and fixed-volume Wick
estimates needed for this construction. Attributed classical results supply
the infinite-volume target and its reconstruction as a relativistic quantum
field in one space and one time dimension; we verify their hypotheses at the
points of use. Finite width counts whole random processes, each using
countably many random variables globally. The common limiting sequence is
existential, with no effective end-to-end parameter schedule or computational
advantage established. The result gives a precise route from these two neural
architectures to the specified interacting field.
\end{abstract}

% Current integration guide; source markers feed PROOF-MAP.json.
% MAP: guide-current | S:10-46 S:398-470 P:522-557 US:3-17 UP:19-27 UD:1-30 D:1-30
\clearpage
\section*{Current guide: every dyadic accuracy}
\label{uda:guide}
\addcontentsline{toc}{section}{Current guide: every dyadic accuracy}
A process neuron is a random pattern on the plane. The construction adds
independent patterns and changes their probability weights through a local
quartic interaction. Its limiting law is the original interacting field.
The new result makes a quantitative part of that construction explicit:
after choosing an allowed finite collection of measurements, every dyadic
accuracy $2^{-q}$, $q\ge3$, has terminating finite outer stages and a
sufficient count of the original process neurons. The separately accepted
UDC implementation produces and independently checks the symbolic prescription
when its configured finite capacities suffice.

% MAP: guide-domain | S:10-46 S:416-470 P:522-557 US:3-15
The field, target and two architectures are unchanged. The measurement
family consists of an ordered one to four unweighted translates of the
original product bump, with canonical dyadic centers in $[-8,8]^2$ and
denominator exponents at most sixteen. Repetitions, singular observation
laws, and overlapping, touching or separated supports retain the full
ambient dimension. A separate $1/5$ request remains available. The guarantee
controls the joint bounded-Lipschitz law of these chosen measurements.

% MAP: guide-mechanism | S:351-424 P:229-436 P:438-557
The proof follows the joint law of a partial field and the next replacement
direction under the actual interaction weight. Coercivity controls field
moments; density domination controls direction moments; a centered fourth-order
comparison controls adjacent normalized laws. The argument keeps every Bell
and observation-product term and then sums the replacements. Volume-dependent
Nelson and physical coupling estimates select the later stages. All five
arrows end at the same original pressure-identified field.

% MAP: guide-width | S:249-266 S:398-414 S:426-461 P:438-520 P:559-749
Write $a=2^{-p}$, $\epsilon=2^{-t}$ and $\eta=2^{-d}$ for the regulator,
stabilizer and per-arrow allocation. Appendix~\ref{uda:current} gives
the complete prescription, proofs and exact source qualifications. Its
$N_{\mathrm{new}}=2^{e_{\mathrm{new}}}$ is strictly below the older sufficient
$2^{2^{w_{\mathrm{old}}}}$ at the same stages for every allowed $q$.
The exponents depend on the selected volume, and the proof retains
$\lceil t/2\rceil$ at both parities. On the eleven older settings the
stages and sharper width agree exactly with SVA/SVC.

% MAP: guide-version | S:444-484 P:687-749 US:3-17 UP:19-27 D:1-30 UD:1-30
The earlier reading guides and appendices are preserved as dated results
with their named versions. Their eleven-setting domains and statements
about then-pending software describe those versions. The current extension
is UDA-STAGE-RELATIVE-v2, with the separately accepted
UDC-dyadic-accuracy-v1 finite-resource implementation. The mathematical
every-$q$ theorem and the software's success or resource refusal have
different scopes. Section~\ref{udc:contract} specifies the current request,
certificate, source roots, verifier and capacities. Acceptance of the two
source packets is an internal exact-version decision; this manuscript
integration has its own fresh review gate.

% MAP: guide-limits | S:10-21 S:70-193 S:463-484 P:767-776 US:13-17 UC:19-35
The finite-positive-stage strong-$L^4$ laws, every bounded-bilinear identity,
physical common space, classical/source and NR1 qualifications remain
literal premises. Strict AP6 supplies the numerical target route; the
original fixed $g_*$ remains unknown. Counts are astronomical sufficient
whole-plane process-neuron counts, with countably many internal variables.
Canonical least inner indices, a global finite-weight simulator, practical
or necessary costs, full-field rates, broader models and priority remain
separate. Public v1 is unchanged.
\clearpage
\clearpage
\section*{Current reading guide: 28 October 2026}
\label{sva:guide}
\addcontentsline{toc}{section}{Current reading guide: 28 October 2026}

A process neuron is a random pattern on the plane. Adding independent
patterns and changing their probability weights through a local quartic
interaction produces the laws studied here. Removing inner approximations
and enlarging the interaction box reaches the original interacting field.

This revision improves the sufficient width when a supported finite
collection of measurements and an accuracy are chosen first. The sharper
relative estimate now follows the accuracy-selected outer stages, giving
$N_{\rm new}=2^{16\lceil t/2\rceil+20p+4d+544}$, strictly below the
older sufficient $N_{\rm old}=2^{2^{w_{\rm old}}}$ for every supported
request. Here $a=2^{-p}$ is the regulator, $\varepsilon=2^{-t}$ the
stabilizer, and $\eta=2^{-d}$ the per-arrow allocation. Theorem~\ref{sva:width}
proves the enclosure, ceiling and odd-$t$ rounding argument. The four
later errors, actual field and target are unchanged. Widths remain
astronomical sufficient counts of whole processes, each with countably
many internal spatial variables.

The domain is precise: ordered tuples of one to four unweighted canonical
dyadic translates of the original product bump, centers in $[-8,8]^2$
with denominator exponent at most sixteen, and exactly the eleven
tolerances $1/5,2^{-3},\ldots,2^{-12}$. Repetitions, singular observation
covariances, overlapping, touching and separated supports all keep their
original ambient dimension. For the four distinct corners
$(-8,-8),(-8,8),(8,-8),(8,8)$ at tolerance $2^{-12}$, the sufficient
new width is $2^{15668170148}$; the same-stage old width is
$2^{2^{1958531694}}$. Their outer stages and error allocation coincide.
The odd-$t$ unshifted scalar example at $1/5$ has sufficient width
$2^{13265754396}$; Section~\ref{sva:examples} gives its separate stages.

A separately versioned symbolic backend implements this rule.
Proposition~\ref{sva:software} gives the accepted SVC contract: independent
arithmetic verification, source/code bindings, parser limits and explicit
refusal/resource behavior. Proof, software identity, finite checks and
internal model review have separate roles. Checksums establish bytes;
classical and random-element premises retain their qualifications.

The key premise is the complete strong-$L^4$ reference-law interface at
every selected stage: strong measurability, zero Bochner means, global
sign symmetry and matching for every bounded bilinear form. Point-pair
covariance alone does not supply that interface. The physical own-variance
Wick and full neural/Gaussian normalizers have separate roles. Strict
AP6, P1--P6/source/classical assumptions, NR1 sharper-author versus
larger-blind-core scopes and the unknown original fixed $g_*$ remain.

The qualitative construction in Theorem~\ref{vqal:seq} retains its
inclusive cap $0<u/m^2\le1/8000$ under its named premises and its
all-later-test quantifiers. The present finite-observation schedules use
the strict quantitative cap and do not compute that construction's
canonical least inner indices. Theorem~\ref{vqal:mom} retains local
$H^{-2}$ moment convergence; the separate SEQ-HS result retains
physical/target $H^{-s}$ moment finiteness for every real $s>1$.

The accepted 107-page baseline's mathematical blocks and bibliography are
preserved. Its 27 and 26 October guides and fixed-$1/5$/old-ACC results
are dated checkpoints; Appendix~\ref{sva:current} gives the current result.
The preserved SVA-era unlaunched-SVC status describes its mathematics-only
checkpoint; the later accepted SVC decision supersedes that implementation
status. Full-field rates, practical/necessary complexity, broader models
and priority remain separate. Public v1 is unchanged.

This integration candidate awaits fresh source, proof and visual review.
The reversible delta, block audit and local proof map accompany it;
the coordinator owns acceptance.

\clearpage
\section*{Historical revision guide: 27 October 2026}
\label{vqal:guide}
\addcontentsline{toc}{section}{Historical revision guide: 27 October 2026}

A process neuron is a random pattern on the plane. Adding independent
patterns gives a reference field; weighting their probabilities by a local
quartic interaction turns it into an interacting law. We remove the four
inner approximations, then enlarge the interaction box. Reflection
positivity survives at every finite stage across the two central coordinate
lines and permits the classical quantum interpretation of the limiting law.
The detailed construction and classical attribution remain in the original
paper below.

The accepted qualitative statement is Theorem~\ref{vqal:seq}. Under the
literal accepted original, source and classical premises in
Appendix~\ref{vqal:current}, it extends the original construction to
$0<u/m^2\le1/8000$, including equality. One unchanged deterministic
sequence serves both original architectures and every compact measurement,
polynomial, correlation test and permitted reflection list chosen afterwards.
The local-moment proof is given in Theorem~\ref{vqal:mom}. MOM formally
concerns the original local $H^{-2}$ seminorms; the separately derived
SEQ-HS statement gives all-real-$s>1$ physical/target moment finiteness,
without asserting moment convergence at other $s$.

The quantitative question fixes measurements and an accuracy first.
For a nonzero scalar bump at $m=1$, $u=10^{-10}$, $\ell=1000$, a relative
normalization argument reduces one sufficient width to $2^{549756142138}$.
For ordered tuples of one to four supported dyadic bump translates, the
separate vector result gives $2^{549756142144}$ at accuracy $1/5$.
Both are immense counts of whole random processes. Each process retains
countably many internal spatial variables. The exact scopes, outer stages,
error allocations and proofs are in
Proposition~\ref{vqal:widths} and Table~\ref{vqal:comparison}.

The bounded accuracy search of Theorem~\ref{vqal:acc} uses the older
full-normalizer vector backend. It supports exactly $1/5$ and
$2^{-3},\ldots,2^{-12}$ for the stated typed family. It selects a
sufficient tuple for each request. At this 27 October checkpoint, the
sharper fixed-$1/5$ vector result remained separate. The later selected-stage
combination is now given in Appendix~\ref{sva:current}. All quantitative target
uses retain the strict cap $0<u/m^2<1/8000$ and literal AP6/P1--P6
premises. The originally fixed $g_*$ remains numerically unknown.

Theorem~\ref{vqal:link} proves that both routes reach the identical
original Fr\"ohlich--Simon law and compares them through their same
physical box. It provides a sufficient physical-box label for one
represented observation when certified source/norm data are supplied.
It leaves the original four canonical least indices existential. The
finite-observation-to-field reverse inference fails in general, as shown
by an explicit smooth-field kernel example. The construction's global
neural Schwartz extension, practical or necessary complexity, wider
models/dimensions and historical priority remain separate questions.

The $91$-page accepted private baseline is preserved mathematically and
its complete bibliography is unchanged. The original $g_*$ theorem and
26 October revision guide remain below as the historical route; the
new appendix states precisely which later premises yield the inclusive
qualitative cap. Full classical/source, NR1 author-versus-blind-core,
reconstruction and earlier adverse-version qualifications remain visible.
All accepted input results are internally reviewed at exact conditional
scopes. This integrated source/PDF is a new private candidate for fresh
source, proof and rendered-editorial review; its own acceptance is pending.
The published v1 remains unchanged.

The following private revision guide is the retained 26 October guide.
Its references lead to the earlier quantitative appendix. The current
results begin in Appendix~\ref{vqal:current}; the
\href{run:PROOF-MAP.md}{public proof map} identifies the mathematical locations
in this edition. Original exact-version source bindings and the historical
integration checks are retained in the private research archive.

\clearpage
\section*{Historical private revision guide: 26 October 2026}
\label{qrs:guide}
\addcontentsline{toc}{section}{Historical private revision guide: 26 October 2026}

This private post-v1 candidate preserves the accepted construction's
mathematical statements, proofs and equations. The abstract and main theorem
continue to describe the original observable-independent sequence, with its
original $g_*$ assumption and noneffective selection. The new quantitative
and reconstruction material is collected in Appendix~\ref{qrs:supplement}.

The subsequent chain proves local shell-response decay, converts it into a
finite interaction-box error, and controls all five approximation errors in
one fixed finite-observation metric. For a nonzero smooth product bump at
$m=1$, $u=10^{-10}$ and $\ell=1000$, it specifies an exact sufficient tuple
for both process-neuron architectures with error at most $65/4096<1/5$.
An earlier sufficient tuple and its slightly smaller certified error sum
remain visible. The optimized tuple has smaller parameters at the same
requested tolerance; its width remains astronomical.

The explicit numerical coupling uses a separately proved pressure-target
route under its named classical and finite-interface premises. It supplies
no numerical lower bound for the originally fixed $g_*$. The original
universal sequence remains intact. Corrected REC2 applies attributed
Fr\"ohlich--Simon/Osterwalder--Schrader reconstruction to that
pressure-identified target, keeping the precise half-plane observable domain.

The historical quantitative assertions have exact-version application maps
in the private research archive. The public
\href{run:PROOF-MAP.md}{PROOF-MAP} and
\href{run:COMPANION-INDEX.md}{companion index} identify the mathematical
locations and the limits of this public export. Earlier in-house reviews
retain their exact scope. The 26 October candidate status is historical;
version-specific review is summarized in the accompanying REVIEW document.
The published v1 artifacts remain unchanged.

Finite process-neuron width counts whole random processes, each with internal
whole-plane randomness. The certificate is a sufficient mathematical bound,
with sampling cost and necessary complexity left separate. Classical target
construction, clustering and quantum reconstruction retain attribution.

\clearpage
\tableofcontents
\clearpage

% ===========================================================================
\section{Introduction}\label{sec:intro}

\subsection{The question}\label{sec:intro-question}

Can simple random building blocks be assembled into an interacting quantum
field? We study this question by starting with random patterns on a plane,
combining them, and changing the probabilities of the resulting patterns to
introduce an interaction. The goal is to prove that their probability laws
approach a specified quantum field as the approximations are removed.

The building blocks are \emph{process neurons}: whole-plane random sign fields.
In one architecture a neuron is formed from two independent random sign
processes, one for each coordinate. In the other, a random division of the
plane into rectangular cells receives random signs. A width-$N$ field is a
normalized sum of $N$ such neurons. Before the interaction is introduced the
neurons are independent; reweighting their joint probability law introduces
the interaction. The width counts these random processes. Each process uses
countably many random variables globally, so finite width does not mean a
finite list of parameters encoding the entire plane.

Our target is the massive, weakly coupled two-dimensional Wick-quartic field,
usually denoted $\varphi^4_2$. Here ``quartic'' means that the interaction
involves the fourth power of the field, with the subtractions needed to define
its continuum limit. We first work in a Euclidean plane. Classical
reconstruction then identifies the limiting model with a relativistic quantum
field in one space and one time dimension. The main construction result concerns this particular model and weak-coupling regime.

\subsection{Why the looking glass?}\label{sec:intro-looking-glass}

Reflection gives a way to compare measurements on opposite halves of the
plane. \emph{Reflection positivity} requires certain averages formed from a
measurement and its reflected counterpart to be nonnegative. In the
reconstruction of a quantum theory, this supplies the positivity needed for
the Hilbert-space inner product. The precise condition concerns the full
probability law, including its interactions.

We preserve this condition across the two central coordinate lines at every
finite stage of the construction. The limiting field has full
Euclidean symmetry and reflection positivity across every line, provided by
the classical theory. The looking glass in the title points to this
mathematical role of reflection.

\subsection{From random patterns to a field}\label{sec:intro-route}

There are several approximations to control. We increase the number of
neurons, remove a stabilizing term, enlarge an auxiliary region, and remove a
short-distance regulator. We then enlarge the physical region in which the
quartic interaction acts. These are ordered limiting steps, assembled into
one common sequence for both architectures.

The sequence is chosen once, before selecting measurements. For example,
one may later choose smooth weighted averages of the field over bounded
regions: their joint polynomial moments converge along that same sequence.
The field laws themselves also converge in a space of distributions, which
allows the limiting field to be rougher than an ordinary function. This is
an existence result for the full sequence; it does not give an effective
rule for how many neurons or how large a region achieves a chosen accuracy.

The quantitative supplement in Appendix~\ref{qrs:supplement} addresses a
specified finite observation under separately named premises. It supplies a
concrete sufficient example while retaining this universal sequence's
existential meaning.

The target and the classical tools have a substantial history. Fr\"ohlich
and Simon constructed its infinite-volume states and proved their regularity
\cite{FS1977}; Albeverio, H{\o}egh-Krohn and Zegarlinski proved a weak-coupling
uniqueness theorem that identifies them \cite{AHKZ1989}; Guerra, Rosen and
Simon developed the free-boundary lattice approximation and the correlation
inequalities used below \cite{GRS1976}; and Simon and Griffiths proved the
finite-system inequality we need \cite{SG1973}. Relativistic reconstruction
uses the Osterwalder--Schrader theorem \cite{OS1973,OS1975} in the measure-level
formulation of \cite[Thm.~1.0]{FS1977}. These results supply the classical
inputs to our architecture-specific construction. Section~\ref{sec:prior}
compares that construction with earlier neural-field work.

The next subsection states the result precisely. Readers seeking an overview
can then follow the three layers in Section~\ref{sec:intro-layers} and the
contribution summary in Section~\ref{sec:intro-contrib}; the definitions and
proofs develop the details.

\subsection{The result in brief}\label{sec:intro-result}

Fix $m>0$, a length $\ell>0$ and $u>0$ with $0<u/m^{2}<g_{*}$
(Definition~\ref{def:threshold}). For $b\in\{\bprod,\bMon\}$ let
$P^{b}_{N,\varepsilon,\Omega,a;\Lambda}$ be the normalized annealed neural law
of Definition~\ref{def:laws}, a Borel probability on the Polish space
$X=\Hloc$. Theorem~\ref{thm:main} states that there is one deterministic
sequence $(a_n,R_n,\varepsilon_n,N_n)_{n\ge1}$, common to both architectures and
depending only on $m,\ell,u$, a fixed localization system and fixed parameter
enumerations, such that with the axis-parallel physical squares
$\Lambda_n=(-n\ell,n\ell)^{2}$ and auxiliary squares $\Omega_n=[-R_n,R_n]^{2}$
the laws $P^{b}_{n}=P^{b}_{N_n,\varepsilon_n,\Omega_n,a_n;\Lambda_n}$ satisfy:
\begin{enumerate}[label=(\Alph*)]
\item weak convergence on $X$ to the $X$-valued version $\nu$ of the classical
  target, with uniformly bounded and convergent powers of every local seminorm;
\item convergence of all polynomial moments of every finite tuple of real
  compactly supported smooth tests chosen after the sequence;
\item existence and weak distributional convergence of the joint correlation
  distributions of every order on compactly supported joint tests;
\item identification of the limits with restrictions of the classical tempered
  Schwinger distributions;
\item reflection positivity of every finite stage across the two central
  coordinate lines, for bounded half-plane observables and for finite
  polynomial forms in compact joint tests, with convergence of these forms.
\end{enumerate}
Theorem~\ref{thm:complement} records, as attributed classical results with
project-derived verification steps, the Euclidean invariance, translation
mixing and ergodicity, all-line reflection positivity, exponential regularity,
tempered Schwinger distributions, non-Gaussianity and $1+1$-dimensional
Wightman reconstruction of the target.

\subsection{What is proved, what is attributed}\label{sec:intro-layers}

The paper has three layers.
\begin{enumerate}[label=(\roman*)]
\item \emph{Classical inputs} (Section~\ref{sec:classical}). Published theorems
  \ref{prem:C1}--\ref{prem:C5} are imported with attribution, in formulations
  whose hypotheses are matched to the application at the point of use, and
  \ref{prem:C6} lists standard tools. Their proofs are not audited here, and
  the original Osterwalder--Schrader and Streater--Wightman texts were not
  available (Appendix~\ref{app:provenance}).
\item \emph{Foundations proved here.} The process neurons and their
  finite-width central reflection positivity (Section~\ref{sec:endpoint},
  Appendix~\ref{app:E}); the Gaussian reference fields, including the
  finite-region mass tilt and the common Gaussian space with Borel
  distribution-valued versions (Section~\ref{sec:gauss}, Appendix~\ref{app:G});
  the fixed-volume Wick theory (Section~\ref{sec:wick}, Appendix~\ref{app:W});
  the field-space facts (Section~\ref{sec:space}); and the fixed-volume
  chain of normalized laws (Section~\ref{sec:fixed}).
\item \emph{Construction and target.} The reduced continuum lemmas
  (Section~\ref{sec:target}): the continuum limit with its Wick identification,
  the whole-field volume limit with all compact moments, and the target
  properties. The common sequence, the local moments and the joint
  correlations (Sections~\ref{sec:H1}--\ref{sec:K1}). The assembly
  (Section~\ref{sec:proof}).
\end{enumerate}
The main theorem assumes \ref{prem:C1}--\ref{prem:C6} only. An earlier
conditional version of this construction assumed, besides these, nine project
premises; Appendix~\ref{app:dispositions} records for each of them whether it is
now a theorem of this paper or no longer used, and where the replacing
argument is.

\subsection{Contribution and organization}\label{sec:intro-contrib}

Our contribution is an architecture-specific realization and synthesis. In each
of the two architectures the width-$N$ field
$F_N=N^{-1/2}\sum_{i=1}^{N}g_i$ is a normalized sum of $N$ independent copies of a
whole-plane process neuron with $|g|=1/2$, $\E g=0$ and the endpoint covariance
$e^{-\mu|z|_1}/4$, $|z|_1=|z_1|+|z_2|$, $\mu=1/a$. On these fields we place
normalized annealed Gibbs laws: the joint law of all neuron randomness is
reweighted by $e^{-V(F_N)}$ and normalized, where for a centered axis-parallel
physical square $\Lambda$ inside an auxiliary square $\Omega=[-R,R]^{2}$,
\[
V(F_N)=u\int_{\Lambda}\bigl(F_N^{4}-6c_aF_N^{2}+3c_a^{2}\bigr)\,dz
+\varepsilon\int_{\Omega\setminus\Lambda}F_N^{4}\,dz
-\frac{\kappa_a^{2}}{2}\int_{\Omega}F_N^{2}\,dz,\qquad \kappa_a^{2}=a^{-2}-m^{2}.
\]
The fixed-volume limits, the growing-volume target and the joint correlation
statements above are proved for both architectures along one sequence. The
infinite-volume target, its Osterwalder--Schrader properties and its
reconstruction are classical results used with attribution. The
reflection-positivity mechanism, the density-deformation mechanism and the
endpoint covariance kernel all have prior instances (Section~\ref{sec:prior}).
We make no historical-priority claim.

Section~\ref{sec:setting} fixes notation and the order of the limiting stages.
Section~\ref{sec:classical} states the classical inputs and fixes the admissible
pair. Section~\ref{sec:main} states the results. Sections~\ref{sec:space}--\ref{sec:wick}
state the foundations, with proofs in Appendices~\ref{app:E}--\ref{app:W}.
Sections~\ref{sec:fixed}--\ref{sec:K1} prove the construction and
Section~\ref{sec:proof} assembles it. Sections~\ref{sec:prior} and
\ref{sec:limits} discuss prior work and limitations.

% ===========================================================================
\section{Setting, notation and conventions}\label{sec:setting}

This section fixes the objects that appear in the main theorem. It introduces the model parameters and the three covariances, lists the order in which the approximations are removed (Section~\ref{sec:massshift}), and defines the field spaces in which convergence is measured, the two neuron architectures, the interaction potential and the normalized laws. A first reading can concentrate on Definitions~\ref{def:params}, \ref{def:neuralfield}, \ref{def:potentials} and~\ref{def:laws} together with the list of stages; the remaining conventions can be consulted where they are cited.

\subsection{Parameters, geometry and covariances}\label{sec:params}

Points of the plane are $z=(z_1,z_2)$; $|z|$ is the Euclidean norm,
$|z|_1=|z_1|+|z_2|$ and $|z|_\infty=\max(|z_1|,|z_2|)$. In the reconstruction
(Theorem~\ref{thm:complement}) Euclidean time is $z_2$, following
\cite[p.~495]{FS1977}. $|A|$ is Lebesgue measure, $\cL(\R^{2})$ the Lebesgue
$\sigma$-algebra and $\cB(\cdot)$ a Borel $\sigma$-algebra.
$\N=\{1,2,3,\dots\}$. The symbol $N$ alone always denotes a width.

\begin{definition}[Parameters and physical geometry]\label{def:params}
Fix a mass $m>0$, a length $\ell>0$ and a quartic coefficient $u>0$. A
\emph{regulator} is a real $a$ with $0<a<1/m$; put $\mu=1/a$,
$\kappa_a^{2}=a^{-2}-m^{2}>0$ and $r_a=1-a^{2}m^{2}\in(0,1)$; the value $a=0$
denotes the unregulated case. A \emph{stabilizer} is $\varepsilon\ge0$.
\begin{enumerate}[label=(\roman*)]
\item An \emph{auxiliary square} is $\Omega=[-R,R]^{2}$, $R>0$.
\item A \emph{physical square} inside $\Omega$ is a centered
  \textbf{axis-parallel} square $\Lambda=(-l,l)^{2}$ or $\Lambda=[-l,l]^{2}$ with
  $0<l\le R$. Then $\bar\Lambda=[-l,l]^{2}\subseteq\Omega$, and both $\Lambda$ and
  $\Omega$ are invariant under the two central coordinate reflections
  $\theta_1(z_1,z_2)=(-z_1,z_2)$ and $\theta_2(z_1,z_2)=(z_1,-z_2)$.
\item The physical squares of the main theorem are
  $\Lambda_n=(-n\ell,n\ell)^{2}$, $n\ge1$, with closures
  $\bar\Lambda_n=[-n\ell,n\ell]^{2}$.
\end{enumerate}
Open and closed physical squares differ by the Lebesgue-null boundary
$\partial\Lambda$; every integrated object below is unchanged when $\Lambda$ is
replaced by $\bar\Lambda$ (Lemma~\ref{lem:openclosed}).
\end{definition}

The axis-parallel requirement in (ii) is a hypothesis, not a convention:
centering alone does not give the two reflection invariances used for
finite-stage positivity (Remark~\ref{rem:rotated}). Rotated physical squares are
not part of this definition; the general positivity theorem of
Section~\ref{sec:endpoint} applies to every region and potential that are
invariant under both central reflections, whatever their shape.

The Fourier convention is $\hat f(p)=\int e^{-ip\cdot z}f(z)\,dz$, extended to
$L^{2}$ by Plancherel, with inversion $C(z)=(2\pi)^{-2}\int e^{ip\cdot z}\hat
C(p)\,dp$ and $\int|f|^{2}=(2\pi)^{-2}\int|\hat f|^{2}$; for $g\in L^{1}(\R^{2})$ we write
$\cI g(x)=(2\pi)^{-2}\int e^{ip\cdot x}g(p)\,dp$. With $p=(w,q)$,
\begin{align}
\hat\Cfree(p)&=\hat C_0(p)=(m^{2}+|p|^{2})^{-1}, &
\hat C_a(w,q)&=(m^{2}+w^{2}+q^{2}+a^{2}w^{2}q^{2})^{-1},\label{eq:covariances}\\
\Cend(z)&=\tfrac14 e^{-\mu|z|_{1}}, &
\hCend(w,q)&=\frac{\mu^{2}}{(\mu^{2}+w^{2})(\mu^{2}+q^{2})}.
\nonumber
\end{align}
$\Cfree=(-\Delta+m^{2})^{-1}$ is the free covariance, $C_a$ the regulated
physical covariance and $\Cend$ the endpoint covariance. The Wick variance $c_a$
is the variance $C_a(0)$ of the regulated Gaussian physical reference field, the
centered Gaussian field with covariance $C_a$ (Theorem~\ref{thm:G-point}):
\begin{equation}\label{eq:ca}
c_a=(2\pi)^{-2}\int_{\R^{2}}\hat C_a(p)\,dp
=(2\pi)^{-1}\int_{0}^{\infty}\bigl[(m^{2}+q^{2})(1+a^{2}q^{2})\bigr]^{-1/2}dq\in(0,\infty).
\end{equation}
It is strictly decreasing in $a$, satisfies $c_a>1/4$ for every regulator
(compare the integrand in \eqref{eq:ca} with its value at $a=1/m$, where the right-hand side is $1/4$; see Lemma~\ref{lem:W-kernels}(K4)), and tends to $+\infty$ as $a\downarrow0$ by monotone
convergence. The endpoint variance is $\Cend(0)=1/4$; the two variances are
never interchanged.

\subsection{The mass-shift identity and the order of stages}\label{sec:massshift}

The neurons produce the endpoint covariance $\Cend$, whereas the target is built on the massive free covariance $\Cfree$. The identity below shows how this gap is bridged for the Gaussian references: subtracting $\kappa_a^{2}$ from the reciprocal of the endpoint multiplier gives the reciprocal of the multiplier of $C_a$, and $1/\hat C_a$ differs from $1/\hat\Cfree$ only by the term $a^{2}w^{2}q^{2}$ in \eqref{eq:covariances}. The negative quadratic term $-T_{\Omega,a}$ of the potential (Definition~\ref{def:potentials}) implements this shift (stages (3) and (5) below). The numbered list records the order in which the approximations are removed and the result controlling each stage.

Since $1/\hCend=\mu^{2}+w^{2}+q^{2}+a^{2}w^{2}q^{2}$ and $a\mu=1$,
\begin{equation}\label{eq:massshift}
\frac{1}{\hCend}-\kappa_a^{2}=\frac{1}{\hat C_a},\qquad
\hat C_a=\frac{\hCend}{1-\kappa_a^{2}\hCend}=\sum_{k\ge0}\kappa_a^{2k}(\hCend)^{k+1}.
\end{equation}
This is an identity between Fourier multipliers of \emph{Gaussian reference
covariances}. It is used in the following precise order of stages; it is not a
formula for the covariance of a finite-width law or of a quartically tilted law.
\begin{enumerate}[label=(\arabic*)]
\item \emph{Finite width.} The neural law $P^{b}_{N,\varepsilon,\Omega,a;\Lambda}$
  reweights the bounded, non-Gaussian field $F_N$ by the full potential $V$,
  whose negative quadratic term acts on all of $\Omega$.
\item \emph{Width limit.} As $N\to\infty$ these laws converge to the full Gibbs
  tilt $\nu^{v}_{\varepsilon,\Omega,a;\Lambda}$ of the Gaussian endpoint
  reference $G_\mu$ (covariance $\Cend$) by the same potential
  (Proposition~\ref{prop:width}); we assert neither that this tilt is Gaussian
  nor that it is not.
\item \emph{Pure quadratic tilt on finite $\Omega$.} The tilt of $G_\mu$ by
  $e^{T_{\Omega,a}}$ alone is the centered Gaussian law $G'_{\Omega,a}$ on the
  whole plane with covariance $C'_{\Omega,a}$ (Theorem~\ref{thm:G-tilt}), and
  $\Cend\le C'_{\Omega,a}\le C_a$ pointwise. The full tilt is the residual
  quartic tilt of $G'_{\Omega,a}$ (Proposition~\ref{prop:masstilt}).
\item \emph{Stabilizer removal} $\varepsilon\downarrow0$ at fixed $a,\Omega,\Lambda$
  (Proposition~\ref{prop:stab}).
\item \emph{Auxiliary-volume removal} $\Omega\uparrow\R^{2}$ at fixed $a,\Lambda$:
  $C'_{\Omega,a}\to C_a$ locally uniformly (Theorem~\ref{thm:G-diff}); this is
  where the multiplier \eqref{eq:massshift} becomes the covariance of the
  Gaussian reference, and the residual laws converge to the regulated physical
  law $\nu_{a,\Lambda}$ (Proposition~\ref{prop:aux}).
\item \emph{Regulator removal} $a\downarrow0$ at fixed $\Lambda$
  (Proposition~\ref{prop:reg}); the anisotropy $a^{2}w^{2}q^{2}$ of the Gaussian
  reference disappears.
\item \emph{Physical volume} $\Lambda_n\uparrow\R^{2}$, through classical
  continuum input (Section~\ref{sec:target}).
\end{enumerate}
The interacting laws at every stage are non-Gaussian in general, and their
covariances are not given by \eqref{eq:massshift}.

\subsection{Field spaces and localization}\label{sec:fieldspace-def}

Two spaces are needed. At every finite stage the neural field is a bounded function, so the quartic potential is an ordinary integral of a function in $E$. The limiting law lives on distributions (Theorem~\ref{thm:V1}(iii)), so convergence is stated in the space $X$ of local $H^{-2}$ distributions. Its seminorm $p_j$ cuts the field off with $\chi_j$, regards the result as a periodic distribution on the torus $\T_j$, and measures it by a weighted Fourier sum; every compact smooth smearing $v(f)$ is controlled by one such seminorm through \eqref{eq:duality}.

\begin{definition}[Field spaces]\label{def:fieldspace}
$E=\Lfour$ consists of real locally fourth-power integrable functions modulo
null functions, with the metric
$d_E(v,w)=\sum_{j\ge1}2^{-j}\min(1,\|v-w\|_{L^{4}((-j,j)^{2})})$.
$X=\Hloc$ consists of the real distributions $v$ on $\R^{2}$ with
$\chi v\in H^{-2}(\R^{2})$ for every $\chi\in C_c^{\infty}(\R^{2})$. The map
$\iota:E\to X$ sends a function class $v$ to $f\mapsto\int vf$. The
\emph{localization system} consists of:
\begin{itemize}
\item real cutoffs $\chi_j\in C_c^{\infty}(\R^{2})$, $j\ge1$, with
  $0\le\chi_j\le1$, $\chi_j=1$ on $[-j,j]^{2}$,
  $\supp\chi_j\subset(-j-\tfrac12,j+\tfrac12)^{2}$ and $\chi_{j+1}=1$ on
  $\supp\chi_j$;
\item the square tori $\T_j$ of side $2(j+1)$, $|\T_j|=4(j+1)^{2}$, with
  fundamental cell $\cC_j=(-(j+1),j+1)^{2}$; a distribution supported in
  $\supp\chi_j$ is regarded as a periodic distribution on $\T_j$;
\item the orthonormal basis $e_{j,k}(z)=|\T_j|^{-1/2}\exp(i\pi k\cdot z/(j+1))$,
  $k\in\Z^{2}$, of $L^{2}(\T_j)$, and the weights
  $\lambda_{j,k}=1+\pi^{2}|k|^{2}/(j+1)^{2}$, so that
  $(1-\Delta)\conj{e_{j,k}}=\lambda_{j,k}\conj{e_{j,k}}$ and
  $\lambda_{j,-k}=\lambda_{j,k}$.
\end{itemize}
For $r\in\R$ and a periodic distribution $h$ on $\T_j$ put
$\|h\|_{H^{r}(\T_j)}^{2}=\sum_{k}\lambda_{j,k}^{r}|h(\conj{e_{j,k}})|^{2}\in[0,\infty]$.
The local seminorms, the metric and the bounded-Lipschitz distance are
\begin{equation}\label{eq:metric}
\begin{gathered}
p_j(v)=\|\chi_j v\|_{H^{-2}(\T_j)},\qquad
d(v,w)=\sum_{j\ge1}2^{-j}\min(1,p_j(v-w)),\\
\beta(\rho,\eta)=\sup_{|H|\le1,\ \mathrm{Lip}_d(H)\le1}|\rho(H)-\eta(H)|,
\end{gathered}
\end{equation}
where $(\chi_jv)(g)=v(\chi_jg)$ for smooth periodic $g$, so that
$p_j(v)^{2}=\sum_k\lambda_{j,k}^{-2}|v(\chi_j\conj{e_{j,k}})|^{2}$, with $v$
extended complex-linearly. $X$ carries the topology of $d$ and its Borel
$\sigma$-algebra.
\end{definition}

Proposition~\ref{prop:space} proves that $E$ and $X$ are Polish, that $d$ is a
complete metric, that compact smooth evaluations $v\mapsto v(f)$ are continuous
and generate the Borel $\sigma$-algebra of $X$, that $\iota$ is a continuous
injection, and that $\beta$ metrizes weak convergence. A \emph{weak law
statement on $X$} always means convergence of $\int H\,d\rho$ for every real
bounded continuous $H:X\to\R$.

\subsection{The two process-neuron architectures}\label{sec:neurons-def}

Both architectures produce a random sign field with values $\pm1/2$, mean zero and the same endpoint covariance $\Cend(z-z')=\tfrac14e^{-\mu|z-z'|_1}$ (Theorems~\ref{thm:E-product}(b) and~\ref{thm:E-plane}(c)). The product neuron multiplies two independent telegraph signals, one along each coordinate, each flipping at the points of a Poisson process. The Mondrian neuron colors the cells of a random rectangular partition with independent signs; two points lie in a common cell with probability $e^{-\mu|z-z'|_1}$ (Theorem~\ref{thm:E-box}(e) with $s=\mu$). Because the covariances agree, both architectures have the same Gaussian reference in the width limit (Proposition~\ref{prop:width}).

The neurons are defined exactly; Section~\ref{sec:endpoint} states their
properties, proved in Appendix~\ref{app:E}. Random objects live on a probability
space $(\Gamma,\cA,\PP)$; $\mathrm{Exp}(c)$ has density $ce^{-ct}$ on
$(0,\infty)$, $U(0,1)$ is uniform on $(0,1)$, and a \emph{uniform sign} is
uniform on $\{-1,1\}$.

\begin{definition}[Product (telegraph) neuron]\label{def:product}
Let $\cN$ be a Poisson process on $\R$ of intensity $\mu/2$ and $\epsilon$ an
independent uniform sign. The stationary symmetric telegraph process of flip
rate $\mu/2$ is $\sigma(t)=\epsilon(-1)^{\cN((0,t])}$ for $t\ge0$ and
$\sigma(t)=\epsilon(-1)^{\cN((t,0])}$ for $t<0$ (set $\sigma=\epsilon$ on the null
event where $\cN$ is not locally finite or charges $0$). With $\tau$ an
independent copy, the product neuron is $g(z)=\tfrac12\sigma(z_1)\tau(z_2)$.
\end{definition}

\begin{definition}[Boxes, colored partitions and the canonical Mondrian]\label{def:mondrian-box}
A \emph{box} is $C=[a_1,b_1]\times[c_1,d_1]$ with $a_1\le b_1$, $c_1\le d_1$, and
half-perimeter $p(C)=(b_1-a_1)+(d_1-c_1)$. For a box $C$ with nonempty interior,
$\Part^{\mathrm{col}}(C)$ is the set of finite sets of pairs $(R,\xi)$, $R$ a
closed box with nonempty interior and $\xi\in\{-1,1\}$, such that the boxes $R$
have pairwise disjoint interiors and union $C$; it is encoded by finite counting
measures on $\R^{4}\times\{-1,1\}$ (Lemma~\ref{lem:E-partspaces}). For boxes
$B\subseteq C$ with nonempty interiors the restriction is
$\resB_B(P)=\{(R\cap B,\xi):(R,\xi)\in P,\ \mathrm{int}(R\cap B)\neq\emptyset\}$;
images under translations and under $\theta_1,\theta_2$ act on closed boxes. The
\emph{lower-left representative} of $R=[a_1,b_1]\times[c_1,d_1]$ is
$R^{\mathrm{ll}}=[a_1,b_1)\times[c_1,d_1)$.

Let $\mathrm{Tr}$ be the set of finite words over $\{0,1\}$ (root $\emptyset$)
and let $(E_w,D_w,X_w,\xi_w)_{w\in\mathrm{Tr}}$ be independent, $E_w\sim\mathrm{Exp}(1)$,
$D_w,X_w\sim U(0,1)$, $\xi_w$ a uniform sign. For a box $C$ and budget $s\ge0$
put $R_\emptyset=C$, $b_\emptyset=0$; a born node $w$ with box
$R_w=[a_1,b_1]\times[c_1,d_1]$ dies at $d_w=b_w+E_w/p(R_w)$ ($d_w=\infty$ if
$p(R_w)=0$). If $d_w<s$, $w$ splits: if $D_w<(b_1-a_1)/p(R_w)$ at the vertical
line $x=a_1+X_w(b_1-a_1)$, otherwise at the horizontal line $y=c_1+X_w(d_1-c_1)$;
its children $w0,w1$ receive the two closed pieces (left/lower, right/upper) and
are born at $d_w$. Born nodes with $d_w\ge s$ are leaves. The output
$\Pi^{\mathrm{can}}(C,s)=\{(R_w,\xi_w):w\text{ a leaf}\}$ has law
$\Mcol(C,s)$. Thus each cell with side lengths $w',h'$ splits at rate $w'+h'$,
orthogonally to the first axis with probability $w'/(w'+h')$, at a uniform
position, and the daughters evolve independently until the budget is exhausted.
\end{definition}

\begin{definition}[Whole-plane Mondrian neuron; closed cells and half-open
representatives]\label{def:mondrian}
Let $B_n=[-n,n]^{2}$. A \emph{consistent colored Mondrian tessellation} of budget
$\mu$ is a sequence $P=(P_n)_{n\ge1}$ of random elements
$P_n\in\Part^{\mathrm{col}}(B_n)$ on one probability space with
$P_n\sim\Mcol(B_n,\mu)$ for every $n$ and, almost surely,
$\resB_{B_n}(P_{n+1})=P_n$ for all $n$ (the \emph{consistency event}). For
$z\in\R^{2}$ let $n_z=\min\{n\ge1:z\in[-n,n)^{2}\}$, and for $n\ge n_z$ let
$(R_n(z),\xi_n(z))$ be the cell of $P_n$ with $z\in R_n(z)^{\mathrm{ll}}$. The
\emph{closed geometric cell} is $R(z)=\bigcup_{n\ge n_z}R_n(z)$ and the
\emph{half-open representative} is $H(z)=\bigcup_{n\ge n_z}R_n(z)^{\mathrm{ll}}$.
The \emph{Mondrian neuron} is the field
$g(z)=\tfrac12\xi_{n_z}(z)$, defined for every $z$ and every outcome
(lower-left convention). An \emph{admissible boundary convention} is any
$\cB(\R^{2})\otimes\cA$-measurable $g'$ such that on the consistency event
$g'(z)=\xi/2$ for some cell $(R,\xi)$ of $P_{n_z+1}$ with $z\in R$. Maps act on
tessellations through closed cells; half-open representatives of an image are
obtained by applying the lower-left convention to its closed cells.
\end{definition}

\begin{definition}[Neural fields]\label{def:neuralfield}
For $b\in\{\bprod,\bMon\}$, $N\ge1$ and independent copies
$g^{b}_{1,a},\dots,g^{b}_{N,a}$ of the neuron of Definition~\ref{def:product} or
\ref{def:mondrian} with $\mu=1/a$, on the product $(\Gamma_b,\PP_b)$ of $N$
copies of its probability space, the neural field of width $N$ is
\begin{equation}\label{eq:neuralfield}
F^{b}_{N,a}=N^{-1/2}\sum_{i=1}^{N}g^{b}_{i,a},\qquad |F^{b}_{N,a}|\le\sqrt N/2 .
\end{equation}
\end{definition}

\begin{remark}[What the width counts]\label{rem:width}
The width $N$ counts process neurons. Each neuron uses countably many random
variables; on every bounded region it has almost surely finitely many telegraph
jump lines or Mondrian cell edges, but that number is random, unbounded and
grows with the region. A finite-width stage is therefore not a random field
encoded by finitely many global parameters for the whole plane. For background
on the Mondrian process see \cite{RoyTeh2008}; for stationary tessellations
stable under iteration see \cite{NagelWeiss2005}. These are background
attributions; every property used below is proved in Appendix~\ref{app:E}.
\end{remark}

\subsection{Potentials, Wick and coupling conventions}\label{sec:potentials}

The potential has three parts with separate jobs. The Wick-ordered quartic $W_{a,\Lambda}$ on the physical square is the interaction that survives in the limit; its subtraction $-6c_ax^{2}+3c_a^{2}$ lets the interaction converge as $a\downarrow0$, although $c_a\to\infty$ (Theorem~\ref{thm:W}(a)). The negative quadratic $-T_{\Omega,a}$ implements the mass shift of Section~\ref{sec:massshift}. The stabilizing quartic $\varepsilon J_{\Omega,\Lambda}$ on $\Omega\setminus\Lambda$ keeps the weight $e^{-V}$ bounded, see \eqref{eq:Mfull}, which the width limit uses; it is removed at fixed $a,\Omega,\Lambda$ (Proposition~\ref{prop:stab}).

\begin{definition}[Integrated potentials]\label{def:potentials}
For $v\in E$, a physical square $\Lambda$, an auxiliary square $\Omega\supseteq\bar\Lambda$,
a regulator $a$ and $\varepsilon\ge0$, put
\begin{align}
W_{a,\Lambda}(v)&=u\int_{\Lambda}\bigl(v^{4}-6c_a v^{2}+3c_a^{2}\bigr)\,dz,&
J_{\Omega,\Lambda}(v)&=\int_{\Omega\setminus\Lambda}v^{4}\,dz,\nonumber\\
T_{\Omega,a}(v)&=\frac{\kappa_a^{2}}{2}\int_{\Omega}v^{2}\,dz,&
V_{\varepsilon,\Omega,a;\Lambda}&=W_{a,\Lambda}+\varepsilon J_{\Omega,\Lambda}-T_{\Omega,a}.
\label{eq:potential}
\end{align}
Equivalently $V_{\varepsilon,\Omega,a;\Lambda}(v)=\int_\Omega v_V(z,v(z))\,dz$ with
the Borel integrand
\begin{equation}\label{eq:vV}
v_V(z,x)=\ind_\Lambda(z)\,u(x^{4}-6c_ax^{2}+3c_a^{2})+\varepsilon\ind_{\Omega\setminus\Lambda}(z)\,x^{4}-\tfrac12\kappa_a^{2}x^{2}.
\end{equation}
The negative quadratic term acts on all of $\Omega$, including $\Lambda$. These
are integrated potentials of ordinary functions; they are never evaluated as
pointwise polynomials of an arbitrary distribution.
\end{definition}

\begin{remark}[Wick and coupling conventions]\label{rem:wick}
At a positive regulator, $x^{4}-6c_ax^{2}+3c_a^{2}$ is the fourth Wick polynomial
at the variance $c_a=C_a(0)$ of the regulated Gaussian physical reference field
(the centered Gaussian field with covariance $C_a$): own-mass, own-regulator Wick
ordering of that Gaussian reference. The endpoint variance $1/4$ is never a Wick
constant, and no finite-width or interacting law is asserted to have
variance~$c_a$. At $a=0$ the interaction is the common-space fourth-chaos $L^{2}$ Wick
limit $U_{0,\Lambda}$ constructed in Section~\ref{sec:wick}; its precise
construction statement is Theorem~\ref{thm:W}(b). The coefficient $u$
multiplies the Wick quartic with no factorial divisor, as in \cite{FS1977}
(polynomial $\lambda X^{4}$) and \cite{AHKZ1989} ($\lambda\int{:}P(\phi){:}_0$,
$P(X)=X^{4}$); the finite-volume neural construction of \cite{Frank2026v2} uses
$\lambda/4!$, so its $\lambda/24$ plays the role of $u$. In natural units the
field is dimensionless in two dimensions; $m$, $\mu$, the telegraph flip rate
$\mu/2$ and the Mondrian budget $\mu$ have units of inverse length; $u$ and
$\kappa_a^{2}$ of inverse length squared; $u/m^{2}$ is the dimensionless
coupling, and the dilation of Section~\ref{sec:target} preserves it.
\end{remark}

\subsection{The normalized laws}\label{sec:laws}

Definition~\ref{def:laws} introduces the laws met along the stages of Section~\ref{sec:massshift}. The neural laws (a) are the objects of the main theorem, and the full Gaussian tilt (b) is their width limit. For $\varepsilon>0$, the residual laws (c) rewrite that tilt as a quartic tilt of the Gaussian law $G'_{\Omega,a}$; they also define its stabilizer-free version at $\varepsilon=0$. The physical laws (d) arise when the auxiliary square is removed; at $a=0$ they are the continuum laws on a fixed square. Item (e) names the growing-volume laws and the target. Remark~\ref{rem:chain} lists the four fixed-volume arrows joining (a)--(d).

The following objects are well defined by results of
Sections~\ref{sec:endpoint}--\ref{sec:wick} (finiteness and positivity of
normalizers, existence of Gaussian versions and of the $L^{2}$ Wick limit); the
proofs of those results, in Appendices~\ref{app:E}--\ref{app:W}, do not use
this definition.

\begin{definition}[Laws]\label{def:laws}
Fix a regulator $a$, a physical square $\Lambda$ and an auxiliary square
$\Omega\supseteq\bar\Lambda$.
\begin{enumerate}[label=(\alph*)]
\item \emph{Neural laws} ($\varepsilon>0$). With
  $Z^{b}_{N}=\E_b\exp[-V_{\varepsilon,\Omega,a;\Lambda}(F^{b}_{N,a})]$,
  \begin{equation}\label{eq:neurallaw}
  P^{b}_{N,\varepsilon,\Omega,a;\Lambda}(B)=\frac{\E_b\bigl[\ind_{B}(\iota F^{b}_{N,a})\,
  e^{-V_{\varepsilon,\Omega,a;\Lambda}(F^{b}_{N,a})}\bigr]}{Z^{b}_{N}},\qquad B\in\cB(X).
  \end{equation}
  These are \emph{annealed} laws: the weight reweights the joint law of all
  neuron randomness. (The same formula is meaningful at $\varepsilon=0$ for
  finite $N$, Theorem~\ref{thm:E-RP}(a); that finite-width extension is not
  used below.)
\item \emph{Full Gaussian tilt} ($\varepsilon>0$). Let $G_\mu$ be the law on
  $C(\R^{2})$ of the continuous centered Gaussian field $\gamma_\mu$ with
  covariance $\Cend$ (Theorem~\ref{thm:G-endpoint}). With
  $Z_{\mathrm{full}}=\E_{G_\mu}e^{-V(\gamma_\mu)}$,
  $\nu^{v}_{\varepsilon,\Omega,a;\Lambda}(B)
  =\E_{G_\mu}[\ind_B(\iota\gamma_\mu)e^{-V(\gamma_\mu)}]/Z_{\mathrm{full}}$.
\item \emph{Mass tilt and residual laws} ($\varepsilon\ge0$). Let
  $\Det_{\Omega,a}=\E_{G_\mu}e^{T_{\Omega,a}(\gamma_\mu)}$ and
  $G'_{\Omega,a}=\Det_{\Omega,a}^{-1}e^{T_{\Omega,a}}G_\mu$, the centered
  Gaussian law on $C(\R^{2})$ with covariance $C'_{\Omega,a}$
  (Theorem~\ref{thm:G-tilt}). Put
  $z_{\varepsilon,\Omega,a;\Lambda}=\E_{G'_{\Omega,a}}e^{-W_{a,\Lambda}-\varepsilon J_{\Omega,\Lambda}}$,
  $d\Res_{\varepsilon,\Omega,a;\Lambda}=z^{-1}_{\varepsilon,\Omega,a;\Lambda}
  e^{-W_{a,\Lambda}-\varepsilon J_{\Omega,\Lambda}}\,dG'_{\Omega,a}$ and
  $\nu^{(2)}_{\Omega,a;\Lambda}=\iota_{*}\Res_{0,\Omega,a;\Lambda}$.
\item \emph{Physical laws} ($0\le a<1/m$). On the common Gaussian space
  $(\Gamma_0,\cF_0,\PP_0)$ of Theorem~\ref{thm:G-common}, let
  $\Phi_a:\Gamma_0\to X$ be the Borel version of Theorem~\ref{thm:G-versions}
  and, for $a>0$, $\tilde\phi_a$ the continuous modification of
  Theorem~\ref{thm:G-point}. Put $U_{a,\Lambda}=W_{a,\Lambda}(\tilde\phi_a)$ for
  $a>0$, $U_{0,\Lambda}=L^{2}(\PP_0)\text{-}\lim_{a\downarrow0}U_{a,\Lambda}$
  (Theorem~\ref{thm:W}(a)), $Z_{a,\Lambda}=\E_0e^{-U_{a,\Lambda}}$ and
  \begin{equation}\label{eq:physlaw}
  \nu_{a,\Lambda}(B)=\E_0\bigl[\ind_B(\Phi_a)\,e^{-U_{a,\Lambda}}\bigr]/Z_{a,\Lambda}.
  \end{equation}
\item \emph{Continuum finite-volume laws and the target.} $\nu_n=\nu_{0,\Lambda_n}$.
  The classical finite-volume laws $\nu_{m,u,\Lambda}$ and the classical target
  $\nu^{S}$ are probabilities on the real tempered distributions $\cS'(\R^{2})$
  with the cylinder $\sigma$-algebra $\cC$ (Definition~\ref{def:classical-laws},
  Theorem~\ref{thm:LR}); $\nu=\nu^{X}_\infty$ is the $X$-valued version of
  $\nu^{S}$ (Theorem~\ref{thm:V1}).
\end{enumerate}
\end{definition}

\begin{table}[ht]
\centering
\small
\begin{tabular}{@{}p{0.2\textwidth}p{0.34\textwidth}p{0.38\textwidth}@{}}
\toprule
Symbol & Meaning & Bounds and traps\\
\midrule
$Z^{b}_{N}$ & neural normalizer & $e^{-|\Omega|M_N}\le Z^b_N\le e^{|\Omega|M_N}$, $M_N=\sup\{|v_V(z,x)|:z\in\Omega,|x|\le\sqrt N/2\}$; not $\ge1$ in general\\
$Z_{\mathrm{full}}$ & full Gaussian normalizer & $Z_{\mathrm{full}}=\Det_{\Omega,a}\,z_{\varepsilon,\Omega,a;\Lambda}$ for $\varepsilon>0$; no monotonicity in $\Omega$ is asserted (Remark~\ref{rem:chain})\\
$\Det_{\Omega,a}$ & quadratic-tilt normalizer & $1\le\Det_{\Omega,a}\le\exp(\kappa_a^{2}|\Omega|/(8a^{2}m^{2}))$; nondecreasing in $\Omega$\\
$z_{\varepsilon,\Omega,a;\Lambda}$ & residual normalizer & $e^{-3u|\Lambda|c_a^{2}-3\varepsilon|\Omega\setminus\Lambda|c_a^{2}}\le z_\varepsilon\le e^{6u|\Lambda|c_a^{2}}$\\
$Z_{a,\Lambda}$ & physical normalizer & $1\le Z_{a,\Lambda}\le\Xi_{1}(\Lambda)$, $0\le a<1/m$; physical laws only; no uniformity in $\Lambda$\\
\bottomrule
\end{tabular}
\caption{Normalizers. The bound $Z\ge1$ holds only for the physical laws on the
common Gaussian space.}
\label{tab:normalizers}
\end{table}

\subsection{Tests, distributions and pairings}

\begin{definition}[Tests and complex-linear pairings]\label{def:tests}
For an open set $U$ of a Euclidean space, $\cD(U)=C_c^{\infty}(U;\C)$ carries
its standard inductive-limit (LF) topology over the Fr\'echet spaces $\cD_K(U)$,
$K\subset U$ compact, with seminorms
$\|F\|_{C^{M}}=\max_{|\alpha|\le M}\sup|\partial^{\alpha}F|$. The dual
$\cD'(U)$ consists of continuous \emph{complex-linear} functionals; pairings are
linear in the test. A real distribution $v$ extends by $v(f+ig)=v(f)+iv(g)$. For
$f_1,\dots,f_k\in\cD(\R^{2})$ put
$(f_1\otimes\cdots\otimes f_k)(z_1,\dots,z_k)=\prod_r f_r(z_r)$. For
$F\in\cD(((-j,j)^{2})^{k})$, regarded as a smooth periodic function on
$\T_j^{k}$ after extension by zero,
\begin{equation}\label{eq:Ajk}
A_{j,k}(F)=\Bigl\|\prod_{r=1}^{k}(1-\Delta_{z_r})F\Bigr\|_{L^{2}(\T_j^{k})}.
\end{equation}
\end{definition}

\subsection{Reflections and half-plane observables}

The positivity statements compare an observable with its mirror image. An observable is $\Sigma^{\theta}_{+}$-measurable when it depends only on smearings of the field by tests supported in the open half-plane $H^{\theta}_{+}$; then $F(\theta v)$ is the same observable applied to the reflected field, and so depends on the opposite half-plane. Inequality \eqref{eq:boundedRP} requires the average of $\conj{F(\theta v)}F(v)$ to be nonnegative. The form $\Pi_\rho(f)$ is the analogous expression for polynomial observables built from compact joint tests, used in clause~(E2) of Theorem~\ref{thm:main}.

\begin{definition}[Central reflections, half-planes and positivity]\label{def:reflections}
For $\theta\in\{\theta_1,\theta_2\}$ let $H^{\theta}_{+}$ be the open positive
half-plane ($\{z_1>0\}$ for $\theta_1$, $\{z_2>0\}$ for $\theta_2$), define
$(\theta v)(f)=v(f\circ\theta)$ for $v\in X$, and let $\Sigma^{\theta}_{+}$ be the
$\sigma$-algebra on $X$ generated by the evaluations $v\mapsto v(f)$, real
$f\in C_c^{\infty}(H^{\theta}_{+})$. A Borel probability $\rho$ on $X$ is
\emph{reflection positive across $\theta$ for bounded observables} if
\begin{equation}\label{eq:boundedRP}
\int_{X}\conj{F(\theta v)}\,F(v)\,\rho(dv)\ge0
\end{equation}
for every bounded complex $\Sigma^{\theta}_{+}$-measurable $F$. For $J\ge0$,
$f_0\in\C$ and $f_r\in\cD((H^{\theta}_{+})^{r})$, $1\le r\le J$, put
$(\theta_r\conj{f_r})(z_1,\dots,z_r)=\conj{f_r(\theta z_1,\dots,\theta z_r)}$ (value
$\conj{f_0}$ at $r=0$) and
\begin{equation}\label{eq:polyform}
\cO_f(v)=\sum_{r=0}^{J}v^{\otimes r}(f_r),\qquad
\Pi_{\rho}(f)=\sum_{r,s=0}^{J}\cS^{\rho}_{r+s}\bigl((\theta_r\conj{f_r})\otimes f_s\bigr),
\end{equation}
where $v^{\otimes r}$ and $\cS^{\rho}_{k}$ are the tensor evaluations and
correlation distributions of Section~\ref{sec:K1}, $v^{\otimes 0}(c)=c$ and
$\cS^{\rho}_{0}(c)=c$.
\end{definition}

% ===========================================================================
\section{Classical inputs}\label{sec:classical}

The inputs answer distinct questions. \ref{prem:C1} supplies the infinite-volume states, their Osterwalder--Schrader properties and local regularity, and convergence of free-boundary finite-volume laws once the pressure is differentiable. \ref{prem:C2} gives uniqueness at weak coupling, which makes the two states coincide; Definition~\ref{def:threshold} turns its existential constants into the regime $0<u/m^{2}<g_{*}$. \ref{prem:C3} provides reconstruction and a kernel theorem; \ref{prem:C4} provides lattice approximation and correlation inequalities, and \ref{prem:C5} a third-cumulant inequality, both used for the volume-uniform bounds of Theorem~\ref{thm:V1}. \ref{prem:C6} lists standard tools.

The following published results are used with attribution. Their proofs are
not audited here. Each formulation records exactly the content used; its
application to the present objects is verified at the point of use. Page
references are to the printed pages; Appendix~\ref{app:provenance} records for
each page whether the classical-input review read it as a page image or as
supplied text.

\emph{Conventions shared by \ref{prem:C1}--\ref{prem:C4}.} For a mass $M>0$,
$P_M$ is the centered Gaussian probability on $(\cS'(\R^{2}),\cC)$ with covariance
$\langle f,(-\Delta+M^{2})^{-1}g\rangle$ (the free measure $\mu_0$ of
\cite[p.~496]{FS1977}); for open $U$, $\cC_U$ is generated by $\psi\mapsto\psi(f)$
with real $f\in\cS$, $\supp f\subset U$. The interaction is
$\lambda\int_\Lambda{:}\phi^{4}{:}$ with whole-plane, same-mass Wick ordering and no
factorial; ``free boundary'' means the whole-plane Gaussian with the interaction
restricted to the bounded region and zero external field
\cite[pp.~496--497]{FS1977}, \cite[(4), p.~685]{AHKZ1989},
\cite[(IX.5a), pp.~311--312]{GRS1976}. In these sources ${:}\phi^{4}{:}$ denotes
Gaussian Wick normal ordering at the whole-plane reference covariance; for
degree four this is the fourth-chaos variable realized as the $L^{2}$ limit of
own-variance Wick quartics of mollified fields. Lemma~\ref{lem:classical-wick}
proves that this classical variable exists and has positive finite exponential
normalizers using only the free Gaussian covariance and \ref{prem:C6}. The
source finite-volume convention is applied only after these provisos are
proved. The separate identification with the common-space regulator limit
is Proposition~\ref{prop:identify}; it is not an input to the restricted
classical route.

\begin{classical}[Fr\"ohlich--Simon]\label{prem:C1}
Fix $M>0$ and $P(X)=\lambda X^{4}$, $\lambda>0$ \cite{FS1977}.
\begin{enumerate}[label=(\roman*)]
\item \cite[Thm.~4.1, p.~509]{FS1977}: the Section-4 limits $\nu_{+}$ and
  $\nu_{-}$ at zero external field are probabilities on $(\cS',\cC)$ satisfying
  all Osterwalder--Schrader axioms except a priori clustering; used are proper
  Euclidean invariance and Osterwalder--Schrader positivity, in the measure
  form of \cite[p.~495, (a)--(b)]{FS1977} or for moments of positive-time
  (possibly time-ordered) tests; both readings are treated in
  Section~\ref{sec:target}.
\item \cite[\S6 definition, pp.~513--514, Remarks~1--2 and Thm.~6.1; Thm.~7.2,
  p.~517; (1.2b), p.~498]{FS1977}: $\nu_+$ and $\nu_-$ are ultraregular: for every
  bounded open $B$, the restriction of $\nu_\pm$ to $\cC_B$ has a density $f_B$ with
  respect to $P_M$ lying in every $L^{p}(P_M)$, $p<\infty$, and
  $f_B=e^{-U_B}e^{-\alpha_\infty|B|}g_{\partial B}$ with $g_{\partial B}>0$ almost
  everywhere, in every $L^{p}$, and measurable for the boundary $\sigma$-algebra,
  which by \cite[p.~514, Remark~2]{FS1977} is either the intersection of
  $\cC_R$ over open $R\supseteq\partial B$, or generated by $\phi(f)$ with
  $f\in H_{-1}$ of compact support in $\partial B$. Perimeter estimates are
  stated for rectangles.
\item \cite[Cor.~4.3(ii), p.~510]{FS1977}: if $\nu_+=\nu_-$, the pressure
  $\alpha_\infty(P-hX)$ is differentiable at $h=0$ (only this implication is used).
\item \cite[Thm.~5.2, p.~512]{FS1977}: if the pressure is differentiable at
  $h=0$, the classical boundary-condition finite-volume measures, including
  the free ones, converge as $\Lambda$ increases to $\R^{2}$ (ordinary
  exhaustion, not the iterated convention of \cite[p.~503]{FS1977}) to
  $\nu_+=\nu_-$. It is used only along centered squares, and only for the
  characteristic functions of finitely many compact smooth tests; the
  conclusion is the same under the moment or Laplace-functional readings of
  convergence (Lemma~\ref{lem:readings}).
\end{enumerate}
\end{classical}

The two boundary algebras in C1(ii) are denoted BC1 (the neighborhood-intersection algebra) and BC2 (the boundary-supported $H_{-1}$ algebra, with the completion specified below). The next remark specifies the two positivity readings.

\begin{remark}[Literal alternative readings]\label{rem:C1-readings}
The positivity alternatives above are retained separately. In the measure
reading (b$_M$), the inequality is for every complex
$F\in L^2(\cS',\Sigma_+,\nu_\pm)$. In the weaker reading (b$_P$), it is
for every polynomial $F=\sum_r c_r\prod_{i=1}^{n_r}\psi(f_{r,i})$,
$c_r\in\C$, $n_r\ge0$, real Schwartz $f_{r,i}$, whose individual product
supports lie in $0<(z_1)_2<\cdots<(z_{n_r})_2$; empty products are constants.
For C1(iv), the alternatives are convergence of the characteristic value
for every real compact $g$, convergence of every moment of $\psi(g)$, or
convergence of $\E e^{t\psi(g)}$ for every real $t$. The first supplies all
frequencies by replacing $g$ with $tg$. In BC2 the boundary algebra is
completed by $P_M$-null sets; its variables for boundary-supported real
$H^{-1}$ tests are defined by their Gaussian $L^2$ limits of Schwartz
approximants. Both BC1 and BC2, both positivity readings and all three
convergence readings are treated below.
\end{remark}

\begin{classical}[Albeverio--H{\o}egh-Krohn--Zegarlinski]\label{prem:C2}
\cite{AHKZ1989} work on a set $\Omega_A\subset\cS'(\R^{2})$ with the traces of the
weak-Borel sets; for closed $D$, $\cB_D=\bigcap_{\Lambda'\supseteq D\text{ open}}\cB_{\Lambda'}$
and $\Sigma_{\Lambda^{c}}=\cB_{\Lambda^{c}}\cap\Omega_A$ \cite[p.~684]{AHKZ1989}; the
specification consists of kernels $E^{\eta}_{\Lambda^{c}}$ with (a) each
$E^{\eta}_{\Lambda^{c}}$ a probability with $E^{\eta}_{\Lambda^{c}}(F)=F(\eta)$ for
$\Sigma_{\Lambda^{c}}$-measurable $F$, (b) $\eta\mapsto E^{\eta}_{\Lambda^{c}}(F)$
$\Sigma_{\Lambda^{c}}$-measurable, (c) compatibility; a Gibbs measure satisfies
$\mu(E^{\cdot}_{\Lambda^c}F)=\mu(F)$ for all $F\in L^{1}(\mu)$
\cite[(7)--(8), p.~686]{AHKZ1989}; $\cG_{\mathrm{ur}}$ is the ultraregular Gibbs
class \cite[(10)--(12), p.~686]{AHKZ1989}.
\begin{enumerate}[label=(\roman*)]
\item \emph{Membership} \cite[pp.~686--687]{AHKZ1989}, citing
  \cite[Thm.~7.2]{FS1977}: both states $\nu_+$ and $\nu_-$ of \ref{prem:C1} at
  mass $m_0$ belong to $\cG_{\mathrm{ur}}$ (\cite{AHKZ1989} prints the maximal
  measures without sign suffix; \cite[Thm.~7.2]{FS1977} names both signs).
\item \emph{Regime and uniqueness} \cite[Theorem, p.~687; after (47), p.~691]{AHKZ1989}:
  there are $M_0,\lambda_0>0$ such that for every $m_0>M_0$ and $0<\lambda<\lambda_0$
  the class $\cG_{\mathrm{ur}}$ for $\lambda X^{4}$ at mass $m_0$ is a singleton.
\item \emph{Boundary-condition independence} \cite[(13)--(14), p.~687; completed
  p.~691]{AHKZ1989}: in that regime, for $\mu,\tilde\mu\in\cG_{\mathrm{ur}}$ and every
  polynomial function $F$ measurable for $\Sigma_{\Lambda_0}$, $\Lambda_0$ bounded,
  \[
  \lim_{\Lambda\in F_0}(\mu\otimes\tilde\mu)\bigl|E^{\eta}_{\Lambda^{c}}F-E^{\tilde\eta}_{\Lambda^{c}}F\bigr|=0,
  \]
  where $F_0$ is the increasing sequence of centered squares. ``Polynomial
  function'' is read as $Q(\phi(f_1),\dots,\phi(f_k))$ with $Q$ a real polynomial
  and real $f_i\in\cS$ supported in a bounded open $\Lambda_0$; this reading is
  supported by \cite[p.~689, after (31)]{AHKZ1989}, which localizes $F$ in a
  finite union of unit cubes (Appendix~\ref{app:provenance}).
\end{enumerate}
\end{classical}

\begin{definition}[Fixed admissible pair, threshold and regime]\label{def:threshold}
Fix once and for all one pair $(M_0,\lambda_0)$ of \ref{prem:C2}(ii), a reference
mass $M_{\mathrm{ref}}>M_0$ and $0<\lambda_*\le\lambda_0$, and put
$g_{*}=\lambda_*/M_{\mathrm{ref}}^{2}>0$. The \emph{regime} consists of the pairs
$(m,u)$ with $m>0$, $u>0$ and $0<u/m^{2}<g_{*}$. This instantiates an existential
statement: no numerical value of $M_{\mathrm{ref}}$, $\lambda_*$ or $g_*$ is
supplied and no larger coupling range is inferred. For $(m,u)$ in the regime put
$s=m/M_{\mathrm{ref}}$ and $\lambda=u/s^{2}=uM_{\mathrm{ref}}^{2}/m^{2}$; then
$0<\lambda<\lambda_*$.
\end{definition}

\begin{classical}[Reconstruction and tempered kernels]\label{prem:C3}
\begin{enumerate}[label=(\roman*)]
\item \cite[Thm.~1.0, first assertion, and conditions (a)--(d), p.~495]{FS1977}:
  let $\mu$ be a probability on $(\cS'(\R^{2}),\cC)$ with (a) invariance under
  proper Euclidean motions; (b) $\int\conj{\Theta F}\,F\,d\mu\ge0$ for
  $\Sigma_+$-measurable $F$, where $\Theta$ reflects the last coordinate $z_2$ and
  $\Sigma_+$ is generated by $\phi(f)$, $f\in\cS$, $\supp f\subset\{z_2>0\}$,
  and the inequality is required for every $F\in L^2(\Sigma_+,\mu)$; (c) a
  norm $|||\cdot|||$ continuous on $\cS$ such that $f\mapsto\int e^{\phi(f)}d\mu$ is
  uniformly bounded and $|||\cdot|||$-continuous on $\{|||f|||\le1\}$; (d)
  ergodicity of the translation group of $\R^{2}$. Then the moments of $\mu$ are
  the Schwinger functions of a unique relativistic quantum field theory
  satisfying all the Wightman axioms, listed on \cite[p.~494]{FS1977} as
  temperedness, Poincar\'e invariance, positivity, spectrum condition, locality
  and clustering; by \cite[p.~495, Remark~2]{FS1977}, (d) implies uniqueness of
  the physical vacuum. \cite{FS1977} states this as a consequence of
  \cite{OS1973,OS1975}, whose bodies were not available.
\item \cite[p.~494]{FS1977}: the Wightman functions are boundary values of
  functions holomorphic in the permuted extended tube, which contains the
  Euclidean points with pairwise distinct arguments; the Schwinger functions are
  their values there and coincide with the moments. Uniqueness of a function
  holomorphic in a tube with vanishing boundary values is used as a standard
  attributed fact \cite{StreaterWightman1964} (not available).
\item \cite[Lecture~15, p.~2, Thm.~1]{Melrose2016Lecture15}: $K\in\cS'(\R^{N}\times\R^{n})$
  corresponds bijectively to linear maps $L:\cS(\R^{n})\to\cS'(\R^{N})$ with
  $\psi\mapsto(L\psi)(\varphi)$ continuous for each $\varphi$, through
  $(L\psi)(\varphi)=K(\varphi\otimes\psi)$; the pairing is bilinear.
\end{enumerate}
Not used: \cite[Thm.~4.4]{FS1977}, an identification of an
Osterwalder--Schrader semigroup with a Wightman energy, the spectrum condition as
an input, and tube uniqueness as an input to (d). Clustering, ergodicity and
the vacuum statements of the target are obtained in the order of
Theorem~\ref{thm:CT}: Euclidean mixing first, reconstruction afterwards.
\end{classical}

Below, C3(ii') is shorthand for the tube-uniqueness assertion contained in C3(ii).

\begin{classical}[Guerra--Rosen--Simon, free boundary condition]\label{prem:C4}
Fix $m,u>0$ and a centered square $\Lambda$ of side $L$; take meshes
$\delta=L/(2k+1)$, so that $\partial\Lambda$ lies midway between lattice sites
\cite[\S IX.1, pp.~309--310]{GRS1976}. Let $I_\delta=\{i\in\Z^{2}:\delta i\in\Lambda\}$,
$A_\delta=(m^{2}\delta^{2}+4)I-\mathrm{adjacency}$ on $\Z^{2}$
\cite[(IX.2)]{GRS1976}, $K_\delta=(A_\delta^{-1})|_{I_\delta\times I_\delta}$ and
$B_\delta=K_\delta^{-1}$ (the free precision, \cite[(IX.4c), p.~311]{GRS1976}, there
written $A^{F}_\Lambda$), whose off-diagonal entries are nonpositive; the diagonal
$c_\delta$ of $K_\delta$ is the whole-lattice variance. The normalized lattice law
has density proportional to
$\exp[-\tfrac12q^{T}B_\delta q-u\delta^{2}\sum_{i\in I_\delta}(q_i^{4}-6c_\delta q_i^{2}+3c_\delta^{2})]$
\cite[pp.~311--312]{GRS1976}. For real $h\in C_c^{\infty}(\Lambda)$ put
$Y_\delta(h)=\delta^{2}\sum_ih(\delta i)q_i$ \cite[(IX.19), p.~316]{GRS1976}.
\begin{enumerate}[label=(\roman*)]
\item \emph{Fixed-volume convergence} \cite[(IX.19)--(IX.21), p.~316; Lemmas~IX.3
  (with the remark covering the sharp indicator of $\Lambda$) and IX.4, p.~317;
  review to p.~319]{GRS1976}: on a common Gaussian realization of mass $m$, the
  lattice smearings converge to the continuum smearings and the lattice
  interaction densities to the continuum density $e^{-u\int_\Lambda{:}\phi^4{:}}$
  in every finite Gaussian $L^{p}$, with positive normalizer limits; the
  smeared Schwinger functions converge.
\item \emph{Correlation inequalities} \cite[\S IX.2(i),(iii), p.~324]{GRS1976},
  boundary condition F, zero external field: for nonnegative real
  $h_1,\dots,h_r\in C_c^{\infty}(\Lambda)$,
  $\E_{\nu_\Lambda}\prod_i\phi(h_i)\ge0$, and for nonnegative
  $g\in C_c^{\infty}(\Lambda)$ the truncated four-point function smeared with $g$
  is nonpositive.
\end{enumerate}
Not used: \cite[Thm.~IX.5, p.~320]{GRS1976} (periodic and Neumann) and
Thm.~IX.7, and any free--Neumann comparison. \emph{Attribution limit.}
\cite[pp.~316--319]{GRS1976} review the free case using its references [29],
[62] and [28], and sketch a uniform exponential estimate (p.~317); these
antecedents were not available, and the complete classical proof was not audited.
\end{classical}

\begin{classical}[Simon--Griffiths]\label{prem:C5}
\cite[Thm.~1, p.~148; (4) and Thm.~2, (5), p.~151]{SG1973}: for finitely many real
variables with density proportional to
$\prod_iF_i(s_i)\exp(\sum_{i\ne j}a_{ij}s_is_j+\sum_ih_is_i)$, where
$F_i(s)=e^{-a_is^{4}-b_is^{2}}$ with $a_i>0$ and $b_i\in\R$, the ordered-pair
coefficients satisfy $a_{ij}>0$, and $h_i\ge0$, every mixed third cumulant
$\langle s_is_js_k\rangle+2\langle s_i\rangle\langle s_j\rangle\langle s_k\rangle
-\langle s_i\rangle\langle s_js_k\rangle-\langle s_j\rangle\langle s_is_k\rangle
-\langle s_k\rangle\langle s_is_j\rangle$ (repeated indices allowed) is
nonpositive. This is a finite-system statement.
\end{classical}

\begin{classical}[Standard tools]\label{prem:C6}
Standard analysis and probability, non-exhaustively and always with
hypotheses checked where used, including: Fubini--Tonelli, dominated and
monotone convergence, Fatou, H\"older, Minkowski, Jensen (also conditional),
Markov, Chernoff, Paley--Zygmund, the portmanteau and continuous-mapping
theorems, L\'evy's continuity theorem, the method of moments, Prokhorov's
theorem on Polish spaces, Dynkin's $\pi$--$\lambda$ and the monotone-class
theorems, Doob--Dynkin factorization and countable dependence of functions
measurable for generated $\sigma$-algebras, the Lusin--Souslin Borel-image
theorem, conditional expectations, martingale convergence for increasing
filtrations, polynomial density in $L^{2}$ of a finite-dimensional law with
finite exponential moments, Fourier series and Parseval on tori and Fej\'er
kernels, compact Sobolev embeddings on a fixed torus, Banach--Steinhaus on
Fr\'echet spaces, density of finite sums of product tests, Gaussian likelihoods
and the Cameron--Martin shift for finite-dimensional Gaussian vectors, the
Hermite expansion, real and complex $L^2$ duality on finite measure spaces,
Bochner integration in $L^2$, approximate identities, the layer-cake formula,
uniqueness of Fourier transforms of finite measures, and the Schur product
theorem. The following standard formulations are stated explicitly:
\begin{enumerate}[label=(\alph*)]
\item \emph{Hypercontractivity} (Nelson): for $Y$ in the $k$-th homogeneous Wiener
  chaos of a Gaussian Hilbert space and real $r\ge2$,
  $\|Y\|_{r}\le(r-1)^{k/2}\|Y\|_{2}$.
\item \emph{Wick products and chaos}: for centered jointly Gaussian
  $\xi_i,\eta_i$, $\E[{:}\xi_1\cdots\xi_k{:}\,{:}\eta_1\cdots\eta_l{:}]=\delta_{kl}
  \sum_{\pi\in S_k}\prod_i\E[\xi_i\eta_{\pi(i)}]$; the $k$-th chaos is the closed span of
  ${:}\xi_1\cdots\xi_k{:}$, chaoses of different order are orthogonal, chaos of
  order $\ge1$ has mean zero, and finite products of jointly Gaussian variables
  with dense span generate a dense subspace of each chaos.
\item \emph{Hausdorff--Young and Plancherel} on $\R^{2}$ with the stated Fourier
  convention.
\item \emph{Continuity criterion} (Kolmogorov--Chentsov--Lo\`eve): a real process on
  $\R^{d}$ with $\E|Y(z)-Y(z')|^{\alpha}\le c|z-z'|^{d+\beta}$ has a version with
  continuous paths for every outcome.
\item \emph{Kolmogorov extension} for consistent finite-dimensional laws, in
  particular centered Gaussian processes with a given positive semidefinite
  kernel, and for sequences of standard Borel spaces.
\item \emph{Isonormal processes}: every separable real Hilbert space carries an
  isonormal process.
\item \emph{Operators}: Neumann series; Hilbert--Schmidt operators and trace class
  ($\trace T^{*}T=\|T\|^{2}_{\mathrm{HS}}$); the spectral theorem for compact
  self-adjoint operators.
\item \emph{Point processes}: restrictions of a homogeneous Poisson process on $\R$
  to disjoint Borel sets are independent; $t\mapsto-t$ preserves its law; a fixed
  point is almost surely not an atom; $\E(-1)^{\mathrm{Poi}(c)}=e^{-2c}$. Finite
  counting measures on a Polish space form a standard Borel space.
\end{enumerate}
\end{classical}

% ===========================================================================
\section{Main results}\label{sec:main}

Two devices keep the main theorem compact. The enumerations \eqref{eq:enumerations} fix in advance the countable sets from which the regulator, the auxiliary half-side, the stabilizer and the width are chosen. The discrepancy $\dist_n$ of \eqref{eq:distn} adds to the bounded-Lipschitz distance $\beta$ the differences in expected values of the first $n$ powers of the first $n$ local seminorms, so one number controls weak closeness and moment closeness at stage $n$. In Theorem~\ref{thm:main}, clause~(S) asserts the sequence, (0) that its laws are well defined, (A)--(D) the convergence and identification statements, and (E) finite-stage positivity.

Fix once and for all the parameter enumerations
\begin{equation}\label{eq:enumerations}
a(k)=\frac{1}{2mk},\qquad R(k)=k,\qquad \varepsilon(k)=\frac1k,\qquad N(k)=k,\qquad k\in\N .
\end{equation}
For $n\ge1$ and probabilities $\rho,\eta$ on $X$ with finite moments of $p_j^{r}$
for $j,r\le n$, put
\begin{equation}\label{eq:distn}
\dist_n(\rho,\eta)=\beta(\rho,\eta)+\sum_{j=1}^{n}\sum_{r=1}^{n}\bigl|\rho(p_j^{r})-\eta(p_j^{r})\bigr| .
\end{equation}

\begin{theorem}[Main theorem]\label{thm:main}
Fix the admissible pair of Definition~\ref{def:threshold}, and $m>0$, $\ell>0$,
$u>0$ with $0<u/m^{2}<g_{*}$. Assume the classical inputs
\ref{prem:C1}--\ref{prem:C5} in the formulations of Section~\ref{sec:classical}
and the standard tools \ref{prem:C6}. No other hypothesis is made. Fix the
localization system of Definition~\ref{def:fieldspace} and the enumerations
\eqref{eq:enumerations}.

\smallskip\noindent\textbf{(S) One sequence.} There is one deterministic
sequence $(a_n,R_n,\varepsilon_n,N_n)_{n\ge1}$, determined by $m,\ell,u$, the
localization system and the enumerations alone, and the same for both
architectures, with
\begin{gather*}
0<a_n<\min(1/n,1/(2m)),\quad a_n<a_{n-1}/2;\qquad
R_n\in\N,\quad R_n\ge\max(n,n\ell),\quad R_n>R_{n-1};\\
0<\varepsilon_n<1/n,\quad \varepsilon_n<\varepsilon_{n-1}/2;\qquad
N_n\in\N,\quad N_n>n,\quad N_n>N_{n-1}
\end{gather*}
(conditions with index $n-1$ only for $n>1$), such that with
$\Omega_n=[-R_n,R_n]^{2}$ and $P^{b}_{n}=P^{b}_{N_n,\varepsilon_n,\Omega_n,a_n;\Lambda_n}$,
\begin{equation}\label{eq:mainseq}
\dist_n(P^{b}_{n},\nu_n)<2^{-n}\qquad\text{for both $b$ and every $n\ge1$.}
\end{equation}
Apart from the discrepancies $\dist_n$, which are determined by the localization
system, no observable, test function, polynomial or functional enters this
choice; in particular none of the objects quantified in (A)--(E) does. The
physical index $n$ is never passed to a subsequence.

\smallskip\noindent\textbf{(0) Laws.} Every $P^{b}_{n}$ is the Borel probability
\eqref{eq:neurallaw} on $X$ with positive finite normalizer, and every power of
every $p_j$ is integrable under it. The limit $\nu$ is the $X$-valued version of
the classical target $\nu^{S}$.

\smallskip\noindent\textbf{(A) Whole field and local moments.} For both $b$,
$P^{b}_{n}\Rightarrow\nu$ weakly on $X$; for all $j,r\ge1$,
$\sup_{n\ge1,b}P^{b}_{n}(p_j^{r})<\infty$ and $P^{b}_{n}(p_j^{r})\to\nu(p_j^{r})$.

\smallskip\noindent\textbf{(B) Every later compact polynomial tuple.} For every
$k\ge1$, all real $f_1,\dots,f_k\in C_c^{\infty}(\R^{2})$ (repetitions allowed) and
every real or complex polynomial $Q$ on $\R^{k}$ of any degree, all expectations
below are absolutely finite,
\[
\E_{P^{b}_{n}}Q(v(f_1),\dots,v(f_k))\to\E_{\nu}Q(v(f_1),\dots,v(f_k)),\qquad
\sup_{n,b}\E_{P^{b}_{n}}|Q(v(f_1),\dots,v(f_k))|^{q}<\infty\ (q>0),
\]
the same powers are $\nu$-integrable, and the laws of $(v(f_1),\dots,v(f_k))$
converge weakly.

\smallskip\noindent\textbf{(C) Joint compact-test correlations.} Put
$\cL=\{\nu\}\cup\{P^{b}_{n}:n\ge1,\ b\in\{\bprod,\bMon\}\}$ and
$M_{j,r}=\sup_{\rho\in\cL}\int p_j^{r}\,d\rho$, $M_{j,0}=1$. For every $k\ge1$:
\begin{enumerate}[label=(C\arabic*),leftmargin=3em]
\item for every $v\in X$ there is a unique $v^{\otimes k}\in\cD'((\R^{2})^{k})$
  with $v^{\otimes k}(f_1\otimes\cdots\otimes f_k)=\prod_r v(f_r)$ for
  $f_r\in\cD(\R^{2})$; it is symmetric, real in the sense
  $\conj{v^{\otimes k}(F)}=v^{\otimes k}(\conj F)$, $v\mapsto v^{\otimes k}(F)$ is
  continuous on $X$, and for $\supp F\subset((-j,j)^{2})^{k}$,
  $|v^{\otimes k}(F)|\le p_j(v)^{k}A_{j,k}(F)\le3^{k}|\T_j|^{k/2}p_j(v)^{k}\|F\|_{C^{2k}}$;
\item for $\rho\in\cL$, $\cS^{\rho}_{k}(F)=\int_X v^{\otimes k}(F)\,\rho(dv)$ is
  absolutely defined, lies in $\cD'((\R^{2})^{k})$, is the unique distribution
  extending $(f_1,\dots,f_k)\mapsto\int\prod_r v(f_r)\,\rho(dv)$, and
  $\sup_{\rho\in\cL}|\cS^{\rho}_{k}(F)|\le M_{j,k}A_{j,k}(F)$ for
  $\supp F\subset((-j,j)^{2})^{k}$, with $M_{j,k}<\infty$;
\item with $\cS^{b}_{n,k}=\cS^{P^{b}_{n}}_{k}$ and $\cS_{\infty,k}=\cS^{\nu}_{k}$,
  $\cS^{b}_{n,k}(F)\to\cS_{\infty,k}(F)$ for every $F\in\cD((\R^{2})^{k})$, that
  is, in $\sigma(\cD'((\R^{2})^{k}),\cD((\R^{2})^{k}))$.
\end{enumerate}

\smallskip\noindent\textbf{(D) Target identification.} $\cS_{\infty,k}$ is the
restriction to $\cD((\R^{2})^{k})$ of the tempered Schwinger distribution
$\Schw_k$ of $\nu^{S}$ (Theorem~\ref{thm:CT}(v)).

\smallskip\noindent\textbf{(E) Finite-stage central reflection positivity.} Fix
$\theta\in\{\theta_1,\theta_2\}$.
\begin{enumerate}[label=(E\arabic*),leftmargin=3em]
\item Every $P^{b}_{n}$ satisfies \eqref{eq:boundedRP} for every bounded complex
  $\Sigma^{\theta}_{+}$-measurable $F$.
\item For every $J\ge0$, $f_0\in\C$ and $f_r\in\cD((H^{\theta}_{+})^{r})$, and every
  $\rho\in\cL$: $\cO_f$ and $\cO_f\circ\theta$ lie in $L^{2}(\rho)$ with separate
  bounds, every term of $\Pi_{\rho}(f)$ is absolutely finite and
  $\Pi_{\rho}(f)=\int\conj{\cO_f(\theta v)}\cO_f(v)\,\rho(dv)$; for $\rho=P^{b}_{n}$
  it is a nonnegative real number, and $\Pi_{P^{b}_{n}}(f)\to\Pi_{\nu}(f)\ge0$.
\end{enumerate}
\end{theorem}

\begin{theorem}[Target properties and reconstruction; attributed classical
results with verified hypotheses]\label{thm:complement}
Fix the admissible pair and $(m,u)$ in the regime. Assume
\ref{prem:C1}--\ref{prem:C3} and \ref{prem:C6}. The target $\nu^{S}$ of
Theorem~\ref{thm:LR} has the properties (i)--(vii) of Theorem~\ref{thm:CT}:
Euclidean and field-sign invariance; reflection positivity across every line for
square-integrable positive-half-plane observables; strong mixing under
translations, hence clustering in this Euclidean sense and translation
ergodicity; the exponential regularity of \ref{prem:C3}(i)(c) for an explicit
Schwartz-continuous norm; tempered symmetric Schwinger distributions;
non-Gaussianity with local $L^{2}(B^{n})$ kernels and a nonzero even truncated
Schwinger distribution of order at least four on tests in disjoint
Euclidean-time slabs; and, by \cite[Thm.~1.0]{FS1977} applied only after
ergodicity has been established, the scalar relativistic Wightman theory in
$1+1$ dimensions that it determines, with the Wightman axioms listed in
\ref{prem:C3}(i), a unique physical vacuum, and a field that fails Gaussian
factorization. The classical existence and reconstruction theorems are
attributed; the verification of their hypotheses for this target is
project-derived and appears in Section~\ref{sec:target}.
\end{theorem}

Every property of Theorem~\ref{thm:complement} is proved along the restricted
classical route in Lemma~\ref{lem:classical-wick}, the source portion of
Theorem~\ref{thm:LR}, Lemmas~\ref{lem:target-symmetry}--\ref{lem:kernels}
and Theorem~\ref{thm:CT}. This route uses only \ref{prem:C1}--\ref{prem:C3}
and \ref{prem:C6}; the separate common-space identification for the neural
construction uses the Gaussian and Wick foundations. Sign symmetry, including
its use in the even-cumulant witness, has the same restricted source route.

\begin{remark}[Quantifier order; one unchanged sequence]\label{rem:quantifiers}
The order of quantifiers is: fix the admissible pair, hence $g_{*}$ and the
regime; fix $(m,\ell,u)$ in the regime; fix the localization system and the
enumerations; then the sequence of (S) is determined, once, for both
architectures. Only afterwards are the following quantified, each
independently and without reselection of parameters: every $H\in C_b(X)$; every
$j,r$; every $k$, every finite real compact smooth tuple and every real or
complex polynomial of any degree; every $q>0$; every order $k$ and every joint
test $F\in\cD((\R^{2})^{k})$; every reflection $\theta$, every $J$ and every list $f$.
Clauses (A)--(E) use exactly the sequence of (S), with no subsequence. No
countability restriction is placed on the later objects. The constants
$M_{j,k}$ are suprema over all stages and over $\nu$; they are finite by
Proposition~\ref{prop:H1conv} and do not depend on $n$, $b$ or $F$.
\end{remark}

\begin{remark}[What is not claimed]\label{rem:notclaimed}
No effective end-to-end rate or schedule for the full chain of limits and no
parameter cost: the $\sqrt a$ bounds of Theorem~\ref{thm:W}(d),(f) and
Proposition~\ref{prop:reg} are upper bounds at fixed physical volume, with
volume-dependent constants, and they give no rate for the width, stabilizer,
auxiliary-volume or physical-volume limits. No numerical value of $g_{*}$, and
no algorithm deciding the selection predicates. No convergence on Schwartz or
other noncompact tests at the neural stages, no strong-dual $\cS'$ convergence
or tightness, no compact-$L^{2}$ or bounded Borel local tests, no local total
variation or uniform density bound, and no interchange of the ordered limits.
No finite-stage Euclidean invariance or clustering statement and no finite-stage
reflection positivity across lines other than the two central coordinate lines;
nothing is asserted about those other lines either way. No bound uniform over all
neural parameter choices: uniformity holds along the selected sequence only.
No independent neural-estimate construction of the thermodynamic limit (the
volume limit uses classical continuum input), no reconstruction of the neural
approximants, no new interaction, dimension or architecture class, no
scattering statement, no quantitative mass gap, no computational advantage, no
training or sampling procedure, and no priority. ``Interacting'' means the
non-Gaussian quantum field of Theorem~\ref{thm:complement}.
\end{remark}

% ===========================================================================
\section{Field spaces}\label{sec:space}

Sections~\ref{sec:space}--\ref{sec:wick} contain the foundations; none of them uses an interacting law as an input (Remark~\ref{rem:noncircular}). This section establishes that $X$ supports the probabilistic arguments used later. It is Polish, so Prokhorov's theorem applies and $\beta$ metrizes weak convergence; compact smooth smearings are continuous and controlled by one seminorm through \eqref{eq:duality}; the potentials are continuous on $E$, the reflections on $X$, and the map $L:C(\R^{2})\to X$ for uniform convergence on compacts. Readers mainly interested in the construction can treat Proposition~\ref{prop:space} as a toolbox and return to it where it is cited.

\begin{proposition}[Field-space facts]\label{prop:space}
\begin{enumerate}[label=(\roman*)]
\item $(E,d_E)$ is Polish. If $F:\R^{2}\times\Gamma\to\R$ is
  $\cB(\R^{2})\otimes\cA$-measurable and every path is locally fourth-power
  integrable (for instance $|F|\le r$), then $\omega\mapsto[F(\cdot,\omega)]$ is a
  Borel $E$-valued random element.
\item For $v\in X$ and $j\ge1$, $p_j(v)<\infty$; $d$ is a complete separable metric
  on $X$, so $X$ is Polish.
\item \emph{Duality.} For $v\in X$, $j\ge1$ and $f\in C_c^{\infty}((-j,j)^{2};\C)$,
  with $\|f\|_{H^{2}(\T_j)}$ computed for the zero extension,
  \begin{equation}\label{eq:duality}
  |v(f)|\le p_j(v)\,\|f\|_{H^{2}(\T_j)} .
  \end{equation}
  Hence every compact smooth evaluation is $d$-continuous. The Borel
  $\sigma$-algebra of $X$ equals the $\sigma$-algebra generated by the compact
  smooth evaluations, and equals the one generated by the countably many real
  and imaginary parts of $v\mapsto v(\chi_j\conj{e_{j,k}})$, $j\ge1$, $k\in\Z^{2}$.
\item $\iota:E\to X$ is a continuous injection.
\item $\beta$ metrizes weak convergence of Borel probabilities on $X$.
\item For bounded measurable $D\subset\R^{2}$ and $k\in\{2,4\}$, $v\mapsto\int_Dv^{k}$ is
  continuous on $E$; hence $W_{a,\Lambda}$, $J_{\Omega,\Lambda}$, $T_{\Omega,a}$,
  $V_{\varepsilon,\Omega,a;\Lambda}$ and their exponentials are continuous on $E$,
  and also on $C(\R^{2})$ with the topology of uniform convergence on compacts.
\item \emph{Reflections.} For $\theta\in\{\theta_1,\theta_2\}$ and $v\in X$,
  $\theta v\in X$, $\theta\theta v=v$, $(\theta v)(f)=v(f\circ\theta)$, and
  $p_j(\theta v)\le C_jp_{j+1}(v)$ with
  $C_j=\|\chi_j\|_\infty+2\sum_{a=1,2}\|\partial_a\chi_j\|_\infty+\|\Delta\chi_j\|_\infty$; in
  particular $\theta$ is continuous on $X$. For $w\in E$, $\theta\iota w=\iota(w\circ\theta)$.
\item \emph{The map $L$.} $(LG)(f)=\int G\,(1-\Delta)f$ defines a linear map
  $L:C(\R^{2})\to X$ (the distribution $(1-\Delta)G$) with
  $p_j(LG)\le(2j+1)C_j\sup_{[-j-1/2,j+1/2]^{2}}|G|$; $L$ is continuous for uniform
  convergence on compacts.
\item \emph{A countable separating family.} Let $\cD_0$ be the set of real and
  imaginary parts of $\chi_nP$, $n\ge1$, $P$ a trigonometric polynomial on $\T_n$
  with coefficients in $\Q+i\Q$. Two real distributions that agree on $\cD_0$ are
  equal.
\item \emph{Countable factorization.} If $F$ is bounded, complex and
  $\Sigma^{\theta}_{+}$-measurable, there are real $f_1,f_2,\dots\in C_c^{\infty}(H^{\theta}_{+})$
  and a Borel $g:\R^{\N}\to\C$ with $|g|\le\sup|F|$ such that $F(v)=g((v(f_j))_j)$ and
  $F(\theta v)=g((v(f_j\circ\theta))_j)$ for every $v\in X$; in particular
  $v\mapsto\conj{F(\theta v)}F(v)$ is bounded and Borel, and
  $\Sigma^{\theta}_{+}\subset\cB(X)$.
\end{enumerate}
\end{proposition}

\begin{proof}
\emph{(i)} Restriction embeds $E$ isometrically as the closed subspace of compatible
families in $\prod_jL^{4}((-j,j)^{2})$ with the product metric; each factor is a
separable Banach space, so $E$ is complete and separable. For a fixed Borel
representative $s_0$ of an element of $E$, Tonelli makes
$\omega\mapsto\int_{(-j,j)^{2}}|F(z,\omega)-s_0(z)|^{4}dz$ measurable, so
$\omega\mapsto d_E([F(\cdot,\omega)],s_0)$ is measurable; in a separable metric space
every open set is a countable union of open balls.

\emph{(iii), duality.} Let $f\ne0$ be supported in $(-j,j)^{2}$ and let $\tilde f$
be the $2(j+1)$-periodic extension of $f|_{\cC_j}$; it is smooth on $\T_j$, and
$f_k=\int_{\cC_j}\conj{e_{j,k}}f$ satisfies $\sum_k\lambda_{j,k}^{2}|f_k|^{2}=\|(1-\Delta)\tilde f\|^{2}_{L^{2}(\T_j)}<\infty$.
The Fourier series of $\tilde f$ converges in $C^{\infty}(\T_j)$, so
$\chi_j\sum_{|k|\le K}f_ke_{j,k}\to\chi_j\tilde f=f$ in $C^{\infty}$ with supports in
$\supp\chi_j$ (because $\chi_j=1$ on $[-j,j]^{2}\supset\supp f$). Since $v$ is a
distribution, $v(f)=\lim_K\sum_{|k|\le K}f_kv(\chi_je_{j,k})$, and
$v(\chi_je_{j,k})=v(\chi_j\conj{e_{j,-k}})$. Cauchy--Schwarz with the weights
$\lambda_{j,k}^{\pm1}$ and $\lambda_{j,-k}=\lambda_{j,k}$ give \eqref{eq:duality},
provided $p_j(v)<\infty$, which is proved next.

\emph{(ii), finiteness.} Let $u=\chi_jv\in H^{-2}(\R^{2})$ and choose a real
$\zeta\in C_c^{\infty}((-j-\tfrac12,j+\tfrac12)^{2})$ with $\zeta=1$ near $\supp\chi_j$. For a
trigonometric polynomial $P=\sum_{|k|\le K}c_ke_{j,k}$, $\zeta P$ is a compactly
supported smooth function and $\sum_kc_ku(\zeta e_{j,k})=u(\zeta P)$, with
$u(\zeta e_{j,k})=v(\chi_je_{j,k})$. By the product rule
$(1-\Delta)(\zeta P)=\zeta(1-\Delta)P-2\nabla\zeta\cdot\nabla P-P\Delta\zeta$, and Parseval on
$\T_j$ gives $\|\zeta P\|_{H^{2}(\R^{2})}\le C(\zeta)(\sum_k\lambda_{j,k}^{2}|c_k|^{2})^{1/2}$
with $C(\zeta)=\|\zeta\|_\infty+2\sup|\nabla\zeta|+\|\Delta\zeta\|_\infty$. Taking
$c_k=\lambda_{j,k}^{-2}\conj{u(\zeta e_{j,k})}$ in
$|u(\zeta P)|\le\|u\|_{H^{-2}}\|\zeta P\|_{H^{2}}$ gives
$\sum_{|k|\le K}\lambda_{j,k}^{-2}|v(\chi_je_{j,k})|^{2}\le C(\zeta)^{2}\|u\|^{2}_{H^{-2}}$ for all
$K$; reindexing $k\mapsto-k$ yields $p_j(v)\le C(\zeta)\|\chi_jv\|_{H^{-2}(\R^{2})}$.

\emph{(ii), metric and completeness.} Each $p_j$ is a seminorm, so $d$ is a
pseudometric; if $d(v,w)=0$ then $p_j(v-w)=0$ for all $j$ and \eqref{eq:duality}
gives $v=w$. Let $(v_n)$ be $d$-Cauchy. Then $(\chi_jv_n)_n$ is Cauchy in the
Hilbert space of periodic distributions with finite $H^{-2}(\T_j)$ norm. For
$f\in C_c^{\infty}((-j,j)^{2})$, $v_n(f)$ is Cauchy by \eqref{eq:duality}; define
$v(f)=\lim_nv_n(f)$, which does not depend on $j$. On each $\cD_K$,
$|v(f)|\le\sup_np_j(v_n)\,\|f\|_{H^{2}(\T_j)}\le3|\T_j|^{1/2}\sup_np_j(v_n)\|f\|_{C^{2}}$, so $v$
is a distribution. Fatou for series gives
$p_j(v_n-v)^{2}\le\liminf_{l}p_j(v_n-v_l)^{2}\to0$. Finally $v\in X$: for
$\chi\in C_c^{\infty}((-j,j)^{2})$ and $g\in\cS(\R^{2})$, $(\chi v)(g)=v(\chi g)$ and
\eqref{eq:duality} with the product rule give
$|v(\chi g)|\le p_j(v)C(\chi)\|g\|_{H^{2}(\R^{2})}$, so $\chi v\in H^{-2}(\R^{2})$. Hence
$d(v_n,v)\to0$.

\emph{(ii), separability, and (iii), Borel generation.} The map
$J(v)=\bigl((\lambda_{j,k}^{-1}v(\chi_j\conj{e_{j,k}}))_{k}\bigr)_{j}$ takes values in
$\prod_j\ell^{2}(\Z^{2};\C)$, and $p_j(v-w)=\|J_j(v)-J_j(w)\|_{\ell^{2}}$; thus
$d(v,w)=\sum_j2^{-j}\min(1,\|J_j(v)-J_j(w)\|)$ is the restriction of a product
metric on a separable space, and $X$ is separable. Let $\cE$ be the
$\sigma$-algebra generated by the real and imaginary parts of the localized
Fourier evaluations. For fixed $w$, $v\mapsto d(v,w)$ is a countable expression in
$\cE$-measurable functions, so open balls, and by separability all open sets,
lie in $\cE$. The Fourier tests $\chi_j\conj{e_{j,k}}$ have compact support, so
$\cE$ is contained in the $\sigma$-algebra of compact smooth evaluations, which is
contained in $\cB(X)$ because each evaluation is continuous: if
$d(v,w)<2^{-j}$ then $p_j(v-w)\le2^{j}d(v,w)$, and \eqref{eq:duality} applies for
$\supp f\subset(-j,j)^{2}$.

\emph{(iv)} For $v\in E$, $(\iota v)(\chi_j\conj{e_{j,k}})=\langle e_{j,k},\chi_jv\rangle_{L^{2}(\T_j)}$, so
Parseval and $\lambda_{j,k}\ge1$ give
$p_j(\iota v)\le\|\chi_jv\|_{L^{2}}\le(2j+2)^{1/2}\|v\|_{L^{4}((-j-1,j+1)^{2})}$ by H\"older;
$\chi\iota v=\chi v\in L^{2}\subset H^{-2}$. A locally integrable function with zero
integral against every compact smooth test vanishes almost everywhere.

\emph{(v)} Every Borel probability on the Polish space $X$ is tight. If
$\beta(\rho_n,\rho)\to0$, the bounded Lipschitz functions
$(1-k\,\dist_d(\cdot,F))_{+}/\max(1,k)$ give $\limsup\rho_n(F)\le\rho(F)$ for closed
$F$, hence weak convergence. Conversely, let $\rho_n\Rightarrow\rho$ and fix
$\delta>0$; choose a compact $K$ with $\rho(K)>1-\delta$; its open $\eta$-neighbourhood
$K^{\eta}$ has $\rho_n(K^{\eta})>1-2\delta$ for large $n$. Choose a finite $\eta$-net of
$K$ and quantize the values of functions $H$ with $|H|\le1$, $\mathrm{Lip}_d(H)\le1$
at the net points to accuracy $\eta$; this yields finitely many classes, with one
representative each. A function and its representative differ by at most $3\eta$
on $K$ and $5\eta$ on $K^{\eta}$, so
$\beta(\rho_n,\rho)\le10\eta+6\delta+\max|\rho_n(H_i)-\rho(H_i)|$ over the finitely many
representatives; let $n\to\infty$ and then $\eta,\delta\to0$.

\emph{(vi)} $|x^{k}-y^{k}|\le k|x-y|(|x|+|y|)^{k-1}$ and H\"older on $D$; the
inclusion $C(\R^{2})\to E$ is continuous.

\emph{(vii)} For a smooth periodic $f$ on $\T_j$ put
$\zeta=(\chi_j\circ\theta)(f\circ\theta)$, supported in $(-j-\tfrac12,j+\tfrac12)^{2}$
where $\chi_{j+1}=1$. Then $(\chi_j\theta v)(f)=v(\zeta)=(\chi_{j+1}v)(\zeta)$, and the
product rule with $\|\partial^{\alpha}(f\circ\theta)\|_{L^{2}(\cC_j)}=\|\partial^{\alpha}f\|_{L^{2}(\T_j)}$
and $\|f\|_2,\|\partial_af\|_2\le\|f\|_{H^{2}(\T_j)}$ give
$|(\chi_j\theta v)(f)|\le C_jp_{j+1}(v)\|f\|_{H^{2}(\T_j)}$; taking the supremum over
trigonometric polynomials $f$ with $\|f\|_{H^{2}(\T_j)}\le1$ gives the bound. For
$\chi\in C_c^{\infty}$, $\chi\theta v=\theta((\chi\circ\theta)v)$ and the Fourier transform
of $\theta w$ is $\hat w\circ\theta$ with $|\theta p|=|p|$, so $\theta v\in X$. The identity
for $w\in E$ is the change of variables $z\mapsto\theta z$.

\emph{(viii)} For real $\chi\in C_c^{\infty}$, testing against $f$ gives
$\chi\cdot LG=(1-\Delta)(\chi G)+2\sum_i\partial_i(G\partial_i\chi)-G\Delta\chi$, so $LG\in X$.
With $\chi=\chi_j$, $(LG)(\chi_j\conj{e_{j,k}})=\lambda_{j,k}\langle e_{j,k},\chi_jG\rangle
+2i\tfrac{\pi}{j+1}\sum_ik_i\langle e_{j,k},G\partial_i\chi_j\rangle-\langle e_{j,k},G\Delta\chi_j\rangle$,
and $\lambda_{j,k}^{-2}(\pi|k_i|/(j+1))^{2}\le1$, Minkowski in $\ell^{2}$ and Parseval give
the bound, since all functions involved are supported in $[-j-\tfrac12,j+\tfrac12]^{2}$
of area $(2j+1)^{2}$. Continuity follows from dominated convergence for the
series defining $d$, with summable majorant $2^{-j}$.

\emph{(ix)} For real $f\in C_c^{\infty}$ choose $n$ with $\supp f\subset[-n,n]^{2}$; the
periodic extension $\tilde f$ of $f|_{\cC_n}$ is smooth, its Fourier series
converges in $C^{\infty}(\T_n)$, and truncation and rational rounding give
trigonometric polynomials $P_l$ with $\chi_nP_l\to\chi_n\tilde f=f$ in $C^{\infty}$ with
supports in $\supp\chi_n$; hence $\mathrm{Re}(\chi_nP_l)\to f$ in that sense and
$v(\mathrm{Re}(\chi_nP_l))\to v(f)$.

\emph{(x)} The sets in $\Sigma^{\theta}_{+}$ that are measurable with respect to the
evaluations at some countable subfamily of $C_c^{\infty}(H^{\theta}_{+};\R)$ form a
$\sigma$-algebra containing the generators, hence all of $\Sigma^{\theta}_{+}$. A
bounded $F$ is a uniform limit of simple functions; collecting countably many
tests gives $F=\tilde g\circ(v\mapsto(v(f_j))_j)$ by Doob--Dynkin, and
$g=\tilde g\ind_{\{|\tilde g|\le\sup|F|\}}$ still represents $F$. By (vii),
$F(\theta v)=g(((\theta v)(f_j))_j)=g((v(f_j\circ\theta))_j)$; both factors are Borel
by (iii).
\end{proof}

\begin{remark}[Use and scope of these facts]\label{rem:space-scope}
In (viii), $LG$ is the distribution $(1-\Delta)G$, not $\iota G$. It is used only
for the Gaussian versions $\Phi_a=L\tilde\psi_a$, where $\tilde\psi_a$ carries the
extra spectral multiplier $(1+|p|^{2})^{-1}$, so that $\Phi_a(f)=\phi_a(f)$
(Theorem~\ref{thm:G-versions}). Items (ii) (separability and completeness) and
(iii) (Borel generation) are proved here directly; they
re-express arguments of the earlier width record and of a Gaussian-foundations
lemma whose Borel-law part lies outside that record's accepted scope, and they
require current-version review in this text (Appendix~\ref{app:provenance}).
Completeness is used for Prokhorov compactness and for (v).
\end{remark}

% ===========================================================================
\section{Process neurons and finite-width central reflection positivity}\label{sec:endpoint}

Conditional factorization gives positivity for the interacting law at finite width. For product neurons, given the complete tuple of boundary data of all $N$ neurons, the two half-plane sides are independent with equal conditional laws (Theorem~\ref{thm:E-product}(d)); for Mondrian neurons, the trace decomposition of Theorem~\ref{thm:E-trace}, followed by approximation from finite boxes (Proposition~\ref{prop:E-sides}), plays this role. The reflection-invariant additive potential splits into half-plane contributions and is absorbed into the side observables (Lemma~\ref{lem:E-factor}). Theorem~\ref{thm:E-RP} states the result, and Corollary~\ref{cor:E-P2} applies it to the tilt of the main theorem.

This section states the properties of the neurons and the finite-width
positivity theorem; all proofs are in Appendix~\ref{app:E}. For open
$U\subseteq\R^{2}$, $L^{1}_{\mathrm{loc}}(U)$ is the Polish Fr\'echet space of classes of
locally integrable functions on $U$, and
$\Ball_r(U)=\{s\in L^{1}_{\mathrm{loc}}(U):|s|\le r\text{ a.e.}\}$, a closed subset.
For $\theta\in\{\theta_1,\theta_2\}$ write $H_+=H^{\theta}_+$. The \emph{one-neuron sides} of
a neuron $g$ across $\theta$ are $s^{\theta}_{+}=g|_{H_+}$ and $s^{\theta}_{-}=(g\circ\theta)|_{H_+}$;
the \emph{$N$-neuron sides} are $S^{\theta}_{+}=F_N|_{H_+}$ and
$S^{\theta}_{-}=(F_N\circ\theta)|_{H_+}$, random elements of $\Ball_{\sqrt N/2}(H_+)$. For a
product neuron the \emph{boundary data} are $b^{\theta_1}=(\epsilon_\sigma,\tau)$ and
$b^{\theta_2}=(\sigma,\epsilon_\tau)$, where $\epsilon_\sigma$, $\epsilon_\tau$ are the initial
signs of $\sigma$, $\tau$.

\begin{theorem}[Product neuron]\label{thm:E-product}
\begin{enumerate}[label=(\alph*)]
\item $\sigma$ is c\`adl\`ag with $|\sigma|=1$, $\E\sigma(t)=0$,
  $\E\sigma(t)\sigma(t')=e^{-\mu|t-t'|}$, and its finite-dimensional laws are those of
  the stationary symmetric two-state Markov chain with flip rate $\mu/2$.
\item $g$ is $\cB(\R^{2})\otimes\cA$-measurable, $|g|=1/2$ everywhere, $\E g(z)=0$,
  $\E g(z)g(z')=\Cend(z-z')$, and almost surely $g$ is constant on each cell of the
  locally finite grid formed by the jump points of $\sigma$ and of $\tau$.
\item Given $b^{\theta_1}$, the one-neuron sides across $\theta_1$ are conditionally
  independent with equal conditional laws; likewise across $\theta_2$ given
  $b^{\theta_2}$.
\item For independent product neurons $g_1,\dots,g_N$, the $N$-neuron sides
  $S^{\theta}_{\pm}$ are conditionally independent with equal conditional laws given
  the \emph{complete} tuple $(b^{\theta}_1,\dots,b^{\theta}_N)$. Conditioning on the data
  of a single neuron does not suffice when $N\ge2$.
\end{enumerate}
\end{theorem}

\begin{theorem}[Box Mondrian]\label{thm:E-box}
Let $C$ be a box with nonempty interior and $s\ge0$.
\begin{enumerate}[label=(\alph*)]
\item Almost surely finitely many nodes are born; the mean number of born nodes is
  at most $2e^{sp(C)}-1$; $\Pi^{\mathrm{can}}(C,s)\in\Part^{\mathrm{col}}(C)$ is a Borel
  function of $(C,s)$ and the ingredients.
\item $\Mcol(C,s)$ satisfies the first-step recursion of
  Definition~\ref{def:E-rec}, and every family of laws satisfying that recursion
  with an admissible daughter assignment coincides with $\Mcol$.
\item $T_*\Mcol(C,s)=\Mcol(TC,s)$ for translations $T$ and for $T\in\{\theta_1,\theta_2\}$.
\item $(\resB_B)_*\Mcol(C,s)=\Mcol(B,s)$ for every box $B\subseteq C$ with nonempty
  interior.
\item For $z,z'\in C$, the probability that some cell contains both is
  $e^{-s|z-z'|_1}$; a fixed $z\in\mathrm{int}\,C$ is almost surely interior to a cell.
\end{enumerate}
\end{theorem}

\begin{theorem}[Core-exterior construction and trace decomposition]\label{thm:E-trace}
For boxes $K\subseteq C$ the core-exterior construction of
Definition~\ref{def:E-coreext} produces almost surely a finite colored partition
of $C$ with law $\Mcol(C,s)$, whose restriction to a core $K$ with nonempty
interior is the canonical partition of $K$. For $C=[-L,L]\times J_0$ and
$K=\{0\}\times J_0$, the colored partition of the right half-box and the
$\theta_1$-mirror image of that of the left half-box are, given the trace data
(positions and auxiliary times of the horizontal cuts of straddling cells and
colors of the final straddling cells), conditionally independent with the same
conditional law; the same holds with coordinates exchanged. Times and positions
are both needed in this conditioning (Appendix~\ref{app:E}).
\end{theorem}

\begin{theorem}[Whole-plane Mondrian neuron]\label{thm:E-plane}
\begin{enumerate}[label=(\alph*)]
\item A consistent colored Mondrian tessellation of budget $\mu$ exists; the law
  of $(P_n)_n$ is unique and invariant under translations and under
  $\theta_1,\theta_2$ acting on closed cells; for a box $B$ with nonempty interior,
  $\resB_B(P_n)$ ($B\subseteq B_n$) does not depend on $n$ almost surely and has law
  $\Mcol(B,\mu)$. The lower-left representatives of the cells of any element of
  $\Part^{\mathrm{col}}(B_n)$ partition $[-n,n)^{2}$.
\item Almost surely each closed cell $R(z)$ is a bounded closed box with nonempty
  interior, equal to $R_n(z)$ for large $n$; the distinct closed cells cover
  $\R^{2}$, have disjoint interiors and are locally finite; $H(z)=R(z)^{\mathrm{ll}}$,
  and the distinct half-open representatives partition $\R^{2}$ and are locally
  finite.
\item $g$ is $\cB(\R^{2})\otimes\cA$-measurable, $|g|=1/2$ everywhere, almost surely
  constant on each $H(z)$, $\E g(z)=0$ and $\E g(z)g(z')=\Cend(z-z')$.
\item Every admissible convention agrees with $g$ off the union of cell
  boundaries, a null set for every outcome of the consistency event; for open
  $U$ the law of $g|_U$ in $L^{1}_{\mathrm{loc}}(U)$ is the same for every realization
  and every admissible convention; $g\circ\theta_j\overset{d}{=}g$ and
  $g(\cdot+x)\overset{d}{=}g$ in $L^{1}_{\mathrm{loc}}(\R^{2})$.
\end{enumerate}
\end{theorem}

For $b\in\{\bprod,\bMon\}$, $N\ge1$, a bounded Borel $\Omega$ and an
$\cL(\R^{2})\otimes\cB(\R)$-measurable $v:\R^{2}\times\R\to\R$, bounded on each
$\Omega\times[-r,r]$, put $U=\int_\Omega v(z,F_N(z))\,dz$, $Z=\E e^{-U}$,
$P^{v}_{N}(B)=\E[\ind_B(\iota F_N)e^{-U}]/Z$ and
$M_N=\sup\{|v(z,x)|:z\in\Omega,|x|\le\sqrt N/2\}$, with the convention
$\sup\emptyset=0$; if $|\Omega|=0$, then $U=0$ and $Z=1$.

\begin{theorem}[Finite-width annealed normalized reflection positivity]\label{thm:E-RP}
Let $b\in\{\bprod,\bMon\}$, $N\ge1$, $\theta\in\{\theta_1,\theta_2\}$, $\Omega$ bounded Borel
with $\theta\Omega=\Omega$, and $v$ as above with $v(\theta z,x)=v(z,x)$ for all $z\in\Omega$,
$x\in\R$. Then:
\begin{enumerate}[label=(\alph*)]
\item $U$ is a real random variable, $|U|\le|\Omega|M_N$ and
  $e^{-|\Omega|M_N}\le Z\le e^{|\Omega|M_N}$; $P^{v}_{N}$ is a Borel probability on $X$
  determined by the law of one neuron, independently of the realization and of the
  admissible boundary convention;
\item $P^{v}_{N}$ is reflection positive across $\theta$ for bounded observables,
  in the sense of \eqref{eq:boundedRP};
\item $\E[e^{-U}\conj{\Phi(S^{\theta}_{-})}\Phi(S^{\theta}_{+})]\ge0$ for every bounded complex
  Borel $\Phi$ on $L^{1}_{\mathrm{loc}}(H_+)$; and for every bounded complex $Y$
  measurable for $\sigma(F_N(z):z\in H_+)$, written $Y=\tilde Y((F_N(z_k))_{k\in\N})$,
  the reflected variable $Y^{\theta}=\tilde Y((F_N(\theta z_k))_k)$ is almost surely
  independent of the representation and $\E[e^{-U}\conj{Y^{\theta}}Y]\ge0$.
\end{enumerate}
\end{theorem}

\begin{corollary}[Two central lines; the axis-parallel manuscript tilt]\label{cor:E-P2}
Let $\Lambda=[-l_1,l_1]\times[-l_2,l_2]$ and $\Omega=[-w_1,w_1]\times[-w_2,w_2]$ with
$0\le l_i\le w_i$: closed axis-parallel rectangles, \emph{both centered}. Let
$q_{\mathrm{in}},q_{\mathrm{out}}$ be real polynomials, $\kappa\in\R$, and
$v(z,x)=\ind_\Lambda(z)q_{\mathrm{in}}(x)+\ind_{\Omega\setminus\Lambda}(z)q_{\mathrm{out}}(x)-\kappa^{2}x^{2}/2$.
Then for both architectures and every $N\ge1$ the conclusions (a)--(c) of
Theorem~\ref{thm:E-RP} hold across $\theta_1$ and across $\theta_2$.
\emph{Manuscript tilt.} If $\Lambda$ is a physical square in $\Omega=[-R,R]^{2}$ in
the sense of Definition~\ref{def:params}, the integrand $v_V$ of \eqref{eq:vV}
agrees off $\partial\Lambda\times\R$ with the displayed $v$ for the rectangles
$\bar\Lambda=[-l,l]^{2}$, $\Omega=[-R,R]^{2}$, with
$q_{\mathrm{in}}(x)=u(x^{4}-6c_ax^{2}+3c_a^{2})$, $q_{\mathrm{out}}(x)=\varepsilon x^{4}$,
$\kappa=\kappa_a$; for every outcome $U=V_{\varepsilon,\Omega,a;\Lambda}(F_N)$, and
$P^{v_V}_{N}=P^{b}_{N,\varepsilon,\Omega,a;\Lambda}$. Hence every neural law
\eqref{eq:neurallaw} is reflection positive across both central coordinate lines;
this covers $\Lambda_n=(-n\ell,n\ell)^{2}$ with $n\ell\le R$.
\end{corollary}

\begin{corollary}[General potentials invariant under both reflections]\label{cor:E-general}
If $\theta_1\Omega=\theta_2\Omega=\Omega$ and $v$ is $\cL(\R^{2})\otimes\cB(\R)$-measurable,
bounded on each $\Omega\times[-r,r]$ and invariant under both central reflections on
$\Omega$, then (a)--(c) of Theorem~\ref{thm:E-RP} hold across both central lines,
for both architectures and every $N\ge1$.
\end{corollary}

\begin{remark}[Geometry, measurability and scope]\label{rem:rotated}
(i) \emph{Centering alone is not enough.} For
$O=5^{-1/2}\bigl(\begin{smallmatrix}2&-1\\1&2\end{smallmatrix}\bigr)$, the centered
rotated square $O((-1,1)^{2})$ lies in $\Omega=[-2,2]^{2}$, and with $m=u=\varepsilon=1$,
$a=1/2$ the integrand built on it satisfies $v_V(z,0)=3uc_a^{2}>0=v_V(\theta_jz,0)$ for
both $j$ on a disc of positive area. This shows only that the invariance
hypothesis of Theorem~\ref{thm:E-RP} fails for that tilt; no failure of reflection
positivity is claimed for it or for any rotated square.
(ii) \emph{The general clause keeps its own hypotheses.} A centered square rotated
by an odd multiple of $\pi/4$, $\{|z_1|+|z_2|<c\}$, is invariant under both
central reflections; a potential of the form \eqref{eq:vV} built on it, with
closure inside an invariant $\Omega$, satisfies the hypotheses of
Corollary~\ref{cor:E-general}, which therefore applies to it, although it is not
a physical square in the sense of Definition~\ref{def:params}.
(iii) \emph{Measurability.} ``Measurable potential'' means
$\cL(\R^{2})\otimes\cB(\R)$-measurable. Under the broader reading ``Lebesgue
measurable on $\R^{3}$'' the integrated potential can be undefined: for a
non-measurable $A\subset(0,1)^{2}$ symmetrized under both reflections and
$v(z,x)=\ind_{A}(z)\ind\{x=x_0\}$ with $x_0\in\{0,1/(2\sqrt N)\}$ of the parity of
$N$, $v$ is Lebesgue-null on $\R^{3}$, but on an event of positive probability
$F_N=x_0$ on $\Omega$ and $z\mapsto v(z,F_N(z))$ is not measurable.
(iv) \emph{Scope.} Positivity is asserted across $\theta$ only when $\Omega$ and $v$ are
$\theta$-invariant; for the two-line statements, invariance under both reflections
of both the inner and the outer region is used. Nothing is asserted for
shifted or rotated lines, for Euclidean invariance of tilted laws, or for
limits in $N$, $\varepsilon$, $\Omega$ or $a$ by this theorem.
\end{remark}

% ===========================================================================
\section{Gaussian reference fields}\label{sec:gauss}

Gaussian fields enter in two roles. First, they serve as references for the intermediate laws of the chain: the width limit replaces the untilted neural reference field by the continuous endpoint field $\gamma_\mu$ (Theorem~\ref{thm:G-endpoint}); its quadratic tilt on a bounded $\Omega$ is again Gaussian, with a covariance $C'_{\Omega,a}$ between $\Cend$ and $C_a$ (Theorems~\ref{thm:G-tilt} and~\ref{thm:G-order}); and $C'_{\Omega,a}\to C_a$ locally uniformly as $\Omega$ grows (Theorem~\ref{thm:G-diff}). Second, Theorems~\ref{thm:G-common}--\ref{thm:G-versions} realize all regulated fields and the free field on one probability space, so that the regulator can be removed through $L^{2}$ estimates for random variables.

Proofs are in Appendix~\ref{app:G}. Throughout, $a$ is a regulator unless $a=0$
is stated, and $\Omega$ is a bounded measurable set.

\begin{theorem}[Endpoint Gaussian]\label{thm:G-endpoint}
For every $\mu>0$ there is a $\cB(\R^{2})\otimes\cF$-measurable
$\gamma_\mu:\R^{2}\times\Gamma_\mu\to\R$ with continuous paths for every outcome, which
is a centered Gaussian process with covariance $\Cend(z-z')$ and
$\E|\gamma_\mu(z+h)-\gamma_\mu(z)|^{2}=(1-e^{-\mu|h|_1})/2\le\mu|h|_1/2$. Its law $G_\mu$ on
$C(\R^{2})$ is the unique Borel probability under which the coordinate process is
centered Gaussian with this covariance.
\end{theorem}

\begin{theorem}[Regulated covariance and mass shift]\label{thm:G-shift}
\begin{enumerate}[label=(\alph*)]
\item $\hat C_a\le\hat C_0\le m^{-2}$,
  $\hat C_a(w,q)\le m^{-2}(1+a^{2}w^{2})^{-1}(1+a^{2}q^{2})^{-1}$, $\hat C_a\in L^{1}$, and
  \eqref{eq:ca} holds with $c_a\in(0,\infty)$.
\item \eqref{eq:massshift} holds pointwise, with $0\le\kappa_a^{2}\hCend\le r_a$.
\item $s_k=\kappa_a^{2k}(\Cend)^{*(k+1)}$ is continuous and positive with
  $\sup s_k\le r_a^{k}/4$ and $\|s_k\|_1=\mu^{-2}r_a^{k}$; $\sum_ks_k$ converges uniformly
  and in $L^{1}$ to $C_a$. Hence $C_a$ is continuous, $0<\Cend\le C_a\le c_a$,
  $\|C_a\|_1=m^{-2}$ and $\|C_a\|_2^{2}\le1/(4\pi m^{2})$.
\item As operators on $L^{2}(\R^{2})$: $\|\Cend\|=\mu^{-2}$,
  $C_a=\Cend(I-\kappa_a^{2}\Cend)^{-1}$, and $0\le\langle f,C_af\rangle\le m^{-2}\|f\|_2^{2}$.
\item For every $h\in\R^{2}$: $2(c_a-C_a(h))\le\varpi_a(h)^{2}
  :=(2\pi m^{2}a^{3})^{-1}\sum_{i=1,2}|h_i|[\log(1+4a^{2}/h_i^{2})+2]$ (summand $0$ if $h_i=0$).
\item For $0<\alpha<1$: $2(c_a-C_a(h))\le M'_{a,\alpha}|h|^{\alpha}$ for all $h$, with
  $M_{a,\alpha}=[2(2a)^{1-\alpha}/(1-\alpha)+2]/(\pi m^{2}a^{3})$ and
  $M'_{a,\alpha}=\max(M_{a,\alpha},2c_a)$.
\end{enumerate}
\end{theorem}

\begin{theorem}[Restricted endpoint covariance]\label{thm:G-K}
$K_\Omega f(x)=\int_\Omega\Cend(x-y)f(y)\,dy$ on $L^{2}(\Omega;\R)$ is positive, self-adjoint
and trace class with $\trace K_\Omega=|\Omega|/4$, $\|K_\Omega\|\le\mu^{-2}$ and
$\kappa_a^{2}\|K_\Omega\|\le r_a<1$; it has an orthonormal eigenbasis with eigenvalues
$\lambda_j\in[0,\mu^{-2}]$.
\end{theorem}

\begin{theorem}[Pure quadratic tilt on finite $\Omega$]\label{thm:G-tilt}
\begin{enumerate}[label=(\alph*)]
\item $\Det_{\Omega,a}=\E_{G_\mu}e^{T_{\Omega,a}}=\prod_j(1-\kappa_a^{2}\lambda_j)^{-1/2}
  =\det_F(I-\kappa_a^{2}K_\Omega)^{-1/2}$, with
  $1\le\Det_{\Omega,a}\le\exp(\kappa_a^{2}|\Omega|/(8a^{2}m^{2}))$, nondecreasing in $\Omega$.
\item $c^{\Omega}_0(z,z')=\Cend(z-z')$ and
  $c^{\Omega}_k(z,z')=\kappa_a^{2k}\int_{\Omega^{k}}\Cend(z-x_1)\cdots\Cend(x_k-z')\,dx$ are
  continuous, symmetric and $0\le c^{\Omega}_k\le s_k(z-z')$; $C'_{\Omega,a}=\sum_kc^{\Omega}_k$
  converges uniformly on $\R^{2}\times\R^{2}$, and the operator with this kernel is
  $\Cend(I-\kappa_a^{2}\ind_\Omega\Cend)^{-1}=\Cend+\kappa_a^{2}\Cend R^{*}(I-\kappa_a^{2}K_\Omega)^{-1}R\,\Cend$.
\item $G'_{\Omega,a}=\Det_{\Omega,a}^{-1}e^{T_{\Omega,a}}G_\mu$ is the unique Borel probability
  on $C(\R^{2})$ under which the coordinate process at \emph{all} points of the plane
  is centered Gaussian with covariance $C'_{\Omega,a}$, and
  $\E_{G_\mu}[Fe^{T_{\Omega,a}}]=\Det_{\Omega,a}\E_{G'_{\Omega,a}}[F]$ for Borel $F\ge0$ and
  for Borel $F$ integrable under $G'_{\Omega,a}$.
\end{enumerate}
\end{theorem}

\begin{theorem}[Covariance order]\label{thm:G-order}
$0<\Cend(z-z')\le C'_{\Omega,a}(z,z')\le C_a(z-z')\le c_a$ pointwise;
$0\le\Cend\le C'_{\Omega,a}\le C_a$ as operators; $C'_{\Omega,a}\le C'_{\Omega',a}$ for
$\Omega\subseteq\Omega'$.
\end{theorem}

\begin{theorem}[The covariance difference]\label{thm:G-diff}
Let $\Delta_{\Omega,a}=C_a-C'_{\Omega,a}$ and $Q=\ind_{\R^{2}\setminus\Omega}$.
\begin{enumerate}[label=(\alph*)]
\item $\Delta=\kappa_a^{2}C_aQC'=\kappa_a^{2}C'QC_a=\kappa_a^{2}C'QC'+\kappa_a^{4}C'QC_aQC'$ as
  operators, with the corresponding kernel identities.
\item $\Delta$ is continuous and positive semidefinite as an operator and as a
  kernel.
\item $0\le\Delta(z,z')\le\kappa_a^{2}\int_{\R^{2}\setminus\Omega}C_a(z-x)C_a(x-z')\,dx$; if $\Omega$
  contains every point at distance $<r$ from a compact $K$, then
  $0\le\Delta(z,z')\le\epsilon_a(r)=\kappa_a^{2}\int_{|x|>r}C_a(x)^{2}dx$ for $z,z'\in K$, and
  $\epsilon_a(r)\downarrow0$.
\item $\sup_{K\times K}\Delta_{\Omega,a}\to0$ along every family of bounded $\Omega$ that
  eventually contains each bounded set, for instance $[-R,R]^{2}$, $R\to\infty$.
\item $\Delta(z+h,z+h)+\Delta(z,z)-2\Delta(z+h,z)\le2(c_a-C_a(h))$, and the same for $C'$;
  centered Gaussian processes with covariance $\Delta_{\Omega,a}$ or $C'_{\Omega,a}$ exist and
  have jointly measurable modifications with continuous paths.
\end{enumerate}
\end{theorem}

The spectral Gaussian space: $\mu_0(dp)=(2\pi)^{-2}\hat C_0(p)\,dp$,
$H_{\C}=\{h\in L^{2}(\mu_0;\C):h(-p)=\conj{h(p)}\}$ with the real inner product
$\langle g,h\rangle=\int\conj gh\,d\mu_0$, $\tau_a=(\hat C_a/\hat C_0)^{1/2}$ ($\tau_0=1$),
$e_z(p)=e^{-ip\cdot z}$ and $\rho(p)=(1+|p|^{2})^{-1}$. An isonormal process on a real
Hilbert space is a real-linear isometry into $L^{2}$ whose range consists of
centered jointly Gaussian variables.

\begin{theorem}[One common Gaussian space]\label{thm:G-common}
There is a probability space $(\Gamma_0,\cF_0,\PP_0)$ carrying an isonormal process
$W$ on $H_{\C}$. With $\phi_a(f)=W(\hat f\tau_a)$ for $0\le a<1/m$ and real
$f\in L^{2}(\R^{2})$, $\phi=\phi_0$, the whole family is centered and jointly Gaussian,
$f\mapsto\phi_a(f)$ is real-linear, and
$\E[\phi_a(f)\phi_b(g)]=(2\pi)^{-2}\int\conj{\hat f}\hat g(\hat C_a\hat C_b)^{1/2}dp$.
\end{theorem}

\begin{theorem}[Smearing bounds]\label{thm:G-bounds}
For real $f\in L^{2}$: $\|\phi_a(f)\|_2\le\|\phi(f)\|_2\le m^{-1}\|f\|_2$ and
$\|\phi_a(f)-\phi(f)\|_2\le\tfrac a2\|f\|_2$. For complex $f$ and
$T_a(f)=\phi_a(\mathrm{Re}f)+i\phi_a(\mathrm{Im}f)$: $\E|T_a(f)|^{2}\le m^{-2}\|f\|_2^{2}$ and
$\E|T_a(f)-T_0(f)|^{2}\le\tfrac{a^{2}}4\|f\|_2^{2}$.
\end{theorem}

\begin{theorem}[Point fields and continuous modifications]\label{thm:G-point}
For a regulator $a$, $\phi_a(z)=W(e_z\tau_a)$ is a centered Gaussian process with
$\E[\phi_a(z)\phi_b(z')]=C_{ab}(z-z')$, $C_{ab}=\cI[(\hat C_a\hat C_b)^{1/2}]$, and
$\E|\phi_a(z+h)-\phi_a(z)|^{2}=2(c_a-C_a(h))\le\varpi_a(h)^{2}$. There is a
$\cB(\R^{2})\otimes\cF_0$-measurable $\tilde\phi_a$ with continuous paths for every
outcome and $\PP_0(\tilde\phi_a(z)=\phi_a(z))=1$ for every $z$; for real
$f\in L^{1}\cap L^{2}$, $\int|f||\tilde\phi_a|<\infty$ and $\int f\tilde\phi_a=\phi_a(f)$ almost
surely. The exceptional null sets depend on $(a,z)$ and on $(a,f)$; for countable
families of regulators, points and tests there is one common null set; no common
null set for uncountably many $a$, and no measurability or continuity in $a$, is
asserted.
\end{theorem}

\begin{theorem}[Borel $X$-valued versions, including $a=0$]\label{thm:G-versions}
For each $0\le a<1/m$, $\psi_a(z)=W(e_z\tau_a\rho)$ has a modification
$\tilde\psi_a$ with continuous paths, and $\Phi_a:=L(\tilde\psi_a(\cdot,\omega))$ (the map
$L$ of Proposition~\ref{prop:space}(viii)) is a Borel map $\Gamma_0\to X$ with:
\begin{enumerate}[label=(\alph*)]
\item $\Phi_a(f)=\phi_a(f)$ almost surely for every real $f\in C_c^{\infty}(\R^{2})$, and
  $\Phi_a(f)=T_a(f)$ for complex $f$;
\item $\Phi_a=\iota\tilde\phi_a$ outside a null set depending on $a$, for $a>0$;
\item if $\Phi':\Gamma_0\to X$ satisfies $\Phi'(f)=\phi_a(f)$ almost surely for every
  $f\in\cD_0$ (Proposition~\ref{prop:space}(ix)), then $\Phi'=\Phi_a$ almost surely;
\item for countable families of regulators and tests the identities in (a), (b)
  hold off one null set;
\item $\E_0p_j(\Phi_a)^{2}\le m^{-2}B_j$ and, for $a>0$, $\E_0p_j(\Phi_a-\Phi_0)^{2}\le\tfrac{a^{2}}4B_j$,
  with $B_j=(\|\chi_j\|_2^{2}/|\T_j|)\sum_{k}\lambda_{j,k}^{-2}<\infty$; hence
  $\|d(\Phi_a,\Phi_0)\|_{L^{2}(\PP_0)}\le h(a):=\sum_j2^{-j}\min(1,\tfrac a2B_j^{1/2})\to0$;
\item for real $f_1,\dots,f_n\in C_c^{\infty}$ the vector $(\Phi_a(f_i))_i$ is centered
  Gaussian with covariance $(2\pi)^{-2}\int\conj{\hat f_i}\hat f_l\hat C_a\,dp$.
\end{enumerate}
No interaction functional, Wick power, weight or normalizer enters these
constructions.
\end{theorem}

\begin{corollary}[Laws of the versions on $X$]\label{cor:G-laws}
For $0\le a<1/m$ the Borel law of $\Phi_a$ on $X$ is the unique Borel probability
whose finite-dimensional laws on compact smooth tests are those of (f); at $a=0$
it is the centered Gaussian law with covariance $\Cfree$.
\end{corollary}

\begin{proof}
By Proposition~\ref{prop:space}(iii) the finite-dimensional cylinder events of
compact smooth evaluations form a $\cap$-stable generator of $\cB(X)$; apply
Dynkin's theorem.
\end{proof}

\begin{remark}[Scope of the Gaussian foundations]\label{rem:G-scope}
Theorems~\ref{thm:G-endpoint}--\ref{thm:G-versions} use only the Gaussian space;
they assert nothing about interaction weights, Wick powers, normalizers,
interacting laws or reflection positivity. Corollary~\ref{cor:G-laws} uses the
Borel generation of Proposition~\ref{prop:space}(iii), which is proved in this
paper. The interaction objects on the same space are supplied by
Theorem~\ref{thm:W} below. Theorem~\ref{thm:G-shift}(e) is proved for every
regulator $0<a<1/m$; nothing is claimed at $a=1/m$.
\end{remark}

% ===========================================================================
\section{Fixed-volume Wick theory}\label{sec:wick}

This section asks whether the interaction has a limit when the regulator is removed. The free field is distribution-valued, so its fourth power has no pointwise meaning, and the Wick constant $c_a$ diverges; nevertheless the Wick-ordered integrals $U_{a,\Lambda}$ converge in $L^{2}$ to a fourth-chaos variable $U_{0,\Lambda}$ (Theorem~\ref{thm:W}(a),(b)). The key estimate is the bound $\E_0e^{-pU_a}\le\Xi_p$, uniform in the regulator (Theorem~\ref{thm:W}(e)), which controls the weights and normalizers. These weight and regulator estimates have volume-dependent constants (Remark~\ref{rem:W-scope}); the volume-uniform moment bounds needed later come from classical input (Section~\ref{sec:target}).

Proofs are in Appendix~\ref{app:W}. Fix $m,u>0$, a physical square $\Lambda$ with
half-side $l>0$ (open or closed; the boundary is null), and $A=|\Lambda|=4l^{2}$.
There is no weak-coupling restriction. For $c\ge0$ let
${:}x^{4}{:}_c=x^{4}-6cx^{2}+3c^{2}$; for a centered Gaussian variable $\xi$ put
${:}\xi^{4}{:}={:}\xi^{4}{:}_{\Var\xi}$ (own-variance Wick power; zero if $\Var\xi=0$).
Let $H_1=W(H_{\C})$ be the Gaussian Hilbert space of Theorem~\ref{thm:G-common},
$H^{{:}4{:}}$ its fourth Wiener chaos, $\cN$ the $\PP_0$-null sets and
$\cF_{\mathrm{free}}=\sigma(\phi(f):f\in C_c^{\infty}(\R^{2};\R))\vee\cN$.

\begin{definition}[Interactions and admissible cutoffs]\label{def:W-objects}
For a regulator $a$, $U_a=U_{a,\Lambda}=u\int_\Lambda{:}\tilde\phi_a(z)^{4}{:}_{c_a}\,dz$
pathwise, and $U_0=L^{2}(\PP_0)\text{-}\lim_{a\downarrow0}U_a$ through all real regulators
(existence is part of Theorem~\ref{thm:W}). An \emph{admissible cutoff family} is a
family $(\chi_\epsilon)_{\epsilon>0}$ of measurable functions with
$\chi_\epsilon(-p)=\conj{\chi_\epsilon(p)}$, $\sup_\epsilon\|\chi_\epsilon\|_\infty<\infty$,
$\chi_\epsilon\to1$ almost everywhere and $\int|\chi_\epsilon|^{2}\hat C_0<\infty$; put
$\phi^{\chi}_\epsilon(z)=W(e_z\chi_\epsilon)$ and $N^{\chi}_\epsilon=u\int_\Lambda{:}\phi^{\chi}_\epsilon(z)^{4}{:}\,dz$
(Bochner integral in $L^{2}(\PP_0)$). Examples: smooth mollifiers
$\chi_\epsilon(p)=\conj{\hat\eta(\epsilon p)}$ with real $\eta\in C_c^{\infty}$, $\int\eta=1$, for
which $\phi^{\chi}_\epsilon(z)=\phi(\eta_\epsilon(z-\cdot))$; and sharp cutoffs
$\ind_{\{|p|\le1/\epsilon\}}$. The constants are
$\kappa_0=e(1+\sqrt2)$, $\kappa_2=(2\pi)^{-1}(17\pi/16)^{1/2}$,
$M_6=(2\pi)^{-5/3}(5\pi)^{5/6}m^{-1/3}$, $K_\Lambda=u(384A\kappa_2M_6^{3})^{1/2}$,
$K_\sharp=\max\{12uA,(2/e)K_\Lambda(\kappa_0/m)^{1/2}\}$,
$T(K)=\exp(-\exp((\pi/2)(K/(12uA))^{1/2}-1))$ and
$\Xi_p=e^{pK_\sharp}+\int_{K_\sharp}^{\infty}pe^{pK}T(K)\,dK$.
\end{definition}

For $\theta\in\{\theta_1,\theta_2\}$ let $\Theta$ be the reflection of the Gaussian space
induced by $h\mapsto h\circ\theta$ on $H_{\C}$ (Appendix~\ref{app:W}): it is linear,
multiplicative and law-preserving, with $\Theta\phi_a(f)=\phi_a(f\circ\theta)$ and
$\Theta\phi_a(z)=\phi_a(\theta z)$. Put
$\cA^{\theta,a}_{+}=\sigma(\phi_a(f):f\in C_c^{\infty}(H^{\theta}_{+};\R))\vee\cN$.

\begin{theorem}[Fixed-volume Wick theory]\label{thm:W}
On the common Gaussian space of Theorem~\ref{thm:G-common}:
\begin{enumerate}[label=(\alph*)]
\item $U_a\in L^{2}(\PP_0)$ lies in $H^{{:}4{:}}$, equals the Bochner integral of
  $z\mapsto u{:}\phi_a(z)^{4}{:}$, and satisfies $U_a\ge-6uAc_a^{2}$ pathwise; the limit
  $U_0$ exists and lies in $H^{{:}4{:}}$.
\item \emph{Construction statement.} For every admissible cutoff family,
  $N^{\chi}_\epsilon\to U_0$ in $L^{2}(\PP_0)$; $U_0$ is the unique $Y\in H^{{:}4{:}}$ with
  $\E_0[Y{:}\phi(f_1)\cdots\phi(f_4){:}]=4!\,u\int_\Lambda\prod_{j\le4}(\Cfree*f_j)(z)\,dz$ for
  all real $f_j\in C_c^{\infty}$; and $U_0$ is $\cF_{\mathrm{free}}$-measurable. The
  construction uses only the Gaussian variables $W(h)$, $h\in H_{\C}$, their
  own-variance Wick polynomials, Bochner integrals in $L^{2}(\PP_0)$ and
  $L^{2}(\PP_0)$-limits; it uses no point variable at $a=0$ and forms no pointwise
  product of the distribution-valued free field. Nothing is asserted about other
  representations of $U_0$.
\item $\E_0U_a=0$ for $0\le a<1/m$.
\item $\|U_a-U_0\|_2\le K_\Lambda\sqrt a$ for every regulator $a$.
\item For every real $p\ge1$, $\Xi_p<\infty$ and $\E_0e^{-pU_a}\le\Xi_p$ for $0\le a<1/m$,
  including $a=0$.
\item For real $p\ge1$ and regulators $a$,
  $\|e^{-U_a}-e^{-U_0}\|_p\le2(2p-1)^{2}K_\Lambda\Xi_{2p}^{1/(2p)}\sqrt a$.
\item $1\le Z_a:=\E_0e^{-U_a}\le\Xi_1$ for $0\le a<1/m$.
\item For $\theta\in\{\theta_1,\theta_2\}$, a regulator $a$ and $F\in L^{2}(\cA^{\theta,a}_{+})$:
  $\E_0[\conj{\Theta F}Fe^{-U_a}]\ge0$.
\item For $\theta\in\{\theta_1,\theta_2\}$ and bounded $\cA^{\theta,0}_{+}$-measurable $F$:
  $\E_0[\conj{\Theta F}Fe^{-U_0}]\ge0$.
\item For $0\le a<1/m$ and any measurable $\Phi':\Gamma_0\to X$ with $\Phi'(f)=\phi_a(f)$
  almost surely for real $f\in C_c^{\infty}$: $B\mapsto\E_0[\ind_B(\Phi')e^{-U_a}]/Z_a$ is
  a Borel probability on $X$, reflection positive across $\theta_1$ and $\theta_2$
  for bounded observables.
\item No weak-coupling restriction is used; no statement is uniform in $\Lambda$.
\end{enumerate}
Moreover, for $K\ge K_\sharp$: $\sup_{0<a<1/m}\PP_0(U_a\le-K)\le T(K)$ and
$\PP_0(U_0<-K)\le T(K)$.
\end{theorem}

\begin{corollary}[The physical laws]\label{cor:W-laws}
For $0\le a<1/m$, $\nu_{a,\Lambda}$ of Definition~\ref{def:laws}(d), built with
$\Phi_a=L\tilde\psi_a$ of Theorem~\ref{thm:G-versions}, is a Borel probability on $X$
with $1\le Z_{a,\Lambda}\le\Xi_1(\Lambda)$; it is reflection positive across both central
coordinate lines for bounded observables; and any Borel $\Phi'$ satisfying the
identity of Theorem~\ref{thm:G-versions}(a) gives the same law. At $a=0$ this is
the finite-volume Wick-quartic law on $X$; no interface hypothesis remains.
\end{corollary}

\begin{proof}
Theorem~\ref{thm:G-versions}(a) supplies the hypothesis of Theorem~\ref{thm:W}(j)
for every $0\le a<1/m$, including $a=0$. Uniqueness: Theorem~\ref{thm:G-versions}(c)
gives $\Phi'=\Phi_a$ almost surely.
\end{proof}

\begin{remark}[Scope of the Wick theory]\label{rem:W-scope}
Every constant depends on $\Lambda$ through $A=|\Lambda|$: $K_\Lambda\propto uA^{1/2}m^{-1/2}$,
$K_\sharp\ge12uA$ and $\Xi_p\ge e^{12puA}$. The only volume-uniform statement is
$Z_{a,\Lambda}\ge1$. The $\sqrt a$ bounds are fixed-volume upper bounds, not practical
convergence rates. The Gaussian input is used only through
Theorems~\ref{thm:G-common}, \ref{thm:G-point} and \ref{thm:G-versions}; the weight,
the normalizer and reflection positivity enter only after that input is fixed.
Exceptional null sets are those of Theorems~\ref{thm:G-point} and
\ref{thm:G-versions}: one per fixed regulator value and test, one for countable
families, none for uncountably many regulators. The classical inputs
\ref{prem:C6}(a)--(c) are used in their stated forms; their original sources were
not inspected.
\end{remark}

% ===========================================================================
\section{The fixed-volume chain}\label{sec:fixed}

Throughout this section a physical square $\Lambda$ (Definition~\ref{def:params}) is
fixed; only the width, the stabilizer, the auxiliary square and the regulator
move, one at a time, with all other parameters fixed. Every statement holds for
all $m,u>0$; no weak-coupling restriction is used, and no constant is uniform in
the parameters held fixed. The inputs are the theorems of
Sections~\ref{sec:space}--\ref{sec:wick}. Appendix~\ref{app:provenance} records the argument lineage.

\begin{lemma}[Open and closed physical squares]\label{lem:openclosed}
The neural, Gaussian, residual and physical laws built with $\Lambda=(-l,l)^{2}$
coincide with those built with $\bar\Lambda=[-l,l]^{2}$.
\end{lemma}

\begin{proof}
For $v\in E$ the integrals of $v^{4}-6c_av^{2}+3c_a^{2}$ and of $v^{4}$ over $\Lambda$ and
$\bar\Lambda$ (and over the corresponding parts of $\Omega$) coincide because
$\partial\Lambda$ is Lebesgue-null; $T_{\Omega,a}$ does not involve $\Lambda$. Hence every
ordinary-field law coincides. On the common Gaussian space $U_{a,\Lambda}=U_{a,\bar\Lambda}$
for every regulator, so their $L^{2}(\PP_0)$ limits coincide.
\end{proof}

The four arrows below remove the width, the stabilizer, the auxiliary square and the regulator, in that order. The width arrow combines a central limit theorem for finitely many smearings with compactness in $E$, obtained from the increment bounds \eqref{eq:incr4}; the bounded continuous weight $e^{-V}$ then passes to the limit. The stabilizer arrow is a monotone limit of weights, the auxiliary-square arrow couples two Gaussian fields whose covariances differ by $\Delta_{\Omega,a}$, and the regulator arrow uses the estimates of Theorem~\ref{thm:W}. Remark~\ref{rem:chain} summarizes the resulting chain.

\subsection{The width arrow}

\begin{proposition}[Width arrow]\label{prop:width}
Fix $b\in\{\bprod,\bMon\}$, a regulator $a$, $\varepsilon>0$, a physical square $\Lambda$ and
$\Omega=[-R,R]^{2}\supseteq\bar\Lambda$. Then $Z^{b}_{N}$ and $Z_{\mathrm{full}}$ are positive and
finite, $Z^{b}_{N}\to Z_{\mathrm{full}}$, and
$P^{b}_{N,\varepsilon,\Omega,a;\Lambda}\Rightarrow\nu^{v}_{\varepsilon,\Omega,a;\Lambda}$ weakly on $X$ as
$N\to\infty$. The limit is the same full Gibbs tilt of the Gaussian reference for
both architectures.
\end{proposition}

\begin{proof}
Write $F_N=F^{b}_{N,a}$, $Q_N$ for its law on $E$ and $Q$ for the law of $\gamma_\mu$ on
$E$; both are Borel by Proposition~\ref{prop:space}(i), since
$|F_N|\le\sqrt N/2$ and $\gamma_\mu$ is continuous and jointly measurable
(Theorems~\ref{thm:E-product}, \ref{thm:E-plane}, \ref{thm:G-endpoint}).

\emph{Moments.} By Theorems~\ref{thm:E-product}(b) and \ref{thm:E-plane}(c), each
neuron has $|g|=1/2$ everywhere, $\E g=0$ and $\E g(z)g(z')=\Cend(z-z')$. For
$r=|z-z'|_1$ put $b(r)=\tfrac12(1-e^{-\mu r})\le\mu r/2$. Then
$d_i=g_i(z)-g_i(z')\in\{0,\pm1\}$, $\E d_i=0$ and $\E d_i^{4}=\E d_i^{2}=b(r)$; independence
across neurons gives
\begin{equation}\label{eq:incr4}
\E|F_N(z)-F_N(z')|^{4}=\frac{b(r)}N+3\Bigl(1-\frac1N\Bigr)b(r)^{2}\le\frac{\mu r}2+\frac{3\mu^{2}r^{2}}4,
\qquad \E F_N(z)^{4}=\frac{3-2/N}{16}\le\frac3{16}.
\end{equation}
For $\gamma_\mu$ the increment has variance $b(r)$ and fourth moment $3b(r)^{2}$, and
$\E\gamma_\mu(z)^{4}=3/16$, so the same bounds hold.

\emph{Compact containment in $E$.} For a box $B$ and the cell-average projection
$P_\delta$ onto a partition of $B$ into squares of $\ell^{1}$-diameter at most $\delta$,
Jensen on each cell $C$ gives
$\int_C|v-|C|^{-1}\!\int_Cv|^{4}\le|C|^{-1}\int_C\int_C|v(x)-v(y)|^{4}$, so by
\eqref{eq:incr4} and Tonelli
$\E\|v-P_\delta v\|^{4}_{L^{4}(B)}\le|B|(\mu\delta/2+3\mu^{2}\delta^{2}/4)$ and
$\E\|v\|^{4}_{L^{4}(B)}\le3|B|/16$ under every $Q_N$ and under $Q$. Fix $\rho>0$. For
$B_j=(-j,j)^{2}$ choose $R_j$ with Markov tail $\PP(\|v\|_{L^{4}(B_j)}>R_j)\le\rho2^{-j-1}$, and
for $l\ge1$ a dyadic projection $P_{j,l}$ with
$\PP(\|v-P_{j,l}v\|_{L^{4}(B_j)}>2^{-l})\le\rho2^{-j-l-1}$, uniformly in $N$. The set
$K_j=\{v:\|v\|_4\le R_j,\ \|v-P_{j,l}v\|_4\le2^{-l}\ \forall l\}\subset L^{4}(B_j)$ is closed and
totally bounded (each $P_{j,l}K_j$ is bounded in a finite-dimensional space), hence
compact; $K=\{v\in E:v|_{B_j}\in K_j\ \forall j\}$ is compact in $E$, and
$\inf_NQ_N(K)\ge1-\rho$, $Q(K)\ge1-\rho$.

\emph{Finite vectors.} For real $f_1,\dots,f_r\in C_c^{\infty}$, the vectors
$Y_i=(\int g_if_1,\dots,\int g_if_r)$ are independent, identically distributed, bounded
and centered, with covariance $\Sigma_{kl}=\iint f_k(z)\Cend(z-z')f_l(z')\,dz\,dz'$
(Fubini). A Taylor expansion of $\E e^{it\cdot Y_i/\sqrt N}$ and L\'evy's continuity
theorem give $(F_N(f_1),\dots,F_N(f_r))\Rightarrow(\gamma_\mu(f_1),\dots,\gamma_\mu(f_r))$.

\emph{Identification on $E$.} Let $\cD_E$ be a countable family of real compact
smooth tests, dense in each $\cD_K$; the evaluations $v\mapsto\int vd$, $d\in\cD_E$, are
continuous on $E$ and separate its points. Let $H\in C_b(E)$, $\eta>0$, and $K$ as
above. Around each $x\in K$ choose a relatively open $U_x\ni x$ in $K$ on which
$|H-H(x)|<\eta$. Some finite-coordinate cylinder
$\{v\in K:|\int(v-x)d_i|<\epsilon_x,\ i\le r_x\}$ lies in $U_x$; otherwise a sequence in the
closed set $K\setminus U_x$ would agree ever more closely with $x$ on the first $k$
coordinates, and a convergent subsequence would have the same coordinates as
$x$, contradicting separation. Cover $K$ by finitely many half-size cylinders,
let $\pi$ be the finite vector of all coordinates used, choose continuous
$\alpha_i\ge0$ on its range space supported in the full cylinders and positive on
the half-size ones, put $S=\sum\alpha_i$, $A=\pi(K)$ and
$G(y)=\sum_iH(x_i)\alpha_i(y)/(S(y)+\dist(y,A))$. Then $G$ is continuous,
$\|G\|_\infty\le\|H\|_\infty$ and $\sup_K|H-G\circ\pi|\le\eta$. Hence
$|Q_N(H)-Q(H)|\le|Q_N(G\circ\pi)-Q(G\circ\pi)|+2\eta+4\|H\|_\infty\rho$, and the finite-vector
convergence gives $Q_N\Rightarrow Q$ on $E$.

\emph{Full weight.} Completing squares in $x^{2}$ gives
$ux^{4}-b_ax^{2}+3uc_a^{2}\ge3uc_a^{2}-b_a^{2}/(4u)$ with $b_a=6uc_a+\kappa_a^{2}/2$ on $\Lambda$, and
$\varepsilon x^{4}-\kappa_a^{2}x^{2}/2\ge-\kappa_a^{4}/(16\varepsilon)$ on $\Omega\setminus\Lambda$, so
\begin{equation}\label{eq:Mfull}
0<e^{-V_{\varepsilon,\Omega,a;\Lambda}}\le M_{\mathrm{full}}
=\exp\bigl\{|\Lambda|[b_a^{2}/(4u)-3uc_a^{2}]+|\Omega\setminus\Lambda|\kappa_a^{4}/(16\varepsilon)\bigr\}.
\end{equation}
$V$ is continuous on $E$ (Proposition~\ref{prop:space}(vi)) and finite on every
element. Applying $Q_N\Rightarrow Q$ to the bounded continuous functions $e^{-V}$ and
$He^{-V}$ gives $Z^{b}_{N}\to Z_{\mathrm{full}}\in(0,\infty)$ and convergence of the normalized
tilts on $E$; composing with the continuous $\iota$ gives the statement on $X$.
\end{proof}

\subsection{The mass tilt and the residual laws}

\begin{proposition}[Whole-field mass-tilt identities]\label{prop:masstilt}
Fix a regulator $a$, a physical square $\Lambda$ and $\Omega=[-R,R]^{2}\supseteq\bar\Lambda$. For
$\varepsilon\ge0$, the residual weight satisfies
$0<e^{-W_{a,\Lambda}-\varepsilon J_{\Omega,\Lambda}}\le M_{\mathrm{res}}=e^{6u|\Lambda|c_a^{2}}$ on
$C(\R^{2})$, and
$e^{-3u|\Lambda|c_a^{2}-3\varepsilon|\Omega\setminus\Lambda|c_a^{2}}\le z_{\varepsilon,\Omega,a;\Lambda}\le M_{\mathrm{res}}$.
For $\varepsilon>0$, $Z_{\mathrm{full}}=\Det_{\Omega,a}z_{\varepsilon,\Omega,a;\Lambda}$ and
$\nu^{v}_{\varepsilon,\Omega,a;\Lambda}=\iota_*\Res_{\varepsilon,\Omega,a;\Lambda}$. These identities
concern the whole field, including observations outside $\Omega$.
\end{proposition}

\begin{proof}
$x^{4}-6cx^{2}+3c^{2}=(x^{2}-3c)^{2}-6c^{2}$ gives $W_{a,\Lambda}\ge-6u|\Lambda|c_a^{2}$, and
$J\ge0$. Under $G'_{\Omega,a}$ the coordinate $\gamma'(z)$ is centered Gaussian with
variance $r(z)=C'_{\Omega,a}(z,z)\in[1/4,c_a]$ (Theorem~\ref{thm:G-order}), so
$\E W_{a,\Lambda}=3u\int_\Lambda(r-c_a)^{2}\le3u|\Lambda|c_a^{2}$ and
$\E J=3\int_{\Omega\setminus\Lambda}r^{2}\le3|\Omega\setminus\Lambda|c_a^{2}$; Jensen gives the lower
bound. The functionals are continuous on $C(\R^{2})$, hence Borel. Since
$e^{-V}=e^{T_{\Omega,a}}e^{-W-\varepsilon J}$, Theorem~\ref{thm:G-tilt}(c) with
$F=e^{-W-\varepsilon J}\ind_B\circ\iota\ge0$ gives
$\E_{G_\mu}[\ind_B(\iota\gamma_\mu)e^{-V}]=\Det_{\Omega,a}\E_{G'}[\ind_B(\iota\gamma')e^{-W-\varepsilon J}]$;
$B=X$ gives the normalizer identity.
\end{proof}

\subsection{The stabilizer and auxiliary-square arrows}

\begin{proposition}[Stabilizer arrow]\label{prop:stab}
With $a,\Omega,\Lambda$ fixed, $\|\Res_{\varepsilon,\Omega,a;\Lambda}-\Res_{0,\Omega,a;\Lambda}\|_{\mathrm{var}}\to0$ as
$\varepsilon\downarrow0$, where $\|\rho-\eta\|_{\mathrm{var}}=\sup_{|H|\le1\text{ Borel}}|\rho(H)-\eta(H)|$; hence
$\|\nu^{v}_{\varepsilon,\Omega,a;\Lambda}-\nu^{(2)}_{\Omega,a;\Lambda}\|_{\mathrm{var}}\to0$,
$z_\varepsilon\to z_0$ and $Z_{\mathrm{full}}\to\Det_{\Omega,a}z_0$.
\end{proposition}

\begin{proof}
Under $G'_{\Omega,a}$ the weights $h_\varepsilon=e^{-W-\varepsilon J}$ increase pointwise to $h_0$ as
$\varepsilon\downarrow0$ and are bounded by $M_{\mathrm{res}}$, so
$\|h_\varepsilon-h_0\|_1=z_0-z_\varepsilon\to0$ and
$\|h_\varepsilon/z_\varepsilon-h_0/z_0\|_1\le2\|h_\varepsilon-h_0\|_1/z_0$. The variation distance of the path
laws equals this $L^{1}$ distance, and a measurable pushforward does not increase
it.
\end{proof}

\begin{proposition}[Auxiliary-square arrow]\label{prop:aux}
With a regulator $a$ and $\Lambda$ fixed, along every increasing family of squares
$\Omega=[-R,R]^{2}\supseteq\bar\Lambda$ with $R\to\infty$,
$\beta(\nu^{(2)}_{\Omega,a;\Lambda},\nu_{a,\Lambda})\le e^{3u|\Lambda|c_a^{2}}\bigl(2e_\Omega+M_{\mathrm{res}}r_\Omega\bigr)\to0$
and $|Z_{a,\Lambda}-z_{0,\Omega,a;\Lambda}|\le e_\Omega\to0$, where $e_\Omega$ and $r_\Omega$ are defined in
the proof. For each $\Omega$ a separate coupling is used.
\end{proposition}

\begin{proof}
Fix $\Omega$ and write $\Delta=\Delta_{\Omega,a}$. On a product space let $\gamma'$ have law
$G'_{\Omega,a}$ and let $\psi_\Omega$ be an independent centered Gaussian field with
covariance $\Delta$ and continuous paths (Theorem~\ref{thm:G-diff}(e)); independence
of the two continuous paths follows from their countable dense point
evaluations. Put $X_\Omega=\gamma'+\psi_\Omega$ and $Y_\Omega=\gamma'$. Then $X_\Omega$ is a continuous
centered Gaussian field with covariance $C'+\Delta=C_a$, so its law on $C(\R^{2})$ is
that of $\tilde\phi_a$ (Theorem~\ref{thm:G-point}). On the common Gaussian space
$\Phi_a=\iota\tilde\phi_a$ almost surely (Theorem~\ref{thm:G-versions}(b)) and
$U_{a,\Lambda}=W_{a,\Lambda}(\tilde\phi_a)$, so
$\nu_{a,\Lambda}(B)=\E[\ind_B(\iota X_\Omega)e^{-W(X_\Omega)}]/\E e^{-W(X_\Omega)}$ and
$Z_{a,\Lambda}=\E e^{-W(X_\Omega)}$, while
$\nu^{(2)}_{\Omega,a;\Lambda}(B)=\E[\ind_B(\iota Y_\Omega)e^{-W(Y_\Omega)}]/z_0$.

Let $\delta_\Omega(K)=\sup_{z\in K}\Delta(z,z)$ for compact $K$. By Parseval,
$p_j(\iota\psi)\le\|\chi_j\psi\|_{L^{2}}$, so
$\E p_j(\iota X_\Omega-\iota Y_\Omega)^{2}\le\|\chi_j\|_2^{2}\delta_\Omega(\supp\chi_j)$ and
$\E d(\iota X_\Omega,\iota Y_\Omega)\le r_\Omega:=\sum_j2^{-j}\min(1,\|\chi_j\|_2\delta_\Omega(\supp\chi_j)^{1/2})$,
which tends to $0$ by Theorem~\ref{thm:G-diff}(d) and dominated convergence.
For $p_c(x)=x^{4}-6cx^{2}+3c^{2}$,
$|p_c(x)-p_c(y)|\le|x-y|(|x|+|y|)(x^{2}+y^{2}+6c)\le|x-y|\bigl[2(|x|^{3}+|y|^{3})+6c(|x|+|y|)\bigr]$.
Each $X_\Omega(z)$, $Y_\Omega(z)$ has variance at most $c=c_a$, so the bracket has $L^{2}$
norm at most $L_c=(4\sqrt{15}+12)c^{3/2}$, and Cauchy--Schwarz gives
$\E|W(X_\Omega)-W(Y_\Omega)|\le u|\Lambda|L_c\,\delta_\Omega(\bar\Lambda)^{1/2}$. Both interactions are at
least $-6u|\Lambda|c^{2}$, where $x\mapsto e^{-x}$ is $M_{\mathrm{res}}$-Lipschitz, so
$\E|e^{-W(X_\Omega)}-e^{-W(Y_\Omega)}|\le e_\Omega:=M_{\mathrm{res}}u|\Lambda|L_c\delta_\Omega(\bar\Lambda)^{1/2}\to0$.
For $|H|\le1$ and $\mathrm{Lip}_d(H)\le1$, adding and subtracting the cross terms gives
$|\nu^{(2)}(H)-\nu_{a,\Lambda}(H)|\le[2e_\Omega+\E(e^{-W(X_\Omega)}d(\iota X_\Omega,\iota Y_\Omega))]/z_0$, and
$z_0\ge e^{-3u|\Lambda|c^{2}}$ by Proposition~\ref{prop:masstilt}.
\end{proof}

\subsection{The regulator arrow}

\begin{proposition}[Regulator arrow]\label{prop:reg}
With $\Lambda$ fixed, $Z_{a,\Lambda}\to Z_{0,\Lambda}\ge1$ and
$\beta(\nu_{a,\Lambda},\nu_{0,\Lambda})\le2e_a+\Xi_2(\Lambda)^{1/2}h(a)\to0$ as $a\downarrow0$, where
$e_a=\|e^{-U_{a,\Lambda}}-e^{-U_{0,\Lambda}}\|_{L^{1}(\PP_0)}\le2K_\Lambda\Xi_2(\Lambda)^{1/2}\sqrt a$ and $h(a)$
is as in Theorem~\ref{thm:G-versions}(e).
\end{proposition}

\begin{proof}
Theorem~\ref{thm:W}(f) with $p=1$ bounds $e_a$, so $|Z_a-Z_0|\le e_a$, and $Z_a,Z_0\ge1$ by
Theorem~\ref{thm:W}(g). For $|H|\le1$, $\mathrm{Lip}_d(H)\le1$, with $w_a=e^{-U_a}$,
$|\E[H(\Phi_a)w_a-H(\Phi_0)w_0]|\le e_a+\E[w_0d(\Phi_a,\Phi_0)]\le e_a+\Xi_2^{1/2}h(a)$ by
Cauchy--Schwarz, Theorem~\ref{thm:W}(e) with $p=2$ and
Theorem~\ref{thm:G-versions}(e); the normalizer correction is at most $e_a$.
\end{proof}

\begin{remark}[The chain of laws]\label{rem:chain}
For fixed parameters the five laws
$P^{b}_{N,\varepsilon,\Omega,a;\Lambda}$, $\nu^{v}_{\varepsilon,\Omega,a;\Lambda}=\iota_*\Res_\varepsilon$,
$\nu^{(2)}_{\Omega,a;\Lambda}=\iota_*\Res_0$, $\nu_{a,\Lambda}$ and $\nu_{0,\Lambda}$ are joined by the four
arrows of Propositions~\ref{prop:width}, \ref{prop:stab}, \ref{prop:aux} and
\ref{prop:reg}; there is no mass change, no factorial change, no omitted
quadratic term and no extra normalization between them. The quadratic-tilt
normalizer $\Det_{\Omega,a}$ is nondecreasing in $\Omega$ (Theorem~\ref{thm:G-tilt}(a)).
No monotonicity in $\Omega$ is asserted for the full Gaussian normalizer
$Z_{\mathrm{full}}=\Det_{\Omega,a}z_{\varepsilon,\Omega,a;\Lambda}$, and none holds in general at a fixed
$\varepsilon>0$. For $\Omega_1=\bar\Lambda$ the outer integral $J_{\Omega_1,\Lambda}$ vanishes,
so $Z_{\mathrm{full}}$ for $\Omega_1$ does not depend on $\varepsilon$. For a strictly larger
centered square $\Omega_2$, the continuous endpoint Gaussian is almost surely
nonzero at an interior point of $\Omega_2\setminus\bar\Lambda$, so $J_{\Omega_2,\Lambda}>0$ almost
surely; since $e^{-W_{a,\Lambda}}e^{T_{\Omega_2,a}}\le e^{6u|\Lambda|c_a^{2}}e^{T_{\Omega_2,a}}$ is
integrable, dominated convergence gives $Z_{\mathrm{full}}\to0$ for $\Omega_2$ as
$\varepsilon\to\infty$. Only the residual normalizers $z_{0,\Omega,a;\Lambda}$, which satisfy
$e^{-3u|\Lambda|c_a^{2}}\le z_{0,\Omega,a;\Lambda}\le e^{6u|\Lambda|c_a^{2}}$ for every $\Omega$
(Proposition~\ref{prop:masstilt}), enter the auxiliary arrow. The identification of $\nu_{0,\Lambda}$ with
the classical free-boundary law is Proposition~\ref{prop:identify}.
\end{remark}

% ===========================================================================
\section{The continuum target: reduced lemmas}\label{sec:target}

The growth of the physical square is handled with classical input in three steps. Proposition~\ref{prop:identify} identifies the continuum law $\nu_{0,\Lambda}$ on compact smooth tests with the classical free-boundary law, so the classical theorems apply. Theorem~\ref{thm:LR} uses dilation (Lemma~\ref{lem:dilation}) to move $(m,u)$ to the reference mass, where \ref{prem:C2} makes the two classical states coincide and \ref{prem:C1} gives convergence for finitely many smearings. Theorem~\ref{thm:V1} upgrades this to weak convergence on $X$ with all compact moments, and Theorem~\ref{thm:CT} records the properties of the target.

This section proves, directly from \ref{prem:C1}--\ref{prem:C6} and the
foundations of Sections~\ref{sec:space}--\ref{sec:fixed}, the three continuum
statements used by the construction: the continuum limit with its Wick
identification (Theorem~\ref{thm:LR}), the whole-field volume limit with all
compact moments (Theorem~\ref{thm:V1}) and the target properties
(Theorem~\ref{thm:CT}). Appendix~\ref{app:provenance} records the argument lineage.

\subsection{Classical finite-volume laws and the Wick identification}\label{sec:LR}

For $\psi\in\cS'(\R^{2})$ and a real even nonnegative $\rho_1\in C_c^{\infty}(\R^{2})$ with
$\int\rho_1=1$, put $\rho_\delta=\delta^{-2}\rho_1(\cdot/\delta)$ and
$\psi_\delta(z)=\psi(\rho_\delta(z-\cdot))$.

\begin{lemma}[Classical Wick existence and finite normalizers from standard tools]
\label{lem:classical-wick}
Let $r,b>0$ and $B\subset\R^2$ be bounded measurable, with $A=|B|$.
On the free Gaussian space $(\cS',\cC,P_r)$ the own-variance mollified
quartics $b\int_B{:}\psi_\delta(z)^4{:}\,dz$ have a mollifier-independent
limit $U^S_{r,b,B}$ in the fourth homogeneous Gaussian chaos in $L^2(P_r)$.
It has mean zero and, for every real $t>0$,
\[
 \E_{P_r}e^{-tU^S_{r,b,B}}<\infty,\qquad
 1\le Z^S_{r,b,B}:=\E_{P_r}e^{-U^S_{r,b,B}}<\infty.
\]
The normalized density defines the classical finite-volume law in the
convention of Section~\ref{sec:classical}. If $B$ is a centered axis-parallel
square, this law is invariant under field sign and $\theta_1,\theta_2$.
This lemma uses only the free covariance and \ref{prem:C6}; it does not use
the Gaussian or Wick foundations of Sections~\ref{sec:gauss}--\ref{sec:wick}.
Every constant below is fixed-volume. If $A=0$, the interaction is zero and
all conclusions hold with $Z^S=1$.
\end{lemma}

\begin{proof}
Assume $A>0$ and write $c(p)=(r^2+|p|^2)^{-1}$ and $C=\cI c$.
Polar integration and Hausdorff--Young give
\begin{equation}\label{eq:classical-L4}
 \int_{\R^2}c(p)^{4/3}dp=3\pi r^{-2/3},\qquad
 \|C\|_4\le M_4(r):=(2\pi)^{-3/2}(3\pi)^{3/4}r^{-1/2}.
\end{equation}
For the real even mollifier fixed above, the cross-covariance of
$\psi_\delta(z)$ and $\psi_\epsilon(z')$ is
$C_{\delta\epsilon}(z-z')=(\rho_\delta*C*\rho_\epsilon)(z-z')$.
Approximate-identity convergence gives $C_{\delta\epsilon}\to C$ in $L^4$;
equivalently apply Hausdorff--Young and dominated convergence to
$\hat\rho_1(\delta p)\hat\rho_1(\epsilon p)c(p)$.
The multipliers are bounded uniformly, converge pointwise to one and are
dominated by a constant times $c\in L^{4/3}$. The same argument works with
two different real mollifiers of integral one. Since
\[
 \|f^4-g^4\|_1\le\|f-g\|_4
 (\|f\|_4^3+\|f\|_4^2\|g\|_4+\|f\|_4\|g\|_4^2+\|g\|_4^3),
\]
the fourth powers converge in $L^1$. The Wick covariance identity gives
\begin{equation}\label{eq:classical-chaos}
 \E J_\delta J_\epsilon
 =4!b^2\int_{B^2}C_{\delta\epsilon}(z-z')^4\,dz\,dz',
 \qquad J_\delta=b\int_B{:}\psi_\delta(z)^4{:}\,dz.
\end{equation}
For integrable $H$, $\int_{B^2}|H(z-z')|\le A\|H\|_1$.
Thus $J_\delta$ is Cauchy in $L^2$, with the same limit for every such
mollifier. The integrands are continuous $L^2$ fourth-chaos variables;
their Bochner integrals equal the pathwise integrals by Fubini, using finite
Gaussian moments on $B$. Chaos closure proves fourth-chaos membership and
mean zero. No point variable or product of distributions at zero cutoff is
formed.

\emph{A cutoff with a lower bound.} Fix a real radial
$\chi\in C_c^\infty(\R^2)$ with $0\le\chi\le1$, $\chi=1$ on the unit ball
and $\chi=0$ outside the ball of radius two. For
$R\ge R_0:=\max(1,r)$, put $\chi_R(p)=\chi(p/R)$,
$\psi_R(z)=\psi(\cI\chi_R(z-\cdot))$ and
$U_R=b\int_B{:}\psi_R(z)^4{:}_{c_R}\,dz$. The inverse transform is a real
Schwartz function, so these are ordinary smoothed evaluations of $\psi$;
their paths are smooth and jointly measurable. Their variance and lower bound
are
\begin{equation}\label{eq:classical-cutoff-variance}
 c_R=(2\pi)^{-2}\int\chi_R^2c
 \le\frac1{4\pi}\log(1+4R^2/r^2),\qquad
 U_R\ge-6bA c_R^2.
\end{equation}
The last inequality is $(x^2-3c)^2-6c^2\ge-6c^2$.
Let $C_R^1=\cI(\chi_Rc)$, $C_R^2=\cI(\chi_R^2c)$ and
$a_4=(2\pi)^{-3/2}(3\pi)^{3/4}$. Polar integration of the tail gives,
for $j=1,2$,
\begin{equation}\label{eq:classical-covariance-tail}
 \|C-C_R^j\|_4\le a_4R^{-1/2},\qquad
 \|C_R^j\|_4\le M_4(r).
\end{equation}
Indeed $|1-\chi_R^j|\le\ind_{\{|p|>R\}}$ and
$\int_{|p|>R}c^{4/3}\le3\pi R^{-2/3}$.
Taking the mollifier limit in the Gaussian cross-covariances yields
\begin{align}
 \|U^S_{r,b,B}-U_R\|_2^2
 &=4!b^2\int_{B^2}[C^4+(C_R^2)^4-2(C_R^1)^4](z-z')\,dz\,dz'
 \nonumber\\
 &\le288 b^2 A M_4(r)^3a_4 R^{-1/2}.
 \label{eq:classical-interaction-tail}
\end{align}
Here each fourth-power difference is bounded by four times $M_4(r)^3$
times its $L^4$ difference, and the coefficients of the two differences
are one and two. Put $a_2=b[288 A M_4(r)^3a_4]^{1/2}$. Both interactions
lie in the fourth chaos, so
\begin{equation}\label{eq:classical-hypercontractivity}
 \|U^S_{r,b,B}-U_R\|_p\le(p-1)^2a_2R^{-1/4}\quad(p\ge2).
\end{equation}

\emph{Negative tails and every needed normalization.} The following explicit
choices make the fixed-volume character of the constants unambiguous:
\[
 c_0=\frac1{4\pi}\log(1+4R_0^2/r^2),\qquad
 d=\frac{\pi}{\sqrt{24bA}},\qquad R(K)=R_0e^{d\sqrt K},
\]
\[
 K_0=\max\{1,32bA c_0^2,(16\log2/d)^2,2e a_2\}.
\]
For $K\ge K_0$, the logarithmic variance bound implies
$c_{R(K)}\le c_0+d\sqrt K/(2\pi)$ and hence
$6bA c_{R(K)}^2\le12bA c_0^2+3bA d^2K/\pi^2\le K/2$.
On $\{U^S_{r,b,B}<-K\}$ the difference from $U_{R(K)}$ therefore has
absolute value greater than $K/2$. Set $p=R(K)^{1/16}\ge2$. Markov and
\eqref{eq:classical-hypercontractivity} give
\begin{equation}\label{eq:classical-negative-tail}
 P_r(U^S_{r,b,B}<-K)
 \le\left[\frac{2a_2}{K}R(K)^{-1/8}\right]^{R(K)^{1/16}}
 \le\exp[-R(K)^{1/16}].
\end{equation}
The bracket is at most $e^{-1}$ by the choice of $K_0$.
For every fixed $t>0$, the layer-cake formula bounds the exponential moment
by
\begin{equation}\label{eq:classical-normalizer}
 \E e^{-tU^S_{r,b,B}}
 \le e^{tK_0}+\int_{K_0}^\infty
       te^{tK}\exp[-R_0^{1/16}e^{d\sqrt K/16}]\,dK<\infty.
\end{equation}
The double exponential in $\sqrt K$ dominates every $e^{tK}$: for example
$e^{d\sqrt K/16}/K\to\infty$. Jensen and $\E U^S_{r,b,B}=0$ give $Z^S\ge1$.
Thus both the $L^2$ and finite-normalizer provisos in the source convention
are proved before that convention is used. These estimates are separate
from the constants $K_\Lambda,\Xi_p$ of Theorem~\ref{thm:W}.

Finally $P_r$ is invariant under field sign and Euclidean isometries, by
its centered Gaussian covariance and cylinder-law uniqueness. With radial
$\chi_R$, the interaction $U_R$ on a centered axis-parallel square is
unchanged by sign and $\theta_1,\theta_2$. The same identities pass in
$L^2$ to $U^S$, then to its exponential and its normalized law. This
symmetry proof uses no foundation interaction or normalizer.
\end{proof}


\begin{lemma}[Transfer between Gaussian realizations]\label{lem:transfer-gauss}
Let $(\Gamma_0,\cF_0,\PP_0,W)$ be the common space of Theorem~\ref{thm:G-common} at
mass $m$ and $P_m$ the free measure on $(\cS',\cC)$ of Section~\ref{sec:classical}.
There is a unique map $\Xi$ from $L^{0}(\Gamma_0,\cF^{\cS}_{\mathrm{free}},\PP_0)$,
$\cF^{\cS}_{\mathrm{free}}=\sigma(\phi(f):f\in\cS(\R^{2};\R))\vee\cN$, to $L^{0}(\cS',\cC,P_m)$ with
$\Xi[g((\phi(f_j))_j)]=g((\psi(f_j))_j)$ for every sequence of real Schwartz $f_j$ and
every Borel $g:\R^{\N}\to\C$. It is linear and multiplicative, preserves joint laws
of countable families, is an isometry of every $L^{r}$, commutes with
$L^{r}$-limits and with Bochner integrals in $L^{2}$, and maps
${:}\xi^{4}{:}_{\Var\xi}$ to ${:}(\Xi\xi)^{4}{:}_{\Var\xi}$.
\end{lemma}

\begin{proof}
For real Schwartz $f,g$,
$\E_0[\phi(f)\phi(g)]=(2\pi)^{-2}\int\conj{\hat f}\hat g\hat C_0=\E_{P_m}[\psi(f)\psi(g)]$
(Theorem~\ref{thm:G-common}), so for any two sequences the vectors
$((\phi(f_j)),(\phi(f'_j)))$ and $((\psi(f_j)),(\psi(f'_j)))$ are centered Gaussian on
$\R^{\N}\times\R^{\N}$ with the same covariances, hence have the same law. If
$g((\phi(f_j)))=g'((\phi(f'_j)))$ almost surely, the Borel set $\{g\ne g'\}$ therefore
has measure zero for the $\psi$-vectors too, so $\Xi$ is well defined. Every
$Y\in L^{0}(\cF^{\cS}_{\mathrm{free}})$ has such a representation: for a countable family
$\cD_S$ dense in $\cS$, $\|\phi(f)-\phi(d)\|_2\le m^{-1}\|f-d\|_2$
(Theorem~\ref{thm:G-bounds}), so each $\phi(f)$ is an almost sure limit of
$\phi(d_k)$, and Doob--Dynkin applies. Merging families gives linearity and
multiplicativity; law preservation gives the isometries; a bounded linear map
commutes with Bochner integrals; variances are preserved.
\end{proof}

\begin{definition}[Classical finite-volume laws]\label{def:classical-laws}
For a physical square $\Lambda$, the mollified own-variance Wick integrals
$N^{S}_\delta=u\int_\Lambda{:}\psi_\delta(z)^{4}{:}\,dz$ (pathwise integrals of the continuous
functions $z\mapsto\psi_\delta(z)$, equal to the corresponding Bochner integrals in
$L^{2}(P_m)$) converge in $L^{2}(P_m)$ as $\delta\downarrow0$
(Lemma~\ref{lem:classical-wick}); put $U^{S}_{m,u,\Lambda}=\lim_\delta N^{S}_\delta$,
$Z^{S}_{m,u,\Lambda}=\E_{P_m}e^{-U^{S}}$ and
$d\nu_{m,u,\Lambda}=(Z^{S})^{-1}e^{-U^{S}}dP_m$. The normalizer is positive and
finite by Lemma~\ref{lem:classical-wick}, independently of the project Wick
foundations. By the convention paragraph of
Section~\ref{sec:classical}, this is the free-boundary finite-volume law of
\cite[(1.1), pp.~496--497]{FS1977} for $P(X)=uX^{4}$ at mass $m$, and the
continuum law of \cite[\S IX]{GRS1976} with boundary condition F.
\end{definition}

\begin{proposition}[Identification of the zero-regulator law]\label{prop:identify}
For every $m,u>0$ and physical square $\Lambda$:
\begin{enumerate}[label=(\roman*)]
\item $U^{S}_{m,u,\Lambda}$ exists, does not depend on $\rho_1$, and
  $U^{S}=\Xi U_{0,\Lambda}$; also $Z^{S}_{m,u,\Lambda}=Z_{0,\Lambda}\in[1,\Xi_1(\Lambda)]$ and
  $\|d\nu_{m,u,\Lambda}/dP_m\|_{L^{2}(P_m)}\le\Xi_2(\Lambda)^{1/2}$;
\item for every finite tuple of real $f_1,\dots,f_r\in C_c^{\infty}(\R^{2})$, the law of
  $(v(f_1),\dots,v(f_r))$ under $\nu_{0,\Lambda}$ equals the law of
  $(\psi(f_1),\dots,\psi(f_r))$ under $\nu_{m,u,\Lambda}$;
\item $\nu_{m,u,\Lambda}$ is invariant under $\psi\mapsto-\psi$ and under $\theta_1$, $\theta_2$.
\end{enumerate}
\end{proposition}

\begin{proof}
(i) The family $\chi_\delta(p)=\conj{\hat\rho_1(\delta p)}$ is admissible
(Definition~\ref{def:W-objects}), with $\phi^{\chi}_\delta(z)=\phi(\rho_\delta(z-\cdot))$, so by
Theorem~\ref{thm:W}(b) the corresponding Bochner integrals $N_\delta$ on $\Gamma_0$
converge to $U_0=U_{0,\Lambda}$ in $L^{2}(\PP_0)$. Each $N_\delta$ and $U_0$ is
$\cF^{\cS}_{\mathrm{free}}$-measurable (Theorem~\ref{thm:W}(b)), and
$\Xi N_\delta=N^{S}_\delta$ by Lemma~\ref{lem:transfer-gauss}. Hence $N^{S}_\delta\to\Xi U_0$ in
$L^{2}(P_m)$; the limit does not depend on $\rho_1$ because $U_0$ does not.
Multiplicativity and law preservation give $\E_{P_m}e^{-pU^{S}}=\E_0e^{-pU_0}$, and
Theorem~\ref{thm:W}(e),(g) give the bounds. (ii) By
Theorem~\ref{thm:G-versions}(a), $\Phi_0(f_i)=\phi(f_i)$ almost surely; for bounded
Borel $G$, $\E_0[G(\phi(f))e^{-U_0}]=\E_{P_m}[G(\psi(f))e^{-U^{S}}]$ by the lemma. (iii)
$P_m$ is invariant under $\psi\mapsto-\psi$ and under $\theta_j$, and the mollified
quartics are even. By (i) we may take $\rho_1$ radial; then
$(\theta_j\psi)_\delta(z)=\psi_\delta(\theta_jz)$, and since $\Lambda$ is $\theta_j$-invariant (a centered
axis-parallel square) the change of variables $z\mapsto\theta_jz$ shows that
$N^{S}_\delta\circ\theta_j=N^{S}_\delta$; hence $U^{S}\circ\theta_j=U^{S}$ almost surely.
\end{proof}

\begin{lemma}[Dilation]\label{lem:dilation}
For $s>0$ let $(L_sf)(x)=s^{-2}f(x/s)$ and $(T_s\psi)(f)=\psi(L_sf)$ on $\cS'$. Then
$(T_s)_*P_M=P_{sM}$; for $s=m/M$, $\lambda=u/s^{2}$ and every physical square $\Lambda$,
$U^{S}_{m,u,\Lambda}\circ T_s=U^{S}_{M,\lambda,s\Lambda}$ $P_M$-almost surely and
$\nu_{m,u,\Lambda}=(T_s)_*\nu_{M,\lambda,s\Lambda}$. With $(A_b\psi)(f)=\psi(f(\cdot-b))$ and
$(R\psi)(f)=\psi(f\circ R^{-1})$ for orthogonal $R$: $A_bT_s=T_sA_{sb}$ and $RT_s=T_sR$;
$T_s$ commutes with $\theta_1,\theta_2$ and maps $\cC_{H^{\theta}_{+}}$ into itself.
\end{lemma}

\begin{proof}
This proof uses the classical $L^2$ limits and finite normalizers of
Lemma~\ref{lem:classical-wick}, without the common-space identification.
$\widehat{L_sf}(p)=\hat f(sp)$, so
\[
\begin{aligned}
\Var_{P_M}\psi(L_sf)&=(2\pi)^{-2}\int|\hat f(sp)|^{2}(M^{2}+|p|^{2})^{-1}dp\\
&=(2\pi)^{-2}\int|\hat f(q)|^{2}(s^{2}M^{2}+|q|^{2})^{-1}dq.
\end{aligned}
\]
Next,
$L_s[\rho_\delta(x-\cdot)]=\rho_{s\delta}(sx-\cdot)$, so $(T_s\psi)_\delta(x)=\psi_{s\delta}(sx)$; since
$T_s\psi$ has law $P_m$ under $P_M$, the Wick constants agree,
$\Var_{P_m}\phi_\delta(x)=\Var_{P_M}\psi_{s\delta}(sx)$. The change of variables $y=sx$ in the
Bochner integral gives
$u\int_\Lambda{:}(T_s\psi)_\delta(x)^{4}{:}\,dx=\lambda\int_{s\Lambda}{:}\psi_{s\delta}(y)^{4}{:}\,dy$; as
$T_s$ is measure preserving from $(\cS',P_M)$ to $(\cS',P_m)$, the $L^{2}$ limits agree, and
the normalized laws correspond. The intertwining relations follow from
$L_s[f(\cdot-b)]=(L_sf)(\cdot-sb)$ and $L_s(f\circ R^{-1})=(L_sf)\circ R^{-1}$; $L_s$ preserves
supports in $\{z_2>0\}$ and $\{z_1>0\}$.
\end{proof}

\begin{lemma}[Readings of the convergence in \ref{prem:C1}(iv)]\label{lem:readings}
Let $\mu_n,\mu$ be probabilities on $(\cS',\cC)$ and $g\in C_c^{\infty}$ real, and suppose
$\E_\mu e^{t\psi(g)}<\infty$ for all real $t$. If
$\E_{\mu_n}e^{it\psi(g)}\to\E_\mu e^{it\psi(g)}$ for every real $t$, or all moments of
$\psi(g)$ converge, or $\E_{\mu_n}e^{t\psi(g)}\to\E_\mu e^{t\psi(g)}$ for all real $t$, then the
laws of $\psi(g)$ converge weakly. For the states of \ref{prem:C1} at mass $M$,
$\E e^{t\psi(g)}\le\|f_B\|_{L^{2}(P_M)}\exp(t^{2}\langle g,(-\Delta+M^{2})^{-1}g\rangle)$ for
every real $t$ and every bounded open $B\supseteq\supp g$, by \ref{prem:C1}(ii) and
Cauchy--Schwarz.
\end{lemma}

\begin{proof}
Characteristic functions: L\'evy's continuity theorem, applied with convergence
at every real frequency $t$ to the characteristic function of the law of $\psi(g)$
under $\mu$. Moments: a law with finite exponential moments is determined by its
moments, and moment convergence then implies weak convergence. Laplace
transforms: convergence of moment generating functions on a neighbourhood of $0$
implies weak convergence.
\end{proof}

\begin{remark}[One frequency does not suffice]\label{rem:readings}
Convergence of $\E_{\mu_n}e^{i\psi(g)}$ at the single frequency $t=1$ does not imply
weak convergence of the laws of $\psi(g)$: for a nonzero real $g\in C_c^{\infty}$ and
the regular tempered distribution $h=2\pi g/\int g^{2}$, the probabilities
$\mu_n=\delta_h$ and $\mu=\delta_0$ have equal characteristic values at frequency
one, and $\mu$ has every exponential moment, while the laws of $\psi(g)$ are
$\delta_{2\pi}$ and $\delta_0$. An earlier version of Lemma~\ref{lem:readings} admitted
that single-frequency alternative; it is withdrawn. Where the
characteristic-function reading of \ref{prem:C1}(iv) is used (proof of
Theorem~\ref{thm:LR}), it is applied to $g=\sum_jt_jL_sf_j$ for every $t\in\R^{r}$,
which supplies every real frequency.
\end{remark}

\begin{theorem}[Continuum limit and target; reduced lemma LR-R]\label{thm:LR}
Fix the admissible pair of Definition~\ref{def:threshold} and $(m,u)$ in the regime;
let $s=m/M_{\mathrm{ref}}$ and $\lambda=uM_{\mathrm{ref}}^{2}/m^{2}\in(0,\lambda_*)$. The two states
$\nu_\pm$ of \ref{prem:C1}(i) for $\lambda X^{4}$ at mass $M_{\mathrm{ref}}$ coincide; call the
common state $v$, and put $\nu^{S}=(T_s)_*v$. Then for every $\ell>0$ and every finite
tuple of real $f_1,\dots,f_r\in C_c^{\infty}(\R^{2})$,
\[
\Law_{\nu_{m,u,\Lambda_n}}(\psi(f_1),\dots,\psi(f_r))\Rightarrow\Law_{\nu^{S}}(\psi(f_1),\dots,\psi(f_r)),
\qquad \Lambda_n=(-n\ell,n\ell)^{2},
\]
and therefore, by Proposition~\ref{prop:identify}(ii), the laws of
$(v(f_1),\dots,v(f_r))$ under $\nu_n=\nu_{0,\Lambda_n}$ converge to the same limit. The
same $\nu^{S}$ serves all tuples and all $\ell$.
\end{theorem}

\begin{proof}
\phantomsection\label{proof:LR-source}
At mass $M_{\mathrm{ref}}>M_0$ and $0<\lambda<\lambda_*\le\lambda_0$, both $\nu_+$ and $\nu_-$ belong to
$\cG_{\mathrm{ur}}$ (\ref{prem:C2}(i)), which is a singleton (\ref{prem:C2}(ii)); the
states coincide as measures on $\cC$ (only $\cC$-measurable observables are used, so
only the restrictions of the AHKZ members to $\cC$ enter). By \ref{prem:C1}(iii)
the pressure is differentiable at zero external field, and by \ref{prem:C1}(iv)
the free measures $\nu_{M_{\mathrm{ref}},\lambda,\Lambda'}$ converge to $v$ along every exhaustion by
centered squares $\Lambda'$, on characteristic functions of compact smooth tests
(under the moment or Laplace readings, Lemma~\ref{lem:readings} gives weak
convergence of the law of $\psi(g)$, hence convergence of $\E e^{i\psi(g)}$, for every
real compact smooth $g$). For $t\in\R^{r}$ put
$g=\sum_jt_jL_sf_j$, a real compact smooth test. By Lemma~\ref{lem:dilation},
$\E_{\nu_{m,u,\Lambda_n}}e^{i\sum_jt_j\psi(f_j)}=\E_{\nu_{M_{\mathrm{ref}},\lambda,s\Lambda_n}}e^{i\psi(g)}\to\E_ve^{i\psi(g)}
=\E_{\nu^{S}}e^{i\sum_jt_j\psi(f_j)}$, since $s\Lambda_n=(-ns\ell,ns\ell)^{2}$ is an exhaustion by
centered squares. Since $t$ ranges over all of $\R^{r}$, this is convergence of the
joint characteristic functions at every frequency; the limit is the
characteristic function of $\Law_{\nu^{S}}(\psi(f_1),\dots,\psi(f_r))$, and L\'evy's
continuity theorem concludes. The classical finite-volume laws used here
exist with finite normalizers by Lemma~\ref{lem:classical-wick}; that direct
lemma, dilation and the readings argument use only \ref{prem:C6} and the
stated classical sources.
\phantomsection\label{proof:LR-X}
Finally Proposition~\ref{prop:identify}(ii) supplies the separate identification
with $\nu_n$ on the common project space. Only this last bridge uses the
Gaussian/Wick foundations.
\end{proof}

\begin{remark}[Scope of Theorem~\ref{thm:LR}]\label{rem:LR-scope}
The classical target and finite-joint classical limit use only
\ref{prem:C1}, \ref{prem:C2}, \ref{prem:C6} and the direct classical
Lemma~\ref{lem:classical-wick}. Its exponential normalizers follow from
the explicit negative-tail proof, separately from $L^2$ convergence. The
identification with the project common-space law additionally uses
Proposition~\ref{prop:identify} and the Gaussian/Wick foundations. No
neural law, stagewise neural bridge or midpoint approximation is used.
The convergence asserted is finite-joint weak convergence for compact smooth
tests; no moment, density, total-variation or rate conclusion is drawn from
\ref{prem:C1}(iv).
\end{remark}

\subsection{The volume limit on $X$}\label{sec:V1}

\begin{theorem}[Volume estimates and whole-field limit; reduced lemma V1-R]\label{thm:V1}
\begin{enumerate}[label=(\roman*)]
\item \emph{(E, fixed volume.)} For all $m,u>0$, every physical square $\Lambda$, every
  nonnegative real $g\in C_c^{\infty}(\Lambda)$ and real $t$,
  $\E_{\nu_{m,u,\Lambda}}e^{t\psi(g)}\le\exp\bigl(t^{2}\E_{\nu_{m,u,\Lambda}}\psi(g)^{2}/2\bigr)$.
\item \emph{(E, uniform.)} Let $(m,u)$ be in the regime, $g\ge0$ real compact smooth
  and $N_g\ge2$ an integer such that $\supp g\subset\Lambda_n$ for $n\ge N_g$ (the
  choice $N_g\ge2$ makes the early-volume maximum below nonempty). With $R_g>0$ such that
  $\sup_{n\ge N_g}\nu_n(|v(g)|>R_g)<1/12$ and $V_g=2R_g^{2}$: for $n\ge N_g$, real smooth $f$
  with $|f|\le g$ and real $t$, $\E_{\nu_n}e^{tv(f)}\le e^{V_gt^{2}/2}$, and
  $\E_{\nu_n}|v(f)|^{p}\le2(2V_g)^{p/2}\Gamma(1+p/2)$ for $p>0$. Including the finitely many
  $n<N_g$, $\sup_{n\ge1,|f|\le g}\E_{\nu_n}|v(f)|^{p}\le C_{p,g}<\infty$, with
  $C_{p,g}=\max\{2(2V_g)^{p/2}\Gamma(1+\tfrac p2),\ [2^{p}\Gamma(p+\tfrac12)/\sqrt\pi]^{1/2}m^{-p}\|g\|_2^{p}
  \max_{n<N_g}\Xi_2(\Lambda_n)^{1/2}\}$.
\item \emph{(F.)} For $j\ge1$ and $1<s<2$,
  $\sup_n\E_{\nu_n}\|\chi_jv\|^{2}_{H^{-s}(\T_j)}\le B_{j,s}<\infty$, and every finite positive
  moment of this norm is bounded uniformly in $n$. The family $(\nu_n)$ is tight on
  $X$. There is a Borel probability $\nu=\nu^{X}_\infty$ on $X$ and a $\cC$-measurable map
  $\mathcal R$ from a set of full $\nu^{S}$-measure into $X$ with
  $\mathcal R(\psi)(f)=\psi(f)$ for all compact smooth $f$ and $\mathcal R_*\nu^{S}=\nu$
  ($\nu$ \emph{represents} $\nu^{S}$); and $\nu_n\Rightarrow\nu$ weakly on $X$. The target is
  almost surely locally $H^{-s}$ for every $s>1$.
\item \emph{(M.)} For every finite tuple of real compact smooth tests and every
  real or complex polynomial $Q$, $\E_{\nu_n}Q(v(f_1),\dots)\to\E_\nu Q(v(f_1),\dots)$; all
  these expectations are finite, and every fixed positive power of $|Q|$ has
  expectation bounded uniformly in $n$.
\end{enumerate}
\end{theorem}

\begin{proof}
\emph{(i)} Fix $\Lambda$ of side $L$ and meshes $\delta=L/(2k+1)$ as in \ref{prem:C4}. By
\ref{prem:C4} the lattice density is
$\exp[-\sum_i(a_iq_i^{4}+b_iq_i^{2})+\sum_{i<j}J_{ij}q_iq_j]$ up to normalization, with
$a_i=u\delta^{2}>0$, $b_i=(B_\delta)_{ii}/2-6u\delta^{2}c_\delta\in\R$ and $J_{ij}=-(B_\delta)_{ij}\ge0$. That the
off-diagonal entries are nonpositive can be checked directly: with $I$ the sites
in $\Lambda$ and $O$ the others, $A_{OO}^{-1}=\sum_k(4+m^{2}\delta^{2})^{-1-k}\mathrm{adj}_{OO}^{k}$ has
nonnegative entries, $A_{IO}\le0$, and $B_\delta=A_{II}-A_{IO}A_{OO}^{-1}A_{OI}$. In the
ordered-pair form of \ref{prem:C5} the pair coefficients are $a_{ij}=a_{ji}=J_{ij}/2$.
For $\epsilon\in(0,1]$ replace them by $(J_{ij}+\epsilon)/2$; every fixed polynomial moment
and the normalizer converge as $\epsilon\to0$ by dominated convergence with the
integrable majorant
$|q|^{k}\prod_ie^{-a_iq_i^{4}-b_iq_i^{2}}\exp(\sum_{i<j}(J_{ij}+1)|q_i||q_j|+\sum_ih_i|q_i|)$, so
the third-cumulant inequality of \ref{prem:C5} holds at $\epsilon=0$. Let
$Y_\delta=\delta^{2}\sum_ig(\delta i)q_i$ and $F_\delta(t)=\log\E_\delta e^{tY_\delta}$, smooth by quartic
confinement. For $t\ge0$ the $t$-tilt adds the fields $h_i=t\delta^{2}g(\delta i)\ge0$, and
$F'''_\delta(t)$ is the third cumulant of $Y_\delta$, a nonnegative combination of mixed
third cumulants, hence $\le0$. Symmetry gives $F_\delta(0)=F'_\delta(0)=0$ and
$F''_\delta(0)=v_\delta=\E_\delta Y_\delta^{2}$; integrating three times,
$F_\delta(t)\le v_\delta t^{2}/2$ for $t\ge0$, and for all $t$ by evenness. By
\ref{prem:C4}(i), on the common Gaussian realization the smearings and the
interaction densities converge in every Gaussian $L^{p}$ with positive normalizer
limits; for bounded continuous $H$ the difference of weighted expectations is
controlled by the $L^{1}$ density error plus bounded convergence in probability,
so the normalized lattice laws of $Y_\delta$ converge weakly to that of $\psi(g)$ under
$\nu_{m,u,\Lambda}$ (Definition~\ref{def:classical-laws}), and $v_\delta\to\E\psi(g)^{2}$. For the
nonnegative continuous $y\mapsto e^{ty}$, the portmanteau inequality gives
$\E e^{t\psi(g)}\le\liminf_\delta\E_\delta e^{tY_\delta}\le e^{t^{2}\E\psi(g)^{2}/2}$.

\emph{(ii)} By Theorem~\ref{thm:LR} the laws of $v(g)$ under $\nu_n$ converge, hence
are tight, and $R_g$ exists. Let $n\ge N_g$ and $v_n=\E_{\nu_n}v(g)^{2}$; the laws on compact
smooth tests are those of $\nu_{m,u,\Lambda_n}$ (Proposition~\ref{prop:identify}). By
\ref{prem:C4}(ii) and the vanishing of odd moments,
$\E v(g)^{4}\le3v_n^{2}$. If $v_n>0$, Paley--Zygmund gives
$\nu_n(v(g)^{2}>v_n/2)\ge\tfrac14v_n^{2}/\E v(g)^{4}\ge1/12$; if $v_n>2R_g^{2}$ this event lies in
$\{|v(g)|>R_g\}$, a contradiction. So $v_n\le V_g$, and (i) gives the bound for $f=g$.
For $|f|\le g$ put $a=(g+f)/2$, $b=(g-f)/2\ge0$; then
$\E v(g)^{2r}-\E v(f)^{2r}=2\sum_{k\text{ odd}}\binom{2r}{k}\E[v(a)^{2r-k}v(b)^{k}]\ge0$ by the first
Griffiths inequality \ref{prem:C4}(ii), and comparing the series of
$\E\cosh(tv(f))=\E e^{tv(f)}$ and $\E\cosh(tv(g))$ gives the bound for $f$. Chernoff gives
$\nu_n(|v(f)|>z)\le2e^{-z^{2}/(2V_g)}$, and tail integration the moment bound. For
$n<N_g$, H\"older with $\|d\nu_{m,u,\Lambda_n}/dP_m\|_2\le\Xi_2(\Lambda_n)^{1/2}$
(Proposition~\ref{prop:identify}(i)) and the Gaussian moments with
$\Var_{P_m}\psi(f)\le m^{-2}\|g\|_2^{2}$ give the second entry of $C_{p,g}$.

\emph{(iii)} Let $a_{j,k}(v)=v(\chi_j\conj{e_{j,k}})$ and $g_j=|\T_j|^{-1/2}\chi_j$; the real and
imaginary parts of $\chi_j\conj{e_{j,k}}$ are dominated by $g_j$. Choose an integer
$N_j\ge2$ with $\supp\chi_j\subset\Lambda_n$ for $n\ge N_j$, and use $N_{g_j}=N_j$ in (ii).
For $n\ge N_j$, (ii) gives $\E_{\nu_n}|a_{j,k}|^{2}\le2V_{g_j}$; for $n<N_j$,
$\E_{\nu_n}|a_{j,k}|^{2}\le\sqrt3\,\Xi_2(\Lambda_n)^{1/2}m^{-2}\|\chi_j\|_2^{2}/|\T_j|$. Hence
$\sup_{n,k}\E_{\nu_n}|a_{j,k}|^{2}\le A_j/|\T_j|$ with
$A_j=\max\{2V_{g_j}|\T_j|,\sqrt3m^{-2}\|\chi_j\|_2^{2}\max_{n<N_j}\Xi_2(\Lambda_n)^{1/2}\}$. The shell
$|k|_\infty=r$ has $8r$ points with $\lambda_{j,k}\ge\pi^{2}r^{2}/(j+1)^{2}$, so
$S_{j,s}=\sum_k\lambda_{j,k}^{-s}<\infty$ exactly when $s>1$, and Tonelli gives
$\sup_n\E_{\nu_n}\|\chi_jv\|^{2}_{H^{-s}(\T_j)}\le B_{j,s}:=A_jS_{j,s}/|\T_j|$. For $p\ge2$, (ii)
gives $\|a_{j,k}\|_{L^{p}(\nu_n)}\le2C_{p,g_j}^{1/p}$ uniformly, and Minkowski in $L^{p/2}$
bounds $\E_{\nu_n}\|\chi_jv\|^{p}_{H^{-s}}$ by $2^{p}C_{p,g_j}S_{j,s}^{p/2}$; for $p<2$ use Jensen.

Fix $\epsilon>0$ and $K_\epsilon=\{v\in X:\|\chi_jv\|^{2}_{H^{-3/2}(\T_j)}\le2^{j}(B_{j,3/2}+1)/\epsilon\ \forall j\}$.
Each extended norm is a supremum of continuous finite sums, hence lower
semicontinuous, so $K_\epsilon$ is closed; Markov gives $\nu_n(X\setminus K_\epsilon)\le\epsilon$. Let
$(v_l)\subset K_\epsilon$. For each $j$ the localized sequence is bounded in $H^{-3/2}(\T_j)$,
and the embedding into $H^{-2}(\T_j)$ is compact (finite Fourier truncations, and
a tail bounded by the $H^{-3/2}$ radius times $\sup\lambda^{-1/2}$ on the tail); a
diagonal subsequence makes $(\chi_jv_l)_l$ convergent in $H^{-2}(\T_j)$ for every $j$,
hence $d$-Cauchy, hence convergent in $X$ by completeness
(Proposition~\ref{prop:space}(ii)), with limit in $K_\epsilon$. Thus $K_\epsilon$ is compact
and $(\nu_n)$ is tight. By Prokhorov every subsequence has a further subsequence
converging weakly to some $\mu$; by continuity of evaluations and
Theorem~\ref{thm:LR}, $\mu$ has the finite-dimensional compact-test laws of $\nu^{S}$.

Let $J_X:X\to\R^{\N}$ list the real and imaginary parts of all $a_{j,k}$; it is
continuous and injective (Proposition~\ref{prop:space}(iii)), so by Lusin--Souslin
$J_X(X)$ is Borel and $J_X^{-1}$ is Borel on it. Let $J_S$ be the same evaluations on
$\cS'$. Equal finite-dimensional laws give $(J_X)_*\mu=(J_S)_*\nu^{S}$, so
$\nu^{S}(J_S\in J_X(X))=1$. Put $\mathcal R=J_X^{-1}\circ J_S$ on that set. For $\psi$ in it,
$\mathcal R(\psi)$ and $\psi$ have the same localized Fourier coefficients, and the
expansion in the proof of \eqref{eq:duality}, valid for every distribution, gives
$\mathcal R(\psi)(f)=\psi(f)$ for compact smooth $f$; and $\mathcal R_*\nu^{S}=\mu$. Since
$\cB(X)=\sigma(J_X)$, $(J_X)_*\mu$ determines $\mu$; all subsequential limits coincide, and
$\nu_n\Rightarrow\nu:=\mathcal R_*\nu^{S}$. Lower semicontinuity and the portmanteau inequality
transfer the $H^{-s}$ bounds to $\nu$.

\emph{(iv)} Choose a real nonnegative compact smooth $g$ with $|f_i|\le g$ for all
$i$. For a polynomial of degree $d$, $|Q(x)|\le c_Q(1+\sum_i|x_i|^{d})$, so (ii) bounds
$\sup_n\E_{\nu_n}|Q|^{q}$ for every $q>0$. The finite-vector laws converge, the target
has the corresponding moments by the portmanteau inequality, and continuous
truncation at level $T$ with tails at most a constant times $T^{1-q}$ ($q>1$),
uniformly in $n$ and at the target, gives the convergence.
\end{proof}

\begin{remark}[Scope of Theorem~\ref{thm:V1}]\label{rem:V1-scope}
Part (i) holds for all $m,u>0$. Every volume-uniform statement uses the weak-coupling
regime only through the tightness supplied by Theorem~\ref{thm:LR}. The early
volumes use their individual finite density norms; no bound uniform in $n$ on
the densities $d\nu_{m,u,\Lambda_n}/dP_m$ is used. The volume limit is classical
continuum input to the construction, not a neural-estimate thermodynamic limit.
\end{remark}

\subsection{Target properties and reconstruction}\label{sec:CT}

Clause~(D) of Theorem~\ref{thm:main} identifies the limiting correlations with Schwinger distributions, and reconstruction \ref{prem:C3}(i) produces a relativistic quantum field from a measure satisfying four conditions: Euclidean invariance, time-reflection positivity, an exponential regularity bound and translation ergodicity. Theorem~\ref{thm:CT} verifies them in parts (i), (ii), (iv) and (iii), proves temperedness in (v) and non-Gaussianity in (vi), and applies reconstruction only afterwards, in (vii). Non-Gaussianity is the sense in which the reconstructed field is interacting (Remark~\ref{rem:CT-order}).

Throughout, $(m,u)$ is in the regime, $s$, $\lambda$, $v$ and $\nu^{S}=(T_s)_*v$ are as in
Theorem~\ref{thm:LR}, and $M=M_{\mathrm{ref}}$. Time is $z_2$; $\Theta$ is the reflection
$z_2\mapsto-z_2$ acting on observables, $\Sigma_+=\cC_{\{z_2>0\}}$, and $\tau_r$ shifts an
observable by $r$ in the positive time direction.

\begin{lemma}[Determining family]\label{lem:determining}
There is a countable $D\subset C_c^\infty(\R^2;\R)$, independent of all
parameters and dense in $\cS(\R^2;\R)$, with
$\sigma\{T\mapsto T(d):d\in D\}=\cC$. Two probabilities on $\cC$ with
the same finite-dimensional $D$-laws coincide.
\end{lemma}
\begin{proof}
Fix real smooth $\eta,\zeta$ with $\eta=1$ on $[-1,1]^2$,
$\supp\eta\subset(-2,2)^2$, $\zeta=1$ on $[-2,2]^2$ and
$\supp\zeta\subset(-3,3)^2$. Put $\eta_n(x)=\eta(x/n)$ and
$\zeta_n(x)=\zeta(x/n)$. On the numerical torus with cell $(-3n,3n)^2$
let $Q_n$ be the real trigonometric polynomials with rational coefficients
in the sine/cosine basis of frequencies $\pi k/(3n)$, $k\in\Z^2$. Then
$D=\bigcup_{n\ge1}\{\zeta_n P:P\in Q_n\}$ is countable, real compact
smooth and independent of the model and localization parameters.
For every real Schwartz $f$, the product rule and decay outside $[-n,n]^2$
give $\|(1-\eta_n)f\|_{(N)}\le C_N n^{-1}\|f\|_{(2N+1)}$ for $N\ge1$.
Thus $\eta_n f\to f$ in $\cS$. At fixed $n$, the zero periodic extension
of $\eta_n f$ is smooth; Fourier sums and rational approximation of their
finite coefficients approximate it in $C^\infty$. Multiplication by
$\zeta_n$ fixes a compact support and the limit $\zeta_n\eta_n f=\eta_n f$.
Choose $P_n$ so that $\zeta_nP_n-\eta_nf$ is less than $1/n$ in the first
$n$ Schwartz seminorms. Then $d_n=\zeta_nP_n\to f$ in $\cS$.
Every $T\in\cS'$ satisfies $T(f)=\lim_nT(d_n)$, so all Schwartz evaluations
are measurable for $\sigma(D)$; the reverse inclusion is immediate.
Finite $D$-cylinders form a generating $\pi$-system, so probability
uniqueness follows by $\pi$--$\lambda$.
\end{proof}

\begin{lemma}[Source-defined target and symmetries]\label{lem:target-symmetry}
The source-defined $v$ and $\nu^S=(T_s)_*v$ of Theorem~\ref{thm:LR} are
well defined from \ref{prem:C1}, \ref{prem:C2}, \ref{prem:C6} alone,
with the definitions of the admissible pair and dilation. They are invariant
under field sign and the central coordinate reflections, and $\nu^S$ is
invariant under the full Euclidean group. These conclusions hold under
each admitted reading of \ref{prem:C1}. No foundation input is used here.
\end{lemma}
\begin{proof}
At $M=M_{\mathrm{ref}}>M_0$ and
$0<\lambda=uM^2/m^2<\lambda_*\le\lambda_0$, \ref{prem:C2}(i),(ii)
make the two Gibbs extensions of the C1 states the same member $\mu$ of
$\cG_{\mathrm{ur}}$. Their cylinder restrictions therefore coincide as $v$.
This identity uses neither a finite-volume neural law nor a volume-limit
estimate. By \ref{prem:C1}(iii) the specified pressure is differentiable.
Lemma~\ref{lem:classical-wick} has independently proved the two provisos
that license the classical free-boundary convention. Apply \ref{prem:C1}(iv)
to those laws along centered axis-parallel square exhaustions.
For its characteristic reading, apply the assertion to $tg$ for every
real $t$; L\'evy's theorem gives the scalar weak limit. For the moment
reading, \ref{prem:C1}(ii) and Gaussian Cauchy--Schwarz give
\[
 \E_v e^{t\psi(g)}\le\|f_B\|_2
    e^{t^2\langle g,(-\Delta+M^2)^{-1}g\rangle}<\infty
\]
for every real compact smooth $g$, all real $t$, and bounded open $B$
containing its support. Moment determinacy and convergence then give that
same scalar weak limit. For the Laplace reading, the moment-generating
function theorem applies on a neighborhood of zero. Applying these arguments
to every real linear combination of a fixed tuple, and then the
finite-dimensional L\'evy theorem, gives joint compact-test convergence in
all three readings. This is exactly the classical portion of
Theorem~\ref{thm:LR}; the later common-space identification is not needed.

Lemma~\ref{lem:classical-wick} gives finite-volume field-sign and central
reflection invariance. Apply joint convergence to the union of a compact
tuple and its transformed tuple; the symmetries pass to the corresponding
$v$-laws. For real Schwartz tests, compact cutoffs converge in $\cS$ and
their evaluations converge pointwise on every tempered distribution.
Bounded convergence of characteristic functions extends the identities to
Schwartz cylinders; $\pi$--$\lambda$ gives equality on $\cC$.
The positivity readings of C1 are unused in this step, and the three
convergence readings were treated separately above.

Proper Euclidean invariance of $v$ is directly \ref{prem:C1}(i)(a).
The elementary identities $A_bT_s=T_sA_{sb}$ and $RT_s=T_sR$ of
Lemma~\ref{lem:dilation} transport it to $\nu^S$. That lemma's Gaussian and
Wick dilation identities use the direct construction of
Lemma~\ref{lem:classical-wick} and no foundation. Dilation commutes with
sign and $\theta_1,\theta_2$, so those symmetries transfer too. Proper
motions and one central reflection generate every Euclidean isometry.
In particular translation invariance has been proved before any translated
local-density norm is used.
\end{proof}

\begin{lemma}[Local densities of the target]\label{lem:localdensity}
For every bounded open $B$, the restriction of $\nu^S$ to $\cC_B$ has
density $F_B=f_{sB}\circ T_{1/s}$ with respect to $P_m$, and
$\|F_B\|_{L^p(P_m)}=\|f_{sB}\|_{L^p(P_M)}<\infty$ for $1\le p<\infty$.
For an open centered axis-parallel square,
\begin{equation}\label{eq:local-square-density}
 F_B=e^{-U^S_{m,u,B}-\alpha_\infty|sB|}g_B,
 \qquad g_B=g_{\partial(sB)}\circ T_{1/s}>0\quad P_m\text{-a.s.},
\end{equation}
where $g_B$ is measurable for either stated boundary convention, with that
convention retained. Moreover $F_{B+k}=F_B\circ A_k$ $P_m$-a.s., so their
finite $L^p$ norms agree. For real $f_1,\ldots,f_r$ supported in $B$ and
$t\in\R^r$, writing $q=\sum_i t_i f_i$,
\begin{equation}\label{eq:local-exponential}
 \E_{\nu^S}e^{\psi(q)}\le\|F_B\|_2
       e^{\langle q,\Cfree q\rangle}<\infty.
\end{equation}
\end{lemma}
\begin{proof}
$T_r$ is a cylinder-measurable linear bijection with inverse $T_{1/r}$;
$L_r$ acts continuously on $\cS$ and scales test supports from $U$ to $rU$.
For bounded $\cC_B$-measurable $H$, Gaussian dilation gives
\[
 \int HF_B\,dP_m
 =\int(H\circ T_s)f_{sB}\,dP_M
 =\int H\circ T_s\,dv=\int H\,d\nu^S.
\]
The same change of Gaussian law gives the exact equality of $p$th absolute
integrals, with no probability-norm Jacobian. For a square, compose the
C1(ii) factorization with $T_{1/s}$ and use
$U^S_{M,\lambda,sB}\circ T_{1/s}=U^S_{m,u,B}$ from the direct mollified
dilation in Lemma~\ref{lem:dilation}. The pressure factor is
$\alpha_\infty|sB|$ at mass $M$, and positivity of $g_B$ transports with
Gaussian null sets.

\emph{BC1.} For every open neighborhood $V$ of $\partial B$, $sV$ is an
open neighborhood of $\partial(sB)$. The boundary factor pulled back by
$T_{1/s}$ is therefore $\cC_V$-measurable for every such $V$, hence
measurable for their intersection.
\emph{BC2.} For real Schwartz $h$ and $r>0$,
\begin{equation}\label{eq:boundary-dilation}
 \|L_rh\|_{H^{-1}}^2
 =(2\pi)^{-2}\int\frac{|\hat h(q)|^2}{r^2+|q|^2}\,dq
 \le\max(1,r^{-2})\|h\|_{H^{-1}}^2.
\end{equation}
Thus dilation extends continuously on $H^{-1}$ and scales distribution
support by $r$. Gaussian evaluation at mass $M$ is continuous
$H^{-1}\to L^2(P_M)$, since its multiplier is at most
$\max(1,M^{-2})(1+|p|^2)^{-1}$; similarly at mass $m$.
For $h$ supported in $\partial(sB)$ and real Schwartz $h_j\to h$ in
$H^{-1}$, the exact evaluation identity
$\psi(h_j)\circ T_{1/s}=\psi(L_{1/s}h_j)$ passes in $L^2(P_m)$ to the
extended Gaussian evaluations. The transformed test is supported in
$\partial B$. Gaussian null sets also pull back to Gaussian null sets.
This proves transport of the completed BC2 algebra, without identifying it
with the uncompleted BC1 algebra.

\emph{Translated norms, after invariance.}
Lemma~\ref{lem:target-symmetry} already proves translation invariance of
$\nu^S$; the Gaussian law is stationary by its covariance. For bounded
$\cC_{B+k}$-measurable $H$,
\[
 \int H(F_B\circ A_k)\,dP_m
 =\int(H\circ A_{-k})F_B\,dP_m
 =\int H\circ A_{-k}\,d\nu^S=\int H\,d\nu^S.
\]
Local Radon--Nikodym uniqueness gives $F_{B+k}=F_B\circ A_k$ and hence
the exact norm equality. No part of the preceding invariance proof used it.
Finally monotone convergence extends the local density identity to the
nonnegative exponential. Cauchy--Schwarz gives
$\E F_Be^{\psi(q)}\le\|F_B\|_2(\E_{P_m}e^{2\psi(q)})^{1/2}
=\|F_B\|_2e^{\langle q,\Cfree q\rangle}$, with
$\langle q,\Cfree q\rangle\le m^{-2}\|q\|_2^2$.
\end{proof}

\begin{lemma}[Local $L^2$ kernels]\label{lem:kernels}
If a probability $\pi$ on $(\cS',\cC)$ has a density
$F\in L^2(P_m)$ on $\cC_B$, $B$ bounded open, then for every integer
$n\ge1$ its local moment and truncated kernels on $B^n$ are represented
by real functions in $L^2(B^n)$, with complex-linear test pairings.
\end{lemma}
\begin{proof}
Choose the local measurable version $F\ge0$ with $\E_{P_m}F=1$.
Local moments exist by Cauchy--Schwarz and Gaussian polynomial moments;
no moment outside $B$ is required. Plancherel and polar integration give
\begin{equation}\label{eq:local-covariance-L2}
 \|\Cfree\|_{L^2(\R^2)}^2=\frac1{4\pi m^2},\qquad
 R_B^2:=\|\Cfree(x-y)\|_{L^2(B^2)}^2\le\frac{|B|}{4\pi m^2}.
\end{equation}
The restricted covariance operator $K$ on real $L^2(B)$ is positive
self-adjoint with $\|K\|\le m^{-2}$ by its Fourier form; its integral
kernel is $\Cfree(x-y)$ by Plancherel, first on smooth tests and then by
continuity. Empty $B$ is vacuous, so assume $|B|>0$.

For a finite sum $Q$ of real compact smooth product tests on $B^k$,
define $W_k(f_1\otimes\cdots\otimes f_k)={:}\psi(f_1)\cdots\psi(f_k){:}$
linearly. The Wick covariance identity yields
\begin{equation}\label{eq:local-Wick-map}
 \|W_k(Q)\|_{L^2(P_m)}^2
 =k!\langle\mathrm{Sym}\,Q,K^{\otimes k}\mathrm{Sym}\,Q\rangle
 \le k!m^{-2k}\|Q\|_{L^2(B^k)}^2.
\end{equation}
This also proves independence of the product-sum representation. Density
extends $W_k$ boundedly to $L^2(B^k)$. The functional
$Q\mapsto\E_{P_m}[FW_k(Q)]$ has norm at most
$\sqrt{k!}m^{-k}\|F\|_2$, so real $L^2$ duality produces a symmetric
$h_k\in L^2(B^k)$ representing it, with that norm bound. Normalize the
finite Lebesgue measure to a probability when applying C6 duality, then
undo the normalization; the displayed bound is unchanged. Put $h_0=1$.

For a partial pairing $P$ of $\{1,\ldots,n\}$, write $U(P)$ for the
unpaired indices. The ordinary moment kernel is
\begin{equation}\label{eq:local-moment-kernel}
 M_n(x)=\sum_P\left(\prod_{\{i,j\}\in P}\Cfree(x_i-x_j)\right)
              h_{|U(P)|}(x_{U(P)}).
\end{equation}
Each factor uses a disjoint set of coordinates, so its $L^2$ norm is
$R_B^{|P|}\|h_{|U(P)|}\|_2$. In particular
\[
 \|M_n\|_2\le\|F\|_2\sum_{p=0}^{\lfloor n/2\rfloor}
 \frac{n!}{2^pp!(n-2p)!}R_B^p\sqrt{(n-2p)!}\,m^{-(n-2p)}.
\]
To identify it, expand
$e^{\sum_i t_iX_i}=e^{\frac12\sum_{i,j}t_it_j\E X_iX_j}
 e^{\sum_i t_iX_i-\frac12\sum_{i,j}t_it_j\E X_iX_j}$ for the finite
Gaussian tuple $X_i=\psi(f_i)$ and compare the coefficient $t_1\cdots t_n$.
It gives each partial pairing once, with the remaining factors Wick ordered.
Multiply by $F$, integrate and use the representation of $h_k$ to obtain
$\E_\pi\prod_i\psi(f_i)=\int M_n\prod_i f_i$ for product tests.
Every integral is absolute by the displayed $L^2$ bounds. Product-test
density gives the unique distribution on $B^n$.

For a set partition $\pi_0$ with blocks $A$, define
\begin{equation}\label{eq:local-cumulant-kernel}
 M_n^T(x)=\sum_{\pi_0}(-1)^{|\pi_0|-1}(|\pi_0|-1)!
                    \prod_{A\in\pi_0}M_{|A|}(x_A).
\end{equation}
Every block product is an $L^2$ tensor product on disjoint coordinates,
of norm $\prod_A\|M_{|A|}\|_2$. The finite cumulant partition identity
and Fubini identify its product-test pairing with the joint cumulant.
Density again gives uniqueness. Real kernels extend complex-linearly,
as required; no point field of the unregulated distribution is used.
\end{proof}


\begin{theorem}[Target properties; reduced lemma CT-R]\label{thm:CT}
\begin{enumerate}[label=(\roman*)]
\item $\nu^{S}$ is invariant under the full Euclidean group of $\R^{2}$ and under
  $\psi\mapsto-\psi$.
\item For every line $L$, with $\Theta_L$ the reflection in $L$ and $\Sigma^{L}_{+}$ generated by
  smearings supported in one open half-plane bounded by $L$:
  $\int\conj{\Theta_LF}\,F\,d\nu^{S}\ge0$ for every $F\in L^{2}(\Sigma^{L}_{+},\nu^{S})$ and every
  $\Sigma^{L}_{+}$-measurable $F$ with $\conj{\Theta_LF}F$ integrable.
\item For $G,G'\in L^{2}(\cS',\cC,\nu^{S})$, $\E[G\,(G'\circ A_b)]\to\E G\,\E G'$ as $|b|\to\infty$
  (strong mixing). Consequently every $E\in\cC$ with $\nu^{S}(E\triangle A_b^{-1}E)=0$ for all
  $b\in\R^{2}$ has $\nu^{S}(E)\in\{0,1\}$ (translation ergodicity), and
  $\E[\conj{\Theta F}\,\tau_rG]\to\conj{\E F}\,\E G$ for $F,G\in L^{2}(\Sigma_+)$.
\item With a real nonnegative smooth partition of unity $\sum_{k\in\Z^{2}}\vartheta_k=1$,
  $\vartheta_k=\vartheta_0(\cdot-k)$, $\cA(f)=\sum_k\|\vartheta_kf\|_2$ is a norm continuous on $\cS$, and
  with $K=\|F_B\|_{L^{2}(P_m)}\ge1$ for a square $B\supseteq\supp\vartheta_0$ and
  $D=3^{1/4}K^{1/2}/m$: $\E e^{t\psi(f)}\le K\exp(t^{2}\cA(f)^{2}/m^{2})$,
  $\|\psi(f)\|_{L^{2}}\le D\cA(f)$, and
  $|\E e^{\psi(f)}-\E e^{\psi(g)}|\le2DK^{1/2}e^{2/m^{2}}\cA(f-g)$ for $\cA(f),\cA(g)\le1$. Hence
  condition (c) of \ref{prem:C3}(i) holds with $|||\cdot|||=\cA$.
\item For every $k$ there is a unique $\Schw_k\in\cS'(\R^{2k})$ with
  $\Schw_k(f_1\otimes\cdots\otimes f_k)=\int\prod_r\psi(f_r)\,d\nu^{S}$ for $f_r\in\cS(\R^{2};\C)$ (complex-linear
  extension); it is symmetric and real on real tests.
\item $\nu^{S}$ is not Gaussian; its local moment and truncated kernels lie in
  $L^{2}(B^{n})$ for bounded $B$; and some even truncated Schwinger distribution of
  order at least four is nonzero on real tests supported in pairwise disjoint
  time slabs.
\item By \ref{prem:C3}(i), applied with (i), (ii), (iv) and (iii): the moments of
  $\nu^{S}$ are the Schwinger functions of a unique relativistic quantum field theory
  in $1+1$ dimensions satisfying all the Wightman axioms (temperedness,
  Poincar\'e invariance, positivity, spectrum condition, locality, clustering),
  with a unique physical vacuum; its field fails Gaussian factorization.
\end{enumerate}
\end{theorem}

\begin{proof}
All parts of this proof use only \ref{prem:C1}--\ref{prem:C3} and
\ref{prem:C6}, with the source-defined target identity. The proofs of
Lemmas~\ref{lem:classical-wick}, \ref{lem:dilation},
\ref{lem:target-symmetry}, \ref{lem:localdensity} and \ref{lem:kernels}
have supplied the needed classical objects directly. Their use here does
not import the separate common-space Gaussian/Wick foundations.

\phantomsection\label{proof:CT-i}
\emph{(i), invariance.} This is Lemma~\ref{lem:target-symmetry}, proved
without translated-density norms, local kernels, mixing or reconstruction.

\phantomsection\label{proof:CT-ii}
\emph{(ii), full measurable reflection positivity.}
Under the measure reading C1(i)(b$_M$), the $L^2$ time-reflection inequality
for $v$ transfers directly by positive dilation: $T_s$ commutes with
$\Theta$ and its pullback preserves the positive-half-plane algebra.
Under the weaker reading C1(i)(b$_P$), positive dilation first transfers
exactly the inequality for polynomials whose individual monomials have
strictly ordered positive-time supports. Reflection is an $L^2(\nu^S)$
isometry by (i).

Let $B$ be bounded open and let $V_n(G)$ be the polynomial associated to
a finite sum $G$ of compact smooth product tests on $(B\cap H_+)^n$.
The local $2n$-point moment kernel of Lemma~\ref{lem:kernels} gives
\begin{equation}\label{eq:RP-polynomial-map}
 \|V_n(G)\|_{L^2(\nu^S)}^2
 =M_{2n}(\conj G\otimes G)
 \le\|M_{2n}\|_{L^2(B^{2n})}\|G\|_{L^2(B^n)}^2.
\end{equation}
Thus this map extends boundedly from product sums to $L^2((B\cap H_+)^n)$.
Products of real smooth tests with pairwise disjoint positive-time supports
span a dense subspace there. Indeed the union of equal-time hyperplanes
has Lebesgue measure zero; its complement has a countable cover by
relatively compact product boxes with disjoint positive time projections.
A function orthogonal to all smooth products on those boxes vanishes
almost everywhere on each box, hence on the whole complement. For each
such product, reorder the commuting field factors into increasing time
order. It is then exactly a monomial admitted in (b$_P$).
The bounded map above therefore approximates every compact
positive-half-plane polynomial in $L^2(\nu^S)$ by admitted polynomials.

Every finite vector of compact evaluations has finite exponential moments
in all real directions by \eqref{eq:local-exponential}. The polynomial
density theorem of C6 makes its polynomials dense in its $L^2$ space.
For every positive-half-plane Schwartz test, compact spatial cutoffs
converge in $\cS$, keep support in the open half-plane, and converge in
evaluation pointwise on $\cS'$. Thus compact evaluations generate the
stated $\Sigma_+$. Countable dependence supplies a countable family for
each observable; conditional-expectation approximation for increasing
finite families and truncation now make the compact polynomials dense
in $L^2(\Sigma_+,\nu^S)$. If $F_j\to F$ there, reflection isometry and
Cauchy--Schwarz give
\[
 \left|\E\conj{\Theta F_j}F_j-\E\conj{\Theta F}F\right|
 \le\|F_j-F\|_2(\|F_j\|_2+\|F\|_2).
\]
Positivity therefore holds for the full $L^2$ class under both readings.
For a measurable $F$ whose reflected product is integrable, use
$F_j=F\ind_{\{|F|\le j\}}$: the reflected products converge pointwise
and are dominated by $|\Theta F|\,|F|$. Bounded positivity and dominated
convergence prove the additional integrable-form clause. For any line,
conjugate the time-line inequality by a Euclidean isometry using (i).
This concerns the target; the finite neural stages retain only their two
stated central coordinate lines.

\phantomsection\label{proof:CT-iii}
\emph{(iii), Euclidean mixing and ergodicity before reconstruction.}
At mass $M$ use the unique Gibbs extension $\mu$ of $v$ supplied by C2.
Let $D\in F_0$ be its increasing square family and
$K_DF(\eta)=E^\eta_{D^c}F$. For an integrable local polynomial $F$,
integrability follows from C1(ii) and Gaussian moments, and the Gibbs
identity gives $\mu(K_DF)=\mu(F)$. If $H$ is bounded and
$\Sigma_{D^c}$-measurable, the kernel's point-fixing property applied to
$H$ and $|H|^2$ shows that $H$ equals the constant $H(\eta)$ under each
kernel probability (apply to real and imaginary parts). Hence
$K_D(HF)=H K_DF$ wherever the integral of $|F|$ is finite. The Gibbs
identity for $|F|$ gives that finiteness $\mu$-a.e.; its identity for
$HF$ yields
\begin{equation}\label{eq:mixing-kernel-factor}
 \mu(HF)=\mu(HK_DF).
\end{equation}
For a real local polynomial $F$, C2(iii) with both Gibbs measures equal
to $\mu$ and Jensen imply
\begin{equation}\label{eq:mixing-kernel-L1}
 \|K_DF-\mu F\|_1
 \le\iint|K_DF(\eta)-K_DF(\eta')|\,\mu(d\eta)\mu(d\eta')\longrightarrow0.
\end{equation}
Complex polynomials follow by real and imaginary parts. Thus
\[
 |\mu(FH)-\mu F\,\mu H|\le\|H\|_\infty\|K_DF-\mu F\|_1.
\]
For a fixed bounded compact cylinder $G$, $H=G\circ A_b$ is
$\Sigma_{D^c}$-measurable for all sufficiently large $|b|$: its finitely
many translated compact supports lie outside $\bar D$, hence in every
open neighborhood used in the closed-set outside algebra. This condition
holds uniformly over the direction of $b$. First choose $D$ sufficiently
far along $F_0$, then let $|b|$ grow. Translation invariance of $v$ gives
$\mu(G\circ A_b)=\mu G$, so mixing holds for each such $F,G$.
Local polynomials are dense in $L^2(v)$ by the local exponential moments,
finite-dimensional polynomial density and increasing finite-coordinate
conditional expectations. Compact evaluations generate $\cC$ by
Lemma~\ref{lem:determining}. Bounded compact cylinders are also dense by
the same conditional expectations and truncation. Translations preserve
$L^2(v)$ norms, so Cauchy--Schwarz extends the limit to all complex
$G,G'\in L^2(v)$, in every direction.

The identity $A_bT_s=T_sA_{sb}$ transfers this strong mixing to $\nu^S$,
since $s>0$ implies $|sb|\to\infty$. For an invariant event $E$ as in
the statement, take both observables to be $\ind_E$. The mixing product
equals $\nu^S(E)$ for every $b$ and tends to $\nu^S(E)^2$; its probability
is therefore zero or one. The half-plane clustering formula follows by
taking the first observable $\conj{\Theta F}$ and translating the second
in positive time; (i) gives $\E\Theta F=\E F$.
No reconstruction theorem, tube uniqueness, spectrum condition,
unique-vacuum assertion, FS Theorem~4.4 or OS-energy identification has
entered this argument. Condition C3(i)(d) is established at this point.

\phantomsection\label{proof:CT-iv}
\emph{(iv), the exact norm and exponential regularity.}
Fix a real nonnegative $\vartheta_0\in C_c^\infty$ with
$\sum_{k\in\Z^2}\vartheta_0(\cdot-k)=1$, put
$\vartheta_k=\vartheta_0(\cdot-k)$ and
$\cA(f)=\sum_k\|\vartheta_kf\|_2$. The partition identity gives a norm.
For each integer $N>2$, bounded support of $\vartheta_0$ gives
$\|\vartheta_kf\|_2\le c_N(1+|k|)^{-N}\|f\|_{(N)}$; summation proves
finiteness and Schwartz continuity. Bounded overlap of the translated
supports, the product rule and arbitrarily higher Schwartz decay show
that the partition sums converge to $f$ in every Schwartz seminorm.
Choose a bounded open square $B\supset\supp\vartheta_0$ and set
$K=\|F_B\|_2\ge\E F_B=1$, $D=3^{1/4}K^{1/2}/m$.
The translated-density norm equality has already been proved after
invariance. For real $g$ supported in $B+k$, the local density and free
Gaussian formulas give
\begin{equation}\label{eq:Anorm-local}
 \E e^{t\psi(g)}\le K e^{t^2\|g\|_2^2/m^2},\qquad
 \E\psi(g)^2\le K(\E_{P_m}\psi(g)^4)^{1/2}
                  \le\sqrt3 K m^{-2}\|g\|_2^2.
\end{equation}
Minkowski bounds each finite partition sum in $L^2$ by
$D\sum_k\|\vartheta_kf\|_2$. The same estimate for its tails makes it
Cauchy in $L^2$; its pointwise Schwartz-evaluation limit is $\psi(f)$.
Thus $\|\psi(f)\|_2\le D\cA(f)$ for every real Schwartz $f$.

Let $a_k=\|\vartheta_kf\|_2$ and, when $\cA(f)>0$, put
$\alpha_k=a_k/\cA(f)$ for the nonzero terms. For any finite set $J$,
generalized H\"older with weights $\alpha_k$ (including a factor one with
the unused weight) and \eqref{eq:Anorm-local} imply
\begin{align*}
 \E\exp\left(t\sum_{k\in J}\psi(\vartheta_kf)\right)
 &\le K^{\sum_J\alpha_k}
       \exp\left(t^2m^{-2}\sum_J a_k^2/\alpha_k\right)\\
 &\le K\exp(t^2\cA(f)^2/m^2).
\end{align*}
Fatou passes this to the pointwise full sum. The zero-norm case is
immediate. For $\cA(f),\cA(g)\le1$, use
$|e^x-e^y|\le|x-y|(e^x+e^y)$ and Cauchy--Schwarz, with
$\|e^{\psi(f)}\|_2\le K^{1/2}e^{2/m^2}$, to obtain exactly
\begin{equation}\label{eq:Anorm-continuity}
 |\E e^{\psi(f)}-\E e^{\psi(g)}|
 \le2DK^{1/2}e^{2/m^2}\cA(f-g).
\end{equation}
The unit ball has exponential bound $Ke^{1/m^2}$ and this Lipschitz
constant, so C3(i)(c) holds. For later use,
$\E e^{|\psi(f)|/\cA(f)}\le2Ke^{1/m^2}=:K'$ when $f\ne0$, whence
$\|\psi(f)\|_k\le(k!K')^{1/k}\cA(f)$ for integers $k\ge1$.

\phantomsection\label{proof:CT-v}
\emph{(v), tempered Schwinger distributions.}
For a $k$-linear form $\mathsf M$ on $\cS(\R^2;\C)^k$ with
$|\mathsf M(f_1,\ldots,f_k)|\le C\prod_i\|f_i\|_{(N)}$, there is a
unique tempered kernel. Here are the required reduction details. For
$k=1$ it is the definition. Inductively, at a fixed last test $\phi$,
apply the $(k-1)$-variable assertion to obtain $\mathsf L(\phi)$.
Kernel uniqueness makes $\mathsf L$ linear. For a tensor sum $\Psi_j$,
$\phi\mapsto\mathsf L(\phi)(\Psi_j)$ is continuous. Approximate an
arbitrary Schwartz $\Psi$ by these sums; the values converge pointwise
because each $\mathsf L(\phi)$ is a distribution. Banach--Steinhaus on
the Fr\'echet space $\cS(\R^2)$ makes the limiting functional continuous.
C3(iii), with $N=2(k-1)$ and $n=2$, now gives the joint tempered kernel.
Product-test density gives uniqueness.
Apply this to $\mathsf M_k(f_1,\ldots,f_k)=\E\prod_i\psi(f_i)$.
H\"older and (iv) give $|\mathsf M_k|\le k!K'\prod_i\cA(f_i)$ for real
tests, and $\cA(f)\le C_\cA\|f\|_{(3)}$. Splitting complex tests into
real and imaginary parts gives the required complex multilinear bound.
Kernel uniqueness transfers symmetry and reality of the moments to
$\Schw_k\in\cS'_\C(\R^{2k})$.

\phantomsection\label{proof:CT-vi-coordinate}
\emph{(vi), a local quartic conditional density excludes Gaussianity.}
Take an open centered axis-parallel square $B$ and nonzero real
$h\in C_c^\infty(B)$. Put
\[
 k=(-\Delta+m^2)h,\quad q=\int hk=\int(|\nabla h|^2+m^2h^2)>0,
 \quad X=\psi(k)/q,\quad \eta=\psi-Xh.
\]
Under $P_m$, $\Var X=1/q$ and
$\mathrm{Cov}(X,\psi(f))=q^{-1}\int hf$, so
$\mathrm{Cov}(X,\eta(f))=0$ for every real Schwartz $f$.
Gaussianity and the monotone-class theorem make $X$ independent of
the full residual evaluation algebra. Its law is
$\gamma_q(dx)=(q/(2\pi))^{1/2}e^{-qx^2/2}dx$.

The boundary factor $g_B$ in \eqref{eq:local-square-density} has a
residual-measurable version. For BC1 choose a boundary neighborhood
disjoint from $\supp h$; each generating evaluation there is unchanged
by replacing $\psi$ by $\eta$. For BC2, boundary-supported
$f\in H^{-1}$ has $\langle f,h\rangle=0$. If real Schwartz
$f_j\to f$ in $H^{-1}$, then $\langle f_j,h\rangle\to0$ by
$h\in H^1$, so the defining $L^2(P_m)$ limit of $\psi(f_j)$ is also
the limit of $\eta(f_j)$. Gaussian completion changes versions only
on null sets. This covers each stated boundary algebra separately.

For the fixed mollifier, $\psi_\delta=\eta_\delta+Xh_\delta$,
$h_\delta=\rho_\delta*h$. Let $H_j(x;1/q)$ be the monic Hermite
polynomials for $\gamma_q$. At each positive mollifier scale,
\[
 J_\delta=u\int_B{:}\psi_\delta^4{:}
       =\sum_{j=0}^4 A_{j,\delta}(\eta)H_j(X;1/q).
\]
Its residual coefficients are square integrable by finite Gaussian
moments and Minkowski on bounded $B$. Independence and Hermite
orthogonality give
\begin{equation}\label{eq:interior-Hermite}
 \|J_\delta-J_\epsilon\|_2^2
 =\sum_{j=0}^4j!q^{-j}\|A_{j,\delta}-A_{j,\epsilon}\|_2^2.
\end{equation}
The direct $L^2$ Wick limit therefore gives residual coefficients $A_j$
with
\begin{equation}\label{eq:interior-quartic}
 U^S_{m,u,B}=\sum_{j=0}^4A_j(\eta)H_j(X;1/q),\qquad
 A_4=u\int_Bh^4=:c>0.
\end{equation}
This is an actual degree-four polynomial in the scalar $X$, with
deterministic leading coefficient $c$; it is obtained without multiplying
unregulated distributions.

Suppose $\nu^S$ were Gaussian. The local measurable density $F_B$
depends on a countable family $\psi(f_j)$, $f_j\in C_c^\infty(B;\R)$;
include $k$ in this list. Write $Y=(\eta(f_j))_j$ and
$F_B=b(X,Y)$. Both $X$ and each $Y_j$ are local linear evaluations.
Under $P_m$, $X$ is independent of $Y$, and the target law of $(X,Y)$
has density $b(x,y)$ relative to $\gamma_q(dx)P_Y(dy)$, by the local
Radon--Nikodym identity. The full residual variables used by
$A_0,\ldots,A_4,g_B,Y$ depend on a countable residual family too.
They need not all be measurable for the smaller local family $Y$.
Independence of $X$ from that full residual family and Fubini give, for
almost every full residual configuration $v_0$,
\begin{equation}\label{eq:conditional-quartic-log}
 \log b(x,Y(v_0))
 =-\sum_{j=0}^4A_j(v_0)H_j(x;1/q)
   -\alpha_\infty|sB|+\log g_B(v_0)
\end{equation}
for Lebesgue-a.e. $x$. This uses the square factorization, not a global
target density. In particular $b>0$ almost everywhere. The integral
$\int b(x,y)\gamma_q(dx)$ is positive and finite for $P_Y$-a.e. $y$
by positivity and $\E F_B=1$; the target marginal of $Y$ is equivalent
to $P_Y$. The conditional target density of $X$ at such $y$ is
proportional to $b(x,y)e^{-qx^2/2}$.

Under the assumed Gaussian target, every finite vector $(X,Y_1,\ldots,Y_j)$
is Gaussian. Finite Gaussian regression gives conditional characteristic
functions $e^{itm_j-t^2v_j/2}$, where
$m_j=\E_{\nu^S}[X\mid Y_1,\ldots,Y_j]$ converges in $L^2$ by the
increasing-filtration theorem and the deterministic $v_j\ge0$ decrease.
Applying the same conditional-expectation theorem to $e^{itX}$ identifies
the conditional characteristic function given $Y$ as
$e^{it\E[X\mid Y]-t^2v/2}$, $v=\lim v_j\ge0$, first for rational $t$
and then by continuity. Thus it is a Gaussian conditional law, possibly
degenerate. Its exceptional $Y$ set is null for $P_Y$ by equivalence.
Intersect with the full-measure residual set in
\eqref{eq:conditional-quartic-log}. A positive Lebesgue density excludes
zero variance. For positive variance its log density is quadratic, whereas
\eqref{eq:conditional-quartic-log} plus $-qx^2/2$ has leading coefficient
$-c<0$ in degree four. Equality a.e. would make the polynomials identical,
which is impossible. This proves non-Gaussianity while preserving the
local/full-residual distinction and the fixed boundary trace.

\phantomsection\label{proof:CT-vi-witness}
The local moment and truncated $L^2$ kernels for every bounded open $B$
and every order are Lemmas~\ref{lem:localdensity} and \ref{lem:kernels}.
Field-sign invariance annihilates every odd joint cumulant. If all even
compact cumulants of order at least four vanished as well, the
moment-cumulant formula would leave only pair partitions. Local exponential
moments and moment determinacy would make every compact linear evaluation
centered Gaussian, hence every finite compact vector jointly Gaussian.
For any real Schwartz test take compact cutoffs converging in $\cS$;
their evaluations converge pointwise on $\cS'$, hence weakly in law.
Centered Gaussian laws are closed under such weak convergence: unbounded
variances would force a discontinuous limiting characteristic function at
zero; a bounded-variance subsequence has a Gaussian limit. Apply this to
each real linear combination to obtain Gaussian Schwartz vectors,
contradicting the preceding argument. This passage needs no global neural
Schwartz estimate.

Consequently some even $n\ge4$ has a nonzero compact joint cumulant.
Choose bounded open $B$ containing its supports. Its real kernel
$M_n^T\in L^2(B^n)$ is nonzero. The union of hyperplanes
$x_{i,2}=x_{j,2}$ is Lebesgue null. The complement has a countable cover
by relatively compact product boxes $U_1\times\cdots\times U_n$ whose
time projections lie in pairwise disjoint bounded open intervals.
At least one box has nonzero kernel restriction; otherwise the squared
integral would vanish on that countable cover. Density of real smooth
product tests in its $L^2$ space then gives individual
$f_i\in C_c^\infty(U_i;\R)$ with
\begin{equation}\label{eq:disjoint-time-witness}
 S_n^T(f_1,\ldots,f_n)=\int M_n^T(x)\prod_i f_i(x_i)\,dx\ne0.
\end{equation}
Thus the witness is a real tensor product in disjoint time slabs at the
same even order $n\ge4$, rather than only a general joint test or a
distribution on coincident times.

\phantomsection\label{proof:CT-vii}
\emph{(vii), reconstruction and its precise non-Gaussian consequence.}
For $\nu^S$ on $(\cS',\cC)$, condition (a) of C3(i) is (i), condition
(b) is the full $L^2$ time-line assertion of (ii) under either positivity
reading, condition (c) is the exact norm and estimates of (iv), and
condition (d) was proved by Euclidean mixing in (iii).
Apply C3(i) only now. It gives the stated unique scalar relativistic
theory in $1+1$ dimensions, the listed Wightman axioms and unique physical
vacuum; (v) identifies its tempered symmetric Schwinger representatives.

Only after this application use C3(ii),(ii'). On a common ordered tube,
form the truncated holomorphic $n$-point function by the finite cumulant
partition sum of the holomorphic moment functions, preserving the order
within each block. Every factor is holomorphic there and uses disjoint
variables from the other block factors. Its distribution boundary is the
truncated Wightman distribution: the block products are tensor products,
so no undefined multiplication of coincident-variable distributions is
involved. At pairwise distinct Euclidean arguments its values are the
truncated Schwinger function by C3(ii) and the same partition identity.
If its Wightman boundary vanished, C3(ii') would force the holomorphic
function to vanish and hence its analytic continuation at those Euclidean
points. Permute the disjoint time slabs of
\eqref{eq:disjoint-time-witness} into one strict order; symmetry of the
Schwinger kernels preserves that nonzero pairing. Vanishing analytic
values would make the pairing zero, a contradiction. The truncated
Wightman distribution at this same even order $n\ge4$ is therefore
nonzero. In particular the field fails Gaussian factorization at some
order at least three, exactly the claimed consequence. No scattering or
quantitative mass-gap conclusion follows from this argument.
\end{proof}


\begin{remark}[Order of the argument]\label{rem:CT-order}
Ergodicity, condition (d), is established in (iii) by a Euclidean mixing
argument from the Gibbs identity and the boundary-condition independence of
\ref{prem:C2}, before \cite[Thm.~1.0]{FS1977} is applied. The spectrum condition,
clustering in the Wightman sense and uniqueness of the vacuum are consequences
of the reconstruction, not inputs to (d). In particular neither a unique-vacuum
statement nor an ergodic form of the reconstruction theorem is used to prove
ergodicity. A historical route through \cite[Thm.~4.4]{FS1977} and an
identification of the Osterwalder--Schrader energy with a Wightman energy is
recorded, with its first unsupported step, in Appendix~\ref{app:historical}; it is
not used. The main theorem uses only part (v) of Theorem~\ref{thm:CT}, in clause
(D). The word ``interacting'' means the non-Gaussian quantum field of (vii);
no scattering or quantitative mass-gap statement is made.
\end{remark}

% ===========================================================================
\section{Local moments and the common sequence}\label{sec:H1}

To turn the iterated limits into one sequence serving all later observables of Theorem~\ref{thm:main}, each arrow is first upgraded to carry all moments of the local seminorms (Proposition~\ref{prop:H1arrows}), using uniform moment bounds and Lemma~\ref{lem:transfer}. Proposition~\ref{prop:H1seq} then chooses at stage $n$ parameters that make each of four discrepancies smaller than $2^{-n}/4$, and the triangle inequality gives \eqref{eq:mainseq}. Since each compact smearing is controlled by one seminorm through \eqref{eq:duality}, clause~(B) follows (Proposition~\ref{prop:H1later}).

The moment-transfer and selection arguments use the field-space, Gaussian, Wick, fixed-volume and continuum-volume results proved above, including Theorem~\ref{thm:V1}. Their provenance is recorded in Appendix~\ref{app:provenance}.

Throughout, $S_j=\sum_{k\in\Z^{2}}\lambda_{j,k}^{-2}$, which is finite because the shell
$|k|_\infty=r\ge1$ has $8r$ points with $\lambda_{j,k}\ge1+\pi^{2}r^{2}/(j+1)^{2}$; thus
$B_j=\|\chi_j\|_2^{2}S_j/|\T_j|$ in Theorem~\ref{thm:G-versions}(e). Put
\[
K_j=\|\chi_j\|_2/2,\qquad\tau_j=m^{-2}B_j,\qquad
\sigma_j=\frac{\|\chi_j\|_1^{2}S_j}{4|\T_j|},\qquad
L_p=\sum_{q=1}^{p}q^{2p},\qquad (2p-1)!!=\prod_{i=1}^{p}(2i-1).
\]

\subsection{Moment bounds for the reference fields}

\begin{lemma}[Hilbert sums]\label{lem:hilbertsum}
Let $(a_k)_{k\in\Z^{2}}$ be complex random variables, $w_k\ge0$ and $p\ge1$. Then
$\E\bigl(\sum_kw_k|a_k|^{2}\bigr)^{p}\le\bigl(\sum_kw_k\|a_k\|_{2p}^{2}\bigr)^{p}$. In particular,
for a Borel random element $\Xi$ of $X$, with $a_k=\Xi(\chi_j\conj{e_{j,k}})$,
$\E p_j(\Xi)^{2p}\le\bigl(\sum_k\lambda_{j,k}^{-2}\|\Xi(\chi_j\conj{e_{j,k}})\|_{2p}^{2}\bigr)^{p}$.
\end{lemma}

\begin{proof}
Minkowski's inequality in $L^{p}$ for finite partial sums,
$\||a_k|^{2}\|_p=\|a_k\|_{2p}^{2}$, and monotone convergence.
\end{proof}

\begin{corollary}[Neural reference moments]\label{cor:neuralref}
For $b\in\{\bprod,\bMon\}$, a regulator $a$, $N\ge1$, $j\ge1$ and $p\in\N$:
\begin{equation}\label{eq:H1-original-reference}
 \E_b p_j(\iota F^b_{N,a})^{2p}\le L_p K_j^{2p}.
\end{equation}
The separate bound $\E_b p_j(\iota F^b_{N,a})^{2p}\le L_p\sigma_j^p$
also holds. Both are uniform over every finite $N$ and every regulator $a>0$;
$K_j$ and $\sigma_j^{1/2}$ are distinct constants.
\end{corollary}

\begin{proof}
For a fixed spatial point $z$, $X_i=g^b_{i,a}(z)$ are independent,
centered and bounded by $1/2$. In the $2p$-slot expansion of
$\E|N^{-1/2}\sum_iX_i|^{2p}$, every singleton label has zero contribution.
A surviving tuple has $q\le p$ distinct labels; there are at most
$N^qq^{2p}$ tuples with $q$ labels, each term bounded by $2^{-2p}$ in
absolute value. Since $N^{q-p}\le1$,
\begin{equation}\label{eq:H1-pointwise}
 \E_b|F^b_{N,a}(z)|^{2p}\le L_p/4^p.
\end{equation}
Joint measurability and $|F_N|\le\sqrt N/2$ make all spatial integrals
finite for fixed $N$. Integral Minkowski in $L^p$, including $p=1$, gives
\[
 \left\|\int\chi_j(z)^2|F_N(z)|^2dz\right\|_{L^p}
 \le\int\chi_j(z)^2(\E|F_N(z)|^{2p})^{1/p}dz
 \le L_p^{1/p}\|\chi_j\|_2^2/4.
\]
Parseval and $\lambda_{j,k}\ge1$ give
$p_j(\iota F_N)^2\le\|\chi_jF_N\|_2^2$. Raising to power $p$ proves
\eqref{eq:H1-original-reference} with the original $K_j$ exactly.
This uses only positive-regulator pointwise neurons; no point field at
$a=0$ appears.

For the separate $\sigma_j$ bound, let $f\in L^1(\R^2;\C)$,
$Y_i=\int g_if$ and $c=\|f\|_1/2$. The $Y_i$ are independent,
identically distributed, centered (Fubini and $\E g=0$), and $|Y_i|\le c$ because
$|g_i|=1/2$; put $s^{2}=\E|Y_1|^{2}\le c^{2}$. Expanding
$\E|\sum_{i\le N}Y_i|^{2p}=\E(\sum_iY_i)^{p}(\sum_i\conj{Y_i})^{p}$ over index tuples, a term
vanishes unless each index occurring in it occurs at least twice. A term with
exactly $q$ distinct indices, with multiplicities $n_i\ge2$ summing to $2p$, is
bounded by $\prod_i\E|Y_i|^{n_i}\le\prod_ic^{n_i-2}s^{2}=c^{2p-2q}s^{2q}$; there are at most
$\binom Nqq^{2p}\le N^{q}q^{2p}$ such tuples and $q\le p$. Hence
$\E|N^{-1/2}\sum_iY_i|^{2p}\le\sum_{q\le p}q^{2p}N^{q-p}c^{2p-2q}s^{2q}\le L_pc^{2p}$. For
$f=\chi_j\conj{e_{j,k}}$ we have $(\iota F_N)(f)=N^{-1/2}\sum_iY_i$ and
$\|f\|_1\le\|\chi_j\|_1|\T_j|^{-1/2}$, so $\|(\iota F_N)(f)\|_{2p}^{2}\le L_p^{1/p}\|\chi_j\|_1^{2}/(4|\T_j|)$;
Lemma~\ref{lem:hilbertsum} concludes.
\end{proof}

\begin{lemma}[Gaussian localized trace]\label{lem:gausshilbert}
Let $\Xi$ be a Borel random element of $X$ such that for every complex
$f\in C_c^{\infty}$ the real and imaginary parts of $\Xi(f)$ are centered and jointly
Gaussian with $\E|\Xi(f)|^{2}\le m^{-2}\|f\|_2^{2}$. Then $\E p_j(\Xi)^{2p}\le(2p-1)!!\,\tau_j^{p}$ for all
$j,p\ge1$. This applies to $\Phi_a$ under $\PP_0$ for $0\le a<1/m$, to $\iota\gamma'$ under
$G'_{\Omega,a}$ and to $\iota\gamma_\mu$ under $G_\mu$.
\end{lemma}

\begin{proof}
For $T=X'+iY'$ with $X',Y'$ real centered Gaussian,
$\|T\|_{2p}^{2}=\|X'^{2}+Y'^{2}\|_p\le\|X'\|_{2p}^{2}+\|Y'\|_{2p}^{2}=((2p-1)!!)^{1/p}\E|T|^{2}$. With
$\E|\Xi(\chi_j\conj{e_{j,k}})|^{2}\le m^{-2}\|\chi_j\|_2^{2}/|\T_j|$, Lemma~\ref{lem:hilbertsum} gives
the bound. For $\Phi_a$ the hypothesis is Theorem~\ref{thm:G-versions}(a) with
Theorem~\ref{thm:G-bounds}. For a continuous centered Gaussian field $\gamma$ with
covariance kernel $C$ and complex $f\in C_c^{\infty}$, $\int\gamma f$ is an $L^{2}$-limit of
Riemann sums, so its real and imaginary parts are jointly Gaussian, and
$\E|\int\gamma f|^{2}=\langle f,Cf\rangle\le m^{-2}\|f\|_2^{2}$ when $0\le C\le C_a$ as operators
(Theorems~\ref{thm:G-order} and \ref{thm:G-shift}(d)); this covers $C'_{\Omega,a}$ and
$\Cend$.
\end{proof}

\begin{lemma}[Density bounds]\label{lem:weights}
Fix a regulator $a$, a physical square $\Lambda$ and $\Omega=[-R,R]^{2}\supseteq\bar\Lambda$. For
Borel $G:X\to[0,\infty]$:
\begin{equation}\label{eq:wfull-finite}
\E_{P^{b}_{N,\varepsilon,\Omega,a;\Lambda}}G
\le\frac{M_{\mathrm{full}}}{Z_N^b}\,\E_bG(\iota F^b_{N,a})
\qquad(\varepsilon>0,\ N\ge1).
\end{equation}
In particular, before any eventual-width denominator replacement, every
finite $N$, $j,p\in\N$, both architectures and every fixed regulator satisfy
\begin{equation}\label{eq:H1-original-weighted}
 P^b_{N,\varepsilon,\Omega,a;\Lambda}(p_j^{2p})
 \le\frac{M_{\mathrm{full}}L_pK_j^{2p}}{Z_N^b}\qquad(\varepsilon>0).
\end{equation}
Only after these exact finite-width bounds, the eventual-width and other density bounds are
\begin{align}
\E_{P^{b}_{N,\varepsilon,\Omega,a;\Lambda}}G&\le\frac{2M_{\mathrm{full}}}{Z_{\mathrm{full}}}\,\E_b\,G(\iota F^{b}_{N,a})
\quad(\varepsilon>0,\ N\ge N_0),\label{eq:wfull}\\
\E_{\iota_*\Res_{\varepsilon,\Omega,a;\Lambda}}G&\le M_{\mathrm{res}}\,e^{3u|\Lambda|c_a^{2}+3\varepsilon|\Omega\setminus\Lambda|c_a^{2}}\,
\E_{G'_{\Omega,a}}G(\iota\gamma')\quad(\varepsilon\ge0),\label{eq:wres}\\
\E_{\nu_{a',\Lambda}}G&\le\Xi_2(\Lambda)^{1/2}\bigl(\E_0G(\Phi_{a'})^{2}\bigr)^{1/2}\quad(0\le a'<1/m),\label{eq:wphys}
\end{align}
where $M_{\mathrm{full}}$ is given by \eqref{eq:Mfull}, $M_{\mathrm{res}}=e^{6u|\Lambda|c_a^{2}}$, and
$N_0$ depends on $(b,\varepsilon,\Omega,a,\Lambda)$ and is such that $Z^{b}_{N}\ge Z_{\mathrm{full}}/2$ for
$N\ge N_0$. At $\varepsilon=0$ the factor in \eqref{eq:wres} does not depend on $\Omega$.
\end{lemma}

\begin{proof}
\eqref{eq:wfull-finite} is the exact normalized expectation bound from
$0<e^{-V}\le M_{\mathrm{full}}$ and $Z_N^b>0$, valid at every finite width.
Combining it with \eqref{eq:H1-original-reference} gives
\eqref{eq:H1-original-weighted}. Only afterwards, \eqref{eq:wfull} uses
$Z^{b}_{N}\to Z_{\mathrm{full}}>0$
(Proposition~\ref{prop:width}). \eqref{eq:wres}: the density of $\Res_\varepsilon$ with respect to
$G'_{\Omega,a}$ is $e^{-W-\varepsilon J}/z_\varepsilon\le M_{\mathrm{res}}/z_\varepsilon$, and $z_\varepsilon$ is bounded below
in Proposition~\ref{prop:masstilt}. \eqref{eq:wphys}: $Z_{a',\Lambda}\ge1$ and
Cauchy--Schwarz with Theorem~\ref{thm:W}(e) for $p=2$.
\end{proof}

\begin{lemma}[Unbounded moment transfer]\label{lem:transfer}
Let $\rho_k\Rightarrow\rho$ be Borel probabilities on $X$, $j,r\ge1$, and
$\sup_k\rho_k(p_j^{2r})<\infty$. Then $\rho(p_j^{2r})\le\liminf_k\rho_k(p_j^{2r})$, and
$\rho_k(H)\to\rho(H)$ for every continuous $H:X\to\C$ with $|H|\le c(1+p_j^{r})$.
\end{lemma}

\begin{proof}
For $T>0$, $\min(p_j^{2r},T)$ is bounded and continuous; let $T\to\infty$ by monotone
convergence. With $H_T=H\min(1,T/|H|)$, bounded and continuous,
$|\rho_k(H)-\rho_k(H_T)|\le\rho_k(|H|^{2})/T\le2c^{2}(1+\rho_k(p_j^{2r}))/T$, and the same for $\rho$; let
$k\to\infty$, then $T\to\infty$.
\end{proof}

\subsection{Moments along the arrows}

\begin{proposition}[Moments along the arrows]\label{prop:H1arrows}
Fix $j,r\ge1$, a physical square $\Lambda$ and, where they occur, $b$, a regulator $a$,
$\varepsilon>0$ and $\Omega=[-R,R]^{2}\supseteq\bar\Lambda$. Every law below has finite moments of
every power of $p_j$, and:
\begin{enumerate}[label=(\roman*)]
\item \emph{width:} $P^{b}_{N,\varepsilon,\Omega,a;\Lambda}(p_j^{r})\to\nu^{v}_{\varepsilon,\Omega,a;\Lambda}(p_j^{r})$ as $N\to\infty$;
\item \emph{stabilizer:} $\nu^{v}_{\varepsilon,\Omega,a;\Lambda}(p_j^{r})\to\nu^{(2)}_{\Omega,a;\Lambda}(p_j^{r})$ as $\varepsilon\downarrow0$;
\item \emph{auxiliary square:} $\nu^{(2)}_{\Omega,a;\Lambda}(p_j^{r})\to\nu_{a,\Lambda}(p_j^{r})$ as $R\to\infty$, and
  $\sup_R\nu^{(2)}_{\Omega,a;\Lambda}(p_j^{2r})\le M_{\mathrm{res}}e^{3u|\Lambda|c_a^{2}}(2r-1)!!\,\tau_j^{r}$;
\item \emph{regulator:} $\nu_{a,\Lambda}(p_j^{r})\to\nu_{0,\Lambda}(p_j^{r})$ as $a\downarrow0$, and
  $\sup_{0\le a<1/m}\nu_{a,\Lambda}(p_j^{2r})\le\Xi_2(\Lambda)^{1/2}((4r-1)!!)^{1/2}\tau_j^{r}$;
\item \emph{physical volume:} $\nu_n(p_j^{r})\to\nu(p_j^{r})$, with $\sup_n\nu_n(p_j^{2r})<\infty$.
\end{enumerate}
Consequently, for every fixed $n$, $\dist_n$ between consecutive laws of each
arrow tends to $0$ along the arrow.
\end{proposition}

\begin{proof}
(i) By \eqref{eq:wfull} and Corollary~\ref{cor:neuralref},
$\sup_{N\ge N_0}P^{b}_{N}(p_j^{2r})\le2M_{\mathrm{full}}L_rK_j^{2r}/Z_{\mathrm{full}}$
by the exact finite-width inequality \eqref{eq:H1-original-weighted}; for $N<N_0$ each law has
$p_j(\iota F_N)\le\|\chi_j\|_2\sqrt N/2$. Apply Lemma~\ref{lem:transfer} with
Proposition~\ref{prop:width}. (ii) By Proposition~\ref{prop:masstilt},
$\nu^{v}_\varepsilon=\iota_*\Res_\varepsilon$; the densities $e^{-W-\varepsilon J}/z_\varepsilon$ converge pointwise to
$e^{-W}/z_0$ and are bounded by $M_{\mathrm{res}}e^{3u|\Lambda|c_a^{2}+3|\Omega\setminus\Lambda|c_a^{2}}$ for $\varepsilon\le1$
(Proposition~\ref{prop:stab}), and $p_j(\iota\gamma')^{r}\in L^{1}(G'_{\Omega,a})$ by
Lemma~\ref{lem:gausshilbert}; dominated convergence applies. (iii) \eqref{eq:wres}
at $\varepsilon=0$, Lemma~\ref{lem:gausshilbert}, Proposition~\ref{prop:aux} and
Lemma~\ref{lem:transfer}. (iv) \eqref{eq:wphys} with $G=p_j^{2r}$,
Lemma~\ref{lem:gausshilbert} for $\E_0p_j(\Phi_a)^{4r}$, Proposition~\ref{prop:reg} and
Lemma~\ref{lem:transfer}. (v) Since $\lambda_{j,k}\ge1$,
$p_j(v)\le\|\chi_jv\|_{H^{-3/2}(\T_j)}$, whose moments are bounded uniformly in $n$ by
Theorem~\ref{thm:V1}(iii); $\nu_n\Rightarrow\nu$ there, and Lemma~\ref{lem:transfer}
applies. The last sentence combines (i)--(v) with the convergences in $\beta$ of
Propositions~\ref{prop:width}, \ref{prop:stab} (variation distance dominates
$\beta$), \ref{prop:aux}, \ref{prop:reg} and Theorem~\ref{thm:V1}(iii) with
Proposition~\ref{prop:space}(v).
\end{proof}

\subsection{The common sequence}

\begin{proposition}[Selection of the common sequence]\label{prop:H1seq}
Put $\eta_n=2^{-n}/4$ and $\Omega_k=[-k,k]^{2}$. Define recursively, for $n=1,2,\dots$
(conditions involving index $n-1$ only for $n>1$):
\begin{enumerate}[label=(\arabic*)]
\item $a_n=a(k)$ for the least $k\in\N$ with $a(k)<\min(1/n,1/(2m))$, $a(k)<a_{n-1}/2$ and
  $\dist_n(\nu_{a(k),\Lambda_n},\nu_n)<\eta_n$;
\item $R_n=R(k)=k$ for the least $k\in\N$ with $k\ge\max(n,n\ell)$, $k>R_{n-1}$ and
  \[
  \dist_n(\nu^{(2)}_{\Omega_k,a_n;\Lambda_n},\nu_{a_n,\Lambda_n})<\eta_n;
  \]
\item $\varepsilon_n=\varepsilon(k)=1/k$ for the least $k\in\N$ with $1/k<1/n$, $1/k<\varepsilon_{n-1}/2$ and
  \[
  \dist_n(\nu^{v}_{1/k,\Omega_n,a_n;\Lambda_n},\nu^{(2)}_{\Omega_n,a_n;\Lambda_n})<\eta_n,
  \]
  where $\Omega_n=\Omega_{R_n}$;
\item $N_n=N(k)=k$ for the least $k\in\N$ with $k>n$, $k>N_{n-1}$ and
  \[
  \dist_n(P^{b}_{k,\varepsilon_n,\Omega_n,a_n;\Lambda_n},\nu^{v}_{\varepsilon_n,\Omega_n,a_n;\Lambda_n})<\eta_n
  \]
  for both $b$.
\end{enumerate}
Each least index exists, and the resulting sequence satisfies all conditions
and \eqref{eq:mainseq} of Theorem~\ref{thm:main}(S). It is determined by $m,\ell,u$,
the localization system and the enumerations \eqref{eq:enumerations}.
\end{proposition}

\begin{proof}
In each step the order constraints hold for all large $k$, and by
Proposition~\ref{prop:H1arrows} the discrepancy tends to $0$ as $k\to\infty$ (in (1)
because $a(k)\downarrow0$, in (2) along the increasing squares $\Omega_k$, in (3) because
$1/k\downarrow0$, and in (4) for each of the two architectures); so some $k$ satisfies
all conditions and the least one is well defined. In (2), $k\ge n\ell$ gives
$\bar\Lambda_n\subseteq\Omega_k$, so every law in (2)--(4) is defined, and $a_n<1/(2m)$ is a
regulator. Every law entering $\dist_n$ has finite moments of $p_j^{r}$
(Proposition~\ref{prop:H1arrows}). Since $\dist_n$ satisfies the triangle inequality,
$\dist_n(P^{b}_{n},\nu_n)<4\eta_n=2^{-n}$. The laws involved are determined by
$m,u,\ell$, the neuron definitions and the parameters, and $\dist_n$ by the
localization system; no other object enters.
\end{proof}

\begin{remark}[The selection is not an algorithm]\label{rem:nonalgorithmic}
The least-index rule makes the sequence a well-defined deterministic function
of the data listed in Proposition~\ref{prop:H1seq}. It is not an algorithm: the
predicates involve exact values of $\beta$ and of moments of non-explicit laws,
and no procedure deciding them, no bound on the selected indices and no
end-to-end rate is given. Selecting the four parameters in the order $a_n,R_n,\varepsilon_n,N_n$ follows
the order of the iterated limits of Section~\ref{sec:fixed}; no interchange of
limits is used or claimed.
\end{remark}

\subsection{Whole-field convergence, moments and later tests}

\begin{proposition}[Clauses (0) and (A)]\label{prop:H1conv}
For both architectures every $P^{b}_{n}$ is a Borel probability on $X$ with
positive finite normalizer and $p_j(v)\le\|\chi_j\|_2\sqrt{N_n}/2$ on its support;
$\beta(P^{b}_{n},\nu)\to0$; and for all $j,r\ge1$, $P^{b}_{n}(p_j^{r})\to\nu(p_j^{r})$ and
$M_{j,r}=\sup_{\rho\in\cL}\int p_j^{r}\,d\rho<\infty$.
\end{proposition}

\begin{proof}
Lemma~\ref{lem:RPmatch} below and Theorem~\ref{thm:E-RP}(a) give the Borel law and
the normalizer bounds; Parseval gives $p_j(\iota F)\le\|\chi_jF\|_{L^{2}}$. Next,
$\beta(P^{b}_{n},\nu)\le\dist_n(P^{b}_{n},\nu_n)+\beta(\nu_n,\nu)<2^{-n}+\beta(\nu_n,\nu)\to0$ by
Theorem~\ref{thm:V1}(iii) and Proposition~\ref{prop:space}(v). For $n\ge\max(j,r)$,
$|P^{b}_{n}(p_j^{r})-\nu_n(p_j^{r})|<2^{-n}$, and $\nu_n(p_j^{r})\to\nu(p_j^{r})<\infty$ by
Proposition~\ref{prop:H1arrows}(v). A convergent sequence of finite numbers is
bounded, and the finitely many stages with $n<\max(j,r)$ have finite moments.
\end{proof}

\begin{proposition}[Clause (B)]\label{prop:H1later}
Let real $f_1,\dots,f_k\in C_c^{\infty}(\R^{2})$ and a real or complex polynomial $Q$ of
degree $d$ be given after the sequence. Then the expectations in
Theorem~\ref{thm:main}(B) are absolutely finite, uniformly bounded, and converge;
the laws of $(v(f_1),\dots,v(f_k))$ converge weakly.
\end{proposition}

\begin{proof}
The evaluations are continuous on $X$ (Proposition~\ref{prop:space}(iii)),
so the continuous-mapping theorem gives weak convergence of the vectors,
independently of $Q$. If $Q$ is constant, the assertions involving $Q$
follow directly from probability normalization. Otherwise its degree
$d\ge1$. Choose $j$ with all supports
in $(-j,j)^{2}$. By \eqref{eq:duality},
$|v(f_i)|\le p_j(v)\|f_i\|_{H^{2}(\T_j)}$, so
\[
|Q(v(f_1),\dots,v(f_k))|\le c_Q(1+p_j(v)^{d})
\]
and, for
$q>0$, $|Q|^{q}\le c_{Q,q}(1+p_j^{\lceil dq\rceil})$; Proposition~\ref{prop:H1conv} bounds these
uniformly over $\cL$. Lemma~\ref{lem:transfer} with
$H=Q(v(f_1),\dots,v(f_k))$ and $r=d$ gives the convergence of expectations.
\end{proof}

\subsection{Finite-stage reflection positivity}

\begin{lemma}[The stages satisfy the positivity hypotheses]\label{lem:RPmatch}
For every $n$, $\Lambda_n=(-n\ell,n\ell)^{2}$ is a physical square in $\Omega_n=[-R_n,R_n]^{2}$
(Definition~\ref{def:params}); both $\bar\Lambda_n$ and $\Omega_n$ are centered closed
axis-parallel squares invariant under $\theta_1$ and $\theta_2$; and
$P^{b}_{n}=P^{v_V}_{N_n}$ with the integrand $v_V$ of \eqref{eq:vV} for
$(\varepsilon_n,\Omega_n,a_n;\Lambda_n)$, which agrees off the null set $\partial\Lambda_n\times\R$ with the integrand
of Corollary~\ref{cor:E-P2} for the rectangles $\bar\Lambda_n=[-n\ell,n\ell]^{2}$ and $\Omega_n$.
\end{lemma}

\begin{proof}
$R_n\ge n\ell$ by Proposition~\ref{prop:H1seq}; $\varepsilon_n>0$ and $a_n$ is a regulator. The
remaining assertions are the manuscript-tilt part of Corollary~\ref{cor:E-P2}.
\end{proof}

\begin{proposition}[Clause (E1)]\label{prop:H1RP}
For $\theta\in\{\theta_1,\theta_2\}$, every $n$ and both architectures, $P^{b}_{n}$ satisfies
\eqref{eq:boundedRP} for every bounded complex $\Sigma^{\theta}_{+}$-measurable $F$.
\end{proposition}

\begin{proof}
Lemma~\ref{lem:RPmatch} and Corollary~\ref{cor:E-P2}, which gives
Theorem~\ref{thm:E-RP}(b) across both central lines; the integrand is
$\theta$-invariant because both the inner square and the outer square are.
\end{proof}

% ===========================================================================
\section{Joint compact-test correlations}\label{sec:K1}

Correlation functions involve products of field values at several points, which have no pointwise meaning for a distribution. This section uses instead the tensor power $v^{\otimes k}$, a distribution in $k$ plane variables that equals $\prod_rv(f_r)$ on product tests; it is built from localized Fourier coefficients and bounded by $p_j(v)^{k}$ times a norm of the test (Proposition~\ref{prop:tensor}). Averaging over a law gives the correlation distributions, the moment bounds of Section~\ref{sec:H1} give their convergence, and the same objects form the polynomial positivity expressions of clause~(E2) (Proposition~\ref{prop:polyRP}).

For $\kappa=(\kappa_1,\dots,\kappa_k)\in(\Z^{2})^{k}$ put
$e_{j,\kappa}(z_1,\dots,z_k)=\prod_re_{j,\kappa_r}(z_r)$ and $\lambda_{j,\kappa}=\prod_r\lambda_{j,\kappa_r}$.

\begin{lemma}[Fourier facts]\label{lem:fourier}
Let $F\in\cD(((-j,j)^{2})^{k})$, extended by zero to a smooth periodic function on
$\T_j^{k}$, and $F_\kappa=\int_{\cC_j^{k}}\conj{e_{j,\kappa}}F$. Then the Fourier series
$\sum_\kappa F_\kappa e_{j,\kappa}$ converges to $F$ in $C^{\infty}(\T_j^{k})$,
\begin{equation}\label{eq:coeffnorm}
\sum_\kappa\lambda_{j,\kappa}^{2}|F_\kappa|^{2}=A_{j,k}(F)^{2},\qquad
A_{j,k}(F)\le3^{k}|\T_j|^{k/2}\|F\|_{C^{2k}} .
\end{equation}
\end{lemma}

\begin{proof}
$\prod_r(1-\Delta_{z_r})e_{j,\kappa}=\lambda_{j,\kappa}e_{j,\kappa}$, so Parseval on $\T_j^{k}$ gives the
identity; repeated integration by parts makes $F_\kappa$ decay faster than every power
of $|\kappa|$, which gives convergence with all derivatives. Expanding the product
of the operators $1-\partial_1^{2}-\partial_2^{2}$ gives $3^{k}$ terms, each a derivative of order at
most $2k$, and $\|\cdot\|_{L^{2}(\T_j^{k})}\le|\T_j|^{k/2}\|\cdot\|_\infty$.
\end{proof}

\begin{lemma}[Support-preserving product density]\label{lem:density}
Let $U\subseteq\R^{2}$ be open and $F\in\cD(U^{k})$. There are $j$, a real $\zeta\in\cD(U\cap(-j,j)^{2})$
and finite sums $F_l$ of products $f_1\otimes\cdots\otimes f_k$ with $f_r\in\cD(U)$ and
$\supp f_r\subseteq\supp\zeta$, such that $F_l\to F$ in $C^{M}$ for every $M$, all supports lie in
$(\supp\zeta)^{k}$, and $A_{j',k}(F_l-F)\to0$ for every $j'\ge j$.
\end{lemma}

\begin{proof}
Let $K_0\subset U$ be the union of the coordinate projections of $\supp F$, choose $j$
with $K_0\subset(-j,j)^{2}$ and a real $\zeta\in\cD(U\cap(-j,j)^{2})$ with $\zeta=1$ near $K_0$. Then
$F=\zeta^{\otimes k}F$, and with the Fourier partial sums $S_lF$ on $\T_j^{k}$,
$F_l=\zeta^{\otimes k}S_lF=\sum_{|\kappa|_\infty\le l}F_\kappa\bigotimes_r(\zeta e_{j,\kappa_r})$ converges to $F$ in every
$C^{M}$ by Lemma~\ref{lem:fourier} and the Leibniz rule; \eqref{eq:coeffnorm} gives the
last assertion.
\end{proof}

\begin{proposition}[Tensor evaluations; clause (C1)]\label{prop:tensor}
For $v\in X$ and $F\in\cD(((-j,j)^{2})^{k})$ put
$v^{\otimes k}(F)=\sum_\kappa F_\kappa\prod_rv(\chi_je_{j,\kappa_r})$. The series converges absolutely, does
not depend on $j$, and defines the unique $v^{\otimes k}\in\cD'((\R^{2})^{k})$ with
$v^{\otimes k}(f_1\otimes\cdots\otimes f_k)=\prod_rv(f_r)$; it is symmetric, satisfies
$\conj{v^{\otimes k}(F)}=v^{\otimes k}(\conj F)$, $v\mapsto v^{\otimes k}(F)$ is continuous on $X$, and
$|v^{\otimes k}(F)|\le p_j(v)^{k}A_{j,k}(F)\le3^{k}|\T_j|^{k/2}p_j(v)^{k}\|F\|_{C^{2k}}$.
\end{proposition}

\begin{proof}
Cauchy--Schwarz with the weights $\lambda_{j,\kappa}^{\pm1}$, \eqref{eq:coeffnorm},
$v(\chi_je_{j,\kappa_r})=v(\chi_j\conj{e_{j,-\kappa_r}})$ and $\lambda_{j,-\kappa_r}=\lambda_{j,\kappa_r}$ give absolute
convergence and the bound. For $F=f_1\otimes\cdots\otimes f_k$ the coefficients factor and the
series is $\prod_r\sum_{\kappa_r}(f_r)_{\kappa_r}v(\chi_je_{j,\kappa_r})=\prod_rv(f_r)$ by the expansion in the
proof of \eqref{eq:duality}; this holds for every admissible $j$. For general $F$,
Lemma~\ref{lem:density} gives product sums $F_l\to F$ with
$A_{j,k}(F_l-F)\to0$ and $A_{j',k}(F_l-F)\to0$, so the values for $j$ and $j'\ge j$ agree. The
bound shows continuity on each $\cD_K$, so $v^{\otimes k}\in\cD'$; uniqueness follows from the
density of product sums. Symmetry and the reality relation hold on product
tests (for the latter because $v$ is real) and extend by density. For $v,w\in X$,
telescoping gives
$|v^{\otimes k}(F)-w^{\otimes k}(F)|\le kA_{j,k}(F)\,p_j(v-w)\max(p_j(v),p_j(w))^{k-1}$, and
$p_j(v-w)\le2^{j}d(v,w)$ when $d(v,w)<2^{-j}$.
\end{proof}

\begin{remark}[Meaning of the tensor evaluation]\label{rem:tensor-meaning}
$v^{\otimes k}$ is the $k$-fold tensor power of the distribution $v$; for background on
tensor products of distributions see \cite{Zou2022Math218}, cited for orientation
only. Every property used here is proved in Proposition~\ref{prop:tensor}. The
evaluation is defined for every $v\in X$, with no exceptional set.
\end{remark}

\begin{lemma}[Complexification, reflection and products]\label{lem:complexify}
Let $v\in X$, $\theta\in\{\theta_1,\theta_2\}$, $G\in\cD((\R^{2})^{r})$ and $H\in\cD((\R^{2})^{s})$. Then
(a) $(\theta v)^{\otimes r}(G)=v^{\otimes r}(G\circ\theta^{\times r})$, where
$(G\circ\theta^{\times r})(z_1,\dots,z_r)=G(\theta z_1,\dots,\theta z_r)$;
(b) $v^{\otimes r}(G)\,v^{\otimes s}(H)=v^{\otimes(r+s)}(G\otimes H)$;
(c) if $G\in\cD((H^{\theta}_{+})^{r})$, then $v\mapsto v^{\otimes r}(G)$ is $\Sigma^{\theta}_{+}$-measurable;
(d) for $G\in\cD((H^{\theta}_{+})^{r})$, $\conj{(\theta v)^{\otimes r}(G)}=v^{\otimes r}(\theta_r\conj G)$ with
$(\theta_r\conj G)(z)=\conj{G(\theta z_1,\dots,\theta z_r)}$.
\end{lemma}

\begin{proof}
On product tests, (a) is $(\theta v)(f)=v(f\circ\theta)$ (Proposition~\ref{prop:space}(vii),
extended complex-linearly) and (b) is multiplicativity. Both sides of (a) are
continuous in $G$ for $A_{j,r}$ with supports in a fixed compact, since
$p_j(\theta v)\le C_jp_{j+1}(v)$ and $\theta$ preserves $\T_j$ and the Laplacian; for (b),
$A_{j,r+s}(G\otimes H)=A_{j,r}(G)A_{j,s}(H)$. Lemma~\ref{lem:density} concludes. For (c), the
approximants of Lemma~\ref{lem:density} with $U=H^{\theta}_{+}$ give $v^{\otimes r}(G)$ as a pointwise
limit of finite sums of products $\prod_iv(f_i)$ with $f_i\in\cD(H^{\theta}_{+})$, whose real and
imaginary parts are real tests in $C_c^{\infty}(H^{\theta}_{+})$. (d) combines (a) with the
reality relation of Proposition~\ref{prop:tensor}.
\end{proof}

\begin{proposition}[Correlation distributions; clause (C2)]\label{prop:correlations}
For $\rho\in\cL$ and $F\in\cD((\R^{2})^{k})$, $\cS^{\rho}_{k}(F)=\int v^{\otimes k}(F)\,\rho(dv)$ is absolutely
defined, $\cS^{\rho}_{k}\in\cD'((\R^{2})^{k})$ is the unique distribution with
$\cS^{\rho}_{k}(f_1\otimes\cdots\otimes f_k)=\int\prod_rv(f_r)\,\rho(dv)$, and
$\sup_{\rho\in\cL}|\cS^{\rho}_{k}(F)|\le M_{j,k}A_{j,k}(F)$ for $\supp F\subset((-j,j)^{2})^{k}$.
\end{proposition}

\begin{proof}
$v\mapsto v^{\otimes k}(F)$ is continuous, hence Borel, with
$|v^{\otimes k}(F)|\le p_j^{k}A_{j,k}(F)$, and $M_{j,k}<\infty$ (Proposition~\ref{prop:H1conv}).
Linearity is clear; the bound and \eqref{eq:coeffnorm} give continuity on each
$\cD_K$; uniqueness is the density of product sums (Lemma~\ref{lem:density}).
\end{proof}

\begin{proposition}[Kernel convergence; clause (C3)]\label{prop:kernelconv}
For every $k$ and $F\in\cD((\R^{2})^{k})$, $\cS^{b}_{n,k}(F)\to\cS_{\infty,k}(F)$ for both
architectures.
\end{proposition}

\begin{proof}
$H(v)=v^{\otimes k}(F)$ is continuous with $|H|\le A_{j,k}(F)p_j^{k}$; $P^{b}_{n}\Rightarrow\nu$ and
$\sup_nP^{b}_{n}(p_j^{2k})<\infty$ (Proposition~\ref{prop:H1conv}); apply
Lemma~\ref{lem:transfer}.
\end{proof}

\begin{proposition}[Target identification; clause (D)]\label{prop:identification}
For every $k$, $\cS_{\infty,k}=\Schw_k|_{\cD((\R^{2})^{k})}$, with $\Schw_k$ the tempered Schwinger
distribution of $\nu^{S}$ (Theorem~\ref{thm:CT}(v)).
\end{proposition}

\begin{proof}
For $f_r\in\cD(\R^{2})$, $\nu=\mathcal R_*\nu^{S}$ and $\mathcal R(\psi)(f)=\psi(f)$
(Theorem~\ref{thm:V1}(iii), extended complex-linearly) give
$\cS_{\infty,k}(f_1\otimes\cdots\otimes f_k)=\int\prod_r\psi(f_r)\,d\nu^{S}=\Schw_k(f_1\otimes\cdots\otimes f_k)$. The
restriction of $\Schw_k$ to $\cD$ is a distribution because $\cD\hookrightarrow\cS$ is continuous;
two distributions that agree on product tests agree (Lemma~\ref{lem:density}).
\end{proof}

\begin{proposition}[Polynomial reflection positivity; clause (E2)]\label{prop:polyRP}
Let $\theta\in\{\theta_1,\theta_2\}$, $J\ge0$, $f_0\in\C$, $f_r\in\cD((H^{\theta}_{+})^{r})$ and $\rho\in\cL$. Then
$\cO_f,\cO_f\circ\theta\in L^{2}(\rho)$ with
$\|\cO_f\|_{L^{2}(\rho)}\le|f_0|+\sum_{r=1}^{J}A_{j,r}(f_r)M_{j,2r}^{1/2}$ and
$\|\cO_f\circ\theta\|_{L^{2}(\rho)}\le|f_0|+\sum_{r=1}^{J}C_j^{r}A_{j,r}(f_r)M_{j+1,2r}^{1/2}$ for $j$ with all
supports in the corresponding cubes; every term of $\Pi_\rho(f)$ is absolutely finite;
$\Pi_\rho(f)=\int\conj{\cO_f(\theta v)}\,\cO_f(v)\,\rho(dv)$; $\Pi_{P^{b}_{n}}(f)\ge0$; and
$\Pi_{P^{b}_{n}}(f)\to\Pi_\nu(f)\ge0$.
\end{proposition}

\begin{proof}
The $L^{2}$ bounds follow from Proposition~\ref{prop:tensor},
$p_j(\theta v)\le C_jp_{j+1}(v)$ and Minkowski. By Lemma~\ref{lem:complexify}(b),(d),
$\conj{(\theta v)^{\otimes r}(f_r)}\,v^{\otimes s}(f_s)=v^{\otimes(r+s)}((\theta_r\conj{f_r})\otimes f_s)$ for every $v$;
integrating gives the terms of \eqref{eq:polyform}, which are absolutely finite by
Cauchy--Schwarz, and summing gives the integral formula. For $\rho=P^{b}_{n}$, $\cO_f$ is
$\Sigma^{\theta}_{+}$-measurable (Lemma~\ref{lem:complexify}(c)), and the clipped observables
$\clip_T\cO_f=\cO_f\min(1,T/|\cO_f|)$ are bounded and $\Sigma^{\theta}_{+}$-measurable, with
$(\clip_T\cO_f)(\theta v)=\clip_T(\cO_f(\theta v))$. Proposition~\ref{prop:H1RP} gives
$\int\conj{(\clip_T\cO_f)(\theta v)}(\clip_T\cO_f)(v)\,P^{b}_{n}(dv)\ge0$; as $T\to\infty$ the integrands
converge pointwise and are dominated by the integrable $|\cO_f(\theta v)||\cO_f(v)|$, so
$\Pi_{P^{b}_{n}}(f)\ge0$. Each term converges by Proposition~\ref{prop:kernelconv}, so
$\Pi_{P^{b}_{n}}(f)\to\Pi_\nu(f)$, which is therefore real and nonnegative.
\end{proof}

\begin{remark}[What the correlation statements mean]\label{rem:K1-scope}
Convergence in (C3) is convergence of values on each fixed compactly supported
smooth joint test, that is, in $\sigma(\cD'((\R^{2})^{k}),\cD((\R^{2})^{k}))$; no Schwartz or other
noncompact test is used at the neural stages, and no strong-dual convergence is
claimed. The bounds $M_{j,k}$ hold along the selected sequence and at $\nu$. The
positivity of $\Pi_\nu(f)$ is obtained as a limit; it is consistent with, but not
derived from, Theorem~\ref{thm:CT}(ii).
\end{remark}

% ===========================================================================
\section{Proof of the main theorem}\label{sec:proof}

\begin{lemma}[One set of limiting objects]\label{lem:identity}
The law $\nu$ in clauses (0), (A)--(C) and (E) is $\nu=\mathcal R_*\nu^{S}$ of
Theorem~\ref{thm:V1}(iii), with $\nu^{S}$ the target of Theorem~\ref{thm:LR}. Any Borel
probability on $X$ whose finite-dimensional laws on real compact smooth tests are
those of $\nu^{S}$ equals $\nu$. The law $\nu_n$ in \eqref{eq:mainseq} is $\nu_{0,\Lambda_n}$ of
Definition~\ref{def:laws}(d),(e).
\end{lemma}

\begin{proof}
The cylinder events of compact smooth evaluations form a $\cap$-stable generator of
$\cB(X)$ (Proposition~\ref{prop:space}(iii)); apply Dynkin's theorem.
\end{proof}

\begin{proof}[Proof of Theorem~\ref{thm:main}]
The classical inputs enter only through Theorems~\ref{thm:LR}, \ref{thm:V1} and
\ref{thm:CT} and Lemma~\ref{lem:localdensity}; the foundations of
Sections~\ref{sec:space}--\ref{sec:wick} are proved in this paper.
\begin{center}
\small
\begin{tabular}{@{}lp{0.62\textwidth}@{}}
\toprule
Clause & Proved by\\
\midrule
(S) & Proposition~\ref{prop:H1seq}, from Propositions~\ref{prop:width}--\ref{prop:reg}, \ref{prop:H1arrows} and Theorem~\ref{thm:V1}\\
(0) & Proposition~\ref{prop:H1conv}, Lemma~\ref{lem:RPmatch}, Theorem~\ref{thm:E-RP}(a), Theorem~\ref{thm:V1}(iii)\\
(A) & Proposition~\ref{prop:H1conv}\\
(B) & Proposition~\ref{prop:H1later}\\
(C1)--(C3) & Propositions~\ref{prop:tensor}, \ref{prop:correlations}, \ref{prop:kernelconv}\\
(D) & Proposition~\ref{prop:identification}, Theorem~\ref{thm:CT}(v)\\
(E1) & Proposition~\ref{prop:H1RP} (Corollary~\ref{cor:E-P2})\\
(E2) & Proposition~\ref{prop:polyRP}\\
\bottomrule
\end{tabular}
\end{center}
Every clause uses the one sequence of Proposition~\ref{prop:H1seq} with no
subsequence, and the objects quantified in (A)--(E) are chosen after it
(Remark~\ref{rem:quantifiers}). Theorem~\ref{thm:complement} is
Theorem~\ref{thm:CT}.
\end{proof}

\begin{remark}[No circularity]\label{rem:noncircular}
The order of proof is: field spaces (Section~\ref{sec:space}); neurons and their
positivity, Gaussian references and the Wick theory
(Sections~\ref{sec:endpoint}--\ref{sec:wick}, Appendices~\ref{app:E}--\ref{app:W}),
none of which uses an interacting law; the fixed-volume chain
(Section~\ref{sec:fixed}); the continuum lemmas (Section~\ref{sec:target}), where
Theorem~\ref{thm:LR} uses \ref{prem:C1}, \ref{prem:C2} and the Wick identification,
Theorem~\ref{thm:V1} uses \ref{prem:C4}, \ref{prem:C5} and Theorem~\ref{thm:LR}, and
the restricted proof of Theorem~\ref{thm:CT} uses \ref{prem:C1}--\ref{prem:C3},
\ref{prem:C6}, the direct classical Wick lemma and the source portion of
Theorem~\ref{thm:LR}, without its common-space bridge; then
Sections~\ref{sec:H1}--\ref{sec:K1}. Ergodicity of the target is proved in
Theorem~\ref{thm:CT}(iii) before \cite[Thm.~1.0]{FS1977} is applied in (vii). No
statement uses a later one.
\end{remark}

% ===========================================================================
\section{Contribution and relation to prior work}\label{sec:prior}

\subsection{Contribution}\label{sec:prior-contrib}

The quadratic term enters through a Gaussian identity. After the width limit,
the quadratic tilt of the Gaussian endpoint reference over $\Omega$ is again a
Gaussian law, with covariance $C'_{\Omega,a}$; as $\Omega$ grows, $C'_{\Omega,a}$ converges
locally uniformly to $C_a$, whose Fourier multiplier is
$(m^{2}+|p|^{2}+a^{2}p_1^{2}p_2^{2})^{-1}$ by the mass-shift identity
$1/\hCend-\kappa_a^{2}=1/\hat C_a$. The regulator $a$ therefore controls the anisotropy of
this Gaussian reference covariance, and the anisotropy disappears as $a\to0$. The
quartic interaction is separate from this identity: the finite-width laws and
the interacting laws are non-Gaussian, and their covariances are not given by it.

For both architectures we prove the following. First, at fixed physical volume
and for every $m,u>0$, the laws converge successively in the width, the
stabilizer, the auxiliary volume and the regulator to the free-boundary
Wick-quartic law on the physical square, together with all positive integer
powers of the local seminorms. Second, for $0<u/m^{2}<g_{*}$, one deterministic
sequence of all parameters, fixed before any observable is chosen and never
passed to a subsequence, converges weakly on $H^{-2}_{\mathrm{loc}}(\R^{2})$ to the
$H^{-2}_{\mathrm{loc}}$-valued version of the classical infinite-volume
weak-coupling Wick-quartic measure; the same sequence gives convergence, with
uniform moment bounds, of every polynomial in finitely many compactly supported
smooth smearings, and convergence of every compact joint-test correlation
distribution. Third, every stage of this sequence is reflection positive across
the two central coordinate lines $z_1=0$ and $z_2=0$, both for bounded
half-plane observables and for compact joint-test polynomial forms. We prove no
reflection-positivity statement for other lines, and we claim neither Euclidean
invariance nor clustering of the finite stages.

These statements rest on attributed classical inputs and on the foundations
proved here: the endpoint neurons and their finite-width central reflection
positivity, the Gaussian reference constructions, and the fixed-volume Wick
estimates. The infinite-volume target, its Osterwalder--Schrader properties and
its reconstruction are classical results used with attribution.

\subsection{Relation to prior work}\label{sec:prior-work}

Frank \cite{Frank2026v2} gives a constructive neural-network approach to the
finite-volume Wick-ordered $\phi^{4}_{2}$ measure at every coupling. His networks are
quenched random-feature shells with Gaussian coefficients and a smooth spatial
cutoff; for almost every draw of the features, along a prescribed nested width
schedule, the joint laws of finitely many Schwartz smearings converge in either
order of the width and ultraviolet limits, with explicit width thresholds and
width-uniform ultraviolet tail bounds. His coefficient $\lambda/4!$ corresponds to
our $u$. Our construction differs in its architectures, in using annealed rather
than quenched laws, in its sharp-square interaction with a stabilizer and a mass
tilt, in its finite-stage reflection positivity, and in including a
physical-volume limit, which we obtain by combining the fixed-volume neural
limits with classical infinite-volume results. Frank's construction stops at
finite volume and leaves open whether neural constructions reach infinite volume
themselves; we do not claim that an analogous classical thermodynamic step could
not be attached to it. At fixed volume, his width-uniform estimates imply, by a
standard uniform-integrability argument (Remark~\ref{rem:frank-moments}),
convergence of all polynomial moments of finitely many smeared fields for his
laws, so we make no novelty claim based on moment convergence. Whether his
estimates also yield fixed-volume convergence of joint compact-test kernels or of
the field on $H^{-2}_{\mathrm{loc}}$ requires estimates that we have not
established, chiefly a bound on localized smeared fields summed over all
frequency shells uniformly in the localized Fourier mode, followed by tightness
and an $H^{-2}_{\mathrm{loc}}$-valued realization of his infinite-cutoff field. We
base no novelty claim on the kernel or topology labels either.

The annealed reweighting is the parameter-density deformation
$P=P_G\,e^{-\lambda\int\mathcal{O}}$ of Demirtas et al.
\cite{DemirtasHalversonMaitiSchwartzStoner2023}; Zhang
\cite{Zhang2026OptimalArchitecture} writes the same normalized reweighting of a
free ensemble. In the parts of the former work that we consulted, the
deformation is formal and at infinite width. We realize it at finite width for
the two architectures above, with finite positive normalization and with the
limits stated above.

The reflection-positivity argument uses conditional factorization. For product neurons, conditioning on the complete boundary-data tuple of all $N$ neurons
makes the two half-plane sides independent with equal conditional laws. For Mondrian neurons, the square identity is first proved in symmetric finite boxes
using the extended trace, including positions, auxiliary times and colors (Corollary~\ref{cor:E-tracebox}). Whole-plane half-plane positivity then follows
by $L^2$ approximation of bounded observables (Proposition~\ref{prop:E-sides}). In both cases, a reflection-invariant additive potential is absorbed into the
side observables, preserving positivity (Lemma~\ref{lem:E-factor} and Theorem~\ref{thm:E-RP}). This instantiates the conditional-factorization principle
underlying the one-dimensional parameter-splitting mechanism of Ferko and Halverson \cite{FerkoHalverson2025QMNN}.
They also show that symmetric Markov processes are reflection
positive, a connection they attribute to Nelson and others, and that reflection
positivity survives the linear combinations and pointwise neural maps they
consider, for example of independent inputs and with random maps drawn
independently of the input. Reflection positivity of Markov fields and lattice
models is classical \cite{FILS1978}. We instantiate the mechanism in two
dimensions for these architectures, across the two central coordinate lines.

Harvey \cite{Harvey2026v2} shows that, for $d\ge2$, a pointwise ensemble whose
two-point function is Euclidean invariant, $S_2(gx,gy)=S_2(x,y)$ for all
translations, rotations and reflections $g$, and continuous, and which is
two-point reflection positive across one hyperplane, has a constant two-point
function; he presents this as an answer, for that finite-variance class, to the
question of reflection positivity in $d\ge2$ raised in
\cite{FerkoHalversonMutchler2026v1}. The invariance hypothesis concerns the
two-point function, not the full law. Our finite stages have bounded values. We
do not claim that the two-point functions of our normalized interacting
finite-stage laws are Euclidean invariant, and we do not determine whether they
are; we therefore assert neither that the theorem applies to these stages nor
that it does not, and the statements we prove do not conflict with it. He notes
that products of coordinate cosine networks have the product kernel
$\prod_ie^{-m|x_i-y_i|}$ with two-point reflection positivity across coordinate
hyperplanes, giving up exact rotation invariance. Our construction does not use
rotation invariance either: the endpoint covariance $e^{-\mu|z|_1}/4$ of the
neurons is not rotation invariant. This is a property of that reference
covariance; it is not a statement about the two-point function of the
interacting finite-stage laws. What we add at finite width is reflection
positivity of the full normalized law for bounded observables across the two
central coordinate lines, rather than two-point positivity alone.

Ferko and Mutchler \cite{FerkoMutchler2026v1} prove obstructions for specified
finite-width classes: finite expansions in fixed deterministic features,
including distributional neurons, cannot reproduce a nontrivial
K\"all\'en--Lehmann covariance; ensembles with bounded mollified coincident
variance, such as networks with locally bounded pointwise second moments, are
incompatible with a nontrivial K\"all\'en--Lehmann representation, which
temperedness, Euclidean invariance and reflection positivity would furnish; and
finite single-layer ridge networks cannot reproduce all correlation functions of
an Osterwalder--Schrader field with a nontrivial K\"all\'en--Lehmann measure, a
result they state is specific to single-layer networks. They also stress that a
parameter count alone cannot define finiteness, since a single Borel parameter
can encode any field law. Our finite stages have bounded values and are not
claimed to be Euclidean invariant, so they do not conflict with the variance
obstruction, and the ridge theorem does not describe product-telegraph or
Mondrian neurons. Ferko, Halverson and Mutchler
\cite{FerkoHalversonMutchler2026v1} show that every law on $\cS'(\R^{d})$ admits a
neural-network representation with countably many parameters, or formally with a
single one; we therefore do not claim representability of the target as a
contribution.

For fields built from neurons see Halverson, Maiti and Stoner
\cite{HalversonMaitiStoner2021} and Halverson \cite{Halverson2021Neurons}. Zhang
\cite{Zhang2026OptimalArchitecture} and Ageev and Ageeva
\cite{AgeevAgeeva2026Transformer} study adjacent architecture and cost questions
that are not addressed here.

\begin{remark}[Fixed-volume moments for the laws of \cite{Frank2026v2}]\label{rem:frank-moments}
The moment sentence above rests on the following short derivation, which uses
only estimates stated in \cite{Frank2026v2}; the conclusion is not stated there.
Let $X$ be the pairing of the cutoff field of \cite{Frank2026v2} with a Schwartz
function $f$. The normalizer there is at least one \cite[Lemma~2.8]{Frank2026v2},
so Cauchy--Schwarz gives
$\E_\mu|X|^{q}\le(\E_\gamma e^{-2V})^{1/2}(\E_\gamma|X|^{2q})^{1/2}$, where $\gamma$ is the Gaussian
prior. The first factor is the normalizer at doubled coupling, which
\cite{Frank2026v2} bounds uniformly in the shell cutoff and the width index on
the almost-sure feature event. Under the prior, $X$ is centered Gaussian with
variance bounded, for every feature realization, by the deterministic bound on
the shell pairing coefficients \cite[(6.38)]{Frank2026v2}. Gaussian moments
therefore bound $\sup\E|X|^{q}$ over the width index and the cutoff for every $q$,
likewise for finite tuples, and Fatou passes the bounds to the limit laws. The
weak convergence statements of \cite[Lemma~7.5, Thm.~6.6(c), Thm.~7.6,
Cor.~7.8]{Frank2026v2}, together with uniform integrability, then give convergence
of every polynomial moment along each iterated limit. We do not claim the
corresponding statements for joint compact-test kernels or for the topology of
$H^{-2}_{\mathrm{loc}}$; they would require the estimates named above.
\end{remark}

% ===========================================================================
\section{Limitations}\label{sec:limits}

The construction is existential. It provides no effective end-to-end parameter
schedule, computational cost, sampling guarantee or numerical weak-coupling
threshold: the parameter sequence is shown to exist but is not computed. The
fixed-volume Wick step includes explicit $\sqrt a$ bounds whose constants depend
on the volume; these are fixed-volume upper bounds, not practical convergence
rates.

This limitation describes the original common sequence. The private
post-v1 Appendix~\ref{qrs:supplement} gives finite-observation bounds,
conditional elementary selection and an explicit pressure-target example,
with each result's separate hypotheses and source qualifications.

The infinite-volume limit, the target, its Euclidean invariance, clustering and
reflection positivity across every line, and its $1+1$-dimensional Wightman
reconstruction are classical results used with attribution. No new interacting
theory and no statement in other dimensions is claimed. The classical inputs
\ref{prem:C1}--\ref{prem:C5} are used in the formulations of
Section~\ref{sec:classical}; their proofs were not audited, the bodies of
\cite{OS1973,OS1975} and \cite{StreaterWightman1964} and the antecedents of the
review in \cite[pp.~316--319]{GRS1976} were not available, and the page-level
access record is in Appendix~\ref{app:provenance}.

Finite-stage reflection positivity is proved across the two central coordinate
lines, for the centered axis-parallel squares used here and, more generally, for
every bounded auxiliary region and every measurable potential, bounded on
bounded field ranges, that are invariant under both central reflections. Nothing
is claimed for other lines. For a rotated physical square nothing is claimed
beyond this invariant case; in particular, no failure of reflection positivity
is asserted for any rotated square. The finite stages are not claimed to be
Euclidean invariant, and no finite-stage clustering statement is made.

The laws are annealed; nothing is claimed about quenched or fixed-feature laws.

Neural convergence is proved for compactly supported test functions only. No
global Schwartz or strong-dual convergence is claimed.

The reflection-positivity mechanism, the density-deformation mechanism and the
endpoint covariance kernel all have prior instances in the cited works.

The Wick interaction at $a=0$ is constructed on one common Gaussian space as an
$L^{2}$ limit in the fourth chaos, for a fixed physical square; its exceptional
null sets are fixed per regulator value and per test, or common to countable
families. Nothing is asserted about other representations of that variable,
about uncountable families of regulators, or about volume-uniform versions of
the Wick estimates beyond $Z_{a,\Lambda}\ge1$.

The review status of this exact version, including that of
Theorem~\ref{thm:main}, is recorded in its companion decision record, not in this
text. Appendix~\ref{app:provenance} identifies the components that are retained,
modified, re-expressed or newly written here, and the current-text review that
the project's review policy requires for each. All review recorded for the
antecedents was internal model review; none of it is independent expert review.

% ===========================================================================
\appendix

\section{Proofs for the process neurons}\label{app:E}

This appendix proves Theorems~\ref{thm:E-product}--\ref{thm:E-RP} and
Corollaries~\ref{cor:E-P2}, \ref{cor:E-general}. Besides \ref{prem:C6} it uses:
(K1) Dynkin's $\pi$--$\lambda$ theorem and the functional monotone-class theorem;
(K2) Doob--Dynkin factorization, and the fact that a function measurable for the
$\sigma$-algebra generated by a family of maps depends on countably many of them;
(K3) if $Y$ is Polish and continuous real $l_1,l_2,\dots$ separate its points, then
$y\mapsto(l_k(y))_k$ is a continuous injection into $\R^{\N}$, hence (Lusin--Souslin) a
Borel isomorphism onto its Borel image, so $\sigma(l_k:k\ge1)=\cB(Y)$; (K4) finite
counting measures on a Polish space form a standard Borel space whose
$\sigma$-algebra is generated by $\nu\mapsto\nu(A)$, and $\nu\mapsto\int f\,d\nu$ is Borel for
bounded Borel $f$; (K5) Kolmogorov extension for sequences of standard Borel
spaces; (K6) the Poisson facts of \ref{prem:C6}(h). For $\theta\in\{\theta_1,\theta_2\}$ write
$H_+=H^{\theta}_{+}$ and $Q_n=Q^{\theta}_n=H_+\cap(-n,n)^{2}$, and $r_N=\sqrt N/2$.

\subsection{Measurability}

\begin{lemma}[Partition spaces]\label{lem:E-partspaces}
Let $C$ be a box with nonempty interior. (i) $\Part^{\mathrm{col}}(C)$ is a Borel subset of
the standard Borel space of finite counting measures on $\R^{4}\times\{-1,1\}$, hence
standard Borel; restriction $\resB_B$, the union map
$\cup(P_0,P_1)=P_0\cup P_1$ (for daughters of a cut), and images under translations
and $\theta_1,\theta_2$ are Borel. (ii) For open $U\subseteq\mathrm{int}\,C$, the field map
$P\mapsto\fld(P)|_U$, $\fld(P)=\sum_{(R,\xi)\in P}(\xi/2)\ind_{R^{\mathrm{ll}}}$, is Borel into $L^{1}(U)$.
(iii) $(z,P)\mapsto\fld(P)(z)$ is $\cB(\R^{2})\otimes\cB(\Part^{\mathrm{col}}(C))$-measurable.
\end{lemma}

\begin{proof}
(i) On the Borel set of counting measures with exactly $k$ atoms of mass one, the
atoms can be listed lexicographically by a Borel map into $(\R^{4}\times\{\pm1\})^{k}$. The
conditions ``each atom codes a box with nonempty interior contained in $C$,
interiors pairwise disjoint, areas summing to $|C|$'' are finitely many Borel
conditions on the list, and they characterize the $k$-cell partitions: the union
of the boxes is closed, and a missing point of $C$ would have a missing relative
neighbourhood of positive area. Restriction, union and the maps act
coordinatewise on the list by Borel operations (maxima and minima of corners,
discarding pieces of zero width or height). (ii) For bounded Borel $f$ on $U$,
$P\mapsto\int_Uf\fld(P)=\int h_f\,dP$ with
$h_f(a,b,c,d,\xi)=(\xi/2)\int_{U\cap[a,b]\times[c,d]}f$, Borel by (K4); the indicators of
rational boxes with closure in $U$ give continuous functionals on $L^{1}(U)$ that
separate points (Lebesgue differentiation), so (K3) applies. (iii) The integrand
$(z,a,b,c,d,\xi)\mapsto(\xi/2)\ind\{a\le z_1<b,\ c\le z_2<d\}$ is Borel, and integrals of jointly
Borel bounded functions against counting measures are jointly measurable (K1).
\end{proof}

\begin{lemma}[Random fields and parameter maps]\label{lem:E-fields}
Let $r\ge0$, $U\subseteq\R^{2}$ open and $\clip_r(x)=\max(-r,\min(r,x))$.
(i) If $F:\R^{2}\times\Gamma\to\R$ is $\cB(\R^{2})\otimes\cA$-measurable with $|F|\le r$, then
$\omega\mapsto[F(\cdot,\omega)|_U]$ is measurable into $\Ball_r(U)$, and for $U=\R^{2}$ into $E$.
(ii) Two such fields with the same finite-dimensional distributions have the same
law of $F|_U$.
(iii) Let $(Y_{\mathrm{sp}},\mathcal Y)$ be a measurable space and $G:\R^{2}\times Y_{\mathrm{sp}}\to\R$ be
$\cB(\R^{2})\otimes\mathcal Y$-measurable. Then $\Gamma^{r}_{G}(y)=[\clip_r(G(\cdot,y))|_U]$ is defined for every $y$,
takes values in $\Ball_r(U)$ and is measurable. If $Y:\Gamma\to Y_{\mathrm{sp}}$ is measurable and
$F(z,\omega)=G(z,Y(\omega))$ for almost every $z\in U$ for every outcome of an almost sure
event, then $[F|_U]=\Gamma^{r}_{G}(Y)$ almost surely, and the law of $[F|_U]$ is the image of
the law of $Y$. If $|G|\le r$, the clip changes nothing.
\end{lemma}

\begin{proof}
(i) For $s_0\in L^{1}_{\mathrm{loc}}(U)$ and compact $K\subset U$,
$\omega\mapsto\int_K|F(z,\omega)-s_0(z)|dz$ is measurable by Tonelli; $L^{1}_{\mathrm{loc}}(U)$ is separable
metric, so open sets are countable unions of balls; the same argument works in
$E$. (ii) Indicators of rational boxes with closure in $U$ generate the Borel
$\sigma$-algebra (K3), finite cylinders form a $\pi$-system (K1), and the law of a bounded
vector $(\int f_{k_1}F,\dots,\int f_{k_m}F)$ is determined by its mixed moments
$\int\cdots\int\prod_jf_{k_j}(z_j)\E\prod_jF(z_j)\,dz$. (iii) Each section of $G$ is Borel and
$\clip_r$ is continuous, so the clipped section is Borel and bounded by $r$. For a
Borel representative $s_0$ and compact $K$, $(z,y)\mapsto\ind_K(z)|\clip_r(G(z,y))-s_0(z)|$ is
product measurable, so $y\mapsto\int_K|\clip_rG(z,y)-s_0(z)|dz$ is measurable (the measurability
part of Tonelli needs $\sigma$-finiteness only in $z$: the sets $D$ whose sections
satisfy that $y\mapsto|D_y\cap K_m|$ is measurable form a Dynkin system containing
rectangles); the metric of $L^{1}_{\mathrm{loc}}(U)$ and separability give measurability of
$\Gamma^{r}_{G}$. If $F=G(\cdot,Y)$ almost everywhere on $U$ and $|F|\le r$, then
$\clip_r(G(\cdot,Y))=F$ almost everywhere on $U$. Without the clip the map need not
be defined: $G(z,1)=|z|^{-2}$ is not locally integrable at $0$.
\end{proof}

\begin{lemma}[Integrated potentials of bounded fields]\label{lem:E-potential}
Let $U\subseteq\R^{2}$ be open, $A\subseteq U$ bounded and Lebesgue measurable, $r\ge0$, and
$v$ an $\cL(\R^{2})\otimes\cB(\R)$-measurable function bounded on $A\times[-r,r]$. Then
$\Psi_v(s)=\int_Av(z,s(z))\,dz$ is well defined and Borel on $\Ball_r(U)$.
\end{lemma}

\begin{proof}
For a representative with values in $[-r,r]$, $z\mapsto(z,s(z))$ is
$(\cL,\cL\otimes\cB)$-measurable, so $z\mapsto h(z,s(z))$ is Lebesgue measurable, and
representatives differ on null sets. For $h=\ind_{A'}(z)\varphi(x)$ with $\varphi$
continuous, the map is continuous on $\Ball_r(U)$: if $s_k\to s$ in $L^{1}_{\mathrm{loc}}(U)$
within $\Ball_r$ and $K\subset U$ is compact with $|A\setminus K|<\delta$, then
$\int_A|s_k-s|\le\int_K|s_k-s|+2r\delta$, so $s_k\to s$ in measure on $A$ and the integrals
converge by uniform continuity of $\varphi$ on $[-r,r]$ (the bound $r$ handles sets
$A$ touching $\partial U$). The class of bounded $h$ for which the map is Borel is a
vector space closed under bounded monotone limits and contains a
multiplicative generating class; (K1) applies to $v\ind_{A\times[-r,r]}$.
\end{proof}

\begin{lemma}[Restrictions and point recovery]\label{lem:E-restrictions}
Fix $\theta$ and $r\ge0$. (i) $r_n:\Ball_r(H_+)\to L^{1}(Q_n)$, $s\mapsto s|_{Q_n}$, is well defined
and continuous, and $r_m=\rho_{n,m}\circ r_n$ for $m\le n$ with the continuous restriction
$\rho_{n,m}$. (ii) The relative Borel $\sigma$-algebra of $\Ball_r(H_+)$ equals $\sigma(r_n:n\ge1)$,
and $\sigma(r_1,\dots,r_n)=\sigma(r_n)$ increases with $n$. (iii) For $z\in H_+$ with discs
$D_j(z)$ of radius $1/j$, $L_z(s)=\limsup_j|D_j(z)|^{-1}\int_{D_j(z)}s$ is Borel on $\Ball_r(H_+)$ with values
in $[-r,r]$, and $L_z(s)=c$ if $s$ has a representative equal to $c$ near $z$.
\end{lemma}

\begin{proof}
(i) $|s|\le r$ and $|Q_n|<\infty$; if $s_j\to s$ in $L^{1}_{\mathrm{loc}}(H_+)$ within $\Ball_r$, choose compact
$K\subset Q_n$ with $|Q_n\setminus K|<\delta$. (Without the bound the map is undefined:
$1/z_1$ is locally integrable on $\{z_1>0\}$ but not on $(0,n)\times(-n,n)$.) (ii) Each
compact $K\subset H_+$ lies in some $Q_n$, and $s\mapsto\|s-s_0\|_{L^{1}(K)}$ is a continuous
function of $r_n(s)$; relatively open sets are countable unions of such balls.
(iii) For large $j$ the closed disc lies in $H_+$, the averages are continuous, and a
$\limsup$ of countably many Borel maps is Borel.
\end{proof}

\subsection{The canonical box Mondrian}

For each finite binary tree $\tau$ (a finite prefix-closed set of words in which every
non-leaf has both children), the event that the born nodes form $\tau$ is defined by
finitely many Borel inequalities in $(C,s)$ and the ingredients of the nodes of $\tau$,
and on it the output is a Borel function of finitely many variables; on the
complement of the union of these countably many events (a null set, by the next
lemma) set the output to $\{(C,\xi_\emptyset)\}$. For a box $C=[a,b]\times[c,d]$ with $p(C)>0$, let
$\kappa_C$ be the law of the line which is vertical with probability $(b-a)/p(C)$ at a
uniform position in $(a,b)$ and horizontal with probability $(d-c)/p(C)$ at a uniform
position in $(c,d)$; for a line $\lambda$ meeting $\mathrm{int}\,C$, $C_0(\lambda)$ and $C_1(\lambda)$ are the
closed left/lower and right/upper pieces.

\begin{lemma}[Nonexplosion]\label{lem:E-nonexplosion}
Almost surely finitely many nodes are born; the mean number of born nodes is at
most $2e^{sp(C)}-1$; a fixed $z\in\mathrm{int}\,C$ is almost surely interior to a cell; and the
event that $z$ lies on a cell boundary is Borel in $\Part^{\mathrm{col}}(C)$.
\end{lemma}

\begin{proof}
A daughter never has larger half-perimeter than its parent, so a born node $w$ at
depth $k$ with strict ancestors $w_0,\dots,w_{k-1}$ has
$s>b_w=\sum_jE_{w_j}/p(R_{w_j})\ge\sum_jE_{w_j}/p(C)$. For a fixed word of length $k$ the
ancestors' clocks are independent $\mathrm{Exp}(1)$, so with $c=sp(C)$ the mean number of born
nodes is at most $\sum_k2^{k}\PP(\Gamma_k\le c)=\sum_k2^{k}\PP(\mathrm{Poi}(c)\ge k)=\sum_j\PP(\mathrm{Poi}(c)=j)(2^{j+1}-1)=2e^{c}-1$,
$\Gamma_k$ a sum of $k$ independent $\mathrm{Exp}(1)$. For a fixed node the cut position is
$a+X_w(b-a)$ or $c+X_w(d-c)$ with $X_w$ uniform and independent of the node's box, so
it passes through $z$ with probability $0$; cell boundaries inside $\mathrm{int}\,C$ lie on
ancestors' cut lines; there are countably many nodes. The boundary event is
Borel by (K4).
\end{proof}

\begin{definition}[First-step recursion]\label{def:E-rec}
Let $I$ be an index set and, for $i\in I$, $C_i$ a box with nonempty interior,
$s_i\ge0$ and $Q_i$ a probability on $\Part^{\mathrm{col}}(C_i)$. A \emph{daughter assignment}
gives, for every $i$, $t\in[0,s_i)$, line $\lambda$ meeting $\mathrm{int}\,C_i$ and $\epsilon\in\{0,1\}$, a
probability $Q_{i,t,\lambda,\epsilon}$ on $\Part^{\mathrm{col}}(C_\epsilon(\lambda))$ which is either $Q_j$ for some
$j\in I$ with $(C_j,s_j)=(C_\epsilon(\lambda),s_i-t)$ or $\Mcol(C_\epsilon(\lambda),s_i-t)$, such that
$(t,\lambda)\mapsto[Q_{i,t,\lambda,0}\otimes Q_{i,t,\lambda,1}](\cup^{-1}A)$ is Borel for Borel $A$. The family
satisfies the \emph{first-step recursion} if for every $i$ and Borel $A$,
\[
Q_i(A)=e^{-s_ip(C_i)}\PP_\xi(\{(C_i,\xi)\}\in A)
+\int_0^{s_i}p(C_i)e^{-p(C_i)t}\!\int\kappa_{C_i}(d\lambda)\,[Q_{i,t,\lambda,0}\otimes Q_{i,t,\lambda,1}](\cup^{-1}A)\,dt,
\]
with $\xi$ a uniform sign.
\end{definition}

\begin{lemma}[First-step identity and characterization]\label{lem:E-recursion}
(i) The laws $\Mcol(C,s)$ satisfy the first-step recursion with daughter laws
$\Mcol(C_\epsilon(\lambda),s-t)$. (ii) Every family satisfying the first-step recursion
with a daughter assignment has $Q_i=\Mcol(C_i,s_i)$ for all $i$. (iii) For $T$ a
translation, $\theta_1$ or $\theta_2$, $T_*\Mcol(T^{-1}C,s)=\Mcol(C,s)$.
\end{lemma}

\begin{proof}
(i) The root ingredients are independent of the two subfamilies indexed by
$0w$ and $1w$, which are canonical families; $d_\emptyset=E_\emptyset/p(C)\sim\mathrm{Exp}(p(C))$ and
$(D_\emptyset,X_\emptyset)$ produce $\lambda\sim\kappa_C$ independent of $d_\emptyset$. On $\{d_\emptyset\ge s\}$ the output
is $\{(C,\xi_\emptyset)\}$; on $\{d_\emptyset<s\}$ the subtree rooted at $\epsilon$, with births shifted by
$d_\emptyset$, is the construction on $C_\epsilon(\lambda)$ with budget $s-d_\emptyset$; Fubini gives the
formula, and joint Borel measurability of the output gives the measurability.
(ii) Write $Q^{(k)}(A)=Q(A\cap\Part^{\mathrm{col}}_k)$ for the $k$-cell part. Cells of the two
daughters have disjoint interiors, so the $k$-cell part of the integrand involves
only $Q^{(k_0)}\otimes Q^{(k_1)}$ with $k_0+k_1=k$, $k_\epsilon\ge1$. Induct on $k$ simultaneously for
all $i$: for $k=1$ only the first term contributes; for $k\ge2$ every daughter law
agrees with $\Mcol$ on fewer than $k$ cells by hypothesis, so $Q^{(k)}_i$ equals the
$k$-cell part of the right side of (i). Both laws are carried by finite partitions
(Lemma~\ref{lem:E-nonexplosion}). (iii) The family $T_*\Mcol(T^{-1}C,s)$, indexed by all
$(C,s)$, satisfies the recursion: $T$ commutes with $\cup$, preserves half-perimeters,
maps $\kappa_{T^{-1}C}$ onto $\kappa_C$ and daughters onto daughters, possibly exchanged, which
does not change $[Q_0\otimes Q_1](\cup^{-1}A)$. Apply (ii).
\end{proof}

\begin{lemma}[Competing clocks, conditional form]\label{lem:E-clocks}
Let $A_1,\dots,A_k$ be independent, $A_j\sim\mathrm{Exp}(\alpha_j)$, and $Z$ a random element
independent of them; $\alpha=\sum\alpha_j$, $\tau=\min A_j$, $J=\mathrm{argmin}\,A_j$. Then
$\PP(J=j)=\alpha_j/\alpha$, and conditionally on $\{J=j\}$ the random elements $\tau$,
$(A_l-\tau)_{l\ne j}$ and $Z$ are independent, with $\tau\sim\mathrm{Exp}(\alpha)$, $A_l-\tau\sim\mathrm{Exp}(\alpha_l)$ and $Z$
distributed as unconditionally.
\end{lemma}

\begin{proof}
$\PP(J=j,\tau\in dt,A_l-\tau>r_l\ (l\ne j),Z\in G)=\alpha_je^{-\alpha_jt}\prod_{l\ne j}e^{-\alpha_l(t+r_l)}\PP(Z\in G)\,dt$.
The residual vector is indexed by $l\ne J$ and is in general not independent of $J$:
with $\alpha_1=1$, $\alpha_2=2$, the losing residual $R$ has $\PP(R>1\mid J=1)=e^{-2}$ and
$\PP(R>1\mid J=2)=e^{-1}$.
\end{proof}

\subsection{The core-exterior construction}

\begin{definition}[Core-exterior construction]\label{def:E-coreext}
Let $K\subseteq C$ be boxes ($K$ possibly degenerate), $C$ with nonempty interior,
$s\ge0$. Ingredients, all independent: a canonical family
$h=(E_w,D_w,X_w,\xi_w)_{w\in\mathrm{Tr}}$ (the \emph{core}) and, for every $w\in\mathrm{Tr}$, side
$\sigma\in\{L,R,D,U\}$ and $k\ge1$, a triple $(E^{\sigma}_{w,k},X^{\sigma}_{w,k},G^{\sigma}_{w,k})$ with
$E\sim\mathrm{Exp}(1)$, $X\sim U(0,1)$ and $G$ a canonical family. Run the canonical
construction on $K$ with $h$ (core boxes $K_w$, births $b_w$, deaths $d_w$, split lines);
$e_w=\min(d_w,s)$. Each born core node carries an \emph{enclosing box}
$C_w(t)=[a_w(t),b_w(t)]\times[c_w(t),d_w(t)]\supseteq K_w=[a'_w,b'_w]\times[c'_w,d'_w]$, $t\in[b_w,e_w]$, with
\emph{gaps} $g^{L}=a'_w-a_w$, $g^{R}=b_w-b'_w$, $g^{D}=c'_w-c_w$, $g^{U}=d_w-d'_w$; initially
$C_\emptyset(0)=C$. On $[b_w,e_w)$ side $\sigma$ jumps at the times
$t^{\sigma}_{w,k}=t^{\sigma}_{w,k-1}+E^{\sigma}_{w,k}/g^{\sigma}(t^{\sigma}_{w,k-1})$, $t^{\sigma}_{w,0}=b_w$ (no jumps while the gap is
$0$). At a jump time $t<e_w$ the gap $g$ becomes $X^{\sigma}_{w,k}g$, and the removed slab (the
\emph{detached box}, of full current extent in the other coordinate) receives
$\Pi^{\mathrm{can}}(\text{slab},s-t)$ built from $G^{\sigma}_{w,k}$: on the slab itself for $\sigma=R,U$, and on its
image under $\theta_1$ (for $L$) or $\theta_2$ (for $D$), mapped back by the same reflection.
If $d_w<s$, the core split line of $K_w$ cuts $C_w(d_w)$ and each child receives the
piece containing its core box (gap $0$ on the side of the cut). The output is
$\Pi^{\mathrm{ce}}(C,K,s)=\{(C_w(e_w),\xi_w):w\text{ a core leaf}\}\cup(\text{all detached colored partitions})$,
with law $\mathcal Q(C,K,s)$.

\emph{Trace data.} For $C=[-L,L]\times J_0$ ($L>0$, $J_0$ a compact interval of positive
length) and $K=\{0\}\times J_0$, all core splits are horizontal, $K_w=\{0\}\times J_w$ and
$g^{D}=g^{U}=0$. Write $Rr=(E^{R}_{w,k},X^{R}_{w,k},G^{R}_{w,k})_{w,k}$ and $Ll$ likewise. The
\emph{straddling cells} are the enclosing boxes, and the trace data are
$\trd(h)=((J_w,b_w,d_w)_{w\text{ born}},(\xi_w)_{w\text{ core leaf}})$: positions and auxiliary times of
the horizontal cuts of straddling cells, and colors of the final straddling
cells. For a side family $Y$ (independent triples with the laws above) define widths
$x_w(t)$: $x_\emptyset(0)=L$, children inherit $x_w(d_w)$, and on $[b_w,e_w)$ the width jumps from
$x$ to $X^{Y}_{w,k}x$ at times $t^{Y}_{w,k}=t^{Y}_{w,k-1}+E^{Y}_{w,k}/x$; for each jump with $t^{Y}_{w,k}<e_w$ put
$S^{Y}_{w,k}=[X^{Y}_{w,k}x,x]\times J_w$, and
\[
\Psi(h,Y)=\{([0,x_w(e_w)]\times J_w,\xi_w):w\text{ core leaf}\}\cup\bigcup_{t^{Y}_{w,k}<e_w}\Pi^{\mathrm{can}}(S^{Y}_{w,k},s-t^{Y}_{w,k};G^{Y}_{w,k}),
\]
which depends on $h$ only through $\trd(h)$. The same definitions with coordinates
exchanged apply to $C=I_0\times[-L,L]$, $K=I_0\times\{0\}$.
\end{definition}

\begin{lemma}[The core-exterior construction]\label{lem:E-coreext}
(i) Almost surely the construction has finitely many events, its output is a
finite colored partition of $C$ and a Borel function of $(C,K,s)$ and the
ingredients; for every born $w$ and $t\in[b_w,e_w]$, $K_w\subseteq C_w(t)$, $C_w(t)\cap K=K_w$,
$C_w(t)$ has nonempty interior, and every detached box is disjoint from $K$.
(ii) The family $\mathcal Q(C,K,s)$ satisfies the first-step recursion, with daughter laws
$\mathcal Q(C_\epsilon(\lambda),K\cap C_\epsilon(\lambda),s-t)$ for a line meeting the relative interior of $K$ in its
cut coordinate, $\mathcal Q(C_\epsilon(\lambda),K,s-t)$ for the daughter containing $K$ and
$\Mcol(C_\epsilon(\lambda),s-t)$ for the other daughter when the line lies in a gap, and
$\Mcol$ for the null set of lines through $\partial K$.
(iii) $\mathcal Q(C,K,s)=\Mcol(C,s)$.
\end{lemma}

\begin{proof}
(i) The core is the canonical construction on $K$. A gap never exceeds $p(C)$, so
inter-jump times are at least $E/p(C)$ and the number of jumps of one gap process
in a window of length at most $s$ is dominated by a $\mathrm{Poi}(sp(C))$ variable; each
detached output is finite (Lemma~\ref{lem:E-nonexplosion}); there are countably many
items. Measurability: decompose by the finite core tree, the finite numbers and
order of side events and the finite detached trees. Invariants by induction: a
jump on side $R$ removes the slab $[b'_w+Xg,b'_w+g]\times[c_w,d_w]\subset\{z_1>b'_w\}$, which misses
$K$ because $C_w\cap K=K_w\subseteq\{z_1\le b'_w\}$; a core split of $K_w$ by $\{z_1=x\}$ gives
$C_{w0}\cap K=K_w\cap\{z_1\le x\}=K_{w0}$; other sides and horizontal splits are analogous; core
splits are strictly inside $K_w$'s extent and gaps are multiplied by $X\in(0,1)$.
(ii) The root clocks are the core clock $A_K=E_\emptyset/p(K)$ (if $p(K)>0$) and the gap
clocks $A_\sigma=E^{\sigma}_{\emptyset,1}/g^{\sigma}(0)$ (if positive), independent of each other and of
all remaining ingredients $Z$; the rates sum to $p(K)+\sum_\sigma g^{\sigma}=p(C)$. By
Lemma~\ref{lem:E-clocks}, $\tau\sim\mathrm{Exp}(p(C))$, $\PP(J=\mathrm{core})=p(K)/p(C)$, $\PP(J=\sigma)=g^{\sigma}(0)/p(C)$.
The first line is a function of $J$ and one uniform position variable (a component
of $Z$): a core line has law $\kappa_K$, a gap line on side $R$ is uniform on $(b',b)$, and so
on; hence $(\tau,\lambda)$ has law $p(C)e^{-p(C)t}dt\,\kappa_C(d\lambda)$, and off a null set $J$ is a
function of $\lambda$. If $\tau\ge s$ the output is $\{(C,\xi_\emptyset)\}$. On a core event at $\tau<s$, the
children's futures are the construction on $(C_\epsilon(\lambda),K_\epsilon,s-\tau)$ driven by distinct
components of $Z$, conditionally independent of $(\tau,\lambda)$ with the ingredient law. On
a gap event on side $\sigma$, the slab receives a canonical output of law
$\Mcol(S,s-\tau)$ (Lemma~\ref{lem:E-recursion}(iii) for $L$, $D$), and the other daughter
$C'$ carries the continuation driven by the family $Y'$ in which the core clock is
replaced by $p(K)(A_K-\tau)$, side $\sigma$ by its shifted sequence with the new gap, and
each other side $\sigma'$ by $g^{\sigma'}(0)(A_{\sigma'}-\tau)$; pathwise every later event of the original
construction is the corresponding event of the construction on $(C',K,s-\tau)$ with
ingredients $Y'$, shifted by $\tau$, and by Lemma~\ref{lem:E-clocks} the scaled residuals are
fresh $\mathrm{Exp}(1)$ variables independent of $\tau$ and of $Z$. Integrating over $(\tau,\lambda)$
gives the recursion. (iii) Apply Lemma~\ref{lem:E-recursion}(ii) to the family indexed
by all triples.
\end{proof}

\begin{corollary}[Restriction and two points]\label{cor:E-restriction}
(i) If $K=B$ has nonempty interior,
\[
\resB_B(\Pi^{\mathrm{ce}}(C,B,s))=\Pi^{\mathrm{can}}(B,s;h),
\]
and $(\resB_B)_*\Mcol(C,s)=\Mcol(B,s)$. (ii) For $z\ne z'$ in $C$ and $K$ the box they span, $z$ and
$z'$ lie in one cell of $\Pi^{\mathrm{ce}}(C,K,s)$ if and only if the core root does not split
before $s$; hence for $\Mcol(C,s)$ the probability that some cell contains both is
$e^{-s|z-z'|_1}$.
\end{corollary}

\begin{proof}
(i) $C_w(e_w)\cap B=K_w$ for core leaves and detached cells miss $B$
(Lemma~\ref{lem:E-coreext}(i)); apply Lemma~\ref{lem:E-coreext}(iii). (ii) If the root does
not split, $z,z'\in K\subseteq C_\emptyset(s)$; if it splits, the line lies strictly between the two
points in one coordinate, and every final cell lies in a detached slab (disjoint
from $K$) or in one closed side of the line. $p(K)=|z-z'|_1$, and the event is Borel.
\end{proof}

\begin{corollary}[Trace decomposition and box positivity]\label{cor:E-tracebox}
Let $C=[-L,L]\times J_0$ and $K=\{0\}\times J_0$. (i) $g^{D}=g^{U}=0$ for all nodes and times; the
right and left gaps evolve as the widths $x_w$ of Definition~\ref{def:E-coreext} driven by
$Rr$ and by $Ll$ respectively. (ii) $\resB_{[0,L]\times J_0}(\Pi^{\mathrm{ce}})=\Psi(h,Rr)$ and
$\theta_1\resB_{[-L,0]\times J_0}(\Pi^{\mathrm{ce}})=\Psi(h,Ll)$. (iii) $h$, $Rr$, $Ll$ are independent,
$Rr\overset{d}{=}Ll$, and $\Psi(h,Y)=\hat\Psi(\trd(h),Y)$ for a Borel $\hat\Psi$. (iv) For $N$ independent
copies and bounded complex Borel $\Phi$ on $\Part^{\mathrm{col}}([0,L]\times J_0)^{N}$,
$\E[\conj{\Phi((\Psi(h^{i},Ll^{i}))_i)}\Phi((\Psi(h^{i},Rr^{i}))_i)]=\E|\bar\Phi(h^{1},\dots,h^{N})|^{2}$ with
$\bar\Phi=\int\Phi((\Psi(h^{i},y^{i}))_i)\prod_iP_{Rr}(dy^{i})$, a function of the trace data. (v) For independent
$\Pi^{1},\dots,\Pi^{N}$ with law $\Mcol(C,s)$, $T^{i}_{+}=\resB_{[0,L]\times J_0}(\Pi^{i})$ and
$T^{i}_{-}=\theta_1\resB_{[-L,0]\times J_0}(\Pi^{i})$, $\E[\conj{\Phi(T^{1}_{-},\dots,T^{N}_{-})}\Phi(T^{1}_{+},\dots,T^{N}_{+})]\ge0$; the
same holds across $\theta_2$ for $C=I_0\times[-L,L]$.
\end{corollary}

\begin{proof}
(i) The vertical extents of $C$ and $K$ coincide initially, core lines are
horizontal ($K$ has zero width) and cut $C_w$ and $K_w$ at the same height, and side
jumps $L$, $R$ do not change vertical extents; at a horizontal straddling cut both
daughters inherit the current widths and then use their own triples. (ii) Final
straddling cells are $[-l_w(e_w),r_w(e_w)]\times J_w$ and restrict to $[0,r_w(e_w)]\times J_w$; $R$-slabs
$[Xr,r]\times J_w$ carry $\Pi^{\mathrm{can}}(\cdot;G^{R})$ and $L$-slabs carry $\theta_1\Pi^{\mathrm{can}}([Xl,l]\times J_w;G^{L})$; cells
of one half-box meet the other in at most a segment and are discarded by the
restriction. (iii) The ingredients are independent by definition, and $\Psi$ uses $h$
only through the born-node data and the leaf colors; measurability as in
Lemma~\ref{lem:E-coreext}(i). (iv) By Fubini,
$\E[\Phi((\Psi(h^{i},Rr^{i}))_i)\mid h]=\bar\Phi(h)$ and $\E[\conj{\Phi((\Psi(h^{i},Ll^{i}))_i)}\mid h]=\conj{\bar\Phi(h)}$, and the
two are conditionally independent given $h$. (v) By Lemma~\ref{lem:E-coreext}(iii) the
law of $(T^{i}_{+},T^{i}_{-})_i$ is that of $N$ independent copies of $(\Psi(h^{i},Rr^{i}),\Psi(h^{i},Ll^{i}))$;
apply (iv). Only the complete trace data are used; the finite-box hand checks of
the endpoint record exhibit loss of conditional independence when either the
auxiliary times or the positions are dropped from the conditioning.
\end{proof}

\begin{proof}[Proof of Theorems~\ref{thm:E-box} and \ref{thm:E-trace}]
Theorem~\ref{thm:E-box}: (a) Lemma~\ref{lem:E-nonexplosion} and the measurability
paragraph before it; (b) Lemma~\ref{lem:E-recursion}(i),(ii); (c)
Lemma~\ref{lem:E-recursion}(iii); (d) Corollary~\ref{cor:E-restriction}(i); (e)
Corollary~\ref{cor:E-restriction}(ii) and Lemma~\ref{lem:E-nonexplosion}.
Theorem~\ref{thm:E-trace}: Lemma~\ref{lem:E-coreext}(iii),
Corollary~\ref{cor:E-restriction}(i) and Corollary~\ref{cor:E-tracebox}.
\end{proof}

\subsection{The whole-plane Mondrian neuron}

\begin{proof}[Proof of Theorem~\ref{thm:E-plane}(a)]
The bonding maps $\pi_n=\resB_{B_n}$ are Borel with
$(\pi_n)_*\Mcol(B_{n+1},\mu)=\Mcol(B_n,\mu)$ (Corollary~\ref{cor:E-restriction}(i)); drawing
$P_m\sim\Mcol(B_m,\mu)$ and setting $P_n=\pi_n\circ\cdots\circ\pi_{m-1}(P_m)$ gives consistent
finite-dimensional laws, and (K5) gives a probability on $\prod_n\Part^{\mathrm{col}}(B_n)$, under
which the Borel set $S$ of consistent sequences has full measure. For any
consistent tessellation $(P_1,\dots,P_m)$ is almost surely the image of $P_m$ under the
bonding maps, so the law is unique (K1). Restriction is transitive
($(R\cap B_n)\cap B=R\cap B$, and a nonempty interior of $R\cap B$ forces one of $R\cap B_n$), so
$\resB_B(P_m)=\resB_B(P_n)$ on $S$, with law $\Mcol(B,\mu)$. For $T=\theta_j$, $(TP)_n=TP_n$ is
consistent because $TB_n=B_n$, with marginals $\Mcol(B_n,\mu)$
(Lemma~\ref{lem:E-recursion}(iii)); for a translation by $x$ and
$m=n+\lceil|x|_\infty\rceil$, $(TP)_n=\resB_{B_n}(TP_m)$ does not depend on the admissible $m$, is
consistent, and has law $\Mcol(B_n,\mu)$; uniqueness gives $TP\overset{d}{=}P$. The last
sentence of (a) is part (A) of the next proof.
\end{proof}

\begin{proof}[Proof of Theorem~\ref{thm:E-plane}(b),(c)]
(A) Let $P\in\Part^{\mathrm{col}}(D)$, $D=[a_D,b_D]\times[c_D,d_D]$ with nonempty interior, and
$z\in[a_D,b_D)\times[c_D,d_D)$. The points $z+e(1,1)$, $0<e<e_0$, lie in $\mathrm{int}\,D$ and avoid the
finitely many cell boundaries (an axis-parallel segment meets the open diagonal
ray in at most one point), so they lie in the interior of one cell $R$, and then
$z\in R^{\mathrm{ll}}$; conversely $z\in R^{\mathrm{ll}}$ forces $z+e(1,1)\in\mathrm{int}\,R$ for small $e$. So the lower-left
representatives partition $[a_D,b_D)\times[c_D,d_D)$. (B) On the consistency event, for
$z\in[-n,n)^{2}$: $R_n(z)=R_{n+1}(z)\cap B_n$, $\xi_n(z)=\xi_{n+1}(z)$ and $R_n(z)^{\mathrm{ll}}\subseteq R_{n+1}(z)^{\mathrm{ll}}$, since
$z$ lies in the representative of $R_{n+1}(z)\cap B_n$, a cell of $P_n$ with the same color;
uniqueness in (A) gives the claim. (C) For every outcome
$g(z)=\sum_n\ind\{n_z=n\}\int hh(z,r)\,P_n(dr)$ with
$hh(z,(a,b,c,d,\xi))=(\xi/2)\ind\{a\le z_1<b,c\le z_2<d\}$, exactly one atom contributing by (A); so $g$
is $\cB(\R^{2})\otimes\cA$-measurable (Lemma~\ref{lem:E-partspaces}(iii)) with $|g|=1/2$. (D) If
$R(z)\not\subseteq z+(-m,m)^{2}$, some $R_n(z)$ contains one of $z\pm me_1$, $z\pm me_2$; the events that
some cell of $P_n$ contains two fixed points increase in $n$ and have probability
$e^{-\mu|y-y'|_1}$ (Corollary~\ref{cor:E-restriction}(ii)); so
$\PP(R(z)\not\subseteq z+(-m,m)^{2})\le4e^{-\mu m}$ and almost surely every $R(q)$, $q\in\Q^{2}$, is bounded. Every $z$
has a rational $q$ in the interior of $R_{n_z}(z)^{\mathrm{ll}}$, and then $R(q)=R(z)$, $H(q)=H(z)$. (E) If
$R(z)\subseteq B_k$, $k\ge n_z$, then $R_n(z)=R_{n+1}(z)$ for $n\ge k$, so $R(z)=R_k(z)$ and $H(z)=R(z)^{\mathrm{ll}}$. If
$w\in H(z)\cap H(z')$, then $R_n(z)=R_n(z')=R_n(w)$ for large $n$; so distinct representatives
are disjoint and cover $\R^{2}$; $R(z)=\mathrm{cl}\,H(z)$ and $\mathrm{int}\,R(z)=\mathrm{int}\,H(z)$; a representative
meeting $B_k$ contains a representative of a cell of $P_{k+1}$, so finitely many meet
$B_k$; $g=\xi(z)/2$ is constant on $H(z)$. (F) Under $\Mcol(B_n,\mu)$ the colors are independent
uniform given the uncolored partition, so $\E g(z)=0$ and
$\E g(z)g(z')=\tfrac14\PP(R_n(z)=R_n(z'))=\tfrac14e^{-\mu|z-z'|_1}$, using that $z,z'$ are almost surely interior
to cells (Lemma~\ref{lem:E-nonexplosion}).
\end{proof}

\begin{proof}[Proof of Theorem~\ref{thm:E-plane}(d)]
On the consistency event, if $z$ is interior to a cell of $P_{n_z+1}$, that cell is the
only one containing $z$, its color is $\xi(z)$, and every admissible convention gives
$g'(z)=g(z)$; the exceptional points lie on countably many segments. Let
$Y_{\mathrm{sp}}=\prod_n\Part^{\mathrm{col}}(B_n)$ and $G(z,(P'_n)_n)=\sum_n\ind\{n_z=n\}\int hh(z,r)\,P'_n(dr)$. Then $G$ is
product measurable and $|G|=1/2$ on all of $\R^{2}\times Y_{\mathrm{sp}}$, consistent or not, because
the representatives of any element of $\Part^{\mathrm{col}}(B_{n_z})$ partition $[-n_z,n_z)^{2}$ (part
(A) above uses no consistency). So $\Gamma^{1/2}_{G}$ of Lemma~\ref{lem:E-fields}(iii) is a total
measurable map with $[g(\cdot,\omega)|_U]=\Gamma^{1/2}_{G}((P_n(\omega))_n)$ for every outcome, and the law of
$g|_U$ is the image of the unique law of $(P_n)_n$. An admissible $g'$ is clipped to
$F'=\clip_{1/2}(g')$, which equals $g'$ on the consistency event and $g$ almost everywhere
there, so $g'|_U$ has the same law. For $T=\theta_j$, the field of $TP$ equals $g\circ\theta_j$ off
a null set on the consistency event and is $G(\cdot,(TP_n)_n)$, so
$g\circ\theta_j\overset{d}{=}g$ because $TP\overset{d}{=}P$; translations are the same with
$(TP)_n=\resB_{B_n}(TP_m)$.
\end{proof}

\subsection{The product neuron}

\begin{proof}[Proof of Theorem~\ref{thm:E-product}]
(a) The Poisson process $\cN$ is almost surely locally finite, has no atom at $0$ and
has unbounded atoms in both directions; then $\sigma$ is piecewise constant and
c\`adl\`ag, and $|\sigma|=1$. $\E\sigma(t)=0$; for $t<t'$, $\sigma(t)\sigma(t')=(-1)^{\cN((t,t'])}$ in each sign
configuration, so $\E\sigma(t)\sigma(t')=e^{-2(\mu/2)|t-t'|}$ (K6); given $\sigma(t_1)$, parities on disjoint
later intervals are independent, which gives the Markov chain with transition
probabilities $(1+e^{-\mu u})/2$.

(b) A right-continuous process is jointly measurable
($\sigma(t)=\lim_k\sigma(2^{-k}\lceil2^{k}t\rceil)$), so $g$ is $\cB(\R^{2})\otimes\cA$-measurable, $|g|=1/2$,
$\E g=0$ and $\E g(z)g(z')=\tfrac14e^{-\mu|z_1-z'_1|}e^{-\mu|z_2-z'_2|}$; $g$ is constant on the boxes
$[t_i,t_{i+1})\times[y_j,y_{j+1})$ of the jump points.

(c) Across $\theta_1$ let $b=(\epsilon_\sigma,\tau)$, $M_+=\cN|_{(0,\infty)}$ and $M_-$ the image of $\cN|_{(-\infty,0)}$
under $t\mapsto-t$. These are random elements of the space $\mathsf N$ of
$\{0,1,\dots,\infty\}$-valued measures on $(0,\infty)$ (with the $\sigma$-algebra of the maps
$m\mapsto m(A)$); by (K6) they are independent with the same law $P_{\mathsf M}$ and independent of
$b$. Let $D_\pm$ be the c\`adl\`ag functions $\R\to\{-1,1\}$ with the evaluation $\sigma$-algebra;
$(t,x)\mapsto x(t)$ is product measurable by right continuity. Put
$Y_{\mathrm{prod}}=\{-1,1\}\times D_\pm\times\mathsf N$, $\mathrm{Good}_+=\{m:m((0,k])<\infty\ \forall k\}$,
$\mathrm{par}(m,t)=(-1)^{m((0,t])}$ on $\mathrm{Good}_+$ and $1$ otherwise, and
$G_{\mathrm{prod}}(z,(e,x,m))=\tfrac12e\,\mathrm{par}(m,z_1)x(z_2)$ for $z_1>0$, $0$ otherwise. Then $G_{\mathrm{prod}}$ is
product measurable with $|G_{\mathrm{prod}}|\le1/2$ everywhere, and Lemma~\ref{lem:E-fields}(iii) gives
a total measurable $\Psi_{\mathrm{prod}}:Y_{\mathrm{prod}}\to\Ball_{1/2}(H_+)$. On the almost sure event of (a),
$s_+=G_{\mathrm{prod}}(\cdot,(\epsilon_\sigma,\tau,M_+))$ on $H_+$, and $s_-(z)=\tfrac12\epsilon_\sigma(-1)^{\cN((-z_1,0])}\tau(z_2)$ differs
from $G_{\mathrm{prod}}(z,(\epsilon_\sigma,\tau,M_-))$ only on countably many vertical lines; so
$[s_\pm]=\Psi_{\mathrm{prod}}(b,M_\pm)$ almost surely, and Fubini gives
$\E[f(s_+)f'(s_-)\mid b]=\bar f(b)\bar f'(b)$ with $\bar f(b)=\int f(\Psi_{\mathrm{prod}}(b,m))P_{\mathsf M}(dm)$.
Across $\theta_2$ exchange $\sigma$ and $\tau$.

(d) For independent neurons, $(b_i)_i$, $(M^{i}_{+})_i$, $(M^{i}_{-})_i$ are independent families, the
last two independent and identically distributed with the same law, and
$S_\pm=N^{-1/2}\sum_i\Psi_{\mathrm{prod}}(b_i,M^{i}_{\pm})$ almost surely; every value of
$N^{-1/2}\sum_i\Psi_{\mathrm{prod}}(y_i)$ lies in $\Ball_{\sqrt N/2}(H_+)$. The same Fubini computation with the
complete tuple $(b_1,\dots,b_N)$ gives conditional independence and equal conditional
laws. With only $b_1$ given and $N=2$, neuron $2$ is shared by both sides:
$\mathrm{Cov}(F_2(t,0),F_2(-t,0)\mid b_1)=e^{-2\mu t}/8>0$ for $t>0$, and point values are almost surely
Borel functions of the sides (proof of Theorem~\ref{thm:E-RP}(c) below), so the sides
are not conditionally independent given $b_1$ alone.
\end{proof}

\subsection{Finite-width annealed reflection positivity}

Fix $b$, $N\ge1$, $\theta$, $\Omega$ and $v$ as in Theorem~\ref{thm:E-RP}; $|F_N|\le r_N$ everywhere.

\begin{lemma}[Projections]\label{lem:E-projections}
If $\cG_1\subseteq\cG_2\subseteq\cdots$ are $\sigma$-algebras, $\cG_\infty=\sigma(\bigcup_n\cG_n)$ and $Y\in L^{2}(\cG_\infty)$, then
$\E[Y\mid\cG_n]\to Y$ in $L^{2}$.
\end{lemma}

\begin{proof}
The conditional expectations are orthogonal projections onto increasing closed
subspaces $V_n$; the events $A\in\cG_\infty$ with $\ind_A$ in the closure of $\bigcup V_n$ form a
Dynkin system containing the $\pi$-system $\bigcup\cG_n$, so the closure is $L^{2}(\cG_\infty)$; for
$W\in V_m$ with $\|Y-W\|<e$ and $n\ge m$, $\|Y-\E[Y\mid\cG_n]\|\le\|Y-W\|$.
\end{proof}

\begin{lemma}[Weight factorization and reduction to the sides]\label{lem:E-factor}
Let $\Psi_+(s)=\int_{\Omega\cap H_+}v(z,s(z))\,dz$ on $\Ball_{r_N}(H_+)$ (Borel by
Lemma~\ref{lem:E-potential}). Then $U=\Psi_+(S_+)+\Psi_+(S_-)$ for every outcome. For bounded
complex $\Sigma^{\theta}_{+}$-measurable $F$, with the representation of
Proposition~\ref{prop:space}(x), $\Phi_F(s)=g((\int f_ks)_k)$ is bounded Borel on $\Ball_{r_N}(H_+)$,
$F(\iota F_N)=\Phi_F(S_+)$, $F(\theta\iota F_N)=\Phi_F(S_-)$, and
$\int_X\conj{F(\theta v)}F(v)\,P^{v}_{N}(dv)=Z^{-1}\E[\conj{\Phi(S_-)}\Phi(S_+)]$ with $\Phi=\Phi_Fe^{-\Psi_+}$.
\end{lemma}

\begin{proof}
The line $\partial H_+$ is null; the substitution $z'=\theta z$, $\theta(\Omega\cap H_+)=\Omega\cap H_-$ and
$v(\theta z',x)=v(z',x)$ give $\Psi_+(S_-)=\int_{\Omega\cap H_-}v(z',F_N(z'))\,dz'$. For the second part,
$\int f_kF_N=\int f_kS_+$ and $\int f_k(F_N\circ\theta)=\int f_kS_-$ because $\supp f_k\subset H_+$, and
$\theta\iota F_N=\iota(F_N\circ\theta)$ (Proposition~\ref{prop:space}(vii)).
\end{proof}

\begin{proposition}[Unweighted side positivity]\label{prop:E-sides}
$(S_+,S_-)\overset{d}{=}(S_-,S_+)$, and $\E[\conj{\Phi(S_-)}\Phi(S_+)]\ge0$ for every bounded complex
Borel $\Phi$ on $\Ball_{r_N}(H_+)$.
\end{proposition}

\begin{proof}
\emph{Product neurons.} By Theorem~\ref{thm:E-product}(c),(d),
$S_\pm=N^{-1/2}\sum_i\Psi_{\mathrm{prod}}(b_i,M^{i}_{\pm})$ with independent families and equal laws of the
last two, which gives the symmetry, and Fubini gives
$\E[\conj{\Phi(S_-)}\Phi(S_+)]=\E|\bar\Phi(b_1,\dots,b_N)|^{2}$ with
$\bar\Phi(b)=\int\Phi(N^{-1/2}\sum_i\Psi_{\mathrm{prod}}(b_i,y_i))\prod_iP_{\mathsf M}(dy_i)$.

\emph{Mondrian neurons.} Step 1: $g_i\circ\theta\overset{d}{=}g_i$ (Theorem~\ref{thm:E-plane}(d)), so
$F_N\circ\theta\overset{d}{=}F_N$, and the continuous map $w\mapsto(w|_{H_+},(w\circ\theta)|_{H_+})$ gives the
symmetry. Step 2: take $\theta=\theta_1$, $B^{+}_{n}=[0,n]\times[-n,n]$, $Q_n=(0,n)\times(-n,n)$. On the
consistency events, off the null cell boundaries, $g_i=\fld(\resB_{B^{+}_{n}}(P^{i}_{n}))$ and
$g_i\circ\theta_1=\fld(\theta_1\resB_{[-n,0]\times[-n,n]}(P^{i}_{n}))$ on $Q_n$; so $r_n(S_\pm)=\Lambda_n((T^{i}_{\pm,n})_i)$ with
the Borel map $\Lambda_n((P'^{i})_i)=N^{-1/2}\sum_i\fld(P'^{i})|_{Q_n}$, and
Corollary~\ref{cor:E-tracebox}(v) with $C=B_n$ gives
$\E[\conj{\Phi_n(r_nS_-)}\Phi_n(r_nS_+)]\ge0$ for bounded Borel $\Phi_n$ on $L^{1}(Q_n)$. Step 3: by
Lemma~\ref{lem:E-restrictions}(ii) the $\sigma$-algebras $\cG_n=\sigma(r_n(S_+))$ increase and generate
$\sigma(S_+)$; with $Y=\Phi(S_+)$ and $\E[Y\mid\cG_n]=\Phi_n(r_nS_+)$ (K2; truncated to $|\Phi_n|\le\sup|\Phi|$),
Lemma~\ref{lem:E-projections} gives $\Phi_n(r_nS_+)\to\Phi(S_+)$ in $L^{2}$, Step 1 gives the same for
the minus side, and Cauchy--Schwarz passes the inequality to the limit.
\end{proof}

\begin{proof}[Proof of Theorem~\ref{thm:E-RP}]
(a) $(z,\omega)\mapsto v(z,F_N(z,\omega))$ is $\cL(\R^{2})\otimes\cA$-measurable, so $U$ is a random
variable with $|U|\le|\Omega|M_N$ and the stated bounds on $Z$ (no bound $Z\ge1$ is claimed).
$F_N$ is a Borel random element of $\Ball_{r_N}(\R^{2})$ and of $E$ (Lemma~\ref{lem:E-fields}), and
on $\Ball_{r_N}$ the inclusion into $E$ is continuous ($|s_k-s|^{4}\le(2r_N)^{3}|s_k-s|$), so $\iota F_N$ is
Borel and $P^{v}_{N}$ is a Borel probability. $U=\Psi_v(F_N)$ (Lemma~\ref{lem:E-potential}), so
$P^{v}_{N}$ depends only on the law of $F_N$, which is the image of the $N$-fold product
of one neuron's law under the continuous sum map; that law is determined by the
finite-dimensional laws for product neurons (Lemma~\ref{lem:E-fields}(ii)) and by
Theorem~\ref{thm:E-plane}(d) for Mondrian neurons. Admissible conventions change
$F_N$ only on a null set of $z$, on the almost sure intersection of the consistency
events, so $U$ and $\iota F_N$ are unchanged almost surely.
(b) Lemma~\ref{lem:E-factor} gives $e^{-U}\conj{\Phi(S_-)}\Phi(S_+)=\conj{\Phi'(S_-)}\Phi'(S_+)$ with
$\Phi'=\Phi e^{-\Psi_+}$; apply Proposition~\ref{prop:E-sides}, and for $X$-observables
Lemma~\ref{lem:E-factor} and $Z>0$.
(c) The first assertion is the computation in (b). For each fixed $z$, almost surely
$F_N$ is constant near $z$ (for Mondrian neurons $z$ is interior to a cell of every $P^{i}_{n}$
with $z\in\mathrm{int}\,B_n$; for product neurons $z_1,z_2$ are not atoms of the locally finite
Poisson processes). Hence, almost surely for all $k$, $L_{z_k}(S_+)=F_N(z_k)$ and
$L_{z_k}(S_-)=F_N(\theta z_k)$ (Lemma~\ref{lem:E-restrictions}(iii)), so $Y=\Phi(S_+)$ and
$Y^{\theta}=\Phi(S_-)$ with the bounded Borel $\Phi=\tilde Y\circ(L_{z_k})_k$. If $\Phi'$ comes from another
representation, $\Phi(S_+)=\Phi'(S_+)$ almost surely, and $S_-\overset{d}{=}S_+$ gives $\Phi(S_-)=\Phi'(S_-)$
almost surely.
\end{proof}

\begin{proof}[Proof of Corollaries~\ref{cor:E-P2} and \ref{cor:E-general}]
Corollary~\ref{cor:E-P2}: each coordinate interval of $\Lambda$ and $\Omega$ is symmetric
about $0$, so both rectangles are invariant under $\theta_1$ and $\theta_2$; $v$ is Borel,
bounded on $\Omega\times[-r,r]$ by $\max_{|x|\le r}(|q_{\mathrm{in}}|+|q_{\mathrm{out}}|+\kappa^{2}x^{2}/2)$ and invariant under both
reflections on $\Omega$; Theorem~\ref{thm:E-RP} applies to each $\theta$ separately. For a
physical square, $\bar\Lambda=[-l,l]^{2}$ and $\Omega=[-R,R]^{2}$ are such rectangles; off
$\partial\Lambda$, $\ind_\Lambda=\ind_{\bar\Lambda}$ and $\ind_{\Omega\setminus\Lambda}=\ind_{\Omega\setminus\bar\Lambda}$, so $v_V$ agrees with the displayed
$v$ off the null set $\partial\Lambda\times\R$, and for every outcome
$\int_\Omega v_V(z,F_N(z))\,dz=V_{\varepsilon,\Omega,a;\Lambda}(F_N)$, which is also the $U$ of the displayed $v$.
Axis-parallelism is used exactly in the identification $\bar\Lambda=[-l,l]^{2}$; for the
rotated square of Remark~\ref{rem:rotated}(i) no conclusion, positive or negative,
is drawn. The hypotheses ``$\Omega$ centered'' and ``$\Lambda$ centered'' are both used: an
outer rectangle that is invariant under $\theta_1$ only does not satisfy the
hypothesis for $\theta_2$, and no claim is made for it.
Corollary~\ref{cor:E-general}: the hypotheses of Theorem~\ref{thm:E-RP} hold for each
$\theta$.
\end{proof}

\begin{remark}[What the endpoint results do not contain]\label{rem:E-notcontained}
The conditional identities of Corollary~\ref{cor:E-tracebox} are proved on the extended
probability space of the core-exterior construction in a box; the auxiliary
times are not functions of the final tessellation, and no whole-plane
auxiliary-time history is constructed. Not needed and not proved: the Poisson law
of trace cut points; reflection positivity across shifted or rotated lines; any
Euclidean invariance of tilted laws; reflection positivity of the limits in $N$,
$\varepsilon$, $\Omega$ or $a$ by these arguments; and reflection positivity, or its failure, for
tilts built on rotated squares that are not invariant under both reflections.
The measurability hypothesis is used exactly in Lemma~\ref{lem:E-potential} and in
part (a) of Theorem~\ref{thm:E-RP}; the broader reading of Remark~\ref{rem:rotated}(iii)
fails.
\end{remark}

% ---------------------------------------------------------------------------
\section{Proofs for the Gaussian reference fields}\label{app:G}

Besides \ref{prem:C6}(c)--(g) the following elementary Gaussian facts are used:
an $L^{2}$-limit of centered jointly Gaussian vectors is centered jointly Gaussian;
uncorrelated jointly Gaussian families are independent; for $\xi\sim N(0,\lambda)$, $\lambda>0$,
$t\lambda<1$: $\E e^{t\xi^{2}/2}=(1-t\lambda)^{-1/2}$, and the law with density proportional to
$e^{tx^{2}/2}$ relative to $N(0,\lambda)$ is $N(0,\lambda/(1-t\lambda))$; a centered Gaussian of variance $v$
has $\E|Y|^{2n}=(2n-1)!!v^{n}$; and \emph{regression}: if $(Y,\xi_1,\dots,\xi_N)$ is centered jointly
Gaussian, $Y\in\R^{n}$, with independent $\xi_j$ of variances $\lambda_j>0$, then
$Y-\sum_j\E[Y\xi_j]\xi_j/\lambda_j$ is independent of $(\xi_j)$ with covariance
$\mathrm{Cov}\,Y-\sum_j\E[Y\xi_j]\E[Y\xi_j]^{T}/\lambda_j$. The path space $C(\R^{2})$ carries the metric of
uniform convergence on compacts; it is separable, its Borel $\sigma$-algebra is generated
by the evaluations at rational points, and two Borel probabilities on it with the
same finite-dimensional laws at rational points coincide.

\begin{lemma}[Spectral point fields]\label{lem:G-S}
Let $\nu$ be an even $\sigma$-finite Borel measure on $\R^{2}$,
$H_\nu=\{h\in L^{2}(\nu;\C):h(-p)=\conj{h(p)}\}$ with $\langle g,h\rangle=\int\conj gh\,d\nu$, $W$ an isonormal process
on $H_\nu$, $\sigma\ge0$ even and Borel with $S_\sigma=\int\sigma^{2}d\nu<\infty$, and $X(z)=W(e_z\sigma)$.
\begin{enumerate}[label=(S\arabic*),start=0]
\item $H_\nu$ is a separable real Hilbert space.
\item $X$ is centered Gaussian with $\E[X(z)X(z')]=k(z-z')$, $k(x)=\int\cos(p\cdot x)\sigma^{2}d\nu$,
  real, even, continuous, $|k|\le k(0)=S_\sigma$; and $\E[X(z)W(h)]=\int e^{ip\cdot z}\sigma h\,d\nu$.
\item $\E|X(z+h)-X(z)|^{2}=\int2(1-\cos(p\cdot h))\sigma^{2}d\nu=2(k(0)-k(h))$.
\item If there are constants $M>0$ and $\vartheta>0$ with $2(k(0)-k(h))\le M|h|^{\vartheta}$ for
  all $h\in\R^{2}$, there is a $\cB(\R^{2})\otimes\cF$-measurable $\tilde X$ with
  continuous paths for every outcome, $\PP(\tilde X(z)=X(z))=1$ for every $z$, which is a
  Borel $C(\R^{2})$-valued random element whose law is the unique Borel law with
  centered Gaussian coordinates of covariance $k(z-z')$.
\item Under (S3), for real $g\in L^{1}(\R^{2})$: $\hat g\sigma\in H_\nu$, $\int|g||\tilde X|<\infty$ almost surely, and
  $\int g\tilde X\,dz=W(\hat g\sigma)$ almost surely.
\end{enumerate}
\end{lemma}

\begin{proof}
(S0) $p\mapsto-p$ preserves $\nu$, so $\langle g,h\rangle$ is real; $H_\nu$ is closed in $L^{2}(\nu)$ (an
a.e.-convergent subsequence preserves the symmetry) and separable. (S1)
$e_z\sigma\in H_\nu$; isometry gives $\E[X(z)X(z')]=\int e^{ip\cdot(z-z')}\sigma^{2}d\nu$, whose sine part vanishes
by evenness; continuity by dominated convergence; all $X(z)$ lie in the range of $W$.
(S2) $|e^{-it}-1|^{2}=2(1-\cos t)$. (S3) Increments are Gaussian with
$\E|X(z+h)-X(z)|^{2n}\le(2n-1)!!M^{n}|h|^{n\vartheta}$; for an integer $n>2/\vartheta$ (which
exists because $\vartheta>0$) the continuity criterion \ref{prem:C6}(d) applies with
$\alpha=2n$, $\beta=n\vartheta-2>0$. The applications below use $(M,\vartheta)=(\mu/\sqrt2,1)$,
$(M'_{a,1/2},1/2)$ and $(1/(4\pi),2)$, all with $M,\vartheta>0$. Joint measurability: the dyadic
approximations $\tilde X(2^{-l}\lfloor2^{l}z\rfloor,\omega)$ are product measurable and converge everywhere by
continuity. Borel measurability and the law follow from the path-space facts.
(S4) $|\hat g|\le\|g\|_1$ and $\hat g(-p)=\conj{\hat g(p)}$; $\E\int|g||\tilde X|\le\|g\|_1S_\sigma^{1/2}$ and
$\E(\int|g||\tilde X|)^{2}\le\|g\|_1^{2}S_\sigma$, which justify Fubini in:
$\E I^{2}=\iint g(z)g(z')k(z-z')=\int|\hat g|^{2}\sigma^{2}d\nu$,
$\E[IW(\hat g\sigma)]=\int g(z)\E[X(z)W(\hat g\sigma)]dz=\int|\hat g|^{2}\sigma^{2}d\nu$ and
$\E W(\hat g\sigma)^{2}=\int|\hat g|^{2}\sigma^{2}d\nu$, with $I=\int g\tilde X$; hence $\E|I-W(\hat g\sigma)|^{2}=0$.
\end{proof}

\begin{proof}[Proof of Theorem~\ref{thm:G-endpoint}]
Take $\nu=(2\pi)^{-2}\,dp$ and $\sigma=(\hCend)^{1/2}$. In one variable
$\int e^{-iwt}e^{-\mu|t|}dt=2\mu/(\mu^{2}+w^{2})$, so $\hCend$ is as in \eqref{eq:covariances},
$\int\hCend\,dp=\pi^{2}$, $S_\sigma=1/4$, and by Fourier inversion $k=\Cend$. Then
$2(k(0)-k(h))=(1-e^{-\mu|h|_1})/2\le\mu|h|_1/2\le\mu|h|/\sqrt2$, and Lemma~\ref{lem:G-S}(S3) with
$\vartheta=1$ gives $\gamma_\mu$ and the law $G_\mu$.
\end{proof}

\begin{proof}[Proof of Theorem~\ref{thm:G-shift}]
(a) $\hat C_a\le\hat C_0\le m^{-2}$; since $a^{2}m^{2}<1$, termwise
$m^{2}(1+a^{2}w^{2})(1+a^{2}q^{2})\le m^{2}+w^{2}+q^{2}+a^{2}w^{2}q^{2}$, which gives the product bound and
$\int\hat C_a\le m^{-2}\pi^{2}/a^{2}$; $\int_\R dw/(A+Bw^{2})=\pi/\sqrt{AB}$ with $A=m^{2}+q^{2}$, $B=1+a^{2}q^{2}$ gives
\eqref{eq:ca}, whose integrand is at most $1/m$ and at most $1/(aq^{2})$.
(b) $1/\hCend=(\mu^{2}+w^{2})(\mu^{2}+q^{2})/\mu^{2}=\mu^{2}+w^{2}+q^{2}+a^{2}w^{2}q^{2}$, and $\hCend\le\mu^{-2}$ gives
$\kappa_a^{2}\hCend\le r_a<1$, so the geometric series converges pointwise.
(c) $s_0=\Cend$ is continuous, positive, $\sup s_0=1/4$, $\|s_0\|_1=\mu^{-2}$; inductively
$s_{k+1}=\kappa_a^{2}\Cend*s_k$ has $\sup s_{k+1}\le r_a\sup s_k$, $\|s_{k+1}\|_1=r_a\|s_k\|_1$ (Tonelli) and is
continuous (continuity of translation in $L^{1}$); $\hat s_k=\kappa_a^{2k}(\hCend)^{k+1}$. The sum
$S=\sum s_k$ converges uniformly and in $L^{1}$; $\hat S=\hat C_a$ by (b), and inversion gives
$S=C_a$. Hence $C_a\ge\Cend>0$, $|C_a|\le C_a(0)=c_a$, $\|C_a\|_1=\hat C_a(0)=m^{-2}$, and
$\|C_a\|_2^{2}=(2\pi)^{-2}\int\hat C_a^{2}\le(2\pi)^{-2}\int(m^{2}+|p|^{2})^{-2}dp=1/(4\pi m^{2})$.
(d) Convolution by $\Cend$ is the multiplier $\hCend$ with norm $\hCend(0)=\mu^{-2}$; Young's
inequality and the uniform-and-$L^{1}$ convergence give $C_a=\sum_k\kappa_a^{2k}(\Cend)^{k+1}$ in operator
norm, $=\Cend(I-\kappa_a^{2}\Cend)^{-1}$; the quadratic form is
$(2\pi)^{-2}\int|\hat f|^{2}\hat C_a\in[0,m^{-2}\|f\|_2^{2}]$.
(e) $2(c_a-C_a(h))=(2\pi)^{-2}\int2(1-\cos(p\cdot h))\hat C_a\,dp$. Use $1-\cos t\le\min(2,|t|)$,
$\min(2,|\alpha+\beta|)\le\min(2,|\alpha|)+\min(2,|\beta|)$, the product bound of (a) and
$\int(1+a^{2}q^{2})^{-1}dq=\pi/a$: the right side is at most
$(2\pi)^{-2}2m^{-2}(\pi/a)[I_a(h_1)+I_a(h_2)]$ with $I_a(s)=\int\min(2,|w||s|)(1+a^{2}w^{2})^{-1}dw$. Splitting at
$|w|=2/|s|$, $I_a(s)\le(|s|/a^{2})[\log(1+4a^{2}/s^{2})+2]$, and the prefactor is
$(2\pi m^{2}a^{3})^{-1}$; this is $\varpi_a(h)^{2}$, for every $h$.
(f) For $y\ge0$ and $\beta\in(0,1]$, $\log(1+y)\le y^{\beta}/\beta$ (the difference has nonnegative
derivative). With $\beta=(1-\alpha)/2$, $|h_i|\log(1+4a^{2}/h_i^{2})\le2(2a)^{1-\alpha}|h_i|^{\alpha}/(1-\alpha)$; for $|h|\le1$,
$2|h_i|\le2|h_i|^{\alpha}$ and $\sum_i|h_i|^{\alpha}\le2|h|^{\alpha}$, giving $M_{a,\alpha}|h|^{\alpha}$; for $|h|\ge1$, $C_a(h)\ge0$ gives
$2(c_a-C_a(h))\le2c_a|h|^{\alpha}$.
\end{proof}

\begin{proof}[Proof of Theorem~\ref{thm:G-K}]
Let $B$ be the $L^{2}$ inverse transform of $(\hCend)^{1/2}$: real, even,
$\|B\|_2^{2}=1/4$, and $B*B=\Cend$ pointwise. $(Tf)(u)=\int_\Omega B(u-x)f(x)\,dx$ is
Hilbert--Schmidt from $L^{2}(\Omega)$ to $L^{2}(\R^{2})$ with $\|T\|_{\mathrm{HS}}^{2}=|\Omega|/4$, and Fubini gives
$\langle Tf,Tg\rangle=\langle f,K_\Omega g\rangle$, so $K_\Omega=T^{*}T$ is positive and trace class with $\trace K_\Omega=|\Omega|/4$.
$\langle f,K_\Omega f\rangle=\langle R^{*}f,\Cend R^{*}f\rangle\le\mu^{-2}\|f\|_2^{2}$, with $R^{*}$ extension by zero. The kernel is
real and symmetric, so the spectral theorem gives a real eigenbasis with
eigenvalues in $[0,\mu^{-2}]$ summing to $|\Omega|/4$.
\end{proof}

\begin{proof}[Proof of Theorem~\ref{thm:G-tilt}]
Work with $\gamma=\gamma_\mu$; $T=T_{\Omega,a}$ is continuous on $C(\R^{2})$. Put
$\xi_j=\int_\Omega\gamma e_j$ for the eigenbasis $(e_j,\lambda_j)$ of $K=K_\Omega$; by Lemma~\ref{lem:G-S}(S4)
every finite vector of point values and $\xi_j$'s is almost surely in the range of
the isonormal process, hence centered jointly Gaussian, with
$\E\xi_j\xi_l=\lambda_j\delta_{jl}$ and $\E[\gamma(z)\xi_j]=\langle g_z,e_j\rangle$, $g_z=\Cend(z-\cdot)|_\Omega$. The $\xi_j$ are independent
$N(0,\lambda_j)$ and, by Parseval for every outcome, $T(\gamma)=(\kappa_a^{2}/2)\sum_j\xi_j^{2}$.
(a) With $T_N$ the partial sums, $e^{T_N}\uparrow e^{T(\gamma)}$ and
$\E e^{T_N}=\prod_{j\le N}(1-\kappa_a^{2}\lambda_j)^{-1/2}$, so $\Det_{\Omega,a}=\prod_j(1-\kappa_a^{2}\lambda_j)^{-1/2}\ge1$; with
$-\log(1-x)\le x/(1-r_a)$ on $[0,r_a]$, $1-r_a=a^{2}m^{2}$ and $\sum\lambda_j=|\Omega|/4$,
$\log\Det\le\kappa_a^{2}|\Omega|/(8a^{2}m^{2})$; monotonicity in $\Omega$ is pointwise.
(b) $0\le c_k(z,z')\le\kappa_a^{2k}(\Cend)^{*(k+1)}(z-z')=s_k(z-z')$; symmetry by relabelling; $\Cend$ is
Lipschitz with constant $\mu/4$ in $|\cdot|_1$, so $c_k$ is Lipschitz; the operator
$S_k=\Cend(\kappa_a^{2}P\Cend)^{k}$, $P=\ind_\Omega$, has kernel $c_k$ (Fubini, using
$\int c_k(x,z')|f(z')|dz'\le\|f\|_2r_a^{k}/(2\mu)$); $\|S_k\|\le\mu^{-2}r_a^{k}$, so
$\sum S_k=\Cend(I-\kappa_a^{2}P\Cend)^{-1}$ in norm, and its kernel is $C'=\sum c_k$ (uniform limit).
Since $P\Cend=R^{*}R\Cend$ and $(\kappa_a^{2}R^{*}R\Cend)^{k+1}=\kappa_a^{2}R^{*}(\kappa_a^{2}K)^{k}R\Cend$, the restriction form
follows.
(c) Fix points $z_1,\dots,z_n$ anywhere in the plane, $Y=(\gamma(z_l))_l$ and
$b_j=\E[Y\xi_j]$. Regression on the $\xi_j$ with $\lambda_j>0$, $j\le N$, and the Gaussian
tilting fact show that under $e^{T_N}\PP/\E e^{T_N}$ the vector $Y$ is centered Gaussian with
covariance $\mathrm{Cov}\,Y+\kappa_a^{2}\sum_{j\le N}b_jb_j^{T}/(1-\kappa_a^{2}\lambda_j)$, and
$\E[e^{it\cdot Y}e^{T_N}]=\E[e^{T_N}]\exp(-\tfrac12t^{T}\Sigma_Nt)$. Dominated convergence (majorant
$e^{T(\gamma)}$) and absolute convergence of the series give, in the limit, the covariance
$\Cend(z_l-z_{l'})+\kappa_a^{2}\langle g_{z_l},(I-\kappa_a^{2}K)^{-1}g_{z_{l'}}\rangle$; expanding the resolvent and using
the iterated kernels, this is $C'(z_l,z_{l'})$. Hence under $G'_{\Omega,a}$ the coordinates at
all points are centered Gaussian with covariance $C'$; uniqueness is the
path-space fact; the change-of-measure identity is the definition of $G'$,
extended to integrable $F$ by splitting into positive and negative parts.
\end{proof}

\begin{proof}[Proof of Theorem~\ref{thm:G-order}]
Pointwise: $c_0=\Cend>0$, $c_k\ge0$ and $c_k\le s_k$ give $\Cend\le C'\le\sum s_k=C_a\le c_a$; each $c_k$ is
nondecreasing in $\Omega$. As operators: $C'-\Cend=\kappa_a^{2}\Cend R^{*}(I-\kappa_a^{2}K_\Omega)^{-1}R\Cend\ge0$ because the
resolvent has spectrum in $[1,(a^{2}m^{2})^{-1}]$; $C_a-C'=\Delta\ge0$ by
Theorem~\ref{thm:G-diff}(b), whose proof uses only Theorem~\ref{thm:G-tilt}(b), the
pointwise part and Theorem~\ref{thm:G-shift}.
\end{proof}

\begin{proof}[Proof of Theorem~\ref{thm:G-diff}]
(a) For $x=(x_1,\dots,x_n)$, $1-\prod_i\ind_\Omega(x_i)=\sum_i[\prod_{l<i}\ind_\Omega(x_l)]\ind_{\Omega^{c}}(x_i)$ (first exit), so
$s_n(z-z')-c_n(z,z')=\sum_{i=1}^{n}\kappa_a^{2}\int_{\Omega^{c}}c_{i-1}(z,x)s_{n-i}(x-z')\,dx$ by Tonelli; summing,
$\Delta(z,z')=\kappa_a^{2}\int_{\Omega^{c}}C'(z,x)C_a(x-z')\,dx$, which is finite and symmetric, so also
$=\kappa_a^{2}\int_{\Omega^{c}}C_a(z-x)C'(x,z')\,dx$; substituting $C_a=C'+\Delta$ gives the third form. As
operators, with $A=I-\kappa_a^{2}\Cend$ and $B=I-\kappa_a^{2}P\Cend$, $C_a=\Cend A^{-1}$, $C'=\Cend B^{-1}$ and
$B-A=\kappa_a^{2}Q\Cend$, so $C_a-C'=\Cend A^{-1}(B-A)B^{-1}=\kappa_a^{2}C_aQC'$; self-adjointness gives the other
forms. (b) Continuity is that of $C_a$ and $C'$. Operator positivity:
$\langle f,\Delta f\rangle=\kappa_a^{2}\|QC'f\|^{2}+\kappa_a^{4}\langle QC'f,C_aQC'f\rangle\ge0$. Kernel positivity: with
$F(x)=\ind_{\Omega^{c}}(x)\sum_lc_lC'(z_l,x)\in L^{1}\cap L^{\infty}$,
$\sum c_lc_{l'}\Delta(z_l,z_{l'})=\kappa_a^{2}\|F\|_2^{2}+\kappa_a^{4}\langle F,C_aF\rangle\ge0$. (c) From the last-exit form and
$C'\le C_a$; if $\Omega$ contains the open $r$-neighbourhood of $K$, then $|z-x|\ge r$ for $z\in K$,
$x\in\Omega^{c}$, and Cauchy--Schwarz gives $\epsilon_a(r)$, which decreases to $0$ by dominated
convergence since $C_a\in L^{2}$. (d) Choose $r$ with $\epsilon_a(r)\le\eta$; every bounded $\Omega$
containing the $r$-neighbourhood of $K$ has $\sup_{K\times K}\Delta\le\eta$; for $[-R,R]^{2}$ this holds once
$R\ge r+\sup_K|z|_\infty$. (e) With $q_k=k(z+h,z+h)+k(z,z)-2k(z+h,z)$, $q_{C_a}=q_{C'}+q_\Delta$ with both
terms nonnegative, so each is at most $2(c_a-C_a(h))\le M'_{a,1/2}|h|^{1/2}$; Kolmogorov
extension gives Gaussian processes with covariances $\Delta$ and $C'$, and the continuity
criterion with $\E|Z(z)-Z(z')|^{12}\le11!!(M'_{a,1/2})^{6}|z-z'|^{3}$ and the dyadic argument give
jointly measurable continuous modifications.
\end{proof}

\begin{proof}[Proof of Theorems~\ref{thm:G-common} and \ref{thm:G-bounds}]
$\mu_0$ is even and $\sigma$-finite, so $H_\C$ is a separable real Hilbert space carrying an
isonormal $W$ (\ref{prem:C6}(f)). For real $f\in L^{2}$, $\hat f\tau_a$ is conjugate-symmetric with
$\int|\hat f\tau_a|^{2}d\mu_0=(2\pi)^{-2}\int|\hat f|^{2}\hat C_a\le m^{-2}\|f\|_2^{2}$; all $\phi_a(f)$ lie in the range of
the single $W$, hence are jointly Gaussian; the covariance is the isometry, with
$\tau_a\tau_b\hat C_0=(\hat C_a\hat C_b)^{1/2}$. For Theorem~\ref{thm:G-bounds}: with $D=m^{2}+|p|^{2}$,
$0\le\hat C_0-\hat C_a=a^{2}w^{2}q^{2}/(D(D+a^{2}w^{2}q^{2}))\le a^{2}w^{2}q^{2}/D^{2}\le a^{2}/4$ (as $4w^{2}q^{2}\le D^{2}$), and
$(1-\tau)^{2}\le1-\tau^{2}$ on $[0,1]$, so
$\|\phi_a(f)-\phi(f)\|_2^{2}=(2\pi)^{-2}\int|\hat f|^{2}(1-\tau_a)^{2}\hat C_0\le(2\pi)^{-2}\int|\hat f|^{2}(\hat C_0-\hat C_a)\le\tfrac{a^{2}}4\|f\|_2^{2}$; complex
tests split into real and imaginary parts.
\end{proof}

\begin{proof}[Proof of Theorem~\ref{thm:G-point}]
Lemma~\ref{lem:G-S} with $\nu=\mu_0$ and $\sigma=\tau_a$, $\int\tau_a^{2}d\mu_0=c_a$ and $k=C_a$: (S1) gives the
covariances, with $C_{ab}=\cI[(\hat C_a\hat C_b)^{1/2}]$ where
$\cI g(x)=(2\pi)^{-2}\int e^{ip\cdot x}g(p)\,dp$ for $g\in L^{1}$ (here
$(\hat C_a\hat C_b)^{1/2}\le\tfrac12(\hat C_a+\hat C_b)\in L^{1}$); (S2) with Theorem~\ref{thm:G-shift}(e) gives the
increment; (S3) with Theorem~\ref{thm:G-shift}(f), $\alpha=1/2$, gives $\tilde\phi_a$; (S4) gives the
smearing identity, the $L^{1}$ and Plancherel transforms agreeing on $L^{1}\cap L^{2}$. Each
modification and each smearing identity has its own null set; countable unions
of null sets are null; $\tilde\phi_a$ is constructed separately for each $a$, so no common
null set over uncountably many regulators arises, and no measurability in $a$ is
asserted.
\end{proof}

\begin{proof}[Proof of Theorem~\ref{thm:G-versions}]
$\sigma=\tau_a\rho$ has $\int\sigma^{2}d\mu_0\le(2\pi)^{-2}\int(m^{2}+|p|^{2})^{-1}(1+|p|^{2})^{-2}dp<\infty$, and
$\E|\psi_a(z+h)-\psi_a(z)|^{2}\le|h|^{2}(2\pi)^{-2}\int|p|^{2}\hat C_0(1+|p|^{2})^{-2}dp\le|h|^{2}/(4\pi)$, so
Lemma~\ref{lem:G-S}(S3) with $\vartheta=2$ gives $\tilde\psi_a$, a Borel $C(\R^{2})$-valued random element,
and $\Phi_a=L\tilde\psi_a$ is Borel because $L$ is continuous
(Proposition~\ref{prop:space}(viii)). (a) For real $f\in C_c^{\infty}$ and $g=(1-\Delta)f$,
$\hat g=(1+|p|^{2})\hat f$ and, for every outcome, $\Phi_a(f)=\int\tilde\psi_ag$; by (S4) this is
$W(\hat g\tau_a\rho)=W(\hat f\tau_a)=\phi_a(f)$ almost surely. (b) For $a>0$, $\Phi_a$ and $\iota\tilde\phi_a$ agree on
the countable family $\cD_0$ of Proposition~\ref{prop:space}(ix) off one null set
(by (a) and Theorem~\ref{thm:G-point}), hence everywhere off it. (c) Same argument
with $\Phi'$. (d) Countable unions of null sets. (e) By (a) for the countably many
tests $\chi_j\conj{e_{j,k}}$, Tonelli and Theorem~\ref{thm:G-bounds},
$\E p_j(\Phi_a)^{2}=\sum_k\lambda_{j,k}^{-2}\E|T_a(\chi_j\conj{e_{j,k}})|^{2}\le m^{-2}B_j$, and the difference bound is
the same with $\tfrac{a^{2}}4$; $B_j<\infty$ by the shell count. Since
$\|\min(1,Y)\|_2\le\min(1,\|Y\|_2)$ for $Y\ge0$, Minkowski for the series gives
$\|d(\Phi_a,\Phi_0)\|_{L^{2}}\le h(a)$, and $h(a)\to0$ by dominated convergence. (f) is (a) with
Theorem~\ref{thm:G-common}.
\end{proof}

\begin{remark}[Scope of the Gaussian proofs]\label{rem:G-proofscope}
These proofs use no interaction functional, Wick power, weight, normalizer or
statement of Theorem~\ref{thm:W}; in particular the construction of $\Phi_0$ does not
use the existence of $U_0$. They establish exactly the Gaussian interface used by
Theorem~\ref{thm:W}(j) and Corollary~\ref{cor:W-laws}: the common isonormal space, the
continuous modifications for regulators, and the Borel versions $\Phi_a$ for
$0\le a<1/m$ with their compact-test identities. That the interacting laws exist
on this space is a consequence of Theorem~\ref{thm:W}, not of this appendix.
Nothing is asserted at the endpoint value $a=1/m$.
\end{remark}

% ---------------------------------------------------------------------------
\section{Proofs for the fixed-volume Wick theory}\label{app:W}

Throughout, $\Lambda$ is a physical square with $A=|\Lambda|$, $0<a<1/m$ unless $a=0$ is stated,
$\|\cdot\|_r$ is the norm of $L^{r}(\PP_0)$ or of $L^{r}(\R^{2},dz)$, $D=m^{2}+|p|^{2}$,
$\sigma_a=\hat C_0-\hat C_a$, $\hat C_{ab}=(\hat C_a\hat C_b)^{1/2}$ and $C_{ab}=\cI\hat C_{ab}$. The standard inputs are
\ref{prem:C6}(a)--(c), Gaussian characteristic functions
$\E e^{iX}=e^{-\Var X/2}$ and $\E\mathrm{He}_4(N(0,1))=0$, Fourier series on tori, and Schur's product
theorem.

\subsection{Gaussian preliminaries}

\begin{lemma}[Wick powers, reflection and the free $\sigma$-algebra]\label{lem:W-prelim}
(i) For $\xi$ in the first chaos $H_1$ with $c=\Var\xi\ge0$, ${:}\xi^{4}{:}={:}\xi^{4}{:}_c\in H^{{:}4{:}}$ (for $c=0$,
$\xi=0$ almost surely and ${:}\xi^{4}{:}=0$), $\E[{:}\xi^{4}{:}{:}\eta^{4}{:}]=4!(\E\xi\eta)^{4}$ for all $\xi,\eta\in H_1$, and
$z\mapsto{:}\phi_a(z)^{4}{:}$ is continuous into $H^{{:}4{:}}$. For $\xi\sim N(0,c)$ and $c'\ge0$,
$\E[\xi^{4}-6c'\xi^{2}+3c'^{2}]=3(c-c')^{2}$, which vanishes exactly at the own variance.
(ii) Let $T$ be a surjective real-linear isometry of $H_1$ and $\cF_H=\sigma(H_1)\vee\cN$. There is a
unique map $\Theta_T$ on $L^{0}(\cF_H)$ with $\Theta_T[g((\xi_j)_j)]=g((T\xi_j)_j)$ for countable $(\xi_j)\subset H_1$ and
Borel $g$; it is linear, multiplicative, commutes with Borel functions (in
particular with conjugation and $e^{-\cdot}$), preserves laws, is an isometry of every
$L^{r}$ and commutes with $L^{r}$-limits and Bochner integrals in $L^{2}$. For
$\theta\in\{\theta_1,\theta_2\}$, $(T_\theta h)(p)=h(\theta p)$ is orthogonal on $H_\C$, and with $T=WT_\theta W^{-1}$,
$\Theta=\Theta_T$ satisfies $\Theta\phi_a(f)=\phi_a(f\circ\theta)$, $\Theta\phi_a(z)=\phi_a(\theta z)$ and
$\Theta{:}\phi_a(z)^{4}{:}={:}\phi_a(\theta z)^{4}{:}$.
(iii) $\{\phi(f):f\in C_c^{\infty}(\R^{2};\R)\}$ is dense in $H_1$, and $\cF_H=\cF_{\mathrm{free}}$.
\end{lemma}

\begin{proof}
(i) \ref{prem:C6}(b) for positive variances; if a variance vanishes both sides vanish
since $|\E\xi\eta|\le(\Var\xi\Var\eta)^{1/2}$. Hence $\E|{:}\xi^{4}{:}-{:}\eta^{4}{:}|^{2}=4!(c_\xi^{4}-2(\E\xi\eta)^{4}+c_\eta^{4})$, and
$z\mapsto e_z\tau_a$ is continuous into $H_\C$. (ii) $H_1$ is separable; every $Y\in L^{0}(\cF_H)$ is
$g((\xi_j)_j)$ by Doob--Dynkin; if $g((\xi_j))=g'((\xi'_j))$ almost surely, the vectors
$((\xi_j),(\xi'_j))$ and $((T\xi_j),(T\xi'_j))$ are centered Gaussian with the same covariances,
hence have the same law, so the Borel set $\{g\ne g'\}$ is null for both. The listed
properties follow by merging countable families. $T_\theta$ is orthogonal because
Lebesgue measure and $\hat C_0$ are $\theta$-invariant; $(f\circ\theta)\hat{}=\hat f\circ\theta$, $\tau_a\circ\theta=\tau_a$ and
$T_\theta(e_z\tau_a)=e_{\theta z}\tau_a$. (iii) For $h\in H_\C$ and $R>0$, $\hat C_0\ge(m^{2}+R^{2})^{-1}$ on $\{|p|\le R\}$, so
$h_R=h\ind_{\{|p|\le R\}}\in L^{2}(dp)$ is Hermitian, $f_R=\cI h_R$ is real with $\phi(f_R)=W(h_R)$, and
$\|h-h_R\|_{H_\C}\to0$; $\|\phi(f)-\phi(f')\|_2\le m^{-1}\|f-f'\|_2$ and density of $C_c^{\infty}$ in $L^{2}$ conclude.
Each element of $H_1$ is an almost sure limit of free-field smearings, and $\cN\subset\cF_{\mathrm{free}}$.
\end{proof}

\begin{lemma}[Kernel facts]\label{lem:W-kernels}
(K1) For $0\le b\le a<1/m$: $0<\hat C_a\le\hat C_{ab}\le\hat C_b\le\hat C_0$, and $\hat C_b-\hat C_{ab}$,
$\hat C_{ab}-\hat C_a$, $\hat C_b-\hat C_a$ lie in $[0,\sigma_a]$. (K2) $\sigma_a=a^{2}w^{2}q^{2}/(D(D+a^{2}w^{2}q^{2}))\in[0,\min(a^{2}/4,\hat C_0)]$.
(K3) $\cI$ satisfies $\|\cI g\|_2=(2\pi)^{-1}\|g\|_{L^{2}(dp)}$ and $\|\cI g\|_6\le(2\pi)^{-5/3}\|g\|_{6/5}$; for bounded
Hermitian $k$, $K[k]=\cI(k\hat C_0)$ is real, $k\mapsto K[k]$ is real-linear,
$\|K[k]\|_2=(2\pi)^{-1}\|k\hat C_0\|_{L^{2}(dp)}$ and $\|K[k]\|_6\le\|k\|_\infty M_6$; in particular $\|C_{ab}\|_6\le M_6$ for
$a,b>0$, and each of $C_b-C_{ab}$, $C_{ab}-C_a$, $C_b-C_a$ has $L^{2}$ norm at most $\kappa_2a$ for
$0<b\le a$. (K4) $c_a$ is strictly decreasing, $c_a\le(2\pi)^{-1}\log(\kappa_0/(am))$, and $c_a>c_{1/m}=1/4$.
(K5) $C_a(t,y)=(2\pi)^{-1}\int e^{iqy}e^{-E_a(q)|t|}(2E_a(q)(1+a^{2}q^{2}))^{-1}dq$ with
$E_a(q)=[(m^{2}+q^{2})/(1+a^{2}q^{2})]^{1/2}$, and $\iint h(\theta z)C_a(z-z')h(z')\,dz\,dz'\ge0$ for real
$h\in C_c^{\infty}(H^{\theta}_{+})$. (K6) $C_{ab}$ is bounded and continuous for $a,b>0$,
$\sup|C_a|=c_a$ and $C_a(\theta x)=C_a(-x)=C_a(x)$.
\end{lemma}

\begin{proof}
(K1) With $x=\hat C_a\le y=\hat C_b$: $\sqrt{xy}\in[x,y]$, $y-\sqrt{xy}\le y-x$, $\sqrt{xy}-x\le y-x\le\sigma_a$.
(K2) Expand $1/D-1/(D+a^{2}w^{2}q^{2})$ and use $w^{2}q^{2}\le D^{2}/4$. (K3) For $s>1$,
$\int(m^{2}+|p|^{2})^{-s}dp=\pi m^{2-2s}/(s-1)$, so $\hat C_0\in L^{2}\cap L^{6/5}$ with
$\|\hat C_0\|_{6/5}=(5\pi m^{-2/5})^{5/6}$; Hausdorff--Young for the unnormalized transform gives the
factor $(2\pi)^{-2}(2\pi)^{1/3}$; the Plancherel and Hausdorff--Young extensions agree almost
everywhere (truncations converge in both spaces); Hermitian symmetry gives
reality. For the differences, $\|\sigma_a\|^{2}_{L^{2}(dp)}\le\int_{|p|\le1/a}(a^{2}/4)^{2}+\int_{|p|>1/a}|p|^{-4}=17\pi a^{2}/16$.
(K4) $\int_\R dw/(\alpha+\beta w^{2})=\pi/\sqrt{\alpha\beta}$; the integrand of \eqref{eq:ca} is strictly decreasing in
$a$ for $q\ne0$; $\int_0^{1/a}dq/\sqrt{m^{2}+q^{2}}=\mathrm{asinh}(1/(am))\le\log((1+\sqrt2)/(am))$ and
$\int_{1/a}^{\infty}dq/(aq^{2})=1$; at $a=1/m$ the integral is $(2\pi)^{-1}\pi/2$ (used only for this
comparison; $1/m$ is not a regulator). (K5) $\hat C_a=[(1+a^{2}q^{2})(w^{2}+E_a^{2})]^{-1}$ and
$(2\pi)^{-1}\int e^{iwt}(w^{2}+E^{2})^{-1}dw=e^{-E|t|}/(2E)$; for $\theta=\theta_1$ and $t,t'>0$ on $\supp h$, Fubini
gives $(2\pi)^{-1}\int dq\,[2E_a(1+a^{2}q^{2})]^{-1}|\int h(t,y)e^{-E_a(q)t}e^{iqy}dt\,dy|^{2}\ge0$; $\hat C_a$ is symmetric in
$(w,q)$, which gives $\theta_2$. (K6) $\hat C_{ab}\ge0$ is integrable and invariant under $p\mapsto\theta p$ and
$p\mapsto-p$.
\end{proof}

\subsection{The interaction and its \texorpdfstring{$L^{2}$}{L2} limit}

\begin{proof}[Proof of Theorem~\ref{thm:W}(a), (c), (d) and the existence of $U_0$]
\emph{Pathwise equals Bochner.} $z\mapsto{:}\tilde\phi_a(z,\omega)^{4}{:}$ is continuous for every outcome, so
$U_a$ is a finite pathwise integral, measurable by Fubini. $Y(z)=u{:}\phi_a(z)^{4}{:}$ is
continuous into $H^{{:}4{:}}$ with $\E Y(z)^{2}=4!u^{2}c_a^{4}$, so the Bochner integral $B=\int_\Lambda Y$ exists in
$H^{{:}4{:}}$; for bounded $V$, $\E[VU_a]=\int_\Lambda\E[VY(z)]dz=\E[VB]$, so $U_a=B$ almost surely.
\emph{Mean zero and lower bound.} $\E U_a=0$ because the Wick constant is the own
variance; $x^{4}-6cx^{2}+3c^{2}=(x^{2}-3c)^{2}-6c^{2}$ gives $U_a\ge-6uAc_a^{2}$ for every outcome, and the
same for $U^{\pm}_a=u\int_{\Lambda\cap H_\pm}{:}\tilde\phi_a^{4}{:}$. \emph{$L^{2}$ estimate.} For $0<b\le a$, Fubini for Bochner
integrals and Lemma~\ref{lem:W-prelim}(i) give
$\E|U_a-U_b|^{2}=4!u^{2}\iint_{\Lambda^{2}}[C_a^{4}-2C_{ab}^{4}+C_b^{4}](z-z')\le4!u^{2}A(\|C_a^{4}-C_{ab}^{4}\|_1+\|C_b^{4}-C_{ab}^{4}\|_1)$;
$|X^{4}-Y^{4}|\le4|X-Y|(|X|^{3}+|Y|^{3})$, Cauchy--Schwarz and (K3) give
$\|C_a^{4}-C_{ab}^{4}\|_1\le8\kappa_2M_6^{3}a$, so $\E|U_a-U_b|^{2}\le384u^{2}A\kappa_2M_6^{3}a=K_\Lambda^{2}a$.
\emph{Limit.} Along any sequence of regulators tending to $0$ the $U_{a_n}$ are Cauchy;
interlacing shows that all limits coincide; letting $b\downarrow0$ gives $\|U_a-U_0\|_2\le K_\Lambda\sqrt a$
for every regulator; $H^{{:}4{:}}$ is closed, and $\E U_0=\lim\E U_a=0$. A Wick constant $c'\ne c_a$ would
give the nonzero mean $3uA(c_a-c')^{2}$.
\end{proof}

\begin{proof}[Proof of Theorem~\ref{thm:W}(b)]
\emph{Continuity of $Q$.} For bounded Hermitian $k$ put $Q(k)=\iint_{\Lambda^{2}}K[k](z-z')^{4}$. If
$k_n\to k$ almost everywhere with $\sup\|k_n\|_\infty\le\kappa$, then
$|Q(k_n)-Q(k)|\le(4A/\pi)\kappa^{3}M_6^{3}\|(k_n-k)\hat C_0\|_{L^{2}(dp)}\to0$ by dominated convergence with
majorant $4\kappa^{2}\hat C_0^{2}$. \emph{All admissible cutoffs.} Since $\int\hat C_0=\infty$, Fatou gives
$c^{\chi}_\epsilon=\Var\phi^{\chi}_\epsilon(z)\to\infty$, so $c^{\chi}_\epsilon>0$ for small $\epsilon$ (if $c^{\chi}_\epsilon=0$, then $\chi_\epsilon=0$ almost
everywhere and $N^{\chi}_\epsilon=0$). For admissible $\chi,\psi$,
$\E[\phi^{\chi}_\epsilon(z)\phi^{\psi}_\delta(z')]=K[\conj{\chi_\epsilon}\psi_\delta](z-z')$; the regulator family $\tau$ is admissible with
$N^{\tau}_a=U_a$; so $\E|N^{\chi}_\epsilon-U_a|^{2}=4!u^{2}[Q(|\chi_\epsilon|^{2})-2Q(\conj{\chi_\epsilon}\tau_a)+Q(\tau_a^{2})]$. Let $a\downarrow0$ at
fixed $\epsilon$ (existence of $U_0$ and continuity of $Q$), then $\epsilon\downarrow0$ along any sequence:
the limit is $4!u^{2}[Q(1)-2Q(1)+Q(1)]=0$. For a mollifier,
$\chi_\epsilon(p)=\conj{\hat\eta(\epsilon p)}$ gives $\phi^{\chi}_\epsilon(z)=\phi(\eta_\epsilon(z-\cdot))$, with $\sup|\chi_\epsilon|\le\|\eta\|_1$.
\emph{Pairings and uniqueness.} For $a>0$,
$\E[U_a{:}\phi(f_1)\cdots\phi(f_4){:}]=4!u\int_\Lambda\prod_j(C_{a0}*f_j)(z)\,dz$; since
$|1-\tau_a|\hat C_0\le\sigma_a\le a^{2}/4$ and $\hat f_j\in\cS$, $C_{a0}*f_j\to\Cfree*f_j$ uniformly, and the formula
passes to $U_0$. If $Y\in H^{{:}4{:}}$ has the same pairings, $Y-U_0$ is orthogonal to the Wick
monomials of the dense set of Lemma~\ref{lem:W-prelim}(iii), which span a dense
subspace of $H^{{:}4{:}}$ (\ref{prem:C6}(b)); so $Y=U_0$. \emph{Measurability and scope.}
$U_0\in H^{{:}4{:}}\subset L^{2}(\cF_H)=L^{2}(\cF_{\mathrm{free}})$. The construction uses only the variables
$W(h)$, their own-variance Wick polynomials, Bochner integrals and $L^{2}$-limits; it
never uses a point variable at $a=0$ (indeed $e_z\notin L^{2}(\mu_0)$ because $\int\hat C_0=\infty$) and
forms no pointwise product of the distribution-valued free field; nothing is
asserted about other representations of $U_0$.
\end{proof}

\subsection{Nelson tail, weights and normalizers}

\begin{proof}[Proof of the tail bound and of Theorem~\ref{thm:W}(e)]
By \ref{prem:C6}(a) with $k=4$, $\|X\|_r\le(r-1)^{2}\|X\|_2$ for $X\in H^{{:}4{:}}$ and real $r\ge2$. Fix
$K\ge K_\sharp$ and put $s=(K/(12uA))^{1/2}\ge1$ and $a'=(\kappa_0/m)e^{-2\pi s}$.
(D1) $\kappa_0=e(1+\sqrt2)<e^{2}<e^{2\pi}$, so $a'm<1$ and $a'$ is a regulator.
(D2) $c_{a'}\le s$ by (K4), so $U_{a'}\ge-6uAs^{2}=-K/2$ for every outcome.
(D3) For $a'\le a<1/m$, $c_a\le c_{a'}$, so $U_a\ge-K/2$ and $\PP_0(U_a\le-K)=0$.
(D4) For $0<a<a'$, $\{U_a\le-K\}\subseteq\{|U_{a'}-U_a|\ge K/2\}$.
(D5) $X=U_{a'}-U_a\in H^{{:}4{:}}$ has $\|X\|_2\le K_\Lambda a'^{1/2}$, so for real $r\ge2$,
$\PP_0(|X|\ge K/2)\le[(2K_\Lambda/K)(\kappa_0/m)^{1/2}(r-1)^{2}e^{-\pi s}]^{r}$.
(D6) With $r=1+e^{\pi s/2-1}\ge2$, $(r-1)^{2}=e^{\pi s-2}$ and the bracket is
$(2K_\Lambda/K)(\kappa_0/m)^{1/2}e^{-2}\le e^{-1}$ because $K\ge(2/e)K_\Lambda(\kappa_0/m)^{1/2}$.
(D7) $\PP_0(U_a\le-K)\le e^{-r}\le\exp(-\exp((\pi/2)(K/(12uA))^{1/2}-1))=T(K)$.
The coarse regulator depends on $K$ only, so the bound is uniform in $a$. By the
layer-cake formula, $\E_0e^{-pU_a}=\int p\,e^{pK}\PP_0(U_a<-K)\,dK\le\Xi_p$. $\Xi_p<\infty$: with
$c=(\pi/2)(12uA)^{-1/2}$ and $x=c\sqrt K-1$, for $K\ge4/c^{2}$ one has $e^{x}\ge x^{4}/24\ge c^{4}K^{2}/384$, so the
integrand is at most $pe^{-pK}$ for $K\ge\max(K_\sharp,4/c^{2},768p/c^{4})$. At $a=0$ choose $a_j\downarrow0$ with
$U_{a_j}\to U_0$ almost surely; Fatou gives $\E_0e^{-pU_0}\le\Xi_p$ and
$\PP_0(U_0<-K)\le\liminf\PP_0(U_{a_j}<-K)\le T(K)$ (open event only).
\end{proof}

\begin{proof}[Proof of Theorem~\ref{thm:W}(f), (g)]
$|e^{-x}-e^{-y}|\le|x-y|(e^{-x}+e^{-y})$ and H\"older with exponents $(2p,2p)$ give
$\|e^{-U_a}-e^{-U_0}\|_p\le\|U_a-U_0\|_{2p}\|e^{-U_a}+e^{-U_0}\|_{2p}\le(2p-1)^{2}K_\Lambda\sqrt a\cdot2\Xi_{2p}^{1/(2p)}$. Jensen and
$\E U_a=0$ give $Z_a\ge1$, and (e) with $p=1$ gives $Z_a\le\Xi_1$; the lower bound concerns
the physical laws on the common Gaussian space only.
\end{proof}

\subsection{Central reflection positivity at fixed volume}

\begin{proof}[Proof of Theorem~\ref{thm:W}(h), (i)]
\emph{Half-plane localization.} For $z\in H_+$, the mollified tests $\rho_n(\cdot-z)$ lie in
$C_c^{\infty}(H_+)$ for large $n$ and $\phi_a(\rho_n(\cdot-z))\to\phi_a(z)$ in $L^{2}$, so ${:}\phi_a(z)^{4}{:}\in L^{2}(\cA^{\theta,a}_{+})$
and the Bochner integral $U^{+}_a$ lies in $L^{2}(\cA^{\theta,a}_{+})$; $U_a=U^{+}_a+U^{-}_a$ (the line is null),
and $\Theta U^{+}_a=U^{-}_a$ because $\Theta$ commutes with Bochner integrals and $\Lambda$ is
$\theta$-invariant. \emph{Exponentials are dense.} The span of $e^{i\phi_a(f)}$, $f\in C_c^{\infty}(H_+;\R)$, is
dense in $L^{r}(\cA^{\theta,a}_{+})$, $1\le r<\infty$: an annihilating $g\in L^{r'}$ gives finite complex
measures with vanishing Fourier transforms on every finite cylinder, hence
$\E[g\ind_S]=0$ on a generating $\pi$-system and, by Dynkin, on $\cA^{\theta,a}_{+}$.
\emph{Gaussian positivity.} For $G=\sum_kc_ke^{i\phi_a(f_k)}$,
$\E[\conj{\Theta G}G]=\sum_{j,k}\conj{d_j}d_k\exp(M_{jk})$ with $d_k=c_ke^{-\langle f_k,C_af_k\rangle/2}$ and
$M_{jk}=\langle f_j\circ\theta,C_af_k\rangle$; $M$ is real symmetric and positive semidefinite by (K5), (K6), so
$\exp\circ M=\sum_nM^{\circ n}/n!$ is positive semidefinite by Schur's theorem; density extends
this to $G\in L^{2}(\cA^{\theta,a}_{+})$. \emph{(h)} With $e^{-U^{\pm}_a}\le e^{6uAc_a^{2}}$, $G=Fe^{-U^{+}_a}\in L^{2}(\cA^{\theta,a}_{+})$
and $\Theta G=(\Theta F)e^{-U^{-}_a}$, so $\E[\conj{\Theta F}Fe^{-U_a}]=\E[\conj{\Theta G}G]\ge0$. \emph{(i)} For
$G=\sum_jc_je^{i\phi(f_j)}$ and $G_a=\sum_jc_je^{i\phi_a(f_j)}$, (h) gives $\E[\conj{\Theta G_a}G_ae^{-U_a}]\ge0$, and with
$c=\sum|c_j|$, $|e^{ix}-e^{iy}|\le|x-y|$, Theorem~\ref{thm:G-bounds}, (e) with $p=2$ and (f) with $p=1$,
the difference from $\E[\conj{\Theta G}Ge^{-U_0}]$ is at most $c\,\Xi_2^{1/2}(a\sum_j|c_j|\|f_j\|_2+2cK_\Lambda\sqrt a)\to0$. For
bounded $\cA^{\theta,0}_{+}$-measurable $F$, approximate in $L^{4}$ by such $G_n$ and use H\"older with
exponents $(4,4,2)$ and $\|e^{-U_0}\|_2\le\Xi_2^{1/2}$.
\end{proof}

\begin{proof}[Proof of Theorem~\ref{thm:W}(j) and of the identification with Definition~\ref{def:laws}(d)]
For Borel $B$, $\ind_B(\Phi')$ is measurable and $Z_a\in[1,\Xi_1]$, so the formula defines a Borel
probability with $\int H\,d\nu=Z_a^{-1}\E_0[H(\Phi')e^{-U_a}]$ for bounded Borel $H$. For bounded
$\Sigma^{\theta}_{+}$-measurable $F$, Proposition~\ref{prop:space}(x) gives real
$f_j\in C_c^{\infty}(H^{\theta}_{+})$ and a bounded Borel $g$ with $F(v)=g((v(f_j))_j)$ and
$F(\theta v)=g((v(f_j\circ\theta))_j)$. Off one null set, $\Phi'(f)=\phi_a(f)$ for the countably many $f_j$
and $f_j\circ\theta$; there $F(\Phi')=G:=g((\phi_a(f_j))_j)$, which is $\cA^{\theta,a}_{+}$-measurable, and
$F(\theta\Phi')=g((\phi_a(f_j\circ\theta))_j)=\Theta G$ almost surely by Lemma~\ref{lem:W-prelim}(ii). So
$\int\conj{F(\theta v)}F(v)\,\nu(dv)=Z_a^{-1}\E_0[\conj{\Theta G}Ge^{-U_a}]\ge0$ by (h) for $a>0$ and by (i) for $a=0$.
On the common space, Theorem~\ref{thm:G-versions}(a) supplies the hypothesis on
$\Phi'=\Phi_a$ for every $0\le a<1/m$, including $a=0$; two continuous modifications of
$\phi_a$ agree everywhere off one null set (rational points and continuity), so the
pathwise $U_a$ does not depend on the modification; and any other Borel $\Phi'$ with
the same compact-test identities equals $\Phi_a$ almost surely
(Theorem~\ref{thm:G-versions}(c)), hence gives the same law. This is
Corollary~\ref{cor:W-laws}.
\end{proof}

\begin{remark}[Proof of Theorem~\ref{thm:W}(k)]
No step above uses a restriction on $u/m^{2}$, and every constant is computed for
the fixed square $\Lambda$.
\end{remark}

\begin{remark}[Volume dependence and unused routes]\label{rem:W-proofscope}
$K_\Lambda=u(384\kappa_2M_6^{3})^{1/2}A^{1/2}$, $K_\sharp\ge12uA$ and $\Xi_p\ge e^{12puA}$: every quantitative bound is
for a fixed physical square, and the only volume-uniform statement is $Z\ge1$.
A frequency-split tail argument, which gives another family of
regulator-uniform moment bounds, was checked in the antecedent record and is not
used. The proofs use the Gaussian input only through
Theorems~\ref{thm:G-common}, \ref{thm:G-point} and \ref{thm:G-versions} and
Proposition~\ref{prop:space}; the weight, the normalizer and positivity enter
only afterwards.
\end{remark}

% ---------------------------------------------------------------------------
\section{Provenance, review status and source access}\label{app:provenance}

\paragraph{Authorship of this text.} The manuscript and corrections A, AC and
AC2 were written by Claude (Opus~5.5, Anthropic), with retained earlier Codex
argument lineage as listed below, in a project directed by Cody Mitchell.
The AC3 integration author is GPT-6.1 Sol (max effort, OpenAI), adopting the
Sol-authored SB5 and SB6 proof arguments and SMH numerical repair, expanding
their complete derivations and writing the dependency/source separation.
Older Astra-line H1/K1 lineage and Opus authorship remain attributed.
The authorship map records every adopted or rewritten argument; a reviewer
is an author only when its own proof argument is adopted, never merely for
its review verdict. Historical A/AC/AC2 and adverse review bytes are preserved;
all earlier verdicts remain bound to the versions they read. The exact-version review status of this
text is recorded in the companion decision. Every review mentioned below was an
internal review by language models; none is an independent expert review, and
no record is a formal (machine-checked) proof.

\paragraph{Historical PV4 provenance.} The following publication-revision
paragraph describes the retained 23 October baseline. The new private
26 October supplement records its own argument adoptions, source premises
and required exact-version reviews in Appendix~\ref{qrs:supplement} and the
accompanying proof map. The original mathematical blocks are unchanged.

\paragraph{Publication revision.} The present revision adds explanatory prose by Claude (Opus~5.5), with title, opening, integration and precision clarifications by GPT-6 Astra under Cody Mitchell's direction. All numbered mathematical statements and proof bodies are retained verbatim from the accepted construction manuscript. The companion decision binds the internal revision reviews to the exact source and PDF.

\paragraph{Argument lineage in the construction chain.}
The fixed-volume arguments of Section~\ref{sec:fixed} retain those of two earlier
Codex-authored records (F1 and F2 in the table below): the local-$L^{4}$ compact
containment and finite-vector identification of the width record in
Proposition~\ref{prop:width}, and the residual normalization, covariance-defect
coupling, estimates and constants of the fixed-volume record in
Propositions~\ref{prop:masstilt}--\ref{prop:reg}.

Theorems~\ref{thm:LR}, \ref{thm:V1} and~\ref{thm:CT} have their own proofs; no
earlier conditional record is cited for them. Their arguments retain forms from
earlier Codex-authored records (LR, V1 and CT in the table below):
the dilation and characteristic-function route of Theorem~\ref{thm:LR}; the lattice
inequality, Paley--Zygmund bound, Fourier moments, tightness and version of
Theorem~\ref{thm:V1}; and the invariance, regularity, non-Gaussianity, kernel and
determining arguments of Theorem~\ref{thm:CT} other than part~(iii). The
Gaussian-realization transfer (Lemma~\ref{lem:transfer-gauss}) and corrected
readings lemma (Lemma~\ref{lem:readings}) were newly written in the Opus text.
The AC3 integration adopts the Sol-authored SB5/SB6 derivations: the direct
classical Wick/normalizer lemma, explicit determining/density/kernel proofs,
source-only symmetries, full positivity, Euclidean mixing, exact regularity,
and the local conditional-density non-Gaussianity and reconstruction route.
The mixing argument also retains the earlier Claude classical-input lineage.
These are manuscript argument adoptions, not attribution to the reviewers.

Sections~\ref{sec:H1} and~\ref{sec:K1} re-express, for the present text, the
moment-transfer, selection and joint-correlation arguments of two earlier project
records (H1 and K1 in the table below). In those records the fixed-volume arrows,
the Gaussian and Wick estimates and the volume limit were premises; here each
input is a theorem proved in the main text or Appendices~\ref{app:E}--\ref{app:W},
and every place where an input is used is cited. The arguments are otherwise of
the same form.

\paragraph{Antecedent records.} The following records were supplied with
their acceptance decisions; an acceptance decision covers only its literal
scope, and a lemma appearing in a file is not covered merely by being there.
{\small
\begin{longtable}{@{}p{0.2\textwidth}p{0.17\textwidth}p{0.55\textwidth}@{}}
\toprule
Record & Author family & Use in this text\\
\midrule
Endpoint foundations (E) & Claude (Opus~5.5) & Accepted as endpoint foundations under stated classical field-space premises; re-expressed in Section~\ref{sec:endpoint} and Appendix~\ref{app:E}.\\
Gaussian foundations (G) & Claude (Opus~5.5) & Accepted for its statements G1--G11, the interaction consequence of G11 only conditionally; its Borel-law lemma for $X$ is outside that scope and is replaced by Proposition~\ref{prop:space}(ii),(iii), proved here. Re-expressed in Section~\ref{sec:gauss} and Appendix~\ref{app:G}.\\
Wick theory (W) & Claude (Opus~5.5) & Accepted; re-expressed in Section~\ref{sec:wick} and Appendix~\ref{app:W}, including the precise construction statement of Theorem~\ref{thm:W}(b).\\
Classical inputs (C) & Claude (Opus~5.5) & Accepted formulations of \ref{prem:C1}--\ref{prem:C6}, used verbatim in content in Section~\ref{sec:classical}. The application question on the reading of ``polynomial function'' in \ref{prem:C2}(iii) was closed by a later reading of the AHKZ supplement in the reconstruction-coverage review (RC), not by the C author.\\
Dependency route (D) & Claude (Opus~5.5) & Accepted route ``candidate B'': remove four project premises by new direct proofs and prove the other five; implemented here (Appendix~\ref{app:dispositions}).\\
Contribution wording (N) & Claude (Opus~5.5) & Accepted wording, used in Sections~\ref{sec:intro}, \ref{sec:prior} and \ref{sec:limits}, including the derivation of Remark~\ref{rem:frank-moments}, which is N-authored; the Codex-family review RN reviewed that derivation and is not one of its authors.\\
Older records F1 (width), F2 (fixed volume), LR (continuum limit), V1 (volume), CT (target) & Codex & Conditional on the nine project premises; their premise-conditional statements are not cited for any conclusion here. Their argument forms are retained in the current text, which is therefore jointly authored (Codex argument lineage; historical Opus integration, current Sol adoption): F1's local-$L^{4}$ compact containment and finite-vector identification in Proposition~\ref{prop:width}; F2's field-space construction in Proposition~\ref{prop:space} and its residual normalization, covariance-defect coupling, estimates and constants in Propositions~\ref{prop:masstilt}--\ref{prop:reg}; LR's dilation and characteristic-function route in Lemma~\ref{lem:dilation} and Theorem~\ref{thm:LR}; V1's lattice inequality, Paley--Zygmund bound, Fourier moments, tightness and version in Theorem~\ref{thm:V1}; CT's invariance, regularity, non-Gaussianity, kernel and determining arguments in Theorem~\ref{thm:CT} other than part~(iii) and Lemmas~\ref{lem:determining}--\ref{lem:kernels}.\\
Older records H1 (moments), K1 (kernels) & Codex (requested model of the Astra line) & Reviewed by Claude (Opus~5.5) reviewers in hostile, source, scope, blind and matching roles. Their arguments are re-expressed, with rewritten proofs of the inputs, in Sections~\ref{sec:H1}--\ref{sec:K1}; their authorship is unchanged by this integration. AC2 replaced the original bounded Hilbert-sum constant by a distinct scalar $\sigma_j$ bound. SMH recorded that literal mismatch. AC3 adopts SMH's Sol-authored pointwise-scalar/spatial-Minkowski repair to restore the exact original $K_j$ inequality and its finite-width $Z_N^b$ consequence, alongside the separate $\sigma_j$ estimate. The common sequence and compact K1 arguments are unchanged.\\
\bottomrule
\end{longtable}}

\paragraph{Components and the review they require.}
The retained Opus/Astra-line components and joint claims keep their actual
authorship. The adopted SB5/SB6/SMH arguments and new AC3 integration text
add Sol authorship. Their current-version qualifying review must therefore
cover each actual author model, with fresh contexts and the model-specific
gate required by the activated project policy; an author model cannot
qualify its own contribution. The exact allocations, distinct-model fallback
status, findings, actual-proof matches and acceptance decision belong in the
companion records. Historical family labels and earlier review requirements
are preserved there; this paragraph does not promote a historical verdict
or claim independent expert review. Theorems~\ref{thm:main} and
\ref{thm:complement} require coverage of every component chain.

\paragraph{Source access.} The authors of this text (tasks A, AC, AC2 and AC3) read no
original classical source. AC3 read only the supplied frozen manuscript,
statement, companion maps, proof modules and internal review records. The page-level access below is that recorded by the
classical-input and reconstruction-coverage reviews; ``image'' means a page image
was viewed, ``text'' means only supplied text of the page was read. Access
recorded by later reviews of version A is recorded in their own records and is
not credited to the authors of this text.
\begin{center}
\small
\begin{tabular}{@{}p{0.17\textwidth}p{0.36\textwidth}p{0.38\textwidth}@{}}
\toprule
Source & Page images viewed & Supplied text only\\
\midrule
\cite{FS1977} & pp.~495--498, 503, 512--514, 517, 518 & pp.~494, 509--511\\
\cite{AHKZ1989} & pp.~684, 685 (with a detail of eq.~(1)), 686--690 & p.~691 and a supplement\\
\cite{GRS1976} & pp.~310--312, 316, 324 & pp.~317--320\\
\cite{SG1973} & p.~151 & p.~148\\
\cite{Melrose2016Lecture15} & --- & excerpts\\
\bottomrule
\end{tabular}
\end{center}
Further page images of a source of Fr\"ohlich were viewed for an alternative
route that is not used here; that source is not cited. Not available: the
bodies of \cite{OS1973,OS1975}, \cite{StreaterWightman1964}, the antecedents of the
review in \cite[pp.~316--319]{GRS1976}, and the originals behind the Markov and
lattice positivity attributions of Section~\ref{sec:prior}. The neural-network
literature of Section~\ref{sec:prior} was read by the contribution analysts,
partly in excerpts (see the notes in the bibliography); the authors of this text
did not read it.

% ---------------------------------------------------------------------------
\section{Disposition of the earlier premises}\label{app:dispositions}

An earlier conditional assembly of this construction assumed
\ref{prem:C1}--\ref{prem:C6} together with nine project premises P1--P9. Their status
here is as follows; \ref{prem:C1}--\ref{prem:C6} are the classical inputs of
Section~\ref{sec:classical}, in their final accepted formulations.
{\small
\begin{longtable}{@{}p{0.27\textwidth}p{0.17\textwidth}p{0.48\textwidth}@{}}
\toprule
Premise & Status & Where\\
\midrule
P1 endpoint neurons & proved & Theorems~\ref{thm:E-product}--\ref{thm:E-plane}; Appendix~\ref{app:E}\\
P2 central reflection positivity, all finite parameters (with the general invariant-potential form) & proved, without the earlier degree and mass hypotheses & Theorem~\ref{thm:E-RP}, Corollaries~\ref{cor:E-P2}, \ref{cor:E-general}\\
P3 width moments at fixed parameters & removed; not used & local seminorm moments along the width arrow are proved in Proposition~\ref{prop:H1arrows}(i)\\
P4 finite-region Gaussian mass tilt and increment bound & proved & Theorems~\ref{thm:G-shift}--\ref{thm:G-diff}, \ref{thm:G-point}; the endpoint value $a=1/m$ is excluded\\
P5 common Gaussian space & proved, with Borel versions at $a=0$ & Theorems~\ref{thm:G-common}--\ref{thm:G-versions}\\
P6 fixed-volume Wick integrability & proved, with the precise construction statement & Theorem~\ref{thm:W}\\
P7 regulator products & removed; not used & regulator arrow proved directly (Propositions~\ref{prop:reg}, \ref{prop:H1arrows}(iv))\\
P8 stagewise finite-vector bridge & removed by new proofs & Section~\ref{sec:fixed} (arrows on $X$) and Theorem~\ref{thm:LR} (classical continuum limit); the earlier transfer claim is not used\\
P9 exact-cell midpoint convergence & removed; not used & no midpoint law occurs\\
\ref{prem:C1}--\ref{prem:C5} & classical, attributed & Section~\ref{sec:classical}\\
\ref{prem:C6} & standard tools & Section~\ref{sec:classical}\\
\bottomrule
\end{longtable}}

% ---------------------------------------------------------------------------
\section{Notation}\label{app:notation}

{\small
\begin{longtable}{@{}p{0.25\textwidth}p{0.67\textwidth}@{}}
\toprule
Symbol & Meaning\\
\midrule
$m,\ell,u$; $a$, $\mu=1/a$ & mass, length, quartic coefficient (no factorial); regulator and endpoint rate\\
$\kappa_a^{2}=a^{-2}-m^{2}$, $r_a=1-a^{2}m^{2}$ & mass-tilt coefficient; contraction constant\\
$\Cend$, $C_a$, $\Cfree$, $C'_{\Omega,a}$, $\Delta_{\Omega,a}$ & endpoint, regulated, free, tilted covariances; $C_a-C'_{\Omega,a}$\\
$c_a=C_a(0)$ & Wick constant: variance of the regulated Gaussian physical reference field (covariance $C_a$); not the endpoint variance $1/4$, and not asserted to be a variance of any finite-width or interacting law\\
$\Lambda$, $\Lambda_n=(-n\ell,n\ell)^{2}$, $\Omega=[-R,R]^{2}$ & physical squares (axis-parallel, centered); auxiliary square\\
$F^{b}_{N,a}$, $b\in\{\bprod,\bMon\}$ & width-$N$ neural field\\
$V=W+\varepsilon J-T$, $v_V$ & full potential and its integrand\\
$P^{b}_{N,\varepsilon,\Omega,a;\Lambda}$, $\nu^{v}$, $\Res$, $\nu^{(2)}$, $\nu_{a,\Lambda}$, $\nu_n$ & neural, full-tilt, residual, zero-stabilizer, physical and continuum laws\\
$\nu_{m,u,\Lambda}$, $\nu^{S}$, $\nu$ & classical finite-volume law, target on $\cS'$, its $X$-valued version\\
$E$, $X$, $p_j$, $d$, $\beta$, $\dist_n$ & field spaces, local seminorms, metric, bounded-Lipschitz distance, selection discrepancy\\
$(\Gamma_0,\cF_0,\PP_0,W)$, $\phi_a$, $\tilde\phi_a$, $\Phi_a$ & common Gaussian space, smeared and point fields, continuous modification, Borel version\\
$U_{a,\Lambda}$, $Z_{a,\Lambda}$, $K_\Lambda$, $\Xi_p$ & Wick interaction, normalizer, fixed-volume Wick constants\\
$K_j=\|\chi_j\|_2/2$, $\sigma_j$ & distinct neural reference-moment constants; original $K_j$ bound restored in AC3\\
$\theta_1,\theta_2$, $H^{\theta}_{+}$, $\Sigma^{\theta}_{+}$ & central reflections, half-planes, half-plane $\sigma$-algebras\\
$v^{\otimes k}$, $\cS^{\rho}_{k}$, $A_{j,k}$, $\Schw_k$ & tensor evaluation, correlation distributions, coefficient norm, Schwinger distribution\\
\bottomrule
\end{longtable}}

% ---------------------------------------------------------------------------
\section{Historical records and paths not used}\label{app:historical}

\paragraph{The historical ergodicity route.} An earlier target record obtained
ergodicity of the target through \cite[Thm.~4.4]{FS1977} and an identification of
the Osterwalder--Schrader semigroup with a Wightman energy, using the spectrum
condition and uniqueness of the vacuum in the course of proving condition (d) of
\ref{prem:C3}(i). Its first unsupported step was that identification. The route
is not used: condition (d) is proved in Theorem~\ref{thm:CT}(iii) by Euclidean
mixing from \cite[(7)--(8), (13)]{AHKZ1989}, before \cite[Thm.~1.0]{FS1977} is
applied.

\paragraph{Replaced sentences of the earlier assembly.} (i) A compressed sentence
in the setting described the mass-shift identity as if it applied to the
neural laws; it is replaced by the stage order of Section~\ref{sec:massshift},
where the identity concerns Gaussian reference covariances only, the pure
quadratic tilt on a finite auxiliary square has covariance $C'_{\Omega,a}$, and $C_a$
appears only in the auxiliary-volume limit. (ii) A universal sentence stating
that the interaction is never a pointwise fourth power is replaced by the precise
construction statement of Theorem~\ref{thm:W}(b) and Remark~\ref{rem:wick}. (iii) The
description of the zero-regulator interaction as ``the conventional Wick
integral'' is replaced by the same construction statement and by
Proposition~\ref{prop:identify}.

\paragraph{Geometry.} An earlier formulation required only that the physical
square be centered; a rotated centered square then breaks the invariance
hypothesis of Theorem~\ref{thm:E-RP} (Remark~\ref{rem:rotated}(i)). The definition is
restricted to axis-parallel squares, which contain the main sequence unchanged;
no failure of positivity is asserted for any rotated square, and the general
invariant-potential clause is unchanged. A proposed counterexample with an
outer rectangle invariant under only one reflection lies outside the hypotheses,
which require both rectangles to be centered.

\paragraph{Historical SMH comparison and the AC3 additions.}
SMH's comparison of AC2 with original H1 returned MISMATCH: the original
$K_j$ bounds (2.2), (2.4) had been replaced by different constants, and
the generic centered-square wording had been narrowed to axis-parallel
squares. That exact historical verdict remains unchanged. AC3 writes the
separate proposed numerical repair in
\eqref{eq:H1-original-reference} and \eqref{eq:H1-original-weighted}; its
current-version review and historical statement match remain companion
obligations. The geometry disposition above remains: the original generic
orientation wording is not recovered, the actual $\Lambda_n$ and two
central-line conclusions are unchanged, and no new rotated-square claim
is made. OB4's earlier PROVED WEAKER verdict and OB4-003/OB4-010 remain
historical. The source-only Wick/normalizer, sign, kernel and conditional
density proofs now written in this text address their mathematical issues;
the author does not close those findings or accept the integrated result.

\paragraph{Corrections of version A.} This version corrects the following
statements of the earlier version A, whose bytes are preserved unchanged. (i) The
readings lemma admitted convergence of a characteristic function at a single
frequency; that alternative is false (Remark~\ref{rem:readings}) and is replaced by
convergence at every real frequency, which the proof of Theorem~\ref{thm:LR}
already supplies. (ii) A remark asserted that the full Gaussian normalizer grows
with the auxiliary square; that assertion is withdrawn
(Remark~\ref{rem:chain}); monotonicity of $\Det_{\Omega,a}$ and the residual bounds are
unchanged. (iii) The empty-set convention $\sup\emptyset=0$ is restored for $M_N$, and the
early-volume cutoffs of Theorem~\ref{thm:V1} are chosen at least $2$. (iv) The
positive constants $M,\vartheta$ are restored in the continuity hypothesis of
Lemma~\ref{lem:G-S}. (v) The Wick variance $c_a$ is identified throughout as the
variance of the regulated Gaussian physical reference field. (vi) Absolute
``no rate'' wording is narrowed to the absence of an effective end-to-end rate or
schedule, with the fixed-volume $\sqrt a$ bounds kept. (vii) Authorship lineage is
recorded per component (Appendix~\ref{app:provenance}). (viii) An undefined macro
and several overfull displays are corrected without change of any formula.

\paragraph{Withdrawn formulations in the antecedents.} An unclipped
parameter-map clause (false as a total map; replaced by Lemma~\ref{lem:E-fields}(iii));
an unconditional independence assertion for competing clocks (replaced by
Lemma~\ref{lem:E-clocks}); the reading ``Lebesgue measurable on $\R^{3}$'' of the
potential hypothesis (fails; Remark~\ref{rem:rotated}(iii)); a sentence extending
the increment bound to $a=1/m$ (not part of the proved statements and not used);
the stagewise transfer and midpoint premises P8, P9 (not used); and a
frequency-split Nelson route (checked, not used).

% ===========================================================================
\begingroup
\small
% Private post-v1 supplement. Authored integration candidate, 26 October 2026.
% Every result below is an application/restatement of the named companion packet.
\clearpage
\section{Quantitative observation control and the pressure target}
\label{qrs:supplement}

The construction above reaches one interacting field through an
observable-independent sequence. The subsequent analysis asks a more focused
question: for specified smooth measurements and a requested tolerance, which
finite stages suffice? The answer proceeds through locality, a finite
interaction-box estimate, and five compatible errors. An explicit example then
supplies one finite tuple for both process-neuron architectures. A separate
pressure argument identifies that example's target and permits an application
of classical reconstruction at an explicit coupling cap.

This appendix records applications of the complete companion results, rather
than introducing a further synthesis theorem. Its mathematical premises are
the literal accepted statements and restrictions in the packets named below.
The public \href{run:PROOF-MAP.md}{claim-to-proof map} identifies the arguments
contained in this edition. The \href{run:COMPANION-INDEX.md}{companion index}
distinguishes this public manuscript source from the original private proof
packets, decisions and software. Those original packets are not exported here;
their full named statements remain premises where explicitly invoked below.
Historical candidate headers describe their authoring stage. Internal
acceptance and public source availability are separate qualifications.

\subsection{Two parameter routes and one field convention}
\label{qrs:routes}

The original main theorem retains $0<u/m^2<g_*$ and its fixed admissible
reference pair. Its universal sequence, selected before observations, remains
the sequence proved above. The locality and budget packets CL1, FV, EB1 and
EC1 were initially proved on
$0<u/m^2<\min(g_*,1/8000)$. No numerical lower bound for that originally
fixed $g_*$ is inferred here.

The distinct packet ADMIT (ADMIT-PRESSURE-v1) has the literal premises
P1--P6. P1 supplies the finite ferromagnetic DSS/SK-q interfaces, including
Gaussian sites and added mass at every finite positive scale. P2 is the
finite-scale P* bound at its exact coupling, cutoff and torus restrictions.
P3 gives fixed-square common-Gaussian Wick and source transfer with finite
local exponential moments. P4 is fixed-mesh periodization. P5 is the original
Fr\"ohlich--Simon pressure/state/exhaustion theorem, with its converse
pressure-differentiability implication. P6 lists the additional physical-law,
field-space, local-density, DSS/SKQ/P*/H1/MIX/REG/FLOW/TP/T and fixed-stage
NR1/L0 interfaces. These are named mathematical premises with the
qualifications in their complete statements, not conclusions of arithmetic
verification.

In the inclusive range $0<u/m^2\le1/8000$, ADMIT AP2--AP5 identify the
common original Section-4 maximal state of \cite{FS1977}, denoted here by
$\mu$, and its same-distribution $X$ version $\mu^X$. For the quantitative
neural budget, AP6 requires the strict range
$0<u/m^2<1/8000$ and every literal P6 interface. AP6 rebuilds applicability
of the actual five-error argument to that same state; it does not invoke an
old global theorem outside its $g_*$ domain. Its original private proof is ADMIT, Sections
1--12 as indexed in the private acceptance map. This is the route used in the numerical
example and in the strict-cap target interpretation below.

All routes use the whole-plane covariance $(-\Delta+m^2)^{-1}$, Fourier
measure $dp/(2\pi)^2$, the same-mass coefficient-one fourth Wick integral
$J_B$, and the potential $uJ_B$ with no factorial. A bounded square restricts
the interaction while the field occupies the plane. In particular the finite
law is
\begin{equation}\label{qrs:finite-law}
 d\mu_B=Z_B^{-1}e^{-uJ_B}\,dP_m.
\end{equation}
AP4 identifies the original reference/dilation expression as the same
target. Equality of these two specific maximal states is distinct from
uniqueness in a larger class of Gibbs states.

\subsection{Local response and the finite interaction box}
\label{qrs:locality}

Fix real $f=(f_1,\ldots,f_r)\in C_c^\infty(\R^2)^r$, with support union
$K\subset[-R,R]^2$, and let $Y_f(v)=(v(f_i))_i$. The observation distance
throughout this appendix is
\begin{equation}\label{qrs:BL}
 d_{\mathrm{BL}}(\rho,\eta)
 =\sup_{\substack{\|G\|_\infty\le1\\ \mathrm{Lip}_{2}(G)\le1}}
     |\rho(G)-\eta(G)|.
\end{equation}
The two constraints are separate, and the Lipschitz norm is Euclidean on
$\R^r$. Repeated or linearly dependent tests need no density assumption.

For $\Lambda_j=(-j\ell,j\ell)^2$ set $S_j=\Lambda_{j+1}\setminus\Lambda_j$,
$\Delta A_j=4\ell^2(2j+1)$ and $d_j=\dist(K,S_j)$. On the common free
Gaussian space put
\begin{equation}\label{qrs:shell}
 U_{j,t}=u[(1-t)J_{\Lambda_j}+tJ_{\Lambda_{j+1}}],\quad
 W_j=u(J_{\Lambda_{j+1}}-J_{\Lambda_j}),\quad
 d\mu_{j,t}=Z_{j,t}^{-1}e^{-U_{j,t}}dP_m.
\end{equation}
The shell insertion is a difference of fourth-chaos variables, so it uses
no pointwise quartic of a distribution at zero regulator.

\paragraph{Application CL1.}
Under CL1's original classical hypotheses and
$0<u/m^2<\min(g_*,1/8000)$, choose a positive integer $n_0$ with
$K\subset\Lambda_{n_0}$. One constant $C_f<\infty$ works for every
$j\ge n_0$, every $t\in[0,1]$, and every real $G\in C_b^1(\R^r)$ with
the two separate unit bounds:
\begin{equation}\label{qrs:CL1}
 |\mathrm{Cov}_{\mu_{j,t}}(G(Y_f),W_j)|
 \le C_f u\Delta A_j e^{-\kappa d_j},\qquad \kappa=m/256.
\end{equation}
The complete CL1 proof, Sections 1--6, checks the common finite-lattice
parameter intersection, confinement classes and all actual A1--A6/FLOW
premises. It first proves a common torus estimate, then takes the
fixed-mesh torus limit using TP and finally the plane mesh limit using T.
The common bound precedes both pointwise limits. No separate unproved
continuum response or spatial mixing assumption is added.

For reference, the literal coefficient used for the numerical work is the
CL1 coefficient below. Here
$\epsilon_*=\min(1/4,1/m)$ is an auxiliary \emph{lattice cutoff bound},
distinct from the physical stabilizer $\varepsilon$; $s$ is a covariance
scale, and $C_H,C_{H2}$ are the H1 and MIX coefficients after their actual
premises have been discharged:
\begin{align}
 \bar F&=(2R+2\epsilon_*)\Big(\sum_i\|f_i\|_\infty^2\Big)^{1/2},
 &\bar\delta&=\frac{2}{m^2\ell^2}\log\frac1{1-e^{-4m\ell}},
 \label{qrs:CL1-data}\\
 g(s)&=1+\bar\delta+\frac\pi{16}\log(1+(m^2s)^{-1}),
 &b_{CT}&=1+\frac{81\kappa}{5m},\nonumber\\
 C_*&=\frac{486\kappa}{5m}+\frac{6561\kappa^2}{25m^2}
       +b_{CT}^2\left(\frac{4\sqrt{2\pi}\kappa}{m}
                         +\frac{3300u}{m^2}\right),\nonumber\\
 M(s)&=32\ell^{-2}\left(\frac{C_H^2g(s)^3}{12}
                          +\frac{u^2C_{H2}^2g(s)^7}{9}\right),\nonumber\\
 C_f&=e^{8\sqrt2\kappa/m}b_{CT}^2\bar F
 \left[e^{C_*}\int_0^1\sqrt{M(s)}\,ds
 +\frac{e^{2C_*}\sqrt{M(1)}}{m^2(1+m^2)}\right].
 \label{qrs:CL1-constant}
\end{align}
The constant $81/5=16.2$ retains the smoothed Combes--Thomas weight
restriction. All source, Gaussian/Wick, DSS/BD/BBD/Griffiths--Simon,
SK-q and P* qualifications are those of the companion statements.

\paragraph{Application FV (CL1-FV1, version 2).}
Let $\rho_j$ be the law of $Y_f$ under \eqref{qrs:finite-law} for
$B=\Lambda_j$, and let $\rho$ be the observation law of the original
target. On the same old coupling range, for $n\ge n_0$ and $n\ell\ge R$,
\begin{align}
 q&=e^{-\kappa\ell},\qquad
 T_n=q^n\left[\frac{2n+1}{1-q}+\frac{2q}{(1-q)^2}\right],\nonumber\\
 E_n&=4C_fu\ell^2e^{\kappa R}T_n,\qquad
 d_{\mathrm{BL}}(\rho_n,\rho)\le E_n
       =O_{m,u,\ell,K,f}(ne^{-m\ell n/256}).
 \label{qrs:FV}
\end{align}
The enclosure $K\subset[-R,R]^2$ is required when using
$d_j\ge j\ell-R$.

For clarity, the FV application is the full proof of FV1--FV6, Sections
1--7. At each fixed $j$, exponential moments give an integrable envelope
for $e^{-U_{j,t}}$ and $W_je^{-U_{j,t}}$, including both endpoints.
Normalized differentiation is therefore
$\partial_t\E_{\mu_{j,t}}G(Y_f)=-\mathrm{Cov}_{\mu_{j,t}}(G(Y_f),W_j)$.
The endpoints agree as actual densities on the same Gaussian space.
Finite telescoping uses \eqref{qrs:CL1} and the exact sum
$\sum_{j=n}^{M-1}(2j+1)q^j=T_n-T_M$.
Nonnegative unit-mass Euclidean mollification preserves both norms in
\eqref{qrs:BL}; it extends the smooth estimate to every admitted Lipschitz
$G$. The identified target limit is taken for each $G$ before taking the
supremum. These are the precise steps in the companion argument, including
its physical/classical target identification for every $\ell>0$.
Empty/all-zero tuples have zero error directly; their arbitrary-$R$
zero-error convention does not remove the enclosure from a geometric
distance assertion. The literal earlier FV geometry statement remains a
rejected version in the companion record.

\subsection{Five errors between actual normalized laws}
\label{qrs:budget}

Application EB1 uses the old coupling range and the exact inherited
Gaussian/Wick, finite-stage NR1/L0 and FV premises. Fix a physical box
$\Lambda=\Lambda_n$, $0<a<1/m$, a positive stabilizer $\varepsilon$ and
$\Omega=[-R_{\rm aux},R_{\rm aux}]^2\supset\bar\Lambda$.
For one architecture the adjacent whole-field laws are
\begin{equation}\label{qrs:chain}
 P^b_{N,\varepsilon,\Omega,a;\Lambda}
 \longrightarrow Q_\varepsilon\longrightarrow Q_0
 \longrightarrow Q_a\longrightarrow Q_{\rm phys}
 \longrightarrow \nu.
\end{equation}
Here $Q_\varepsilon=\nu^v_{\varepsilon,\Omega,a;\Lambda}$ is the full
endpoint-Gaussian tilt and equals the pushed-forward residual law by the
mass-tilt identity. $Q_0=\nu^{(2)}_{\Omega,a;\Lambda}$ is the residual
zero-stabilizer endpoint, $Q_a=\nu_{a,\Lambda}$ is the physical regulated
law and $Q_{\rm phys}=\nu_{0,\Lambda}$ is its fourth-chaos limit. Every
arrow compares the same observation map under normalized probabilities.

Write $A=|\Lambda|$, $B_{\rm ext}=|\Omega\setminus\Lambda|$,
$c=c_a=C_a(0)$, $\kappa_a^2=a^{-2}-m^2$ and
$F_p=(\sum_i\|f_i\|_p^2)^{1/2}$.
The physical covariance $C_a$ has multiplier
$(m^2+w^2+q^2+a^2w^2q^2)^{-1}$; its variance $c$ differs from the
endpoint variance $1/4$. Put
$\Delta_{\Omega,a}=C_a-C'_{\Omega,a}$,
$\delta_\Omega(H)=\sup_{x\in H}\Delta_{\Omega,a}(x,x)$ and
$L_c=(4\sqrt{15}+12)c^{3/2}$. Then EB1.b--d give
\begin{align}
 S&=\min\{2,6\varepsilon B_{\rm ext}c^2e^{9uAc^2}\},
 \label{qrs:physical-errors}\\
 A_{\rm aux}&=\min\{2,e^{9uAc^2}
       [2uAL_c\sqrt{\delta_\Omega(\bar\Lambda)}
                  +F_1\sqrt{\delta_\Omega(K)}]\},\nonumber\\
 R_{\rm phys}&=\min\{2,\Xi_2(\Lambda)^{1/2}
                      [4K_\Lambda\sqrt a+(a/2)F_2]\}.
 \nonumber
\end{align}
The constants $K_\Lambda,\Xi_p$ are exactly those of the original Wick
theory and EC1 below. Variation in the stabilizer bound is the convention
$\sup_{|H|\le1}$ with maximum two. The auxiliary proof couples independent
Gaussian covariance defects and controls both required compact sets.
The regulator proof uses the actual common-space Wick weights and positive
normalizers at fixed physical volume.

For completeness the exact sufficient width coefficient is given here.
Choose a bounded box $D$ containing $\Omega$ and every test support, put
$O=|\Omega|$, $d_D=|D|$, $F=(\sum_i\|f_i\|_{4/3,D}^2)^{1/2}$,
$\alpha=\min(u,\varepsilon)$, $U=\max(u,\varepsilon)$ and
$B_N=\kappa_a^2/2+6uc$. In EB1's equal-square case
$\Lambda_0=\Lambda_1=\Lambda$, set
\begin{align}
 L_N&=B_N^2O/(2\alpha),
 &v_1(x)&=4Ux^3+2B_N\sqrt O\,x,\nonumber\\
 v_2(x)&=12Ux^2+2B_N\sqrt O,
 &v_3(x)&=24Ux,\quad v_4(x)=24U,\nonumber\\
 e_0&=1,\quad e_1=v_1,\quad e_2=v_1^2+v_2,
 &e_3&=v_1^3+3v_1v_2+v_3,\nonumber\\
 e_4&=v_1^4+6v_1^2v_2+3v_2^2+4v_1v_3+v_4.
 \label{qrs:width-polynomials}
\end{align}
For a nonnegative-coefficient polynomial $P(x)=\sum p_kx^k$, define
\begin{align}
 b_0&=1,\quad b_k=(k/(2\alpha e))^{k/4}\ (k>0),
 &\mathcal T(P)&=e^{L_N}\sum p_kb_k,\nonumber\\
 M_Z&=\mathcal T(e_4),
 &M_A&=\mathcal T\left(\sum_{j=0}^4\binom4j F^je_{4-j}\right),\nonumber\\
 z&=e^{-3UO/16-3uc^2A},
 &C_{\rm law}&=\frac{d_D(M_A+M_Z)}{96z},\nonumber\\
 D_{\rm law}&=\frac43(3\sqrt6 C_{\rm law})^{1/4}
                       (2\sqrt r)^{3/4}.
 \label{qrs:width-coefficient}
\end{align}
For $r\ge1$ the width term is $W=\min(2,D_{\rm law}N^{-1/4})$.
Its complete proof is EB1/NLAW, Sections 1--7: fourth-order Banach
replacement uses full sign symmetry, matching covariance and the exact
coefficient $d_D/(96N)$; positivity gives a normalization cost $1/z$;
L0 smooths in ambient observation space. It proves both the $C_b^4$ law
bound $\min(2,C_{\rm law}/N)$ and this all-width BL bound. NR1 itself is
a fixed-stage response theorem; the law consequence is proved in EB1,
with NR1's sharper-author/larger-blind-core qualification retained.

With $E_n$ from \eqref{qrs:FV}, the complete EB1.f application is
\begin{equation}\label{qrs:five-errors}
 d_{\rm BL}\bigl(\mathcal L_f(P^b_{N,\varepsilon,\Omega,a;\Lambda_n}),
                        \mathcal L_f(\nu)\bigr)
 \le \min\{2,W+S+A_{\rm aux}+R_{\rm phys}+E_n\}.
\end{equation}
The EB1 integration paragraph and PHYS P1--P6 contain the full adjacent-law
proofs, the local-$H^{-2}$ metric bridge and the original-sequence
application. In particular the unchanged universal sequence satisfies
$\min(2,L_{f,J}2^{-n}+E_n)$ in this observation distance, where
$L_{f,J}=2^J\max(2,(\sum_i\|f_i\|_{H^2(\T_J)}^2)^{1/2})$ and every support
lies in $(-J,J)^2$. This gives no relation between $N_n$ and $n$.
For the explicit-cap use, AP6 supplies the same actual-law budget with
$\nu$ replaced by the identical pressure-identified target $\mu^X$,
under its own strict cap and every P1--P6 premise.

\subsection{Elementary enclosures and conditional selection}
\label{qrs:effective}

The five clauses of EC1-v2 give elementary enclosures and conditional
selection, with their complete proofs in GROWTH Sections 1--5. For example,
\begin{align}
 c_a&\le\bar c(a)=\frac{\operatorname{asinh}(1/(ma))+1}{2\pi},
 \label{qrs:variance}\\
 \int_{|x|>s}C_a(x)^2\,dx
 &\le\min\left\{\frac1{4\pi m^2},
                  \frac{e^{-\sqrt2ms}}{2\pi m^2}\right\}.
 \label{qrs:tail}
\end{align}
These estimates keep $a>0$ fixed in the auxiliary-volume step.
Multiplication of \eqref{qrs:tail} by $\kappa_a^2$ bounds both covariance
defects when $\Omega$ contains the Euclidean $s$-neighborhood of the two
required compacts; the centered-square sufficient condition is
$0<s<R_{\rm aux}-\max(n\ell,R)$.

The literal Nelson constants admit the following enclosure for $A>0$ and
every fixed $p\ge1$:
\begin{align}
 K_\Lambda^2&=\frac{75\sqrt{85}}{2\pi^3m}u^2A,
 &K_\sharp&=\max\{12uA,(2/e)K_\Lambda\sqrt{e(1+\sqrt2)/m}\},\nonumber\\
 \beta&=\pi^2/(48uA),
 &H_p&=\max\{K_\sharp,(4/\beta)[\log_+(32p/(e\beta))]^2\},\nonumber\\
 \Xi_p&=e^{pK_\sharp}
    +\int_{K_\sharp}^\infty pe^{pH}
                    e^{-\exp(\sqrt{\beta H}-1)}\,dH\nonumber\\
 &\le e^{pH_p}+e^{-pH_p}\le2e^{pH_p}.
 \label{qrs:Nelson}
\end{align}
Here $\log_+x=\max(0,\log x)$. The symbol $\beta$ in this local formula
is a positive scalar, distinct from the whole-field probability metric.
No volume-uniform exponential moment is asserted.

EC1.e is conditional computability for \emph{one fixed finite tuple}.
Its additional numerical inputs are computable $m,u,\ell$, an actual
certificate for the original coupling gate, a valid $n_0$, a certified
upper bound for $C_f$, a support enclosure and upper bounds for $F_1,F_2,F$.
With positive allocations to the five errors it selects, in order,
\begin{equation}\label{qrs:selection-order}
 n,\qquad a,\qquad R_{\rm aux},\qquad\varepsilon,\qquad N.
\end{equation}
This reverses the order in which limits are taken. Its strict interval
searches terminate without exact-real equality decisions. EC1 by itself
supplies neither the old $g_*$ certificate nor a numerical locality
coefficient; ADMIT and SCALE/OPT furnish the separate fixed-example inputs.
The literal width majorant contains
$\exp(B_N^2O/(2\alpha)+3UO/16+3uc_a^2A)$ and the full positive polynomial.
EC1.d characterizes growth of this sufficient rule. It supplies no necessary
width or complexity lower bound.

\subsection{One explicit nonzero observation and two sufficient tuples}
\label{qrs:certificate}

Fix $m=1$, $u=10^{-10}$ and $\ell=1000$. The test is the single nonzero
compact smooth product bump
\begin{equation}\label{qrs:bump}
 f(x_1,x_2)=b(x_1)b(x_2),\qquad
 b(t)=\begin{cases}\exp[-1/(1-t^2)],&|t|<1,\\0,&|t|\ge1.\end{cases}
\end{equation}
It has support $[-1,1]^2$, $R=1$, $n_0=1$ and the certified norm bounds
$F_1\le4$, $F_2\le2$, $F\le3$.
The numerical ratio is strictly below $1/8000$. The target is precisely
$\mu^X$ from AP3--AP6. All finite interfaces, $K_{\rm cov}=1$, source
qualifications and both process-neuron architectures are retained.

NS1 uses SCALE's $C_{H2}\le2^{1800}$ and $C_f\le2^{1810}$.
OPT proves a changed bound for the same finite-interface H2 row,
$C_{H2}\le2^{384}$, and the \emph{literal} coefficient in
\eqref{qrs:CL1-constant} satisfies $C_f\le2^{360}$.
The proof is OPT Sections 1--9: finite original-law moment and shift
estimates, mixed real finite differences through derivative order eight,
connected Bell partitions with both overlap charges and every bridge
binomial, the actual source integral, and a separate H1 coefficient
evaluation. The complete finite domain includes all allowed lattice
parameters, $0<s\le1$, both profile endpoints, coincident and
zero-profile sites. The coefficients are conclusions on the literal
conditional interfaces, with AP6 supplying their target applicability.

One displayed Wick-factor estimate was omitted from the frozen OPT author
proof. The current OPTREF/OPTMATCH derivations complete that step from the
supplied moments. Explicitly, with $g=1+\bar\Gamma_s$ and
$|\eta|\le g$, those moments give
$\E\phi^{2q}\le(2q-1)!!g^q$. Since
$7!!=105<2^8$, $15!!=2027025<8^8$ and
$23!!=316234143225<32^8$, the triangle inequality in $L^8$ yields
\begin{equation}\label{qrs:Wick-L8}
 H_0=1,\quad \|H_1\|_8<2\sqrt g,\quad
 \|H_2\|_8<9g,\quad \|H_3\|_8<38g^{3/2},\qquad
 \|H_j\|_8\le8^jg^{j/2}\ (0\le j\le3).
\end{equation}
Here $H_1=\phi$, $H_2=\phi^2-\eta$ and
$H_3=\phi^3-3\eta\phi$. The largest required degree is three.
Thus the relation uses an explicit argument and no added moment premise;
the original omission remains visible in the frozen author proof.

\begin{table}[!htbp]
\centering\small
\caption{Two sufficient tuples for the same observation, target and requested
$1/5$ tolerance. Each nonspatial error is bounded by $1/256$ in both rows.
The optimized tuple uses a looser spatial allocation.}
\label{qrs:tuple-table}
\begin{tabular}{@{}p{0.24\textwidth}p{0.33\textwidth}p{0.33\textwidth}@{}}
\toprule
Quantity & NS1 baseline & OPT sufficient tuple\\
\midrule
$C_{H2}$ upper bound & $2^{1800}$ & $2^{384}$\\
$C_f$ upper bound & $2^{1810}$ & $2^{360}$\\
$n$ & $512$ & $74$\\
$a$ & $2^{-16777216}$ & $2^{-16384}$\\
$R_{\rm aux}$ & $18446744073710063617$ & $137439027473$\\
$\varepsilon$ & $2^{-4611686018427387904}$ & $2^{-68719476736}$\\
$N$ & $2^{(2^{4611686018494496904})}$ & $2^{(2^{68719542354})}$\\
Spatial error upper & $2^{-736}$ & $1/4096$\\
Certified total & $1/64+2^{-736}<1/5$ & $65/4096<1/5$\\
\bottomrule
\end{tabular}
\end{table}

For OPT use $\Omega=D=[-R_{\rm aux},R_{\rm aux}]^2$ and
$\Lambda_0=\Lambda_1=\Lambda_{74}$. Sections 10--17 of its complete proof
establish, in the order of \eqref{qrs:selection-order},
\begin{equation}\label{qrs:OPT-errors}
 E_{74}\le\frac1{4096},\qquad
 R_{\rm phys},A_{\rm aux},S,W\le\frac1{256},\qquad
 d_{\rm BL}(\mathcal L_f(P^b),\mathcal L_f(\mu^X))
    \le\frac{65}{4096}<\frac15.
\end{equation}
The NS1 complete proof gives the other row, including its five actual
errors. Both applications retain the normalizers, physical variance,
improper Nelson integral, two covariance defects and full width polynomial.
The new tuple is smaller at the same requested tolerance; its certified
error sum is slightly larger. This comparison asserts neither improved
accuracy, optimality nor runtime speedup.

Every displayed power denotes one finite mathematical number. The
certificate language has canonical integers, exact dyadics $2^z$ and
recursive integer-valued powers of two. It represents
$2^{(2^{68719542354})}$ without expanding the integer. Width counts
whole-plane process neurons with their internal randomness; it remains
astronomical and supplies no practical sampler.

\paragraph{Software assurance and its resource scope.}
The accepted OPT executable verifies selected finite arithmetic and frozen
source bindings. The binding resource erratum qualifies the original
``whole-command alarm'' wording: API deadline checks are cooperative;
the POSIX alarm starts after argument parsing and covers the protected
verification phase (including a selection certificate write), then is
cancelled before a final verification result-file write, serialization and
standard output. Output-device latency can therefore exceed the configured
phase limit. No hard whole-process time or memory guarantee is supplied.
The selected author monitor's 385-bit watched values differ from the
independently implemented expanded-polynomial calculation's 131349-bit
high-water value. Neither computation materializes a giant power or neuron
field. Code, pin tables and stable source files remain external trust
roots; the package is not a formal proof kernel or an atomic adversarial
filesystem guarantee. The NS1 rational representation attempt that exceeded
its original 500000-bit ceiling remains a recorded inconclusive attempt.
Historical candidate flags inside frozen code/certificates are governed by
the separate current decisions and the binding erratum.

\subsection{Reconstruction at the explicit pressure cap}
\label{qrs:reconstruction}

The original reconstruction theorem above retains its original hypotheses.
Here the exact application is REC-EXPLICIT-CAP-v2 (packet REC2), clauses
(a)--(g), with its complete proof Sections 1--10. For the core clauses
(a)--(f), assume $0<u/m^2\le1/8000$, AP-source/AP-version and the literal
classical inputs C1/C3/C6 in REC2's statement. These include the two
specific original FS Section-4 states; their local positive Gaussian
$L^p$ densities under either boundary-algebra reading BC1 or BC2; proper
Euclidean invariance; either the full measure OS reading M or the
strictly time-ordered polynomial reading P; qualitative polynomial
clustering from \cite[Thm.~4.4, p.~511]{FS1977};
\cite[Thm.~1.0, p.~495]{FS1977}'s reconstruction theorem; and the named
kernel and tube-uniqueness inputs C3(ii),(iii). BC1 is the intersection
of neighborhood algebras of the boundary. BC2 uses boundary-supported
compact $H^{-1}$ evaluations obtained by Gaussian $L^2$ approximation and
Gaussian-null completion. These source premises and their access limits
remain literal; no new numerical density norm is assumed.

\paragraph{Regularity, mixing and positivity before reconstruction.}
For every real Schwartz $f$ and real $t$, REC2.b gives
\begin{align}
 1\le\E_\mu e^{t\phi(f)}
 &\le e^{3t^2\|f\|_2^2/(2m^2)},
 &N_m(f)&=\frac{\sqrt3}{m}\|f\|_2,
 \label{qrs:Schwartz}\\
 \E_\mu\phi(f)&=0,
 &\|\phi(f)\|_{L^2(\mu)}&\le N_m(f).\nonumber
\end{align}
This continuous Schwartz norm satisfies, on $N_m(f),N_m(g)\le1$,
\begin{equation}\label{qrs:unitball}
 \E_\mu e^{\phi(f)}\le e^{1/2},\qquad
 |\E_\mu e^{\phi(f)}-\E_\mu e^{\phi(g)}|
       \le2eN_m(f-g).
\end{equation}
For nonzero real $f$ and integers $k\ge1$,
\begin{equation}\label{qrs:moment}
 \E_\mu|\phi(f)|^k\le2e^{1/2}k!N_m(f)^k.
\end{equation}
These are target estimates. They make no global Schwartz error assertion
for finite neural fields.

REC2 Sections 1--4 match the exact pressure state to the classical
clustering theorem, extend compact source bounds to Schwartz tests by
cutoffs and Fatou, and use finite exponential moments to prove compact
cylinder-polynomial density in full $L^2(\mu)$. Classical polynomial
clustering and translation isometry therefore imply
\begin{equation}\label{qrs:mixing}
 \E_\mu[F(G\circ A_b)]\longrightarrow\E_\mu F\,\E_\mu G
 \quad (|b|\to\infty),\qquad F,G\in L^2(\mu).
\end{equation}
Invariant events then have probability zero or one. This establishes
ergodicity before reconstruction; it uses no reconstructed vacuum as a
premise. The same mixing statement holds on $X$ under $\mu^X$.

For any line $L$ and either one of its open half-planes let
$\Sigma_{L+}$ be its evaluation algebra. REC2.d is exactly
\begin{equation}\label{qrs:OS}
 \E_\mu[\overline{\Theta_L F}F]\ge0,
 \qquad F\in L^2(\Sigma_{L+},\mu).
\end{equation}
It also holds for $\Sigma_{L+}$-measurable $F$ whose reflected product is
absolutely integrable, and on the corresponding $X$ algebras. The
\emph{same-half-plane measurability} is essential for both assertions.
Under reading M it begins with the measure inequality. Under reading P,
REC2 Sections 5--6 first construct local $L^2$ moment kernels, approximate
positive-time polynomial tensors away from null equal-time hyperplanes,
and then use full half-plane $L^2$ density. Reflection isometry passes
positivity to $L^2$ limits. Truncation and dominated convergence give the
stated same-half-plane integrable extension. Euclidean invariance transfers
the result to any line. These complete arguments replace no hypothesis
by a model verdict.

\paragraph{Kernels and the attributed quantum consequence.}
REC2.e supplies symmetric tempered Schwinger kernels at every order,
with local moment and truncated kernels in $L^2(B^k)$ for each bounded
$B$. Its full Sections 5 and 7 use the Gaussian Wick map and local
$L^2$ density, the ordinary pairing and cumulant partition identities,
and the stated Schwartz-kernel theorem. AP5's non-Gaussianity implies
some even truncated pairing of order $k\ge4$ is nonzero. Local $L^2$
regularity rules out a witness supported solely on coincident-time
hyperplanes: real product tests can be chosen in pairwise disjoint
Euclidean-time slabs.

At this point REC2 Section 8 has established the four hypotheses of
\cite[Thm.~1.0]{FS1977}: proper Euclidean invariance, full half-plane
$L^2$ OS positivity, the continuous norm with
\eqref{qrs:unitball}, and translation ergodicity. That attributed
consequence of \cite{OS1973,OS1975} reconstructs the unique scalar
relativistic Wightman theory for these Schwinger data in $1+1$ dimensions,
with its unique physical vacuum. The retained analytic tube-identification
and boundary-uniqueness input then transfers the separated nonzero even
truncated Schwinger pairing to a nonzero truncated Wightman distribution.
Interaction here means failure of Gaussian factorization. No quantitative
mass gap or scattering result is inferred.

REC2.g additionally requires the \emph{strict} cap
$0<u/m^2<1/8000$ and every unchanged AP6/P1--P6 interface. It identifies
the numerical budget's target as this reconstructed field's Euclidean
law. Hence it applies to both rows of Table~\ref{qrs:tuple-table}, with
their existing errors, metric and widths. It supplies no extension of
the old universal neural sequence.

\paragraph{Retained correction and source boundaries.}
The unrestricted measurable-$F$ integrable OS clause in REC-v1 was false.
Choose nonzero real $h$ supported in the positive-time half-plane and put
$q=h-h\circ\theta$. The retained two-half-plane observable $F=\phi(q)$
satisfies $\Theta F=-F$ and has form $-\E F^2<0$. Its bounded sign
observable has reflected form $-1$ almost surely. Positive local Gaussian
densities justify nondegeneracy; these witnesses show why integrability
alone is insufficient.
The version-2 half-plane restriction is a genuine domain correction.
The old GAP/PROVED WEAKER assessments and CLOSED historical gate remain
in the companion provenance; only REC2's current exact statement is used
here. Source access is statement/application verification, with the
original classical bodies and antecedents attributed as recorded. The
kernel/tube premise and BC1/BC2 readings remain explicit. Neither arbitrary
Gibbs-class uniqueness nor the AHKZ boundary estimate is newly proved.

\subsection{Positioning and proof navigation}
\label{qrs:positioning}

The comparison is architecture specific. Frank's version 2 already gives
a rigorous finite-volume Wick-quartic neural construction with quenched
features, a finite width per shell, quantitative ultraviolet and
finite-observation bounds, commuting width/UV limits and target
identification \cite{Frank2026v2}. Removing its UV cutoff retains
infinitely many shells. Its coefficient-marginal total-variation statement
concerns a fixed finite-dimensional coefficient marginal; it does not say
that all ultraviolet/infinite-dimensional field laws converge in total
variation. The original Remark~\ref{rem:frank-moments} records the
smeared-polynomial-moment overlap. Generic smeared-moment convergence is
therefore not used as a novelty separator.

Generic neural representations and interacting network/path-integral
reparametrizations have predecessors
\cite{FerkoHalversonMutchler2026v1,HashimotoHironoMaedaTotsukaYoshinaka2024}.
Finite-width non-Gaussianity and density deformation also precede this
work \cite{HalversonMaitiStoner2021,Halverson2021Neurons,DemirtasHalversonMaitiSchwartzStoner2023}.
Roberts and Yaida's effective-network treatment, based on research with
Boris Hanin, relates a small quartic deformation and connected four-point
function in a perturbative inverse-width expansion, and gives the leading
quartic coupling for its specified second-layer marginal
\cite[Secs.~1.3 and 4.2]{RobertsYaida2022}. This is bounded
effective-network context and supplies no construction premise. The
preprint metadata names Roberts, Yaida and Hanin; the book cover names
Roberts and Yaida and credits Hanin's collaboration. The bibliography
retains that credit distinction, with fresh manuscript source adjudication
still required for this proposed addition.

Conditional-factorization and symmetric-Markov positivity mechanisms have
prior instances \cite{FerkoHalverson2025QMNN,FILS1978}. Harvey's
coordinate-product example already has the product-exponential reference
kernel and coordinate two-point positivity \cite{Harvey2026v2}.
The construction here verifies the full normalized annealed law's
positivity for every bounded half-plane observable at the two central
coordinate lines, under its invariant geometry. The finite-stage
interacting two-point law's Euclidean invariance is not established by
anisotropy of the free reference. The exact obstruction hypotheses in
the cited neural no-go results retain their separate scope.

Quantitative stochastic-quantization approximations and fixed-torus
equilibrium convergence also have established results
\cite{MaZhu2020Galerkin,MaWangYang2024FullDiscrete,TsatsoulisWeber2018Gap}.
Their norms, initial data, dynamical/equilibrium targets and fixed volumes
differ from the five-error infinite-volume observation problem here.
The bounded contribution is the specified annealed product/Mondrian
process ensembles, their full central-coordinate law positivity, and
this traceable sufficient observation certificate under named premises.
The comparison establishes neither priority nor efficiency.

\paragraph{Complete companion map and actual attribution.}
Table~\ref{qrs:proof-table} summarizes the applications and historical
complete-proof locations. The adjacent public
\href{run:PROOF-MAP.md}{PROOF-MAP} and machine-readable JSON index this
edition's source blocks. The original 23-packet index, complete private
statement/proof/decision bindings and archive paths are not public attachments. SOURCE is a corrected application/source audit, not an
additional theorem. Multi-result packets retain exclusions: in particular
SM's E4/$r_0$ supplement and CR Corollary, blanket P3 $p=p_1$ wording,
blind-only REG extensions and stronger blind-only coefficients are not
adopted. A later theorem's separate acceptance does not retroactively
accept an earlier excluded claim.

The baseline authorship and source-access record in
Appendix~\ref{app:provenance} is preserved as historical PV4 provenance.
This supplement's integration is by GPT-6.1 Sol at max effort under
Cody Mitchell's direction. It adopts the exact complete packet arguments,
retaining their individual Sol/Astra/Opus lineage, including the
OPTREF/OPTMATCH Wick-factor derivation. Verdicts alone confer no argument
authorship. The packet decisions record internal model review under their
actual policies, including fresh-session Sol-only policy 1.7 where used;
they are not independent expert review or formal machine-checked proofs.
That 26 October integration was an authored private candidate at the time.
The public REVIEW document records the later accepted manuscript baseline and
the scope of this edition's review. Version 1 remains separately archived.

\clearpage
\begin{longtable}{@{}p{0.12\textwidth}p{0.32\textwidth}p{0.47\textwidth}@{}}
\caption{Applications and historical complete-proof locations. Original
path/hash/line bindings are retained in the private research archive; the
public PROOF-MAP indexes this edition's source blocks.}\label{qrs:proof-table}\\
\toprule
Packet & Application in this appendix & Complete argument to apply\\
\midrule
\endfirsthead
\toprule Packet & Application in this appendix & Complete argument to apply\\
\midrule\endhead
CL1 & \eqref{qrs:CL1}--\eqref{qrs:CL1-constant} & CL1-two-limit assembly,
Sections 1--6; every prerequisite substitution and both ordered limits.\\
FV & \eqref{qrs:FV} & CL1-FV1-v2, FV1--FV6, Sections 1--7;
normalization, endpoints, telescoping, BL bridge and exact target.\\
EB1 & \eqref{qrs:physical-errors}--\eqref{qrs:five-errors} & Integration
clause map; complete NLAW Sections 1--7 and PHYS P1--P6.\\
EC1 & \eqref{qrs:variance}--\eqref{qrs:selection-order} & EC1-v2/GROWTH
a--e; complete Sections 1--5.\\
ADMIT & Explicit-cap target, AP2--AP6 & Complete AP1--AP6 pressure/source/
version/budget proof with literal P1--P6.\\
SCALE; NS1 & Baseline row of Table~\ref{qrs:tuple-table} & SCALE's coefficient
proof; NS1 Sections 1--10, especially 3--9 for every actual error.\\
OPT & Optimized row and \eqref{qrs:OPT-errors} & OPT Sections 1--17;
OPTMATCH's explicit missing-Wick-display derivation; resource erratum.\\
REC2 & \eqref{qrs:Schwartz}--\eqref{qrs:OS}, kernels and reconstruction &
REC2(a)--(g), Sections 1--10, using only its version-2 half-plane domain.\\
Support & Literal finite interfaces & NR1, SM, SKQ, WFSP, H1, H1PREQ,
GQ, H2RED, T, REG, MIX, TP, FLOW and qualified SOURCE, at their exact
indexed scopes.\\
\bottomrule
\end{longtable}

\clearpage

% MAN integration candidate, 27 October 2026. Frozen accepted inputs only.
% This file is also inlined, byte-for-byte, in neural-construction.tex.
% New labels use vqal:. Exact statement/proof ranges are in PROOF-MAP.json.
\clearpage
\section{27 October construction, vector accuracy and linkage}
\label{vqal:current}

This appendix brings the construction and the later quantitative results
to the same field law. The central distinction is between fixing a
sequence before measurements and selecting a sufficient finite tuple for
specified measurements. The first gives one construction of the field;
the second gives explicit, extremely large approximation bounds. Their
comparison is a theorem about probability laws, even when the parameters
of the first sequence have not been computed.

The seven inputs MOM, SEQ, NORM, SELECT, VEC, ACC and LINK have separate
current internal acceptance decisions at their literal conditional scopes.
The present integration is a candidate for fresh review. Original author
headers in their frozen proofs record the time of writing; they are
preserved alongside the later decisions. This appendix adopts their proof
arguments, with GPT-6.1 Sol integration authorship; the earlier Opus,
Astra-line and Sol argument lineage in Appendix~\ref{app:provenance}
remains. No internal review is described as an outside expert review or
as a machine-checked proof.

\subsection{Sources, laws and the two coupling domains}
\label{vqal:premises}

All field spaces, Fourier conventions, neuron laws, own-variance Wick
quartics, potentials, annealed probabilities and central reflections are
those of Definitions~\ref{def:laws} and~\ref{def:reflections} and the
preceding construction. In particular $u$ multiplies the Wick quartic
without a factorial, the free covariance is $(-\Delta+m^2)^{-1}$, and a
finite interaction square restricts the interaction of a whole-plane
field. The endpoint neuron covariance $e^{-|z|_1/a}/4$ has point variance
$1/4$; it is distinct from the physical covariance $C_a$ and its variance
$c_a$. Let $\nu^S=v_{m,u}$ be the common original Fr\"ohlich--Simon
Section-4 state, and put $\nu=\mu^X=R_*\nu^S$ on $X$. The realization $R$
preserves every compact smooth evaluation simultaneously on one
$\nu^S$-full measurable set. In this appendix $\mu=1/a$ in the neuron
definitions continues to mean inverse regulator, rather than a probability
law. We use $\nu_n=\nu_{0,\Lambda_n}$, $\Lambda_n=(-n\ell,n\ell)^2$,
for the actual original physical laws throughout.

The source assumptions behind the explicit cap are mathematical premises,
not consequences of a numerical certificate. For clarity their operational
content is recorded here. P1--P6 below are the accepted ADMIT premises;
the named interfaces retain their full source qualifications.

\paragraph{P1--P2: finite positive scales.}
On the shifted torus $\epsilon(\Z+1/2)^2/L\Z^2$, use the inner product
$\epsilon^2\sum$, $A=-\Delta^\epsilon+m^2$, the \emph{plane} lattice Wick
variance $c_\epsilon$, and profiles $q_x\in[0,1]$. P1 supplies DSS/SK-q
for the density proportional to
$\exp[-u\epsilon^2\sum q_x{:}\psi_x^4{:}_{c_\epsilon}]$
relative to $P_{A^{-1}}$, including Gaussian sites. At every finite real
field $H$ it supplies
\begin{equation}\label{vqal:P1}
 0\le\operatorname{Cov}_{\nu^H}(\psi_x,\psi_y)
       \le E_\nu(\psi_x\psi_y).
\end{equation}
The same statements hold after adding physical mass $1/s$, for every
finite $s>0$. P2 assumes
$0<\epsilon\le\min(1/4,18/(5m))$,
$L\ge\max(1,\sqrt2/m)$, $L/\epsilon\in2\N$, $L/\epsilon\ge4$,
and $0<u\le m^2/8000$. Core P* with $S_{\max}=\infty$, $\Phi=1$,
$\kappa=0$ gives
\begin{equation}\label{vqal:P2}
 \max_x\epsilon^2\sum_y S_{s,q}(x,y)
 \le {s\over1+m^2s}+{2s\over2+m^2s},\qquad
 d\mu_{s,q}\propto e^{-u\epsilon^2\sum q_x{:}\psi_x^4{:}_{c_\epsilon}}
                  dP_{(A+1/s)^{-1}}.
\end{equation}
This is a finite-positive-scale statement, with no original-law endpoint
included as an assumption.

\paragraph{P3--P4: fixed-square transfer.}
For each fixed $B=(-b,b)^2$, the commensurate shifted plane mesh
$\epsilon=b/N\downarrow0$ is on a common Gaussian space with
\begin{equation}\label{vqal:P3}
 \epsilon^2\sum_{x\in B}{:}\psi_x^4{:}_{c_\epsilon}\longrightarrow J_B,
 \qquad \epsilon^2\sum_x f(x)\psi_x\longrightarrow\phi(f)
 \quad\hbox{in Gaussian }L^2.
\end{equation}
Here $f$ is any real compact smooth test. For each fixed finite $p\ge1$,
the expectations of $\exp(-puJ_B^\epsilon)$ are uniformly finite for
sufficiently small $\epsilon$. The continuum normalizer is positive finite,
with finite Gaussian $L^p$ density at every finite $p$. At fixed mesh P4
gives entrywise torus-to-plane Gaussian covariance convergence on finite
site sets through admissible even periods; the plane Wick variance is
unchanged. These are fixed-square statements for all $m,u>0$.

\paragraph{P5: pressure and the original source states.}
For $P_h(X)=uX^4-hX$, use the normalized free pressure of
\cite[Proposition~2.1 and the remark after (2.3), pp.~502--503]{FS1977}.
Its actual Corollary~4.3(ii), p.~510, equates differentiability at zero in
$h$ with equality of the two original maximal states. The converse is
part of this premise; reversing the narrower old
\ref{prem:C1}(iii) would not supply it. Theorem~5.2, p.~512, then gives
ordinary free-boundary exhaustion convergence to that common state, with
compact-test generating functionals read at every real frequency.
These are attributed source facts whose full original proofs are not
reproved here.

\paragraph{P6: the additional physical and finite interfaces.}
Retain the all-$m,u>0$ physical/classical Wick-law identity, the complete
$X$ field-space and Borel-determining interfaces, and the original
Gaussian/common-space Wick laws proved in the paper. Retain the
all-semibounded-polynomial local-density and positive boundary
factorization of \ref{prem:C1}(ii), at mass $m$:
\begin{equation}\label{vqal:P6density}
 d v_{m,u}|_{\cC_B}
  =e^{-uJ_B-\alpha_m(uX^4)|B|}g_{\partial B}\,dP_m,
 \qquad g_{\partial B}>0,
\end{equation}
with the two recorded boundary-algebra readings and finite local Gaussian
$L^p$ densities. Also retain the exact DSS/SKQ/P*/H1-PREQ/H1/MIX/REG/
smoothed-CT/FLOW/TP/T interfaces, including the later source-normalization
correction; FLOW has its own cap $u\le m^2/8000$ and TP/T are fixed-square
transfers. Finally retain the raw positive-stage NR1/L0 interfaces,
\begin{equation}\label{vqal:P6raw}
 \delta_Z\le {d_D M_Z\over96N},\qquad
 \delta_A\le {d_D M_A\over96N},\qquad Z_Q\ge z,
\end{equation}
the full bilinear reference covariance matching, independent Gaussian
covariance-defect coupling, and physical common-space Wick bounds and
normalizers. L0 has its full ambient Euclidean operator-norm meaning
stated below. NR1's exact sharper author constant and larger isolated-proof
core keep their separate accepted scopes. Old globally $g_*$-qualified
CL1/FV/EB1 conclusions are not used as expanded-domain assumptions.
P6 includes neither full AHKZ-class uniqueness nor its averaged boundary
estimate \ref{prem:C2}(iii).

The accepted consequences used in the following proofs are precisely
these. AP1 supplies the finite original-law row bound $3/m^2$ and source
bound. AP2 transfers that bound to the actual classical free law for every
real compact smooth source $h$, supported anywhere in the plane, every
$t\in\R$, and every physical square:
\begin{equation}\label{vqal:AP2}
 0\le\log E_{\widehat\nu_B}e^{t\phi(h)}
          \le {3t^2\over2m^2}\|h\|_2^2.
\end{equation}
It also holds for the Gaussian $L^2$ source extension. AP3 gives
$0\le\alpha_m(uX^4-hX)-\alpha_m(uX^4)\le3h^2/(2m^2)$,
pressure differentiability and equality of the original FS states, with
their actual compact-joint free-volume limit. AP4 identifies the same
state as $(T_s)_*v_{M,u/s^2}$ for every $M>0$, $s=m/M$, using the
original dilation. AP5 gives the physical/classical joint-law identity,
the same-distribution realization $\nu=R_*v_{m,u}$, full
$\nu_n\Rightarrow\nu$ on $X$, fixed compact polynomial moment/UI
convergence, and non-Gaussianity from the positive local quartic density.
These AP1--AP5 premises have domain
\begin{equation}\label{vqal:inclusive}
                         0<u/m^2\le1/8000.
\end{equation}
They do not assume any $p_j$ moment conclusion. AP6 separately ports the
five actual-law error formulas of Appendix~\ref{qrs:supplement} to this
same target under P1--P6 and the \emph{strict} domain
\begin{equation}\label{vqal:strict}
                         0<u/m^2<1/8000.
\end{equation}
The original fixed $g_*$ remains numerically unknown; no statement here
deduces $10^{-10}<g_*$.

For target kernels and reconstruction we retain exactly REC2.a--f, as
applied in Appendix~\ref{qrs:supplement}: source-defined target,
Schwartz regularity, qualitative $L^2$ mixing before reconstruction,
half-plane positivity, tempered/local kernels and attributed reconstruction.
Its added FS Theorem~4.4 clustering premise is distinct from the old
restricted C1--C3,C6 route, whose mixing used C2(iii). Both OS M/P readings,
both boundary readings, local-density, half-plane, kernel/tube/boundary and
reconstruction premises remain literal. REC.g and quantitative uses of
AP6 remain strict. Hashes and internal reviews certify neither unavailable
primary classical bodies nor stronger source assertions.

\subsection{All local moments and the original common sequence at the cap}
\label{vqal:qualitative}

\begin{theorem}[MOM and the separately derived SEQ-HS finiteness]
\label{vqal:mom}
Assume the original field-space, fixed-volume and physical/classical
interfaces and AP2--AP5 under P1--P6, at \eqref{vqal:inclusive}.
Use the original $p_j$, modes $e_{j,k}$, weights $\lambda_{j,k}$ and
volume $V_j=4(j+1)^2$. For $s>1$ put
\begin{align}\label{vqal:hsconstants}
 h_{j,s}(v)^2&=\sum_k\lambda_{j,k}^{-s}
                 |v(\chi_j\overline e_{j,k})|^2,&
 c_j&=\|\chi_j\|_2^2/V_j,\nonumber\\
 S_{j,s}&=\sum_k\lambda_{j,k}^{-s}
 \le1+8((j+1)/\pi)^{2s}\left(1+{1\over2s-2}\right),&
 D_p&={6\over m^2}(2\Gamma(p+1))^{1/p}\quad(p\ge1).
\end{align}
For every fixed integer $j\ge1$, real $s>1$ and real $r>0$, with $p=\max(1,r/2)$,
\begin{equation}\label{vqal:hsmom}
 \sup_{\rho\in\{\nu_n:n\ge1\}\cup\{\nu\}}
       \rho(h_{j,s}^{r})\le(D_pc_jS_{j,s})^{r/2}<\infty.
\end{equation}
In the formal MOM scope $s=2$, $h_{j,2}=p_j$. Put
$B_j=c_jS_{j,2}$ and $C_{j,r}=(D_{\max(1,r/2)}B_j)^{r/2}$.
Then, for every real $p\ge1$ and $r>0$,
\begin{align}\label{vqal:momconvergence}
 \sup_n\nu_n(p_j^{2p})&\le2\Gamma(p+1)(6B_j/m^2)^p,&
 \nu_n(p_j^r)&\longrightarrow\nu(p_j^r),\nonumber\\
 \sup_\rho\rho(p_j^r1_{\{p_j^r>L\}})&\le C_{j,2r}/L,&
 \nu_n(H)&\longrightarrow\nu(H)
\end{align}
for every continuous complex $H$ with $|H(v)|\le c(1+p_j(v)^r)$.
All expectations are absolutely finite. Every early physical volume is
included. The all-$s>1$ assertion is SEQ-HS finiteness; moment convergence
at $s\ne2$ is not asserted, and $h_{j,s}$ need not be finite or continuous
on all of $X$ when $1<s<2$.
\end{theorem}

\begin{proof}
The exact scalar physical/classical law identity transfers
\eqref{vqal:AP2}, including its unbounded nonnegative exponential.
With $\sigma_h^2=3\|h\|_2^2/m^2$, Chernoff at $t=x/\sigma_h^2$ and
at the opposite sign gives
\begin{equation}\label{vqal:scalar-mom}
 \nu_n(|v(h)|>x)\le2e^{-x^2/(2\sigma_h^2)},\qquad
 \nu_n(|v(h)|^q)\le2(2\sigma_h^2)^{q/2}\Gamma(1+q/2).
\end{equation}
The latter follows by Tonelli and $y=x^2/(2\sigma_h^2)$ in the tail
integral for any $q>0$. A zero source is zero directly, without division.
In particular $\|v(h)\|_{2p}^2\le D_p\|h\|_2^2$.

The real and imaginary sources of $a_{j,k}=v(\chi_j\overline e_{j,k})$
are $V_j^{-1/2}\chi_j\cos(\pi k\cdot x/(j+1))$ and
$-V_j^{-1/2}\chi_j\sin(\pi k\cdot x/(j+1))$. Their squared $L^2$
norms sum to $c_j$ at every frequency, including the zero mode.
Minkowski on their two nonnegative squares gives
$\|a_{j,k}\|_{2p}^2\le D_pc_j$, without an independence assumption.
There are $8l$ modes on the $|k|_\infty=l$ shell; integral comparison of
$\sum l^{1-2s}$ proves \eqref{vqal:hsconstants}.

For the finite nonnegative sum $H_M=\sum_{|k|_\infty\le M}
\lambda_{j,k}^{-s}|a_{j,k}|^2$, finite Minkowski gives
$\nu_n(H_M^p)\le(D_pc_jS_{j,s})^p$. Increase $M$ by monotone
convergence to obtain the physical bound. To obtain the target bound,
keep $M$ finite first: $\min(L,H_M^p)$ is bounded continuous on $X$,
so AP5 gives the same upper bound at $\nu$. Increase $L$, then $M$.
This uses no continuity of the infinite $h_{j,s}$. Concavity at
$r/(2p)\le1$ gives \eqref{vqal:hsmom}. In particular the same laws are
concentrated on all real $H^{-s}_{\rm loc}$, $s>1$, by intersecting over
countably many rational $s>1$ and using $h_{j,s}\le h_{j,t}$ for $t<s$.

At $s=2$ the original Hilbert norm is continuous, since
$|p_j(v)-p_j(w)|\le p_j(v-w)$ and
$d(v,w)<2^{-j}\delta$ implies $p_j(v-w)<\delta$ for $0<\delta<1$.
All positive powers are continuous. The order-$2r$ bound gives the tail
inequality in \eqref{vqal:momconvergence}. Apply AP5 to
$\min(p_j^r,L)$, then let $L$ increase, to prove moment convergence.
For the complex-growth conclusion use radial truncation
$\pi_L(z)=z\min(1,L/|z|)$, $\pi_L(0)=0$.
The uniform second moment of $H$ is at most
$2c^2(1+C_{j,2r})$, so the truncation error is bounded by that quantity
divided by $L$. Bounded-continuous convergence and then $L\to\infty$
prove the last assertion. No physical subsequence or uniform-volume
Gaussian density bound has been used.
\end{proof}

\begin{theorem}[CAPSEQ: one unchanged sequence]
\label{vqal:seq}
Under the premises of Theorem~\ref{vqal:mom}, the original accepted
all-positive-$m,u$ fixed-volume interfaces and the full REC2.a--f
premises for target identification, all clauses (S),(0),(A)--(E) of
Theorem~\ref{thm:main} hold at \eqref{vqal:inclusive}, with the target
$\nu=R_*v_{m,u}$. The sequence is exactly the original four least-index
recursion of Proposition~\ref{prop:H1seq}, with
$a(k)=1/(2mk)$, $R(k)=k$, $\varepsilon(k)=1/k$, $N(k)=k$,
$\eta_n=2^{-n}/4$ and the unchanged
\begin{equation}\label{vqal:dist}
 \operatorname{dist}_n(\rho,\eta)
 =\beta_X(\rho,\eta)+\sum_{1\le j,r\le n}
                    |\rho(p_j^r)-\eta(p_j^r)|.
\end{equation}
It obeys, for every physical $n$, both architectures simultaneously,
\begin{align}\label{vqal:S}
 &0<a_n<\min(1/n,1/(2m)),\quad a_n<a_{n-1}/2,\nonumber\\
 &R_n\in\N,\quad R_n\ge\max(n,n\ell),\quad R_n>R_{n-1},\nonumber\\
 &0<\varepsilon_n<1/n,\quad\varepsilon_n<\varepsilon_{n-1}/2,
 \quad N_n\in\N,\quad N_n>n,\quad N_n>N_{n-1},\nonumber\\
 &\operatorname{dist}_n(P_n^b,\nu_n)<2^{-n},\qquad
 \Omega_n=[-R_n,R_n]^2,
\end{align}
with predecessor conditions omitted at $n=1$.

Every selected law is a Borel probability with positive finite normalizer
and every positive local moment. Both full sequences converge weakly on
$X$, their $p_j^r$ expectations converge for every real $r>0$, and their
joint compact-test correlations converge in $\sigma(\cD',\cD)$ at
every order, including nondecomposable joint tests. Every real/complex
polynomial of every later real compact smooth tuple has converging
expectations and uniformly bounded, uniformly integrable positive powers.
The deterministic tensors, correlation distributions, target restrictions
of tempered Schwinger kernels, and all constants $M_{j,r}$ are exactly
those in (C),(D) and Section~\ref{sec:K1}. Finite laws are reflection
positive across each central coordinate line, for bounded half-plane
observables and the full polynomial forms (E2), with both $L^2$ estimates
of that clause. These choices occur after the sequence, without
reselection, countability restriction on later tests, or a subsequence.

On the overlap with the original fixed $0<u/m^2<g_*$, every selected
index equals the original index. All clauses here preserve that original
sequence. The quantitative premise AP6 still requires \eqref{vqal:strict}.
\end{theorem}

\begin{proof}
The fixed-volume arrows in Section~\ref{sec:fixed} hold at every
$m,u>0$. Their moment arrows follow from the reference bounds and exact
normalizer domination in Section~\ref{sec:H1}: at the neural stage the
density bound is $M_{\rm full}/Z_N^b$ for every finite $N$, becoming
$2M_{\rm full}/Z_{\rm full}$ only eventually at fixed parameters.
Residual and physical bounds keep their fixed-volume dependence.
Truncate a continuous function of $p_j$ growth, apply the weak arrow,
and use a higher uniform moment to remove the truncation, as in
Lemma~\ref{lem:transfer}. Theorem~\ref{vqal:mom} supplies the previously
missing physical-volume seminorm-power arrow at the new cap.

At fixed $n$, each arrow therefore tends to zero in the full
$\operatorname{dist}_n$. In the original order, first choose the least
regulator index satisfying its discrepancy and size/predecessor rules;
then the least auxiliary integer; then the least positive reciprocal
stabilizer; finally the least width satisfying the discrepancy for
\emph{both} architectures. At each step the admissible integers contain
an entire tail. The intersection of the two width tails still contains
a tail, as do the finitely many extra size constraints. Thus every
least index exists. The four strict errors are each $<\eta_n$, and
their triangle gives \eqref{vqal:S} with the full moment sum.

Each early neural law has
$p_j(\iota F_N^b)\le\|\chi_j\|_2\sqrt N/2$ by Parseval and the
field bound. Its potential is bounded on the bounded field range in
finite $\Omega$, so its exact $Z_N^b$ is positive finite. For any fixed
integer $j,r$, the corresponding term of
$\operatorname{dist}_n(P_n^b,\nu_n)$ is $<2^{-n}$ once
$n\ge\max(j,r)$. Physical moment convergence and the finitely many
early laws give finite $M_{j,r}$ over both sequences and $\nu$.
Also $\beta_X(P_n^b,\nu)\le2^{-n}+\beta_X(\nu_n,\nu)\to0$.
The same continuous-growth truncation extends moments to every real
positive order. Compact duality $|v(f)|\le\|f\|_{H^2(\T_j)}p_j(v)$
then proves all later polynomial/power assertions, including complex,
constant, vanishing and dependent cases.

For completeness the correlation mechanism is deterministic and has no
coupling restriction. For $F\in\cD(((-j,j)^2)^k)$ its Fourier tensor
series is absolutely bounded by $A_{j,k}(F)p_j(v)^k$ through weighted
Cauchy--Schwarz. Moreover
\begin{equation}\label{vqal:tensorproof}
 |v^{\otimes k}(F)-w^{\otimes k}(F)|
 \le kA_{j,k}(F)p_j(v-w)\max(p_j(v),p_j(w))^{k-1}.
\end{equation}
The $k=1$ last factor is one. The product-density construction and
gluing in Section~\ref{sec:K1} give all arbitrary joint tests,
conjugations and reflected products, as well as
$A_{j,k}(F)\le3^kV_j^{k/2}\|F\|_{C^{2k}}$.
Uniform higher moments allow integration and continuous-growth transfer
of \eqref{vqal:tensorproof}. REC2's same-target tempered kernels agree
on compact products; product density gives (D) on all compact joint
tests. Compact evaluations generate $\operatorname{Borel}(X)$, so no
second target law is introduced.

The centered axis-parallel squares in \eqref{vqal:S} obey the original
E.P2/E.T5 hypotheses. Replacing the open physical square by its closure
changes no bounded-field potential integral. Hence bounded half-plane
RP holds at each stage across $\theta_1$ and $\theta_2$. For
$\cO_f=\sum_{r=0}^Jv^{\otimes r}(f_r)$, the two separate bounds are
\begin{align}\label{vqal:RPnorms}
 \|\cO_f\|_{L^2(\rho)}
 &\le|f_0|+\sum_{r=1}^JA_{j,r}(f_r)M_{j,2r}^{1/2},\nonumber\\
 \|\cO_f\circ\theta\|_{L^2(\rho)}
 &\le|f_0|+\sum_{r=1}^JC_j^r
                          A_{j,r}(f_r)M_{j+1,2r}^{1/2}.
\end{align}
The reflection estimate is the deterministic one of the original field
space; no invariance of $\rho$ is assumed. Half-plane product density
gives the required measurability. Apply bounded RP to the radial
truncation of $\cO_f$, then let the truncation increase using these
two $L^2$ bounds and Cauchy--Schwarz. Every reflection-form term is
absolutely integrable; the tensor identities identify it with the
corresponding compact correlation. Their finite sum converges to the
real nonnegative target form. At $J=0$ the form is $|f_0|^2$ directly.

Finally AP3/AP4 and compact-cylinder uniqueness identify the original
target on the old-domain overlap. The four predicates have the identical
laws, normalizers, enumerations, discrepancies, allocations and
predecessor rules there. Induction in $n$, and within each $n$ in
$a,R,\varepsilon,N$ order, makes the admissible sets and their least
members identical. No bound for a least index, decision procedure or
new value of $g_*$ is obtained.
\end{proof}

\subsection{A finite family with full ambient vector geometry}
\label{vqal:family}

For the numerical results fix $m=1$, $u=10^{-10}$, $\ell=1000$.
The strict gate follows from $8000<10^{10}$ through AP6; the old fixed
$g_*$ is not part of that numerical check. Define
\begin{equation}\label{vqal:bumpfamily}
 b(t)=\begin{cases}e^{-1/(1-t^2)},&|t|<1,\\0,&|t|\ge1,\end{cases}
 \qquad f_h(x,y)=b(x-h_1)b(y-h_2).
\end{equation}
The supported tuple is ordered, unweighted and has $1\le r\le4$.
Every center coordinate is the canonical dyadic $p/2^k\in[-8,8]$,
$0\le k\le16$: zero and integers have $k=0$, and $p$ is odd when
$k>0$. Decimal numerators have no plus sign, leading zero or negative
zero. Repetition preserves order and multiplicity. There are no input
amplitudes, signs, dilations, rotations or custom tests. Write
\begin{equation}\label{vqal:familynorms}
 K=\bigcup_i(h_i+[-1,1]^2),\quad R=1+\max_{i,j}|h_{ij}|\le9,
 \quad c_r=\lceil\sqrt r\rceil,\quad
 F_q=\left(\sum_i\|f_{h_i}\|_q^2\right)^{1/2}.
\end{equation}
The observation law $\mathcal L_f(\rho)$ is the pushforward by
$O_f(v)=(v(f_{h_i}))_i$ into the original ambient $\R^r$.
Write $\beta_X=\beta$ for the original field-law metric.
The observation distance is \eqref{qrs:BL}, with separate value and Euclidean
Lipschitz bounds. Let $\mathcal U_4(r)$ consist of real $C^4$ functions
on $\R^r$ with all full multilinear operator norms of orders $0$ through
$4$ at most one.

\begin{proposition}[SELECT geometry, source and covariance statements]
\label{vqal:selectgeometry}
Under the literal SELECT/AP2--AP6 and Gaussian/coupling premises,
$K\subset[-9,9]^2\subset\Lambda_1$ with open containment, and
\begin{equation}\label{vqal:normbounds}
 F_1\le4\sqrt r\le4c_r,\quad F_2\le2\sqrt r\le2c_r,
 \quad F=F_{4/3}\le3\sqrt r\le3c_r\le6.
\end{equation}
Use the identical chosen H1 coefficient $C_H<2^{16}$ and
$C_{H2}=2^{384}$ in \eqref{qrs:CL1-data}--\eqref{qrs:CL1-constant}.
At these physical parameters its literal coefficient satisfies
\begin{equation}\label{vqal:Cf}
 C_f(f;R)={\sqrt r(4R+1)\over5}C_f(f_0;1),\qquad
 C_f(f_0;1)\le2^{360},\quad
 C_f(f;R)\le\bar C_f:={2^{360}c_r(4R+1)\over5}<2^{364}.
\end{equation}
The separate spatial factor $e^{R/256}$ remains in $E_n$.

Let $H_{ij}=\langle f_{h_i},f_{h_j}\rangle_{L^2}$.
Writing $h_*(d)=\int b(t)b(t-d)dt$, one has
$H_{ij}=h_*(h_{i1}-h_{j1})h_*(h_{i2}-h_{j2})$ and
$0\le H\le4rI$. An entry is positive exactly when both coordinate
separations are $<2$, and is zero if either is $\ge2$.
Under every physical free law and the same target, for all $t\in\R^r$,
\begin{equation}\label{vqal:gramsources}
 1\le E e^{t\cdot O_f}\le e^{3t^THt/2},\qquad
 E O_f=0,\qquad 0\le\operatorname{Cov}(O_f)\le3H\le12rI.
\end{equation}
These are quadratic-form inequalities. They are not entrywise covariance
formulas for the interacting law, or source bounds for finite neural laws.

For the Gaussian observation covariances,
$0\le\Sigma_a\le\Sigma_0\le H$ and
$0\le\Sigma_{\rm end}\le a^2H$.
For a positive covariance defect $\Delta$ with
$\delta_\Omega(K)=\sup_{x\in K}\Delta(x,x)$, its observation covariance
$H_\Delta$ obeys
\begin{equation}\label{vqal:defect}
 |(H_\Delta)_{ij}|\le\|f_i\|_1\|f_j\|_1\delta_\Omega(K),\qquad
 \operatorname{tr}H_\Delta\le F_1^2\delta_\Omega(K).
\end{equation}
The actual Gaussian coupling therefore has mean observation displacement
at most $F_1\sqrt{\delta_\Omega(K)}$. The physical regulator coupling has
mean displacement at most $(a/2)F_2$. Repeated or dependent observations
are included with their original ambient dimension throughout.
\end{proposition}

\begin{proof}
Every derivative of the interior bump is a polynomial in
$(1-t^2)^{-1}$ times its exponential, which tends to zero at the two
endpoints. Thus the zero extension is smooth, has support $[-1,1]$ and
supremum $e^{-1}$. Translation preserves the Lebesgue norms; area four
and $|f_h|\le1$ give the individual bounds $4,2,4^{3/4}<3$.
Sum their squares over all coordinates. H\"older gives the full
observation operator bound $|O_f(v)|_2\le F\|v\|_4$.

In the literal CL1 expression all factors except
$(2R+1/2)\sqrt r e^{-2}$ are identical to the base bump's
$(5/2)e^{-2}$. Their ratio proves \eqref{vqal:Cf}, even if their common
factor vanishes. The upper uses $c_r\le2$, $R\le9$ and $74<80$.
This compares the displayed coefficients with identical H1/H2 choices;
it does not compare independently minimized constants.

Fubini gives the Gram formula. The overlap integrand is positive on an
open rectangle exactly in the stated positive case. For any $t$,
$t^THt=\|\sum t_if_i\|_2^2$, while $\operatorname{tr}H\le4r$.
Apply AP2 to this real compact source. The upper passes to the target
by bounded truncations of the positive exponential and weak convergence.
Sign symmetry and the resulting exponential moments give zero mean and
Jensen's lower bound. Differentiating the scalar MGF at zero in direction
$t$, justified by its two-sided exponential bound, gives the covariance
inequality. A zero Gram entry describes deterministic support geometry;
it supplies no interacting independence.

The physical Fourier multipliers satisfy
$0\le\widehat C_a\le\widehat C_0\le1$, proving their Loewner bounds
by Plancherel on $\sum t_if_i$. The positive endpoint kernel integrates
to $a^2$; Young's inequality proves its bound. Positivity of $\Delta$
gives $|\Delta(x,y)|\le\sqrt{\Delta(x,x)\Delta(y,y)}$ and hence
\eqref{vqal:defect}. Sum the diagonal and use Cauchy--Schwarz in the
actual coupling. The physical common-space bound is
$E|(\Phi_a-\Phi_0)(f_i)|^2\le(a^2/4)\|f_i\|_2^2$; summing gives the
last assertion. Identical centers give identical coordinates, possibly
singular laws; none of these proofs uses an inverse, density or rank.
\end{proof}

\subsection{Relative normalization and the sharper fixed-accuracy width}
\label{vqal:relative}

We now give the full argument behind NORM and VEC's change to the first
arrow. Its improvement comes from normalizing each replacement locally
through a common zero-amplitude partition. It uses the actual joint law
of background and replacement direction.

Take either of the following outer tuples, keeping their separate
observation domains:
\begin{equation}\label{vqal:outer}
 \begin{array}{c|c|c|c|c}
 \text{case}&n&p\ (a=2^{-p})&R_{\rm aux}&t\ (\varepsilon=2^{-t})\\\hline
 \text{scalar NORM}&74&16384&137439027473&68719476736\\
 \text{vector VEC}&75&16384&137439028473&68719476736
 \end{array}
\end{equation}
The scalar case has the single unshifted $f_0$; the vector case has the
whole typed family above, at accuracy $1/5$. In each case
$D=\Omega=[-R_{\rm aux},R_{\rm aux}]^2$,
$\Lambda=(-1000n,1000n)^2$, $A=|\Lambda|$, $O=d_D=|\Omega|$,
and $E_D=L^4(D;\R)$. Restrict the original neuron and endpoint Gaussian
to $D$. They are strongly measurable, Bochner centered, globally sign
symmetric, have covariance $C_{\rm end}$, and match in expectation for
\emph{every bounded bilinear form} on $E_D$. This last full matching is
the literal finite-reference premise; point-pair matching alone is not
substituted for it. Their point variances are $1/4$, and fourth moments
are at most $3/16$.

Let $V$ be the exact potential in Definition~\ref{def:potentials} on
$D$, with physical $c=c_a$, and write $w=e^{-V}$, $x=\|v\|_4$.
For independent neuron/Gaussian copies set
\begin{align}\label{vqal:joint}
 S_j&=N^{-1/2}\left(\sum_{i=1}^j\gamma_i+
                              \sum_{i=j+1}^{N-1}g_i^b\right),\nonumber\\
 v&=S_j+\theta h/\sqrt N,\quad Z_{j,h}(\theta)=Ew(v),\qquad
 d\Pi_{j,h,\theta}=w(v)\,d\Law(S_j,h)/Z_{j,h}(\theta),
 \quad 0\le j<N,\quad0\le\theta\le1.
\end{align}
The direction $h$ is independent of $S_j$ \emph{before} tilt and has
either neuron or Gaussian law. Under $\Pi$ it is generally dependent on
$v$. Using the $v_k,e_k$ of \eqref{qrs:width-polynomials}, with
$B=\kappa_a^2/2+6uc$, define
\begin{align}\label{vqal:relativeconstants}
 \alpha&=\min(u,\varepsilon),\quad U=\max(u,\varepsilon),
 &L&=B^2O/(2\alpha),\quad J_z=3UO/16+3uc^2A,\nonumber\\
 H&=L+J_z,\quad X=[2(H+4)/\alpha]^{1/4},
 &Y&=2d_D^{1/4}(H+10)^{1/2},\nonumber\\
 P_F(x)&=2e_4(x)+\sum_{q=1}^4\binom4q F^qe_{4-q}(x),
 &K_1&=v_1(X)Y,\nonumber\\
 N_0&=\max(1,\lceil16K_1^2\rceil),
 &K_4(F)&=P_F(X)Y^4/6.
\end{align}
The scalar $Y$ here bounds direction moments; $O_f$ denotes the
observation. All polynomials have nonnegative coefficients and
$\deg P_F\le12$. An upper for $c$ or $F$ can be used in these
majorants without changing the actual potential.

\begin{theorem}[NORM/VEC relative replacement at the specified stages]
\label{vqal:relative-theorem}
Assume their literal reference, Wick, finite-interface and strict AP6
premises. Every amplitude bridge in \eqref{vqal:joint} has positive
finite $Z$, point variance at most $1/4$, fourth moment at most $3/16$,
and
\begin{align}\label{vqal:jointbounds}
 w(v)&\le e^{L-\alpha x^4/2},\quad Z\ge e^{-J_z},
 &{d\Pi\over d\Law(S_j,h)}&\le e^H,\nonumber\\
 \Pi(x^4>z)&\le\min(1,e^{H-\alpha z/2}),
 &\|x\|_{L^q(\Pi)}\le X,\quad\|\|h\|_4\|_{L^q(\Pi)}
     &\le Y\quad(1\le q\le16).
\end{align}
Before tilt $E\exp(\|\gamma\|_4^2/(4\sqrt{d_D}))<4$.
The logarithmic partition has a $C^1$ extension past the endpoints and
\begin{equation}\label{vqal:ratios}
 (\log Z)'=-N^{-1/2}E_\Pi DV(v)[h],\quad
 |(\log Z)'|\le K_1/\sqrt N,
 \quad {Z_{j,h}(\theta)\over Z_{j,g}(1)}<2\quad(N\ge N_0).
\end{equation}
More precisely the neuron ratio is at most $e^{K_1/\sqrt N}$, and the
Gaussian ratio at most $e^{2K_1/\sqrt N}$; the zero-amplitude partitions
agree.

For $G\in\mathcal U_4(r)$, center at its normalized Gaussian endpoint
$\eta_j$ and put $\Phi_j(v)=[G(O_f(v))-\eta_j]w(v)$. Then
$|\eta_j|\le1$, $\Phi_j\in C_b^4(E_D)$ and
\begin{equation}\label{vqal:centered}
 \|D^4\Phi_j(v)\|_{\rm op}\le w(v)P_F(x),\qquad
 E_\Pi[P_F(x)\|h\|_4^4]\le P_F(X)Y^4.
\end{equation}
Each adjacent normalized replacement costs at most $K_4(F)/N^2$.
Consequently the actual neural and full Gaussian tilted laws satisfy,
for $N\ge N_0$,
\begin{align}\label{vqal:relativeBL}
 |P_N^b(G\circ O_f)-Q_\varepsilon(G\circ O_f)|
 &\le\min(2,K_4(F)/N),\nonumber\\
 d_{\rm BL}(\mathcal L_f(P_N^b),\mathcal L_f(Q_\varepsilon))
 &\le\min\left(2,{\sqrt6K_4(F)\over N\delta^3}
                                   +2\sqrt r\,\delta\right)
 \quad(0<\delta\le1),\nonumber\\
 d_{\rm BL}(\mathcal L_f(P_N^b),\mathcal L_f(Q_\varepsilon))
 &\le\min(2,D_rN^{-1/4}),\qquad
 D_r^4={2048\sqrt6\over27}r^{3/2}K_4(F).
\end{align}
This includes all singular and repeated observation tuples in the stated
family; their dimension factors remain.
\end{theorem}

\begin{proof}
At almost every point a bridge is a sum of independent centered
variables, with coefficients $1/\sqrt N$ and $\theta/\sqrt N$.
Their squared coefficient sum is $(N-1+\theta^2)/N\le1$.
Expanding the fourth moment leaves the single-index and two-pair terms,
bounded by $(3/16)(\sum a_i^2)^2\le3/16$. This includes $N=1$,
$\theta=0,1$ and both extreme $j$. Tonelli on bounded $D$ gives the
needed potential integrability. H\"older and completing the square give
$V(v)\ge\alpha x^4-B\sqrt O x^2\ge\alpha x^4/2-L$.
Dropping the negative quadratic only from an upper bound gives
$EV\le J_z$; Jensen gives $Z\ge e^{-J_z}$, hence the density and tail
bounds. No neural or full-tilt normalizer is assumed at least one.

If $T$ has tail $\le\min(1,e^{a-z})$ for $a\ge0$, its moments are
bounded by those of $a+E_1$ with a scalar unit exponential $E_1$.
Tail integration proves the comparison and $EE_1^k=k!$.
Apply it to $T=\alpha x^4/2$. Minkowski in $L^4$ gives
$\|T\|_4\le H+(4!)^{1/4}<H+4$, hence $\|x\|_{16}\le X$.
For a neuron $\|h\|_4=d_D^{1/4}/2$ deterministically, also after tilt.
For Gaussian $h$, spatial/probability Minkowski and its variance $1/4$
give $E\|h\|_4^{2k}\le d_D^{k/2}(k/2)^k$ for integers $k\ge2$;
at $k=1$, $E\|h\|_4^2\le\sqrt{3d_D}/4\le\sqrt{d_D}/2$.
The nonnegative exponential series and $k!\ge(k/e)^k$ give
\begin{equation}\label{vqal:gaussianexp}
 E e^{\|h\|_4^2/(4\sqrt{d_D})}
 \le\sum_{k\ge0}(e/8)^k<4.
\end{equation}
The actual joint density cap therefore yields the tilted tail
$\Pi(T_h>z)\le\min(1,e^{H+2-z})$ for
$T_h=\|h\|_4^2/(4\sqrt{d_D})$.
The same scalar comparison in $L^8$ gives
$\|T_h\|_8<H+10$, proving the $Y$ bound. Lyapunov gives every real
order through sixteen. Dependence after tilt is fully allowed.

Quartic H\"older gives the full Fr\'echet derivative bounds
$\|D^kV\|_{\rm op}\le v_k(x)$, $k\le4$.
The exponential chain rule gives $\|D^kw\|_{\rm op}\le w e_k(x)$,
with the exact fourth Bell coefficients $1,6,3,4,1$ in
\eqref{qrs:width-polynomials}. Each weighted derivative is globally
bounded: its majorant is $e^{L-\alpha x^4/2}$ times a finite polynomial,
and $\sup x^ke^{-\alpha x^4/2}=(k/(2\alpha e))^{k/4}$ for $k>0$.
Dominated differentiation under expectation is thus bounded by a
constant times $\|h\|_4/\sqrt N$, integrable before tilt.
It gives a continuous derivative on a real neighborhood of $[0,1]$.
Under the same joint $\Pi$, H\"older at orders four and two gives
$E_\Pi|DV(v)[h]|\le(4UX^3+2B\sqrt O X)Y=K_1$.
Integrate along the neuron branch, and along both branches through their
equal zero-amplitude partitions. Their costs are respectively one and
two copies of $K_1/\sqrt N$. For $N\ge N_0$ the latter exponential is
at most $\sqrt e<2$, proving \eqref{vqal:ratios}.

The observation operator bound is $F$. The fourth product rule, and
$|G-\eta_j|\le2$, give \eqref{vqal:centered}; the two $e_4$ copies
are essential. If $P_F(x)=\sum_{k=0}^{12}a_kx^k$, H\"older at total
order $k+4\le16$ gives $E_\Pi x^k\|h\|_4^4\le X^kY^4$,
including $k=0$ directly. This uses joint moments, never a product of
dependent marginal expectations.

Taylor expansion at $S_j$ through order three has fourth remainder
\begin{equation}\label{vqal:taylor}
 N^{-2}\int_0^1{(1-\theta)^3\over6}
 D^4\Phi_j(S_j+\theta h/\sqrt N)[h,h,h,h]d\theta.
\end{equation}
Globally bounded derivatives and finite untilted fourth moments justify
all conditional expectations and interchanges. Given $S_j$, the two
directions have matching constant terms, zero first and third terms by
Bochner centering/global sign symmetry, and matching second terms by
the full bounded-bilinear premise. Divide their difference by
$Z_{j,g}(1)$; Gaussian-endpoint centering makes this exactly the
difference of normalized expectations. Each of the two remainders uses
its own partition ratio $\le2$. Since the integral weight is $1/24$,
their sum is at most $4P_F(X)Y^4/(24N^2)=K_4(F)/N^2$.
Successive normalized endpoints agree, and the final normalized Gaussian
sum has the original Gaussian law. The $N$ terms telescope to $K_4/N$.

Literal L0 uses a standard Gaussian in the \emph{ambient} $\R^r$.
It gives $\|G-G_\delta\|_\infty\le\sqrt r\delta$ and full derivative
bounds through four at most $\sqrt6\delta^{-3}$ for $0<\delta\le1$.
Apply the smooth estimate to $G_\delta/(\sqrt6\delta^{-3})$ and keep
both smoothing errors. With $A_0=\sqrt6K_4$, $B_0=2\sqrt r$, the
interior minimizer is $(3A_0/(B_0N))^{1/4}$ when it is at most one.
Its value has fourth power $(2048\sqrt6/27)r^{3/2}K_4/N$.
Otherwise that expression exceeds $(4/3)B_0>2$ and the probability cap
suffices. Both cases retain $N\ge N_0$. No covariance is inverted.
\end{proof}

\begin{proposition}[Exact scalar and vector sufficient widths]
\label{vqal:widths}
At the separate stages \eqref{vqal:outer}, the accepted NORM and VEC
enclosures are
\begin{equation}\label{vqal:headline}
 \begin{array}{c|c|c|c|c}
 &N_0\text{ strictly below}&K_4\text{ strictly below}
 &N\text{ sufficient}&\text{target upper}\\\hline
 \text{NORM, scalar}&2^{274878071042}&2^{549756142097}
 &2^{549756142138}&65/4096\\
 \text{VEC, }1\le r\le4&2^{274878071042}&2^{549756142100}
 &2^{549756142144}&129/8192
 \end{array}
\end{equation}
Both target uppers are $<1/5$ for both architectures. The vector result
uses $n=75$ and all its ambient factors; it does not transplant the
scalar $n=74$ constant. At the identical vector outer stages the old
SELECT width $2^{2^{68719542354}}$ is sufficient with the same upper
$129/8192$. The comparison is between specified sufficient integers;
it states neither a necessary width nor an execution-time improvement.
\end{proposition}

\begin{proof}
Put $p=16384$, $t=2^{36}$, $l=t+4p+78$,
$x_*=t/2+p+21$, $y_*=t/2+2p+61$, $k_*=2t+5p+127$.
The variance bound in \eqref{qrs:variance} gives
$c\le\bar c<(p+3)/6\le p$ from $\pi>3$,
$\operatorname{asinh}(2^p)<\log(3\,2^p)<p+2$.
In either case $R_{\rm aux}<2^{38}$, $O<2^{78}$,
$6up<1$, $B\le2^{2p}$ and $\alpha=\varepsilon$, $U=u$.
Consequently
\begin{equation}\label{vqal:majorants}
 L<2^l,\quad J_z<2^{79}\le2^l,\quad H+10,H+4<2^{l+2},
 \quad X<2^{x_*},\quad Y<2^{y_*},\quad K_1<2^{k_*-2}.
\end{equation}
The strong last inequality gives
$16K_1^2<2^{2k_*}$, so its ceiling is at most that integer;
$N_0\le2^{2k_*}<2^{2k_*+4}=2^{274878071042}$.
This explicitly retains the ceiling step.

Set $W_0=2^{3x_*+5}$. All four $v_k(X)<W_0$:
$v_1(X)<2^{3x_*+1}$, the terms of $v_2$ have upper powers
$4+2x_*$ and $2p+40$, and $v_3,v_4$ have upper powers
$5+x_*$ and $5$. The Bell bounds are $1,W_0,2W_0^2,5W_0^3,15W_0^4$.
Thus $P_F(X)\le16W_0^4[1+(1+F)^4]$.
For scalar $F\le3$ its coefficient is $4112<2^{13}$; for vector
$F\le6$ it is $38432<2^{16}$. Multiplying by $Y^4/6$ and relaxing
the divisor upward gives respectively
\begin{equation}\label{vqal:k4exponents}
 K_4<2^{12x_*+4y_*+33}=2^{549756142097},\qquad
 K_4(F)<2^{12x_*+4y_*+36}=2^{549756142100}.
\end{equation}
Scalar smoothing at $\delta=2^{-10}$ with its width in
\eqref{vqal:headline} costs $<3\,2^{-11}+2^{-9}=7/2048<1/256$.
Vector smoothing at $\delta=2^{-11}$ costs
$<3\,2^{-11}+4\,2^{-11}=7/2048<1/256$.
The width threshold is satisfied in both cases.

Here is the outer-arrow check, retaining its actual volume and vector
factors. At $n=75$, $A=22500000000$, $uA=9/4$;
the scalar $n=74$ values are smaller. In \eqref{qrs:Nelson},
$K_\Lambda^2<14u^2A<1$, $K_\sharp<32$ and
$\beta>1/12$. Since $e^6>67591/175>384$, $H_2<2048$ and
$\sqrt{\bar\Xi_2}<2^{4097}$. The vector regulator term is therefore
$<2^{4097}(4\,2^{-8192}+2\,2^{-16384})\le1/256$;
the scalar norm is smaller. With $j_0=2^{34}$,
$J_0=9uA\bar c^2<(81/4)2^{28}$ and $e^{J_0}<2^{j_0}$, using
$\log2>1/2$. For auxiliary radius $s=2^{37}$ the strict geometric
margin is $R_{\rm aux}-1000n-s=1$ for both compacts.
Equation~\eqref{qrs:tail} gives
$\sqrt{\rho_a(s)}<2^{p-s/2}$. As $L_{\bar c}<32p^2$ and vector
$F_1\le8$, the full bracket is $<256p^2$. Thus
\begin{equation}\label{vqal:outererrors}
 A_{\rm aux}<2^{-51539591132}\le1/256,\qquad
 S<2^{-51539607443}\le1/256.
\end{equation}
The latter retains $B_{\rm ext}<O<2^{78}$ and the full common Wick
exponential. For the spatial tail, the degree-eight exponential partial
sum at $125/32$ is
$432528188408407789/8866461766385664>32$; hence $q<1/32$.
For vector $R\le9$, $e^{R/256}\le256/247<2$, and the $n=75$
bracket at $q=1/32$ is $149856/961<256$. Therefore
$E_{75}<2^{364-10+8-375}=2^{-13}$.
For scalar $C_f\le2^{360}$, $R=1$, the $n=74$ bracket is $<256$,
giving $E_{74}<2^{-12}$. Adding the four inner $1/256$ allocations
gives the two exact totals in \eqref{vqal:headline}.

For comparison, the old vector polynomial has coefficient sum
$P(1)\le38432V^4$ at $V=2^{2p+41}$, degree at most twelve, and
$b_k\le2^{3t+12}$. Its full coefficient obeys
$C_{\rm law}<2^{2^{l+2}+q_0+78}$ with
$q_0=4(2p+41)+3t+28=206158561472$.
The full exponential normalizer is retained in this inequality.
The ambient prefactor in $D_{\rm law}^4$ is $<2^{11}$, so
$N\ge2^{43}C_{\rm law}$ suffices for width error $\le1/256$.
Since $q_0+78+43<2^{38}<2^{l+2}$, the SELECT width with outer exponent
$l+4=68719542354$ satisfies that bound. The new vector exponent lies
strictly between $2^{39}$ and $2^{40}$ and is smaller than
$2^{68719542354}$. This proves the exact sufficient-width comparison.
\end{proof}

\subsection{Supported accuracy search with the old vector backend}
\label{vqal:accuracy}

ACC selects stages for the typed family in \eqref{vqal:bumpfamily},
using the \emph{old full-normalizer vector width coefficient}
\eqref{qrs:width-polynomials}--\eqref{qrs:width-coefficient}.
Its supported tolerance strings are exactly
\begin{equation}\label{vqal:tolerances}
          \{\texttt{1/5}\}\ \cup\ \{\texttt{2\^{}-q}:3\le q\le12\}.
\end{equation}
The sharper VEC backend above is a separate fixed-stage, fixed-$1/5$
result. This section does not apply it to ACC's changing stages.

For a positive rational $x$, write $\operatorname{clog}_2(x)$ for the
least nonnegative integer $h$ with $x\le2^h$. The following gives the
complete deterministic selection rule; all displayed operations are
integer or rational operations. Let
\begin{align}\label{vqal:acc-spatial}
 d&=\min\{i\ge0:5\,2^{-i}\le\tau\},\qquad\delta=2^{-d},\nonumber\\
 \bar E(n)&={4\over10000}\bar C_f{256\over256-R}\,2^{-5n}
                   \left[{(2n+1)32\over31}+{64\over961}\right],\nonumber\\
 n&=\min\{1\le i\le80:\bar E(i)\le\delta\},\quad
 A=4000000n^2,\quad U_A=uA,\quad b_N^0={3\over16U_A},\nonumber\\
 L_{10}(k)&=\sum_{j=0}^{10}{k^j\over j!},\quad
 k=\min\{1\le i\le7:L_{10}(i)>32/b_N^0\},\nonumber\\
 H_N&=1+\left\lceil\max(12U_A,4,4k^2/b_N^0)\right\rceil.
\end{align}
Here $b_N^0$ is a Nelson lower bound, distinct from the width polynomial
weight $b_0=1$. Next select
\begin{align}\label{vqal:acc-outer}
 p&=\min\{i\ge4:\ i\text{ even},\quad
       2H_N+3-i/2\le-d-1,\quad2H_N+2-i\le-d-1\},\nonumber\\
 a&=2^{-p},\quad j_e=\lceil18U_Ap^2\rceil,\quad
 h_A=\operatorname{clog}_2(64U_Ap^2+4c_r),\nonumber\\
 h_s&=\min\{i\ge1:j_e+h_A+p-2^{i-1}\le-d\},\quad s=2^{h_s},\nonumber\\
 R_{\rm aux}&=1000n+s+1,\quad O=4R_{\rm aux}^2,
 \quad o=\operatorname{clog}_2(O),\nonumber\\
 h_S&=\operatorname{clog}_2(6Op^2),\quad
 t=\max(34,j_e+h_S+d),\quad\varepsilon=2^{-t}.
\end{align}
Finally put
\begin{align}\label{vqal:acc-width}
 h_R&=\operatorname{clog}_2(2R_{\rm aux}),\quad v_e=2p+h_R+2,\nonumber\\
 q_0&=4v_e+3t+28,\quad l=t+4p+o,
 \quad H_{\rm small}=q_0+o+11+4d.
\end{align}
Require $\operatorname{bitlength}(H_{\rm small})\le l+2$.
Start $w=l+2$ and increase until this guard holds and $w>l+2$;
successful output has $w=l+3$. Set $N=2^{2^w}$. The powers defining
$a,\varepsilon,N$ are inert mathematical expressions; the program
does not expand their denominators or widths. These minima refer to the
displayed sufficient predicates, rather than canonical CAPSEQ least indices.

\begin{theorem}[ACC termination and accuracy at its literal domain]
\label{vqal:acc}
Under the accepted SELECT/OPT/EB1/EC1/AP6 and P1--P6 premises, every
typed tuple and every tolerance in \eqref{vqal:tolerances} has the
selection above. It satisfies $n\le80$, $k\le7$, $H_N<4096$,
even $p<20000$, $h_s\le40$, $t<2^{40}$ and $w<2^{40}$.
All selected watched integer/rational witnesses fit the stated
$1024$-bit ceiling and fewer than $20000$ logical guard steps suffice
for the specified arithmetic/source-inventory loops. These are finite
operation and named-witness assertions, rather than hardware-time or
process-memory bounds. The actual adjacent-law errors satisfy
\begin{equation}\label{vqal:acc-result}
 E_n\le\bar E(n)\le\delta,\quad
 R_{\rm phys},A_{\rm aux},S,W\le\delta,\qquad
 d_{\rm BL}(\mathcal L_f(P^b),\mathcal L_f(\nu))
       \le\bar E(n)+4\delta\le5\delta\le\tau.
\end{equation}
The same selected stages work for both architectures and all the
supported repeated, overlapping, touching, separated and singular tuples.
\end{theorem}

\begin{proof}
At $\tau=1/5$ the allocation gives $d=5$; at $\tau=2^{-q}$ it gives
$d=q+3$, so $5\le d\le15$. The exponential partial sum used in the
preceding proof gives $q_{\rm spatial}<1/32$.
Also $e^{R/256}\le256/(256-R)<2$. Positivity in the shell formula
therefore proves $E_n\le\bar E(n)$. At $n=80$ its bracket is
$159776/961<256$, giving $\bar E(80)<2^{364-10+8-400}=2^{-38}$.
Thus the box search stops, with $1000n\ge R$.

Here $U_A=n^2/2500\le64/25<3$ and
$K_\Lambda^2<14u^2A<1$. Consequently
$K_\sharp\le\max(12U_A,4)$ and the actual Nelson
$\beta=\pi^2/(48U_A)>b_N^0>1/16$.
Since $L_{10}(7)>512>32/b_N^0$, the $k$ search succeeds.
For its chosen $k$, $e^k>L_{10}(k)>32/b_N^0$ bounds the logarithm
in $H_2$ by $k$. Thus $H_2<H_N\le3137<4096$ and
$\sqrt{\bar\Xi_2}<2^{2H_N+1}$.
The full vector regulator expression is then
$<2^{2H_N+3-p/2}+2^{2H_N+2-p}$.
The two predicates bound each term by $2^{-d-1}$.
The even value $4H_N+2d+8$ suffices, so $p\le12586<20000$.
As above $c_a\le\bar c<p$ and $6up<1$.

The complete Wick exponential obeys
$e^{9U_A\bar c^2}<2^{j_e}$, using $\log2>1/2$;
$L_{\bar c}<32p^2$. Both required compact sets have strict auxiliary
margin $R_{\rm aux}-1000n-s=1$, so the defect tail gives
$\sqrt{\rho_a(s)}<2^{p-s/2}$.
The full auxiliary expression is therefore
$<2^{j_e+h_A+p-s/2}\le2^{-d}$.
For termination, $j_e\le21600000000<2^{35}$,
$64U_Ap^2+4c_r<2^{37}$, and
$j_e+h_A+p+d<2^{35}$. Hence $h_s\le36$ already suffices, and
$R_{\rm aux}<2^{37}$, $O<2^{76}$, $o\le76$.
As $6Op^2<2^{109}$, $h_S\le109$ and $t<2^{36}$.
The choice $t\ge34$ gives $\varepsilon<u$; retaining
$B_{\rm ext}<O$, the stabilizer is
$<2^{j_e+h_S-t}\le2^{-d}$.

For the old width backend use $B_N<2^{2p}$ and
$\sqrt O=2R_{\rm aux}\le2^{h_R}$. Its $v_1,v_2$ coefficient sums
are at most $2^{v_e}$; those of $v_3,v_4$ are smaller.
The Bell sums are bounded by $1,V,2V^2,5V^3,15V^4$ for $V=2^{v_e}$.
Keeping the extra pure $e_4$ and all five observation terms gives
$P(1)\le38432V^4<2^{16}V^4$.
For each degree $1\le j\le12$, the scalar weight
$(j/(2\alpha e))^{j/4}<2^{3t+12}$; $b_0=1$ obeys the same upper.
Thus the full $Q_{\rm poly}<2^{q_0}$.

The normalizer exponent has all three positive terms:
$L_N<2^l$, $3uO/16\le2^{l-1}$ and
$3U_A\bar c^2\le2^{l-1}$. The two latter comparisons are the
explicit guards $o\le l-1$ and
$\operatorname{clog}_2(3U_Ap^2)\le l-1$.
They hold since $l\ge51$, while $3U_Ap^2<2^{32}$.
Therefore
\begin{equation}\label{vqal:acc-normalizer}
 C_{\rm law}<2^{2^{l+2}+q_0+o},\qquad
 D_{\rm law}^4<2^{11}C_{\rm law}.
\end{equation}
Here $(2048\sqrt6/27)r^{3/2}<2^{11}$ retains ambient $r\le4$.
It is enough to have $N\ge2^{11+4d}C_{\rm law}$.
The proven bounds give $0<H_{\rm small}<2^{38}$ and
$\operatorname{bitlength}(H_{\rm small})\le38\le l+2$.
Thus the width search succeeds at $w=l+3$, since
\begin{equation}\label{vqal:acc-symbolic}
 \log_2(2^{11+4d}C_{\rm law})
       <2^{l+2}+H_{\rm small}<2^{l+3}=\log_2N.
\end{equation}
Also $l<2^{37}$ and $w<2^{40}$. This proves the width allocation with
the full old normalizer and polynomial. The triangle on the same five
actual observation laws proves \eqref{vqal:acc-result}.

For the stated finite-resource theorem, there are at most $16,80,7,6292,36,2$
allocation, box, Nelson, regulator, auxiliary and width predicates,
respectively. Each $\operatorname{clog}_2$ loop has at most $110$
iterations; fixed checks add fewer than $100$ guards and the fixed
source/core inventory adds $203$. The resulting bound is below $8000$,
hence below $20000$. Center denominators are at most $2^{16}$;
$\bar C_f$ has numerator below $384$ bits and denominator at most
$5\,2^{16}$. The factors in $\bar E(n)$ have denominator bit bounds
$19,25,10,14,401$, with the last from $2^{5n}$, so even their unreduced
product is below $469$ bits and its numerator below $450$ bits.
Adding $4\delta$ keeps the final components below $512$ bits. The fixed
Nelson partial sums use fewer than $128$ bits, the other named selected
quantities fewer than $256$, and the largest guarded logarithm shift at
most $110$ bits. The coefficient power itself uses $365$ bits.
These prove the named $1024$-bit sufficiency, with no assertion about
every uninstrumented intermediate or the work of hashing and I/O.
\end{proof}

For the four-corner request $(-8,-8),(-8,8),(8,-8),(8,8)$ and
$\tau=2^{-12}$, the accepted selected arithmetic records
\begin{align}\label{vqal:acc-example}
 &n=75,\quad p=6954,\quad s=4294967296,\quad
 t=1958503808,\quad w=1958531694,\nonumber\\
 &\bar E(n)+4\delta={919201379\over6076595200000}<{1\over4096}.
\end{align}
This is the historical ACC tuple, with $N=2^{2^{1958531694}}$. Its old
backend remains unchanged. The later same-stage sharper width is in
\eqref{sva:corner-example}; VEC retains its separate fixed-$1/5$ stages.

\subsection{Actual same-law linkage and what selection computes}
\label{vqal:linkage}

Write $C_n^b=P_n^b$ for the canonical law of
Theorem~\ref{vqal:seq}. For an independently selected admissible inner
tuple $\vartheta=(a,R_{\rm aux},\varepsilon,N)$ at the same physical
box, let $F_n^c(\vartheta)$ be its neural law, and set
$I_n(\vartheta)=W+S+A_{\rm aux}+R_{\rm phys}$.
The architecture choices $b,c$ are independent. The physical label $n$
specifies $\Lambda_n$; it is distinct from a width and from the four
enumeration indices in the least-index recursion.

\begin{theorem}[LINK metric, same-box comparison and physical-label modulus]
\label{vqal:link}
For a real compact smooth tuple $f$ choose $J$ so all supports are
compactly contained in $(-J,J)^2$, and define
\begin{equation}\label{vqal:L}
 A_{f,J}=\left(\sum_i\|f_i\|_{H^2(\T_J)}^2\right)^{1/2},\qquad
 L_{f,J}=2^J\max(2,A_{f,J}).
\end{equation}
For all Borel probabilities on $X$,
\begin{equation}\label{vqal:pushforward}
 d_{\rm BL}(\mathcal L_f(\rho),\mathcal L_f(\eta))
       \le\min(2,L_{f,J}\beta_X(\rho,\eta))
       \le\min(2,L_{f,J}\operatorname{dist}_n(\rho,\eta))
\end{equation}
when the latter is defined. The first inequality needs no moment
assumption or $J\le n$. Empty/all-zero tuples have distance zero.
There is no reverse bound tending to zero with this one finite-observation
discrepancy for either field metric, even among moment-finite smooth-field
point masses.

Under the literal CAPSEQ and strict AP6 premises, for $n\ge n_0$ and
$n\ell\ge R$, all target symbols below denote the identical original
FS state version $\nu$:
\begin{align}\label{vqal:linkbounds}
 d_{\rm BL}(\mathcal L_f(C_n^b),\mathcal L_f(\nu))
       &\le\min(2,L_{f,J}2^{-n}+E_n),\nonumber\\
 d_{\rm BL}(\mathcal L_f(F_n^c(\vartheta)),\mathcal L_f(\nu))
       &\le\min(2,I_n(\vartheta)+E_n),\nonumber\\
 d_{\rm BL}(\mathcal L_f(C_n^b),\mathcal L_f(F_n^c(\vartheta)))
       &\le\min(2,L_{f,J}2^{-n}+I_n(\vartheta)).
\end{align}
Under the full REC2 interpretation premises this is also its Euclidean
law. At equality in the qualitative cap, CAPSEQ and REC2.a--f remain
available but this AP6/REC.g quantitative use is refused.

With the additional certified computable positive $m,u,\ell$, an
admissible supplied $n_0$, a computable support enclosure $R$, finite
computable uppers $\bar C_f\ge C_f$ and $\bar A\ge A_{f,J}$, and
positive rational $\tau$, one can compute a sufficient physical label
$n_\tau$ so that every $n\ge n_\tau$ has
\begin{equation}\label{vqal:physicalmodulus}
 \bar L2^{-n}+4\bar C_fu\ell^2e^{\kappa R}
 q^n\left[{2n+1\over1-q}+{2q\over(1-q)^2}\right]<\tau,
 \quad\bar L=2^J\max(2,\bar A),\quad q=e^{-m\ell/256}.
\end{equation}
This output is a sufficient box label, without a claim of leastness;
it does not compute $a_n,R_n,\varepsilon_n,N_n$ or their least indices.
\end{theorem}

\begin{proof}
Duality in the original torus convention gives
$|O_f(v)-O_f(w)|_2\le A_{f,J}p_J(v-w)$.
For a unit BL function the difference after composition is at most
$\min(2,A_{f,J}p_J(v-w))$, which is at most
$\max(2,A_{f,J})\min(1,p_J(v-w))\le L_{f,J}d(v,w)$.
Thus $(G\circ O_f)/L_{f,J}$ is in the unit BL class on $X$.
The definition of $\beta_X$ and probability diameter two prove
\eqref{vqal:pushforward}; nonnegative moment discrepancies give its
last inequality for every $n$.

For the exact reverse obstruction choose $r+1$ nonzero smooth bumps
with disjoint compact supports inside $(-1,1)^2$. The linear map from
their coefficients to their $r$ integrals against the tests has a
nonzero kernel. Its combination $h$ is nonzero, smooth compact,
$O_f(h)=0$ and $p_1(h)>0$. For
$\rho=\delta_0$, $\eta_M=\delta_{Mh}$ all local moments are finite
and the observed laws coincide, while
\begin{equation}\label{vqal:nonreverse}
 \beta_X(\delta_0,\delta_{Mh})=d(0,Mh)>0,\qquad
 \operatorname{dist}_n(\delta_0,\delta_{Mh})\ge Mp_1(h)
 \quad(n\ge1).
\end{equation}
The equality uses the admissible test $v\mapsto d(v,0)\in[0,1]$.
These point masses demonstrate a general metric non-implication; they
are not asserted to be physical-model laws or to prove undecidability.

AP3--AP5 identify the physical/classical finite joint laws and their
target with those in CAPSEQ. Compact evaluations determine Borel
probabilities on $X$, so the AP6 target is the same law. CAPSEQ gives
$\beta_X(C_n^b,\nu_n)<2^{-n}$; apply the pushforward bound and the
AP6 spatial arrow. For the selected tuple, sum its four inner arrows
to $\nu_n$, then its last arrow to $\nu$. For the last inequality
of \eqref{vqal:linkbounds}, insert $\mathcal L_f(\nu_n)$ directly
between the two neural laws. Its spatial-to-target error is absent
because this path ends at the same finite physical law. No upper
bounds have been subtracted. The old $g_*$ is neither estimated nor
used to enlarge an old globally qualified theorem.

For the computable modulus choose certified rational $q\le\bar q<1$,
$L_0\ge\bar L$, $B_0\ge4\bar C_fu\ell^2e^{\kappa R}$.
Strict positive bounds for $\ell$ and an upper for $R$ give a safe
integer $n_{\rm start}\ge n_0$ with $n_{\rm start}\ell\ge R$.
The geometric identity
$T_n(z)=\sum_{j=n}^\infty(2j+1)z^j$ shows it is increasing in $z$,
decreasing in $n$ and tends to zero for $0<z<1$.
Search exact rational $L_0\,2^{-n}+B_0T_n(\bar q)<\tau$ from this
starting label. It terminates, and monotonicity proves the assertion
for every subsequent label. No equality of computable reals has to be
decided. The canonical law at that label remains the existential
least-index law; the label calculation does not evaluate its inner
parameters or field discrepancy.
\end{proof}

\begin{corollary}[Literal bump at physical label 74]
\label{vqal:linkbump}
For the unshifted scalar bump at the fixed physical parameters, $J=2$
is admissible and $A_{f,2}\le374$, $L_{f,2}\le1496$. Write
$P_{\rm OPT}^c$ and $P_{\rm RB4}^c$ for the two separate accepted
scalar tuples at $\Lambda_{74}$, with their original and relative
widths. For all independent architecture choices $b,c$,
\begin{align}\label{vqal:concretelink}
 d_{\rm BL}(\mathcal L_f(C_{74}^b),\mathcal L_f(\nu))
       &\le1/4096+1496\,2^{-74}<1/2048,\nonumber\\
 d_{\rm BL}(\mathcal L_f(P_{\rm OPT/RB4}^c),\mathcal L_f(\nu))
       &\le65/4096<1/5,\nonumber\\
 d_{\rm BL}(\mathcal L_f(C_{74}^b),\mathcal L_f(P_{\rm OPT/RB4}^c))
       &\le1/64+1496\,2^{-74}<1/32,\nonumber\\
 d_{\rm BL}(\mathcal L_f(P_{\rm OPT}^b),\mathcal L_f(P_{\rm RB4}^c))
       &\le1/128.
\end{align}
No concrete inner tuple here is claimed to equal the canonical record.
\end{corollary}

\begin{proof}
For $r_t=(1-t^2)^{-1}$ direct differentiation gives
$b''=(4t^2r_t^4-2r_t^2-8t^2r_t^3)e^{-r_t}$.
The maxima of $r^ke^{-r}$ give
$\|b''\|_\infty\le8/e^2+216/e^3+1024/e^4<93$.
On the area-four support, $\|f\|_2\le2$ and each pure second
derivative has $L^2$ norm at most $186$.
The exact Fourier convention yields
$\|f\|_{H^2(\T_2)}=\|(1-\Delta)f\|_2\le374$, and hence
$L_{f,2}\le4\,374=1496$. The supports are compactly contained in
$(-2,2)^2$, as required.
The scalar tuples share all four outer parameters. Each first-arrow
error is at most $1/256$, with RB4 satisfying its $N\ge N_0$
threshold by \eqref{vqal:headline}; the other three inner errors are
at most $1/256$ each, and $E_{74}\le1/4096$.
Apply \eqref{vqal:linkbounds}. Since $1496<2^{11}$, its small term
is $<2^{-63}$, proving both strict bounds.
For the last comparison insert the identical actual
$Q_\varepsilon$ between the two laws and add just their two
first-arrow errors. This remains valid for different architectures.
\end{proof}

\subsection{Comparison, software domains and proof navigation}
\label{vqal:navigation}

Table~\ref{vqal:comparison} keeps each certificate at its literal
domain. A smaller upper bound is a statement about these sufficient
rules. All widths remain astronomical process-neuron counts with
countably many internal spatial variables.

{\small
\begin{longtable}{@{}p{0.13\textwidth}p{0.20\textwidth}p{0.17\textwidth}p{0.25\textwidth}p{0.13\textwidth}@{}}
\caption{Exact sufficient-bound comparisons at $m=1$, $u=10^{-10}$,
$\ell=1000$. The NS1 and OPT rows retain their baseline scope.
ACC's example is one supported request.}\label{vqal:comparison}\\
\toprule
Result & Observations / accuracy & Outer stages & Sufficient width & Target upper\\
\midrule\endfirsthead
\toprule Result & Observations / accuracy & Outer stages & Sufficient width & Target upper\\
\midrule\endhead
NS1 & scalar base / $1/5$ & $n=512$, $p=2^{24}$ &
$2^{2^{4611686018494496904}}$ & $1/64+2^{-736}$\\
OPT & scalar base / $1/5$ & $n=74$, $p=16384$ &
$2^{2^{68719542354}}$ & $65/4096$\\
NORM & scalar base / $1/5$ & same OPT outer tuple &
$2^{549756142138}$ & $65/4096$\\
SELECT & typed $1$--$4$ / $1/5$ & $n=75$, $p=16384$ &
$2^{2^{68719542354}}$ & $129/8192$\\
VEC & same typed family / $1/5$ & same SELECT outer tuple &
$2^{549756142144}$ & $129/8192$\\
ACC & four corners / $2^{-12}$ & $n=75$, $p=6954$ &
$2^{2^{1958531694}}$ & $\frac{919201379}{6076595200000}$\\
\bottomrule
\end{longtable}}

\paragraph{Supported and refused software domains.}
The NORM executable is a fixed scalar-case verifier. SELECT and the
separately versioned VEC executable support only the canonical ordered
dyadic bump family at exactly $1/5$; ACC supports that family and exactly
the eleven tolerance strings in \eqref{vqal:tolerances}. Their strict
schemas refuse altered model data, custom/weighted/signed tests,
noncanonical dyadics, unsupported coordinates/counts or tolerances,
floating/nonfinite values, duplicate keys and out-of-contract expression
trees. Geometry, multiplicity, coefficient enclosures and stages are
derived rather than accepted from caller assertions. Required source
bytes and current decisions are hash-bound; a relocated root must supply
those identical bytes. Missing source, changed source, failed witnesses
and resource exhaustion supply no successful numerical certificate.
This is an account of the accepted software contract, not a claim that
the integration author reran their original full source trees.

VEC binds $174$ scientific source files; ACC binds $195$. Those original
source inventories and their recorded selected arithmetic/replays remain
in the frozen companion packets. The present author verified all $189$
assigned input files, including every one of the seven statement/proof/
decision triples. Source inventories inside a packet are evidence of its
original run, rather than a claim that all their historical paths are
present in this smaller integration packet. Complete unavailable primary
classical bodies remain attributed premises.

The independently implemented ACC arithmetic verifier does not call its
stage search; it shares the declared strict parser/source/guard utilities.
The accepted repaired implementation checks the entire success envelope
and refuses false candidate/availability labels. A well-typed historical
timing record is metadata, without authenticated historical measurement.
Its cooperative core timing excludes request parsing, includes
verification-certificate parsing, and excludes serialization and the
whole process. Named watched values and logical guards provide no
all-intermediate, memory, hostile-filesystem-atomicity or hard latency
guarantee. Default ACC ceilings are $1024$ watched bits, $30000$ core
milliseconds, $20000$ guards, $n\le80$, $p\le20000$, $h_s\le40$,
and $t,w\le2^{40}$. Stricter ceilings may stop valid requests.
Its six outcomes distinguish supported success, unsupported input,
numerical source gap, source-binding refusal, resource exhaustion and
refused certificate; every non-success has no certificate. VEC has
corresponding invalid-input, source-mismatch, missing-source,
inequality-failure and resource-limit refusals. Reviewed code, manifests
and external freezes remain trust roots, rather than a formal proof kernel.

The accepted ACC decision records $110$ actual search/verify pairs,
maximum $4081$ search guards, $307$ verifier guards and $393$ watched
bits. These are observed finite-run facts, distinct from the theorem's
supported-domain termination and limits. The sharper vector and variable
accuracy backends have not been silently combined.

\paragraph{Retained adverse evidence and boundaries.}
The original MOM actual-match overall MISMATCH is retained: its mathematics
matched, while two external-openness metadata flags were separately
repaired and fresh MOMMETAFIX resolved that defect. The old SELECTNUM
disagreement record, $314/315$ checks, and its binding ancillary manifest
refusal-label correction remain. ACC's original hostile GAP, float
serialization and outer-envelope gaps remain with their fixes; the later
three-file timing repair changed the contract/code wording while its
proof, statement, schema, tests and $195$ source records stayed unchanged.
Earlier rejected REC-v1, witnesses and CLOSED gate, partial/incomplete
versions, source-access exceptions, the rational resource-inconclusive
route, and all old author/reviewer scope distinctions remain historical
evidence. They are neither erased nor substituted for the current decisions.

The Gaussian four-corner obstruction is still useful: the product neuron
at $(0,0),(a,0),(0,a),(a,a)$ has fourth product $1/16$, while its
Gaussian covariance-pairing sum is $(2e^{-2}+e^{-4})/16<1/16$.
The excess persists as $1/N$ for the iid normalized neuron sum.
Thus Gaussian joint-moment domination of finite neural laws cannot
replace the joint coercivity/density argument above.

The established neural topology is weak $X$ convergence and compact
smooth joint-correlation convergence in $\sigma(\cD',\cD)$.
Target Schwartz regularity does not supply global neural Schwartz or
strong-dual control. The effective-selection obstruction consists of the
missing all-field/moment discrepancy predicates and leastness decisions:
one finite observation law cannot determine them by
\eqref{vqal:nonreverse}. This is a limitation of the proved implications,
not a model-specific impossibility theorem. Full AHKZ-class uniqueness,
the old boundary estimate, practical/necessary/minimal complexity,
sampling/training, other models/dimensions, historical priority and
public research release remain separate. Finite-stage RP is exactly
across the two central coordinate lines in the centered axis-parallel
geometry. The target's full symmetry, arbitrary-line half-plane
positivity and quantum reconstruction retain their separate classical
premises and attribution.

\paragraph{How to audit this revision.}
The public \href{run:PROOF-MAP.md}{proof map} gives this edition's source
block locations and hashes. The original private audit retains the exact
frozen statement/proof line ranges, packet hashes and baseline reuse. The full new proofs are in this appendix and the earlier
paper sections it explicitly invokes; the private source packets give
additional original details and provenance. Mathematical reasoning is
not replaced by an operational ledger. The baseline bibliography is
unchanged. All $277$ original mathematical environments, including
nested displays, are byte-identical, and a reversible source delta
recovers the exact baseline after removing the two added text regions
and restoring its date. The preservation audit, compilation/reference
checks and all-page render review are separate local evidence. The
current integration's fresh source/quantifier/editorial review,
actual-proof match and coordinator decision remain required.

\clearpage
\section{Sharper relative accuracy at the selected stages}
\label{sva:current}

The accepted SVA-ACC-RELATIVE-v1 mathematics carries the relative estimate
through all supported accuracy-selected stages. The separately accepted
SVC-relative-accuracy-v1 implementation realizes the symbolic rule.
This appendix states their precise content and mathematical reasoning;
it preserves the earlier results as dated checkpoints. Its mathematical
implications are conditional on the literal source and finite-law
interfaces below. The integrated presentation awaits its own fresh review.

\subsection{Domain, notation and the actual unchanged laws}
\label{sva:domain}

Fix $m=1$, $u=10^{-10}$, $\ell=1000$, $K_{\rm cov}=1$ and the original
coefficient-one, same-mass, own-variance Wick quartic, with no factorial.
The field state space is the original complete $X=H^{-2}_{\rm loc}$ and
$\nu^X=\nu$ is the same Fr\"ohlich--Simon target of
Section~\ref{vqal:premises}. All P1--P6/AP, Gaussian/Wick,
DSS/BD/BBD/Griffiths--Simon/SK-q/P*, H1/MIX/REG/smoothed-CT/FLOW/TP/T,
source/classical/coefficient and NR1 qualifications there retain their
original scopes. Numerical target identification uses strict AP6,
$0<u/m^2<1/8000$. The original fixed $g_*$ and $10^{-10}<g_*$ remain
numerically unknown. The sharper NR1 author result and larger isolated
blind-proof core remain distinct; excluded NR2 E4/r0 is not used.

The bump and ordered observations are exactly
\eqref{vqal:bumpfamily}: $b(x)=\exp[-1/(1-x^2)]$ for $|x|<1$, zero
otherwise, and $f_h(z)=b(z_1-h_1)b(z_2-h_2)$. There are $r\in\{1,2,3,4\}$
centers. Each coordinate is a canonical $P/2^k\in[-8,8]$,
$0\le k\le16$, with $k=0$ for integers and zero, and odd numerator
when $k>0$. Decimal numerators have no plus, leading zero or negative
zero. The tolerance strings are exactly \eqref{vqal:tolerances}.
Additional observation amplitudes or signed coefficients, dilations and
rotations are absent. Order and multiplicity are retained. Put
\begin{align}\label{sva:geometry}
 K&=\bigcup_i[h_{i1}-1,h_{i1}+1]\times[h_{i2}-1,h_{i2}+1],
 &R&=1+\max_{i,j}|h_{ij}|,\quad c_r=\lceil\sqrt r\rceil,\nonumber\\
 F_q&=\left(\sum_i\|f_{h_i}\|_q^2\right)^{1/2},
 &F&=F_{4/3},\quad\bar F=3c_r,\nonumber\\
 \bar C_f&=2^{360}c_r(4R+1)/5,
 &\eta&=2^{-d}.
\end{align}
The full ambient observation is $O_f(v)=(v(f_{h_i}))_{i=1}^r$.
The metric is the supremum over real $G$ with $\|G\|_\infty\le1$
and Euclidean Lipschitz constant at most one separately on $\R^r$.
The class $\mathcal U_4(r)$ uses full Euclidean multilinear operator
norms through order four, all at most one, including absolute value at
order zero. H\"older gives $\|O_f\|\le F$.

Use exactly the outer minima \eqref{vqal:acc-spatial}--\eqref{vqal:acc-outer},
with their $\delta$ renamed $\eta$ here. Thus $H_N$ is $H_{\rm Nel}$,
$j_e$ is $j_{\exp}$, $h_A$ is $h_{\rm aux}$ and $h_S$ is
$h_{\rm stab}$. All of $n,p,a,s,R_{\rm aux},O,o,t,\varepsilon$ keep
their original meanings. Define
$\Lambda=(-1000n,1000n)^2$, $D=\Omega=[-R_{\rm aux},R_{\rm aux}]^2$,
$A=4000000n^2$ and $E=L^4(D;\R)$. Here $O=|D|$, not the observation
operator. The accepted SVA proof, applied to this unchanged ACC selection,
gives the following whole-domain bounds; their derivation is recalled
in the proof of Theorem~\ref{sva:width} below:
\begin{align}\label{sva:stagebounds}
 5&\le d\le15,&1&\le n\le80,& U_A=uA&\le64/25<3,\nonumber\\
 k&\le7,&H_N&\le3137,&4\le p&\le12586<20000\quad(p\text{ even}),\nonumber\\
 h_s&\le36,&R_{\rm aux}&<2^{37},& O&<2^{76},\quad o\le76,\nonumber\\
 h_{\rm stab}&\le109,&34\le t&<2^{36},&\varepsilon&<u,\quad6up<1.
\end{align}
Also $1\le R\le9$, $K\subset(-1000,1000)^2$ openly,
$F_1\le4c_r\le8$, $F_2\le2c_r\le4$, $F\le\bar F\le6$.
With identical chosen H1 and H2 coefficients, the literal CL1 formula
\eqref{qrs:CL1-constant} gives
$C_f(\text{tuple};R)=\sqrt r(4R+1)C_f(\text{base};1)/5\le\bar C_f<2^{364}$;
its separate factor $\exp(R/256)$ remains separate. The source flow
variable in that formula is distinct from the auxiliary radius $s$.

\paragraph{Full selected-stage reference premise.}
At every selected $(a,D)$ assume the original EB1/NR1 finite-reference
interface in its complete form. The physical multiplier and variance are
\begin{equation}\label{sva:physicalvariance}
 {1\over1+w^2+q^2+a^2w^2q^2},\qquad
 c=c_a={1\over2\pi}\int_0^\infty
 {dq\over\sqrt{(1+q^2)(1+a^2q^2)}},\qquad \kappa_a^2=a^{-2}-1,
\end{equation}
with Fourier measure $dw\,dq/(2\pi)^2$ and
$c\le\bar c=[\operatorname{asinh}(1/a)+1]/(2\pi)<p$.
The endpoint $\gamma$ is a centered, strongly measurable Gaussian
$E$-valued random element with zero Bochner mean, global sign symmetry
and covariance $e^{-|z-z'|_1/a}/4$. For each original architecture,
$g^b|_D$ is strongly measurable, Bochner centered, globally sign
symmetric, with $|g^b(z)|=1/2$ a.e. almost surely and that covariance.
The product neuron uses two independent stationary symmetric telegraph
processes of rate $1/(2a)$ and their product divided by two. The Mondrian
neuron uses the projectively consistent whole-plane axis-parallel
partition at budget $1/a$, independent uniform cell signs, half-valued
cells and the original lower-left half-open representative. A rectangle
of sides $w,h$ splits at rate $w+h$, vertically with probability
$w/(w+h)$, at a uniform position, and daughters evolve independently for
the remaining budget. Original existence, consistency and measurability
are premises; boundary changes agree as $E$ elements. For \emph{every}
bounded bilinear $B_0:E\times E\to\R$ assume
\begin{equation}\label{sva:bilinear}
 E B_0(g^b,g^b)=E B_0(\gamma,\gamma).
\end{equation}
Independent copies exist on product spaces. This full assumption is
separate from point covariance. It is assumed at selected stages rather
than inferred by applying the fixed-$n=75$ VEC theorem elsewhere.

The actual potential and full normalized laws are
\begin{align}\label{sva:actuallaws}
 V(v)&=u\int_\Lambda(v^4-6cv^2+3c^2)
       +\varepsilon\int_{\Omega\setminus\Lambda}v^4
       -{\kappa_a^2\over2}\int_\Omega v^2,\quad w(v)=e^{-V(v)},\nonumber\\
 F_N^b&=N^{-1/2}\sum_{i=1}^N g_i^b,\quad
 P_N^b={w\,\Law(F_N^b)\over Ew(F_N^b)},\qquad
 Q_\varepsilon={w\,\Law(\gamma)\over Ew(\gamma)}.
\end{align}
These are restrictions of the original whole-plane laws, not a new
reference model. The actual $V$ always retains $c_a$. All downstream
laws and their common Gaussian/Wick couplings in
\eqref{qrs:chain}--\eqref{qrs:physical-errors} remain literal premises:
$Q_\varepsilon\to Q_0\to Q_a\to Q_{\rm phys}\to\nu^X$ has error bounds
$S,A_{\rm aux},R_{\rm phys},E_n$. The physical Wick normalizer
$Z_a=E e^{-U_{a,\Lambda}}\ge1$ uses $EU_{a,\Lambda}=0$ on its physical
Gaussian space. It is separate from each full normalizer in
\eqref{sva:actuallaws}. AP5/AP6 identifies the same physical-to-target
limit. Neither source hashes nor this restatement reprove these inputs.

Literal ambient L0 uses a standard Gaussian $\xi$ on the original
$\R^r$: $G_\zeta(y)=E G(y-\zeta\xi)$ has approximation error
$\sqrt r\zeta$, absolute value and first derivative bounded by one,
and $\|D^qG_\zeta\|_{\rm op}\le\sqrt{(q-1)!}\zeta^{1-q}$ for
$q\ge2$. For $0<\zeta\le1$, all four norms are bounded by
$\sqrt6\zeta^{-3}$. Singular covariances retain this ambient dimension.

\subsection{Relative comparison under that interface}
\label{sva:relative}

For any $c_*\ge c$ used only in positive majorants define
\begin{align}\label{sva:constants}
 \alpha&=\varepsilon,\quad U=u,\quad B=\kappa_a^2/2+6uc_*,
 &L&={B^2O\over2\alpha},\quad J_z={3UO\over16}+3uc_*^2A,
 \quad H=L+J_z,\nonumber\\
 v_1(x)&=4Ux^3+2B\sqrt O x,&v_2(x)&=12Ux^2+2B\sqrt O,\nonumber\\
 v_3(x)&=24Ux,&v_4(x)&=24U,\nonumber\\
 e_0&=1,\quad e_1=v_1,\quad e_2=v_1^2+v_2,
 &e_3&=v_1^3+3v_1v_2+v_3,\nonumber\\
 e_4&=v_1^4+6v_1^2v_2+3v_2^2+4v_1v_3+v_4,
 &P_F(x)&=2e_4(x)+\sum_{q=1}^4\binom4q F^qe_{4-q}(x),\nonumber\\
 X_{\rm rel}&=[2(H+4)/\alpha]^{1/4},
 &Y_{\rm rel}&=2O^{1/4}(H+10)^{1/2},\nonumber\\
 K_1&=v_1(X_{\rm rel})Y_{\rm rel},
 &N_0&=\max(1,\lceil16K_1^2\rceil),\quad
 K_4(F)=P_F(X_{\rm rel})Y_{\rm rel}^4/6.
\end{align}
The scalar moment bounds $X_{\rm rel},Y_{\rm rel}$ are distinct from
$X=H^{-2}_{\rm loc}$ and any physical coupling fields. All polynomial
coefficients are nonnegative and $\deg P_F\le12$.

\begin{theorem}[SVA.a--c at every selected stage]\label{sva:relative-theorem}
Under the complete premises above, fix either architecture, integer
$N\ge1$ and $0\le j<N$. Form $S_j$, $v$, $Z_{j,h}(\theta)$ and the
actual joint law $\Pi_{j,h,\theta}$ exactly as in \eqref{vqal:joint},
with $h$ independent before tilt and of neuron or Gaussian law. Then the
untilted point variance is at most $1/4$, its fourth moment at most
$3/16$, every $Z$ is positive finite, and, for $z\ge0$,
\begin{align}\label{sva:moments}
 w(v)&\le e^{L-\alpha\|v\|_4^4/2},\quad Z\ge e^{-J_z},
 &{d\Pi\over d\Law(S_j,h)}&\le e^H,\nonumber\\
 \Pi(\|v\|_4^4>z)&\le\min(1,e^{H-\alpha z/2}),
 &\|\|v\|_4\|_{L^q(\Pi)}&\le X_{\rm rel},\quad
 \|\|h\|_4\|_{L^q(\Pi)}\le Y_{\rm rel}\quad(1\le q\le16).
\end{align}
Before tilt $E\exp(\|\gamma\|_4^2/(4\sqrt O))<4$.
Each $Z$ and $\log Z$ extends as $C^1$ past $[0,1]$ with
\begin{align}\label{sva:partitionratios}
 (\log Z)'&=-N^{-1/2}E_\Pi DV(v)[h],
 &| (\log Z)'|&\le K_1/\sqrt N,\nonumber\\
 {Z_{j,g}(\theta)\over Z_{j,g}(1)}&\le e^{K_1/\sqrt N},
 &{Z_{j,\gamma}(\theta)\over Z_{j,g}(1)}&\le e^{2K_1/\sqrt N}<2
 \quad(N\ge N_0).
\end{align}
The zero-amplitude partitions agree. All derivatives of $w$ through
four are globally bounded, with $\|D^kV(v)\|\le v_k(\|v\|_4)$ for
$1\le k\le4$ and $\|D^kw(v)\|\le w(v)e_k(\|v\|_4)$ for $0\le k\le4$.

For $G\in\mathcal U_4(r)$, let $m_j$ be its actual normalized Gaussian
endpoint expectation and put $\Phi_j=(G\circ O_f-m_j)w$. Then
$|m_j|\le1$, $\Phi_j\in C_b^4(E)$ and
\begin{equation}\label{sva:derivative}
 \|D^4\Phi_j(v)\|_{\rm op}\le w(v)P_F(\|v\|_4),\qquad
 E_\Pi[P_F(\|v\|_4)\|h\|_4^4]\le P_F(X_{\rm rel})Y_{\rm rel}^4.
\end{equation}
For every integer $N\ge N_0$, each adjacent normalized replacement has
error at most $K_4(F)/N^2$ and the actual laws satisfy
\begin{align}\label{sva:BL}
 |P_N^b(G\circ O_f)-Q_\varepsilon(G\circ O_f)|
 &\le\min(2,K_4(F)/N),\nonumber\\
 d_{\rm BL}(\mathcal L_f(P_N^b),\mathcal L_f(Q_\varepsilon))
 &\le\min\left(2,{\sqrt6K_4(F)\over N\zeta^3}
                         +2\sqrt r\zeta\right)\quad(0<\zeta\le1),\nonumber\\
 d_{\rm BL}(\mathcal L_f(P_N^b),\mathcal L_f(Q_\varepsilon))
 &\le\min(2,D_rN^{-1/4}),\qquad
 D_r^4={2048\sqrt6\over27}r^{3/2}K_4(F).
\end{align}
All supported repetitions, singularities and support geometries retain
order, multiplicity and ambient dimension.
\end{theorem}

\begin{proof}
We carry the argument of Section~\ref{vqal:relative} to these stages
using the explicitly assumed full interface, not its fixed-stage theorem.
At a.e. point the $N$ independent centered summands have variance $1/4$
and fourth moments at most $3/16$. Their squared coefficient sum is
$(N-1+\theta^2)/N\le1$. Expansion of the fourth power yields
$E(\sum\beta_i\xi_i)^4\le(3/16)(\sum\beta_i^2)^2$; only single-index
and two-pair terms survive. Tonelli gives integrability on finite $D$.
H\"older and completion of the square give
$V(v)\ge\alpha x^4-B\sqrt O x^2\ge\alpha x^4/2-L$, $x=\|v\|_4$.
Dropping negative quadratic terms only in an upper bound gives
$EV\le J_z$; Jensen proves $Z\ge e^{-J_z}$. The resulting actual
joint density is bounded by $e^{H-\alpha x^4/2}$, proving its tail.

For a nonnegative $T$ with tail at most $\min(1,e^{a-z})$, $a\ge0$,
tail integration bounds its positive integer moments by those of
$a+E_1$ for a unit exponential scalar; Minkowski gives
$\|T\|_k\le a+(k!)^{1/k}$. At $T=\alpha x^4/2$, $k=4$ this gives
$\|x\|_{16}\le X_{\rm rel}$. For Gaussian directions, the
spatial/probability Minkowski inequality gives
$E\|\gamma\|_4^{2k}\le O^{k/2}(k/2)^k$ for $k\ge2$; at $k=1$,
$E\|\gamma\|_4^2\le\sqrt{3O}/4\le\sqrt O/2$.
The positive series and $k!\ge(k/e)^k$ give
$E\exp(\|\gamma\|_4^2/(4\sqrt O))\le\sum_{k\ge0}(e/8)^k<4$.
The actual joint density cap therefore bounds the tilted tail of
$T_\gamma=\|\gamma\|_4^2/(4\sqrt O)$ by
$\min(1,e^{H+2-z})$. At $k=8$, $(8!)^{1/8}<8$ gives
$\|T_\gamma\|_8<H+10$, hence the $Y_{\rm rel}$ bound.
For a neuron, $\|g^b\|_4=O^{1/4}/2$ deterministically, even after tilt.
Lyapunov gives every real order through sixteen. No post-tilt
independence is used.

Quartic H\"older gives the full Fr\'echet $v_k$ bounds. The exponential
partition rule gives the Bell polynomials in \eqref{sva:constants},
including coefficients $1,6,3,4,1$ at order four. A finite polynomial
times $e^{L-\alpha x^4/2}$ is bounded globally; for $k>0$ its monomial
maximum is $(k/(2\alpha e))^{k/4}$. Thus $w$ has bounded derivatives
through four. Differentiating $Z$ is dominated by a constant times
$\|h\|_4/\sqrt N$, integrable before tilt, on a real neighborhood of
the amplitude interval. H\"older under the same actual $\Pi$ gives
$E_\Pi|DV(v)[h]|\le(4UX_{\rm rel}^3+2B\sqrt O X_{\rm rel})Y_{\rm rel}=K_1$.
Integrating the logarithmic derivative along one branch, or along both
through their common zero-amplitude partition, proves the two ratios.
For $N\ge N_0$, $2K_1/\sqrt N\le1/2$, so the second ratio is at most
$\sqrt e<2$.

The fourth product rule uses $|G-m_j|\le2$ and
$\|D^q(G\circ O_f)\|\le F^q$, giving exactly $P_F$ with its second
pure $e_4$ copy and all observation binomials. Its nonnegative monomial
$x^k$ has $k\le12$. H\"older at total order $k+4\le16$ gives
$E_\Pi x^k\|h\|_4^4\le X_{\rm rel}^kY_{\rm rel}^4$; at $k=0$ use
the fourth direction moment directly. This proves the mixed-moment
bound without multiplying dependent marginal expectations.

Taylor through order three has the remainder \eqref{vqal:taylor}.
Conditioning before tilt on $S_j$, zero Bochner means cancel first
terms, global sign symmetry cancels third terms, and the full
bounded-bilinear premise cancels second terms for
$D^2\Phi_j(S_j)$. Strong measurability, separability and bounded
derivatives justify evaluations, expectations and interchanges. Centering
at $m_j$ makes the difference divided by $Z_{j,g}(1)$ exactly the
normalized endpoint difference. Each of its two remainders has partition
ratio at most two, and the integral weight is $1/24$. Their sum is at
most $4P_F(X_{\rm rel})Y_{\rm rel}^4/(24N^2)=K_4(F)/N^2$.
The $N$ actual normalized mixtures telescope. The final normalized sum
of independent Gaussian $E$-elements has exactly the original Gaussian
law, since continuous-linear-functional laws determine Borel laws on
separable $E$. The endpoints are the actual $P_N^b,Q_\varepsilon$.

Ambient L0 contributes two errors $\sqrt r\zeta$ and rescales its
smooth test by $\sqrt6\zeta^{-3}$, proving \eqref{sva:BL}.
For $A_0=\sqrt6K_4$, $B_0=2\sqrt r$, the interior minimizer is
$(3A_0/(B_0N))^{1/4}$. Its value has fourth-power coefficient
$(2048\sqrt6/27)r^{3/2}K_4$. When this minimizer exceeds one, that
value exceeds $(4/3)B_0\ge8/3>2$, and the probability cap suffices.
Both cases retain $N\ge N_0$; no covariance inverse or rank reduction
occurs.
\end{proof}

\subsection{Complete logarithmic enclosure and same-stage width}
\label{sva:enclosure}

\begin{theorem}[SVA.d--f: all supported requests]\label{sva:width}
At each unchanged selected stage, set $c_*=p$ only in the majorants and
\begin{align}\label{sva:exponents}
 T&=\lceil t/2\rceil,&x_e&=T+p+21,&y_e&=T+2p+61,\nonumber\\
 k_1&=4T+5p+127,&n_0&=8T+10p+258,&k_4&=16T+20p+532,\nonumber\\
 \zeta&=2^{-(d+3)},&e_{\rm new}&=k_4+4d+12
              =16T+20p+4d+544,&N_{\rm new}&=2^{e_{\rm new}}.
\end{align}
Here $T$ is a scalar half-exponent, distinct from any source transfer
interface or field. Then
\begin{align}\label{sva:enclosure-result}
 X_{\rm rel}&<2^{x_e},\quad Y_{\rm rel}<2^{y_e},
 &K_1&<2^{k_1-2},\quad16K_1^2<2^{2k_1},\nonumber\\
 N_0&\le2^{2k_1}<2^{n_0},
 &K_4(F)&\le K_4(\bar F)<2^{k_4},\quad N_{\rm new}>N_0,\nonumber\\
 d_{\rm BL}(\mathcal L_f(P_{N_{\rm new}}^b),
                \mathcal L_f(Q_\varepsilon))&<7\eta/8<\eta.
\end{align}
The four other errors remain $E_n\le\bar E(n)\le\eta$ and
$R_{\rm phys},A_{\rm aux},S\le\eta$ at the same stages. Thus
\begin{equation}\label{sva:target}
 d_{\rm BL}(\mathcal L_f(P_{N_{\rm new}}^b),\mathcal L_f(\nu^X))
 \le\bar E(n)+4\eta\le5\eta\le\tau.
\end{equation}
For the identical old comparison \eqref{vqal:acc-width},
$w_{\rm old}=l+3$, and
\begin{equation}\label{sva:order}
 e_{\rm new}<2^{40}<2^{w_{\rm old}},\qquad
 N_{\rm new}<N_{\rm old}=2^{2^{w_{\rm old}}}.
\end{equation}
These compare sufficient mathematical whole-plane process-neuron
integers at identical model, tests, architectures, target, outer stages
and error allocation.
\end{theorem}

\begin{proof}
The ACC predicates give \eqref{sva:stagebounds} analytically throughout
the domain. To display the outer dependencies, $q=\exp(-125/32)<1/32$
follows from the positive degree-eight exponential partial sum
$432528188408407789/8866461766385664>32$. The shell bracket increases
in $q$; $\exp(R/256)\le256/(256-R)<2$ gives $E_n\le\bar E(n)$.
At $n=80$ the bracket is $159776/961<256$ and
$\bar E(80)<2^{-38}\le\eta$. The actual Nelson bounds give
$K_\Lambda^2<14u^2A<1$, $K_\sharp\le\max(12U_A,4)$,
$b_{\rm Nelson}>3/(16U_A)>1/16$ and
$\log_+(64/(e b_{\rm Nelson}))<k\le7$, so $H_2<H_N\le3137$.
Hence $\sqrt{\bar\Xi_2}<2^{2H_N+1}$ and
\begin{equation}\label{sva:regulator}
 R_{\rm phys}<2^{2H_N+3-p/2}+2^{2H_N+2-p}\le\eta.
\end{equation}
The even candidate $4H_N+2d+8$ passes both predicates, giving the bound
on $p$. The variance estimate follows from
$\operatorname{asinh}(2^p)<\log(3\,2^p)<p+2$ and $\pi>3$.

For completeness the original physical Wick exponential and full
vector bracket satisfy $\exp(9U_A\bar c^2)<2^{j_e}$ and
$2U_A L_{\bar c}+F_1<64U_Ap^2+4c_r\le2^{h_A}$.
The strict margin $R_{\rm aux}-\max(1000n,R)-s=1$ includes the
open $s$-neighborhoods of both required compacts. The covariance-defect
bound gives $\sqrt{\rho_a(s)}<2^{p-s/2}$. Thus
\begin{equation}\label{sva:remaining-errors}
 A_{\rm aux}<2^{j_e+h_A+p-s/2}\le\eta,\qquad
 S<2^{j_e+h_{\rm stab}-t}\le\eta.
\end{equation}
The existence bounds $j_e\le21600000000<2^{35}$,
$64U_Ap^2+4c_r<2^{37}$ and $j_e+h_A+p+d<2^{35}$ imply $h_s\le36$.
They yield $O<2^{76}$, $h_{\rm stab}\le109$, $t<2^{36}$.
The condition $t\ge34$ ensures $\varepsilon<u$. Every area,
variance, geometry and stabilizer dependence used next is retained.

Relax $O<2^{76}$ to $O<2^{78}$ and $\sqrt O<2^{39}$. Since
$\kappa_a^2=2^{2p}-1$ and $6up<1$, $B<2^{2p-1}+1\le2^{2p}$.
Put $l_c=t+4p+78\ge128$. Directly from \eqref{sva:constants},
\begin{align}\label{sva:H-enclosure}
 L&<2^{t+4p+77}=2^{l_c-1},\nonumber\\
 J_z&<2^{78}+9p^2<2^{78}+2^{32}<2^{79}\le2^{l_c},\nonumber\\
 H&<2^{l_c+1},\qquad H+4,H+10<2^{l_c+2},\nonumber\\
 X_{\rm rel}^4&<2^{2t+4p+81},\qquad
 Y_{\rm rel}^4<2^{2t+8p+242}.
\end{align}
For both parities of $t$, $4T\ge2t$, and the two fourth-power slacks are
$4T-2t+3\ge3$ and $4T-2t+2\ge2$. Therefore
$X_{\rm rel}<2^{x_e}$ and $Y_{\rm rel}<2^{y_e}$. Odd $t$ rounds
upwards; no $\lfloor t/2\rfloor$ is used.

As $4u<1$ and $2p+40\le2x_e$,
\begin{align}\label{sva:K1-enclosure}
 v_1(X_{\rm rel})&<2^{3x_e}+2^{2p+40+x_e}<2^{3x_e+1},\nonumber\\
 K_1&<2^{3x_e+y_e+1}=2^{4T+5p+125}=2^{k_1-2},\nonumber\\
 16K_1^2&<2^{2k_1},\qquad
 N_0\le2^{2k_1}<2^{2k_1+4}=2^{n_0}.
\end{align}
The right side $2^{2k_1}$ is an integer. A real number strictly below
an integer has ceiling at most that integer; the same-exponent ceiling
comparison is weak, followed by a strict four-bit slack. This handles
$N_0=\max(1,\lceil16K_1^2\rceil)$ exactly.

Let $W_0=2^{3x_e+5}\ge1$. All four $v_j(X_{\rm rel})<W_0$.
For $v_2$ its two terms have upper powers $2x_e+4$ and
$2p+40\le2x_e$; for $v_3,v_4$ use $x_e+5$ and $5$.
The Bell bounds are $1,W_0,2W_0^2,5W_0^3,15W_0^4$, each at most
$16W_0^4$. Retaining the centered second $e_4$, every binomial term
and $\bar F\le6$ gives
\begin{align}\label{sva:Bell-enclosure}
 P_F(X_{\rm rel})&\le P_{\bar F}(X_{\rm rel})
 \le16W_0^4\left[2+\sum_{q=1}^4\binom4q6^q\right]\nonumber\\
 &=38432W_0^4<2^{12x_e+36},\nonumber\\
 K_4(F)&\le K_4(\bar F)<2^{12x_e+4y_e+36}
                         =2^{16T+20p+532}=2^{k_4}.
\end{align}
Here $38432=16(1+7^4)<2^{16}$; the scalar coefficient $4112$ is not
used. We only enlarge the upper by dropping the divisor six in $K_4$.
These whole-domain analytic enclosures retain the actual tilt, volume
and variance dependence; finite stage lists do not prove them.

Now $e_{\rm new}-n_0=8T+10p+4d+286>0$, so $N_{\rm new}>N_0$.
With $\sqrt6<3$ and $\sqrt r\le2$, the two terms in \eqref{sva:BL}
obey
\begin{equation}\label{sva:allocation}
 {\sqrt6K_4(F)\over N_{\rm new}\zeta^3}<3\eta/8,
 \qquad2\sqrt r\zeta\le\eta/2.
\end{equation}
Their sum is $<7\eta/8<\eta$. The four unchanged later actual-law
arrows and the triangle inequality give \eqref{sva:target}. Strict
AP6 identifies the same $\nu^X$; an unknown $g_*$ inequality is not
inserted into that step.

For the old guard, $h_R\le38$, $p<20000$, $t<2^{36}$, $o\le76$
give $v_e<40040$ and $0<H_{\rm small}<3\,2^{36}+160335<2^{38}$.
As $l=t+4p+o\ge50$, its bit-length guard holds. The first old candidate
$l+2$ fails its explicit strict predicate; the next succeeds, hence
$w_{\rm old}=l+3$. This retains the original old full-normalizer bound
as its inherited accepted premise. Since $T\le2^{35}$,
$e_{\rm new}<2^{39}+400604<2^{40}$, while
$w_{\rm old}\ge53>40$. This proves \eqref{sva:order} at identical
outer stages and allocations. It proves no necessary or minimal cost,
measured speedup or effective common-sequence least-index rule.
\end{proof}

\subsection{Exact examples and the separately versioned implementation}
\label{sva:examples}

For four \emph{distinct} corner centers $(-8,-8),(-8,8),(8,-8),(8,8)$
and tolerance $2^{-12}$ the accepted example is
\begin{align}\label{sva:corner-example}
 d&=15,\quad n=75,\quad p=6954,\quad s=4294967296,\nonumber\\
 t&=1958503808,\quad T=979251904,\quad
 e_{\rm new}=15668170148,\quad w_{\rm old}=1958531694,\nonumber\\
 \bar E(n)+4\eta&={919201379\over6076595200000}<{1\over4096},
 \qquad N_{\rm new}=2^{15668170148}<2^{2^{1958531694}}=N_{\rm old}.
\end{align}
The historical tuple $(0,0),(8,8),(0,0),(-8,-8)$ has three distinct
centers and four ordered slots, including a repeated observation. Its
same radius and ambient dimension give the same selected numerical
values, while the ordered tuple, request bytes and semantic identity
remain distinct.
For the single unshifted bump at $1/5$, the accepted odd-$t$ example is
\begin{equation}\label{sva:odd-example}
 n=73,\quad p=6574,\quad t=1658202793,\quad T=829101397,
 \quad e_{\rm new}=13265754396,\quad w_{\rm old}=1658229159.
\end{equation}
Here $d=5$ and the upward rounding in \eqref{sva:exponents} matters.
These examples use their own selected stages; they are distinct from the
older fixed-$1/5$ NORM/VEC rows in Table~\ref{vqal:comparison}.

\begin{proposition}[Accepted SVC-relative-accuracy-v1 software contract]
\label{sva:software}
The exact accepted SVC version implements Theorems~\ref{sva:relative-theorem}
and~\ref{sva:width} under their conditional premises and its bounded
input/source/resource contract. A successful request retains the exact
ordered tuple and unchanged outer minima, returns the derived geometry,
source/Nelson/regulator/auxiliary/stabilizer and rational five-error
witnesses, and encodes $N$ as the inert node
$\operatorname{pow2}(\operatorname{integer}(e_{\rm new}))$.
The units of its width cap are $e_{\rm new}=\log_2N$.
The old comparison is separately typed, with units
$w_{\rm old}=\log_2\log_2N_{\rm old}$, and never controls new selection.

A separately implemented arithmetic verifier binds exact request bytes,
source closure and eight current core/version roots, reconstructs every
stage and nested witness with its own arithmetic, and checks exact
member sets, types, canonical spelling and scope flags. Success verifies
bytes and arithmetic under the accepted mathematical premises. Matching
changes to code/proof roots define a different, unaccepted version.
\end{proposition}

\noindent\emph{Public availability: the following proof describes the original
private SVC package and proof map, which are not exported here.}

\begin{proof}
The strict request schema, parser and source contract are supplied with
the frozen code; the exact locations and hashes are in the proof map.
The parser rejects duplicate members, floats/nonfinite tokens, invalid
UTF-8/BOM, boolean integer fields, unknown members and unsupported
coordinate/tolerance syntax. Canonical integer/rational arithmetic
preserves order and multiplicity. It admits the domain in
Section~\ref{sva:domain}, not user coefficients or expressions.

The generator searches the first admissible $d,n,k,p,h_s$ in the stated
predicates. The verifier independently uses $d=5$ at $1/5$ and $d=q+3$
at $2^{-q}$, a rearranged rational shell bound, a Nelson term recurrence,
$p=4H_N+2d+8$ and $h_s=1+\operatorname{clog}_2(j_e+h_A+p+d)$.
The preceding even $p-2$ fails the first regulator predicate; the
preceding dyadic auxiliary exponent fails its predicate. Independent
minimal-power calculations give $h_A,o,h_{\rm stab},h_R$.
No generator selection functions are used by the verifier.

With $T=t\mathbin{//}2+(t\bmod2)$ the verifier independently calculates
$k_1=3x_e+y_e+3$, $n_0=2k_1+4$, $k_4=12x_e+4y_e+36$ and
$e_{\rm new}=k_4+4d+12$. It checks both parity slacks, the integer-ceiling
threshold, coefficient $38432$, joint moment order sixteen, the separate
physical variance flag, smoothing fractions and exact rational total.
It verifies the typed old guard and strict order without materializing
$N$, $2^p$, $2^t$ or $2^{w_{\rm old}}$.

The manifest has exactly 208 direct source rows: 195 inherited runtime
bodies and thirteen new mathematical/old-contract provenance bodies.
Every fixed safe path is bounded-read below the explicit source root,
checked for size/hash and required material aliases; symlink escape
refuses. Accepted SELECT/OPT/EB1/EC1/ADMIT and SVA decisions/gates retain
their exact material relations. The three runtime files, source manifest,
schema, software statement/proof and input contract are freshly hashed
as eight core bindings. External review and freeze supply a separate
version trust root. These checks establish identity and availability,
not signatures, classical truth or hostile-filesystem atomicity.

Request and result byte caps are respectively 32768 and 131072. Parser
nesting is sixteen; integer tokens have at most twenty characters
including sign, strings 4096, member names 256, and value/container nodes
5000. Each source read is capped at one MiB, aggregate source bytes at
sixteen MiB. All eight resource members are mandatory. Their exact
defaults and allowed intervals are shown in Table~\ref{sva:resources}.
A stricter valid cap can refuse before returning a certificate.

The clock begins after request parsing/validation; verifier certificate
parsing is inside its measured interval. Source reads, selected ticks
and named integer/Fraction watches poll time. Equality at the millisecond
ceiling passes; excess, invalid/nonfinite/boolean clock values or
backward clocks refuse. Values must be finite monotone integers or floats,
and an integer clock has absolute value at most $10^{12}$.
Named watches cover guarded source-coefficient,
spatial/allocation/auxiliary powers; selected rational geometry, Nelson,
outer and width exponents; old comparison integers; and the exact total.
Logical steps count the documented read/predicate/search ticks. Ordinary
Python intermediates and fixed small shifts are not universally watched.
Serialization/output, whole-process time, memory and filesystem atomicity
are outside the contract. Typed historical measurement metadata is
unauthenticated; the verifier measures its own interval.

The exact outcomes are supported-domain success, certificate verified,
unsupported input, numerical source gap, source binding refusal, resource
exhaustion and refused certificate. Refusals contain no numerical
certificate; CLI exit is zero only on success and two for defined
refusals. Interpreter/hardware/output-device failures remain outside this
exception contract. Executed finite tests and monkeypatch traps support
the current implementation correspondence, separately from the analytic
proof above and the upstream statement-only proof. The accepted decision
governs retained earlier candidate headers in those frozen files.
\end{proof}

\begin{table}[htbp]
\centering\small
\begin{tabular}{lll}\toprule
Resource member (prefix \texttt{max\_})&Default&Allowed interval\\\midrule
\texttt{watched\_bits}&1024&16--4096\\
\texttt{milliseconds}&30000&1--30000\\
\texttt{steps}&20000&1--50000\\
\texttt{n}&80&1--80\\
\texttt{regulator\_exponent}&20000&4--20000\\
\texttt{auxiliary\_radius\_exponent}&40&1--64\\
\texttt{stabilizer\_exponent}&$2^{40}$&$34$--$2^{50}$\\
\texttt{width\_exponent}&$2^{40}$&$1$--$2^{50}$\\\bottomrule
\end{tabular}
\caption{Exact SVC selected-resource contract. The last member bounds
$e=\log_2N$. The older \texttt{max\_width\_outer\_exponent} member is
incompatible and refuses. Named watches and cooperative time have the
specific scope stated in the proof.}\label{sva:resources}
\end{table}

\subsection{Source, history and construction boundaries}
\label{sva:boundaries}

The SVA mathematical decision, accepted at its exact conditional scope,
retains its original statement that SVC was then unlaunched. The later
accepted SVC decision supersedes only that implementation status for the
separately reviewed version. Its author/hostile/numerical/code-matching
records remain distinct from SVA's hostile, complete statement-only,
arithmetic and actual-proof records. All are internal fresh-session
model review, not an outside-expert or statistical-independence credential.
Earlier adverse attempts, source gaps, adapted versus unchanged replays
and findings/dispositions remain recorded under their original identities.
The original private companion index recorded which exact bodies were local.

In the historical private integration, the curated inherited companion
tree resolved the 23-packet index and its ninety indexed bindings locally,
alongside the seven later baseline packets and the SVA/SVC source trees.
That historical availability closed a dated assignment gap. The historical
tree did not recursively bundle every quoted path or supply unavailable
primary classical proofs. Those private source trees and their original
availability map are not attachments to this public edition. The public
proof map gives this edition's source locations and hashes. Hashes record
exact bytes only.

The deterministic common sequence, its four noneffective least inner
indices, quantifiers over all later tests, correlations and reflection
lists, and both
original architectures remain as in Theorem~\ref{vqal:seq}.
The quantitative schedule fixes observations first. MOM proves local
$H^{-2}$ moment convergence/UI; the separately derived SEQ-HS statement
gives physical/target $H^{-s}$ moment finiteness for real $s>1$ at its
literal scope. The inclusive qualitative cap and strict AP6 numerical
cap remain separate. The present single-exponent widths are sufficient
whole-plane process-neuron counts; practical/minimal/necessary complexity,
full-field rates, global neural Schwartz control, broader domains and
priority are separate questions. Public v1 is unchanged.

The standalone source contains the full present argument and every
preserved baseline mathematical block. The original private companion proof map binds each
new block and reading-guide headline to accepted statement and full-proof
clauses. The original reversible editorial delta, exact preservation audit,
build, reference/numerical receipts and all-page inspection record are
retained as private integration evidence. Fresh source/proof/visual review and coordinator
acceptance remain separate.

% Exact source appendix embedded in the standalone manuscript.
\clearpage
\begingroup
\normalsize
\renewcommand{\theHequation}{uda.\arabic{equation}}
% Current manuscript integration of accepted conditional UDAv2; no manuscript self-acceptance.
% Exact accepted sources and contributing draft provenance are bound in PROOF-MAP.json.
% MAP: status | S:1-6 P:1-7 C:3-8
\section{A conditional schedule for every dyadic accuracy}
\label{uda:current}
This appendix gives all six clauses of the internally accepted conditional
UDA-STAGE-RELATIVE-v2 and their complete proof. Its source statement and proof
are the exact bytes bound by the current decision. The manuscript presentation
awaits its own fresh source and actual-proof reviews. Every integer below is a
finite mathematical integer; software resource refusal does not make the
mathematical definition partial. Section~\ref{udc:contract} supplies the separately
accepted UDC implementation contract.

% MAP: model | S:10-21 P:11-19 C:18-21
\subsection{Model, observations and the literal conditional interface}
Fix $m=1$, $u=10^{-10}$, $\ell=1000$ and $K_{\mathrm{cov}}=1$.
The model is the original whole-plane, coefficient-one, same-mass,
own-variance Wick interaction $u{:}\phi^4{:}$, with no factorial.
The original complete state space is $X=H^{-2}_{\mathrm{loc}}$, and
$\nu^X$ is the original pressure-identified common Fr\"ohlich--Simon state.
The literal P1--P6, Gaussian/Wick, DSS/BD/BBD/Griffiths--Simon/SK-q/$P^*$,
H1/MIX/REG/smoothed-CT/FLOW/TP/T and classical/source premises retain their
original finite scopes. The quantitative bridge uses strict AP6,
$0<u/m^2<1/8000$. The original fixed $g_*$ and the inequality
$10^{-10}<g_*$ remain unknown and unused. NR1's sharper author constants and
larger statement-only blind core retain their separate original scopes;
excluded NR2 E4/r0 material is absent. The inherited source interface is
spelled out in Section~\ref{uda:source}; its source bodies are not newly proved.

% MAP: tests | S:23-46 P:21-39
Let
\[
 b(t)=\begin{cases}e^{-1/(1-t^2)},&|t|<1,\\0,&|t|\ge1,\end{cases}
 \qquad f_h(z)=b(z_1-h_1)b(z_2-h_2).
\]
A request is an ordered tuple $(h_1,\ldots,h_r)$ with $r\in\{1,2,3,4\}$.
Each coordinate is a canonical dyadic $P/2^k\in[-8,8]$, $0\le k\le16$:
integers and zero have $k=0$, and $P$ is odd if $k>0$.
Decimal numerators have no plus sign, leading zeros or negative zero.
Order and repetitions remain; amplitudes, signs, dilations and rotations
are unchanged. The primary accuracy is $\tau=2^{-q}$ for every integer
$q\ge3$; $\tau=1/5$ is a separate compatibility case. Set
\begin{align*}
 K&=\bigcup_{i=1}^r[h_{i1}-1,h_{i1}+1]\times[h_{i2}-1,h_{i2}+1],&
 R&=1+\max_{i,j}|h_{ij}|,\\
 c_r&=\lceil\sqrt r\rceil,&
 F_j&=\left(\sum_{i=1}^r\|f_{h_i}\|_j^2\right)^{1/2},
 \qquad F=F_{4/3},\quad\overline F=3c_r.
\end{align*}
For field laws $\rho$, let $\mathcal L_f(\rho)=(\mathcal O_f)_*\rho$, where
$\mathcal O_f(v)=(v(f_{h_i}))_{i=1}^r$ takes values in the full Euclidean
$\R^r$. The bounded-Lipschitz distance is
\[
 d_{\mathrm{BL}}(\rho,\sigma)
 =\sup_{\|G\|_\infty\le1,\ \operatorname{Lip}_{\R^r}(G)\le1}
 |\rho(G)-\sigma(G)|.
\]
The two constraints are separate. Write $\mathcal U_4(r)$ for real $C^4$
functions on $\R^r$ satisfying
\[\max_{0\le j\le4}\sup\|D^jG\|_{\mathrm{op}}\le1.\]
order zero means absolute value, and higher orders use full multilinear
Euclidean operator norms.

% MAP: coefficient | S:48-68 P:40-47
The literal CL1 coefficient uses the same chosen $C_H<2^{16}$ and
$C_{H2}=2^{384}$ in every expression. With $\kappa=1/256$,
$\epsilon_*=1/4$ and $c_{\mathrm{CT}}=81/5$, define
\begin{align*}
 \beta_{\mathrm{CT}}&=1+c_{\mathrm{CT}}\kappa,
 &\Gamma_G&=4\sqrt{2\pi}\,\kappa,\\
 C_*&=6c_{\mathrm{CT}}\kappa+c_{\mathrm{CT}}^2\kappa^2
       +\beta_{\mathrm{CT}}^2(\Gamma_G+3300u),\\
 g_{\mathrm{CL1}}(v)&=1+\frac{2}{10^6}\log\frac1{1-e^{-4000}}
                  +\frac\pi{16}\log(1+1/v),\quad0<v\le1,\\
 M(v)&=32\ell^{-2}\left[\frac{C_H^2g_{\mathrm{CL1}}(v)^3}{12}
                     +\frac{u^2 C_{H2}^2g_{\mathrm{CL1}}(v)^7}{9}\right],\\
 J_{\mathrm{CL1}}&=e^{8\sqrt2\kappa}\beta_{\mathrm{CT}}^2
 \left[e^{C_*}\int_0^1\sqrt{M(v)}\,dv
            +\frac{e^{2C_*}\sqrt{M(1)}}2\right],\\
 C_f(\mathbf h;R)&=(2R+1/2)\sqrt r\,e^{-2}J_{\mathrm{CL1}},&
 C_f(\mathrm{base};1)&=\tfrac52 e^{-2}J_{\mathrm{CL1}}.
\end{align*}
Assume the accepted literal base enclosure $C_f(\mathrm{base};1)\le2^{360}$
with its exact finite-source premises, and put
$\overline C=2^{360}c_r(4R+1)/5$. The separate factor $e^{R/256}$ remains.
No independently optimized coefficient replaces the source integral.

% MAP: covariance | S:70-92
For every fixed $0<a<1$ and finite $D=\Omega=[-R_{\mathrm{aux}},R_{\mathrm{aux}}]^2$
containing the physical interaction-square closure and every test support,
assume the original finite-law interface below. Its finite-positive-stage
provenance is EB1/NR1/P6. Bounded SVA or fixed-$n75$ VEC acceptance is not
an all-stage premise, and no all-$q$ conclusion is assumed. The physical
covariance and variance are
\begin{align*}
 C_a(z)&=(2\pi)^{-2}\int_{\R^2}
 \frac{e^{i(\omega,\xi)\cdot z}}{1+\omega^2+\xi^2+a^2\omega^2\xi^2}
 \,d\omega\,d\xi,\\
 c=c_a&=\frac1{2\pi}\int_0^\infty
 \frac{d\xi}{\sqrt{(1+\xi^2)(1+a^2\xi^2)}},
 &\kappa_a^2&=a^{-2}-1,\\
 c&\le\overline c(a)=\frac{\operatorname{asinh}(1/a)+1}{2\pi}.
\end{align*}
On the separable real Banach space $E=L^4(D)$, $\gamma$ is the original
centered continuous Gaussian endpoint, strongly measurable, with zero
Bochner mean, global sign symmetry and covariance
$C_{\mathrm{end}}(z,z')=e^{-|z-z'|_1/a}/4$. Its point variance is $1/4$.
The Fourier integration variable $\xi$ above is local and does not denote
an observation or the standard Gaussian vector used for smoothing.

% MAP: neurons | S:94-116
For $b\in\{\mathrm{product},\mathrm{Mondrian}\}$, $g^b$ is the restriction
of the original whole-plane process neuron. The product neuron is
$g(z)=\sigma(z_1)\tau(z_2)/2$ for independent stationary symmetric
telegraph processes of flip rate $1/(2a)$. The Mondrian neuron uses the
projectively consistent whole-plane axis-parallel partition at budget
$1/a$, independently uniformly signed cells, and point values half the
cell sign. A rectangle of sides $w,h$ splits at rate $w+h$, vertically
with probability $w/(w+h)$ and otherwise horizontally, at a uniform
coordinate; daughters evolve independently for the remaining budget.
Use the original lower-left half-open representative. Retain the original
existence, projective-consistency and measurability premises. Each $g^b$
is strongly measurable and $E$-valued, half-valued almost everywhere almost
surely, has zero Bochner mean, global sign symmetry and covariance
$C_{\mathrm{end}}$. Boundary representatives agree in $E$. For \emph{every}
bounded bilinear form $B_0:E\times E\to\R$, assume
\[
 \E B_0(g^b,g^b)=\E B_0(\gamma,\gamma).
\]
Independent copies may be combined on product probability spaces. These
are the original full laws and their restrictions; point-pair covariance
does not define this full interface.

% MAP: potential | S:118-142
For $\Lambda=(-1000n,1000n)^2$ with $\overline\Lambda\subset\Omega$, the
actual potential, weight and normalized laws are
\begin{align*}
 V(v)&=u\int_\Lambda(v^4-6cv^2+3c^2)
       +\epsilon\int_{\Omega\setminus\Lambda}v^4
       -\frac{\kappa_a^2}{2}\int_\Omega v^2,&w(v)&=e^{-V(v)},\\
 F_N^b&=N^{-1/2}\sum_{i=1}^N g_i^b,&
 P_N^b&=\frac{w\,\Law(F_N^b)}{\E w(F_N^b)},\qquad
 Q_\epsilon=\frac{w\,\Law(\gamma)}{\E w(\gamma)}.
\end{align*}
Their observation laws are exactly the original whole-plane neural and
full-Gaussian-tilt laws. The physical $c$ remains in $V$; a positive
majorant $c_*$ never changes it. Width counts whole-plane process neurons,
each with countably many internal variables. The ambient L0 premise says
that for unit BL $G$ on $\R^r$ and standard Gaussian $\xi\in\R^r$,
\begin{align*}
 G_\zeta(y)&=\E G(y-\zeta\xi),&
 \|G-G_\zeta\|_\infty&\le\sqrt r\,\zeta,\\
 \|G_\zeta\|_\infty&\le1,&\|DG_\zeta\|_{\mathrm{op}}&\le1,\\
 \|D^jG_\zeta\|_{\mathrm{op}}&\le\sqrt{(j-1)!}\,\zeta^{1-j},&&j\ge2.
\end{align*}
For $0<\zeta\le1$, every norm through order four is at most
$\sqrt6\zeta^{-3}$, even for singular observation covariance.

% MAP: endpoints | S:144-175 C:22-24
The full Gaussian tilt $Q_\epsilon$ has the original residual/mass-tilt
representation. Its zero-stabilizer residual endpoint is $Q_0$, the
positive-regulator physical law is $Q_a$, and its own-variance
zero-regulator fourth-chaos limit is $Q_{\mathrm{phys}}$. On the actual
common Gaussian space, $Q_a$ has density $e^{-U_{a,\Lambda}}/Z_a$,
$\E U_{a,\Lambda}=0$ and $Z_a\ge1$. The compact-test laws of
$Q_{\mathrm{phys}}$ are the original classical free finite-volume laws.
AP5/AP6 identify the same original target $\nu^X$. These are inherited
conditional identifications, with the literal conventions of
Section~\ref{uda:source}. For $A=|\Lambda|>0$, $U_A=uA$,
$O=|\Omega|$ and $B_{\mathrm{ext}}=O-A$, put
\begin{align*}
 K_\Lambda^2&=\frac{75\sqrt{85}}{2\pi^3}u^2 A,&
 \kappa_0&=e(1+\sqrt2),\\
 K_{\mathrm{sharp}}&=\max\{12U_A,(2/e)K_\Lambda\sqrt{\kappa_0}\},&
 b_{\mathrm{Nelson}}&=\frac{\pi^2}{48U_A},\\
 T_\Lambda(H)&=\exp[-\exp(\sqrt{b_{\mathrm{Nelson}}H}-1)],\\
 \Xi_2&=e^{2K_{\mathrm{sharp}}}
       +\int_{K_{\mathrm{sharp}}}^\infty2e^{2H}T_\Lambda(H)\,dH,\\
 H_2&=\max\left\{K_{\mathrm{sharp}},
 \frac4{b_{\mathrm{Nelson}}}
 \left[\log_+\frac{64}{e b_{\mathrm{Nelson}}}\right]^2\right\},&
 \overline\Xi_2&=2e^{2H_2},\quad\Xi_2\le\overline\Xi_2,\\
 \rho_a(s)&=\frac{\kappa_a^2}{4\pi}\min\{1,2e^{-\sqrt2s}\},&
 q_{\mathrm{sp}}&=e^{-1000/256},\\
 T_n&=q_{\mathrm{sp}}^n\left[\frac{2n+1}{1-q_{\mathrm{sp}}}
               +\frac{2q_{\mathrm{sp}}}{(1-q_{\mathrm{sp}})^2}\right].
\end{align*}
Here $e$ is Euler's number and $\log_+x=\max(0,\log x)$.
EC1.a--c apply at every fixed finite positive physical area; $\Xi_2$ is
not volume-uniform. With $\Delta=C_a-C'_{\Omega,a}$ positive semidefinite,
$\delta_\Omega(H)=\sup_{z\in H}\Delta(z,z)$ is at most $\rho_a(s)$
whenever the strict Euclidean $s$-neighborhood of $H$ lies in $\Omega$.

% MAP: errors | S:177-193
The actual successive errors $Q_\epsilon\to Q_0\to Q_a\to
Q_{\mathrm{phys}}\to\nu^X$, in the stated observation metric (variation
for the first arrow), are bounded by
\begin{align*}
 S&=\min\{2,6\epsilon B_{\mathrm{ext}}c^2e^{9U_Ac^2}\},\\
 A_{\mathrm{aux}}&=\min\left\{2,e^{9U_Ac^2}
 \left[2U_A(4\sqrt{15}+12)c^{3/2}\sqrt{\delta_\Omega(\overline\Lambda)}
                 +F_1\sqrt{\delta_\Omega(K)}\right]\right\},\\
 R_{\mathrm{phys}}&=\min\{2,\sqrt{\Xi_2}[4K_\Lambda\sqrt a+(a/2)F_2]\},\\
 E_n&=4C_f(\mathbf h;R)u\ell^2e^{R/256}T_n,
 \qquad n\ge1,\quad1000n\ge R.
\end{align*}
The strict source geometry remains mandatory. On the physical common
space, $\E|w_a-w_0|\le2K_\Lambda\sqrt{\Xi_2}\sqrt a$,
$\|w_0\|_2\le\sqrt{\Xi_2}$ and
$\E|(\Phi_a-\Phi_0)(f)|^2\le(a^2/4)\|f\|_2^2$.
Full neural/Gaussian normalizers, physical $Z_a\ge1$ and the residual
Jensen normalizer are distinct. Strict AP6 supplies this original-target
route without using or establishing $10^{-10}<g_*$.

% MAP: relative-constants | S:195-225
\subsection{Relative constants and the complete old coefficient}
At finite positive stages with $0<\epsilon<u$ and $c_*\ge c$ used only
in positive majorants, define
\begin{align*}
 \alpha&=\epsilon,& U&=u,\\
 B&=\kappa_a^2/2+6uc_*,&L&=B^2O/(2\alpha),\\
 J_z&=3UO/16+3uc_*^2 A,&H&=L+J_z.
\end{align*}
\begin{align*}
 v_1(x)&=4Ux^3+2B\sqrt O\,x,&v_2(x)&=12Ux^2+2B\sqrt O,\\
 v_3(x)&=24Ux,&v_4(x)&=24U.
\end{align*}
\begin{align*}
 e_0&=1,&e_1&=v_1,&e_2&=v_1^2+v_2,\\
 e_3&=v_1^3+3v_1v_2+v_3,\\
 e_4&=v_1^4+6v_1^2v_2+3v_2^2+4v_1v_3+v_4.
\end{align*}
\begin{align*}
 P_F(x)&=2e_4(x)+\sum_{j=1}^4\binom4j F^j e_{4-j}(x),\\
 X_{\mathrm{rel}}&=[2(H+4)/\alpha]^{1/4},&
 Y_{\mathrm{rel}}&=2O^{1/4}(H+10)^{1/2},\\
 K_1&=v_1(X_{\mathrm{rel}})Y_{\mathrm{rel}},&
 N_0&=\max\{1,\lceil16K_1^2\rceil\},\\
 K_4(F)&=P_F(X_{\mathrm{rel}})Y_{\mathrm{rel}}^4/6.
\end{align*}
Derivative norms use $E=L^4(D)$. The scalar moment bounds
$X_{\mathrm{rel}},Y_{\mathrm{rel}}$ are distinct from the field space $X$.
For $P(x)=\sum_k a_kx^k$ with nonnegative coefficients, put
\begin{align*}
 \beta_0&=1,&\beta_k&=\left(\frac{k}{2\alpha e}\right)^{k/4}\ (k>0),\\
 \mathcal T(P)&=e^L\sum_k a_k\beta_k,\\
 M_Z&=\mathcal T(e_4),&
 M_A&=\mathcal T\left(\sum_{j=0}^4\binom4jF^je_{4-j}\right),\\
 z&=e^{-J_z},&C_{\mathrm{old}}&=\frac{O(M_A+M_Z)}{96z},\\
 D_{\mathrm{old}}^4&=\frac{2048\sqrt6}{27}r^{3/2}C_{\mathrm{old}}.
\end{align*}
One may replace $F$ by $\overline F$ and $c$ by $p$ only in positive
majorants. The complete old coefficient retains coercivity and both Jensen terms.

% MAP: selection | S:227-255
\subsection{Universal rational stage prescription}
For positive rational $v$, let $\operatorname{clog}_2(v)$ be the least
integer $h\ge0$ with $v\le2^h$ (the source's \texttt{ceil\_log2}).
For $\tau=2^{-q}$ set $d=q+3$; for $\tau=1/5$ set $d=5$.
Let $\eta=2^{-d}$ and, in the displayed order, define
\begin{align*}
 \overline{\exp R}&=\frac{256}{256-R},\\
 \overline E(i)&=\frac4{10000}\overline C\,\overline{\exp R}\,2^{-5i}
              \left[(2i+1)\frac{32}{31}+\frac{64}{961}\right],\\
 n&=\min\{i\ge1:\overline E(i)\le\eta\},\\
 A&=4000000n^2,\quad U_A=uA,\quad b_0=3/(16U_A),\\
 h_K&=\min\{h\ge0:14u^2 A\le2^{2h}\},\\
 \overline K&=2^{h_K},\\
 L_{10}(k)&=\sum_{j=0}^{10}\frac{k^j}{j!},\\
 k&=\min\{i\ge1:L_{10}(i)>32/b_0\},\\
 H_N&=1+\left\lceil\max\{12U_A,4\overline K,4k^2/b_0\}\right\rceil,\\
 p&=4H_N+2h_K+2d+8,\\
 a&=2^{-p},\\
 j&=\lceil18U_Ap^2\rceil,\\
 B_{\mathrm{aux}}&=64U_Ap^2+4c_r,\\
 h_{\mathrm{aux}}&=\operatorname{clog}_2(B_{\mathrm{aux}}),\\
 h_s&=1+\operatorname{clog}_2(j+h_{\mathrm{aux}}+p+d),\\
 s&=2^{h_s},\\
 R_{\mathrm{aux}}&=1000n+s+1,\\
 O&=4R_{\mathrm{aux}}^2,\\
 o&=\operatorname{clog}_2(O),\\
 h_{\mathrm{stab}}&=\operatorname{clog}_2(6Op^2),\\
 t&=\max\{34,j+h_{\mathrm{stab}}+d\},\quad\epsilon=2^{-t},\\
 \Lambda&=(-1000n,1000n)^2,\\
 \Omega=D&=[-R_{\mathrm{aux}},R_{\mathrm{aux}}]^2.
\end{align*}
Use $c_*=p$ in relative majorants and put
\begin{align*}
 o_+&=\max\{o,78\},\\
 T_{1/2}&=\lceil t/2\rceil,\\
 x&=T_{1/2}+p+\left\lceil\frac{o_++3}{4}\right\rceil,\\
 y&=T_{1/2}+2p+\left\lceil\frac{3o_++8}{4}\right\rceil,\\
 k_1&=3x+y+3,\\
 n_0&=2k_1+4,\quad k_4=12x+4y+36,\\
 \zeta&=2^{-d-3},\\
 e_{\mathrm{new}}&=k_4+4d+12,
 \qquad N_{\mathrm{new}}=2^{e_{\mathrm{new}}}.
\end{align*}
The notation $o_+,T_{1/2}$ translates the source's \texttt{oplus,T\_half}.

% MAP: old-selection | S:257-267
At these same outer stages, define
\begin{align*}
 h_{\mathrm{root}}&=\operatorname{clog}_2(2R_{\mathrm{aux}}),&
 v_{\mathrm{old}}&=2p+h_{\mathrm{root}}+2,\\
 q_{\mathrm{old}}&=4v_{\mathrm{old}}+3t+28,&
 l_{\mathrm{old}}&=t+4p+o,\\
 H_{\mathrm{small}}&=q_{\mathrm{old}}+o+11+4d.
\end{align*}
Start at $w=l_{\mathrm{old}}+2$ and increment until both
$w>l_{\mathrm{old}}+2$ and
$\operatorname{bitlen}(H_{\mathrm{small}})\le l_{\mathrm{old}}+2$ hold.
Call the result $w_{\mathrm{old}}$ and put
$N_{\mathrm{old}}=2^{2^{w_{\mathrm{old}}}}$.
For positive integer $M$, $\operatorname{bitlen}(M)$ is the least positive
$h$ with $M<2^h$. These are sufficient-width rules, rather than necessary
costs or canonical inner least-index definitions.

% MAP: legacy-selection | S:269-322
\subsection{The explicitly defined legacy comparator}
For $\tau\in\{1/5\}\cup\{2^{-q}:3\le q\le12\}$, use the same ordered
tuple, $r,R,c_r,\overline C,\overline E$ and $L_{10}$. The suffix $L$
denotes a separately defined legacy rule. Its termination, sufficiency
and equality with the universal rule will be proved, rather than assumed.
Define in order
\begin{align*}
 d_L&=\min\{i\ge0:5\cdot2^{-i}\le\tau\},\\
 \eta_L&=2^{-d_L},\\
 n_L&=\min\{1\le i\le80:\overline E(i)\le\eta_L\},\\
 A_L&=4000000n_L^2,\\
 U_{AL}&=uA_L,\quad b_{0L}=3/(16U_{AL}),\\
 k_L&=\min\{1\le i\le7:L_{10}(i)>32/b_{0L}\},\\
 H_L&=1+\lceil\max\{12U_{AL},4,4k_L^2/b_{0L}\}\rceil,\\
 p_L&=\min\left\{\begin{array}{l}i\ge4\text{ even}:\\
 2H_L+3-i/2\le-d_L-1,\\
 2H_L+2-i\le-d_L-1\end{array}\right\},
 \\
 a_L&=2^{-p_L},\\
 j_L&=\lceil18U_{AL}p_L^2\rceil,\\
 B_{\mathrm{aux},L}&=64U_{AL}p_L^2+4c_r,\\
 h_{\mathrm{aux},L}&=\operatorname{clog}_2(B_{\mathrm{aux},L}),\\
 h_{sL}&=\min\{i\ge1:j_L+h_{\mathrm{aux},L}+p_L-2^{i-1}\le-d_L\},\\
 s_L&=2^{h_{sL}},\\
 R_{\mathrm{aux},L}&=1000n_L+s_L+1,\\
 O_L&=4R_{\mathrm{aux},L}^2,\\
 o_L&=\operatorname{clog}_2(O_L),\\
 h_{\mathrm{stab},L}&=\operatorname{clog}_2(6O_Lp_L^2),\\
 t_L&=\max\{34,j_L+h_{\mathrm{stab},L}+d_L\},\\
 \epsilon_L&=2^{-t_L},\\
 \Lambda_L&=(-1000n_L,1000n_L)^2,\\
 D_L=\Omega_L&=[-R_{\mathrm{aux},L},R_{\mathrm{aux},L}]^2.
\end{align*}
For the legacy old-width rule, put
\begin{align*}
 h_{\mathrm{root},L}&=\operatorname{clog}_2(2R_{\mathrm{aux},L}),\\
 v_{\mathrm{old},L}&=2p_L+h_{\mathrm{root},L}+2,\\
 q_{\mathrm{old},L}&=4v_{\mathrm{old},L}+3t_L+28,\\
 l_{\mathrm{old},L}&=t_L+4p_L+o_L,\\
 H_{\mathrm{small},L}&=q_{\mathrm{old},L}+o_L+11+4d_L.
\end{align*}
Start at $w=l_{\mathrm{old},L}+2$ and increment by one until both
$w>l_{\mathrm{old},L}+2$ and
$\operatorname{bitlen}(H_{\mathrm{small},L})\le l_{\mathrm{old},L}+2$ hold.
The selected value is $w_{\mathrm{old},L}$, with
$N_{\mathrm{old},L}=2^{2^{w_{\mathrm{old},L}}}$. The legacy SVC relative
exponents are
\begin{align*}
 T_L&=\lceil t_L/2\rceil,\\
 x_L&=T_L+p_L+21,\\
 y_L&=T_L+2p_L+61,\\
 k_{1L}&=4T_L+5p_L+127,\\
 n_{0L}&=8T_L+10p_L+258,\\
 k_{4L}&=16T_L+20p_L+532,\\
 \zeta_L&=2^{-d_L-3},\\
 e_L&=k_{4L}+4d_L+12,\\
 N_L&=2^{e_L}.
\end{align*}
These mathematical rational/integer definitions remain meaningful if a
software run refuses for a resource limit. They identify no request or
certificate with another and assert no software execution.

% MAP: clause-a | S:324-349
\subsection{Six exact candidate clauses}
\begin{theorem}[UDA.a: finite selection and four downstream errors]
\label{uda:clause-a}
Every allowed request has finite selected stages, with termination witnesses
\[
 n\le\left\lceil\frac{d+356}{4}\right\rceil,\quad
 k\le\lceil32/b_0\rceil,\quad
 h_K\le\lceil14u^2A\rceil,\quad
 \operatorname{clog}_2(v)\le\lceil v\rceil.
\]
The last bound holds for every positive argument used. All subsequent
formulas are finite compositions. The integer $p$ is the first even
integer at least four satisfying both
\[
 2H_N+h_K+3-p/2\le-d-1,\qquad2H_N+2-p\le-d-1,
\]
and $h_s$ is the first integer at least one satisfying
$j+h_{\mathrm{aux}}+p-2^{h_s-1}\le-d$. Moreover
\begin{gather*}
 1\le R\le9,\quad K\subset\Lambda_1\text{ (open)},\quad
 F_1\le4c_r\le8,\quad F_2\le2c_r\le4,\quad F\le\overline F\le6,\\
 C_f(\mathbf h;R)\le\overline C<2^{364},\qquad
 R_{\mathrm{aux}}-\max(1000n,R)-s=1,\\
 0<a<1,\quad0<\epsilon<u,\quad c\le\overline c(a)<p,\\
 K_\Lambda<\overline K,\quad H_2<H_N,\quad
 \sqrt{\overline\Xi_2}<2^{2H_N+1},\\
 E_n\le\overline E(n)\le\eta,\qquad
 R_{\mathrm{phys}}\le\eta,\quad A_{\mathrm{aux}}\le\eta,\quad S\le\eta.
\end{gather*}
\end{theorem}

% MAP: clause-b | S:351-396
\begin{theorem}[UDA.b: arbitrary finite-stage joint laws and derivatives]
\label{uda:clause-b}
At any fixed finite stage satisfying the premises above, for either
architecture, $N\ge1$, $0\le j\le N-1$, and an independent direction
$h$ with neuron or Gaussian law, let
\begin{align*}
 S_j&=N^{-1/2}\left[\sum_{i=1}^j\gamma_i+
                         \sum_{i=j+1}^{N-1}g_i^b\right],\\
 v&=S_j+\theta h/\sqrt N,&Z_{j,h}(\theta)&=\E w(v),\\
 \Pi_{j,h,\theta}&=w(v)\Law(S_j,h)/Z_{j,h}(\theta),&&0\le\theta\le1.
\end{align*}
The untilted $v$ has almost-everywhere point variance at most $1/4$ and
fourth moment at most $3/16$. Every $Z$ is positive finite. Writing
$x(v)=\|v\|_4$, for $z\ge0$ and every real $1\le q_0\le16$,
\begin{gather*}
 w(v)\le e^{L-\alpha x(v)^4/2},\quad Z\ge e^{-J_z},\quad
 0<\frac{d\Pi}{d\Law(S_j,h)}\le e^{H-\alpha x(v)^4/2}\le e^H,\\
 \Pi(x(v)^4>z)\le\min\{1,e^{H-\alpha z/2}\},\\
 \|x(v)\|_{L^{q_0}(\Pi)}\le X_{\mathrm{rel}},\quad
 \|\|h\|_4\|_{L^{q_0}(\Pi)}\le Y_{\mathrm{rel}},\quad
 \E e^{\|\gamma\|_4^2/(4\sqrt O)}<4.
\end{gather*}
Each $Z$ and $\log Z$ extends $C^1$ to a neighborhood of $[0,1]$, with
\[
 (\log Z)'=-N^{-1/2}\E_\Pi DV(v)[h],\qquad
 |(\log Z)'|\le K_1/\sqrt N.
\]
The zero-amplitude partitions agree. For $N\ge N_0$,
\[
 \frac{Z_{j,g}(\theta)}{Z_{j,g}(1)}\le e^{K_1/\sqrt N},\qquad
 \frac{Z_{j,\gamma}(\theta)}{Z_{j,g}(1)}\le e^{2K_1/\sqrt N}<2.
\]
The potential is $C^4(E)$, with
$\|D^kV(v)\|_{\mathrm{op}}\le v_k(x(v))$ for $1\le k\le4$ and
$\|D^kw(v)\|_{\mathrm{op}}\le w(v)e_k(x(v))$ for $0\le k\le4$.
All derivatives of $w$ through four are globally bounded. For
$G\in\mathcal U_4(r)$, define
\begin{align*}
 m_j&=\frac{\E[G(\mathcal O_f(S_j+\gamma/\sqrt N))
                       w(S_j+\gamma/\sqrt N)]}{Z_{j,\gamma}(1)},\\
 \Phi_j(v)&=[G(\mathcal O_f(v))-m_j]w(v).
\end{align*}
Then $|m_j|\le1$, $\Phi_j\in C_b^4(E)$, $\deg P_F\le12$, and under
each actual joint $\Pi$,
\[
 \|D^4\Phi_j(v)\|_{\mathrm{op}}\le w(v)P_F(x(v)),\qquad
 \E_\Pi[P_F(x(v))\|h\|_4^4]\le P_F(X_{\mathrm{rel}})Y_{\mathrm{rel}}^4.
\]
This is a joint-law estimate, without a product of tilted marginal expectations.
\end{theorem}

% MAP: clause-c | S:398-414
\begin{theorem}[UDA.c: normalized laws, ambient BL and selected width]
\label{uda:clause-c}
At the arbitrary finite stages of UDA.b and for every $N\ge N_0$, each
actual adjacent normalized replacement has $\mathcal U_4$ error at most
$K_4(F)/N^2$, and
\begin{align*}
 |P_N^b(G\circ\mathcal O_f)-Q_\epsilon(G\circ\mathcal O_f)|
 &\le\min\{2,K_4(F)/N\},\quad G\in\mathcal U_4(r),\\
 d_{\mathrm{BL}}(\mathcal L_f(P_N^b),\mathcal L_f(Q_\epsilon))
 &\le\min\{2,\sqrt6K_4(F)/(N\zeta^3)+2\sqrt r\zeta\},\quad0<\zeta\le1,\\
 d_{\mathrm{BL}}(\mathcal L_f(P_N^b),\mathcal L_f(Q_\epsilon))
 &\le\min\{2,D_rN^{-1/4}\},\qquad
 D_r^4=\frac{2048\sqrt6}{27}r^{3/2}K_4(F).
\end{align*}
At every selected stage, including both parities of $t$,
\begin{gather*}
 X_{\mathrm{rel}}<2^x,\quad Y_{\mathrm{rel}}<2^y,\quad
 K_1<2^{k_1-2},\quad16K_1^2<2^{2k_1},\\
 N_0\le2^{2k_1}<2^{n_0}<N_{\mathrm{new}},\quad
 K_4(F)\le K_4(\overline F)<2^{k_4},\\
 d_{\mathrm{BL}}(\mathcal L_f(P_{N_{\mathrm{new}}}^b),
                  \mathcal L_f(Q_\epsilon))<7\eta/8<\eta.
\end{gather*}
\end{theorem}

% MAP: clause-d | S:416-424
\begin{theorem}[UDA.d: actual target at every dyadic accuracy]
\label{uda:clause-d}
For both original architectures, every integer $q\ge3$ and every allowed
ordered request,
\[
 d_{\mathrm{BL}}(\mathcal L_f(P_{N_{\mathrm{new}}}^b),\mathcal L_f(\nu^X))
 \le\overline E(n)+4\eta\le5\eta\le2^{-q}.
\]
The separate $1/5$ case has the same conclusion with right side $\tau$.
The five actual arrows are
\[
 P_{N_{\mathrm{new}}}^b\longrightarrow Q_\epsilon\longrightarrow Q_0
 \longrightarrow Q_a\longrightarrow Q_{\mathrm{phys}}\longrightarrow\nu^X.
\]
\end{theorem}

% MAP: clause-e | S:426-442
\begin{theorem}[UDA.e: complete old coefficient and same-stage order]
\label{uda:clause-e}
At all selected stages, the complete $C_{\mathrm{old}}$ gives, for every
$N\ge1$ and $G\in\mathcal U_4(r)$,
\begin{gather*}
 |P_N^b(G\circ\mathcal O_f)-Q_\epsilon(G\circ\mathcal O_f)|
 \le\min\{2,C_{\mathrm{old}}/N\},\\
 d_{\mathrm{BL}}(\mathcal L_f(P_N^b),\mathcal L_f(Q_\epsilon))
 \le\min\{2,D_{\mathrm{old}}N^{-1/4}\}.
\end{gather*}
With $c_*=p$ and $\overline F$ in positive majorants,
\begin{gather*}
 C_{\mathrm{old}}<2^{\,2^{l_{\mathrm{old}}+2}+q_{\mathrm{old}}+o},\qquad
 H_{\mathrm{small}}\le7l_{\mathrm{old}}+47<2^{l_{\mathrm{old}}+2},\\
 w_{\mathrm{old}}=l_{\mathrm{old}}+3,\qquad
 d_{\mathrm{BL}}(\mathcal L_f(P_{N_{\mathrm{old}}}^b),
                  \mathcal L_f(Q_\epsilon))\le\eta,\\
 e_{\mathrm{new}}\le12l_{\mathrm{old}}+553<2^{l_{\mathrm{old}}+3}
 =2^{w_{\mathrm{old}}},\qquad N_{\mathrm{new}}<N_{\mathrm{old}}.
\end{gather*}
The sufficient mathematical integers are compared at identical stages,
allocations, model, source premises, architectures, tests and target.
\end{theorem}

% MAP: clause-f | S:444-461
\begin{theorem}[UDA.f: exact old-domain compatibility]
\label{uda:clause-f}
For $\tau\in\{1/5\}\cup\{2^{-q}:3\le q\le12\}$ the explicit legacy
rules terminate. Every common outer variable
\[
\begin{gathered}
 d,\eta,n,A,U_A,b_0,k,H_N,p,a,j,B_{\mathrm{aux}},h_{\mathrm{aux}},h_s,s,
 R_{\mathrm{aux}},O,o,h_{\mathrm{stab}},t,\epsilon,\\
 \Lambda,D,\Omega,h_{\mathrm{root}},v_{\mathrm{old}},q_{\mathrm{old}},
 l_{\mathrm{old}},H_{\mathrm{small}},w_{\mathrm{old}}
\end{gathered}
\]
equals its displayed legacy counterpart, with $H_N$ corresponding to $H_L$
and $U_A$ to $U_{AL}$. Also
\[
 (T_{1/2},x,y,k_1,n_0,k_4,\zeta,e_{\mathrm{new}},N_{\mathrm{new}})
 =(T_L,x_L,y_L,k_{1L},n_{0L},k_{4L},\zeta_L,e_L,N_L),
\]
and $N_{\mathrm{old}}=N_{\mathrm{old},L}$. In particular
\begin{gather*}
 h_K=0,\quad o_+=78,\quad n\le80,\quad k\le7,\quad H_N\le3137,\\
 p\le12586,\quad h_s\le36,\quad O<2^{76},\quad o\le76,\quad t<2^{36},\\
 x=T_{1/2}+p+21,\qquad y=T_{1/2}+2p+61,\\
 k_1=4T_{1/2}+5p+127,\qquad n_0=8T_{1/2}+10p+258,\\
 k_4=16T_{1/2}+20p+532,\qquad
 e_{\mathrm{new}}=16\lceil t/2\rceil+20p+4d+544.
\end{gather*}
This mathematical compatibility identifies no old/new request or certificate
and supplies no software extension.
\end{theorem}

% MAP: theorem-scope | S:463-470 P:767-776 C:34-35
All clauses retain ordered repeats, overlap, touching, separation and singular
laws in the full ambient $\R^r$. They assert no independence of interacting
observations. Widths are sufficient, possibly astronomical, whole-plane
process-neuron counts, each retaining countably many internal variables.
Practical, minimal or necessary complexity, finite global scalar weights,
measured speedup, canonical least inner indices, an observable-independent
effective common sequence, global neural Schwartz control, whole-field rates,
larger models/dimensions/test families, novelty priority and public release
remain outside these clauses.

% Full proof translation. Source ranges are the registered PROOF.md, unchanged.
% MAP: proof-intro | P:1-19 S:1-21 C:18-21
\subsection{Complete conditional proof}
\label{uda:proof}
The following is the author derivation of the exact frozen conditional
statement. Its premise is the explicit finite interface above and the literal
inherited qualifications below. The original source truths and NR1
author-versus-blind-core scopes are retained. In particular, the full laws
are supplied at every finite positive stage by EB1/NR1/P6; bounded SVA and
fixed-$n75$ VEC acceptances are not extrapolated. Strict AP6 applies because
$0<10^{-10}<1/8000$ and identifies the original FS $\nu^X$ without the old
$g_*$ inequality. Finite diagnostics are evidence, rather than premises.

% MAP: proof-1-smooth | P:21-28 S:23-36
\subsubsection{Observation geometry and the literal coefficient}
For $|t|<1$, each derivative of $b(t)$ is $e^{-1/(1-t^2)}$ times a rational
function whose endpoint poles have finite order. For every integer $k\ge0$,
$x^ke^{-x}\to0$ as $x\to\infty$, for example by the positive exponential
series. All derivatives therefore extend by zero at both endpoints.
Thus $b$ and its translates are real, nonzero and compact smooth. Since
$\max b=e^{-1}$, $f_h\le e^{-2}<1$, and its support is the stated closed
square of area four. Hence $1\le R\le9$, all supports lie in $[-9,9]^2$,
and $K\subset\Lambda_1=(-1000,1000)^2$, with the latter square open.

% MAP: proof-1-norms | P:30-39 S:34-46
Translation invariance and the support/height bounds give
$\|f_h\|_1\le4$, $\|f_h\|_2\le2$ and
$\|f_h\|_{4/3}\le4^{3/4}<3$, because $64<81$.
Taking the square sum over the actual ordered $r$ coordinates gives
$F_1\le4\sqrt r\le4c_r$, $F_2\le2c_r$ and
$F\le3c_r=\overline F\le6$. H\"older's inequality coordinate by coordinate gives
\begin{equation}
 |\mathcal O_f(v)|_2^2\le\|v\|_4^2\sum_i\|f_{h_i}\|_{4/3}^2
 =F^2\|v\|_4^2.\tag{1.1}
\end{equation}
Repeated tests remain in this sum. No support separation, covariance rank
or inverse matrix enters the argument.

% MAP: proof-1-ratio | P:40-47 S:48-68
The same $J_{\mathrm{CL1}}$ and chosen H1/H2 coefficients occur in both
literal source coefficients, so their ratio is exactly
\[
 \frac{C_f(\mathbf h;R)}{C_f(\mathrm{base};1)}
 =\frac{\sqrt r(4R+1)}5.
\]
The base enclosure gives $C_f\le\overline C$. Since
$c_r(4R+1)/5\le74/5<16$, $\overline C<2^{364}$.
The separate factor $e^{R/256}$ is retained. This proves the observation
and coefficient portion of UDA.a.

% MAP: proof-2-accuracy | P:49-54 S:229-230
\subsubsection{Universal outer selection and the four actual errors}
For $q\ge3$, $d=q+3$ is the least integer with
$5\cdot2^{-d}\le2^{-q}$, since $4<5\le8$. For $\tau=1/5$ the least
integer is $d=5$. Thus $d\ge5$, $\eta=2^{-d}>0$ and $5\eta\le\tau$.
There is no upper bound on $d$ in proof steps 2--11.

% MAP: proof-2-spatial | P:56-82 S:168-185 S:232-235
The positive degree-eight exponential partial at $125/32$ is
\[
 \sum_{j=0}^8\frac{(125/32)^j}{j!}
 =\frac{432528188408407789}{8866461766385664}>32.
\]
Therefore $0<q_{\mathrm{sp}}=e^{-125/32}<1/32$.
For $0\le v<1$, the exponential series is termwise at most the geometric
series, so $e^v\le(1-v)^{-1}$. In particular
$e^{R/256}\le256/(256-R)<2$. Both positive shell bracket terms increase
with $q_{\mathrm{sp}}\in(0,1)$; substituting $1/32$ gives exactly
$(2n+1)32/31+64/961$. Since $4u\ell^2=4/10000$ and $C_f\le\overline C$,
\begin{equation}E_n\le\overline E(n).\tag{2.1}\end{equation}
The rational bracket is $(1984n+1056)/961<4n$ for $n\ge1$.
Also $(4/10000)\overline{\exp R}<8/10000<2^{-10}$ and
$\overline C<2^{364}$. Thus
\begin{equation}
 \overline E(n)<2^{356}n2^{-5n}\le2^{356-4n},\qquad n\ge1,\tag{2.2}
\end{equation}
where $n\le2^n$ follows by induction. The positive integer
$n_w=\lceil(d+356)/4\rceil$ is a termination witness. The defining set
for $n$ is nonempty and its minimum is obtained by finite exact rational
comparisons. Furthermore $\overline E$ decreases strictly: the ratio of
successive positive linear brackets is below two, while $2^{-5n}$
contributes $1/32$. At the selected $n$, $1000n\ge R$, so the actual AP6
geometry and (2.1) give $E_n\le\overline E(n)\le\eta$.
The physical area $A=4000000n^2$ is positive finite; $n\le80$ is not used.

% MAP: proof-2-nelson | P:84-112 S:156-175 S:237-240
For positive rational $v$, $2^h\ge h+1$ for integers $h\ge0$.
Hence $\operatorname{clog}_2(v)$ terminates no later than
$h=\lceil v\rceil$. Likewise $4^h\ge h+1$ makes
$\lceil14u^2 A\rceil$ a termination witness for $h_K$.
These rational bounds may be loose. Using $\sqrt{85}<10$ and $\pi>3$,
\[
 K_\Lambda^2=\frac{75\sqrt{85}}{2\pi^3}u^2 A
 <14u^2 A\le\overline K^2.
\]
Thus $K_\Lambda<\overline K$. Since $e>2$ and $\sqrt2<2$, the second
$K_{\mathrm{sharp}}$ branch is
$2K_\Lambda\sqrt{(1+\sqrt2)/e}<4\overline K$. Consequently
$K_{\mathrm{sharp}}\le\max\{12U_A,4\overline K\}$. Also
\[
 b_{\mathrm{Nelson}}=\frac{\pi^2}{48U_A}>b_0=\frac3{16U_A}>0,
 \qquad \frac{64}{e b_{\mathrm{Nelson}}}<\frac{32}{b_0}.
\]
The integer $k_w=\lceil32/b_0\rceil\ge1$ satisfies
$L_{10}(k_w)\ge1+k_w>32/b_0$, so $k\le k_w$ is a finite minimum.
The positive exponential remainder gives
$e^k>L_{10}(k)>32/b_0>64/(e b_{\mathrm{Nelson}})$, hence
$\log_+(64/(e b_{\mathrm{Nelson}}))<k$. Therefore the original,
volume-dependent Nelson bound obeys
\begin{align}
 H_2&\le\max\{12U_A,4\overline K,4k^2/b_0\}<H_N,\notag\\
 \sqrt{\Xi_2}&\le\sqrt{\overline\Xi_2}=\sqrt2 e^{H_2}
                         <2^{2H_N+1}.\tag{2.3}
\end{align}
Here $\log2>1/2$, because $e^{1/2}=\sqrt e<\sqrt3<2$, so
$e^z\le2^{2z}$ for $z\ge0$. The strict extra unit in the integer $H_N$
preserves strictness even when the first maximum branch is attained.
This uses the actual $\Xi_2$ at this physical area and its original
common-space fourth-chaos/$L^2$ interface. Finite-regulator stability and
Jensen alone would not establish the zero-regulator bound, and no
volume-uniform $\Xi_2$ is claimed.

% MAP: proof-2-p | P:114-125 S:241-241 S:333-337
The selected even integer $p=4H_N+2h_K+2d+8$ is at least four.
The actual regulator bound with $F_2\le4$ gives
\begin{equation}
 R_{\mathrm{phys}}<2^{2H_N+h_K+3-p/2}+2^{2H_N+2-p}.\tag{2.4}
\end{equation}
The first exponent is exactly $-d-1$; the second is
$-2H_N-2h_K-2d-6\le-d-1$. Their sum is at most $\eta$.
The first two-summand predicate is equivalent to
$p\ge4H_N+2h_K+2d+8$. This even value satisfies the second predicate,
so it is exactly the first even sufficient $p\ge4$.
Thus $0<a=2^{-p}<1$ and $R_{\mathrm{phys}}\le\eta$.
The actual $K_\Lambda$ contribution has not been replaced by one.

% MAP: proof-2-variance | P:127-137 S:79-86 S:118-130 S:987-994
For $p\ge4$, $\sqrt{1+2^{2p}}<2\cdot2^p$, whence
$\operatorname{asinh}(2^p)<\log(3\cdot2^p)<p+2$, using
$\log2<1$ and $\log3<2$. Since $2\pi>6$,
\begin{equation}c\le\overline c(a)<(p+3)/6\le p.\tag{2.5}\end{equation}
Only positive majorants use $p$; the physical $c$ remains in $V$ and the
Wick ordering. Finite-regulator Wick stability is the pointwise identity
$z^4-6cz^2+3c^2=(z^2-3c)^2-6c^2\ge-6c^2$.
The physical/residual weights are therefore bounded by $e^{6U_Ac^2}$,
the physical Jensen normalizer is at least one, and the residual Jensen
normalizer is at least $e^{-3U_Ac^2}$, exactly as in the original
conditional interface. None requires $U_A<3$.

% MAP: proof-2-aux | P:139-166 S:242-247 S:173-183
Let $j=\lceil18U_Ap^2\rceil$. By (2.5) and $\log2>1/2$,
$e^{9U_Ac^2}<2^j$. Since $\sqrt{15}<4$ and $p\ge1$,
$(4\sqrt{15}+12)c^{3/2}<32p^2$. The complete auxiliary bracket satisfies
\begin{equation}
 2U_A(4\sqrt{15}+12)c^{3/2}+F_1
 <64U_Ap^2+4c_r=B_{\mathrm{aux}}\le2^{h_{\mathrm{aux}}}.\tag{2.6}
\end{equation}
All arguments before their $\operatorname{clog}_2$ operations are finite
positive rationals. In particular $Z_s=j+h_{\mathrm{aux}}+p+d>1$, and
$h_s=1+\operatorname{clog}_2(Z_s)$ is exactly the first integer at least
one with $2^{h_s-1}\ge Z_s$. Thus $s=2^{h_s}\ge2$ is finite, and
\[
 R_{\mathrm{aux}}-\max(1000n,R)-s=1,
\]
because $R\le9<1000n$. A Euclidean displacement of length less than $s$
changes each coordinate by less than $s$. The strict $s$-neighborhoods
of both $\overline\Lambda$ and $K$ lie in $\Omega$, so the original
covariance-defect bound applies to both. Discarding the minimum in
$\rho_a(s)$ gives
\begin{equation}
 \sqrt{\rho_a(s)}\le\sqrt{\frac{\kappa_a^2}{2\pi}}e^{-s/\sqrt2}
 <2^p e^{-s/2}<2^{p-s/2}.\tag{2.7}
\end{equation}
The final strict inequality uses $e>2$, $s>0$ and no upper bound on $p$.
Multiply by the same physical exponential and both vector terms in (2.6):
\begin{equation}
 A_{\mathrm{aux}}<2^{j+h_{\mathrm{aux}}+p-s/2}\le2^{-d}=\eta.\tag{2.8}
\end{equation}

% MAP: proof-2-stabilizer | P:168-180 S:245-247 S:180-180
Since $O=4R_{\mathrm{aux}}^2>A$, $B_{\mathrm{ext}}=O-A>0$.
With $h_{\mathrm{stab}}=\operatorname{clog}_2(6Op^2)$ and
$t=\max\{34,j+h_{\mathrm{stab}}+d\}$, the actual stabilizer error obeys
\begin{equation}S<2^{j+h_{\mathrm{stab}}-t}\le\eta.\tag{2.9}\end{equation}
Also $\epsilon=2^{-t}\le2^{-34}<u$, because $2^{34}>10^{10}$.
Thus $\alpha=\epsilon$ and $U=u$ in every width formula.
All later stages are finite compositions of these integers and rational
searches, and $\operatorname{clog}_2(v)\le\lceil v\rceil$ covers each
remaining search. No equality oracle, canonical inner index or expansion
of an infinite sequence is used. This proves the rest of UDA.a, including
the four actual downstream errors.

% MAP: proof-3-point | P:182-202 S:351-370
\subsubsection{Arbitrary bridges, coercivity and full normalizers}
Fix any finite positive stage satisfying the premises, with $0<\epsilon<u$;
proof steps 3--8 use no selected-stage or bounded-$q$ consequence.
Fix an architecture, $N\ge1$, $j$, a branch $h$ and $\theta\in[0,1]$.
Before tilt, $v$ is a sum of $N$ independent centered process random
elements, $N-1$ with coefficient $1/\sqrt N$ and one with coefficient
$\theta/\sqrt N$. At almost every point $z$ each unscaled variance is
$1/4$; neuron and Gaussian fourth moments are $1/16$ and $3/16$,
respectively. Independence of summands and centering give, for real $\beta_i$,
\begin{align}
 \E\left(\sum_i\beta_i\xi_i(z)\right)^2&=\frac14\sum_i\beta_i^2,\notag\\
 \E\left(\sum_i\beta_i\xi_i(z)\right)^4
 &=\sum_i\beta_i^4\E\xi_i(z)^4
          +\frac6{16}\sum_{i<k}\beta_i^2\beta_k^2\notag\\
 &\le\frac3{16}\left(\sum_i\beta_i^2\right)^2.\tag{3.1}
\end{align}
Here the squared coefficient sum is $(N-1+\theta^2)/N\le1$.
This includes $N=1$, both endpoint amplitudes and every extreme $j$.
No spatial independence or Gaussian domination of neuronal joint moments
is used. Tonelli on finite $D$ gives a finite fourth $E$-norm moment.

% MAP: proof-3-coercivity | P:204-227 S:197-212 S:365-367
Write $x=\|v\|_4$ and take $c_*\ge c$. The quartic coefficient in $V$
lies in $[\alpha,U]$, the negative quadratic coefficient is pointwise at
most $B$, and the constant is nonnegative. H\"older and completion of the
square in $x^2$ give
\begin{equation}
 V(v)\ge\alpha x^4-B\sqrt O\,x^2
       \ge\alpha x^4/2-\frac{B^2O}{2\alpha}
       =\alpha x^4/2-L.\tag{3.2}
\end{equation}
Therefore $0<w(v)\le e^{L-\alpha x^4/2}\le e^L$.
Every $Z$ is finite and strictly positive, as the expectation of a strictly
positive finite random variable. Every term of $V$ is integrable by (3.1).
Discarding negative quadratic terms only in the upper bound yields
$\E V(v)\le3UO/16+3uc^2A\le J_z$. Jensen gives $Z\ge e^{-J_z}$.
For the actual joint law $\Pi$ of $(S_j,h)$ it follows that
\begin{equation}
 \frac{d\Pi}{d\Law(S_j,h)}\le e^{H-\alpha x(v)^4/2}\le e^H,
 \qquad\Pi(x(v)^4>z)\le\min\{1,e^{H-\alpha z/2}\}.\tag{3.3}
\end{equation}
These are full neural/Gaussian-tilt partitions. The physical Wick
normalizer $Z_a\ge1$ is neither used here nor asserted for $Z$.

% MAP: proof-4-field | P:229-245 S:368-370
\subsubsection{Joint tilted field and direction moments through sixteen}
If $T\ge0$ and $\PP(T>z)\le\min\{1,e^{a-z}\}$ with $a\ge0$, then
$T$ is stochastically dominated by $a+E_1$, where $E_1$ has unit
exponential law. This is a distribution-function comparison, rather than
a joint independence assumption. Integrating
$qz^{q-1}\PP(T>z)$, and integrating by parts for exponential moments,
gives for every positive integer $q$
\begin{equation}\|T\|_q\le a+(q!)^{1/q}.\tag{4.1}\end{equation}
For $T=\alpha x(v)^4/2$, (3.3) gives $a=H$. At order four,
$\|T\|_4\le H+24^{1/4}<H+4$, so
\begin{equation}
 \|x(v)\|_{L^{16}(\Pi)}\le[2(H+4)/\alpha]^{1/4}=X_{\mathrm{rel}}.\tag{4.2}
\end{equation}
Lyapunov's inequality gives every real order between one and sixteen.

% MAP: proof-4-gaussian | P:247-265 S:370-370
For the Gaussian direction, Minkowski exchanges
$L^{2k}(\text{probability})$ and $L^4(D)$ for $k\ge2$.
Its point variance is $1/4$, hence
\begin{equation}
 \E\|\gamma\|_4^{2k}
 \le O^{k/2}\E|\gamma(z)|^{2k}
 =O^{k/2}\frac{(2k-1)!!}{4^k}
 \le O^{k/2}(k/2)^k.\tag{4.3}
\end{equation}
The equality is the constant point-moment value; the first comparison is
an inequality and needs no spatial independence. At $k=1$, use
$\E\|\gamma\|_4^2\le\sqrt{\E\int_D\gamma^4}
=\sqrt{3O}/4\le\sqrt O/2$. The positive exponential series for
$T_\gamma=\|\gamma\|_4^2/(4\sqrt O)$ therefore gives
\begin{equation}
 \E e^{T_\gamma}\le1+\sum_{k\ge1}\frac{k^k}{8^kk!}
 \le\sum_{k\ge0}(e/8)^k<8/5<4.\tag{4.4}
\end{equation}

% MAP: proof-4-direction | P:266-281 S:368-370
The factorial bound $k!\ge(k/e)^k$ follows by integrating $\log$ on
$[1,k]$. Tonelli and positivity justify the series interchange.
Under the actual joint $\Pi$, density domination (3.3) gives
$\E_\Pi e^{T_\gamma}<4e^H$. Markov's inequality and $\log4<2$ imply
$\Pi(T_\gamma>z)\le\min\{1,e^{H+2-z}\}$.
Apply (4.1) at order eight and use $(8!)^{1/8}<8$:
$\|T_\gamma\|_{L^8(\Pi)}<H+10$. Consequently
\begin{equation}
 \|\|\gamma\|_4\|_{L^{16}(\Pi)}
 \le2O^{1/4}(H+10)^{1/2}=Y_{\mathrm{rel}}.\tag{4.5}
\end{equation}
For a neuron, $\|g^b\|_4=O^{1/4}/2$ deterministically, including after
tilt, and the same $Y_{\mathrm{rel}}$ bounds it. Lyapunov gives every
real order through sixteen. Conditional independence of $v,h$ after
tilt is never used.

% MAP: proof-5-potential | P:283-295 S:203-207 S:382-384
\subsubsection{Full Fr\'echet derivatives and differentiable partitions}
Quartic and quadratic integrals on $L^4(D)$ are continuous Fr\'echet
polynomials. For unit $L^4$ directions $h_i$, the first derivative is the
integral of $4k(z)v^3h_1-2b(z)vh_1$, with quartic coefficient $k(z)\le U$
and quadratic coefficient $b(z)\le B$. The constant $3uc^2A$
differentiates to zero. Higher derivatives have coefficients
$12k(z),24k(z),24k(z)$, with a quadratic contribution only at order two.
H\"older for four $L^4$ factors and the factor $\sqrt O$ for two give
\begin{align}
 \|DV(v)\|&\le4Ux^3+2B\sqrt O\,x=v_1(x),\notag\\
 \|D^2V(v)\|&\le12Ux^2+2B\sqrt O=v_2(x),\notag\\
 \|D^3V(v)\|&\le24Ux=v_3(x),\qquad
 \|D^4V(v)\|\le24U=v_4(x).\tag{5.1}
\end{align}
These are full multilinear norms.

% MAP: proof-5-bell | P:297-309 S:205-207 S:382-384
The exponential chain rule, written as a sum over partitions of
$\{1,\ldots,k\}$, expresses $D^kw$ as $w$ times products of derivatives
of $-V$. Norms give exactly $e_0,\ldots,e_4$. At order four the
multiplicities are $1,6,3,4,1$ for
$v_1^4,v_1^2v_2,v_2^2,v_1v_3,v_4$. Thus
\begin{equation}
 \|D^kw(v)\|_{\mathrm{op}}\le w(v)e_k(x),\qquad0\le k\le4.\tag{5.2}
\end{equation}
Each $e_k$ is a nonnegative-coefficient polynomial of degree at most $3k$.
By (3.2), its bound is a polynomial times $e^{L-\alpha x^4/2}$.
Each monomial is globally bounded: for $k>0$,
$\sup_{x\ge0}x^ke^{-\alpha x^4/2}=(k/(2\alpha e))^{k/4}$.
Hence $w$ and all derivatives through four are globally bounded.

% MAP: proof-5-partition | P:311-324 S:372-375
Differentiation of $Z(\theta)$ is dominated by
$M\|h\|_4/\sqrt N$ for a finite deterministic global derivative bound $M$.
The untilted direction has finite first $E$-norm moment. Dominated
convergence therefore gives a continuous derivative on every compact real
$\theta$ interval. In fact $Z$ extends $C^1$ to all real $\theta$ and stays
positive; therefore $\log Z$ also extends $C^1$. On $[0,1]$,
\begin{equation}
 (\log Z)'=-N^{-1/2}\E_\Pi DV(v)[h].\tag{5.3}
\end{equation}
Under the same joint $\Pi$, H\"older at order four for $x^3\|h\|_4$ and
at order two for $x\|h\|_4$, together with (4.2),(4.5), gives
\begin{equation}
 \E_\Pi|DV(v)[h]|\le4UX_{\mathrm{rel}}^3Y_{\mathrm{rel}}
       +2B\sqrt O X_{\mathrm{rel}}Y_{\mathrm{rel}}=K_1.\tag{5.4}
\end{equation}

% MAP: proof-5-ratios | P:326-335 S:377-380
Integrate (5.3) between $\theta$ and one on the neuron branch to obtain
the stated neuron/neuron ratio. At zero both branches have the same
$S_j$, so their partitions agree. Integrating from zero on both branches gives
\begin{equation}
 \frac{Z_{j,\gamma}(\theta)}{Z_{j,g}(1)}
 \le e^{(\theta+1)K_1/\sqrt N}\le e^{2K_1/\sqrt N}.\tag{5.5}
\end{equation}
If $N\ge N_0$, then $N\ge16K_1^2$, so the last expression is at most
$e^{1/2}=\sqrt e<2$. This also bounds the neuron/neuron ratio by two.
The denominator is the actual neuron endpoint partition, rather than a
physical normalizer.

% MAP: proof-6-product | P:337-350 S:384-394
\subsubsection{The complete centered product derivative and mixed moments}
For $G\in\mathcal U_4(r)$, (1.1) bounds
$D^q(G\circ\mathcal O_f)$ by $F^q$ for $1\le q\le4$, and its absolute
value by one. The deterministic $m_j$ lies in $[-1,1]$.
The full fourth product rule for $\Phi_j=(G\circ\mathcal O_f-m_j)w$ gives
\begin{equation}
 \|D^4\Phi_j(v)\|_{\mathrm{op}}
 \le w(v)\left[2e_4(x)+\sum_{q=1}^4\binom4qF^q e_{4-q}(x)\right]
 =w(v)P_F(x).\tag{6.1}
\end{equation}
The leading two copies of $e_4$ account for $|G-m_j|\le2$.
Every observation binomial term and every Bell multiplicity remains.
The lower-order product calculations followed by (3.2) likewise show
$\Phi_j\in C_b^4(E)$.

% MAP: proof-6-holder | P:352-363 S:390-396
Write $P_F(x)=\sum_{k=0}^{12}a_kx^k$ with $a_k\ge0$ and
$r_h=\|h\|_4$. For $k>0$, H\"older with conjugate exponents
$(k+4)/k$ and $(k+4)/4$ gives
\begin{align}
 \E_\Pi x^kr_h^4
 &\le(\E_\Pi x^{k+4})^{k/(k+4)}
          (\E_\Pi r_h^{k+4})^{4/(k+4)}\notag\\
 &\le X_{\mathrm{rel}}^kY_{\mathrm{rel}}^4.\tag{6.2}
\end{align}
At $k=0$ this is the fourth direction-moment bound. Since $k+4\le16$,
all moments used are established. Summing nonnegative coefficients yields
$\E_\Pi[P_F(x)r_h^4]\le P_F(X_{\mathrm{rel}})Y_{\mathrm{rel}}^4$.
The actual joint law is used throughout; dependent marginal expectations
are never multiplied.

% MAP: proof-7-cancellations | P:365-384 S:351-360 S:390-404
\subsubsection{Centered adjacent replacement and the normalized telescope}
Taylor's theorem for $\Phi_j$ at $S_j$ in direction $h/\sqrt N$, through
order three, has remainder
\begin{equation}
 R_h=N^{-2}\int_0^1\frac{(1-\theta)^3}{6}
 D^4\Phi_j(S_j+\theta h/\sqrt N)[h,h,h,h],d\theta.\tag{7.1}
\end{equation}
Global derivative bounds and the finite untilted fourth $E$-norm moment
justify exchanging expectations, conditioning on $S_j$ and integration.
Independence is used before tilt. Conditional on $S_j$, the first terms
vanish by zero Bochner means. The second terms agree by the premise for
every bounded bilinear form, applied to the continuous form
$D^2\Phi_j(S_j)$. The third terms vanish by global sign symmetry against
each continuous homogeneous cubic; the third $E$-moment guarantees
integrability. Constant terms agree. The forms are measurably evaluable
because $E$ is separable, the random elements are strongly measurable,
and derivatives are continuous. The full bilinear match is an explicit
premise, rather than an inference from point-pair covariance.

% MAP: proof-7-adjacent | P:386-401 S:398-404
By definition of $m_j$, $\E\Phi_j(S_j+\gamma/\sqrt N)=0$.
Consequently, for the actual normalized endpoint expectations,
\begin{equation}
 L_g(G)-L_\gamma(G)=\frac{\E R_g-\E R_\gamma}{Z_{j,g}(1)}.\tag{7.2}
\end{equation}
For branch $h$ and amplitude $\theta$, the expected absolute integrand
in (7.1), divided by $Z_{j,g}(1)$, is at most
\[
 \frac{Z_{j,h}(\theta)}{Z_{j,g}(1)}P_F(X_{\mathrm{rel}})Y_{\mathrm{rel}}^4.
\]
For $N\ge N_0$ each ratio is at most two by (5.5), and
$\int_0^1(1-\theta)^3/6\,d\theta=1/24$. The two branch remainders sum to
\begin{equation}
 \frac{4P_F(X_{\mathrm{rel}})Y_{\mathrm{rel}}^4}{24N^2}
 =\frac{K_4(F)}{N^2}.\tag{7.3}
\end{equation}

% MAP: proof-7-endpoints | P:403-414 S:398-404
Let $H_j$ be the law of the normalized $N$-summand mixture with $j$
Gaussian and $N-j$ neuronal summands, then tilted by the same $w$ and
its own normalizer. The neuron endpoint of $S_j$ is $H_j$ and the
Gaussian endpoint is $H_{j+1}$. This follows from the original independent
product copies and equality in law, hence also equality after common tilt
and normalization. Thus $H_0=P_N^b$ is the actual neural law. The normalized
sum of $N$ independent copies of the centered Gaussian $E$-valued element
$\gamma$ has exactly the $\gamma$ law: all finite continuous-linear-functional
laws agree, and these determine Borel probabilities on separable $E$.
Hence $H_N=Q_\epsilon$ is the actual full Gaussian tilt. Summing (7.3)
over $j=0,\ldots,N-1$ gives $K_4(F)/N$. The bound $|G|\le1$ supplies
the cap two. This proves the normalized part of UDA.c and completes UDA.b.

% MAP: proof-8-smoothing | P:416-430 S:133-142 S:402-407
\subsubsection{Ambient smoothing, including singular observation laws}
Apply the ambient $\R^r$ L0 premise to a unit BL $G$.
For $0<\zeta\le1$,
$G_\zeta/(\sqrt6\zeta^{-3})\in\mathcal U_4(r)$.
The two smoothing errors are each at most $\sqrt r\zeta$, so
\begin{equation}
 d_{\mathrm{BL}}(\mathcal L_f(P_N^b),\mathcal L_f(Q_\epsilon))
 \le\min\{2,\sqrt6K_4(F)/(N\zeta^3)+2\sqrt r\zeta\}.\tag{8.1}
\end{equation}
Put $A_0=\sqrt6K_4(F)$, $B_0=2\sqrt r$ and $N_*=3A_0/B_0$.
For $N\ge N_*$, choose $\zeta=(N_*/N)^{1/4}\le1$. This gives
$(4/3)(3A_0)^{1/4}B_0^{3/4}N^{-1/4}$, whose coefficient has fourth
power $(2048\sqrt6/27)r^{3/2}K_4(F)$.
For $N<N_*$ that expression exceeds $(4/3)B_0\ge8/3>2$, so the cap
two gives the same optimized bound. Both cases retain $N\ge N_0$.

% MAP: proof-8-degeneracy | P:432-436 S:463-465
Smoothing stays on the full original $\R^r$. If two tests coincide,
their coordinates coincide under every law, but (1.1), (6.1) and (8.1)
retain both coordinates. Touching, overlap and separation change neither
these arguments nor the moment estimates. No covariance inverse, density,
rank reduction or independence of interacting observations was assumed.

% MAP: proof-9-b | P:438-454 S:249-255
\subsubsection{Stage-dependent logarithmic enclosures for every finite accuracy}
Specialize the arbitrary finite-stage result to the universal selection,
with $c_*=p$. We have $p\ge4$, $t\ge34$, $A<O$, $O\le2^o$ and $d\le t$.
No old-domain bound on $n,U_A,p,O,t$ or $e_{\mathrm{new}}$ is used.
Since $u<1/6$ and $p\le2^{2p-1}$ for $p\ge4$,
\begin{equation}
 B=(2^{2p}-1)/2+6up<2^{2p-1}+p\le2^{2p}.\tag{9.1}
\end{equation}
The polynomial-exponential comparison holds at $p=4$ and persists because
the right side multiplies by four and $(p+1)/p\le5/4$.
Similarly $3p^2<2^{4p}$ holds at four and persists since
$((p+1)/p)^2\le25/16<16$. These are universal integer inequalities;
$6up<1$ is unnecessary and eventually fails.

% MAP: proof-9-moments | P:456-476 S:201-212 S:251-253
Write $l=l_{\mathrm{old}}=t+4p+o$ and $l_+=t+4p+o_+$.
The full coercivity and Jensen terms satisfy
\begin{align*}
 L&=B^2O/(2\alpha)<2^{t+4p+o-1}=2^{l-1},\\
 3uO/16&<O\le2^o\le2^{l-1},\\
 3up^2A&<3up^2O<2^{t+4p+o-1}=2^{l-1}.
\end{align*}
The last bound uses $A<O$, $3up^2<3p^2<2^{4p}$ and $t\ge34$.
Thus $J_z<2^l$ and $H=L+J_z<2^{l+1}\le2^{l_++1}$.
Since $l\ge50$, both $H+4,H+10<2^{l_++2}$.
Every term of the physical-area contribution $3up^2A$ remains, with
$A\le O$; no bounded-$U_A$ shortcut is hidden. Hence
\begin{align}
 X_{\mathrm{rel}}^4&<2^{2t+4p+o_++3},\notag\\
 Y_{\mathrm{rel}}^4=16O(H+10)^2&<2^{2t+8p+3o_++8}.\tag{9.2}
\end{align}
The actual $o$ is first bounded by $o_+$ in the second line.
Since $4T_{1/2}\ge2t$ at both parities of $t$, the definitions of $x,y$
give $X_{\mathrm{rel}}<2^x$, $Y_{\mathrm{rel}}<2^y$.
Strictness survives even when a ceiling has an integer argument.

% MAP: proof-9-k1 | P:478-495 S:409-410
The quadratic derivative coefficient obeys
\begin{equation}
 2B\sqrt O<2^{2p+\lceil o/2\rceil+1},\qquad
 2p+\lceil o/2\rceil+1\le2x.\tag{9.3}
\end{equation}
Indeed $2\lceil(o_++3)/4\rceil\ge(o_++3)/2\ge\lceil o/2\rceil+1$,
and the extra $2T_{1/2}$ is nonnegative. Since $4u<1$,
\begin{align}
 v_1(X_{\mathrm{rel}})&<2^{3x}+2^{2p+\lceil o/2\rceil+1+x}
                         <2^{3x+1},\notag\\
 K_1&<2^{3x+y+1}=2^{k_1-2}.\tag{9.4}
\end{align}
In particular $16K_1^2<2^{2k_1}$. The right side is an integer;
the ceiling of a real number strictly below integer $M$ is at most $M$.
Consequently
\begin{equation}
 N_0=\max\{1,\lceil16K_1^2\rceil\}\le2^{2k_1}
       <2^{2k_1+4}=2^{n_0}.\tag{9.5}
\end{equation}
This handles the ceiling without asserting strictness at that same exponent.

% MAP: proof-9-k4 | P:497-520 S:411-414
Let $W=2^{3x+5}\ge1$. At $X_{\mathrm{rel}}$, (9.3)--(9.4) give $v_1<W$.
For $v_2$ the two terms are bounded by $2^{2x+4}$ and $2^{2x}$; their
sum is below $2^{2x+5}\le W$. Similarly $v_3<2^{x+5}\le W$ and
$v_4<2^5\le W$. The complete Bell polynomials therefore satisfy
\[
 e_0\le1,\quad e_1\le W,\quad e_2\le2W^2,\quad e_3\le5W^3,
 \quad e_4\le15W^4.
\]
Each is at most $16W^4$. Keeping both centered copies of $e_4$ and all
four observation derivative terms, with $F\le\overline F\le6$, gives
\begin{align}
 P_F(X_{\mathrm{rel}})&\le P_{\overline F}(X_{\mathrm{rel}})\notag\\
 &\le16W^4\left[2+\sum_{j=1}^4\binom4j6^j\right]
 =16(1+7^4)W^4\notag\\
 &=38432W^4<2^{12x+36}.\tag{9.6}
\end{align}
Here $38432<2^{16}$ is the full vector coefficient; the smaller scalar
coefficient is not substituted. Combining $Y_{\mathrm{rel}}^4<2^{4y}$
and enlarging only upward by dropping the divisor six yields
\begin{equation}
 K_4(F)\le K_4(\overline F)<2^{12x+4y+36}=2^{k_4}.\tag{9.7}
\end{equation}
All volume, variance, stabilizer and physical-area dependence remains.
There is no use of $q\le12$, $n\le80$, $U_A<3$, $p<20000$,
$O<2^{76}$, $t<2^{36}$ or $e_{\mathrm{new}}<2^{40}$.
This proves UDA.c's universal exponent part from the arbitrary finite-stage
joint-law proof.

% MAP: proof-10-width | P:522-538 S:255-255 S:409-414
\subsubsection{Sufficient width and the actual five-arrow target triangle}
Take $\zeta=2^{-d-3}$, $e_{\mathrm{new}}=k_4+4d+12$ and
$N_{\mathrm{new}}=2^{e_{\mathrm{new}}}$. These are finite positive
mathematical numbers and $0<\zeta\le1$. Exact integer algebra gives
\begin{equation}
 e_{\mathrm{new}}-n_0=6x+2y+4d+38>0.\tag{10.1}
\end{equation}
Thus $N_{\mathrm{new}}>2^{n_0}>N_0$. Apply (8.1) with $\sqrt6<3$,
$\sqrt r\le2$ and (9.7):
\begin{align}
 \frac{\sqrt6K_4(F)}{N_{\mathrm{new}}\zeta^3}
 &<3\cdot2^{k_4-e_{\mathrm{new}}+3d+9}=3\eta/8,\notag\\
 2\sqrt r\zeta&\le4\cdot2^{-d-3}=\eta/2.\tag{10.2}
\end{align}
The actual neural-to-full-Gaussian-tilt distance is below $7\eta/8<\eta$.
This completes the selected-width part of UDA.c; steps 3--8 already prove
all its arbitrary finite-stage clauses.

% MAP: proof-10-target | P:540-557 S:144-193 S:416-424
The four quantities bounded in step 2 belong to the same actual
$Q_\epsilon,Q_0,Q_a,Q_{\mathrm{phys}},\nu^X$ chain in the source premises,
with the same physical and auxiliary squares. They do not depend on $N$.
Physical and full endpoint-reference normalizers have stayed separate;
$Q_\epsilon$ is the full Gaussian tilt, distinct from a physical
zero-regulator law. Strict AP6 identifies the last endpoint with the
original pressure target. On the full original $\R^r$, the triangle gives
\begin{align}
 d_{\mathrm{BL}}(\mathcal L_f(P_{N_{\mathrm{new}}}^b),\mathcal L_f(\nu^X))
 &\le E_{\mathrm{width}}+S+A_{\mathrm{aux}}+R_{\mathrm{phys}}+E_n\notag\\
 &\le\overline E(n)+4\eta\le5\eta\le\tau.\tag{10.3}
\end{align}
The width error has the strict margin (10.2). This holds for every
$q\ge3$ with $d=q+3$ and for the separate $1/5$ case with $d=5$.
It proves UDA.d for the unchanged target, architectures and observations.
The optional probability cap two does not repair an allocation or a failed stage.

% MAP: proof-11-raw | P:559-579 S:426-430
\subsubsection{The complete old coefficient and same-stage comparison}
First derive the old law coefficient at arbitrary finite stages from the
same full reference laws; bounded SVA/ACC acceptance is not extrapolated.
For $\Psi\in C_b^4(E)$ with fourth derivative norm bounded by $M$,
expand around the common $S_j$ through order three. The centering,
every-bilinear match and global sign-symmetry cancellations of step 7
apply. Each fourth remainder is at most $M\|h\|_4^4/(24N^2)$.
Tonelli and the point moments give
$\E\|g\|_4^4=O/16$, $\E\|\gamma\|_4^4=3O/16$.
After $N$ replacements,
\begin{equation}
 |\Law(F_N^b)(\Psi)-\Law(\gamma)(\Psi)|\le\frac{OM}{96N}.\tag{11.1}
\end{equation}
The normalized Gaussian iid sum has exactly the $\gamma$ law on separable
$E$, as above. No tilt, independence after tilt, normalizer premise or
threshold $N_0$ is required in this raw estimate.

% MAP: proof-11-normalization | P:581-611 S:214-225 S:426-430
The polynomial-exponential bound (3.2) and the complete derivatives
(5.1)--(5.2) give
\begin{align}
 \sup\|D^4w\|&\le\mathcal T(e_4)=M_Z,\notag\\
 \sup\|D^4[(G\circ\mathcal O_f)w]\|
 &\le\mathcal T\left(\sum_{j=0}^4\binom4jF^je_{4-j}\right)=M_A.\tag{11.2}
\end{align}
For $k>0$, differentiation shows
$\sup_{x\ge0}x^ke^{-\alpha x^4/2}=\beta_k$, and $\beta_0=1$.
Summing nonnegative coefficients bounds the polynomial supremum from above.
Lower derivatives are also bounded, so (11.1) applies legitimately.
Let $A_N,A_\gamma$ be weighted $G$ expectations and $Z_N,Z_\gamma$ their
actual full partitions. Then $\delta_A\le OM_A/(96N)$ and
$\delta_Z\le OM_Z/(96N)$. Both partitions are at least
$z=e^{-J_z}$ by the separate endpoint-reference Jensen estimate, while
$|A_\gamma/Z_\gamma|\le1$. The exact identity
\[
 \frac{A_N}{Z_N}-\frac{A_\gamma}{Z_\gamma}
 =\frac{A_N-A_\gamma}{Z_N}
 +\frac{A_\gamma}{Z_\gamma}\frac{Z_\gamma-Z_N}{Z_N}
\]
gives normalized error at most $(\delta_A+\delta_Z)/z=C_{\mathrm{old}}/N$,
using one power of $1/z$. It is also at most two. The same ambient L0
smoothing and optimizer as in step 8, with the cap two handling an
optimizer above the allowed smoothing scale, give
\begin{equation}
 d_{\mathrm{BL}}(\mathcal L_f(P_N^b),\mathcal L_f(Q_\epsilon))
 \le\min\{2,D_{\mathrm{old}}N^{-1/4}\},\quad
 D_{\mathrm{old}}^4=\frac{2048\sqrt6}{27}r^{3/2}C_{\mathrm{old}},
 \quad N\ge1.\tag{11.3}
\end{equation}
This proves the first part of UDA.e at all finite stages.
Nonnegative coefficients ensure that replacing $F$ by $\overline F$ and
$c$ by $p$ only enlarges the old coefficient; $V$ itself is unchanged.

% MAP: proof-11-polynomial | P:613-642 S:257-262 S:432-435
At universal selected stages let
$h_{\mathrm{root}}=\operatorname{clog}_2(2R_{\mathrm{aux}})$,
$v=v_{\mathrm{old}}=2p+h_{\mathrm{root}}+2$ and $V_0=2^v$.
Then $\sqrt O=2R_{\mathrm{aux}}\le2^{h_{\mathrm{root}}}$.
By (9.1), the coefficient sums of $v_1,v_2$ are strictly below
$4+2^{2p+h_{\mathrm{root}}+1}$ and $12+2^{2p+h_{\mathrm{root}}+1}$,
respectively. Since $p\ge4$, $h_{\mathrm{root}}\ge0$, each is at most
$V_0$. Those of $v_3,v_4$ are below $24<V_0$.
The Bell coefficient sums are consequently bounded by
$1,V_0,2V_0^2,5V_0^3,15V_0^4$, each at most $16V_0^4$.
The polynomial for $M_A+M_Z$ has coefficient sum
\begin{equation}
 P_{\overline F}(1)\le16V_0^4
 \left[1+\sum_{j=0}^4\binom4j6^j\right]
 =38432V_0^4<2^{4v+16}.\tag{11.4}
\end{equation}
This is the whole vector polynomial, including both $e_4$ copies, of
degree at most twelve. For $1\le k\le12$, $\alpha=2^{-t}$ and $e>2$ give
\[
 \beta_k=\left(\frac{k}{2\alpha e}\right)^{k/4}
 <[2^{t+4}]^{k/4}\le2^{3t+12}.
\]
This also covers $\beta_0=1$. Thus the complete polynomial part is below
$2^{q_{\mathrm{old}}}$, where $q_{\mathrm{old}}=4v+3t+28$.
Step 9 gave the full $H=L+J_z<2^{l_{\mathrm{old}}+1}$ without a bounded-area
shortcut. Since $e^H\le2^{2H}$,
\begin{equation}
 C_{\mathrm{old}}=(O/96)e^H\sum_k a_k\beta_k
 <2^{\,2^{l_{\mathrm{old}}+2}+q_{\mathrm{old}}+o}.\tag{11.5}
\end{equation}
Neither Jensen remainder, the physical-area term nor the $O/96$ prefactor
has been discarded downward. The same-stage old remainder guards also
hold: $o\le l_{\mathrm{old}}-1$ and
$\operatorname{clog}_2(3U_Ap^2)\le l_{\mathrm{old}}-1$, by step 9's strict bound.

% MAP: proof-11-old-width | P:644-667 S:434-437
The vector coefficient in (11.3) is below $2^{11}$:
$(2048\sqrt6/27)<256$ follows from $\sqrt6<5/2$, and $r^{3/2}\le8$.
Hence $N\ge2^{11+4d}C_{\mathrm{old}}$ suffices for old first-arrow error
at most $\eta$. The old symbolic guard must retain both
$2^{l_{\mathrm{old}}+2}$ and $H_{\mathrm{small}}=q_{\mathrm{old}}+o+11+4d$.
Since $O\ge1$ gives $h_{\mathrm{root}}\le o$, and $t\ge d$, exact expansion yields
\begin{align}
 H_{\mathrm{small}}&=8p+4h_{\mathrm{root}}+3t+o+47+4d\notag\\
 &\le8p+5o+7t+47\le7l_{\mathrm{old}}+47.\tag{11.6}
\end{align}
For every integer $l\ge5$, $7l+47<2^{l+2}$: it holds at five and its
strict gap grows because $2^{l+2}>7$. Here $l_{\mathrm{old}}\ge50$.
Thus $\operatorname{bitlen}(H_{\mathrm{small}})\le l_{\mathrm{old}}+2$.
The first candidate $w=l_{\mathrm{old}}+2$ fails the explicit strict
predicate; the next, $w_{\mathrm{old}}=l_{\mathrm{old}}+3$, satisfies it
and the unchanged guard. Equations (11.5)--(11.6) give
\begin{equation}
 \log_2[2^{11+4d}C_{\mathrm{old}}]
 <2^{l_{\mathrm{old}}+2}+H_{\mathrm{small}}
 <2^{l_{\mathrm{old}}+3}=\log_2N_{\mathrm{old}}.\tag{11.7}
\end{equation}
Thus $N_{\mathrm{old}}$ is sufficient at these same stages, with first-arrow
error at most $\eta$ by (11.3). The coefficient is unchanged, and no
necessary-width lower bound is proved.

% MAP: proof-11-order | P:669-686 S:438-442
For integer $a$, $\lceil a/4\rceil\le(a+3)/4$. Expand
\begin{align*}
 e_{\mathrm{new}}
 &=16T_{1/2}+20p+12\left\lceil\frac{o_++3}{4}\right\rceil
     +4\left\lceil\frac{3o_++8}{4}\right\rceil+4d+48\\
 &\le16T_{1/2}+20p+6o_++4d+77.
\end{align*}
Since $T_{1/2}\le(t+1)/2$, $o_+\le o+78$ and $d\le t$,
\begin{equation}
 e_{\mathrm{new}}\le12t+20p+6o+553\le12l_{\mathrm{old}}+553.\tag{11.8}
\end{equation}
For every integer $l\ge8$, $12l+553<2^{l+3}$: it holds at eight and the
gap grows because $2^{l+3}>12$. Hence
\begin{equation}
 e_{\mathrm{new}}<2^{l_{\mathrm{old}}+3}=2^{w_{\mathrm{old}}},\qquad
 N_{\mathrm{new}}=2^{e_{\mathrm{new}}}<2^{2^{w_{\mathrm{old}}}}=N_{\mathrm{old}}.
 \tag{11.9}
\end{equation}
This proves UDA.e universally. It compares sufficient mathematical integers;
there is no fixed bound $e_{\mathrm{new}}<2^{40}$.

% MAP: proof-12-history | P:688-695 S:269-275 C:3-8
\subsubsection{Exact recovery of each formerly supported stage and formula}
Version 2 gives the complete legacy comparator explicitly. The comparisons
here concern those definitions, so their equality needs no external ACC/SVC
statement. Version 1's original statement and proof remain preserved; its
standalone UDA.f omitted the comparator definitions. The repair specifies
the theorem conclusion fully without adding an analytic premise.

% MAP: proof-12-compact | P:698-711 S:277-289 S:451-453
On the old tolerance set, $d\in\{5,\ldots,15\}$. At $n=80$ the spatial
bracket is $159776/961<256$, so
$\overline E(80)<2^{-38}\le\eta$. The unbounded minimum $n$ therefore
equals the old first witness in $1,\ldots,80$.
Then $U_A=n^2/2500\le64/25<3$ and
\[
 14u^2A\le358400000000/10^{20}<1.
\]
Thus $h_K=0$ and $\overline K=1$ exactly. Also $b_0>1/16$ and
$32/b_0<512$. The degree-ten rational partial at seven exceeds $512$,
so the unbounded minimum $k$ is the same old first witness in $1,\ldots,7$.
Since $12U_A<36$ and $4k^2/b_0<4\cdot49\cdot16=3136$, the identical
integer formula gives $H_N\le3137$. The regulator formula reduces to
$p=4H_N+2d+8\le12586$, exactly the legacy first even two-predicate
witness. This proves equality of $n,k,H_N,p$, rather than merely another
sufficient choice.

% MAP: proof-12-outer | P:713-721 S:290-308 S:452-453
At these bounds $j=\lceil18U_Ap^2\rceil\le21600000000<2^{35}$ and
$B_{\mathrm{aux}}<76800000008<2^{37}$.
Hence $h_{\mathrm{aux}}\le37$ and
$j+h_{\mathrm{aux}}+p+d<2^{35}$. The original first auxiliary predicate
gives $h_s\le36$. Step 2 showed that the new formula is exactly this same
minimum, so $h_s,s,R_{\mathrm{aux}},O,o,h_{\mathrm{stab}},t$ recover the
legacy stages exactly. Moreover
$R_{\mathrm{aux}}=1000n+s+1<2^{37}$, $O<2^{76}$, $o\le76$, and thus
$o_+=78$. From $6Op^2<2^{109}$ we get $h_{\mathrm{stab}}\le109$ and
$t=\max\{34,j+h_{\mathrm{stab}}+d\}<2^{36}$.
Every common value uses the same literal predicate and tie rule;
no finite diagnostic-grid inference is required.

% MAP: proof-12-exponents | P:723-733 S:310-316 S:455-458
With $o_+=78$, the ceilings are $\lceil81/4\rceil=21$ and
$\lceil242/4\rceil=61$. Therefore
\begin{align*}
 x&=T_{1/2}+p+21,&y&=T_{1/2}+2p+61,\\
 k_1&=3x+y+3=4T_{1/2}+5p+127,\\
 n_0&=2k_1+4=8T_{1/2}+10p+258,\\
 k_4&=12x+4y+36=16T_{1/2}+20p+532,\\
 e_{\mathrm{new}}&=k_4+4d+12=16\lceil t/2\rceil+20p+4d+544.
\end{align*}
These are exactly the legacy SVA/SVC exponents and width at those identical
stages, including odd-$t$ rounding. Step 11 also gives the identical old
index $w_{\mathrm{old}}=l_{\mathrm{old}}+3$, with the same $l_{\mathrm{old}}$.

% MAP: proof-12-every-variable | P:734-749 S:444-461
For completeness, $d_L=d$, $\eta_L=\eta$ by their definition.
The same spatial predicates and $n\le80$ give $n_L=n$, hence $A_L=A$,
$U_{AL}=U_A$, $b_{0L}=b_0$. The identical finite Nelson predicates and
$k\le7$ give $k_L=k$; $h_K=0$ gives $H_L=H_N$. The first-even predicates
give $p_L=p$, $a_L=a$. Successive substitution gives
\begin{gather*}
 j_L=j,\ B_{\mathrm{aux},L}=B_{\mathrm{aux}},\
 h_{\mathrm{aux},L}=h_{\mathrm{aux}},\ h_{sL}=h_s,\ s_L=s,\\
 R_{\mathrm{aux},L}=R_{\mathrm{aux}},\ O_L=O,\ o_L=o,\
 h_{\mathrm{stab},L}=h_{\mathrm{stab}},\ t_L=t,\ \epsilon_L=\epsilon,\\
 \Lambda_L=\Lambda,\quad D_L=D,\quad\Omega_L=\Omega.
\end{gather*}
Substitution in the old-index definitions gives
$h_{\mathrm{root},L}=h_{\mathrm{root}}$,
$v_{\mathrm{old},L}=v_{\mathrm{old}}$,
$q_{\mathrm{old},L}=q_{\mathrm{old}}$,
$l_{\mathrm{old},L}=l_{\mathrm{old}}$ and
$H_{\mathrm{small},L}=H_{\mathrm{small}}$.
Their guard succeeds by step 11 and the first strict predicate selects
$w_{\mathrm{old},L}=w_{\mathrm{old}}$, so $N_{\mathrm{old},L}=N_{\mathrm{old}}$.
Finally, the displayed legacy relative formulas and the preceding algebra give
\[
 (T_L,x_L,y_L,k_{1L},n_{0L},k_{4L},\zeta_L,e_L,N_L)
 =(T_{1/2},x,y,k_1,n_0,k_4,\zeta,e_{\mathrm{new}},N_{\mathrm{new}}).
\]
Every variable in UDA.f is defined and its equality proved. The separate
$1/5$ case uses $d=5$ without conflating its request with a dyadic one.
This completes UDA.f.

% MAP: proof-13-evidence | P:751-758 C:33-35
\subsubsection{Evidence, retained adverse facts and scope}
Steps 1--12 are symbolic proofs over all declared integers $q$; a finite
case list supplies no universal quantifier. The source author records a
portable audit reconstructing rational stages and checking arithmetic and
enclosures beyond $q=12$, in noncompact volume/variance regimes and at both
$t$ parities. Those recorded runs are author evidence, distinct from fresh
independent review or field-law simulation. Those historical runs are attributed to their recorded roles. Any current
manuscript replay is separately identified in the companion execution receipt.

% MAP: proof-13-counterexample | P:760-765
The product neuron shows why Gaussian joint-moment domination is not a
permitted shortcut. At rectangle corners $(0,0),(a,0),(0,a),(a,a)$ its
fourfold product is exactly $1/16$, whereas the Gaussian pairing sum is
$(2e^{-2}+e^{-4})/16<1/16$. The proof uses one-point moments before
tilt and coercivity/density domination under the actual joint tilt.
The given full bilinear interface remains explicit.

% MAP: proof-13-limits | P:767-776 S:463-484 C:34-42
All original source/classical/strict-AP6/NR1/unknown-$g_*$ qualifications
remain, as do the physical-versus-endpoint covariances and normalizers.
Old-domain bounded quantities and numerical guards are not theorem premises.
Existing ACC/SVC software and certificates keep their original domains
and identities. The mathematical derivation supplies the every-$q$ implication.
The separate UDC contract below gives its finite partial software realization.
Canonical least inner indices, practical/minimal/necessary costs, full-field
rates, global neural Schwartz control, broader models/test families and
novelty/priority remain separate. Both widths count whole-plane process
neurons with countably many internal variables. The accepted source decisions
and the fresh manuscript reviews have distinct scopes.\hfill$\square$


% Literal inherited interface, typeset from the frozen statement only.
% MAP: source-intro | S:472-488 C:37-42
\subsection{Literal inherited source interface and attribution}
\label{uda:source}
The accepted statement includes a statement-only transcription of the
inherited SVA source appendix. The direct finite interfaces stated earlier
are its usable premises. The original private source-availability map identifies
which exact historical bodies were supplied in that integration packet.
Historical paths retain their attribution and source-access limits; supplying
a body does not newly prove it, authenticate its authorship or remove an
original qualification.
Source mesh $\varepsilon$ and flow time $s$ are local to this subsection,
distinct from the outer stabilizer $\epsilon$ and selected auxiliary radius.
Gaussian coupling fields are distinct from the field state space $X$.
The historical global EB1 hypothesis involving $g_*$ is not asserted;
strict AP6 supplies the target route. The old EB1 coefficient is retained
for scope and rederived at every finite stage in proof step 11.
Its bounded SVA acceptance does not supply a universal premise.
The source records the inherited ADMIT statement digest
\par\noindent{\small\texttt{3e361ceb474a571783015c802c75a2314a58005052f07eb46e09bfb18034c997}}\par

% MAP: source-common | S:492-506
For $m>0$, let $P_m$ be the real whole-plane Gaussian free field with
covariance $(-\Delta+m^2)^{-1}$ and Fourier measure $dp/(2\pi)^2$.
For a centered axis-parallel square $B=(-b,b)^2$, $b>0$, let $J_B$ be its
coefficient-one fourth Wick integral on this same Gaussian space, with the
same mass and no factorial. Put
\[
 U_B=uJ_B,\quad Z_B(h)=\E_{P_m}e^{-U_B+h\phi(\mathbf1_B)},\quad
 \nu_B=e^{-U_B}P_m/Z_B(0).
\]
The Gaussian $L^2$ extension defines $\phi(\mathbf1_B)$, and every subsequent
tilt uses that same random variable. The source in $Z_B(h)$ is confined to
$B$; its boundary is Lebesgue null. The following remain named premises
with their original authorship, source qualifications and accepted scopes.

% MAP: source-p1 | S:508-523
\paragraph{P1: finite ferromagnetic DSS and SK-q.}
On the shifted torus
$T=\varepsilon(\Z+1/2)^2/L\Z^2$, with $L/\varepsilon$ even and at least
four, use physical inner product $\varepsilon^2\sum$, operator
$A=-\Delta^\varepsilon+m^2$, and the plane lattice variance
$c_\varepsilon$. For every profile $q:T\to[0,1]$, the law
\[
 d\nu\propto\exp\left[-u\varepsilon^2\sum_xq_x
                {:}\psi_x^4{:}_{c_\varepsilon}\right]dP_{A^{-1}}
\]
satisfies the accepted SK-q inequality. At every finite real physical
field $H$, its tilt by $\exp(\varepsilon^2\sum_xH_x\psi_x)$ satisfies
\[
 0\le\operatorname{Cov}_{\nu^H}(\psi_x,\psi_y)
       \le\E_\nu(\psi_x\psi_y).
\]
The analogous clauses hold with added physical mass $1/s$ for every finite
$s>0$. These finite DSS-q/SKQ clauses include Gaussian sites and retain the
frozen BD/BBD/Griffiths--Simon, Ising-FKG and Robbins--Stirling
qualifications. They have no $g_*$ premise.

% MAP: source-p2 | S:525-543
\paragraph{P2: exact finite-positive-scale $P^*$.}
If
\[
 0<\varepsilon\le\min(1/4,18/(5m)),\quad
 L\ge\max(1,\sqrt2/m),\quad L/\varepsilon\in2\N,\quad
 L/\varepsilon\ge4,\quad0<u\le m^2/8000,
\]
then core $P^*$ using P1, $S_{\max}=\infty$, $\Phi=1$ and $\kappa=0$
gives, for finite $s>0$ and every profile $q:T\to[0,1]$,
\begin{align*}
 \chi_{s,q}&=\max_x\varepsilon^2\sum_y S_{s,q}(x,y),&
 S_{s,q}(x,y)&=\E_{\mu_{s,q}}(\psi_x\psi_y),\\
 d\mu_{s,q}&\propto\exp\left[-u\varepsilon^2\sum_xq_x
                 {:}\psi_x^4{:}_{c_\varepsilon}\right]dP_{(A+1/s)^{-1}},\\
 \chi_{s,q}&\le\frac{s}{1+m^2s}+\frac{2s}{2+m^2s}.
\end{align*}
This combines the accepted $\|S-C_s\|_{1,0}\le\tau_s$ with
$\tau_s=2s/(2+m^2s)$ and the exact Gaussian row sum. P2 includes no
original-law endpoint conclusion.

% MAP: source-p3 | S:545-561
\paragraph{P3: fixed-square Gaussian/Wick transfer.}
At fixed $B=(-b,b)^2$, take commensurate $\varepsilon=b/N\to0$ on the
shifted plane lattice. Accepted T1--T4/GQ/W(b)/W-laws place the plane
Gaussian field, the variables
\[
 J_B^\varepsilon=\varepsilon^2\sum_{x\in B}{:}\psi_x^4{:}_{c_\varepsilon},
 \qquad Y_f^\varepsilon=\varepsilon^2\sum_x f(x)\psi_x,
\]
and their continuum counterparts on one Gaussian space.
For every real compact smooth $f$, $J_B^\varepsilon\to J_B$ and
$Y_f^\varepsilon\to\phi(f)$ in Gaussian $L^2$.
For each fixed finite $p\ge1$,
$\sup_{\varepsilon\ \mathrm{sufficiently\ small}}\E e^{-puJ_B^\varepsilon}<\infty$.
The continuum density has positive finite normalizer and finite Gaussian
$L^p$ density at every finite $p$. These finite-volume statements have no
weak-coupling or $g_*$ premise. Coefficient-one Wick and finite-law
identifications are precisely those of the original construction.

% MAP: source-p4 | S:563-567
\paragraph{P4: fixed-mesh Gaussian periodization.}
At fixed $\varepsilon$, torus Gaussian covariances converge entrywise on
every finite set of fixed sites to the plane lattice covariance as
$L\to\infty$ through admissible even multiples. The plane variance used
for Wick ordering is the same on both sides. This is the accepted H1-PREQ
Gaussian periodization/envelope input.

% MAP: source-p5 | S:569-585
\paragraph{P5: original Fr\"ohlich--Simon source facts.}
Use the original same-mass whole-plane free-boundary convention with
$P_h(X)=uX^4-hX$, $u>0$. The free pressure $\alpha_m(P_h)$ is the
normalized log-partition limit. Its dyadic-square realization is attributed
to \citet[Proposition 2.1 and the remark after (2.3), pp.~502--503]{FS1977}.
The equality of the two original maximal states
$\nu_{m,u,+}=\nu_{m,u,-}$ is equivalent to differentiability of
$h\mapsto\alpha_m(P_h)$ at zero by the original Corollary 4.3(ii), p.~510.
Theorem 5.2, p.~512, then gives ordinary exhaustion convergence of the
classical free finite-volume states to their common state, read via the
compact-test generating functional at every real frequency.
These are attributed published facts; the source reports a source match,
and their original proofs are not reproved here. The historical
\texttt{SOURCE-MAP.md/json} record paths, hashes, locations and primary-PDF
access in the upstream packet; those files are not supplied to this role.
P5 uses the converse in the actual Corollary 4.3(ii); the manuscript's
narrower C1(iii) formulation by itself would not suffice.

% MAP: source-p6-a | S:587-602
\paragraph{P6: additional explicit bridge inputs.}
For AP5--AP6 retain the following frozen accepted inputs at their actual
parameter scopes. First, retain the all-$m,u>0$ physical/classical identity
at historical \texttt{CT.tex:2135--2166}, Gaussian common-space/Wick laws,
and the original complete local $H^{-2}$ field-space, cutoff/Fourier
coordinate, Borel-determining and compactness definitions.
Second, the all-semibounded-polynomial FS C1(ii) local density and positive
boundary factorization at historical \texttt{CT.tex:701--711} give
\[
 dv_{m,u}|_{\mathcal C_B}
 =e^{-uJ_B-\alpha_m(uX^4)|B|}g_{\partial B}\,dP_m,
\]
with positive boundary-measurable $g$ and finite local Gaussian $L^p$
densities. Either recorded boundary-algebra reading is retained.
Third, retain the finite DSS, SKQ, $P^*$, H1-PREQ, H1, MIX, REG, smoothed-CT, FLOW,
TP and T interfaces reproduced at historical \texttt{CL1-ST.md:138--1828},
qualified by their separate decisions and later source-normalization
correction. That file's new CL1 conclusion is excluded as a premise.
FLOW's own cap is $0<u\le m^2/8000$; TP/T1--T4 are fixed-square transfers.

% MAP: source-p6-b | S:603-619
Finally, retain the raw fixed-stage NR1/L0 and Gaussian/common-space/Wick
inequalities of the historical EB-ST accepted-input appendix, at $m,u>0$
and fixed positive regulator/stabilizer/finite geometry. In particular
NR1 gives $\delta_Z\le d_D M_Z/(96N)$,
$\delta_A\le d_D M_A/(96N)$, with $Z_Q\ge z$.
L0 smooths unit BL functions with error at most $\sqrt r\delta$ and
order-through-four operator bounds at most $\sqrt6\delta^{-3}$ for
$0<\delta\le1$. Keep the sharper-author/larger-blind-core scope, Gaussian
mass tilt, independent covariance-defect coupling and actual physical
common-space Wick bounds/normalizers. These are finite raw statements.
Old EB1/FV/CL1 conclusions with global $g_*$ hypotheses are excluded from
this bridge's premises. Prior argument ports give author provenance,
rather than a proof a blind solver must read. The source records the
literal statements, separate decision hashes and ranges in historical
\texttt{support/UPSTREAM-STATEMENTS.md}, \texttt{support/} decisions and
\texttt{SOURCE-MAP.json}; their bodies are unavailable to this role.
P6 excludes AHKZ C2(ii) and C2(iii).

% MAP: source-ap1 | S:621-633
\paragraph{AP1: original finite-lattice endpoint and source bound.}
Under P1--P5 and $0<u\le m^2/8000$, for all lattices, tori and profiles
of P2 the original zero-field law obeys
\[
 \max_x\varepsilon^2\sum_y\E_\nu(\psi_x\psi_y)\le3/m^2.
\]
For real vectors $f$ on $T$ and real $t$,
\[
 0\le\log\E_\nu\exp\left[t\varepsilon^2\sum_xf_x\psi_x\right]
       \le\frac{3t^2}{2m^2}\varepsilon^2\sum_xf_x^2.
\]

% MAP: source-ap2 | S:635-643
\paragraph{AP2: actual free-boundary continuum source bound.}
For centered squares $B$, real compact smooth $f$ and real $t$,
\[
 0\le\log\E_{\nu_B}e^{t\phi(f)}\le\frac{3t^2\|f\|_2^2}{2m^2}.
\]
The same holds for real $f\in L^2(\R^2)$ via the Gaussian extension.
In particular $0\le\log[Z_B(h)/Z_B(0)]\le3h^2|B|/(2m^2)$ for real $h$.

% MAP: source-ap3 | S:645-653
\paragraph{AP3: pressure and the exact original FS target.}
With P5's pressure,
\[
 0\le\alpha_m(uX^4-hX)-\alpha_m(uX^4)\le3h^2/(2m^2),\quad h\in\R.
\]
It is differentiable at zero with derivative zero. The actual two original
FS maximal states coincide; write their common value as $v_{m,u}$.
For every $\ell>0$, finite real compact-smooth tuple $f$, and centered
$\Lambda_n=(-n\ell,n\ell)^2$, their joint test laws under the original
classical free $\nu_{m,u,\Lambda_n}$ converge to those under $v_{m,u}$.

% MAP: source-ap4 | S:655-667
\paragraph{AP4: reference-pair identity and the fixed instance.}
For $M>0$, set $s=m/M$ and $\lambda=u/s^2$.
Then $\lambda/M^2=u/m^2\le1/8000$, and under the manuscript's exact
dilation $(T_s\psi)(f)=\psi(s^{-2}f(\cdot/s))$,
\[
 v_{m,u}=(T_s)_*v_{M,\lambda}.
\]
The source-defined target is therefore the original common FS maximal
state and the same reference/dilation expression as in the construction,
with identical physical finite-volume approximants, rather than an arbitrary
subsequential limit. The fixed $m=1,u=10^{-10}$ has this original-target
existence/equality/free-limit witness under P1--P5 since $10^{-10}<1/8000$.

% MAP: source-ap5 | S:669-677
\paragraph{AP5: physical source identity, field version and interactingness.}
Also assuming P6, the physical zero-regulator finite-volume laws have
exactly the classical $\nu_B$ compact-joint laws. They converge weakly on
the original $X=H^{-2}_{\mathrm{loc}}$ to a unique probability $\nu^X$
representing $v_{m,u}$; all compact-test polynomial moments converge and
are uniformly integrable at every fixed positive order. For every $u>0$
in this ratio range, $v_{m,u}$ is non-Gaussian by the retained positive local
quartic boundary-density structure. This interactingness check is separate
from equality of the two maximal states.

% MAP: source-ap6-a | S:679-699
\paragraph{AP6: strict numerical-budget applicability.}
Also assume P6, $0<u/m^2<1/8000$ and a fixed finite compact-smooth tuple.
Finite-stage definitions/constants are precisely the inherited
\texttt{EB-ST.md:7--55}. For $K\subset[-R_{\mathrm{obs}},R_{\mathrm{obs}}]^2$
use the exact \texttt{CL1-ST.md:47--68} coefficient with
$R=R_{\mathrm{obs}}$, its H1/MIX constants furnished by the separately
accepted finite interfaces after their actual premises are discharged.
The formula has $\epsilon_*=\min(1/4,1/m)$, $K_{\mathrm{cov}}=1$,
$c_{\mathrm{CT}}=16.2$ and $\kappa=m/256$; AP6 itself supplies no
numerical $C_f$. If $K\subset\Lambda_{n_0}$ open, the actual continuum
finite-box interpolation obeys, for $n\ge n_0$, $t\in[0,1]$ and
$G\in C_b^1$ with separately bounded value and Euclidean gradient,
\begin{align*}
 &\left|\operatorname{Cov}_{\nu_{n,t}}
 \left(G(Y_f),u[J_{\Lambda_{n+1}}-J_{\Lambda_n}]\right)\right|\\
 &\hspace{1cm}\le C_f u\,4\ell^2(2n+1)e^{-\kappa\operatorname{dist}(K,S_n)}.
\end{align*}
For $n\ell\ge R_{\mathrm{obs}}$, define
\[
 q=e^{-\kappa\ell},\quad
 T_n=q^n\left[\frac{2n+1}{1-q}+\frac{2q}{(1-q)^2}\right],\qquad
 E_n=4C_f u\ell^2e^{\kappa R_{\mathrm{obs}}}T_n.
\]

% MAP: source-ap6-b | S:701-718
The actual physical observation law then has BL$_1$ distance at most
$E_n$ from the same source target $v_{m,u}$, equivalently $\nu^X$.
For both original process-neuron architectures and every positive finite
stage in EB-ST, the same five-term expression holds:
\begin{align*}
 &d_{\mathrm{BL}}\left(\mathcal L_f(P^b_{N,\epsilon,\Omega,a;\Lambda_n}),
                                 \mathcal L_f(\nu^X)\right)\\
 &\hspace{1cm}\le\min\{2,\min(2,D_{\mathrm{law}}N^{-1/4})
              +S+A_{\mathrm{aux}}+R_{\mathrm{phys}}+E_n\}.
\end{align*}
Here $D_{\mathrm{law}},S,A_{\mathrm{aux}},R_{\mathrm{phys}}$ are the exact
finite-stage formulas of historical \texttt{EB-ST.md:32--44,72--98}, with
the required geometry and physical covariance. BL$_1$ retains its two
separate constraints. With certified coefficient/test/support envelopes,
the EC1.e reverse-stage argument applies using the distinct
\texttt{FS-pressure-target-v1} certificate in place of an old $g_*$
inequality. It supplies no unsupplied envelope and neither amends prior
acceptance nor certifies the original observable-independent sequence.

% MAP: source-eb-intro | S:721-732
\subsubsection{Inherited EB1 finite normalized-law and removal interfaces}
The recorded EB1 statement digest is
\par\noindent{\small\texttt{d6b400d4d47c6f8e28f7344b678b2ed2114504336c1635585a192bcf25319c83}}\par
Definitions and EB1.a--d,f retain their finite-stage scopes, with AP6
supplying the explicit-cap final arrow. EB1.e,g are unused and omitted;
excluded E4/r0 material is absent. The historical global header fixed
$m,u,\ell>0$ with $0<u/m^2<\min(g_*,1/8000)$ and an admissible reference
pair/target. Its $g_*$ condition is preserved as historical scope, rather
than assumed for the present strict-AP6 route. The physical squares
$\Lambda_n=(-n\ell,n\ell)^2$ restrict the interaction; all field laws
occupy the whole plane. The observations, support and BL$_1$ convention
are those already defined. The physical multiplier is
$(m^2+\omega^2+\xi^2+a^2\omega^2\xi^2)^{-1}$, $0<a<1/m$,
$\kappa_a^2=a^{-2}-m^2$, $c=C_a(0)$. The endpoint covariance is
$e^{-|z-z'|_1/a}/4$, a distinct covariance. Stabilizer $\epsilon>0$ and
$\Omega$ are as above. A width is a count of iid product/Mondrian whole-plane
process neurons, each with countably many variables. Lattice mesh,
spectral cutoff and covariance-flow time are not approximation parameters
in this chain.

% MAP: source-eb-width-data | S:734-760
For the fixed-stage width clause allow nested centered physical squares
$\Lambda_0\subset\Lambda_1$, including equality, enclosed by $\Omega$,
and a bounded box $D$ containing $\Omega$ and all supports.
Put $d_D=|D|$, $O=|\Omega|$, $E_D=L^4(D)$, $Y=\mathcal O_f$, and
\begin{align*}
 V_i(v)&=u\int_{\Lambda_i}(v^4-6cv^2+3c^2)
       +\epsilon\int_{\Omega\setminus\Lambda_i}v^4
       -\frac{\kappa_a^2}{2}\int_\Omega v^2,\\
 V_t&=(1-t)V_0+tV_1,\qquad0\le t\le1.
\end{align*}
Let $Q_N^b$ be the law of the normalized neuronal sum restricted to $D$,
and $Q_\gamma$ the endpoint Gaussian law. Write
$w_t=e^{-V_t}$, $Z_Q=Q(w_t)$,
$A_Q(t,G)=Q[G(Y)w_t]$, $\rho_{Q,t}=Y_*(w_tQ/Z_Q)$ and
$L_Q(t,G)=A_Q/Z_Q$. Their $t=i$ endpoints are the actual associated
neural/full-Gaussian-tilt observation laws. The inherited constants are
\begin{align*}
 \alpha&=\min(u,\epsilon),&U&=\max(u,\epsilon),\\
 B_N&=\kappa_a^2/2+6uc,&L_N&=B_N^2O/(2\alpha),\\
 F&=\left(\sum_i\|f_i\|_{4/3,D}^2\right)^{1/2}.
\end{align*}
Use $v_1=4Ux^3+2B_N\sqrt O\,x$, $v_2=12Ux^2+2B_N\sqrt O$,
$v_3=24Ux$, $v_4=24U$, and precisely the $e_0,\ldots,e_4$ Bell
polynomials displayed earlier. With $\beta_0=1$,
$\beta_k=(k/(2\alpha e))^{k/4}$ for $k>0$, put
\begin{align*}
 \mathcal T_N\left(\sum p_kx^k\right)&=e^{L_N}\sum p_k\beta_k\quad(p_k\ge0),\\
 M_Z&=\mathcal T_N(e_4),&
 M_A&=\mathcal T_N\left(\sum_{j=0}^4\binom4jF^je_{4-j}\right),\\
 z&=e^{-3UO/16-3uc^2|\Lambda_1|},&
 C_{\mathrm{law}}&=d_D(M_A+M_Z)/(96z),\\
 D_{\mathrm{law}}&=\tfrac43(3\sqrt6C_{\mathrm{law}})^{1/4}(2\sqrt r)^{3/4}\\
 &=2^{23/8}3^{-5/8}r^{3/8}C_{\mathrm{law}}^{1/4},&
 D'_{\mathrm{law}}&=\sqrt6C_{\mathrm{law}}+2\sqrt r.
\end{align*}
Euler's number $e$ is distinct from $e_j$.

% MAP: source-eb-laws | S:762-778
For one physical $\Lambda$, set $A=|\Lambda|$,
$B_{\mathrm{ext}}=|\Omega\setminus\Lambda|$,
$F_p=(\sum_i\|f_i\|_p^2)^{1/2}$ and
$L_c=(4\sqrt{15}+12)c^{3/2}$. The literal laws are
\begin{gather*}
 P=P^b_{N,\epsilon,\Omega,a;\Lambda},\quad
 Q_\epsilon=\nu^v_{\epsilon,\Omega,a;\Lambda}=\iota_*\operatorname{Res}_\epsilon,
 \quad Q_0=\nu^{(2)}_{\Omega,a;\Lambda}=\iota_*\operatorname{Res}_0,\\
 Q_a=\nu_{a,\Lambda},\quad Q_{\mathrm{phys}}=\nu_{0,\Lambda},
 \quad Q_{\mathrm{target}}=\nu=\nu_\infty^X.
\end{gather*}
The residual law has normalized density $e^{-W_{a,\Lambda}-\epsilon J}$
against $G'_{\Omega,a}$, where
$W=u\int_\Lambda(x^4-6c_ax^2+3c_a^2)$ and
$J=\int_{\Omega\setminus\Lambda}x^4$.
$Q_\epsilon$ is the positive-stabilizer full tilt and $Q_0$ its residual
endpoint. $Q_a$ is the normalized $e^{-U_{a,\Lambda}}$ tilt on the actual
common realization $\Phi_a$; $U_{0,\Lambda}$ is the fourth-chaos limit.
Classical finite-volume compact-test laws equal those of $Q_{\mathrm{phys}}$,
and $\nu$ is the existing target. Mass, factorial and normalization stay fixed.
Define
\begin{align*}
 \Delta_{\Omega,a}&=C_a-C'_{\Omega,a},&
 \delta_\Omega(H)&=\sup_{x\in H}\Delta_{\Omega,a}(x,x),
 \quad\delta_\Omega(\varnothing)=0,\\
 \xi_a(s)&=\kappa_a^2\int_{|x|>s}C_a(x)^2\,dx,\\
 \kappa_0&=e(1+\sqrt2),&\kappa_2&=(2\pi)^{-1}\sqrt{17\pi/16},\\
 M_6&=(2\pi)^{-5/3}(5\pi)^{5/6}m^{-1/3},&
 K_\Lambda&=u\sqrt{384A\kappa_2M_6^3},\\
 K_{\mathrm{sharp}}&=\max\{12uA,(2/e)K_\Lambda\sqrt{\kappa_0/m}\},\\
 T_\Lambda(H)&=\exp[-\exp((\pi/2)\sqrt{H/(12uA)}-1)],\\
 \Xi_p(\Lambda)&=e^{pK_{\mathrm{sharp}}}
    +\int_{K_{\mathrm{sharp}}}^\infty pe^{pH}T_\Lambda(H)\,dH.
\end{align*}

% MAP: source-eb-a | S:780-791
\paragraph{EB1.a: actual fixed-stage neural law.}
For either architecture and every $N\ge1$, uniformly in $t\in[0,1]$,
\begin{align*}
 \sup_{G\in\mathcal U_4(r)}|L_{Q_N^b}(t,G)-L_{Q_\gamma}(t,G)|
 &\le\min\{2,C_{\mathrm{law}}/N\},\\
 d_{\mathrm{BL}}(\rho_{Q_N^b,t},\rho_{Q_\gamma,t})
 &\le\sqrt6C_{\mathrm{law}}/(N\delta^3)+2\sqrt r\delta,
 \quad0<\delta\le1,\\
 \sup_t d_{\mathrm{BL}}(\rho_{Q_N^b,t},\rho_{Q_\gamma,t})
 &\le\min\{2,D_{\mathrm{law}}N^{-1/4}\},\qquad
 D_{\mathrm{law}}\le D'_{\mathrm{law}}.
\end{align*}
All derivative norms are full multilinear operator norms and endpoint
identification is included. The inherited finite clause covers all widths,
singular/repeated/sign-reversed tests, tests outside $\Omega$, constant $G$,
$u=\epsilon$ and empty shells. Zero shell response does not identify neural
and full-tilt laws. Empty/all-zero tuples have zero error directly; $r=0$
does not need smoothing. Uniformity is only for fixed positive parameters
and bounded geometry. The current theorem retains its narrower original
ordered bump family.

% MAP: source-eb-b | S:793-798
\paragraph{EB1.b: stabilizer.}
With variation $\sup_{\text{Borel }|H|\le1}$, whose maximum is two,
\[
 \|Q_\epsilon-Q_0\|_{\mathrm{var}}
 \le\min\{2,6\epsilon B_{\mathrm{ext}}c_a^2e^{9uAc_a^2}\}=S.
\]
This includes $\epsilon=0$ through the residual laws and gives zero when
$B_{\mathrm{ext}}=0$.

% MAP: source-eb-c | S:800-806
\paragraph{EB1.c: auxiliary region.}
In the same tuple metric,
\[
 d_{\mathrm{BL}}(\mathcal L_f(Q_0),\mathcal L_f(Q_a))
 \le\min\left\{2,e^{9uAc_a^2}
 [2uAL_c\sqrt{\delta_\Omega(\overline\Lambda)}
                         +F_1\sqrt{\delta_\Omega(K)}]\right\}.
\]
If $\Omega$ contains all points at Euclidean distance less than $s>0$
from compact $H$, then $\delta_\Omega(H)\le\xi_a(s)$, which tends to
zero at fixed $a$ as $s\to\infty$. Thus the error vanishes for fixed
$a,\Lambda,f$ as $\Omega$ grows. It suffices that
$\overline\Lambda\cup K\subset[-R_H,R_H]^2$ and
$0<s<R_{\mathrm{aux}}-R_H$. This inherited clause asserts no
uniform-in-$a$ or exponential auxiliary rate by itself.

% MAP: source-eb-d | S:808-813
\paragraph{EB1.d: physical regulator.}
At fixed $\Lambda$ and $0<a<1/m$,
\[
 d_{\mathrm{BL}}(\mathcal L_f(Q_a),\mathcal L_f(Q_{\mathrm{phys}}))
 \le\min\{2,\Xi_2(\Lambda)^{1/2}[4K_\Lambda\sqrt a+(a/2)F_2]\}.
\]
It tends to zero at fixed $\Lambda,f$; neither $\Xi_2$ nor $K_\Lambda$
is asserted uniform in $\Lambda$.

% MAP: source-eb-f | S:815-826
\paragraph{EB1.f: compatible budget and literal limits.}
For the fixed tuple, let accepted CL1-FV1 furnish $n_0$ with
$K\subset\Lambda_{n_0}$. For $n\ge n_0$, $n\ell\ge R_{\mathrm{obs}}$,
use $\kappa=m/256$, $q=e^{-\kappa\ell}$ and the same $T_n,E_n$ as AP6.
Choose the EB1.a constants with $\Lambda_0=\Lambda_1=\Lambda_n$.
For every admissible positive $a,\epsilon$, bounded $\Omega$ enclosing
$\overline\Lambda_n$, a box $D$ containing $\Omega$ and supports, and
$N\ge1$, the five-term budget is the AP6-displayed expression with its
historical $\nu$ target. It passes exactly through
$Q_\epsilon,Q_0,Q_a,Q_{\mathrm{phys}},\nu$ and implies the literal
iterated limit: $N$ first, then $\epsilon\to0$, $R_{\mathrm{aux}}\to\infty$,
$a\to0$, $n\to\infty$. A tolerance schedule reverses these dependencies:
$n,a,R_{\mathrm{aux}},\epsilon,N$.
Historical $C_f,g_*$ retain their possibly noneffective source meaning;
there is no arbitrary simultaneous schedule, practical numerical width,
new dimension, stronger coupling or priority assertion in that clause.

% MAP: source-ec-intro | S:830-847
\subsubsection{Inherited EC1.a--c coefficient bounds}
The corrected EC1 statement digest recorded in the source is
\par\noindent{\small\texttt{1406f33ad2370b8f5c59751ec73d06ada0c3a9c159fd75e3ec6b4abbb748b4fe}}\par
These are rules for the physical variance, covariance tail and fixed-square
Nelson constant; the accepted positive-accuracy selector is not an oracle
for the present tuple. Use the exact massive covariance/Fourier convention
at historical \texttt{CONSTRUCTION.tex:303--330},
\[
 C_a(z)=(2\pi)^{-2}\int_{\R^2}
 \frac{e^{ip\cdot z}}{m^2+\omega^2+\xi^2+a^2\omega^2\xi^2}\,dp,
\]
\[m>0,\quad0<a<1/m,\quad\kappa_a^2=a^{-2}-m^2>0,\quad c_a=C_a(0).
\]
Tail norms are Euclidean. The physical covariance differs from the
variance-$1/4$ endpoint covariance. Retain the exact EB1
$K_\Lambda,K_{\mathrm{sharp}},T_\Lambda,\Xi_p$ for $u>0$ and positive
finite area $A$, with $\log_+x=\max(0,\log x)$.
For width clauses use the same possibly equal nested squares, $O>0$,
$d_D>0$, $\alpha=\min(u,\epsilon)>0$, $U=\max(u,\epsilon)$,
$B_N=\kappa_a^2/2+6uc_a$ and tuple norm $F$. No convention changes.

% MAP: source-ec-a | S:849-869
\paragraph{EC1.a: point variance and certified enclosures.}
For $h=ma\in(0,1)$,
\begin{gather*}
 \max\left\{\frac14,\frac{\log(1/h)-(1-h^2)/2}{2\pi}\right\}
 \le c_a\le\overline c(a)=\frac{\operatorname{asinh}(1/h)+1}{2\pi},\\
 \operatorname{asinh}(x)=\log(x+\sqrt{1+x^2}).
\end{gather*}
Hence $c_a=(2\pi)^{-1}\log(1/(ma))+O(1)$ as $a\downarrow0$ at fixed
$m$. Upper-bound uses, including positive-coefficient EB1 widths, may
replace $c_a$ by $\overline c(a)$. For $Q>0$ and a partition
$0=q_0<\cdots<q_M=Q$, define
\begin{align*}
 H(q)&=[(m^2+q^2)(1+a^2q^2)]^{-1/2},\\
 c_{\mathrm{lo}}&=(2\pi)^{-1}\sum_j(q_{j+1}-q_j)H(q_{j+1}),\\
 c_{\mathrm{hi}}&=(2\pi)^{-1}
 \left[\sum_j(q_{j+1}-q_j)H(q_j)+1/(aQ)\right].
\end{align*}
Then $c_{\mathrm{lo}}\le c_a\le c_{\mathrm{hi}}$ and
\[
 c_{\mathrm{hi}}-c_{\mathrm{lo}}
 \le\frac{\max_j(q_{j+1}-q_j)/m+1/(aQ)}{2\pi}.
\]
Thus positive computable $m,a$ satisfying the strict regulator condition
give computability of $c_a$ by certified elementary operations.

% MAP: source-ec-b | S:871-888
\paragraph{EC1.b: exponential squared-covariance tail.}
For $s\ge0$,
\begin{gather*}
 \int_{|z|>s}C_a(z)^2\,dz
 \le\min\{1/(4\pi m^2),e^{-\sqrt2ms}/(2\pi m^2)\},\\
 \xi_a(s)\le\rho_a(s)=\frac{\kappa_a^2}{4\pi m^2}
                         \min\{1,2e^{-\sqrt2ms}\}.
\end{gather*}
At fixed $a$ this vanishes explicitly. The rate depends on $m$ and the
prefactor retains its regulator dependence. The same strict source geometry
is mandatory: if the Euclidean $s$-neighborhood of compact $H$ lies in
$\Omega$, $\delta_\Omega(H)\le\rho_a(s)$. Thus if
$\overline\Lambda\cup K\subset[-R_H,R_H]^2$ and
$0<s<R_{\mathrm{aux}}-R_H$, the auxiliary error is at most
\[
 \min\{2,H_{\mathrm{aux}}\sqrt{\rho_a(s)}\},\qquad
 H_{\mathrm{aux}}=e^{9uA\overline c(a)^2}
 [2uA(4\sqrt{15}+12)\overline c(a)^{3/2}+F_1].
\]
No regulator-uniform defect modulus or interchange of limits is asserted.

% MAP: source-ec-c1 | S:890-908
\paragraph{EC1.c: elementary bounds and enclosures for $\Xi_p$.}
For every fixed $p\ge1$ set
\[
 b=\frac{\pi^2}{48uA},\qquad
 H_p=\max\left\{K_{\mathrm{sharp}},\frac4b
                      \left[\log_+\frac{32p}{eb}\right]^2\right\}.
\]
Then
\[
 e^{pK_{\mathrm{sharp}}}\le\Xi_p(\Lambda)
 \le e^{pH_p}+e^{-pH_p}\le2e^{pH_p}.
\]
These finite elementary values give explicit certified
$\Xi_2^{1/2}$ and fixed-volume regulator bounds. The exact equivalent
formula $K_\Lambda^2=[75\sqrt{85}/(2\pi^3m)]u^2A$ avoids fractional
cube-root powers in numerical evaluation.

% MAP: source-ec-c2 | S:910-924
For $T\ge H_p$ and a partition $K_{\mathrm{sharp}}=x_0<\cdots<x_M=T$,
the lower and upper enclosures for $\Xi_p$ are respectively
\begin{align*}
 &e^{pK_{\mathrm{sharp}}}
   +\sum_j(x_{j+1}-x_j)pe^{px_j}T_\Lambda(x_{j+1}),\\
 &e^{pK_{\mathrm{sharp}}}
   +\sum_j(x_{j+1}-x_j)pe^{px_{j+1}}T_\Lambda(x_j)+e^{-pT}.
\end{align*}
The degenerate interval $T=K_{\mathrm{sharp}}$ has empty sum.
Increasing $T$ and refining the partition makes the enclosure arbitrarily
tight, proving computability for computable $m,u,A,p$ without evaluating
any CL1 source coefficient. At fixed $m,u,p$, as $A$ grows,
\begin{align*}
 \log\Xi_p&\ge12puA,\\
 \log\Xi_p&\le p\max\left\{K_{\mathrm{sharp}},\frac{192uA}{\pi^2}
           \left[\log_+\frac{1536puA}{e\pi^2}\right]^2\right\}+\log2.
\end{align*}
This upper bound grows as $O_{m,u,p}(A\log^2A)$ and concerns the source
Nelson bound, without a volume-uniform moment assertion.

% MAP: source-opt-intro | S:928-942
\subsubsection{Inherited OPT coefficients and physical Gaussian facts}
The accepted OPT statement digest recorded in the frozen source is
\par\noindent{\small\texttt{1c178d728e6ffbf744b19eff30433fff75961b67365dcbfec0ef17f196e61e5b}}\par
Use the same chosen explicit H1 coefficient $C_H<2^{16}$ in every
$C_f$ formula. OPT.a--b supply the usable H2 value $C_{H2}=2^{384}$
in the literal $J$ expression and its scalar $R=1$ coefficient
$C_f\le2^{360}$. Their finite-source assumptions and attribution remain.
The source appendix describes these as inherited inputs for SELECT;
that historical label does not change the present UDA interface. The
finite-lattice variables below are local to that inherited theorem,
and its original scalar tuple is unused. Reproving the earlier coefficient
theorem is outside the present translation.

% MAP: source-opt-a | S:944-951
\paragraph{OPT.a: explicit finite H2 coefficient.}
At the fixed $m,u,\ell,K_{\mathrm{cov}}=1$, on the entire remaining
finite-lattice domain, with actual posterior, actual outer law and literal
A1--A6 from the finite-interface appendix, the physical Wick covariance obeys
\[
 \left\|\varepsilon^2\sum_{z\in S_n^\varepsilon}e^{d_L(x,z)/128}
 \left|\operatorname{Cov}({:}\psi_z^4{:}_c,{:}\psi_x^3{:}_c\mid Y_s)\right|
 \right\|_{L^2(\nu_s)}\le2^{384}g(s)^{7/2}.
\]
This includes every permitted $\varepsilon,n,t,L$, all $0<s\le1$, both
$t$ endpoints, zero-profile sites and coincident $x,z$.
The squared envelope is $2^{768}g(s)^7$, with finite required square-root
integral. No unquantified spatial score constant or translated-moment
premise is added.

% MAP: source-opt-b | S:953-957
\paragraph{OPT.b: literal bump coefficient.}
With the earlier explicit H1 formula, OPT.a's H2 coefficient, $R=1$,
$n_0=1$, $\epsilon_*=1/4$, $\kappa=1/256$ and $c_{\mathrm{CT}}=81/5$,
the literal CL1 expression obeys $C_f\le2^{360}$.
The formula includes its source integral, endpoint, every prefactor and
$[m^2(1+m^2)]^{-1}$. AP6 gives its applicability to the same original
pressure-identified FS target under P1--P6. Earlier acceptances
$C_{H2}\le2^{1800}$ and $C_f\le2^{1810}$ do not by themselves accept
the changed coefficients.

% MAP: source-gaussian-physical | S:961-973
The inherited Gaussian/common-space facts are premises already used in
the EB1 proof, rather than conclusions of SELECT's new source estimate.
They hold for every fixed finite real compact-smooth tuple, including
singular test covariance. At $m=1$, $0<a<1$, the physical $\Phi_a$ has
the multiplier stated earlier, with measure $dp/(2\pi)^2$.
On the accepted common realization,
$\E|(\Phi_a-\Phi_0)(f)|^2\le(a^2/4)\|f\|_2^2$ for every such real $f$.
The continuous Gaussian endpoint has covariance $e^{-|z-w|_1/a}/4$,
distinct from the physical point variance $c_a$.

% MAP: source-gaussian-defect | S:975-985
For each $\Omega$, $C'_{\Omega,a}$ is an actual continuous covariance
kernel with $0\le C'_{\Omega,a}\le C_a$ in positive-semidefinite kernel
order. Thus $\Delta=C_a-C'_{\Omega,a}$ is positive semidefinite.
There exist independent centered continuous Gaussian fields $\gamma',\psi$
with covariances $C',\Delta$. Then $X_c=\gamma'+\psi$ has physical
regulated reference law, and $Y_c=\gamma'$ has the mass-tilt reference law.
The weights $w_X=e^{-W_{a,\Lambda}(X_c)}$ and
$w_Y=e^{-W_{a,\Lambda}(Y_c)}$ yield the actual normalized $Q_a,Q_0$.
Their common upper bound is $e^{6uAc_a^2}$; the $Y$ normalizer is at least
$e^{-3uAc_a^2}$. If the strict Euclidean $s$-neighborhood of compact $H$
lies in $\Omega$,
$\delta_\Omega(H)\le\kappa_a^2\int_{|z|>s}C_a(z)^2\,dz$.

% MAP: source-gaussian-weights | S:987-994
On the physical common Gaussian space $U_{a,\Lambda}$ is the
coefficient-one own-variance Wick interaction, and $U_{0,\Lambda}$ its
fourth-chaos limit. Then $\E U_{a,\Lambda}=0$,
$Z_a=\E e^{-U_{a,\Lambda}}\ge1$. For $w_a=e^{-U_{a,\Lambda}}$,
$w_0=e^{-U_{0,\Lambda}}$,
\[
 \E|w_a-w_0|\le2K_\Lambda\Xi_2(\Lambda)^{1/2}\sqrt a,
 \qquad\|w_0\|_2\le\Xi_2(\Lambda)^{1/2}.
\]
All stated physical moments are finite at each fixed physical square.
$Q_{\mathrm{phys}}$ has the original classical $\nu_B$ compact-test laws,
as AP5 also supplies.

% MAP: source-gaussian-l0 | S:996-1003
Finally, accepted L0 is an ambient Euclidean premise. For $G$ with separate
unit sup and Lipschitz bounds on $\R^r$, standard Gaussian $\xi\in\R^r$
and $G_\delta(y)=\E G(y-\delta\xi)$, $G_\delta$ is smooth,
$\|G-G_\delta\|_\infty\le\sqrt r\delta$,
$\|G_\delta\|_\infty\le1$, $\|DG_\delta\|_{\mathrm{op}}\le1$ and
$\|D^jG_\delta\|_{\mathrm{op}}\le\sqrt{(j-1)!}\delta^{1-j}$ for $j\ge2$.
For $0<\delta\le1$, all norms through four are at most $\sqrt6\delta^{-3}$.
It holds in the original ambient dimension irrespective of covariance rank.


% The accepted UDC source contract and its implementation proof.
% MAP: udc-status | US:1-17 UP:1-3 UC:1-13 UD:1-30
\subsection{The separately versioned finite-resource implementation}
\label{udc:contract}
UDA's every-$q$ statement concerns finite mathematical integers. The
separately accepted UDC-dyadic-accuracy-v1 implements a finite partial
realization: success gives an exact symbolic certificate, and a valid
mathematical request may resource-refuse. Its frozen source files retain
their author-candidate status flags; the external exact-current UDC decision
supplies their internal acceptance. Neither a success flag nor a newly
self-hashed core supplies acceptance. The unchanged older ACC and SVC
interfaces retain their own names, domains and bytes.

% MAP: udc-identities | US:3-9 UC:1-13 UP:23-27 G:141-161 V:139-180 H:15-40
The version identities are:
\begin{description}[leftmargin=1em,style=nextline,itemsep=0.1em]
\item[Request] \nolinkurl{UDC-dyadic-accuracy-request/1}
\item[Certificate] \nolinkurl{UDC-dyadic-certificate/1}
\item[Backend] \nolinkurl{accepted-UDA-STAGE-RELATIVE-v2-bounded-v1}
\item[Source bundle] \nolinkurl{UDC-UDA-STAGE-RELATIVE-v2-sources-20261028-v1}
\end{description}
The following entry-point and source-root descriptions refer to the original
private UDC packet, which is not included in this public document export.
Its original private source map supplies the paths and digests described below.
The command-line entry points are \nolinkurl{dyadic_accuracy_search.py} and
the separate \nolinkurl{verify_dyadic_accuracy.py}. Both use the adjacent \texttt{sources/}
tree by default; \texttt{--source-root} relocates an identical complete
tree. Every declared source body is local and hash-bound, with no live
checkout fallback. The companion source map binds these names to their
actual supplied relative paths and full digests.
In this packet the core root is \nolinkurl{inputs/accepted/UDC/}, with its
direct tree at \nolinkurl{inputs/accepted/UDC/sources/}. The eight bound core
bodies are the producer, verifier, \nolinkurl{dyadic_common.py},
\nolinkurl{SOURCE-PINS.json}, \nolinkurl{STATEMENT-ONLY.md},
\nolinkurl{PROOF.md}, \nolinkurl{INPUT-CONTRACT.md} and
\nolinkurl{INPUT-SCHEMA.json}. The reviewed frozen core supplies the external
version trust root for those exact hashes.

% MAP: udc-request | US:3-7 UC:12-16 UP:5-5 H:86-208 I:1-318
Every fixed request member and member set is exact. The request retains
the original model, metric, architectures and ordered translated-bump domain.
A dyadic tolerance has \texttt and a canonical decimal-string
\texttt{q}, such as \texttt{"13"}; the separate compatibility form is
\texttt{kind=one-fifth}. Decimal $q$ has no sign, leading zero, decimal point,
exponent notation or whitespace. Values below three are unsupported.
Canonical positive text exceeding the chosen representation ceiling
resource-refuses before conversion to an integer or a $q$-dependent
allocation. There is no mathematical upper $q$. Coordinate numerators are
signed canonical decimal strings, with exponent $0,\ldots,16$, odd numerator
for a positive denominator exponent, and exponent zero for integers and zero.
The exact coordinate bound $|P|\le8\,2^k$ uses only fixed small shifts.

% MAP: udc-capacities | UC:14-16 UP:5-5 H:22-40 H:157-208 H:210-268 I:1-318
All ten limits are required. The first four are ordinary bounded integer
fields; each stage limit is a typed integer node with
\texttt and canonical decimal text. A node's text must also fit
\texttt{max\_integer\_digits}. The defaults and inclusive hard ranges are:
\begingroup
\small
\setlength{\parskip}{0pt}
\begin{longtable}{@{}p{0.46\linewidth}p{0.48\linewidth}@{}}
\toprule
Limit & Default; hard range \\
\midrule
\endhead
\texttt{max\_integer\_digits} & $1024$; $1,\ldots,4096$ decimal magnitude digits \\
\texttt{max\_watched\_bits} & $8192$; $16,\ldots,8192$ bits of selected magnitudes \\
\texttt{max\_milliseconds} & $30000$; $1,\ldots,30000$ cooperative milliseconds \\
\texttt{max\_steps} & $20000$; $1,\ldots,50000$ named logical ticks \\
\texttt{max\_n} & $10^{100}$; $1,\ldots,10^{2048}$ \\
\texttt{max\_k} & $10^{100}$; $1,\ldots,10^{2048}$ \\
\texttt{max\_regulator\_exponent} & $10^{400}$; $4,\ldots,10^{2048}$ \\
\texttt{max\_auxiliary\_radius\_exponent} & $8191$; $1,\ldots,8191$ \\
\texttt{max\_stabilizer\_exponent} & $10^{1000}$; $34,\ldots,10^{2048}$ \\
\texttt{max\_width\_exponent} & $10^{1000}$; $1,\ldots,10^{2048}$ \\
\bottomrule
\end{longtable}
\endgroup
Equality at stage, representation, bit, transport, work and timing caps
passes. A conservative small-power preflight can refuse when its
preconstruction bound $\operatorname{bitlen}(b)j$ exceeds the watched-bit
cap even though a particular eventual value could fit. These are selected
software capacities, rather than necessary mathematical resources.
The width cap bounds $e_{\mathrm{new}}=\log_2N_{\mathrm{new}}$;
$w_{\mathrm{old}}$ has the different unit $\log_2\log_2N_{\mathrm{old}}$.

% MAP: udc-parser | US:17-17 UC:16-16 UC:34-34 UP:5-5 UP:27-27 H:70-135 H:137-208
\paragraph{Strict parsing and refusal precedence.}
API requests and certificate results are bytes. Request transport is capped
at $32768$ bytes and certificate-result transport at $131072$ bytes; CLI
reads stop after cap plus one. Before numerical or source work the parser
checks strict UTF-8, no BOM, nesting at most sixteen, at most $5000$
value/container nodes, strings at most $4096$ characters, member names at
most $256$ characters, and integer tokens at most twenty characters including
a sign. Duplicate members, floats, nonfinite tokens, malformed syntax,
trailing data and unsupported types refuse; booleans are distinguished from
integers. Byte capacity is reached before decoding, nesting before recursive
JSON parsing, integer-token capacity during parsing, and the remaining
structure capacities before semantic validation. A reached finite capacity
returns \texttt{RESOURCE\_EXHAUSTED}; it can prevent a later syntax condition
from being reached. Malformed syntax and unsupported mathematical values
return \texttt{UNSUPPORTED\_INPUT} when reached. Certificate parsing preserves
capacity refusal and translates malformed syntax to
\texttt{REFUSED\_CERTIFICATE}. This retains the corrected parser category
without declaring a large valid $q$ outside the theorem.

% MAP: udc-representation-proof | US:7-9 UC:16-16 UP:5-5 H:157-208 H:431-481 G:141-161
Every derived mathematical integer is a typed bounded decimal node;
rationals have typed numerator and positive denominator. Powers are inert
\texttt{pow2} nodes with typed exponents. The tolerance, $a$, $\epsilon$,
$\eta$, $\zeta$ and both widths remain symbolic. No $2^q$, $2^p$, $2^t$,
$N_{\mathrm{new}}$ or $2^{w_{\mathrm{old}}}$ is materialized. Every success
node is watched and checked against its decimal digit ceiling. The hard
$8192$-bit watched bound implies fewer than $2468$ decimal magnitude digits,
below the $4096$-character string cap and the interpreter conversion guard.
The producer reparses its own encoded result before reporting success.
Thus successful output contains neither an oversized bare integer token nor
an unreadable decimal node. A representation failure is an explicit resource
refusal, even when the corresponding mathematical integer exists.

% MAP: udc-factored-proof | UP:7-7 UC:18-18 G:11-23 V:12-20 H:210-268
\paragraph{Exact factored arithmetic.}
Let $a,b>0$ be integers and $z$ a finite signed integer. To decide
$a2^z\le b$ when $z\ge0$, compare
$\operatorname{bitlen}(a)+z$ with $\operatorname{bitlen}(b)$.
Unequal lengths decide the comparison; equal lengths require the guarded
exact integer shift. When $z<0$, compare $\operatorname{bitlen}(a)$ with
$\operatorname{bitlen}(b)-z$, with the same equal-length rule.
The verifier separately computes the displacement
$z+\operatorname{bitlen}(a)-\operatorname{bitlen}(b)$ and uses its sign
before a guarded equal-length comparison. Both rules follow directly from
$2^{h-1}\le M<2^h$ for a positive integer of binary length $h$.
Equality returns true. The guard checks the nonnegative shift exponent plus
one before constructing that power. These exact predicates work even when
the displayed exponent is far beyond a materializable power; no floating
logarithm or astronomical shift is needed.

% MAP: udc-log-proof | UP:9-9 G:25-32 V:22-29 S:227-248
For a positive rational $a/b$, the producer starts at
$h=\max(0,\operatorname{bitlen}(a)-\operatorname{bitlen}(b))$ and increases
once if the factored test $a/b\le2^h$ fails. Binary-length bounds put the
least answer in precisely these two positions; it checks the answer and
predecessor. The verifier independently uses zero for $a\le b$, and otherwise
\begin{equation}\label{udc:clog}
 \operatorname{clog}_2(a/b)
 =\operatorname{bitlen}\!\left(\left\lfloor\frac{a-1}{b}\right\rfloor\right).
\end{equation}
Indeed $\lfloor(a-1)/b\rfloor=\lceil a/b\rceil-1$, whose binary length is
the least $h$ with $a/b\le2^h$. Its predicate and predecessor are also
checked. This proves every non-strict dyadic tie rule in both paths.
In particular $h_K=\lceil\operatorname{clog}_2(14u^2A)/2\rceil$ is the
least integer with $14u^2A\le2^{2h_K}$, including exact powers of four and
arguments at most one.

% MAP: udc-spatial-proof | UP:11-11 G:41-72 V:37-67 S:229-240 P:56-82
\paragraph{Least spatial and Nelson stages.}
Both paths use the exact geometry and $\overline C$ above. The producer writes
$\overline E(i)=C(i)2^{-5i}$ with
\[
 C(i)=\frac{\overline C\,\overline{\exp R}}{2500}
 \left[\frac{32(2i+1)}{31}+\frac{64}{961}\right].
\]
The verifier independently reduces the bracket to $(1984i+1056)/961$.
The spatial predicate is exactly $C(i)2^{d-5i}\le1$. The preceding
mathematical proof shows strict monotonicity and the positive upper witness
$\lceil(d+356)/4\rceil$. The producer bisects from zero to the lesser of
this witness and \texttt{max\_n}; failure of the high predicate is a stage
refusal. The verifier separately doubles from one, clamps at the same finite
cap, and refines its false/true bracket. Each refinement strictly reduces
the positive endpoint gap. The final predicate and predecessor prove the
same least $n$. The selected ratio $\overline E(n)/\eta$ is materialized
only with a separately guarded finite power and watched reduced rational.
This is a capacity check on that ratio, rather than construction of $\eta$.

% MAP: udc-nelson-proof | UP:13-13 G:73-91 V:67-88 S:237-248 P:84-132
At this $n$, $A=4000000n^2$, $U_A=A/10^{10}$ and $b_0=3/(16U_A)$ are
exact rational values. The producer forms $L_{10}(k)$ using direct guarded
nonnegative polynomial powers and factorials. The verifier uses the distinct
recurrence $T_0=1$, $T_j=T_{j-1}k/j$, summing through ten. Both give the
same degree-ten polynomial and strict predicate $L_{10}(k)>32/b_0$.
The finite witness $\lceil32/b_0\rceil$, producer bisection, independent
verifier doubling/refinement, finite \texttt{max\_k} and checked predecessor
prove the same least $k$. Monotonicity follows from the nonnegative
coefficients and positive linear term. Also
$1+\lceil\max(12U_A,4\overline K,4k^2/b_0)\rceil$ equals one plus the
maximum of the individual ceilings. Thus the two implementations use exactly
$H_N$, without substituting the older compact bounds. Both derive
$p=4H_N+2h_K+2d+8$. Its first physical exponent is exactly $-d-1$, its
second is at most $-d-1$, and the preceding even integer fails the first
predicate. The actual $K_\Lambda$ and volume-dependent $\Xi_2$ remain in
the conditional error estimate.

% MAP: udc-outer-proof | UP:15-15 G:92-104 V:88-102 S:241-248 P:145-194
\paragraph{Auxiliary radius and stabilizer.}
At that same $p$, both compute $j=\lceil18U_Ap^2\rceil$ and the complete
$B_{\mathrm{aux}}=64U_Ap^2+4c_r$. By the proved logarithmic rule,
$h_s=1+\operatorname{clog}_2(j+h_{\mathrm{aux}}+p+d)$ is the first integer
at least one satisfying $j+h_{\mathrm{aux}}+p-2^{h_s-1}\le-d$.
Predicate and predecessor are checked before accepting the selected stage.
The guarded $s=2^{h_s}$ gives $R_{\mathrm{aux}}=1000n+s+1$ and
$O=4R_{\mathrm{aux}}^2$. Since $R\le9<1000n$, the strict source margin is
one and $O>A$. Finally
$t=\max(34,j+\operatorname{clog}_2(6Op^2)+d)$ gives the stabilizer exponent
at most $-d$ and $0<\epsilon\le2^{-34}<u$.
At every stage, equality with a cap passes and strict exceed refuses.
The accepted mathematical argument then supplies the actual three
$R_{\mathrm{phys}},A_{\mathrm{aux}},S$ bounds, with their full physical and
common-space premises.

% MAP: udc-relative-proof | UP:17-17 G:105-119 V:102-118 S:249-266 S:398-414 P:438-520
\paragraph{Relative threshold, parity and complete old comparison.}
Both implementations independently reconstruct $o_+=\max(o,78)$,
$T_{1/2}=\lceil t/2\rceil$, $x,y,k_1,n_0,k_4$ and $e_{\mathrm{new}}$.
They verify $4x\ge2t+4p+o_++3$ and
$4y\ge2t+8p+3o_++8$. The proof above keeps both $t$ parities, the actual
joint tilted moments through sixteen, every centered Bell/product term and
the vector coefficient $38432$. Consequently $K_1<2^{k_1-2}$ and
$16K_1^2<2^{2k_1}$. Since the right side is an integer,
$N_0\le2^{2k_1}<2^{n_0}$. The exact positive difference
$e_{\mathrm{new}}-n_0=6x+2y+4d+38$ proves threshold safety.
$K_4<2^{k_4}$ and $\zeta=2^{-d-3}$ give normalized-law error below
$3\eta/8$ and smoothing at most $\eta/2$. The actual first arrow is below
$7\eta/8$. Software reconstruction uses these proved inequalities without
replacing the interacting tilt by independent marginal laws.

% MAP: udc-old-proof | UP:19-19 G:110-119 V:109-118 S:258-266 S:426-442 P:559-685
At exactly those stages each path reconstructs the complete old coefficient
majorants $h_{\mathrm{root}},v_{\mathrm{old}},q_{\mathrm{old}},l_{\mathrm{old}}$
and $H_{\mathrm{small}}$. It checks both positive remainder guards,
$o\le l_{\mathrm{old}}-1$ and
$\operatorname{clog}_2(3U_Ap^2)\le l_{\mathrm{old}}-1$, and checks
$H_{\mathrm{small}}\le7l_{\mathrm{old}}+47$ with
$\operatorname{bitlen}(H_{\mathrm{small}})\le l_{\mathrm{old}}+2$.
Thus the first old candidate fails its explicit strict inequality and
$w_{\mathrm{old}}=l_{\mathrm{old}}+3$ passes both original guards.
The old sufficiency proof above retains coercivity and both Jensen terms.
The new bound $e_{\mathrm{new}}\le12l_{\mathrm{old}}+553$ and
$\operatorname{bitlen}(e_{\mathrm{new}})\le w_{\mathrm{old}}$ prove
$e_{\mathrm{new}}<2^{w_{\mathrm{old}}}$, hence
$2^{e_{\mathrm{new}}}<2^{2^{w_{\mathrm{old}}}}$, without constructing either
count. The obsolete compact $e<2^{40}$ and $6up<1$ guards do not appear.

% MAP: udc-target-proof | UP:21-21 US:9-13 G:127-139 V:124-137 S:416-461 P:522-557 P:687-749
\paragraph{The actual target budget and legacy equality.}
The five certificate arrows are the same
$P_N\to Q_\epsilon\to Q_0\to Q_a\to Q_{\mathrm{phys}}\to\nu^X$.
Their upper factors in units of symbolic $\eta$ are
$(1,1,1,1,\overline E(n)/\eta)$, whose sum is at most five.
The allocated total divided by tolerance is $5/8$ for $d=q+3$, and
$25/32$ for the separate $1/5$ request with $d=5$.
The stronger first-arrow bound is separately retained. UDA.f proves exact
equality of every listed legacy outer variable, both width formulas and all
relative exponents on the eleven old settings. Both software paths check the
bounded consequences and $e=16\lceil t/2\rceil+20p+4d+544$ when that
branch applies. This is mathematical stage equality; the new request and
certificate identities remain distinct. Finite replay cases check executed
paths, while the complete proof supplies the universal quantifier.

% MAP: udc-source-proof | UP:23-23 US:15-15 UC:20-22 H:270-420 V:139-180
\paragraph{Source and core integrity; independent arithmetic.}
The embedded manifest digest binds all $541$ direct source rows, including
each exact relative path, size and hash. Per-source reads are capped at
$1$ MiB and aggregate source reads at $16$ MiB. Resolved paths must remain
inside the supplied root; absolute paths, parent escapes and symlink escapes
refuse. Required original SELECT/OPT/EB1/EC1/ADMIT/SVA gates and material-role
aliases, strict AP6 and unknown-$g_*$ flags, and the exact-current UDA
statement, proof, freeze, decision and open gate are checked among explicit
digest aliases. An unavailable required body is a source gap; changed
size/hash/manifest/identity is a binding refusal; an over-cap read can
resource-refuse first. Every declared body is bundled. These are byte
identity checks, with the externally reviewed frozen core as the version
trust root; a self-rewritten core or manifest is a different, unaccepted
version. Source presence supplies no new analytic truth or authorship proof.

% MAP: udc-verifier-proof | UP:23-23 US:11-15 UC:22-22 G:141-161 V:1-180 H:399-420
The verifier imports only the common syntax, binding, constants and resource
utilities. It imports no producer stage/search/width routine and no older
arithmetic module. Its independent bracket/refinement, Nelson recurrence,
rational-power formula and integer algebra reconstruct the entire expected
typed object. Exact canonical encoded comparison checks every nested witness,
symbolic unit, source/core binding, scope flag, extra or missing member and
Boolean/integer distinction. It freshly binds current code and source bytes.
Historical logical steps, selected peak bits and elapsed time are checked as
typed provenance, with \texttt{historical\_execution\_authenticated=false}.
A fabricated but well-typed historical time may verify; the verifier's own
cooperative measurement is separate. Source/core hashes authenticate neither
historical elapsed time nor external scientific acceptance.

% MAP: udc-resource-proof | UP:25-25 UC:24-30 H:210-268 H:431-481 G:141-161 V:139-180
\paragraph{Selected watches, work and timing scope.}
Every search loop has a finite bracket/refinement variant or fixed degree,
with named logical-step polling. Ticks count selected source/core/manifest
reads, predicates, comparisons, logarithmic minimality checks, search
iterations and polynomial/recurrence terms. They are software actions,
rather than CPU instructions or a complexity bound. Output conversion,
watches and time polls do not add ticks. Strict step exceed refuses.
Selected integer and reduced-rational numerator/denominator watches cover
the parsed stage caps and $q$, all selected stage values, factored comparison
arguments, direct polynomial powers or recurrence terms, radii, area,
relative/old exponents, ratio and typed output. A power-of-two allocation is
guarded by exponent plus one; a degree at most twelve small power uses its
conservative preconstruction bit bound. Fixed small constants and all Python
arithmetic intermediates are not universally instrumented. These watches
supply no all-intermediate or process-memory theorem.

% MAP: udc-timing-proof | UP:25-25 UC:28-30 H:210-268 G:141-161 V:139-180
Cooperative time begins after request parsing and validation in both APIs;
the verifier parses the certificate result after creating its clock.
Source reads, ticks, watches and final polls check elapsed time. Clock reads
must be finite non-Boolean integers/floats, integer magnitude at most
$10^{12}$, callable and monotonically nondecreasing. A bad or backward clock
resource-refuses. Equality at the millisecond limit passes; strict exceed
refuses. An operation is not interrupted between polls. Serialization,
output, interpreter startup and whole-process duration are outside that
interval. Arbitrary interruption, interpreter/hardware/output-device failure
and process-memory guarantees are outside the defined certificate contract.
CLI output uses exclusive creation, preserving prior result bytes on retries.

% MAP: udc-results | US:15-17 UP:23-27 UC:30-34 H:42-59 H:422-459 G:141-161 V:139-194
The producer's \texttt{SUPPORTED\_DOMAIN\_SUCCESS} supplies a conditional
numerical certificate. The verifier's \texttt{CERTIFICATE\_VERIFIED}
reports fresh binding and separate arithmetic. The defined refusals are:
\begin{itemize}[nosep]
\item \texttt{UNSUPPORTED\_INPUT}: malformed or unsupported request.
\item \texttt{NUMERICAL\_SOURCE\_GAP}: required frozen source unavailable.
\item \texttt{SOURCE\_BINDING\_REFUSAL}: changed binding or identity.
\item \texttt{RESOURCE\_EXHAUSTED}: a reached declared finite capacity.
\item \texttt{REFUSED\_CERTIFICATE}: malformed, incompatible or false result.
\end{itemize}
Every defined refusal has
\texttt{numeric\_certificate\_available=false} and omits a numerical
certificate. Success and verification keep self-acceptance false. The
pre-correction parser core, original harness failures and the minor corrected
forty-to-thirty-nine case count remain historical evidence; the current
parser categories and exact current source decision govern this version.

% MAP: udc-witness | E:1-1115 US:3-9 UP:21-21 UD:1-30
\paragraph{A recorded request beyond the old eleven settings.}
The accepted portable example uses the four distinct ordered corners
$(-8,-8),(-8,8),(8,-8),(8,8)$, both architectures, and $q=13$.
Its certificate records
\begin{align*}
 d&=16,&n&=76,&H_N&=2417,&p&=9708,\\
 h_s&=33,&R_{\mathrm{aux}}&=8590010593,&t&=3919396758,\\
 e_{\mathrm{new}}&=31355368832,&w_{\mathrm{old}}&=3919435662.
\end{align*}
The exact spatial ratio is $\overline E(n)/\eta=691456/11411875$.
Thus the symbolic counts are $2^{31355368832}$ and
$2^{2^{3919435662}}$ in their distinct units. This is recorded exact
conditional arithmetic; it is neither an actual field simulation nor a
measurement of the cost of constructing those process neurons. Any replay
performed for this manuscript has its own hashed receipt. The universal
theorem comes from the proof above, independently of this finite witness.

% MAP: udc-final-scope | US:13-17 UC:32-34 UP:25-27 S:10-21 S:463-484 P:767-776 UD:1-42
All physical covariance/Wick/common-space, strong-$L^4$ random-element,
zero-Bochner-mean, sign, every bounded-bilinear and actual endpoint interfaces
retain their original conditional scopes at each selected finite positive
stage. NR1's sharper author versus larger blind core, classical/source truth,
strict AP6 and unknown original fixed $g_*$ remain explicit. Width counts
whole-plane process neurons with countably many internal variables. The
software computes no canonical least inner indices, practical/minimal/necessary
cost, finite global-weight field simulator, full-field rate, global neural
Schwartz control, broader model/test family or priority. Its internal
exact-version acceptance and this manuscript's fresh review are distinct;
the immutable public v1 is unchanged.

\endgroup
% BEGIN PRIVATE COMPOSITE INTEGRATION 20261106
% Private composite update candidate, 6 November 2026.
% CMAP blocks map every new statement/hypothesis/proof block to frozen sources.
\clearpage
\begingroup
% Keep the frozen contents pages and baseline body pagination unchanged.
\renewcommand{\addcontentsline}[3]{}
\section{Renormalized local squares and finite-neural approximation}
\label{cm:update}
\pdfbookmark[0]{Renormalized local squares and finite-neural approximation}{cm-update}
\newcommand{\cmBL}{d_{\mathrm{BL}}}
\newcommand{\cmU}{d_{\mathcal U_4}}
\newcommand{\cmLtwo}{\operatorname{L2plus}}
\newcommand{\cmEXT}{\operatorname{EXT}}

% CMAP guide | CLAIM-CCA.md:1-83;../limits/G6-DECISION.json:70-71;CC1-STATEMENT-ONLY.md:7-96;CA1-STATEMENT-ONLY.md:183-324
\subsection*{Guide to the local-square update}
The preceding all-dyadic theorem controls its stated family of linear
measurements. This addendum integrates the separately accepted local-square
results CC1-SMOOTH-BYPASS-v1 and CA1-v1: a quantitative approximation of the
renormalized square and a comparison of the actual normalized finite-neural
joint law with the original target pair. All named premises remain below.
Internal acceptance of the source results was recorded on 31 October 2026.
Their complete manuscript integration received source and actual-proof matching
review on 6 November; the public REVIEW document states the review limits.

For the same target, physical subtraction and Borel prescription, CC1 gives
\[
 \E_{\nu_\infty}|Q_s-Q_0|\le12L_f\sqrt{s}<36\sqrt{s}.
\]
Coupling the unchanged first coordinate supplies precisely CA1.5b's joint
bounded-Lipschitz input. Under the \emph{union} of S1--S4, H1--H7 and
Appendix S, take
\[
 r(s)=36\sqrt{s},\qquad S_r(\tau)=2^{-2\cmLtwo(36/\tau)}.
\]
The target allocation is $\eta/8$, hence $j=\cmLtwo(288/\eta)$.
CA1 then selects finite positive stages of joint error at most every rational
$0<\eta\le1$. This is conditional mathematical termination. Useful sampling,
expectation-evaluation runtime, finite-bit global simulation,
interaction-resolving gap, matched advantage, priority and larger-model
claims remain separate. G6 remains \textbf{STOP / INCOMPLETE}; no full
G1--G9 capability is promoted. Neither the stronger numerical $L^2$ density
certificate nor a later tail improvement enters this update.

% CMAP conventions | CA1-STATEMENT-ONLY.md:3-5;CC1-STATEMENT-ONLY.md:7-7
\subsection{Conventions and complete retained premises}
Use $z=(t,y)$, $p=(w,q)$ in two Euclidean dimensions,
\[
 \hat g(p)=\int e^{-ip\cdot z}g(z)\,dz,\quad
 d\lambda(p)=dp/(2\pi)^2,\quad m=1,\quad u=10^{-10},\quad\ell=1000.
\]
There is no factorial in $u$. Write $\Lambda=B_n=(-1000n,1000n)^2$,
$n\in\N$, $n\ge1$, and $A_n=4{,}000{,}000n^2$.
Wick ordering is $:x^2:_c=x^2-c$ and $:x^4:_c=x^4-6cx^2+3c^2$.
All finite-dimensional norms, derivative operator norms and Lipschitz
constants are Euclidean. $\mathrm{BL}_1$ imposes separately
$\sup|G|\le1$ and $\operatorname{Lip}(G)\le1$, so its probability distance
is at most two. Source lattice mesh, flow time and profile are distinct
from $a,\beta,s,\eta$ and the shell parameter $q$.

% CMAP H1 | CA1-STATEMENT-ONLY.md:9-34
\paragraph{H1: original target, common Gaussian space and physical laws.}
The original real complete separable space is
$X=H^{-2}_{\mathrm{loc}}(\R^2;\R)$. Smooth $0\le\chi_j\le1$ equal one
on $[-j,j]^2$, are supported in $(-j-1/2,j+1/2)^2$ and satisfy
$\chi_{j+1}=1$ on $\supp\chi_j$. On $\T_j=\R^2/[2(j+1)\Z]^2$, put
\begin{gather*}
 V_j=4(j+1)^2,\quad e_{j,k}(z)=V_j^{-1/2}e^{i\pi k\cdot z/(j+1)},\quad
 \lambda_{j,k}=1+\pi^2|k|^2/(j+1)^2,\\
 p_j(v)^2=\sum_{k\in\Z^2}\lambda_{j,k}^{-2}
       |v(\chi_j\overline{e_{j,k}})|^2,\qquad
 d_X(v,w)=\sum_{j\ge1}2^{-j}\min(1,p_j(v-w)).
\end{gather*}
Compact smooth evaluations are continuous and generate the stated Borel
sigma-algebra. The original strict AP3/AP5/AP6 route identifies
$\nu_\infty$ as the $X$ representation of the unique common original
Fr\"ohlich--Simon maximal state $v_{(1,10^{-10})}$, preserving every such
evaluation, and gives $\nu_n\Rightarrow\nu_\infty$ in $X$.
The strict cap $0<10^{-10}<1/8000$ retains P1--P6 below; the older unknown
$g_*$ route is distinct.

On the real Hermitian isonormal space $(\Gamma,\mathcal F,P)$,
$H=L^2(\hat C_0\,d\lambda)$, use accepted Borel $X$ versions and continuous
positive-$a$ point versions
\begin{gather*}
 \hat C_0=(1+w^2+q^2)^{-1},\quad
 \hat C_a=(1+w^2+q^2+a^2w^2q^2)^{-1},\quad0<a<1,\\
 \tau_a=(\hat C_a/\hat C_0)^{1/2},\quad\tau_0=1,\quad
 \Phi(g)=W_H(\hat g),\quad\Phi_a(g)=W_H(\hat g\tau_a),\\
 c_a=\int\hat C_a\,d\lambda
 =\frac1{2\pi}\int_0^\infty\frac{dq}{\sqrt{(1+q^2)(1+a^2q^2)}},
 \qquad U_{a,n}=u\int_\Lambda(\Phi_a^4-6c_a\Phi_a^2+3c_a^2).
\end{gather*}
$U_{0,n}$ is the accepted Gaussian $L^2$ limit; these are centered fourth
chaos. At every permitted $a\ge0$ define individually
\[
 Z_{a,n}=\E_Pe^{-U_{a,n}},\quad r_{a,n}=e^{-U_{a,n}}/Z_{a,n},\quad
 \mu_{a,n}=r_{a,n}P,\quad\nu_{a,n}=(\Phi_a)_\#\mu_{a,n},\quad\nu_n=\nu_{0,n}.
\]
The actual same-mass, own-variance, whole-plane free-boundary Wick/classical
identification, Gaussian versions, isometry
$\E I_2(K)I_2(L)=2\langle K,L\rangle$, Gaussian moments and homogeneous-chaos
hypercontractivity remain accepted/classical inputs. Version equalities
are fixed-parameter or specified countable-family equalities.

% CMAP H2 | CA1-STATEMENT-ONLY.md:36-50
\paragraph{H2: fixed-box quantitative physical/common-space input HC.}
Define
\begin{gather*}
 F_2=\|f\|_2,\quad H_f=F_2\sqrt{2/\pi},\quad L_f=F_2\sqrt{2\sqrt5/\pi},
 \quad\kappa_2=\frac{\sqrt{17}}{8\sqrt\pi},\quad
 M_6=(2\pi)^{-5/3}(5\pi)^{5/6},\\
 K_n=u\sqrt{384A_n\kappa_2M_6^3}<n/10^6,\quad
 h_2(n)=\max(n^2/100,n^4/125),\quad m_2(n)=1+\lceil4h_2(n)\rceil,\\
 \ell_n=\lceil m_2(n)/2\rceil,\quad M_n=2^{\ell_n},\quad C_n=3+4n/10^6,
 \qquad R_n(a)=M_n\{aF_2/2+(H_f+4K_n)\sqrt a\}.
\end{gather*}
Retain CMP-FV-v1 P1--P4 and literal C2/C4/C5/C6 at every fixed $n$:
\[
 1\le Z_{a,n}\le M_n,\quad\|r_{a,n}\|_2\le M_n,\quad
 \|r_{a,n}-r_{0,n}\|_1\le4K_nM_n\sqrt a.
\]
For the H4 own-smoothed pairs,
\begin{equation}
 \cmBL\bigl(\Law_{\mu_{a,n}}(\Phi_a(f),Y_{a,s}),
             \Law_{\mu_{0,n}}(\Phi(f),Y_s)\bigr)
 \le\min(2,R_n(a)).\tag{HC}\label{cm:HC}
\end{equation}
The common-space density and actual field pushforward are distinct.
Fixed-box raw-square and Borel-prescription qualifications remain;
$Q_0(\Phi_a)=\int f(\Phi_a^2-c_a)$ is not hypothesized.

% CMAP H3 | CA1-STATEMENT-ONLY.md:52-84
\paragraph{H3: finite-linear response, moments and local target density.}
For every real compact smooth $h$ and $t\in\R$,
\[
 \E_\zeta e^{t v(h)}\le e^{(3/2)t^2\|h\|_2^2},\qquad
 \zeta\in\{\nu_n:n\ge1\}\cup\{\nu_\infty\}.
\]
Accepted MOM all-positive local $p_j$ moments, uniform integrability and
polynomial-local-growth expectation transfer remain at their exact scopes.
The unrestricted finite-linear-vector AP6 interface is literal as follows.
For $g=(g_i)$ supported in $K\subset[-R,R]^2$ inside open $B_{n_0}$, set
\begin{gather*}
 \overline F_g=(2R+1/2)(\sum_i\|g_i\|_\infty^2)^{1/2},\quad
 C_g=P_*\overline F_g,\quad q=e^{-125/32},\\
 \Theta_n=q^n\{(2n+1)/(1-q)+2q/(1-q)^2\},\\
 \cmBL(g_\#\nu_n,g_\#\nu_\infty)\le4C_gu\ell^2e^{R/256}\Theta_n
 \quad(n\ge n_0,\ n\ell\ge R).
\end{gather*}
It includes singular/repeated tuples and the norm-preserving BL extension.
The chosen OPT coefficient version has the same $C_H<2^{16}$,
$C_{H2}\le2^{384}$ and $C_{\mathrm{base}}=P_*(5/2)\|f\|_\infty\le2^{360}$.
For literal identification of $P_*$ put
\begin{gather*}
 \bar\delta=2\ell^{-2}\log[1/(1-e^{-4\ell})],\quad
 g_{\mathrm{flow}}(t)=1+\bar\delta+(\pi/16)\log(1+1/t),\\
 c_{\mathrm{CT}}=81/5,\quad b_{\mathrm{CT}}=1+c_{\mathrm{CT}}/256,\quad
 \Gamma_G=4\sqrt{2\pi}/256,\\
 C_*=6c_{\mathrm{CT}}/256+c_{\mathrm{CT}}^2/256^2
       +b_{\mathrm{CT}}^2(\Gamma_G+3300u),\\
 M_{\mathrm{flow}}(t)=32\ell^{-2}
 \{C_H^2g_{\mathrm{flow}}(t)^3/12+u^2C_{H2}^2g_{\mathrm{flow}}(t)^7/9\},\\
 P_*=e^{8\sqrt2/256}b_{\mathrm{CT}}^2
 \{e^{C_*}\int_0^1\sqrt{M_{\mathrm{flow}}(t)}\,dt
        +e^{2C_*}\sqrt{M_{\mathrm{flow}}(1)}/2\}.
\end{gather*}
Source time $t$ is separate from $s$. This general AP6 interface is not an
extension of UDA's restricted ordered one-through-four bump theorem.

Let $D_{\mathrm{loc}}=(-3,3)^2$ and let $\mathcal C_D$ be generated by real
compact smooth evaluations supported inside it. The same target has a
normalized density relative to the \emph{restricted whole-plane mass-one}
Gaussian law:
\begin{gather*}
 k_D\ge0,\quad \E_{P_1}k_D=1,\quad K_D=\|k_D\|_{L^2(P_1)}<\infty,\\
 k_D=e^{-uJ_D-\alpha_1(uX^4)|D|}g_{\partial D},\qquad
 g_{\partial D}>0\ \hbox{a.e.},\quad D=D_{\mathrm{loc}}.
\end{gather*}
The boundary-algebra alternatives and all-semibounded-polynomial scope
remain: the factor is measurable either for the intersection of neighborhood
algebras or for the completed algebra generated by boundary-supported
$H^{-1}$ evaluations. Pressure and boundary factor belong to this density;
no finite-box partition function or Dirichlet Gaussian is substituted.
A numerical upper certificate for $K_D$ is additional only in CA1.5.

% CMAP H4 | CA1-STATEMENT-ONLY.md:86-101
\paragraph{H4: fixed observables and Gaussian composite estimates.}
Define
\begin{gather*}
 b(t)=\begin{cases}e^{-1/(1-t^2)},&|t|<1,\\0,&|t|\ge1,\end{cases}\quad
 f(t,y)=b(t)b(y),\quad I=\int f,\quad\rho=f/I,\\
 R_2=\|\rho\|_2,\quad D_\rho=\|\nabla\rho\|_2,\quad
 \rho_s=s^{-2}\rho(\cdot/s),\quad h_{s,z}(x)=\rho_s(z-x),\quad0<s\le1,\\
 c_s=\int|\hat\rho(sp)|^2\hat C_0(p)\,d\lambda,\quad
 c_{a,s}=\int|\hat\rho(sp)|^2\hat C_a(p)\,d\lambda,\\
 X_{\mathrm{obs}}(v)=v(f),\quad
 Q_s(v)=\int f(z)\{v(h_{s,z})^2-c_s\}\,dz,\\
 Y_{a,s}=\int f\{(\rho_s*\Phi_a)^2-c_{a,s}\},\qquad Y_s=Q_s(\Phi).
\end{gather*}
$Q_s$ is the actual continuous/Borel local $X$ functional. Set
$Q_0(v)=\lim_{j\to\infty}Q_{2^{-2j}}(v)$ on its Borel convergence set,
zero elsewhere, $j\ge1$. This is the same physical local Borel Wick
prescription under $\nu_n,\nu_\infty$. On the Gaussian space,
\[
 Q_0(\Phi)=I_2(\hat f(p+q))\ \hbox{a.s.},\qquad
 \|Y_s-Q_0(\Phi)\|_2\le L_f\sqrt s.
\]
The Gaussian null set and target-local absolute-continuity qualifications
remain. Write $\mathcal L_{\mathrm{target}}=\Law_{\nu_\infty}
(X_{\mathrm{obs}},Q_0)$, with second coordinate $O_2(f)$. Literal envelopes:
\[
 1/27<I<1,\quad F_2<1/2,\quad\|f\|_{4/3}<1,\quad
 R_2\le54,\quad D_\rho\le81,\quad H_f<2,\quad L_f<3.
\]
The weaker $I\le4,F_2\le2$ are allowed. All $h_{s,z}$ on $\supp f$ are
supported in $[-2,2]^2$, strictly inside $D$ and every open $B_n$.
The fixed query is $F(x,y)=\cos((x+y)/2)$.

% CMAP H5 | CA1-STATEMENT-ONLY.md:105-132
\paragraph{H5: original actual source laws RL2/SP1--SP4.}
The product neuron is $g^{\mathrm{product}}(t,y)=\sigma(t)\tau(y)/2$ for
independent stationary symmetric two-sided telegraphs with flip rate $1/(2a)$.
The Mondrian neuron is $g^{\mathrm{Mondrian}}(z)=\eta_{\mathrm{cell}(z)}/2$
from the original stationary projective axis-parallel whole-plane partition
at budget $1/a$, with conditionally independent fair cell signs. A rectangle
of sides $w,h$ splits at rate $w+h$, vertically with probability $w/(w+h)$
and uniform cut; daughters evolve independently for the remaining budget.
Use the lower-left half-open representative modulo area-null boundaries.
Whole-plane existence, stationarity/projectivity, path measurability,
copies and actual architecture/potential identification remain SP1--SP4.

The original centered continuous Gaussian endpoint $\gamma$ has
\[
 C_{\mathrm{end}}(z-z')=e^{-|z-z'|_1/a}/4,\qquad
 \hat C_{\mathrm{end}}=\frac{a^2}{(1+a^2w^2)(1+a^2q^2)}.
\]
For finite $R_{\mathrm{aux}}\ge1000n$, set
$D_{\mathrm{aux}}=[-R_{\mathrm{aux}},R_{\mathrm{aux}}]^2$,
$d=4R_{\mathrm{aux}}^2$, $E=L^4(D_{\mathrm{aux}};\R)$, $K=a^{-2}-1$,
$F_N^{\mathfrak b}=N^{-1/2}\sum_i g_i^{\mathfrak b}$,
$\rho_N^{\mathfrak b}=\Law_E(F_N^{\mathfrak b})$,
$\rho_\gamma=\Law_E(\gamma)$, $\mathfrak b\in\{\mathrm{product},\mathrm{Mondrian}\}$.
These actual strongly measurable $E$ random elements have zero Bochner
means, global sign symmetry, fourth-norm moments $d/16$ and $3d/16$,
respectively, and finite second-norm moments. For \emph{every} bounded real
bilinear form $B_0$, including nonsymmetric ones,
\[
 \E B_0(g^{\mathfrak b},g^{\mathfrak b})=\E B_0(\gamma,\gamma).
\]
Independent full-law copies exist; $N^{-1/2}\sum\gamma_i$ has exactly
$\rho_\gamma$. The finite-rank/every-bilinear route retains accepted NR1
Astra/max attribution and exact-author/larger-blind-core qualification;
RL2 consolidation and normalized-law work retain recorded Sol/max attribution.

For $\beta>0$ define on $E$
\begin{gather*}
 W(v)=u\int_\Lambda(v^4-6c_av^2+3c_a^2),\quad
 J_{\mathrm{out}}(v)=\int_{D_{\mathrm{aux}}\setminus\Lambda}v^4,\\
 V(v)=W(v)+\beta J_{\mathrm{out}}(v)-(K/2)\int_{D_{\mathrm{aux}}}v^2,\quad w=e^{-V},\\
 Z_N^{\mathfrak b}=\rho_N^{\mathfrak b}(w),\quad Z_\gamma=\rho_\gamma(w),\quad
 P_N^{\mathfrak b}=w\rho_N^{\mathfrak b}/Z_N^{\mathfrak b},\quad
 Q_\beta=w\rho_\gamma/Z_\gamma.
\end{gather*}
$T_s$ denotes convolution after zero extension from $D_{\mathrm{aux}}$,
agreeing with the whole-plane convolution on $\supp f$. Put
\begin{gather*}
 B_T(v,h)=\int f(Tv)(Th),\quad
 J_{T,k}(v)=(\int_{D_{\mathrm{aux}}}fv,B_T(v,v)-Ik),\\
 J_s=J_{T_s,c_s},\qquad
 \mathcal L_{\mathrm{neural}}^{\mathfrak b}=(J_s)_\#P_N^{\mathfrak b}.
\end{gather*}
These are full tilted actual source laws with distinct normalizers.
No covariance-to-nonlinear-law identification is assumed.

% CMAP H6 | CA1-STATEMENT-ONLY.md:134-154
\paragraph{H6: literal HP mass tilt, residual laws and defect coupling.}
For the whole-plane endpoint law $G_{\mathrm{end}}$ put
\begin{gather*}
 M_{D_{\mathrm{aux}},a}=\E_{G_{\mathrm{end}}}
       e^{(K/2)\int_{D_{\mathrm{aux}}}v^2}\in(0,\infty),\quad
 G'=e^{(K/2)\int_{D_{\mathrm{aux}}}v^2}G_{\mathrm{end}}/M_{D_{\mathrm{aux}},a},\\
 z_\beta=\E_{G'}e^{-W-\beta J_{\mathrm{out}}},\quad z_0=\E_{G'}e^{-W},\quad
 Q_{\mathrm{res}}^\beta=e^{-W-\beta J_{\mathrm{out}}}G'/z_\beta,\quad
 Q_{\mathrm{res}}^0=e^{-W}G'/z_0.
\end{gather*}
The genuine centered continuous whole-plane Gaussian $G'=G'_{D_{\mathrm{aux}},a}$
has the uniformly convergent covariance series, $x_0=z,x_{j+1}=z'$,
\[
 C'(z,z')=C_{\mathrm{end}}(z-z')+
 \sum_{j\ge1}K^j\int_{D_{\mathrm{aux}}^j}
          \prod_{i=0}^j C_{\mathrm{end}}(x_i-x_{i+1})\,dx_1\cdots dx_j.
\]
The all-Borel whole-path/restriction mass-tilt identity and
$Z_\gamma=M_{D_{\mathrm{aux}},a}z_\beta$ remain; $Q_\beta$ is the
$E$ restriction of $Q_{\mathrm{res}}^\beta$. The physical positive-$a$
law is $e^{-W}G_a/Z_{a,n}$ with its own normalizer.

$\Delta_{D_{\mathrm{aux}},a}=C_a-C'$ is a continuous positive semidefinite
kernel; $C'$ point variances lie in $[0,c_a]$. Independent continuous
centered Gaussian $Y\sim G'$ and $\psi\sim G_\Delta$ exist with
$X_a=Y+\psi\sim G_a$. For compact $H$ put
$\delta_D(H)=\sup_{z\in H}\Delta(z,z)$. If $D_{\mathrm{aux}}$ contains
its Euclidean $r$ neighborhood, $r>0$,
\[
 \delta_D(H)\le\xi_a(r)=K\int_{|x|>r}C_a(x)^2\,dx.
\]
Separate fixed-$D$ couplings suffice. All physical/common-space,
CT G-tilt/G-order/G-diff, classical and source qualifications remain.

% CMAP H7 | CA1-STATEMENT-ONLY.md:156-179
\paragraph{H7: literal normalized NC1 stage estimate.}
Here let $\alpha=\min(u,\beta)$, $A=\max(u,\beta)$, $B=K/2+6uc_a$,
$L=B^2d/(2\alpha)$. These $\alpha,L$ are distinct from pressure and the
CC1 smoothing constant. For $x\ge0$ put
\begin{align*}
 v_1&=4Ax^3+2B\sqrt d\,x,&v_2&=12Ax^2+2B\sqrt d,&v_3&=24Ax,&v_4&=24A,\\
 e_1&=v_1,&e_2&=v_1^2+v_2,&e_3&=v_1^3+3v_1v_2+v_3,\\
 e_4&=v_1^4+6v_1^2v_2+3v_2^2+4v_1v_3+v_4,\\
 q_1&=1+2x,&q_2&=2,&p_1&=q_1,\\
 p_2&=q_1^2+q_2,&p_3&=q_1^3+3q_1q_2,&p_4&=q_1^4+6q_1^2q_2+3q_2^2.
\end{align*}
For a nonnegative-coefficient $P=\sum_jP_jx^j$, let $a_0=1$,
$a_j=[j/(2\alpha e)]^{j/4}$ for $j>0$, $\mathcal T(P)=e^L\sum_jP_ja_j$,
and
\begin{gather*}
 P_Z=e_4,\quad P_J=e_4+4p_1e_3+6p_2e_2+4p_3e_1+p_4,\quad
 M_Z=\mathcal T(P_Z),\quad M_J=\mathcal T(P_J),\\
 z=e^{-3Ad/16-3uc_a^2A_n},\quad C=\frac{d(M_J+M_Z)}{96z},\quad
 D_J=\frac43(3\sqrt6C)^{1/4}(2\sqrt2)^{3/4}.
\end{gather*}
The exact imported conclusions are positive bounded Borel $w$,
$Z_N^{\mathfrak b},Z_\gamma\ge z>0$, all-Borel actual source-law identities
of restrictions and full tilts, and, uniformly in $\mathfrak b$, $N\ge1$,
$0<s\le1$ and every finite common $k$,
\begin{gather*}
 \cmBL((J_{T_s,k})_\#P_N^{\mathfrak b},(J_{T_s,k})_\#Q_\beta)
        \le\min(2,D_JN^{-1/4}),\\
 \cmU((J_{T_s,k})_\#P_N^{\mathfrak b},(J_{T_s,k})_\#Q_\beta)
        \le\min(2,C/N).
\end{gather*}
$\mathcal U_4$ bounds value and derivative operator norms through order four
by one. For $0<\eta_W\le1$,
$N\ge\lceil512\sqrt3 C\eta_W^{-4}\rceil$ suffices for error $\le\eta_W$.
All $a,\beta,n,R_{\mathrm{aux}}$ dependence remains. No $a=0,\beta=0,s=0$
stage is included.

% CMAP S1 | CC1-STATEMENT-ONLY.md:7-9
\paragraph{S1: CC1 model and source qualification.}
CC1 uses exactly H1's model, boxes, Fourier convention, original $X$ Borel
space and common FS state, the actual weak-$X$ limit of the original
same-mass free-boundary Wick-quartic laws. The strict route retains literal
P1--P6 and accepted H0--H5 source hypotheses; the original numerical $g_*$
inequality is uncomputed and unused. P1--P6 include finite ferromagnetic
DSS/SK-$q$ for all finite real physical fields, finite-positive-scale $P_*$
susceptibility, fixed-square Gaussian/Wick exponential transfer, fixed-mesh
periodization, original FS pressure/maximal-state/free-limit facts, and
accepted physical/common-space/local-density/finite-response interfaces.
Individual source, classical, H1/MIX/REG/FLOW/TP/T and NR1 qualifications
remain. RL2/SP1--SP4 and NC1 HP/HC condition the later assembly; CC1 alone
makes no new assertion about those arrows.

% CMAP S2 | CC1-STATEMENT-ONLY.md:11-31
\paragraph{S2: Gaussian reference, local interaction and Nelson input.}
$P_1$ is the centered whole-plane Gaussian with multiplier $\hat C_0$;
only its sigma-algebra is restricted. The common realization and
$\Phi_a,\tau_a,c_a$ are H1's. Let $D=(-3,3)^2$, $A_D=36$ and
$U_D=uJ_D$, the literal coefficient-one mass-one Gaussian fourth Wick
integral, equivalently the Gaussian $L^2$ limit as $a\downarrow0$ of
\[
 U_{a,D}=u\int_D[\Phi_a(z)^4-6c_a\Phi_a(z)^2+3c_a^2]\,dz.
\]
$\E_{P_1}U_D=0$. Retain Gaussian construction, general-rectangle Nelson
estimate, second-chaos isometry, finite-chaos hypercontractivity, Plancherel
and elementary measure/real analysis at their named classical scopes. Put
\begin{gather*}
 \kappa_0=e(1+\sqrt2),\quad K_{A_D}=u\sqrt{384A_D\kappa_2M_6^3},\\
 K^\sharp_{A_D}=\max(12uA_D,(2/e)K_{A_D}\sqrt{\kappa_0}),\qquad
 \beta_{A_D}=\pi/(2\sqrt{12uA_D}).
\end{gather*}
The retained zero-regulator tail is
$P_1(U_D<-v)\le\exp[-\exp(\beta_{A_D}\sqrt v-1)]$ for
$v\ge K^\sharp_{A_D}$. All parameters refer to this same interaction and
whole-plane mass-one reference.

% CMAP S3 | CC1-STATEMENT-ONLY.md:33-46
\paragraph{S3: actual normalized factorization and numerical source.}
Write $\alpha_{\mathrm{pr}}=\alpha_1(uX^4)$ for the original normalized
log-partition pressure. The density is
\[
 k_D=e^{-U_D-36\alpha_{\mathrm{pr}}}g_{\partial D},\quad
 k_D\ge0,\quad\E_{P_1}k_D=1.
\]
The same positive boundary factor and BC1/BC2 readings in H3 remain;
no isolated finite-box normalizer enters. Inherited H5 supplies $K_D<\infty$
without its value. The accepted primary supplement FS1977-v1 matches
Fr\"ohlich--Simon Lemma 7.1 and Theorem 7.2 proof, printed p.~517,
equation (7.1), for the original Section 4 maximal states including their
common value. For an $l$-square $E$, mass $m_0$, and
$q\ge4(1-e^{-m_0l})^{-2}$, every nonnegative bounded $\mathcal C_E$
measurable $H$ satisfies
\[
 \nu_\infty(H)\le e^{-|E|\alpha_1(uX^4)}
                    \|He^{-uJ_E}\|_{L^q(P_1)}.
\]
Here $m_0=1,E=D,l=6$. This attributed fact retains literal same-mass,
pressure, polynomial and state-class scope
\citep[Lemma 7.1 and proof of Theorem 7.2, p.~517]{FS1977}.
The accepted investigation inspected primary text and page image; this
integration transcribes that frozen premise. Its GRS checkerboard
dependency is retained \citep{GRS1976}. The numerical clause is a matched
primary input, not a consequence of the old qualitative H5 alone.

% CMAP S4 | CC1-STATEMENT-ONLY.md:48-58
\paragraph{S4: CC1 observables.}
Use H4's $b,f,I,F_2,\rho,\rho_s,h_{s,z},c_s,Q_s$ and countable Borel
$Q_0$. Set
\[
 r_s(p)=\hat\rho(sp),\quad K_f(p,q)=\hat f(p+q),\quad
 O_2^P=I_2(K_f),\quad \sigma=F_2/\sqrt{2\pi},\quad L_f=F_2\sqrt{2\sqrt5/\pi}.
\]
The isometry has factor two. Under the actual target $Q_0=O_2(f)$.
All smoothed tests stay strictly inside $D$. Subtraction is exactly $c_s$
under target and physical laws. The joint query remains
$X_{\mathrm{obs}}(v)=v(f)$, $F(x,y)=\cos((x+y)/2)$.

% CMAP source-scope | CA1-STATEMENT-ONLY.md:301-324
\subsubsection*{Appendix S: full retained source scope}
These are attributed accepted-source hypotheses with original source and
review qualifications, rather than fresh proofs.

\paragraph{P1.}
On $\mathbb T=\epsilon(\Z+1/2)^2/L\Z^2$, $L/\epsilon$ even and at least
four, use physical inner product $\epsilon^2\sum$, $A=-\Delta^\epsilon+m^2$
and the \emph{plane} variance $c_\epsilon$. For every profile
$q:\mathbb T\to[0,1]$, the law proportional to
$e^{-u\epsilon^2\sum q_x:\psi_x^4:_{c_\epsilon}}dP_{A^{-1}}$ under every
finite real physical field $H$ satisfies
$0\le\operatorname{Cov}_{\nu^H}(\psi_x,\psi_y)\le\E_\nu\psi_x\psi_y$.
The DSS/SK-$q$ clauses with added mass $1/t$ hold at all finite $t>0$.
Retain BD/BBD/Griffiths--Simon, Ising-FKG and Robbins--Stirling qualifications;
no $g_*$ premise is added.

\paragraph{P2.}
For $0<\epsilon\le\min(1/4,18/(5m))$, $L\ge\max(1,\sqrt2/m)$,
admissible even $L/\epsilon\ge4$ and $0<u\le m^2/8000$, the original core
$P_*$ at $S_{\max}=\infty,\Phi=1,\kappa=0$ gives at every finite $t>0$
and profile $q$
\begin{gather*}
 \chi_{t,q}=\max_x\epsilon^2\sum_y\E_{\mu_{t,q}}\psi_x\psi_y
 \le t/(1+m^2t)+2t/(2+m^2t),\\
 d\mu_{t,q}\propto e^{-u\epsilon^2\sum q_x:\psi_x^4:_{c_\epsilon}}
                          dP_{(A+1/t)^{-1}}.
\end{gather*}
It retains $\|S-C_t\|_{1,0}\le2t/(2+m^2t)$; P2 alone supplies no
original-law endpoint conclusion.

\paragraph{P3.}
At fixed $B=(-b,b)^2$, commensurate plane meshes $\epsilon=b/N$ have the
accepted T1--T4/GQ/W($b$)/W-laws common Gaussian realization:
$J_B^\epsilon=\epsilon^2\sum_{x\in B}:\psi_x^4:_{c_\epsilon}$ and
$Y_f^\epsilon=\epsilon^2\sum f(x)\psi_x$ converge in Gaussian $L^2$
to $J_B,\Phi(f)$. For each fixed finite $p\ge1$,
$\sup_{\epsilon\text{ sufficiently small}}\E e^{-puJ_B^\epsilon}<\infty$.
Positive finite continuum normalizers, every finite Gaussian $L^p$
density, coefficient-one Wick conventions and exact finite-law
identifications remain at all $m,u>0$.

\paragraph{P4.}
At fixed mesh torus covariances converge entrywise on finite fixed site sets
to plane covariance through admissible even periods, with the same plane
Wick variance. This is the accepted H1-PREQ periodization/envelope premise.

\paragraph{P5.}
Use the original FS same-mass whole-plane free-boundary convention and
$P_h(X)=uX^4-hX$. Its normalized pressure and dyadic-square realization
are original Proposition 2.1 and the remark after (2.3). Corollary 4.3(ii)'s
converse equates differentiability at $h=0$ with equality of the maximal
states; Theorem 5.2 supplies ordinary-exhaustion convergence to the common
state, in the compact generating-functional reading at every real frequency
\citep{FS1977}. These are attributed published facts; full original FS
proofs were not supplied/read in this integration.

\paragraph{P6.}
Retain (i) the all-$m,u>0$ physical/classical identity and full original
local $H^{-2}$/Borel/compactness definitions; (ii) all-semibounded-polynomial
original FS local $L^p$ density and positive boundary factorization under
either boundary-algebra reading; (iii) exact finite DSS/SKQ/$P_*$/H1-PREQ/
H1/MIX/REG/smoothed-CT/FLOW/TP/T interfaces, individual accepted decisions
and source-normalization correction, including FLOW's $0<u\le m^2/8000$
cap and fixed-square TP/T1--T4 scope; and (iv) raw fixed-stage NR1/L0 and
Gaussian/common-space/Wick normalizers, original author-constant/larger-blind-core
qualification, actual Gaussian mass-tilt/independent-defect coupling and
physical own-variance identification. Old EB1/FV/CL1 with global unknown
$g_*$ are excluded outside their literal domains; AHKZ C2(ii)/(iii) are
not added.

Named FS, AHKZ, GRS and SG, DSS, BD, BBD and Griffiths--Simon,
Gaussian/Wick/Nelson, H1/MIX/FLOW/TP/T, NR1 attribution/review and source-access qualifications
remain. Original FS/OPT full bodies are not a fresh audit here. Frozen
CV1/NC1/CMP/RL2 decisions preserve all inherited adverse findings at
original object, severity and disposition. The original private proof map distinguishes
actual reading in that integration from these transcribed premises.

% CMAP cc-statement | CC1-STATEMENT-ONLY.md:62-96
\subsection{Complete CC1 conclusions}
\begin{theorem}[CC1-SMOOTH-BYPASS-v1, under S1--S4]\label{cm:CC1}
The following conclusions hold for the fixed original target.
\begin{enumerate}[label=\textup{CC1.\arabic*},leftmargin=4.6em]
\item The local Gaussian exponential certificate is
\[
 K_{A_D}<128u,\quad K^\sharp_{A_D}=432u,\quad\beta_{A_D}>1500,\quad
 \E_{P_1}e^{-10U_D}\le19/9<4,\quad\|e^{-U_D}\|_{10}\le4/3.
\]
\item For the literal pressure and boundary factor,
\[
 \alpha_{\mathrm{pr}}\ge0,\qquad
 \|g_{\partial D}\|_{5/4}\le1,\qquad \|k_D\|_{10/9}\le4/3.
\]
This is a different exponent from $K_D=\|k_D\|_2$ and does not certify
the latter numerically.
\item The same Borel $Q_0$ agrees with $O_2^P$ under $P_1$ and its
convergence set has full target probability. For $0<s\le1$ and real $r\ge1$,
\begin{gather*}
 \|Q_s-Q_0\|_{L^r(\nu_\infty)}
 \le(4/3)^{1/r}(10r-1)L_f\sqrt s,\\
 \|Q_0\|_{L^r(\nu_\infty)}
 \le(4/3)^{1/r}(10r-1)\sigma.
\end{gather*}
In particular
$\E_{\nu_\infty}|Q_s-O_2(f)|\le12L_f\sqrt s<36\sqrt s$.
Every fixed positive moment order of $\{Q_s:0<s\le1\}$ is uniformly
integrable under the target.
\item On positive rationals define
$\cmLtwo(x)=\max(0,\lceil\log_2x\rceil)$ by exact integer comparisons.
For positive rational $\tau$, choose
$j=\cmLtwo(36/\tau)$ and $s=2^{-2j}$. Then $0<s\le1$ and the target
$L^1$ smoothing error is $\le\tau$. For an allocation $\epsilon/4$,
$j=\cmLtwo(144/\epsilon)$ suffices. Symbolic dyadic output is permitted.
\item The stronger numerical $L^2$ density norm remains uncertified by
CC1. The Gaussian reference, target, positive boundary factor, pressure
sign and subtraction are unchanged. CC1 alone supplies target-local
smoothing, without raw fixed-$a$ square limit, finite-neural assembly,
sampler, normalized estimator, expectation runtime, matched comparator,
positive interaction gap, full-work benefit or novelty-priority claim.
No full G1--G9 flag is promoted. Its historical acceptance is governed
by the frozen separate review/audit decisions, rather than the original
author candidate header.
\end{enumerate}
\end{theorem}

% CMAP ca1-1 | CA1-STATEMENT-ONLY.md:183-212
\subsection{Complete CA1 conclusions}
\begin{theorem}[CA1.1: actual-law eight-charge budget]\label{cm:CA1-budget}
Assume H1--H7 and Appendix S. Fix $n\ge1$, $0<a<1$, $\beta>0$,
finite $R_{\mathrm{aux}}\ge1000n$, $N\ge1$, $0<s\le1$, $T>0$ finite.
Let $K_f=[-1,1]^2$, $K_*=[-2,2]^2$,
$L_c=(4\sqrt{15}+12)c_a^{3/2}$, and put
\begin{align*}
 E_{\mathrm{width}}&=D_JN^{-1/4},\\
 E_{\mathrm{stab}}&=\min(2,6\beta|D_{\mathrm{aux}}\setminus\Lambda|c_a^2
                                   e^{9uA_nc_a^2}),\\
 E_{\mathrm{aux}}&=\min\bigl(2,e^{9uA_nc_a^2}
  [2uA_nL_c\sqrt{\delta_D(\overline\Lambda)}
  +I\sqrt{\delta_D(K_f)}
  +2I\sqrt{c_a}\sqrt{\delta_D(K_*)}]\bigr),\\
 E_{\mathrm{reg}}&=R_n(a)\le M_nC_n\sqrt a,\qquad
 E_{\mathrm{var}}=d_{a,s}=I(c_s-c_{a,s}),\quad0\le E_{\mathrm{var}}\le27(a/s)^2,\\
 C_s^*&=(9/5)2^{360}\sqrt{1+1/(Is^4)},\qquad
 E_n^*(s)=(1/2500)e^{1/128}C_s^*\Theta_n,\\
 B_s&=3IR_2^2/s^2\le324/s^2,\qquad\Lambda_T=\max(1,2T),\\
 E_{\mathrm{vol}}&=\Lambda_T E_n^*(s),\qquad
 E_{\mathrm{clip}}=4\sqrt2 e^{-T^2/(4B_s)},\qquad
 E_{\mathrm{target}}=K_DL_f\sqrt s .
\end{align*}
For both architectures with the same parameters,
\begin{equation}
 \cmBL(\mathcal L_{\mathrm{neural}}^{\mathfrak b},\mathcal L_{\mathrm{target}})
 \le\min(2,E_{\mathrm{width}}+E_{\mathrm{stab}}+E_{\mathrm{aux}}
       +E_{\mathrm{reg}}+E_{\mathrm{var}}+E_{\mathrm{vol}}
       +E_{\mathrm{clip}}+E_{\mathrm{target}}).
 \tag{CA1-BL}\label{cm:budget}
\end{equation}
The seven successive normalized pair laws, with $r_D$ denoting restriction
to $D_{\mathrm{aux}}$, are
\begin{align*}
 &(J_s)_\#P_N^{\mathfrak b},\\
 &(J_s)_\#Q_\beta=(J_s r_D)_\#Q_{\mathrm{res}}^\beta,\\
 &(J_s r_D)_\#Q_{\mathrm{res}}^0,\\
 &(X_{\mathrm{obs}},Q_s)_\#\nu_{a,n},\\
 &(X_{\mathrm{obs}},Q_s)_\#\nu_n,\\
 &(X_{\mathrm{obs}},Q_s)_\#\nu_\infty,\\
 &(X_{\mathrm{obs}},Q_0)_\#\nu_\infty.
\end{align*}
The regulator arrow has charges $E_{\mathrm{reg}},E_{\mathrm{var}}$ and
the volume arrow has $E_{\mathrm{vol}},E_{\mathrm{clip}}$, each once.
Every law retains its normalizer in H1,H3,H5,H6. Alternatively
$E_{\mathrm{target}}$ may be replaced by any supplied certified joint-law
bound $r(s)$ for the last arrow; the $K_DL_f\sqrt s$ and $r(s)$ options
are alternatives, never summed. For the fixed cosine one may replace
$E_{\mathrm{width}}$ by $C/N$, leaving other charges unchanged.
These bound expectation differences; no Gaussian-comparator separation follows.
\end{theorem}

% CMAP ca1-2 | CA1-STATEMENT-ONLY.md:214-222
\begin{proposition}[CA1.2: subtractions and order]\label{cm:translations}
For each path/restriction,
\begin{gather*}
 J_{T,k_2}=J_{T,k_1}-(0,I(k_2-k_1)),\qquad
 Q_s(\Phi_a)=Y_{a,s}-d_{a,s},\\
 J_{T_s,c_s}=J_{T_s,c_{a,s}}-(0,d_{a,s}),\\
 [B_{\mathrm{id}}(v,v)-Ic_a]-[B_{T_s}(v,v)-Ic_s]
       =\int f[v^2-(T_sv)^2]+I(c_s-c_a).
\end{gather*}
Common translation preserves both distances exactly; translating one law
costs at most $I|k_2-k_1|$. The raw action uses exact $c_a$.
The untilted raw endpoint/neural variance is $1/4$, and its smoothed value is
$c_{\mathrm{end},s}=\int|\hat\rho(sp)|^2\hat C_{\mathrm{end}}\,d\lambda$;
interacting expectations are separate. Endpoint-own to physical subtraction
is the shift $+(0,I(c_{\mathrm{end},s}-c_s))$, absent from
\eqref{cm:budget} because adjacent laws use common $c_s$. No ordering of
$c_s$ and raw $c_a$ is asserted. Successive limits are allowed in the order
$N\to\infty$, $\beta\to0$, $R_{\mathrm{aux}}\to\infty$, $a\to0$ at fixed
$n,s$, $n\to\infty$ at fixed $s$ with clipping controlled, then
$s\to0$ at the target; reverse nested tolerance choices are allowed.
The diagonals below are separately controlled. At fixed $a,N$ the neural
physical-$C_0$-subtracted coordinate tends to $-\infty$ as $s\to0$;
fixed-stage tightness does not follow.
\end{proposition}

% CMAP ca1-3 | CA1-STATEMENT-ONLY.md:224-245
\begin{proposition}[CA1.3: explicit tail and finite rational extension]\label{cm:EXT}
For every $0<a<1$ and $r>0$,
\begin{equation}
 \int|x|^2C_a(x)^2\,dx\le1/16,\qquad
 \xi_a(r)\le K/(16r^2).\tag{CA1-TAIL}\label{cm:tail}
\end{equation}
For integers $n,k\ge1$, $a=2^{-2k}$ and rational
$0<\eta_A,\eta_S,\eta_W\le1$, the following $\cmEXT$ returns finite
$R_{\mathrm{aux}},\beta,N$ with auxiliary, stabilizer and width charges
at most $\eta_A,\eta_S,\eta_W$ for both architectures:
\begin{gather*}
 \bar c=2k+3,\quad\bar E=4^{\lceil9uA_n\bar c^2\rceil},\quad
 \bar H=56uA_n\bar c^2+4+8\bar c,\\
 r=\max(1,\lceil\bar E\bar H/(4a\eta_A)\rceil),\quad
 R_{\mathrm{aux}}=1000n+r+1,\quad d=4R_{\mathrm{aux}}^2,\\
 \beta=\min(u,\eta_S/(6d\bar c^2\bar E)),\quad\alpha=\beta,\quad A=u,\\
 \bar B=K/2+6u\bar c,\quad\bar L=\bar B^2d/(2\alpha),\quad
 \overline Z_{\exp}=3ud/16+3u\bar c^2A_n.
\end{gather*}
Form $\overline P_Z,\overline P_J$ from H7 by replacing $B$ by $\bar B$
and $\sqrt d$ by the exact rational $2R_{\mathrm{aux}}$. Put $t_0=1$,
$t_j=\max(1,j/(2\alpha))^j$ for $1\le j\le12$ and
\begin{gather*}
 \bar C=\frac d{96}4^{\lceil\bar L\rceil+\lceil\overline Z_{\exp}\rceil}
          \sum_{j=0}^{12}(\overline P_Z+\overline P_J)_j t_j,\\
 N=\max(1,\lceil1024\bar C\eta_W^{-4}\rceil).
 \tag{EXT}\label{cm:EXT-formula}
\end{gather*}
Subscripts are polynomial coefficients, absent coefficients zero.
Large powers may be symbolic. Exact $c_a,c_s$ and original laws remain
in the model; $\bar c,\bar C$ only choose sufficient bounds.
\end{proposition}

% CMAP ca1-4 | CA1-STATEMENT-ONLY.md:247-259
\begin{theorem}[CA1.4: analytic convergence using finiteness]\label{cm:diagonal}
For $n\ge1$ take
\begin{gather*}
 s_n=2^{-n},\quad T_n=36\sqrt n/s_n,\quad
 k_n=\max(1,\ell_n+n+\lceil\log_2C_n\rceil,\lceil(3n+5)/4\rceil),\\
 a_n=2^{-2k_n},
\end{gather*}
then use $\cmEXT(n,k_n,2^{-n},2^{-n},2^{-n})$ for
$R_{\mathrm{aux},n},\beta_n,N_n$. Both actual neural laws satisfy
\begin{multline}
 \cmBL(\mathcal L_{\mathrm{neural},n}^{\mathfrak b},\mathcal L_{\mathrm{target}})
 \le\min\bigl(2,5\,2^{-n}+2^{362}n^{3/2}2^{-2n}\\
            +4\sqrt2e^{-n}+3K_D2^{-n/2}\bigr)\longrightarrow0 .
 \tag{CA1-DIAG}\label{cm:diagonal-bound}
\end{multline}
The sequence needs no numerical $K_D$; its unknown finite value remains
in the bound.
\end{theorem}

% CMAP ca1-5 | CA1-STATEMENT-ONLY.md:261-282
\begin{theorem}[CA1.5: conditional numerical sufficient selector]\label{cm:selector}
Supply a certified finite rational $\bar K_D\ge K_D$ for H3's actual
normalized target-local density, with boundary and pressure, and rational
$0<\eta\le1$. Define $\cmLtwo$ as in CC1.4, put $e=\eta/2$, and in order take
\begin{gather*}
 j=\cmLtwo(12\bar K_D/e),\quad s=2^{-2j},\quad
 b_{\mathrm{clip}}=\cmLtwo(32/e),\quad T=36b_{\mathrm{clip}}/s,\\
 n=\max(1,\lceil(364+6j+\cmLtwo(b_{\mathrm{clip}}/e))/4\rceil),\\
 k=\max(1,\ell_n+\cmLtwo(8C_n/e),
                 j+\lceil\cmLtwo(216/e)/4\rceil),\quad a=2^{-2k}.
\end{gather*}
Invoke $\cmEXT(n,k,\eta/8,\eta/8,\eta/4)$. The finite tuple
$(s,T,n,a,R_{\mathrm{aux}},\beta,N)$ gives \eqref{cm:budget} at most $\eta$
for both architectures. Its eight allocations are
\[
\begin{array}{c|cccccccc}
\text{charge}&\text{width}&\text{stab.}&\text{aux.}&\text{reg.}&
 \text{variance}&\text{volume}&\text{clip}&\text{target}\\ \hline
\text{allowance}&\eta/4&\eta/8&\eta/8&\eta/16&
 \eta/16&\eta/8&\eta/8&\eta/8 .
\end{array}
\]
Every operation terminates for finite rational inputs. This is a mathematical
sufficient procedure, without a bounded-resource executable claim.
Without $\bar K_D$, CA1.4 remains analytic and this numerical route stays
conditional. No additional density certificate is extracted from an
evolving sibling or qualitative finiteness.
\end{theorem}

% CMAP ca1-5b | CA1-STATEMENT-ONLY.md:284-297
\begin{theorem}[CA1.5b: alternate certified smoothing bypass]\label{cm:bypass}
Instead of $\bar K_D$, supply an explicit certified $r(s)\ge0$, on
positive rational $s\le1$, such that
\[
 \cmBL((X_{\mathrm{obs}},Q_s)_\#\nu_\infty,\mathcal L_{\mathrm{target}})
 \le r(s),\qquad r(s)\to0,
\]
and a specified terminating algorithm $S_r$ returning on every rational
$0<\tau\le1$ a rational $0<s\le1$ with certified $r(s)\le\tau$.
It concerns exactly the original normalized target and identical physical
$c_s/Q_0$ prescription, retaining all its source premises as additional
hypotheses. Qualitative existence or an unverified program is insufficient;
a certified $L^1$ smoothing modulus suffices.

For rational $0<\eta\le1$ take $e=\eta/2$, $s=S_r(\eta/8)$,
$h=\cmLtwo(1/s)$, and
\begin{gather*}
 b_{\mathrm{clip}}=\cmLtwo(32/e),\quad T=36b_{\mathrm{clip}}/s,\\
 n=\max(1,\lceil(364+3h+\cmLtwo(b_{\mathrm{clip}}/e))/4\rceil),\\
 k=\max(1,\ell_n+\cmLtwo(8C_n/e),
              \lceil(2h+\cmLtwo(216/e))/4\rceil),\quad a=2^{-2k}.
\end{gather*}
Then $\cmEXT(n,k,\eta/8,\eta/8,\eta/4)$ gives actual joint distance
$\le\eta$, with the same eight allocations and $E_{\mathrm{target}}$
replaced by $r(s)$. This terminating conditional selector needs no certified
numerical $L^2$ density norm. H3's conditioning and finiteness remain
in its analytic source basis; the supplied smoothing guarantee is exactly
the additional bypass premise.
\end{theorem}

% CMAP ca1-6 | CA1-STATEMENT-ONLY.md:299-299
\begin{remark}[CA1.6: literal scope]\label{cm:limits}
CA1 concerns the fixed smoothed physical-$C_0$-subtracted joint observation
of the original two process-neuron architectures in the original
two-dimensional massive coefficient-one quartic model. Every finite stage
has positive $a,\beta,s$ and finite geometry and width. Whole-plane
process-neuron counts do not specify a finite global random parameter vector.
The raw growing-volume route, volume-uniform raw Wick Cauchy/UI, finite-bit
generation, useful estimator/sampler, total work, positive cosine gap,
matched interacting-covariance Gaussian comparator, comparative advantage,
priority, broader-QFT and publication claims remain separate. Full G1--G9
capability flags remain false. No historical access violation, missing
proof, parser/software defect, unreviewed implication or raw severity is
repaired by this integration.
\end{remark}

% CMAP cc-proof-1 | derivation-composite-certificate-20261031.md:5-27
\subsection{Complete proof of CC1}
\paragraph{Source identity and local normalization.}
S1 fixes the actual common FS maximal state, represented on original $X$.
The strict P1--P6/AP3/AP5 bridge identifies it with the same physical
free-boundary limit. Both original Section 4 maximal states are among those
covered by FS Theorem 7.2, so S3 applies to this identified target. The older
unknown $g_*$ route is unused. The reference is whole-plane $P_1$ restricted
to $\mathcal C_D$, rather than a Dirichlet reference. Gaussian/common-space
fourth-Wick uniqueness identifies source $J_D$ with the anisotropic
own-$c_a$ limit on $D$. The general fixed-square Gaussian/Wick transfer
and general-rectangle construction apply at this local region independently
of the exhaustion boxes $B_n$.

A $\mathcal C_D$-measurable version of $U_D$ follows by compactly contained
interior squares and compactly supported mollifiers. Fourth-chaos pairing
bounds the squared $L^2$ norm of an omitted boundary layer $E_0$ by
$24u^2|E_0|\|C_0\|_4^4\to0$; the accepted covariance bounds make
$\|C_0\|_4$ finite. Interior mollifier limits are local, and Gaussian $L^2$
uniqueness identifies their limit with $U_D$. This locality does not require
a positive-$a$ common-space representative to be a local functional of $\Phi$.

Set $V_D=e^{-U_D}$ and $g=g_{\partial D}$. The exact density is
$k_D=e^{-36\alpha_{\mathrm{pr}}}V_Dg$, $\E k_D=1$, with both minus signs
and no division by $\E V_D$. For every centered finite square $E$,
$\E U_E=0$ and finite negative-exponential moments give
$Z_E=\E e^{-U_E}\ge e^{-\E U_E}=1$. P5/AP3 identifies pressure as the
same-mass normalized log-partition limit, hence
\[
 \alpha_{\mathrm{pr}}=\lim|E|^{-1}\log Z_E\ge0,\qquad
 0<e^{-36\alpha_{\mathrm{pr}}}\le1.
\]
Existence and source identification of pressure remain premises.

% CMAP cc-proof-2 | derivation-composite-certificate-20261031.md:29-86
\paragraph{Numerical Gaussian exponential certificate.}
The general-rectangle Nelson input has semiboundedness coefficient $6u$
and degree parameter two. Since $\pi>3$,
\[
 \kappa_2<1,\quad
 M_6=(5/(4\pi))^{5/6}<1,\quad384A_D=13824<128^2,
 \quad K_{A_D}<128u.
\]
Also $2<e<3$, $\sqrt2<3/2$ imply $\kappa_0<15/2<8$ and
$(2/e)\sqrt{\kappa_0}<3$. Thus the second candidate for $K^\sharp_{A_D}$
is $<384u$, whereas $12uA_D=432u$, proving
$K^\sharp_{A_D}=432u$. As $432/10^{10}<1/1000^2$,
$\beta_{A_D}>1500$.

Take $h_0=1/100>K^\sharp_{A_D}$. For $v\ge h_0$, let
$R(v)=e^{\beta_{A_D}\sqrt v-1}/v$. Then
\[
 (\log R)'(v)=(\beta_{A_D}\sqrt v-2)/(2v)>0,\qquad R(h_0)>100>20.
\]
The tail therefore gives $P_1(U_D<-v)\le e^{-20v}$ for $v\ge h_0$.
Layer cake cut at $e^{10h_0}$ yields
\begin{align*}
 \E e^{-10U_D}
 &\le e^{10h_0}+\int_{h_0}^{\infty}10e^{10v}P_1(U_D<-v)\,dv\\
 &\le e^{1/10}+e^{-1/10}\le10/9+1=19/9<4 .
\end{align*}
Here $e^x\le(1-x)^{-1}$, $0\le x<1$, follows by termwise comparison of
positive series. Since $(4/3)^{10}>4$,
$\|V_D\|_{10}\le4^{1/10}<4/3$. This proves CC1.1 without numerical
quadrature or a global finite-box density envelope. Its rational arithmetic
and the attributed analytic tail have distinct evidentiary roles.

% CMAP cc-proof-3 | derivation-composite-certificate-20261031.md:88-132
\paragraph{The exact boundary factor and lower exponent.}
For $l=6,m_0=1$, $e>2$ gives
\[
 4(1-e^{-6})^{-2}<4(1-1/64)^{-2}=16384/3969<5.
\]
Thus S3 permits exactly $q=5$:
$\nu_\infty(H)\le e^{-36\alpha_{\mathrm{pr}}}\|HV_D\|_5$.
This is the published chessboard conclusion for the actual Section 4 FS
states, with its retained GRS dependency.

For bounded nonnegative local $G$, let
$A_M=\{1/M\le V_D\le M\}$ and $H_M=GV_D^{-1}\mathbf1_{A_M}$.
The test is bounded and local. Substituting the exact density into the source
inequality and cancelling the positive pressure factor gives
\[
 \E[Gg\mathbf1_{A_M}]\le\|G\mathbf1_{A_M}\|_5\le\|G\|_5.
\]
Since $U_D$ is finite a.s., $A_M$ increases to full measure. Monotone
convergence and truncation extend $\E Gg\le\|G\|_5$ to all nonnegative
local $G\in L^5$. Duality gives $\|g\|_{5/4}\le1$. The factor is a local
equivalence class through $k_D=e^{-36\alpha_{\mathrm{pr}}}V_Dg$ and positive
$V_D$; its original boundary-measurable representative in either reading
is unchanged.

Since $1/(10/9)=1/10+1/(5/4)$, H\"older gives
\[
 \|k_D\|_{10/9}\le e^{-36\alpha_{\mathrm{pr}}}\|V_D\|_{10}\|g\|_{5/4}
 \le4/3.
\]
This proves CC1.2, with no inference of an $L^2$ bound from a smaller
exponent. Equivalently the source inequality and H\"older give
\[
 \nu_\infty(H)\le\|H\|_{10}\|V_D\|_{10}\le(4/3)\|H\|_{10}
\]
first for bounded nonnegative local $H$, then by monotone truncation
whenever its Gaussian $L^{10}$ norm is finite. This justifies unbounded
observables below.

% CMAP cc-proof-4 | derivation-composite-certificate-20261031.md:134-201
\paragraph{Gaussian square, smoothing and the one Borel version.}
Write $C(p)=(1+|p|^2)^{-1}$. Parameter integration
$1/(AB)=\int_0^1[tA+(1-t)B]^{-2}\,dt$ and translation give
\[
 \int_{\R^2}C(p)C(k-p)\,dp
 =\pi\int_0^1[1+t(1-t)|k|^2]^{-1}\,dt\le\pi.
\]
With exactly $d\lambda=dp/(2\pi)^2$, Tonelli and Plancherel show
\[
 \|K_f\|_{H\otimes H}^2
 =\iint|\hat f(p+q)|^2C(p)C(q)\,d\lambda(p)d\lambda(q)
 \le F_2^2/(4\pi).
\]
Thus $O_2^P=I_2(K_f)$ exists and $\|O_2^P\|_2\le\sigma$, including
the isometry's factor two. For $s>0$, $Q_s(\Phi)$ has kernel
$(r_s\otimes r_s)K_f$. First use finite spectral cutoffs, Wick pairing and
Bochner integration; dominated tensor convergence and Gaussian Fubini
identify the limit with the measurable spatial integral. Its centering
is the physical mass-one $c_s$ at every $s$.

Evenness, nonnegativity, unit integral and support of $\rho$ give
$|r_s|\le1$ and $0\le1-r_s(p)\le\min(2,s^2|p|^2)$, by
$1-\cos\theta\le\theta^2/2$ and $|z|^2\le2$ on its support. Denote
inverse transforms of $C,r_sC,r_s^2C$ by $C_0,C_{s0},C_{ss}$. Their
$L^2$ norms are bounded by $M_2=1/(2\sqrt\pi)$. The radial Plancherel
split at $|p|=1/s$ gives, with $D_s=\|C_{s0}-C_0\|_2$,
\begin{align*}
 D_s^2&\le s^4\int_{|p|\le1/s}\frac{|p|^4}{(1+|p|^2)^2}\,d\lambda
       +4\int_{|p|>1/s}|p|^{-4}\,d\lambda
       \le5s^2/(4\pi),\\
 \|C_{ss}-C_{s0}\|_2&\le D_s .
\end{align*}
For integrable $G_0$, translated-product Cauchy--Schwarz gives
$\iint|f(x)f(y)G_0(x-y)|\,dxdy\le F_2^2\|G_0\|_1$.
Wick pairing and $\|A^2-B^2\|_1\le\|A-B\|_2(\|A\|_2+\|B\|_2)$ yield
\begin{align*}
 \|Q_s(\Phi)-O_2^P\|_2^2
 &=2\iint f(x)f(y)(C_{ss}^2+C_0^2-2C_{s0}^2)(x-y)\,dxdy\\
 &\le8F_2^2M_2D_s\le(2\sqrt5/\pi)F_2^2s=L_f^2s.
\end{align*}
For distributional rigor, finite momentum cutoffs converge in tensor norm
by the preceding bound, covariances converge in $L^2$ and their squares
in $L^1$, so the translated-product estimate passes both formula and bound
to the limit. No spatial positivity of $r_s$ or covariance difference is
assumed. Second-chaos hypercontractivity gives for every real $R\ge2$
\[
 \|Q_s(\Phi)-O_2^P\|_R\le(R-1)L_f\sqrt s,\qquad
 \|O_2^P\|_R\le(R-1)\sigma.
\]

For positive $s$, distributional smoothing makes $z\mapsto v(h_{s,z})$
continuous. Its compact integral is a limit of Borel Riemann sums, all
using tests inside $[-2,2]^2$, so $Q_s$ is Borel and $\mathcal C_D$ local.
At $s_j=2^{-2j}$, the sum of squared Gaussian errors is finite.
Tonelli gives almost-sure convergence to $O_2^P$. The convergence event
and limit are local, hence $Q_0(\Phi)=O_2^P$ a.s. Target-local absolute
continuity transfers this same full-probability event. Equality of source
and $X$ representations on compact evaluations transfers the same countable
construction, giving target $O_2(f)=Q_0$ without another subsequence or
observable. Finally $F_2\le2,\pi>3,\sqrt5<3$ imply $\sigma<1,L_f<3$.

% CMAP cc-proof-5 | derivation-composite-certificate-20261031.md:203-228
\paragraph{All orders and numerical target smoothing.}
For real $r\ge1$, put $E_s=Q_s-Q_0$. Its $r$th absolute power is local.
Hypercontractivity at $R=10r$ and the preceding local transfer inequality
give
\[
 \E_{\nu_\infty}|E_s|^r
 \le(4/3)\|E_s\|_{10r,P_1}^r
 \le(4/3)[(10r-1)L_f\sqrt s]^r,
\]
and likewise
$\E_{\nu_\infty}|Q_0|^r\le(4/3)[(10r-1)\sigma]^r$.
Taking roots proves CC1.3; $r=1$ yields $12L_f\sqrt s<36\sqrt s$.
No unknown $K_D$, boundary norm or finite-box $M_n$ occurs.

For $R\ge1$, Minkowski gives
\[
 \sup_{0<s\le1}\|Q_s\|_{L^R(\nu_\infty)}
 \le(4/3)^{1/R}(10R-1)(\sigma+L_f).
\]
For any positive moment order $a_0$, choose $R>\max(a_0,1)$.
Markov/H\"older bounds the $a_0$-moment tail uniformly by a constant times
the threshold to power $-(R-a_0)$. This proves uniform integrability
of $\{|Q_s|^{a_0}:0<s\le1\}$, at the target-local scope.

% CMAP cc-proof-6 | derivation-composite-certificate-20261031.md:230-238
\paragraph{Terminating smoothing selector.}
For positive rational $x$, $\cmLtwo(x)$ is obtained by finite integer
numerator/denominator comparisons with powers of two. At $x=36/\tau$,
$2^j\ge x$, including $j=0$ if $x\le1$. Thus $s=2^{-2j}$ satisfies
$0<s\le1$ and $36\sqrt s=36\,2^{-j}\le\tau$.
Replacing $\tau$ by $\epsilon/4$ gives the stated $144/\epsilon$ formula.
This is exact rational selection with symbolic dyadic output, not
interacting expectation evaluation. Historical finite selector examples
are finite checks separate from this proof. CC1.4 follows.

% CMAP cc-proof-7 | derivation-composite-certificate-20261031.md:240-272
\paragraph{Retained unsuccessful routes and limits.}
The original direct $L^2$ H\"older route gives
\[
 \E_{\nu_\infty}|Q_s-Q_0|\le K_DL_f\sqrt s,\qquad
 K_D=e^{-36\alpha_{\mathrm{pr}}}\|V_Dg\|_2
 \le e^{-36\alpha_{\mathrm{pr}}}\|V_D\|_4\|g\|_4.
\]
Old qualitative finiteness of $\|g\|_4$ gave no numerical certificate.
Normalization alone only fixes $\E V_Dg=e^{36\alpha_{\mathrm{pr}}}$ and
does not bound $\|g\|_4$. Setting $g=1$ or dividing by $\E V_D$ would change
the target. The smaller-exponent proof retains both factors.

Finite-box marginalization gives
$\|\E[r_{0,n}\mid\mathcal C_D]\|_2\le\|r_{0,n}\|_2\le M_n$,
which grows with $n$. The unused raw triangle has sufficient charges
$M_nL_f\sqrt s$ and $2^{362}n^{3/2}s^{-3}2^{-5n}$.
Since $\ell_n\ge2n^4/125$, forcing the former small makes that crude
second majorant grow. This is failure of sufficient majorants, with no
nonconvergence or all-method impossibility implication.

The linear source bound only gives
\[
 \E_{\nu_\infty}|Q_s|\le3I\|\rho_s\|_2^2+Ic_s
 \le4IR_2^2/s^2\le432/s^2,\quad
 \E_{\nu_\infty}|Q_s-Q_t|\le432(s^{-2}+t^{-2}).
\]
These divergent upper bounds, and the explicitly scale-dependent $H^{-2}$
local moment constant, supply no small-scale Wick-square Cauchy modulus.
An unused conditional-boundary proposal would require a correctly normalized
$\E_{\mathrm{cond},\xi}|Q_s-Q_0|\le C(\xi)\sqrt s$ with explicit finite
$\E_{\nu_\infty}C(\xi)$, shifted-Dirichlet Wick conversion, conditional
normalizer and quantitative target-boundary trace control. The linear
$L^2$ source interface supplies none of this boundary control.
Finite-tuple AP6 reconstruction likewise cannot certify the missing
infinite-dimensional density-norm tail by observed stabilization.

The FS source matching retains the GRS checkerboard dependency and the
original classical/P1--P6 scope. Historical source investigators read the
primary FS line/page, without claiming a fresh secondary-proof audit.
Frozen original author candidate headers and operational/access findings
remain historical; their subsequent accepted decisions govern CC1.
There is no new outside-expert, statistical-independence or OS-isolation
claim. The stronger $L^2$ certificate, raw fixed-$a$ limit, assembly from
CC1 alone, useful sampling/work/comparison/gap and priority remain outside
CC1.5. This completes the full positive proof and its retained limits.

% CMAP ca-proof-0 | derivation-composite-assembly-20261031.md:7-30
\subsection{Complete CA1 proof}
\paragraph{P0: target, local prescription and normalization ledger.}
H1 fixes the original pressure-identified target, rather than a new state or
subsequence. Same-mass whole-plane reference, free-boundary interaction and
physical/classical identification are literal source conditions.
$8000<10^{10}$ proves the strict cap, without any old $g_*$ inference.

All $h_{s,z}$ are supported in $[-2,2]^2$, strictly inside $D_{\mathrm{loc}}$,
every $B_n$ and $D_{\mathrm{aux}}$. The accepted CV1 continuity estimate is
\[
 |Q_s(v)-Q_s(w)|\le IA_s^2p_3(v-w)[p_3(v)+p_3(w)],\quad
 A_s^2=s^{-2}R_2^2+2s^{-4}D_\rho^2+s^{-6}\|\Delta\rho\|_2^2 .
\]
Alternatively countable Riemann sums show local Borel measurability.
Zero extension of an $E$ class gives exactly the whole-plane convolution on
$\supp f$, so the linear coordinate and the actual Borel pair agree.

Each probability has its own normalization before observation:
\[
\begin{array}{c|c|c}
 \text{law}&\text{reference and weight}&\text{normalizer}\\ \hline
 P_N^{\mathfrak b}&\rho_N^{\mathfrak b},\ e^{-W-\beta J_{\mathrm{out}}
                                  +(K/2)\int_{D_{\mathrm{aux}}}v^2}
                                  &Z_N^{\mathfrak b}\\
 Q_\beta&\rho_\gamma,\ \text{same full weight}&Z_\gamma\\
 Q_{\mathrm{res}}^\beta&G',\ e^{-W-\beta J_{\mathrm{out}}}&z_\beta\\
 Q_{\mathrm{res}}^0&G',\ e^{-W}&z_0\\
 \nu_{a,n}&G_a,\ e^{-W}&Z_{a,n}\\
 \nu_n&\Phi_\#(e^{-U_{0,n}}P/Z_{0,n})&Z_{0,n}\\
 \nu_\infty&P_1|_{\mathcal C_D},\ k_D&\E_{P_1}k_D=1
\end{array}
\]
The pure Gaussian mass normalizer $M_{D_{\mathrm{aux}},a}$ is separate;
HP gives $Z_\gamma=M_{D_{\mathrm{aux}},a}z_\beta$ and the all-Borel
restriction identity. This does not identify $Q_\beta$ with $\nu_{a,n}$
at finite geometry and $\beta$. Endpoint covariance $C_{\mathrm{end}}$,
full tilt $Q_\beta$ and physical reference $C_a$ remain distinct.
No extra local $Z_D$ is inserted into $k_D$.

% CMAP ca-proof-1 | derivation-composite-assembly-20261031.md:32-69
\paragraph{P1: fixed-stage width and its normalization.}
Use the literal NC1 theorem H7. Nonnegative unit-mass convolution after
zero extension is an $L^4$ contraction. H\"older gives
$\|v\mapsto\int fv\|\le\|f\|_{4/3}<1$ and
$\|B_T\|\le F_2<1/2$ for $T_s$ and $T_{\mathrm{id}}$. For
$x=\|v\|_4$, $J_{T,k}$ therefore has first derivative at most $q_1=1+2x$,
second at most $q_2=2$, higher derivatives zero. The finite subtraction
changes none. For $H(v)=G(J_{T,k}(v))$ with $G\in\mathcal U_4$, the scalar
chain rule gives $p_1,\ldots,p_4$; fourth-order multiplicities $1,6,3$
correspond to four first derivatives, one second and two first, and two
second derivatives.

Differentiation and H\"older give $v_1,\ldots,v_4$ for $V$; in particular
$\|v\|_2^2\le\sqrt d\,x^2$, with quadratic coefficient at most $B$.
The exponential chain rule bounds $D^4w$ by $we_4$, while Leibniz
multiplicities $1,4,6,4,1$ give $D^4(Hw)$ at most $wP_J$. Moreover
\[
 V(v)\ge\alpha x^4-B\sqrt d\,x^2\ge(\alpha/2)x^4-L,\qquad
 \max_{x\ge0}x^je^{-\alpha x^4/2}=[j/(2\alpha e)]^{j/4}\quad(j>0).
\]
The dropped constant $3uc_a^2A_n$ is positive; the $j=0$ maximum is one.
Summing nonnegative coefficients gives $M_Z,M_J$ and bounded lower
derivatives and functions, uniformly in positive $s$ and finite common $k$.

Actual iid neuron signs of magnitude $1/2$ give
$\E F_N(z)^4=(3-2/N)/16\le3/16$, while $\E\gamma(z)^4=3/16$.
Tonelli, dropping nonpositive quadratic terms, gives
$\E V\le3Ad/16+3uc_a^2A_n$ under either reference. Jensen proves the two
own normalizers at least $z$, while quartic damping makes them finite
positive. Their equality is never assumed.

For an $E$-$C^4$ function $\Psi$ with $\sup\|D^4\Psi\|\le M$, successively
replace the $N$ independent normalized neurons by independent normalized
endpoint copies. Integral Taylor remainder is $M\|h\|_4^4/24$.
Zero Bochner mean cancels linear terms, the EVERY-bounded-bilinear identity
at the independent base cancels all quadratic terms, and global sign
symmetry cancels cubics; lower terms are integrable. Each replacement
costs at most $M(d/16+3d/16)/(24N^2)=Md/(96N^2)$.
The telescope ends at the exact Gaussian-sum law $\rho_\gamma$.
Numerator and denominator errors are $dM_J/(96N)$ and $dM_Z/(96N)$.
For $A_N=\rho_N^{\mathfrak b}(Hw)$ and $A_\gamma=\rho_\gamma(Hw)$,
\[
 \frac{A_N}{Z_N^{\mathfrak b}}-\frac{A_\gamma}{Z_\gamma}
 =\frac{A_N-A_\gamma}{Z_N^{\mathfrak b}}
 +\frac{A_\gamma}{Z_\gamma}\frac{Z_\gamma-Z_N^{\mathfrak b}}
                                 {Z_N^{\mathfrak b}} .
\]
Since $|A_\gamma/Z_\gamma|\le1$ and $Z_N^{\mathfrak b}\ge z$, the bound
is $C/N$, with one $z^{-1}$ and both errors.

For joint $G\in\mathrm{BL}_1(\R^2)$, Gaussian convolution at scale
$0<\delta\le1$ approximates it uniformly within $\sqrt2\delta$.
Its first derivative is at most one and its $j$th derivative operator
norm for $2\le j\le4$ is at most $\sqrt{(j-1)!}\delta^{1-j}$:
directional integration by parts and the degree-$(j-1)$ Wick polynomial
$L^2$ bound establish this for arbitrary unit directions. Thus
\[
 \text{width error}\le\sqrt6 C/(N\delta^3)+2\sqrt2\delta .
\]
Optimization gives $D_JN^{-1/4}$, capped separately by two.
Taking $\delta=\eta_W/(4\sqrt2)$ proves the sufficient
$512\sqrt3 C\eta_W^{-4}$ width. These smoothing costs are internal to
$E_{\mathrm{width}}$, distinct from physical $s$ and counted once.
The fixed cosine has derivative norms $\le2^{-j/2}$ and lies in
$\mathcal U_4$ and $\mathrm{BL}_1$, hence its $C/N$ alternative.
This recalls the first actual-law arrow with all NC1/RL2/NR1 qualifications.

% CMAP ca-proof-2 | derivation-composite-assembly-20261031.md:71-94
\paragraph{P2: stabilizer between residual normalized probabilities.}
Put $c=c_a$. Since $x^4-6cx^2+3c^2=(x^2-3c)^2-6c^2$,
$W\ge-6uA_nc^2$ and write $M=e^{6uA_nc^2}$.
Under $G'$, point variance $r(z)\in[0,c]$ gives
\[
 \E_{G'}W=3u\int_\Lambda(r(z)-c)^2\le3uA_nc^2,\quad
 \E_{G'}J_{\mathrm{out}}\le3|D_{\mathrm{aux}}\setminus\Lambda|c^2.
\]
Jensen gives $z_0\ge z_{0,\mathrm{low}}=e^{-3uA_nc^2}$, distinct from P1's
$z$. Let $h_\beta=e^{-W-\beta J_{\mathrm{out}}}$ and $h_0=e^{-W}$.
The nonnegative $\Delta_h=\E_{G'}(h_0-h_\beta)$ satisfies
$\Delta_h\le3\beta M|D_{\mathrm{aux}}\setminus\Lambda|c^2$ by
$1-e^{-t}\le t$. The exact subtraction
\[
 h_\beta/z_\beta-h_0/z_0
 =(h_\beta-h_0)/z_0+h_\beta(1/z_\beta-1/z_0)
\]
has $L^1$ norm at most $2\Delta_h/z_0$. Thus every bounded Borel whole-path
test changes by at most $E_{\mathrm{stab}}$. HP's all-Borel restriction
identity applies on the $Q_\beta$ side and pushforward cannot increase
the bound. This is the single stabilizer charge.

% CMAP ca-proof-3 | derivation-composite-assembly-20261031.md:96-133
\paragraph{P3: auxiliary coupling of action and observation.}
Use independent HP paths $X_a=Y+\psi$, $X_a\sim G_a,Y\sim G'$.
Their $e^{-W}$ tilts are exactly $\nu_{a,n}$ and $Q_{\mathrm{res}}^0$,
with normalizers $Z_{a,n}$ and $z_0$. For $p_c(x)=x^4-6cx^2+3c^2$,
factoring and $(|x|+|y|)(x^2+y^2)\le2(|x|^3+|y|^3)$ give
\[
 |p_c(x)-p_c(y)|
 \le|x-y|\{2(|x|^3+|y|^3)+6c(|x|+|y|)\}.
\]
Centered Gaussian marginals of variance at most $c$ have sixth moment
at most $15c^3$ and second at most $c$. The bracket has $L^2$ norm
$\le L_c=(4\sqrt{15}+12)c^{3/2}$. Cauchy--Schwarz, integration and the
defect diagonal yield
\[
 \E|W(X_a)-W(Y)|\le uA_nL_c\sqrt{\delta_D(\overline\Lambda)},\quad
 \E|e^{-W(X_a)}-e^{-W(Y)}|
 \le M uA_nL_c\sqrt{\delta_D(\overline\Lambda)}.
\]
The latter uses the $M$-Lipschitz bound for $e^{-t}$ at $t\ge-6uA_nc^2$.

Positivity of $f$, PSD covariance and Minkowski give
$\E|X_a(f)-Y(f)|\le I\sqrt{\delta_D(K_f)}$.
For $z\in\supp f$, let $A_z=T_s(X_a-Y)(z)$ and $B_z=T_s(X_a+Y)(z)$.
Unit nonnegative convolution gives
$\|A_z\|_2\le\sqrt{\delta_D(K_s)}$,
$\|B_z\|_2\le2\sqrt c$, where $K_s=[-1-s,1+s]^2\subset K_*$.
Cauchy--Schwarz, without independence of these variables, yields
\[
 \E|B_{T_s}(X_a,X_a)-B_{T_s}(Y,Y)|
 \le2I\sqrt c\sqrt{\delta_D(K_*)}.
\]
The common subtraction cancels, so the expected Euclidean pair difference
is at most the sum of its two coordinate bounds.

For $G\in\mathrm{BL}_1$, write $H_X=G(J_s(X_a))$, $H_Y=G(J_s(Y))$,
$w_X=e^{-W(X_a)},w_Y=e^{-W(Y)}$, $z_X=\E w_X,z_Y=\E w_Y=z_0$.
Direct normalization subtraction gives
\[
 \left|\frac{\E w_YH_Y}{z_Y}-\frac{\E w_XH_X}{z_X}\right|
 \le\frac{2\E|w_X-w_Y|+\E(w_X|H_Y-H_X|)}{z_Y}
 \le\frac{2\E|w_X-w_Y|+M\E\|J_s(Y)-J_s(X_a)\|}{z_{0,\mathrm{low}}}.
\]
Since $M/z_{0,\mathrm{low}}=e^{9uA_nc^2}$, this is exactly
$E_{\mathrm{aux}}$ with probability cap two. Densities, own normalizers
and nonlinear path observations change together.

% CMAP ca-proof-4 | derivation-composite-assembly-20261031.md:135-168
\paragraph{P4: regulator, physical variance and translations.}
HC uses the common-space comparison
\[
 |\E_{\mu_{a,n}}G(A)-\E_{\mu_{0,n}}G(B)|
 \le\|r_{a,n}-r_{0,n}\|_1+M_n\|A-B\|_{L^2(P;\R^2)}
\]
for $A=(\Phi_a(f),Y_{a,s})$, $B=(\Phi(f),Y_s)$.
The density term is $4K_nM_n\sqrt a$, the Gaussian linear difference
at most $aF_2/2$, and the own-smoothed second-chaos difference at most
$H_f\sqrt a$, uniformly in positive $s$. These are literal accepted
CMP/HC bounds for actual anisotropic $\tau_a$. They give $R_n(a)$,
without a finite-box raw smoothing term. Since $a<1,F_2/2\le1,H_f<2$,
$K_n<n/10^6$, $R_n(a)\le M_nC_n\sqrt a$.

With $D_0=1+w^2+q^2$, $T_0=w^2q^2\le D_0^2/4$, algebra gives
\[
 0\le\hat C_0-\hat C_a=\frac{a^2T_0}{D_0(D_0+a^2T_0)}\le a^2/4,
\]
so Plancherel and $\|\rho\|_2=F_2/I$ yield
\[
 0\le d_{a,s}=I\int|\hat\rho(sp)|^2(\hat C_0-\hat C_a)\,d\lambda
 \le F_2^2a^2/(4Is^2)\le27(a/s)^2 .
\]
At every positive-$a$ physical path,
$Q_s(\Phi_a)=Y_{a,s}-d_{a,s}$. Translating one normalized law costs
at most $d_{a,s}$, without an $M_n$ factor. The regulator arrow is
$R_n(a)+d_{a,s}$.

Changing common subtraction $k_1$ to $k_2$ gives
$J_{T,k_2}=J_{T,k_1}-(0,I(k_2-k_1))$. Translation preserves both test
classes, hence common translation preserves distances exactly; one-law
translation costs at most $I|k_2-k_1|$. Using common $c_s$ throughout P1--P3
adds no endpoint-natural-centering charge. The four quantities
$c_{\mathrm{end},s},c_{a,s},c_s,c_a$ differ. Untilting gives
$\E B_{T_s}=Ic_{\mathrm{end},s}$ and $\E B_{\mathrm{id}}=I/4$, with
no interacting-centering consequence. Endpoint-own to physical translation
is $+(0,I(c_{\mathrm{end},s}-c_s))$.
Finally exact subtraction gives the raw-to-smoothed identity in CA1.2,
whose term $I(c_s-c_a)$ has no asserted sign and cannot be replaced by
$d_{a,s}$. The action always uses exact $:v^4:_{c_a}$, including when
upper envelopes later choose stages.

% CMAP ca-proof-5 | derivation-composite-assembly-20261031.md:170-225
\paragraph{P5: direct nonlinear volume, clipping and target smoothing.}
Fix $s>0$. H3 implies $\E_\zeta v(h)^2\le3\|h\|_2^2$ for
$\zeta=\nu_n,\nu_\infty$. For any scalar source $Y$ with the corresponding
subgaussian proxy $\sigma_h^2$, an independent standard Gaussian and
Tonelli give
$\E e^{\lambda Y^2}\le(1-2\lambda\sigma_h^2)^{-1/2}$ for
$0\le\lambda<(2\sigma_h^2)^{-1}$; this does not make $Y$ Gaussian.
Here all $h_{s,z}$ have proxy $\sigma_s^2=3R_2^2/s^2$.

Take a finite rectangular partition $(A_i)$ of $[-1,1]^2$ with $z_i\in A_i$,
exact weights $w_i=\int_{A_i}f$, discarding zero weights. Put
\[
 Z_i(v)=\sqrt{w_i}v(h_{s,z_i}),\quad S_{\mathcal P}=\sum_iZ_i^2,\quad
 Q_{\mathcal P}=S_{\mathcal P}-Ic_s,\quad
 S_s=\int f(z)v(h_{s,z})^2\,dz .
\]
Weights sum exactly to $I$. Jensen with weights $w_i/I$ or $f(z)\,dz/I$
gives $\E_\zeta e^{S_{\mathcal P}/(4B_s)}\le\sqrt2$ and
$\E_\zeta e^{S_s/(4B_s)}\le\sqrt2$, $B_s=I\sigma_s^2$.
Both tails above $T^2$ are bounded by $\sqrt2e^{-T^2/(4B_s)}$.
No field-evaluation independence is used.

If the largest cell diameter is $\delta_{\mathcal P}$, translated
mollifier $L^2$ control and source moments give
\begin{gather*}
 \|h_{s,z}-h_{s,z_i}\|_2\le D_\rho|z-z_i|s^{-2},\\
 \E_\zeta|v(h_{s,z})^2-v(h_{s,z_i})^2|
          \le6R_2D_\rho|z-z_i|s^{-3},\\
 \E_\zeta|Q_s-Q_{\mathcal P}|\le6IR_2D_\rho\delta_{\mathcal P}s^{-3}.
\end{gather*}
Exact total weight cancels subtraction; quadrature tends to zero at each
fixed $n,s,T$, uniformly for these two endpoint laws.

Apply H3 to $g=(f,\sqrt{w_i}h_{s,z_i})$, $R=2,n_0=1$.
The square-sum coefficient is exactly
$e^{-4}[1+1/(Is^4)]$, independently of the number of cells. The same
$P_*$ gives $C_g\le C_s^*$ and AP6 tuple distance $\le E_n^*(s)$.
For joint $G\in\mathrm{BL}_1$ define
$H_T(x,z)=G(x,\min(\|z\|_2^2,T^2)-Ic_s)$.
Its value is at most one and Lipschitz constant at most
$\Lambda_T=\max(1,2T)$, so $H_T/\Lambda_T$ is an admissible AP6 test,
including singular/repeated tuples. Clipping costs at most twice the tail
at each endpoint. Thus
\[
 |\nu_nG(X_{\mathrm{obs}},Q_{\mathcal P})
       -\nu_\infty G(X_{\mathrm{obs}},Q_{\mathcal P})|
 \le\Lambda_T E_n^*(s)+4\sqrt2e^{-T^2/(4B_s)} .
\]
Removing quadrature costs at most
$12IR_2D_\rho\delta_{\mathcal P}s^{-3}\to0$. The right side is independent
of partition and $G$, so the supremum proves the actual nonlinear
$E_{\mathrm{vol}}+E_{\mathrm{clip}}$ arrow. This is a mathematical
observation limit, rather than a numerical work procedure.

For the last arrow the normalized local density and H4 give
\[
 \E_{\nu_\infty}|Q_s-Q_0|
 =\E_{P_1}[k_D|Y_s-Q_0(\Phi)|]
 \le K_D\|Y_s-Q_0(\Phi)\|_2\le K_DL_f\sqrt s.
\]
Only local variables enter. Summability of squared errors at $2^{-2j}$
gives the one Gaussian full-probability convergence set; locality transfers
it by target absolute continuity, and each $\nu_n$ uses the same prescription
by its common-space density. The endpoint is the intended physical Borel
pair, rather than a new pair-law-only variable. A BL test changes by
at most this $L^1$ coordinate difference. Alternatively a certified bound
$r(s)$ for this identical joint-law arrow costs $r(s)$ once, with no
extra density factor.

Taking the triangle over the seven normalized laws proves CA1.1 with
exactly eight charges and cap two. Width smoothing, variance translation,
two-endpoint clipping and target smoothing are counted once. The cosine
belongs to the class and inherits the budget and its sharper width option.

% CMAP ca-proof-6 | derivation-composite-assembly-20261031.md:227-286
\paragraph{P6: actual anisotropic tail and rational EXT.}
Let $A_q=1+q^2$, $B_q=1+a^2q^2$. Then
$\hat C_a(w,q)=(A_q+B_qw^2)^{-1}$ and
$\partial_w\hat C_a=-2B_qw/(A_q+B_qw^2)^2$.
Substitute $w=\sqrt{A_q/B_q}\,t$, then $t=\tan\theta$, to obtain
\[
 \int_\R\frac{w^2\,dw}{(A_q+B_qw^2)^4}
 =\frac{\pi}{16A_q^{5/2}B_q^{3/2}}.
\]
Including exactly the Fourier measure,
\[
 \int|\partial_w\hat C_a|^2\,d\lambda
 =\frac1{16\pi}\int_\R\frac{\sqrt{1+a^2q^2}}{(1+q^2)^{5/2}}\,dq .
\]
The $q$ derivative has the same integral. For $0<a<1$,
\[
 \int|\nabla\hat C_a|^2\,d\lambda
 \le\frac1{8\pi}\int_\R(1+q^2)^{-2}\,dq=1/16.
\]
These integrals establish derivative $L^2$ integrability, while
$\hat C_a\le\hat C_0$ establishes $\hat C_a\in L^2$.
Distributional Fourier differentiation and Plancherel give
$\int|x|^2C_a(x)^2\,dx\le1/16$. Its $L^2$ representative is the
accepted continuous positive-$a$ covariance. The deterministic integral
tail is $\int_{|x|>r}C_a^2\le1/(16r^2)$, proving \eqref{cm:tail}
by HP, without spatial positivity or altered regulator.

For $a=2^{-2k}$, $k\ge1$, the accepted exact variance bound gives
\[
 c_a\le[\operatorname{arsinh}(1/a)+1]/(2\pi)
 \le[\log(3/a)+1]/(2\pi)<2k+3=\bar c,
\]
using $\log3<2,\log2<1,2\pi>1$. Thus $\bar c\ge1$ and
$\bar E\ge e^{9uA_nc_a^2}$ since $e^x\le4^{\lceil x\rceil}$ for $x\ge0$.
Also $L_c\le28\bar c^2$ and, using the weaker $I\le4$,
$2uA_nL_c+I+2I\sqrt{c_a}\le56uA_n\bar c^2+4+8\bar c=\bar H$.
The selected $R_{\mathrm{aux}}=1000n+r+1$ contains the closed Euclidean
$r$ neighborhoods of $\overline\Lambda,K_f,K_*$ with surplus margin.
The tail and $\sqrt K\le1/a$ imply
$\sqrt{\delta_D(H)}\le1/(4ar)$ at each compact, hence
$E_{\mathrm{aux}}\le\bar E\bar H/(4ar)\le\eta_A$.
All these choices precede $\beta$.
Because $|D_{\mathrm{aux}}\setminus\Lambda|\le d$, the chosen positive
$\beta$ gives $E_{\mathrm{stab}}\le6\beta d\bar c^2\bar E\le\eta_S$.
No $\beta$ limit is interchanged with expanding geometry.

For width, $\beta\le u$ gives $\alpha=\beta,A=u$. Increasing $c_a$ to
$\bar c$ increases $B,L$ and all nonnegative polynomial coefficients;
$\sqrt d=2R_{\mathrm{aux}}$ is exact and $z^{-1}\le e^{\overline Z_{\exp}}$.
For $j>0$,
\[
 [j/(2\alpha e)]^{j/4}\le\max(1,j/(2\alpha))^j=t_j:
\]
if $j/(2\alpha)\le1$ the left side is at most one; otherwise use
$e\ge1$ and $j/4\le j$. The maximal polynomial degree is twelve.
Coefficientwise comparison and
$e^{\bar L+\overline Z_{\exp}}\le4^{\lceil\bar L\rceil+\lceil\overline Z_{\exp}\rceil}$
prove $C\le\bar C$. Thus
$N\ge1024\bar C\eta_W^{-4}\ge512\sqrt3 C\eta_W^{-4}$.
Moreover $D_J^4=(8192\sqrt3/27)C<1024C$ directly bounds the displayed
$E_{\mathrm{width}}$ by $\eta_W$. Every coefficient, ceiling and power is
finite rational/integer. The action still uses exact $c_a$ and observation
exact $c_s$; envelopes only choose resources. This proves all CA1.3
claims without a tail or stage-normalizer oracle.

% CMAP ca-proof-7 | derivation-composite-assembly-20261031.md:288-311
\paragraph{P7: analytic diagonal.}
The chosen $k_n$ gives
\[
 2^{\ell_n}C_n2^{-k_n}\le2^{-n},\qquad
 27\,2^{-4k_n+2n}\le27\,2^{-n-5}<2^{-n},
\]
so regulator and variance each cost at most $2^{-n}$ and $0<a_n\le1/4$.
The first six nonnegative Taylor terms of $e^{125/32}$ exceed 32 by
exact rational arithmetic. Thus $q<1/32$, and
\[
 \Theta_n\le4n2^{-5n},\quad C_s^*<2^{364}s^{-2},\quad
 E_n^*(s)\le2^{355}ns^{-2}2^{-5n}.
\]
The last uses $e^{1/128}\le128/127$ and
$4(128/127)/2500<1/512$.
At $T_n=36\sqrt n/s_n$, $\Lambda_T=72\sqrt n/s_n<2^7\sqrt n/s_n$,
and $T_n^2/(4B_{s_n})\ge n$. Consequently
\[
 E_{\mathrm{vol}}\le2^{362}n^{3/2}2^{-2n},\quad
 E_{\mathrm{clip}}\le4\sqrt2e^{-n},\quad
 E_{\mathrm{target}}\le3K_D2^{-n/2}.
\]
EXT bounds each of width, stabilizer and auxiliary by $2^{-n}$.
Adding the two regulator/variance terms gives $5\,2^{-n}$ and proves
\eqref{cm:diagonal-bound}. Every term tends to zero for fixed finite $K_D$.
Each stage is finite positive and works for both architectures.
No value of $K_D$ defines the sequence, although finiteness alone yields no
accuracy-dependent stopping index from this bound.

% CMAP ca-proof-8 | derivation-composite-assembly-20261031.md:313-343
\paragraph{P8: numerical selector under the density certificate.}
A normalized density has $K_D\ge\|k_D\|_1=1$, so $\bar K_D\ge1$.
For rational $x>0$, $x\le2^{\cmLtwo(x)}$ by exact finite comparisons.
With $e=\eta/2$ the selector has $j\ge4$, $b_{\mathrm{clip}}\ge5$,
$0<s\le1$, $T>1/2$ and $\Lambda_T=72b_{\mathrm{clip}}/s$.
The order is $s,T,n,a$, auxiliary box, stabilizer, width.

First $L_f<3$ and the choice of $j$ give
$E_{\mathrm{target}}\le3\bar K_D2^{-j}\le e/4=\eta/8$.
Next $B_s\le324/s^2$ gives $T^2/(4B_s)\ge b_{\mathrm{clip}}^2$.
Using $\sqrt2<2,\exp(1)>2$,
$E_{\mathrm{clip}}\le8\,2^{-b_{\mathrm{clip}}^2}
\le8\,2^{-b_{\mathrm{clip}}}\le e/4$.
For volume use $n\le2^n$:
\[
 E_{\mathrm{vol}}\le2^{362}b_{\mathrm{clip}}s^{-3}n2^{-5n}
 \le2^{362}b_{\mathrm{clip}}2^{6j}2^{-4n}\le e/4
\]
by $4n\ge364+6j+\cmLtwo(b_{\mathrm{clip}}/e)$.
Open-box support and $n\ell\ge2$ already hold at $n=1$.
Writing $L_A=\cmLtwo(8C_n/e)$ and $L_S=\cmLtwo(216/e)$, the
selected $k$ satisfies $k\ge\ell_n+L_A$, $4(k-j)\ge L_S$, hence
\[
 E_{\mathrm{reg}}\le2^{\ell_n}C_n2^{-k}\le e/8=\eta/16,\qquad
 E_{\mathrm{var}}\le27\,2^{-4(k-j)}\le e/8.
\]
The variance charge is outside $M_n$. These five physical-to-target
allowances sum to $e$. EXT at the already fixed $n,k$ gives
$\eta/8,\eta/8,\eta/4$ for auxiliary, stabilizer and width, totaling
the other $\eta/2$. All eight allocations sum exactly to $\eta$ and
the cosine inherits the same guarantee.

Termination consists of finite power-of-two comparisons, rational
$h_2,m_2,\ell_n,C_n$, dyadics $s,a$, finite integer $\bar E,r$, positive
rational $\beta$, finitely many polynomial coefficients of degree at most
twelve, the finite integer exponent $\lceil\bar L\rceil+
\lceil\overline Z_{\exp}\rceil$, rational $\bar C$ and rounded integer $N$.
Unlimited exact arithmetic terminates; a finite symbolic expression
also describes the enormous chosen integer. No finite resource or
performance theorem is asserted. The density certificate is an extra
mathematical input, not verified by this algorithm. Without it CA1.4
and conditional CA1.5 remain; $M_n$, a Dirichlet density, bare boundary
factor or pressure-free expression cannot replace $K_D$.

% CMAP ca-proof-8b | derivation-composite-assembly-20261031.md:345-356
\paragraph{P8b: alternate certified target modulus.}
Under exactly CA1.5b's contract, $S_r(\eta/8)$ terminates with rational
$s\in(0,1]$ and certified last-arrow cost at most $\eta/8$.
Let $h=\cmLtwo(1/s)$, so $s^{-1}\le2^h$. The P7 majorant and choice of
$n$ give
\[
 E_{\mathrm{vol}}\le2^{362}b_{\mathrm{clip}}s^{-3}n2^{-5n}
 \le2^{362}b_{\mathrm{clip}}2^{3h}2^{-4n}\le e/4 .
\]
The same clipping choice gives $E_{\mathrm{clip}}\le e/4$, and
$k\ge\ell_n+\cmLtwo(8C_n/e)$ gives $E_{\mathrm{reg}}\le e/8$.
The condition $4k\ge2h+\cmLtwo(216/e)$ yields
$E_{\mathrm{var}}\le27\,2^{-4k}s^{-2}\le27\,2^{-4k+2h}\le e/8$.
The four non-target physical charges have exactly P8's allowances,
and $r(s)$ occupies the single target allowance. EXT gives the same
other three. The sum is $\eta$. Composition of the terminating $S_r$
with finite rational operations terminates; no monotonicity or rounding
of $r,s$ is used. Dyadic $s$ is unnecessary on this route. This proves
CA1.5b with neither a numerical $L^2$ certificate nor an unproved
integration step.

% CMAP ca-proof-9 | derivation-composite-assembly-20261031.md:358-370
\paragraph{P9: computability, order and retained adverse paths.}
The selector's explicit envelopes are rational/integer once its extra
certificate and accuracy are supplied: $2^{360}$, box/support constants,
$M_n,C_n,\bar c,\bar E,\bar H,\beta$, coefficients, $\bar C,N$.
The accepted OPT $2^{360}$ version remains a source premise, not a new
OPT proof. HP/HC are source-qualified Gaussian/physical theorems rather
than covariance-matching inferences.

Exact positive-stage $c_a,c_s,c_{a,s}$ are specified computable reals.
Frozen accepted CMP C8 procedures bound the $c_a$ integral tail beyond
$q=R$ by $1/(2\pi aR)$, and the $c_s,c_{a,s}$ tails outside $[-R,R]^2$
by $R_2^2/[s^2(1+R^2)]$. Compact elementary quadrature and explicit bump
derivative bounds enclose the rest. These are mathematical procedures,
not executed quadratures at astronomical selected parameters. They are
unnecessary for the rational sufficient selector and compute no normalized
interacting expectation. Nonlinear normalizers remain positive finite
source-law integrals with the stated lower/enclosure bounds; no numerical
values or sampling access follow. CA1.5's extra input is $\bar K_D$;
CA1.5b's is exactly the certified smoothing/selector contract.

Order is essential. Every finite neural path satisfies
$|F_N|\le\sqrt N/2$ area-a.e.; nonnegative unit convolution gives
$B_{T_s}\le IN/4$. But $c_s\to\infty$: on each fixed momentum ball
$|\hat\rho(sp)|^2\to1$, while expanding-ball integrals of $\hat C_0$
diverge logarithmically. Thus its physical-$C_0$-subtracted coordinate
is at most $IN/4-Ic_s\to-\infty$ at fixed $a,N$. Under the continuous
endpoint full tilt the squared convolution has a finite pathwise limit,
with the same divergence of subtraction. Uniform NC1 comparison with a
moving common subtraction supplies no tightness. The stated successive
limits, reverse nested choices and separately estimated diagonals preserve
the regulator/smoothing order, proving the remaining CA1.2 assertions.

The unused triangle through finite-box raw $Q_0$ adds $M_nL_f\sqrt s$.
Its majorant grows as $2^{\mathrm{constant}\,n^4}$, conflicting with the
selected crude smoothed-volume majorant when both are forced to vanish.
This is failure of these sufficient estimates, without nonconvergence
or an all-method lower bound. The actual direct smoothed chain has no
such arrow, so the charge is absent by law choice rather than bookkeeping
omission. Growing-volume raw transfer stays separately open.

Original adverse records remain: CMP GAUSS/AUTHOR MAJOR startup-access
findings and software-v1 NUM DISAGREEMENTS/MAJOR parser-resource defect;
RL2's MAJOR missing-body issue; and all CV1/NC1 raw findings, severities,
assessments, separate confirmations and remaining source/positive
obligations. This integration does not reread unprovided historical raw
logs or declare their original compliance repaired. Historical author
pre-execution hook failures remain distinct from the present author's
own retained command failures. Frozen acceptance decisions supersede
the historical candidate-status headers for accepted mathematics; they
do not erase failures or establish expert/model-error independence.

CA1 establishes conditional actual normalized joint-law assembly and finite
sufficient scales. It establishes no useful estimator, free exact-real
draw, finite-bit source generation, global-action cost, overlap/mixing,
complete runtime, matched Gaussian comparator, positive selected-cosine
gap, comparative benefit, new priority or other model. Full G1--G9 flags
remain false. The original source review and the present derived-manuscript
review are distinct. This completes CA1.1--CA1.6 and all retained limits.

% CMAP integration | CLAIM-CCA.md:26-64;CC1-STATEMENT-ONLY.md:76-94;CA1-STATEMENT-ONLY.md:284-297;derivation-composite-assembly-20261031.md:215-225;derivation-composite-assembly-20261031.md:345-356
\subsection{The same-target CC1--CA1 integration}
\begin{corollary}[Certified composite selector under the union of premises]\label{cm:union}
Assume all S1--S4 and H1--H7/Appendix S for their identical original target,
Gaussian reference, bump, $c_s$ subtraction and one Borel $Q_0$.
CC1 supplies CA1.5b with
\[
 r(s)=36\sqrt s,\qquad S_r(\tau)=2^{-2\cmLtwo(36/\tau)} .
\]
For every rational $0<\eta\le1$, CA1.5b therefore selects finite positive
$s,a,\beta$ and finite $n,R_{\mathrm{aux}},N$ for both architectures with
$\cmBL(\mathcal L_{\mathrm{neural}}^{\mathfrak b},
\mathcal L_{\mathrm{target}})\le\eta$. At its target allocation,
\[
 j=\cmLtwo(288/\eta),\qquad s=2^{-2j},\qquad h=2j.
\]
The alternative CA1.5 density-certificate theorem and the analytic CA1.4
sequence retain their own distinct premises and bounds.
\end{corollary}
\begin{proof}
S1/H1 identify exactly the same original FS state and physical free-boundary
limit, and S2--S4/H3--H4 identify the same restricted whole-plane Gaussian,
pressure/boundary density, bump, physical $c_s$ and countable local
$Q_0$. For any joint $G\in\mathrm{BL}_1$, coupling both observations on
one target field keeps $X_{\mathrm{obs}}$ identical and gives
\[
 |\E_{\nu_\infty}G(X_{\mathrm{obs}},Q_s)
       -\E_{\nu_\infty}G(X_{\mathrm{obs}},Q_0)|
 \le\E_{\nu_\infty}|Q_s-Q_0|\le36\sqrt s.
\]
CC1.4 terminates with $r(S_r(\tau))\le\tau$. Thus all literal CA1.5b
additional premises hold under the union. Take $\tau=\eta/8$ to get the
$288/\eta$ index, and apply its proved selector without adding a second
target charge or replacing the target density norm. Since $s=2^{-2j}$
and $j\ge0$, $\cmLtwo(1/s)=2j$ exactly. Other stages, own normalizers,
eight allocations and positive-stage/order conditions are those already
proved. No later-tail theorem or stronger $L^2$ certificate is invoked.
\end{proof}

% CMAP closure | CLAIM-CCA.md:66-83;../limits/G6-DECISION.json:70-71;CA1-STATEMENT-ONLY.md:299-324
This update integrates accepted conditional mathematics into the
preserved manuscript. Its finite symbolic choices are not executed giant
field stages. The separate G6 STOP / INCOMPLETE decision remains unchanged; earlier
public editions remain separately archived. Historical internal acceptance of source
mathematics, fresh manuscript review, finite algebra checks, local LaTeX
build and visual document QA are separate records in the companion report.
\endgroup

% END PRIVATE COMPOSITE INTEGRATION 20261106
\bibliographystyle{plainnat}
\bibliography{references}
\endgroup

\end{document}
